\beginnumbering
\eledsection*{Adhyāya 39}
\pstart \ledsidenote{\passref{6.39.1}}\skt{अथातो }\edtext{\skt{ज्वरप्रतिषेधं}}{\Afootnote{A K; \skt{ज्वरोत्प}\ H}}\skt{ }\edtext{\skt{अध्यायं}}{\Afootnote{\emph{om.}\ A; \skt{नामाध्यायं}\ K; \skt{त्तिमध्यायं}\ H}}\skt{ व्याख्यास्यामः ॥}\pend

\pstart \ledsidenote{\passref{6.39.2}}\edtext{\skt{अथ}}{\Afootnote{K; \emph{om.}\ A}}\skt{ }\edtext{\skt{खलु}}{\Afootnote{K; \emph{om.}\ A}}\skt{ }\edtext{\skt{वत्ससुश्रुत}}{\Afootnote{K; \skt{यथोवाच भग\-वा\-न्ध\-न्व\-न्तरिः ॥}\ A}}\pend

\pstart\versepar \ledsidenote{\passref{6.39.3}}\skt{येनामृतमपां मध्यादुद्धृतं पूर्व्वजन्मनि । }\pend
\pstart\versepar \skt{}\edtext{\skt{यतोऽमरत्वं}}{\lemma{\skt{ऽमरत्वं}}\Afootnote{A; \skt{मरत्वं}\ K H}}\skt{ }\edtext{\skt{संप्रा\-प्ता\-स्त्रि\-द\-शा\-स्त्रि\-दि\-वे\-श्वराः}}{\lemma{\skt{त्रिदिवेश्वराः}}\Afootnote{K; \skt{त्रिदिवेश्वरात् ॥}\ A; \skt{त्रिदशेश्वरं ॥}\ H}}\skt{ ॥}\pend

\pstart\versepar \ledsidenote{\passref{6.39.4}}\edtext{\skt{शिष्यास्तं}}{\Afootnote{A; \skt{सिष्यास्तन्}\ K; \skt{शिष्यास्तं}\ H}}\skt{ }\edtext{\skt{देवमासीनं}}{\lemma{\skt{आसीनं}}\Afootnote{A K; \skt{आसीन}\ H}}\skt{ }\edtext{\skt{पप्रच्छुः}}{\Afootnote{A; \skt{प्रपच्छस्}\ K; \skt{…म्पप्रच्छुः}\ H}}\skt{ सुश्रुतादयः । }\pend
\pstart\versepar \skt{व्रणस्योपद्रवाः प्रोक्ता व्रणिनामप्यतः परं ॥}\pend

\pstart\versepar \ledsidenote{\passref{6.39.5}}\edtext{\skt{समासतो}}{\Afootnote{H; \skt{समासा}\ A; \skt{सामासतो}\ K}}\skt{ }\edtext{\edtext{\skt{व्यासतश्च}}{\lemma{\skt{व्यासतश्}}\Afootnote{K H; \skt{द्व्यासतश्}\ A}}}{\lemma{\skt{च}}\Afootnote{K H; \skt{चैव}\ A}}\skt{ }\edtext{\skt{ब्रूहि नो}}{\Afootnote{A H; \skt{व्रूहि नो}\ K}}\skt{ भिषजाम्वर । }\pend
\pstart\versepar \skt{उपद्रवेण जुष्टस्य व्रणः कृच्छ्रेण सिध्यति ॥}\pend

\pstart\versepar \ledsidenote{\passref{6.39.6}}\edtext{\skt{उपद्रवाश्च}}{\Afootnote{H; \skt{उपद्रवास्तु}\ A; \skt{उपद्रवस्तु}\ K}}\skt{ व्रणिनः }\edtext{\skt{कृच्छ्रसाध्य्तमो}}{\Afootnote{\skt{कृच्छ्रसाध्याः प्रकी}\ A; \skt{कृच्छ्रसाध्यतमो}\ K H}}\skt{ }\edtext{\skt{मतः}}{\Afootnote{K H; \skt{र्तिताः ।}\ A}}\skt{ । }\pend
\pstart\versepar \skt{}\edtext{\skt{प्रक्षी\-ण\-ब\-ल\-मां\-सस्य}}{\Afootnote{A; \skt{प्रक्षी\-ण\-व\-ल\-मां\-सस्य}\ K; \skt{प्रक्षी\-ण\-व\-ल\-मां\-सस्य्\textbackslash\{\}}\ H}}\skt{ शेष\-धा\-तु\-प\-रि\-क्ष\-यात् ॥}\pend

\pstart \ledsidenote{\passref{6.39.7ab}}\edtext{\skt{तस्मा\-दु\-प\-द्र\-वा\-न्कृ\-त्स्ना\-न्व\-दस्व}}{\Afootnote{\skt{तस्मा\-दु\-प\-द्र\-वा\-न्कृत्स्ना}\ A; \skt{तस्मादुपद्रवां कृत्स्नान्वदस्व}\ K}}\skt{ }\edtext{\edtext{\skt{वदताम्वर}}{\lemma{\skt{वदताम्}}\Afootnote{K; \emph{om.}\ A}}}{\lemma{\skt{वर}}\Afootnote{\skt{न्ब्रूहि नः सचिकित्सितान् ।}\ A; \skt{वरः ॥}\ K}}\skt{ ॥}\pend

\pstart\versepar \ledsidenote{\passref{6.39.8}}\edtext{\edtext{\skt{तेषान्तद्वचनं}}{\lemma{\skt{तेषान्}}\Afootnote{A H; \emph{om.}\ K}}}{\lemma{\skt{तद्वचनं}}\Afootnote{A H; \emph{om.}\ K}}\skt{ }\edtext{\skt{श्रुत्वा}}{\Afootnote{A H; \emph{om.}\ K}}\skt{ }\edtext{\edtext{\skt{प्राब्र\-वी\-द्भि\-ष\-जा\-म्वरः}}{\lemma{\skt{प्राब्र\-वी\-द्भि\-ष\-जाम्}}\Afootnote{A; \emph{om.}\ K; \skt{प्राव्र\-वी\-द्भि\-ष\-जाम्}\ H}}}{\lemma{\skt{वरः}}\Afootnote{A H; \emph{om.}\ K}}\skt{ ॥ }\pend
\pstart\versepar \skt{ज्वरमादौ प्रवक्ष्यामि स रोगा\-नी\-क\-रा\-ट्स्मृतः ।}\pend

\pstart\versepar \ledsidenote{\passref{6.39.9}}\skt{रुद्र\-को\-पा\-ग्नि\-सं\-भूतः सर्व्व\-भू\-त\-प्र\-ता\-पनः ॥ }\pend
\pstart\versepar \skt{}\edtext{\skt{तैस्तै\-र्न्ना\-म\-भि\-र\-न्येषां}}{\lemma{\skt{अन्येषां}}\Afootnote{A K; \skt{अनेकेषां}\ H}}\skt{ सत्त्वा\-ना\-म्प\-रि\-की\-र्त्त्यते ।}\edtext{}{\lemma{\skt{सर्व्व\-भू\-त\-प्र\-ता\-पनः ॥ तैस्तैर्}}\Afootnote{A; \skt{सर्व\-भू\-त\-प्र\-ता\-पनः । तैस्तै}\ K; \skt{सर्व्वभूतः प्रतापनः ॥ तैस्तैर्}\ H}}\pend

\pstart\versepar \ledsidenote{\passref{6.39.10}}\skt{जन्मादौ निधने चैव }\edtext{\skt{भवतीह न}}{\Afootnote{K H; \skt{प्रायो विशति}\ A}}\skt{ }\edtext{\skt{संशयः}}{\Afootnote{K H; \skt{देहिनम् ।}\ A}}\skt{ ॥ }\pend
\pstart\versepar \skt{अतः सर्वविकाराणां }\edtext{\skt{ज्वरो}}{\Afootnote{K H; \skt{अयं}\ A}}\skt{ राजा प्रकीर्त्तितः ।}\pend

\pstart \ledsidenote{\passref{6.39.11}}\skt{ऋते देवमनुष्येभ्यो नान्यो }\edtext{\skt{विषहते तु तं}}{\Afootnote{K H; \skt{विषहते तु तम् । कर्मणा लभते यस्माद्देवत्वं मानुसादपि ॥}\ A}}\skt{ ॥}\pend

\pstart \ledsidenote{\passref{6.39.13}}\edtext{\skt{स्वेदावरोधः}}{\Afootnote{H; \skt{शेषाः सर्वे विपद्यन्ते तैर्यग्योना ज्वरार्दिताः । स्वेदावरोधः}\ A; \skt{स्वेदोपरोधस्}\ K}}\skt{ सन्तापः }\edtext{\edtext{\skt{सर्व्वा\-ङ्ग\-ग्र\-ह\-ण\-म्तथा}}{\lemma{\skt{सर्व्वा\-ङ्ग\-ग्र\-ह\-णम्}}\Afootnote{A H; \skt{सर्व\-ङ्ग\-ग्र\-ह\-णन्}\ K}}}{\lemma{\skt{तथा}}\Afootnote{A K; \skt{भृश ।}\ H}}\skt{ ।}\pend

\pstart\versepar \ledsidenote{\passref{6.39.14}}\edtext{\skt{युगपद्यत्र}}{\Afootnote{K; \skt{विकारा युगप}\ A; \skt{युगपद्यत्र}\ H}}\skt{ }\edtext{\skt{रोगे तु}}{\Afootnote{K H; \emph{om.}\ A}}\skt{ }\edtext{\skt{ज्वर}}{\Afootnote{K H; \skt{द्य}\ A}}\skt{ }\edtext{\edtext{\skt{इत्युपदिश्यते}}{\lemma{\skt{इत्य्}}\Afootnote{K H; \skt{स्मिन्ज्व}\ A}}}{\lemma{\skt{उपदिश्यते}}\Afootnote{K H; \skt{रः स परिकीर्तितः ।}\ A}}\skt{ ॥ }\pend
\pstart\versepar \skt{दोषैः पृथक्समस्तैश्च }\edtext{\edtext{\skt{द्वन्द्वै\-रा\-ग\-न्तु\-रेव}}{\lemma{\skt{द्वन्द्वैर्}}\Afootnote{A; \skt{द्वंद्वर्}\ K; \skt{द्वन्द्वै}\ H}}}{\lemma{\skt{आगन्तुर्}}\Afootnote{A K; \skt{रागन्तुर्}\ H}}\skt{ च ।}\pend

\pstart\versepar \ledsidenote{\passref{6.39.15}}\edtext{\skt{अने\-क\-का\-र\-णो\-त्पन्नः}}{\Afootnote{A K; \skt{अने\-क\-का\-र\-णो…\-त्पन्नः}\ H}}\skt{ }\edtext{\skt{स्मृतस्त्वष्टविधो}}{\lemma{\skt{त्वष्टविधो}}\Afootnote{A H; \skt{त्वष्टविधौ}\ K}}\skt{ ज्वरः ॥ }\pend
\pstart\versepar \skt{दोषाः }\edtext{\skt{प्रकुपिता}}{\Afootnote{K H; \skt{प्रकुपिताः}\ A}}\skt{ }\edtext{\edtext{\skt{गाढ\-म\-प\-थ्या\-हा\-र\-से\-विनः}}{\lemma{\skt{गाढम्}}\Afootnote{K H; \emph{om.}\ A}}}{\lemma{\skt{अप\-थ्या\-हा\-र\-से\-विनः}}\Afootnote{K H; \skt{स्वेषु कालेषु स्वैः प्रकोपणैः ॥}\ A}}\skt{ ।}\pend

\pstart \ledsidenote{\passref{6.39.16}}\skt{व्याप्य देहमशेषेण ज्वर\-मा\-पा\-द\-यन्ति }\edtext{\skt{ह}}{\Afootnote{K H; \skt{हि । दुष्टाः स्वहेतुभिर्दोषाः प्राप्या\-मा\-श\-य\-मू\-ष्मणा ॥}\ A}}\skt{ ॥}\pend

\pstart \ledsidenote{\passref{6.39.19cd}}\edtext{\edtext{\skt{अधि\-मि\-थ्या\-प्र\-यु\-क्तैश्च}}{\lemma{\skt{अधि\-मि\-थ्या\-प्र\-यु\-क्तैश्}}\Afootnote{\skt{मिथ्यातियुक्तैर्}\ A; \skt{अ पि मिथ्याप्रयुक्तैश्}\ K H}}}{\lemma{\skt{च}}\Afootnote{K H; \skt{अपि च}\ A}}\skt{ स्नेहाद्यैः कर्म्मभिर्न्नृणां ।}\pend

\pstart\versepar \ledsidenote{\passref{6.39.20}}\edtext{\skt{विवि\-धा\-द\-भि\-घा\-ताच्च}}{\Afootnote{A K; \skt{विवि\-धा\-वि\-धा\-द\-वि\-घा\-ताच्च}\ H\textsuperscript{ac}; \skt{विवि\-धा\-द\-वि\-घा\-ताच्च}\ H\textsuperscript{pc}}}\skt{ }\edtext{\skt{रोगो\-त्था\-ना\-त्प्र\-पा\-कतः}}{\Afootnote{A K H\textsuperscript{ac}; \skt{रोगो\-त्था\-च्छा\-ना\-त्प्र\-पा\-कतः ॥}\ H\textsuperscript{pc}}}\skt{ ॥ }\pend
\pstart\versepar \skt{श्रमा\-त्क्ष\-या\-द\-जी\-र्ण्णाच्च }\edtext{\skt{विषा\-त्सा\-त्म्य\-र्तु\-प\-र्य\-यात्}}{\Afootnote{A; \skt{विषा\-त्सा\-म\-र्त्तु\-प्र\-य\-यात् ।}\ K; \skt{विषात्साम्य तु पर्ययात् ।}\ H}}\skt{ ।}\pend

\pstart\versepar \ledsidenote{\passref{6.39.21}}\skt{ओष\-धी\-पु\-ष्प\-ग\-न्धाच्च }\edtext{\skt{कोका\-न्न\-क्ष\-त्र\-पी\-डया}}{\Afootnote{\skt{शोका\-न्न\-क्ष\-त्र\-पी\-डया ।}\ A; \skt{काया\-न्न\-क्ष\-त्र\-पी\-डया ।.}\ K; \skt{कोपा\-न्न\-क्ष\-त्र\-पी\-डया ॥}\ H}}\skt{ ॥ }\pend
\pstart\versepar \skt{}\edtext{\skt{अभि\-शा\-पा\-भि\-चा\-राभ्यां}}{\Afootnote{\skt{अभि\-चा\-रा\-भि\-शा\-पाभ्यां}\ A; \skt{अभि\-षा\-पा\-भि\-चा\-राभ्यां}\ K; \skt{अभि\-शा\-पा\-भि\-चा\-रा\-भ्या…म्}\ H}}\skt{ मनो\-भू\-ता\-भि\-श\-ङ्कया ।}\pend

\pstart\versepar \ledsidenote{\passref{6.39.22}}\edtext{\edtext{\skt{स्त्रीणा\-म\-व\-प्र\-जा\-तानां}}{\lemma{\skt{स्त्रीणाम्}}\Afootnote{A K; \skt{स्त्रीणा}\ H}}}{\lemma{\skt{अवप्रजातानां}}\Afootnote{K; \skt{अपप्रजातानां}\ A; \skt{मव\-प्र\-जा\-ता\-नाम्}\ H}}\skt{ प्रजा\-ता\-ना\-न्त\-था\-हितैः ॥ }\pend
\pstart\versepar \skt{}\edtext{\skt{स्तन्यावतणे}}{\Afootnote{\skt{स्तन्यावतरणे चैव}\ A; \skt{स्तन्यावतरणै चैव}\ K; \skt{स्तन्यावतणे चैव}\ H}}\skt{ }\edtext{\skt{ज्वरो}}{\Afootnote{A H; \skt{ज्वरौ}\ K}}\skt{ दोषैः }\edtext{\skt{प्रकुप्यति}}{\Afootnote{K H; \skt{प्रवर्तते ॥}\ A}}\skt{ ।}\pend

\pstart\versepar \ledsidenote{\passref{6.39.23}}\edtext{\skt{तैर्व्वे\-ग\-व\-द्भि\-र्ब्ब\-हुधा}}{\lemma{\skt{ब्बहुधा}}\Afootnote{A H; \skt{वहुधा}\ K}}\skt{ }\edtext{\skt{समुद्भ्रान्तै}}{\Afootnote{\skt{समुद्भ्रान्तैर्}\ A H; \skt{समुभ्राद्भन्तैर्}\ K\textsuperscript{ac}; \skt{समुद्भन्तैर्}\ K\textsuperscript{pc}}}\skt{ }\edtext{\skt{र्व्विमार्ग्गगैः}}{\Afootnote{\skt{विमार्गगैः}\ A K H}}\skt{ ॥ }\pend
\pstart\versepar \skt{}\edtext{\edtext{\skt{विक्षि\-प्य\-मा\-णोऽ\-न्त\-र\-ग्नि\-र्भ्भ\-व\-त्याशु}}{\lemma{\skt{विक्षिप्यमाणो}}\Afootnote{A H; \skt{विक्षिप्यमाणोन्}\ K}}}{\lemma{\skt{ऽन्त\-र\-ग्नि\-र्भ्भ\-व\-त्याशु}}\Afootnote{\skt{ऽन्त\-र\-ग्नि\-र्भ\-व\-त्याशु}\ A; \skt{तर\-ग्नि\-र्भ\-व\-त्या\-शुर्}\ K; \skt{ऽन्तरग्नि र्भ्भवत्याशु}\ H}}\skt{ }\edtext{\edtext{\skt{बहिश्चरः}}{\lemma{\skt{बहिश्}}\Afootnote{A H; \skt{वहि}\ K}}}{\lemma{\skt{चरः}}\Afootnote{A H; \skt{श्चरः ।}\ K}}\skt{ ।}\pend

\pstart\versepar \ledsidenote{\passref{6.39.24}}\skt{रुणद्धि }\edtext{\edtext{\skt{चाप्यपान्धातुं}}{\lemma{\skt{अपान्}}\Afootnote{A K; \skt{अपा}\ H}}}{\lemma{\skt{धातुं}}\Afootnote{A K; \skt{न्धातुं}\ H}}\skt{ यस्मा\-त्त\-स्मा\-ज्ज्व\-रा\-तुरः ॥ }\pend
\pstart\versepar \skt{}\edtext{\edtext{\skt{भव\-त्य\-त्यु\-ष्ण\-गा\-त्रश्च}}{\lemma{\skt{भवत्य्}}\Afootnote{A K\textsuperscript{ac} H; \skt{भव}\ K\textsuperscript{pc}}}}{\lemma{\skt{अत्युष्णगात्रश्च}}\Afootnote{A; \skt{उष्णगात्रश्च}\ K\textsuperscript{ac}; \skt{त्यत्युष्णगात्रश्च}\ K\textsuperscript{pc}; \skt{अत्युष्णगात्रश्च}\ H}}\skt{ }\edtext{\skt{स्विद्यतेन च}}{\Afootnote{K H; \skt{ज्वरितस्तेन चोच्यते ॥}\ A}}\skt{ }\edtext{\skt{सर्व्वशः}}{\Afootnote{K H; \emph{om.}\ A}}\skt{ ॥}\pend

\pstart\versepar \ledsidenote{\passref{6.39.25}}\edtext{\edtext{\skt{श्रमो\-र\-ति\-र्वि\-व\-र्ण्णत्वं}}{\lemma{\skt{श्रमोर्}}\Afootnote{\skt{श्रमो}\ A; \skt{श्रमो}\ K H}}}{\lemma{\skt{अतिर्विवर्ण्णत्वं}}\Afootnote{\skt{ऽर\-ति\-र्वि\-व\-र्णत्वं}\ A; \skt{रत्रिविवर्ण्णत्व}\ K; \skt{रतिर्विवर्ण्णत्वं}\ H}}\skt{ वैरस्यं नयनप्लवः । }\pend
\pstart\versepar \skt{}\edtext{\skt{इच्छाद्वेषौ}}{\Afootnote{A K; \skt{इच्छाद्वेषो}\ H}}\skt{ }\edtext{\edtext{\skt{मुहुश्चापि}}{\lemma{\skt{मुहुश्}}\Afootnote{A H; \skt{वहुश्}\ K}}}{\lemma{\skt{चापि}}\Afootnote{A K; \skt{चापि…}\ H}}\skt{ शीत\-वा\-ता\-त\-पा\-दिषु ॥}\pend

\pstart\versepar \ledsidenote{\passref{6.39.26}}\edtext{\skt{जृम्भाङ्गमर्द्दो}}{\Afootnote{A H; \skt{जृंभोङ्गमर्दो}\ K}}\skt{ गुरुता रोम\-ह\-र्षोऽ\-रु\-चि\-स्तमः । }\pend
\pstart\versepar \skt{अप्रहर्षश्च }\edtext{\skt{शीत च}}{\Afootnote{H; \skt{शीतं च}\ A; \skt{शीतश्च}\ K}}\skt{ }\edtext{\edtext{\skt{भवत्युत्पत्स्यति}}{\lemma{\skt{भवत्य्}}\Afootnote{A K; \skt{भव}\ H}}}{\lemma{\skt{उत्पत्स्यति}}\Afootnote{A K; \skt{त्युत्पत्स्यति}\ H}}\skt{ ज्वरे ॥}\pend

\pstart\versepar \ledsidenote{\passref{6.39.27}}\skt{सामान्यतो }\edtext{\skt{विशेषात्तु}}{\Afootnote{A; \skt{विशेषस्तु}\ K; \skt{विशेषस्तु}\ H}}\skt{ }\edtext{\skt{जृम्भा\-त्य\-र्थ\-न्स\-मी\-र\-णात्}}{\lemma{\skt{जृम्भात्य्}}\Afootnote{H; \skt{जृम्भा}\ A; \skt{जृंभोत्य्}\ K}}\skt{ । }\pend
\pstart\versepar \skt{}\edtext{\skt{पित्ता\-न्न\-य\-न\-र्द्दाहः}}{\lemma{\skt{नयनर्द्दाहः}}\Afootnote{H; \skt{नयनयोर्दाहः}\ A; \skt{नयनयो दाहः ।}\ K}}\skt{ }\edtext{\edtext{\skt{कफा\-द\-न्ना\-रु\-चि\-स्तथा}}{\lemma{\skt{कफा\-द\-न्ना\-रु\-चिस्}}\Afootnote{K; \skt{कफा\-न्न\-न्ना\-भि\-नन्}\ A; \skt{कफा\-द\-न्ना\-रु\-चिस्}\ H}}}{\lemma{\skt{तथा}}\Afootnote{K H; \skt{दनम् ॥}\ A}}\skt{ ॥}\edtext{}{\lemma{\skt{अर्थ\-न्स\-मी\-र\-णात् । पित्तान्}}\Afootnote{H; \skt{त्यर्थं समीरणात् । पित्तान्}\ A; \skt{अर्थं समीरणात् । पित्तान्}\ K}}\pend

\pstart\versepar \ledsidenote{\passref{6.39.28}}\edtext{\skt{सर्व्व\-लि\-ङ्ग\-स\-म\-वायः}}{\Afootnote{A H; \skt{सर्व\-लि\-ङ्ग\-स\-मा\-वा\-यस्}\ K}}\skt{ सर्व्व\-दो\-ष\-प्र\-को\-पजे । }\pend
\pstart\versepar \skt{}\edtext{\edtext{\skt{रूपै\-र\-न्य\-त\-रा\-भ्यान्तु}}{\lemma{\skt{रूपैर्}}\Afootnote{H; \skt{द्वयोर्}\ A; \skt{रुपैर्}\ K}}}{\lemma{\skt{अन्यतराभ्यान्तु}}\Afootnote{H; \skt{द्वयोस्तु रूपेण}\ A; \skt{अन्यराभ्यान्तु}\ K}}\skt{ }\edtext{\skt{संसृष्टे}}{\Afootnote{H; \skt{संसृष्टं}\ A; \skt{संसृष्ट}\ K}}\skt{ द्वन्द्वजं विदुः ॥}\pend

\pstart\versepar \ledsidenote{\passref{6.39.29}}\edtext{\skt{वेपथुर्व्विषमो}}{\Afootnote{\skt{वेपथुर्विषमो}\ A K H}}\skt{ वेगः }\edtext{\skt{कण्ठौ\-ष्ठ\-प\-रि\-शो\-षणं}}{\Afootnote{A; \skt{कण्ठो\-ष्ठ\-प\-रि\-शो\-षणः ।}\ K; \skt{कण्ठौ\-ष्ठ\-प\-रि\-शो\-षणᳵ ।}\ H}}\skt{ । }\pend
\pstart\versepar \skt{निद्रानाशः }\edtext{\skt{क्षयस्तंभो}}{\Afootnote{K; \skt{क्षुतः स्तम्भो}\ A; \skt{क्षयः स्तम्भो}\ H}}\skt{ गात्राणां रौक्ष्यमेव च ॥}\pend

\pstart\versepar \ledsidenote{\passref{6.39.30}}\edtext{\skt{शिरोरुग्गात्र रुग्वक्त्रं}}{\Afootnote{H; \skt{शिरो\-हृ\-द्गा\-त्र\-रु\-ग्वक्त्र}\ A; \skt{शिरो\-रु\-ग्गा\-त्र\-रु\-ग्वक्र}\ K}}\skt{ वैरस्यं }\edtext{\skt{गाढविट्कता}}{\Afootnote{K H; \skt{बद्धविट्कता ।}\ A}}\skt{ । }\pend
\pstart\versepar \skt{}\edtext{\skt{शूलाध्मानौ}}{\Afootnote{H; \skt{जृम्भाध्मान्ं}\ A; \skt{शूलाध्मान}\ K}}\skt{ }\edtext{\skt{जृम्भणञ्च}}{\Afootnote{H; \skt{अ तथा शूलं}\ A; \skt{जृंभणञ्च}\ K}}\skt{ भवत्यनिलजे ज्वरे ॥}\pend

\pstart\versepar \ledsidenote{\passref{6.39.31}}\edtext{\skt{वेग\-स्ती\-क्ष्णो\-म\-ति\-सा\-रश्च}}{\Afootnote{\skt{वेग\-स्ती\-क्ष्णोऽ\-ति\-सा\-रश्च}\ A; \skt{वेगतीक्षौ}\ K; \skt{वेगस्तीक्ष्णोम}\ H\textsuperscript{ac}; \skt{वेग\-स्ती\-क्ष्णो\-म\-ति\-सा\-रश्च}\ H\textsuperscript{pc}}}\skt{ }\edtext{\skt{महादाहो}}{\Afootnote{H; \emph{om.}\ A; \skt{महादाहा}\ K}}\skt{ }\edtext{\skt{निद्राल्पश्च}}{\Afootnote{K H; \skt{निद्राल्पत्वं}\ A}}\skt{ तथा }\edtext{\skt{वमिः}}{\Afootnote{A H; \skt{वमी ॥}\ K}}\skt{ । }\pend
\pstart\versepar \skt{}\edtext{\skt{कण्ठौष्ठम्पाकः}}{\Afootnote{\skt{कण्ठौ\-ष्ठ\-मु\-ख\-ना\-सानां पाकः}\ A H; \skt{कण्ठो\-ष्ठ\-मु\-ख\-ना\-सानां स्वेदः}\ K}}\skt{ }\edtext{\skt{स्वेदश्च}}{\Afootnote{A H; \skt{पाकश्च}\ K}}\skt{ जायते ॥}\pend

\pstart\versepar \ledsidenote{\passref{6.39.32}}\skt{प्रलापो }\edtext{\skt{वक्त्र}}{\Afootnote{K H; \emph{om.}\ A}}\skt{ }\edtext{\skt{कटुका}}{\Afootnote{H; \skt{कटुता}\ A K}}\skt{ }\edtext{\skt{मूर्च्छा}}{\Afootnote{K H; \skt{वक्त्रे मूर्च्छा}\ A}}\skt{ दाहो }\edtext{\skt{मदस्तृषा}}{\Afootnote{A; \skt{मदस्तृषा}\ K H}}\skt{ । }\pend
\pstart\versepar \skt{पीत\-वि\-ण्मू\-त्र\-ने\-त्र\-त्व\-म्पै\-त्तिके भ्रम एव च ॥}\pend

\pstart\versepar \ledsidenote{\passref{6.39.32a}}\edtext{\skt{स्तैमित्यं}}{\Afootnote{H; \skt{स्तेमित्यं}\ K}}\skt{ स्तिमितो वेग आलस्यम्म धुरास्यता । }\pend
\pstart\versepar \skt{शुक्ल\-मू\-त्र\-पु\-री\-षत्वं }\edtext{\skt{स्तम्भ\-स्तृ\-प्ति\-र\-थापि}}{\Afootnote{\skt{स्तंभ\-स्तृ\-प्ति\-र\-थापि}\ K; \skt{स्तम्भ\-स्तृ\-प्ति\-र\-थापि}\ H}}\skt{ च ॥}\pend

\pstart \ledsidenote{\passref{6.39.33}}\skt{गौरवं }\edtext{\skt{शीतमुत्क्लेदो}}{\Afootnote{\skt{शीतमुत्क्लेशो}\ A; \skt{शीतमुत्क्लेदो}\ K; \skt{शीतमुत्क्लेदो…}\ H}}\skt{ }\edtext{\edtext{\skt{रोम\-ह\-र्षोऽ\-ति\-नि\-द्रता}}{\lemma{\skt{रोमहर्षो}}\Afootnote{A H; \skt{रोमहर्षो}\ K}}}{\lemma{\skt{ऽतिनिद्रता}}\Afootnote{H; \skt{ऽतिनिद्रता । स्रोतोरोधो रुगल्पत्वं प्रसेको मधुरास्यता ॥}\ A; \skt{तिनिद्रता ॥}\ K}}\skt{ ।}\pend

\pstart \ledsidenote{\passref{6.39.34}}\edtext{\edtext{\skt{प्रतिश्यायोऽरुचिः}}{\lemma{\skt{प्रतिश्यायो}}\Afootnote{K H; \skt{नात्युष्णगात्रता छर्दि\-र\-ङ्ग\-सा\-दोऽ\-वि\-पा\-कता । प्रतिश्यायो}\ A}}}{\lemma{\skt{ऽरुचिः}}\Afootnote{A H; \skt{रुचिः}\ K}}\skt{ कासः }\edtext{\skt{कफजेऽक्ष्णोश्च}}{\lemma{\skt{ऽक्ष्णोश्च}}\Afootnote{A H; \skt{क्ष्णोश्च}\ K}}\skt{ शुक्लता ।}\pend

\pstart \ledsidenote{\passref{6.39.38}}\edtext{\skt{सर्व्वजः}}{\Afootnote{K H; \skt{स्वेद\-मू\-त्र\-पु\-री\-षा\-णा\-म\-ल्पशः सुचिरात्स्रुतिः । सर्वजे}\ A}}\skt{ }\edtext{\skt{सर्व्वलिङ्ग्स्तु}}{\Afootnote{H; \skt{सर्वलिङ्गानि}\ A; \skt{सर्वलिङ्गस्तु}\ K}}\skt{ विशेषञ्चात्र मे शृणु ।}\pend

\pstart \ledsidenote{\passref{6.39.39}}\edtext{\edtext{\skt{नात्यु\-ष्ण\-शी\-तोऽ\-ल्प\-संज्ञा}}{\lemma{\skt{नात्युष्णशीतो}}\Afootnote{A; \skt{नात्युष्णशीतो}\ K; \skt{नात्युष्णशीतो}\ H}}}{\lemma{\skt{ऽल्पसंज्ञा}}\Afootnote{H; \skt{ऽल्पसंज्ञो}\ A; \skt{ल्पसंज्ञो}\ K}}\skt{ }\edtext{\skt{भ्रान्तप्रेक्षी}}{\Afootnote{A H; \skt{भ्रान्तपेक्षी}\ K}}\skt{ }\edtext{\skt{हतस्वरः}}{\Afootnote{K H; \skt{हतस्वरः । खरजिह्वः शुष्ककण्ठः स्वेद\-वि\-ण्मू\-त्र\-व\-र्जितः ॥}\ A}}\skt{ ॥}\pend

\pstart \ledsidenote{\passref{6.39.40}}\edtext{\skt{श्वसन्निपतितः}}{\Afootnote{H; \skt{सास्रो निर्भुग्नहृदयो भक्तद्वेषी हतप्रभः । श्वसन्निपतितः}\ A; \skt{श्वसन्निपतितिः}\ K}}\skt{ शेते प्रला\-पो\-प\-द्र\-वा\-युतः ।}\pend

\pstart\versepar \ledsidenote{\passref{6.39.41}}\edtext{\edtext{\edtext{\edtext{\edtext{\skt{तम\-भि\-न्या\-स\-मि\-त्या\-हु\-र्ह\-तौ\-ज\-स\-म\-था\-परे}}{\lemma{\skt{अभिन्यासम्}}\Afootnote{A K; \skt{अभिन्यास}\ H}}}{\lemma{\skt{इत्य्}}\Afootnote{A K; \skt{मित्य्}\ H}}}{\lemma{\skt{आहुर्}}\Afootnote{A K; \skt{आहु…र्}\ H}}}{\lemma{\skt{हतौजसम्}}\Afootnote{A; \skt{हतोजसम्}\ K; \skt{हतौजस}\ H}}}{\lemma{\skt{अथापरे}}\Afootnote{A K; \skt{मथापरे ॥}\ H}}\skt{ ॥ }\pend
\pstart\versepar \skt{}\edtext{\skt{सन्निपातज्वरः}}{\Afootnote{K H; \skt{सन्निपातज्वरं}\ A}}\skt{ }\edtext{\skt{कृच्छ्र\-म\-सा\-ध्य\-म\-परे}}{\lemma{\skt{कृच्छ्रम्}}\Afootnote{A K; \skt{कृच्छ्र…म्}\ H}}\skt{ विदुः ।}\pend

\pstart \ledsidenote{\passref{6.39.45}}\edtext{\skt{सप्तमे}}{\Afootnote{K H; \skt{ओजोनिरोधजं तस्य जानीयात्कुशलो भिसक् । सप्तमे}\ A}}\skt{ }\edtext{\skt{दिवसे}}{\Afootnote{K H; \skt{विदसे}\ A}}\skt{ प्राप्ते }\edtext{\skt{दशमे}}{\Afootnote{A K\textsuperscript{pc} H; \skt{र्सदशमे}\ K\textsuperscript{ac}}}\skt{ द्वादशे }\edtext{\skt{पि}}{\Afootnote{K H; \skt{ऽपि}\ A}}\skt{ वा ॥}\pend

\pstart\versepar \ledsidenote{\passref{6.39.46}}\edtext{\skt{पुनर्घोरतरो}}{\lemma{\skt{घोरतरो}}\Afootnote{K H; \skt{धोरतरो}\ A}}\skt{ भूत्वा प्रशमं }\edtext{\skt{यान्ति}}{\Afootnote{K H\textsuperscript{ac}; \skt{याति}\ A H\textsuperscript{pc}}}\skt{ हन्ति वा ॥ }\pend
\pstart\versepar \skt{}\edtext{\skt{द्विदो\-षो\-च्छ्र\-य\-लि\-ङ्गस्तु}}{\Afootnote{H; \skt{द्विदो\-षो\-च्छ्रा\-य\-लि\-ङ्गास्तु}\ A; \skt{द्विदो\-षो\-च्छ्र\-य\-लि\-ङ्गास्तु}\ K}}\skt{ }\edtext{\edtext{\skt{द्वन्द्वजास्त्रिविधः}}{\lemma{\skt{द्वन्द्वजास्}}\Afootnote{A; \skt{द्वंद्वजस्}\ K; \skt{…द्वन्द्वजास्}\ H}}}{\lemma{\skt{त्रिविधः}}\Afootnote{K H; \skt{त्रिविधाः}\ A}}\skt{ }\edtext{\skt{स्मृतः}}{\Afootnote{K H; \skt{स्मृताः ॥}\ A}}\skt{ ।}\pend

\pstart\versepar \ledsidenote{\passref{6.39.47}}\skt{तृष्णा मूर्च्छा भ्रमो दाहः स्वप्ननाशः }\edtext{\skt{शि रोरुजा}}{\Afootnote{A H; \skt{शिरोरुजः ।}\ K}}\skt{ ॥ }\pend
\pstart\versepar \skt{}\edtext{\skt{कण्ठास्यशोषो}}{\Afootnote{A H; \skt{कण्ठस्य शोषो}\ K}}\skt{ }\edtext{\edtext{\skt{वम\-थु\-र्ल्लो\-म\-हर्षो}}{\lemma{\skt{वमथुर्}}\Afootnote{K H; \skt{वमथू}\ A}}}{\lemma{\skt{ल्लोमहर्षो}}\Afootnote{\skt{रोमहर्षो}\ A; \skt{लोमह}\ K H}}\skt{ }\edtext{\edtext{\skt{रुचिस्तमः}}{\lemma{\skt{रुचिस्}}\Afootnote{\skt{ऽरुचिस्}\ A; \skt{र्षोरुचिस्}\ K; \skt{र्षोरुचि}\ H}}}{\lemma{\skt{तमः}}\Afootnote{K; \skt{तथा ॥}\ A; \skt{स्तमः ।}\ H}}\skt{ ।}\pend

\pstart\versepar \ledsidenote{\passref{6.39.48}}\edtext{\skt{पर्व्वभेदश्च}}{\Afootnote{K H; \skt{पर्व्भेदश्च}\ A}}\skt{ }\edtext{\skt{जृम्भा}}{\Afootnote{A H; \skt{जृंभा}\ K}}\skt{ च वात\-पि\-त्त\-ज्व\-रा\-कृतिः ॥ }\pend
\pstart\versepar \skt{}\edtext{\edtext{\skt{स्तैमित्यम्पर्व्वणां}}{\lemma{\skt{स्तैमित्यम्}}\Afootnote{A K; \skt{स्तै…मित्यम्}\ H}}}{\lemma{\skt{पर्व्वणां}}\Afootnote{H; \skt{पर्वणं}\ A; \skt{सर्वाणां}\ K}}\skt{ भेदो निद्रा }\edtext{\skt{गौरव}}{\Afootnote{K H; \skt{गौरवम्}\ A}}\skt{ एव च ।}\pend

\pstart\versepar \ledsidenote{\passref{6.39.49}}\skt{शिरोग्रहः प्रतिश्यायः }\edtext{\skt{कासः}}{\Afootnote{A H; \emph{om.}\ K}}\skt{ }\edtext{\skt{स्वेदः}}{\Afootnote{H; \skt{स्वेद}\ A; \emph{om.}\ K}}\skt{ }\edtext{\skt{प्रवर्त्तते}}{\Afootnote{H; \skt{प्रवर्तनम् ।}\ A; \emph{om.}\ K}}\skt{ ॥ }\pend
\pstart\versepar \skt{}\edtext{\skt{तन्द्रा}}{\Afootnote{H; \skt{सन्ता}\ A; \emph{om.}\ K}}\skt{ }\edtext{\skt{मोहः}}{\Afootnote{H; \skt{पो म}\ A; \emph{om.}\ K}}\skt{ }\edtext{\skt{प्रलापश्च}}{\Afootnote{H; \skt{ध्यवेगश्च}\ A; \emph{om.}\ K}}\skt{ }\edtext{\skt{श्लेष्म\-वा\-त\-ज्व\-रा\-कृतिः}}{\Afootnote{K; \skt{वात\-श्ले\-ष्म\-ज्व\-रा\-कृतिः ॥}\ A; \skt{…श्लेष्म वातज्वराकृतिः ॥}\ H}}\skt{ ॥}\pend

\pstart\versepar \ledsidenote{\passref{6.39.50}}\edtext{\skt{लिप्ततिक्तास्यता}}{\Afootnote{H; \skt{लिप्तातिक्तास्यता}\ A; \skt{लिप्त\-ति\-क्ता\-स्य\-तास्}\ K}}\skt{ तन्द्रा मोहः }\edtext{\edtext{\skt{कासोऽरुचिस्तृषा}}{\lemma{\skt{ऽरुचिस्}}\Afootnote{\skt{ऽरुचि}\ A H; \skt{रुचिस्}\ K}}}{\lemma{\skt{तृषा}}\Afootnote{K; \skt{स्तृषा ।}\ A H}}\skt{ । }\pend
\pstart\versepar \skt{}\edtext{\skt{मुहुर्द्दाहो}}{\Afootnote{A H; \skt{मुहुर्दाहो}\ K}}\skt{ मुहुः }\edtext{\skt{शीतं}}{\Afootnote{A H; \skt{शीत}\ K}}\skt{ श्लेष्म\-पि\-त्त\-ज्व\-रा\-कृतिः ॥}\pend

\pstart \ledsidenote{\passref{6.39.57}}\edtext{\skt{वाता\-धि\-क\-त्वा\-त्प्र\-व\-दन्ति}}{\Afootnote{A H; \skt{वाता\-धि\-का\-त्वा\-त्प्र\-व\-दन्ति}\ K}}\skt{ }\edtext{\edtext{\edtext{\skt{तज्ज्ञा\-स्तृ\-ती\-य\-क\-म्चापि}}{\lemma{\skt{तज्ज्ञास्}}\Afootnote{K; \skt{तज्ज्ञा}\ A H}}}{\lemma{\skt{तृतीयकम्}}\Afootnote{\skt{स्तृतीयकं}\ A H; \skt{तृतीयकश्}\ K}}}{\lemma{\skt{चापि}}\Afootnote{A K; \skt{वापि}\ H}}\skt{ }\edtext{\skt{चतुर्थकम्वा}}{\lemma{\skt{वा}}\Afootnote{K H; \skt{च । औपत्यके मद्यसमुद्भवे च हेतुं ज्वरे पित्तकृतं वदन्ति ॥}\ A}}\skt{ ।}\pend

\pstart \ledsidenote{\passref{6.39.57.add}}\edtext{\skt{सन्ता\-प\-पा\-ना\-भ्य\-स\-म्भ\-वञ्च}}{\Afootnote{H; \skt{सन्ता\-प\-पा\-ना\-त्य\-य\-सं\-भञ्च}\ K}}\skt{ पित्ता\-धि\-क\-त्वा\-त्प्र\-व\-दन्ति तज्ज्ञाः ॥}\pend

\pstart\versepar \ledsidenote{\passref{6.39.58}}\edtext{\skt{प्रलेपकञ्चापि}}{\lemma{\skt{चापि}}\Afootnote{K H; \skt{वात}\ A}}\skt{ }\edtext{\skt{बलासकञ्च}}{\Afootnote{A H; \skt{वलासकञ्च}\ K}}\skt{ कफाधिकत्वेन वदन्ति तज्ज्ञाः । }\pend
\pstart\versepar \skt{}\edtext{\edtext{\skt{मुक्ता\-नु\-ब\-न्ध\-म्वि\-ष\-म\-ज्वरा}}{\lemma{\skt{मुक्तानुबन्धम्}}\Afootnote{H; \skt{मूर्च्छानुबन्धा}\ A; \skt{मुत्वक्त्वानुवन्धं}\ K\textsuperscript{ac}; \skt{मुक्त्वानुवन्धं}\ K\textsuperscript{pc}}}\skt{ ये}}{\lemma{\skt{विषमज्वरा ये}}\Afootnote{K H; \skt{विषमज्वरा ये}\ A}}\skt{ }\edtext{\skt{प्रायेन ते}}{\Afootnote{K; \skt{प्रायेण ते}\ A; \skt{ज्वरे शते}\ H}}\skt{ द्वन्द्व\-स\-मु\-त्थि\-तास्तु ॥}\pend

\pstart\versepar \ledsidenote{\passref{6.39.59}}\skt{त्वक्स्थौ श्लेष्मानिलौ }\edtext{\skt{शीतमादौ}}{\Afootnote{A; \skt{शीतमादौ}\ K H}}\skt{ जनयतो ज्वरे । }\pend
\pstart\versepar \skt{तयोः प्रशान्तयोः }\edtext{\edtext{\skt{पित्तमन्ते}}{\lemma{\skt{पित्तम्}}\Afootnote{H; \skt{पित्त}\ A K}}}{\lemma{\skt{अन्ते}}\Afootnote{H; \skt{मन्ते}\ A K}}\skt{ दाहङ्करोति च ॥}\pend

\pstart\versepar \ledsidenote{\passref{6.39.60}}\skt{करोत्यादौ तथा पित्तं त्वक्स्थं }\edtext{\skt{दाहमिती}}{\lemma{\skt{इती}}\Afootnote{H; \skt{अती}\ A; \skt{इति}\ K}}\skt{ }\edtext{\skt{स्थितिः}}{\Afootnote{K H; \emph{om.}\ A}}\skt{ । }\pend
\pstart\versepar \skt{}\edtext{\skt{तस्मिन्प्रशान्ते}}{\Afootnote{H; \skt{व च । प्रशान्ते}\ A; \skt{तस्मिंप्रशान्ते}\ K}}\skt{ }\edtext{\edtext{\skt{शीतमन्ते}}{\lemma{\skt{शीतम्}}\Afootnote{\skt{कुरु\-त\-स्त\-स्मिं\-श्छी\-तम्}\ A; \skt{त्वितरौ शीत}\ K; \skt{त्वितरौ… शीतम्}\ H}}\skt{ च}}{\lemma{\skt{अन्ते च}}\Afootnote{A H; \skt{मन्ते च}\ K}}\skt{ }\edtext{\skt{कुर्व्वतः}}{\Afootnote{K H; \skt{तावपि ॥}\ A}}\skt{ ॥}\pend

\pstart\versepar \ledsidenote{\passref{6.39.61}}\edtext{\skt{तावेतौ}}{\Afootnote{\skt{द्वावेतौ}\ A; \skt{तावतौ}\ K; \skt{तावेतौ}\ H}}\skt{ }\edtext{\skt{दाहशीतादी}}{\Afootnote{A K; \skt{दाहशीतादी…}\ H}}\skt{ ज्वरौ संसर्ग्गजौ स्मृतौ । }\pend
\pstart\versepar \skt{}\edtext{\skt{दाहपूर्व्वस्तयोः}}{\Afootnote{\skt{दाहपूर्वस्तयोः}\ A H; \skt{दाहपूर्वतयोः}\ K}}\skt{ कष्टः }\edtext{\skt{कृच्छ्रसाध्यतमो}}{\Afootnote{H; \skt{कृच्छ्रसाध्यश्च स}\ A; \skt{कृच्छ्रसाध्यातमो}\ K}}\skt{ }\edtext{\skt{मतः}}{\Afootnote{K H; \skt{स्मृतः ॥}\ A}}\skt{ ॥}\pend

\pstart \ledsidenote{\passref{6.39.62}}\edtext{\skt{रात्र्यह्नोः}}{\Afootnote{K H; \skt{प्रस\-क्त\-श्चा\-भि\-घा\-तो\-त्थ\-श्चे\-त\-ना\-प्र\-भ\-वस्तु यः । रात्र्यह्नोः}\ A}}\skt{ षट्सु कालेषु कीर्त्तितेषु }\edtext{\skt{यथा}}{\Afootnote{K H; \skt{तथा}\ A}}\skt{ पुरा ।}\pend

\pstart\versepar \ledsidenote{\passref{6.39.63}}\edtext{\skt{स्वं}}{\Afootnote{K H; \skt{प्र}\ A}}\skt{ }\edtext{\edtext{\skt{काल\-म्वि\-ष\-मोऽ\-भ्येति}}{\lemma{\skt{कालम्}}\Afootnote{K H; \skt{सह्य}\ A}}}{\lemma{\skt{ऽभ्येति}}\Afootnote{A H; \skt{ह्यति}\ K}}\skt{ }\edtext{\skt{प्रसह्यं}}{\Afootnote{K; \skt{मानवं}\ A; \skt{प्रसक्तं}\ H}}\skt{ }\edtext{\skt{मारुतो}}{\Afootnote{K H; \skt{बहुधा}\ A}}\skt{ ज्वरः ॥ }\pend
\pstart\versepar \skt{स चापि विषमो देहन्न }\edtext{\skt{कदा\-चि\-द्वि\-मु\-ञ्चति}}{\Afootnote{K H; \skt{कदा\-चि\-द्वि\-मु\-ञ्चति ॥}\ A}}\skt{ ।}\edtext{}{\lemma{\skt{ज्वरः ॥ स}}\Afootnote{A H; \skt{ज्वरं । स}\ K}}\pend

\pstart\versepar \ledsidenote{\passref{6.39.64}}\edtext{\skt{यस्मा\-द्गौ\-र\-व\-वै\-वर्ण्ण्यं}}{\Afootnote{H; \skt{ग्लानिगौरव}\ A; \skt{यस्मा\-द्गौ\-र\-व\-वै\-वर्ण्य}\ K}}\skt{ }\edtext{\skt{कार्श्येभ्यो}}{\Afootnote{K H; \skt{कार्श्येभ्यः स}\ A}}\skt{ }\edtext{\skt{न}}{\Afootnote{K H; \skt{यस्मान्न}\ A}}\skt{ }\edtext{\skt{विमुच्यते}}{\Afootnote{K H; \skt{प्रमुच्यते ।}\ A}}\skt{ ॥ }\pend
\pstart\versepar \skt{वेगे तु समतिक्रान्ते }\edtext{\skt{गतो}}{\Afootnote{A; \skt{गतो}\ K H}}\skt{ }\edtext{\skt{यमिति}}{\lemma{\skt{यम्}}\Afootnote{K H; \skt{ऽयम्}\ A}}\skt{ लक्ष्यते ।}\pend

\pstart\versepar \ledsidenote{\passref{6.39.65}}\edtext{\skt{धात्वन्तरस्थो}}{\Afootnote{A H; \skt{धात्वन्तरस्थौ}\ K}}\skt{ }\edtext{\skt{लीनत्वान्न}}{\Afootnote{A K; \skt{लीनत्वा…न्न}\ H}}\skt{ }\edtext{\skt{सौक्ष्म्या\-दु\-प\-ल\-भ्यते}}{\Afootnote{H; \skt{सौक्ष्म्या\-दु\-प\-ल\-भ्यते ।}\ A; \skt{सौक्ष्म्यादुपभ्यते ।}\ K}}\skt{ ॥ }\pend
\pstart\versepar \skt{अल्पदोषे }\edtext{\skt{त्वन}}{\Afootnote{H; \skt{न्धनः}\ A; \skt{न्धन}\ K}}\skt{ }\edtext{\skt{क्षीण}}{\Afootnote{\skt{क्षीणः}\ A; \skt{न्क्षीणां}\ K; \skt{क्षीण…}\ H}}\skt{ }\edtext{\skt{क्षीणेत्वन}}{\Afootnote{H; \skt{क्षीणेन्धन}\ A; \skt{क्षीणोन्धन}\ K}}\skt{ }\edtext{\skt{इपानलः}}{\Afootnote{H; \skt{इवानलः}\ A K}}\skt{ ॥}\pend

\pstart\versepar \ledsidenote{\passref{6.39.66}}\edtext{\edtext{\edtext{\skt{दोषोऽ\-ल्पोऽ\-हि\-त\-सं\-भूतो}}{\lemma{\skt{दोषो}}\Afootnote{A H; \skt{दोषो}\ K}}}{\lemma{\skt{ऽल्पो}}\Afootnote{A H; \emph{om.}\ K}}}{\lemma{\skt{ऽहितसंभूतो}}\Afootnote{A H; \skt{ल्पोहितसंभूतो}\ K}}\skt{ ज्वरोत्सृष्टस्य वा पुनः । }\pend
\pstart\versepar \skt{}\edtext{\edtext{\skt{धातु\-मु\-त्त\-र\-न्य\-तमं}}{\lemma{\skt{धातुम्}}\Afootnote{A K; \skt{धातु}\ H}}}{\lemma{\skt{उत्तरन्यतमं}}\Afootnote{\skt{अन्यतमं}\ A K; \skt{मन्यतमं}\ H}}\skt{ प्राप्य करोति विषमज्वरं ॥}\pend

\pstart\versepar \ledsidenote{\passref{6.39.67}}\edtext{\skt{सततं}}{\Afootnote{A H; \skt{सततᳵ}\ K}}\skt{ }\edtext{\edtext{\skt{रस\-र\-क्त\-स्थ\-म\-न्येद्युः}}{\lemma{\skt{रसरक्तस्थम्}}\Afootnote{K H; \skt{रसरक्तस्थः}\ A}}}{\lemma{\skt{अन्येद्युः}}\Afootnote{H; \skt{सोऽन्येद्युः}\ A; \skt{सान्यैद्युः}\ K}}\skt{ पिशिताश्रितः । }\pend
\pstart\versepar \skt{}\edtext{\edtext{\skt{मेदो\-ग\-त\-स्तृ\-ती\-येऽह्नि}}{\lemma{\skt{मेदोगतस्}}\Afootnote{A H; \skt{मेदोगदस्}\ K}}}{\lemma{\skt{ऽह्नि}}\Afootnote{A H; \skt{ह्नि}\ K}}\skt{ }\edtext{\edtext{\skt{सोऽस्थिमज्जगतः}}{\lemma{\skt{सो}}\Afootnote{H; \emph{om.}\ A; \skt{सो}\ K}}}{\lemma{\skt{ऽस्थिमज्जगतः}}\Afootnote{H; \skt{त्वस्थिमज्जगतः}\ A; \skt{स्थिमज्जगतः}\ K}}\skt{ पुनः ॥}\pend

\pstart \ledsidenote{\passref{6.39.68}}\skt{कुर्याच्चातुर्थकं घोरमन्तकं रोगसंकरं ।}\pend

\pstart\versepar \ledsidenote{\passref{6.39.69}}\skt{सप्ताहम्वा दशाहम्वा द्वाद\-शा\-ह\-म\-थापि वा ॥ }\pend
\pstart\versepar \skt{}\edtext{\skt{सन्तत्या यो}}{\Afootnote{A; \skt{सन्तत्या}\ K H}}\skt{ }\edtext{\skt{विसर्गी}}{\Afootnote{\skt{ऽविसर्गी}\ A; \skt{योविसर्गी}\ K H}}\skt{ स्यात्सन्ततः स निगद्यते ।}\pend

\pstart\versepar \ledsidenote{\passref{6.39.70}}\skt{अहोरात्रे }\edtext{\skt{सन्ततकौ द्वौ}}{\Afootnote{H; \skt{सततको द्वौ}\ A; \skt{सन्ततको द्वौ}\ K}}\skt{ }\edtext{\edtext{\skt{काला\-व\-नु\-व\-र्त्तते}}{\lemma{\skt{कालाव्}}\Afootnote{K H; \skt{काला}\ A}}}{\lemma{\skt{अनुवर्त्तते}}\Afootnote{K H; \skt{वनुवर्तते ।}\ A}}\skt{ ॥ }\pend
\pstart\versepar \skt{}\edtext{\edtext{\edtext{\skt{अन्ये\-द्यु\-ष्क\-स्त्व\-हो\-रा\-त्र\-मे\-क\-कालं}}{\lemma{\skt{अन्येद्युष्कस्}}\Afootnote{\skt{अन्येद्युष्क}\ A; \skt{अन्येदुष्कस्}\ K; \skt{अन्येद्युष्क}\ H}}}{\lemma{\skt{त्वहोरात्रम्}}\Afootnote{K; \skt{स्त्वहोरात्रा}\ A; \skt{स्त्वहोरात्रम्}\ H}}}{\lemma{\skt{एककालं}}\Afootnote{K H; \skt{देककालं}\ A}}\skt{ प्रवर्त्तते ।}\pend

\pstart \ledsidenote{\passref{6.39.71}}\edtext{\edtext{\edtext{\skt{तृती\-य\-क\-स्तृ\-ती\-येऽह्नि}}{\lemma{\skt{तृतीयकस्}}\Afootnote{K; \skt{तृतीयक}\ A H}}}{\lemma{\skt{तृतीये}}\Afootnote{K; \skt{स्तृतीये}\ A H}}}{\lemma{\skt{ऽह्नि}}\Afootnote{A; \skt{ह्नि}\ K; \skt{ऽह्नि…}\ H}}\skt{ }\edtext{\skt{चतुर्थेऽह्नि}}{\lemma{\skt{ऽह्नि}}\Afootnote{A; \skt{ह्नि}\ K H}}\skt{ }\edtext{\skt{चतुर्थकः}}{\Afootnote{K H; \skt{चतुर्थकः । वाते\-नो\-दी\-र्य\-मा\-णाश्च ह्रीयमाणाश्च सर्वतः । एकद्विदोषा मर्त्यानां तस्मि\-न्ने\-वो\-दि\-तेऽ\-हनि ॥}\ A}}\skt{ ॥}\pend

\pstart \ledsidenote{\passref{6.39.68cd}}\edtext{\skt{केचि\-द्भू\-ता\-भि\-षङ्गो त्थं}}{\Afootnote{H; \skt{केचि\-द्भू\-ता\-भि\-ष\-ङ्गोत्थं}\ A; \skt{केचि\-द्भू\-ता\-भि\-ष\-ङ्कोत्थं}\ K}}\skt{ }\edtext{\skt{ब्रुवते}}{\Afootnote{A H; \skt{व्रुवते}\ K}}\skt{ }\edtext{\skt{विषमज्वरं}}{\Afootnote{A H; \skt{विषमज्वरः ।}\ K}}\skt{ ।}\pend

\pstart \ledsidenote{\passref{6.39.75}}\edtext{\skt{विवि\-धे\-ना\-भि\-घा\-तेन}}{\Afootnote{K H; \skt{वेगहानौ प्रशाम्येत यथाम्भः सागरे तथा । विवि\-धे\-ना\-भि\-घा\-तेन}\ A}}\skt{ ज्वरो }\edtext{\skt{यस्तु}}{\lemma{\skt{तु}}\Afootnote{K H; \skt{सं}\ A}}\skt{ }\edtext{\skt{प्रवर्त्तते}}{\Afootnote{A K; \skt{प्रवर्त्तन्ते ॥}\ H}}\skt{ ॥}\pend

\pstart\versepar \ledsidenote{\passref{6.39.76}}\edtext{\skt{यथा दोषप्रकोपस्तु}}{\Afootnote{H; \skt{यथा\-दो\-ष\-प्र\-कोपं तु}\ A K}}\skt{ तथा मन्येत }\edtext{\skt{तं}}{\Afootnote{A H; \skt{त}\ K}}\skt{ }\edtext{\skt{ज्वरं}}{\Afootnote{A H; \skt{ज्ज्वरᳵ।}\ K}}\skt{ । }\pend
\pstart\versepar \skt{श्यावास्यता विषकृते }\edtext{\skt{तथातीसार}}{\Afootnote{K H; \skt{दाहातीसार}\ A}}\skt{ }\edtext{\skt{एव}}{\Afootnote{K H; \emph{om.}\ A}}\skt{ }\edtext{\skt{च}}{\Afootnote{K H; \skt{हृद्ग्रहाः ॥}\ A}}\skt{ ॥}\pend

\pstart\versepar \ledsidenote{\passref{6.39.77}}\edtext{\skt{अभ\-क्त\-रु\-क्पि\-पासा च}}{\Afootnote{A H; \skt{अभ\-क्ता\-रु\-च्पि\-पासा च}\ K}}\skt{ }\edtext{\skt{तोदश्च}}{\lemma{\skt{च}}\Afootnote{K H; \skt{मू}\ A}}\skt{ }\edtext{\skt{सहमूर्च्छ्या}}{\Afootnote{H; \skt{र्च्छा बलक्षयः ।}\ A; \skt{सह मूर्च्छया ।}\ K}}\skt{ । }\pend
\pstart\versepar \skt{ओषधीगन्धजे मूर्च्छा }\edtext{\edtext{\skt{शिरो\-रु\-क्क्व\-य\-थु\-स्तथा}}{\lemma{\skt{शिरो\-रु\-क्क्व\-य\-थुस्}}\Afootnote{\skt{शिरोरुक्वमथुः}\ A; \skt{सिरो\-रु\-क्क्ष\-व\-य\-थुस्}\ K; \skt{शिरोरुक्थ्व यथुस्}\ H}}}{\lemma{\skt{तथा}}\Afootnote{K H; \skt{क्षवः ॥}\ A}}\skt{ ॥}\pend

\pstart \ledsidenote{\passref{6.39.78}}\skt{कामजे }\edtext{\edtext{\skt{चित्त\-वि\-भ्रं\-स\-स्त\-न्द्रा\-ल\-स्य\-म\-भ\-क्त\-रुक्}}{\lemma{\skt{चित्त\-वि\-भ्रं\-स\-स्त\-न्द्रा\-ल\-स्यम्}}\Afootnote{\skt{चित्त\-वि\-भ्रं\-श\-स्त\-न्द्रा\-ल\-स्यम्}\ A K; \skt{चित्तविभ्रं स…स्तन्द्रालस्यम्}\ H}}}{\lemma{\skt{अभक्तरुक्}}\Afootnote{K H; \skt{अरोचकः । हृदये वेदना चास्य गात्रं च परिशुष्यति ॥}\ A}}\skt{ ।}\pend

\pstart\versepar \ledsidenote{\passref{6.39.79}}\edtext{\skt{भयात्प्रलापः}}{\Afootnote{A H; \skt{भयां प्रलापः}\ K}}\skt{ शोकाच्च भवेत्कोपाच्च वेपथुः ॥ }\pend
\pstart\versepar \skt{}\edtext{\skt{अभि\-शा\-पा\-भि\-चा\-राभ्यां}}{\Afootnote{K H; \skt{अभि\-चा\-रा\-भि\-शा\-पाभ्यां}\ A}}\skt{ मोहस्तृष्णा च जायते ॥}\pend

\pstart\versepar \ledsidenote{\passref{6.39.80}}\edtext{\skt{भूता\-भि\-ष\-ङ्गा\-दु\-द्वे\-ग\-हा\-स्य\-रो\-द\-न\-क\-म्पनं}}{\Afootnote{H; \skt{भूता\-भि\-ष\-ङ्गा\-दु\-द्वे\-ग\-हा\-स्य\-क\-म्प\-न\-रो\-द\-नम् ।}\ A; \skt{भूता\-भि\-स\-ङ्गा\-दु\-द्वेगः हास्य\-रो\-द\-न\-कं\-पनᳵ ।}\ K}}\skt{ ॥ }\pend
\pstart\versepar \skt{}\edtext{\skt{श्रम\-क्ष\-या\-भि\-घा\-तोत्थे}}{\Afootnote{K H; \skt{श्रम\-क्ष\-या\-भि\-घा\-तेभ्यो}\ A}}\skt{ देहिनां }\edtext{\skt{कुपितोनिलः}}{\Afootnote{K H; \skt{कुपितोऽनिलः ॥}\ A}}\skt{ ।}\pend

\pstart\versepar \ledsidenote{\passref{6.39.81}}\skt{पूर\-यि\-त्वा\-खि\-ल\-न्देहं }\edtext{\skt{ज्वर\-मा\-पा\-द\-ये\-द्भृशं}}{\lemma{\skt{आपादयेद्भृशं}}\Afootnote{A H; \skt{आपदयेद्भृशᳵ ।}\ K}}\skt{ ॥ }\pend
\pstart\versepar \skt{रोगाणान्तु }\edtext{\skt{समु\-त्था\-ना\-द्वि\-दा\-हा\-जी\-र्ण्ण\-त\-स्तथा}}{\lemma{\skt{समु\-त्था\-ना\-द्वि\-दा\-हा\-जी\-र्ण्ण\-तस्}}\Afootnote{H\textsuperscript{pc}; \skt{समु\-त्था\-ना\-द्वि\-दा\-हा\-ग\-न्तु\-तस्}\ A; \skt{सम\-त्था\-ना\-द्वि\-दा\-हा\-जी\-र्ण्ण\-तस्}\ K; \skt{समु\-त्था\-नां\-द्वि\-दा\-हा\-जी\-र्ण्ण\-तस्}\ H\textsuperscript{ac}}}\skt{ ।}\pend

\pstart\versepar \ledsidenote{\passref{6.39.82}}\edtext{\skt{ज्वरो}}{\Afootnote{A; \skt{ज्वरो}\ K H}}\skt{ }\edtext{\skt{परः}}{\Afootnote{A; \skt{परः}\ K H}}\skt{ }\edtext{\skt{संभवति तैस्तैरन्यैश्च}}{\lemma{\skt{संभवति तैस्तैर्}}\Afootnote{A H; \skt{सम्भवति तैस्तैर्}\ K}}\skt{ हेतुभिः ॥ }\pend
\pstart\versepar \skt{}\edtext{\skt{तस्य}}{\Afootnote{K H; \skt{दोषाणां स तु लिङ्गानि कदा\-चि\-न्ना\-ति\-व\-र्तते ॥}\ A}}\skt{ }\edtext{\skt{रूपाणि}}{\Afootnote{K H; \emph{om.}\ A}}\skt{ }\edtext{\skt{वक्ष्यामि}}{\Afootnote{K H; \emph{om.}\ A}}\skt{ }\edtext{\skt{यथा\-व\-द\-नु\-पू\-र्व्वशः}}{\Afootnote{K H; \emph{om.}\ A}}\skt{ ।}\pend

\pstart\versepar \ledsidenote{\passref{6.39.83}}\skt{गुरुता }\edtext{\skt{हृदयोत्क्लेशः}}{\Afootnote{A K\textsuperscript{pc} H; \skt{हृयोत्क्लेशः}\ K\textsuperscript{ac}}}\skt{ }\edtext{\edtext{\skt{सद\-न\-श्छ\-र्द्य\-रो\-चकौ}}{\lemma{\skt{सदनश्}}\Afootnote{H; \skt{सदनं}\ A; \skt{सीदनश्}\ K}}}{\lemma{\skt{छर्द्यरोचकौ}}\Afootnote{A; \skt{छर्द्यरोचकौ}\ K H}}\skt{ ॥ }\pend
\pstart\versepar \skt{}\edtext{\skt{रसस्थेतु}}{\Afootnote{A K H\textsuperscript{pc}; \skt{रसस्थे}\ H\textsuperscript{ac}}}\skt{ ज्वरे लिङ्गं दैन्यं }\edtext{\skt{चास्योपजायते}}{\Afootnote{A K\textsuperscript{pc} H; \skt{चास्योपजाते ।}\ K\textsuperscript{ac}}}\skt{ ।}\pend

\pstart\versepar \ledsidenote{\passref{6.39.84}}\edtext{\skt{रक्तोनिष्ठीवनं}}{\Afootnote{H; \skt{रक्तनिष्ठीवनं}\ A K}}\skt{ दाहो }\edtext{\edtext{\skt{मोह\-श्छ\-र्द्द\-न\-वि\-भ्रमौ}}{\lemma{\skt{मोहश्}}\Afootnote{H; \skt{स्वेदश्}\ A; \skt{शोफश्}\ K}}}{\lemma{\skt{छर्द्दनविभ्रमौ}}\Afootnote{A H; \skt{छर्दनविभ्रमः ।}\ K}}\skt{ ॥ }\pend
\pstart\versepar \skt{प्रलापः }\edtext{\skt{पिटकास्तृष्णा}}{\lemma{\skt{पिटकास्}}\Afootnote{K H; \skt{पिटिका}\ A}}\skt{ रक्तप्राप्ते ज्वरे नृणां ॥}\pend

\pstart\versepar \ledsidenote{\passref{6.39.85}}\skt{पिण्डिकोद्वेष्टनं तृष्णा सृष्ट\-मू\-त्र\-पु\-री\-षता । }\pend
\pstart\versepar \skt{}\edtext{\skt{ऊष्मा\-न्त\-र्द्दा\-ह\-वि\-क्षेपौ}}{\Afootnote{A; \skt{ऊष्मां\-त\-दा\-ह\-वि\-क्षेपौ}\ K; \skt{ऊष्मा\-न्त\-र्द्दा\-ह\-वि\-क्षेपौ}\ H}}\skt{ }\edtext{\skt{ग्लानिः}}{\Afootnote{A H; \skt{ग्लानि}\ K}}\skt{ }\edtext{\skt{स्यात्मांसगे}}{\Afootnote{K H; \skt{स्यान्मांसगे}\ A}}\skt{ ज्वरे ॥}\pend

\pstart\versepar \ledsidenote{\passref{6.39.86}}\skt{भृशं }\edtext{\skt{स्वेदस्तृषा}}{\Afootnote{A; \skt{स्वेदस्तृषा}\ K H}}\skt{ मूर्च्छा }\edtext{\edtext{\skt{प्रला\-प\-श्छ\-र्द्दि\-रेव}}{\lemma{\skt{छर्द्दिर्}}\Afootnote{K H; \skt{छर्दि}\ A}}\skt{ च}}{\lemma{\skt{एव च}}\Afootnote{K H; \skt{रेव च ।}\ A}}\skt{ । }\pend
\pstart\versepar \skt{}\edtext{\skt{दौर्ग्गन्ध्यारोचकौ}}{\Afootnote{A H; \skt{दौर्गंध्यारोचक}\ K}}\skt{ }\edtext{\edtext{\skt{ग्लानिर्म्मेदःस्थे}}{\lemma{\skt{ग्लानिर्}}\Afootnote{K H; \skt{ग्लानि}\ A}}}{\lemma{\skt{म्मेदःस्थे}}\Afootnote{H; \skt{र्मेदःस्थे}\ A; \skt{मेदस्थे}\ K}}\skt{ चासहिष्णुता ॥}\pend

\pstart\versepar \ledsidenote{\passref{6.39.87}}\edtext{\edtext{\skt{भेदोऽन्त्र}}{\lemma{\skt{भेदो}}\Afootnote{A; \emph{om.}\ K H}}}{\lemma{\skt{ऽन्त्र}}\Afootnote{\skt{ऽस्थ्नां कुञ्चनं(कुञ्}\ A; \skt{भेदोन्त्र}\ K H}}\skt{ }\edtext{\skt{कूजनं}}{\Afootnote{K H; \skt{चजनं)}\ A}}\skt{ श्वासो }\edtext{\edtext{\skt{विरे\-क\-श्छ\-र्द्दि\-रेव}}{\lemma{\skt{छर्द्दिर्}}\Afootnote{K H; \skt{छर्दि}\ A}}\skt{ च}}{\lemma{\skt{एव च}}\Afootnote{K H; \skt{रेव च ।}\ A}}\skt{ । }\pend
\pstart\versepar \skt{विक्षेपणञ्च }\edtext{\skt{गात्रा मेतदस्थिगते}}{\Afootnote{\skt{गात्रा\-णा\-मे\-त\-द\-स्थि\-गते}\ A; \skt{गात्राणां एतदस्थिगते}\ K H}}\skt{ ज्वरे ॥}\pend

\pstart\versepar \ledsidenote{\passref{6.39.88}}\edtext{\skt{तमः}}{\Afootnote{A H; \skt{तम}\ K}}\skt{ प्रवेशनं हिक्का कासः }\edtext{\skt{शैत्यम्वमिस्स्तृषा तथा}}{\lemma{\skt{वमिस्स्तृषा तथा}}\Afootnote{\skt{वमिस्तथा}\ A K H\textsuperscript{ac}; \skt{वमिस्स्तृषा तथा ।}\ H\textsuperscript{pc}}}\skt{ । }\pend
\pstart\versepar \skt{अन्तर्द्दाहो महाश्वासो मर्म्मच्छेदश्च मज्जगे ॥ }\pend
\pstart\versepar \skt{मरणं }\edtext{\skt{प्राप्नुयात्तत्र}}{\Afootnote{H; \skt{प्राप्रुयात्तत्र}\ A; \skt{चाप्नुयान्तत्र}\ K}}\skt{ शुक्रस्थानगते }\edtext{\skt{ज्वरेपरं}}{\Afootnote{H\textsuperscript{pc}; \skt{ज्वरे ॥}\ A; \skt{परᳵ ॥}\ K; \skt{परं ।}\ H\textsuperscript{ac}}}\skt{ ।}\pend

\pstart\versepar \ledsidenote{\passref{6.39.89}}\edtext{\skt{शेफसः}}{\Afootnote{A H; \skt{शेफस}\ K}}\skt{ }\edtext{\skt{स्तब्धता}}{\Afootnote{A H; \skt{स्तव्धता}\ K}}\skt{ }\edtext{\skt{मोक्षः}}{\Afootnote{A H; \skt{मोक्ष}\ K}}\skt{ }\edtext{\skt{शुक्लस्य तु}}{\Afootnote{H; \skt{शुक्रस्य तु}\ A K}}\skt{ विशेषतः ॥ }\pend
\pstart\versepar \skt{}\edtext{\edtext{\skt{दग्धेऽन्धनो}}{\lemma{\skt{दग्धे}}\Afootnote{K H; \skt{दग्ध्वे}\ A}}}{\lemma{\skt{ऽन्धनो}}\Afootnote{H; \skt{न्धनं}\ A; \skt{त्वंनो}\ K}}\skt{ यथा }\edtext{\edtext{\skt{वह्निर्द्धातुं}}{\lemma{\skt{वह्निर्}}\Afootnote{K H; \skt{वह्नि}\ A}}}{\lemma{\skt{द्धातुं}}\Afootnote{K H; \skt{र्धातू}\ A}}\skt{ }\edtext{\skt{हत्वा}}{\Afootnote{K H; \skt{न्हत्वा}\ A}}\skt{ यथा विषं ।}\pend

\pstart\versepar \ledsidenote{\passref{6.39.90}}\skt{कृतकृत्यो व्रजेच्छान्तिन्देहं हत्वा तथा ज्वरः ॥ }\pend
\pstart\versepar \skt{वातपित्तकफो त्थानां ज्वराणां लक्षणं यथा ।}\pend

\pstart\versepar \ledsidenote{\passref{6.39.91}}\skt{तथा }\edtext{\skt{तेषा\-म्भि\-ष\-ग्ब्रू\-या\-द्र\-सा\-दि\-ष्वपि}}{\lemma{\skt{भिष\-ग्ब्रू\-या\-द्र\-सा\-दिष्व्}}\Afootnote{A H; \skt{भिष\-ग्व्रू\-या\-द्र\-सा\-दिष्व्}\ K}}\skt{ बुद्धिमान् ॥ }\pend
\pstart\versepar \skt{समस्तैः }\edtext{\skt{सन्निपातस्य}}{\Afootnote{K H; \skt{सन्निपातेन}\ A}}\skt{ }\edtext{\skt{धातुस्थस्य}}{\Afootnote{K H; \skt{धातुस्थम}\ A}}\skt{ }\edtext{\skt{विनिर्द्दिशेत्}}{\Afootnote{H; \skt{पि निर्दिशेत् ॥}\ A; \skt{विर्निर्दिशेत ।}\ K}}\skt{ ।}\edtext{}{\lemma{\skt{बुद्धिमान् ॥ समस्तैः}}\Afootnote{A H; \skt{वुद्धिमां । समस्तैस्}\ K}}\pend

\pstart\versepar \ledsidenote{\passref{6.39.92}}\skt{द्वन्द्वजं द्वन्द्वजैरेव दोषैश्चापि }\edtext{\skt{वदेत्तु तं}}{\Afootnote{H; \skt{वदेत्कृतम् ।}\ A; \skt{वदेतुतᳵ ।}\ K}}\skt{ ॥ }\pend
\pstart\versepar \skt{गम्भीरस्तु ज्वरो }\edtext{\skt{ज्ञेयो ह्यन्तर्द्दाहेन}}{\lemma{\skt{ज्ञेयो ह्य्}}\Afootnote{A K\textsuperscript{pc} H; \skt{ज्ञेया ज्ञेयो ह्य्}\ K\textsuperscript{ac}}}\skt{ तृष्णया ।}\pend

\pstart\versepar \ledsidenote{\passref{6.39.93}}\skt{आनद्धत्वेन }\edtext{\skt{चात्यर्थं}}{\Afootnote{A; \skt{चात्यर्थं}\ K H}}\skt{ }\edtext{\skt{श्वास\-का\-सो\-द्भ\-वेन च}}{\Afootnote{H; \skt{श्वास\-का\-सो\-द्ग\-मेन च ।}\ A; \skt{श्वास\-का\-सो\-द्ग\-वेन च ।}\ K}}\skt{ ॥ }\pend
\pstart\versepar \skt{हतप्रभेन्द्रियं }\edtext{\skt{क्षाम}}{\Afootnote{H; \skt{क्षीणम्}\ A; \skt{क्षामम्}\ K}}\skt{ }\edtext{\skt{मरो\-च\-क\-नि\-पी\-डितं}}{\Afootnote{H; \skt{अरो\-च\-क\-नि\-पी\-डि\-तम्}\ A K}}\skt{ ।}\pend

\pstart\versepar \ledsidenote{\passref{6.39.94}}\edtext{\edtext{\skt{गम्भी\-र\-स्ती\-क्ष्ण\-वे\-गार्त्ता}}{\lemma{\skt{गम्भीरस्}}\Afootnote{H; \skt{गम्भीर}\ A; \skt{गंभीरस्}\ K}}}{\lemma{\skt{तीक्ष्णवेगार्त्ता}}\Afootnote{K; \skt{तीक्ष्णवेगार्तं}\ A; \skt{तीक्ष्णवेगा र्त्ते}\ H}}\skt{ ज्वरितं परिवर्ज्जयेत् ॥ }\pend
\pstart\versepar \skt{}\edtext{\skt{हीन\-म\-ध्या\-धि\-कै\-र्द्दो\-षै\-स्त्रि\-स\-प्त\-द्वा\-द\-शा\-दिकः}}{\lemma{\skt{त्रिस\-प्त\-द्वा\-द\-शा\-दिकः}}\Afootnote{H; \skt{त्रिस\-प्त\-द्वा\-द\-शा\-हिकः ॥}\ A; \skt{त्रिस\-प्त\-द्वा\-द\-शा\-हिकाः ।}\ K}}\skt{ ।}\pend

\pstart \ledsidenote{\passref{6.39.95}}\skt{ज्वरवेगो }\edtext{\skt{भवेत्तीक्ष्णो}}{\Afootnote{K H; \skt{भवेत्तीव्रो}\ A}}\skt{ यथापूर्व्वं }\edtext{\skt{सुखक्रियः}}{\Afootnote{H; \skt{सुखक्रियः । कालो ह्येष यमश्चैव निय\-ति\-र्मृ\-त्यु\-रेव च ॥}\ A; \skt{सुखक्रिया ॥}\ K}}\skt{ ॥}\pend

\pstart \ledsidenote{\passref{6.39.96}}\edtext{\skt{इति}}{\Afootnote{K H; \skt{तस्मिन्व्यपगते देहाज्जन्मेह पुनरुच्यते । इति}\ A}}\skt{ ज्वराः }\edtext{\skt{समाख्याताः}}{\Afootnote{A H; \skt{समाख्याता}\ K}}\skt{ कर्म्मेदानीं प्रवक्ष्यते ॥ ० ॥}\pend

\pstart\versepar \ledsidenote{\passref{6.39.97}}\skt{ज्वरस्य }\edtext{\skt{पूर्व्वरूपेषु}}{\Afootnote{A H; \skt{पूर्वरूपस्य}\ K}}\skt{ }\edtext{\skt{वर्त्तमानस्य}}{\Afootnote{K H; \skt{वर्तमानेषु}\ A}}\skt{ बुद्धिमान् । }\pend
\pstart\versepar \skt{पाययेत }\edtext{\edtext{\skt{सर्प्पि\-रे\-वा\-च्छ\-न्ततः}}{\lemma{\skt{सर्प्पिर्}}\Afootnote{H; \skt{घृ}\ A; \skt{र्पिर्}\ K}}}{\lemma{\skt{एवाच्छन्}}\Afootnote{K H; \skt{तं स्वच्छं}\ A}}\skt{ }\edtext{\skt{सम्पद्यते}}{\Afootnote{K H; \skt{स लभते}\ A}}\skt{ }\edtext{\skt{सुखं}}{\Afootnote{A H; \skt{सुखᳵ ।}\ K}}\skt{ ॥}\edtext{}{\lemma{\skt{बुद्धिमान् । पाययेत}}\Afootnote{A H; \skt{वुद्धिमान् ॥ पाययेत्स}\ K}}\pend

\pstart\versepar \ledsidenote{\passref{6.39.98}}\edtext{\skt{विधि\-र्म्मा\-रु\-त\-जे\-ष्वेष}}{\lemma{\skt{विधिर्}}\Afootnote{A H; \skt{विधि}\ K}}\skt{ }\edtext{\skt{पैत्तिकेषु}}{\Afootnote{A K\textsuperscript{pc} H; \skt{पैकेषु}\ K\textsuperscript{ac}}}\skt{ }\edtext{\skt{विरेचनम्}}{\Afootnote{A H; \skt{विरेचनᳵ ।}\ K}}\skt{ । }\pend
\pstart\versepar \skt{मृदु प्रच्छर्द्दनं }\edtext{\skt{द्वन्द्वन्कफजेषु}}{\Afootnote{K; \skt{तद्वत्कफजेषु}\ A; \skt{द्वन्द्व कफजेषु}\ H}}\skt{ विधीयते ।}\pend

\pstart\versepar \ledsidenote{\passref{6.39.100}}\edtext{\edtext{\skt{धूमाग्न्योरिव}}{\lemma{\skt{धूमाग्न्योर्}}\Afootnote{H; \skt{रूप\-प्रा\-ग्रू\-प\-योर्}\ A; \skt{धूमाग्श्चर्}\ K}}}{\lemma{\skt{इव}}\Afootnote{K H; \skt{विद्यान्}\ A}}\skt{ }\edtext{\skt{नानात्वं}}{\Afootnote{A H; \skt{नानत्वं}\ K}}\skt{ }\edtext{\skt{विद्या\-त्प्रा\-क्रू\-पि\-रू\-पयोः}}{\Afootnote{H; \skt{वह्निधूमव}\ A; \skt{विद्या\-त्प्रा\-ग्रु\-वि\-भा\-पयाः ।}\ K}}\skt{ ॥ }\pend
\pstart\versepar \skt{}\edtext{\skt{प्रव्यक्तरूपेषु}}{\Afootnote{K H; \skt{त् । प्रव्यक्तरूपेषु}\ A}}\skt{ }\edtext{\skt{हित\-मे\-का\-न्ते\-ना\-प\-त\-र्प्प\-णम्}}{\lemma{\skt{एका\-न्ते\-ना\-प\-त\-र्प्प\-णम्}}\Afootnote{A H; \skt{एका\-न्ते\-ना\-प\-त\-र्पणᳵ ।}\ K}}\skt{ ।}\pend

\pstart \ledsidenote{\passref{6.39.101}}\edtext{\skt{आनद्ध}}{\Afootnote{H; \skt{आमाशयस्थे दोषे तु सोत्क्लेशे वमनं परम् । आनद्धः}\ A; \skt{आनद्धस्}\ K}}\skt{ }\edtext{\skt{स्तिमि\-तै\-र्द्दो\-षै\-र्या\-वन्तं}}{\lemma{\skt{स्तिमितैर्}}\Afootnote{A H; \skt{तिमितैर्}\ K}}\skt{ }\edtext{\skt{काल}}{\Afootnote{K H; \skt{कालम्}\ A}}\skt{ }\edtext{\skt{हृ मातुरः}}{\Afootnote{\skt{आतुरः ॥}\ A; \skt{मातुरः}\ K H}}\skt{ ॥}\pend

\pstart\versepar \ledsidenote{\passref{6.39.102}}\edtext{\skt{तावत्कालन्तु}}{\Afootnote{K H; \skt{कुर्यादनशनं तावत्ततः}\ A}}\skt{ }\edtext{\edtext{\skt{लघ्व\-न्न\-म\-श्नी\-यात्तु}}{\lemma{\skt{लघ्वन्नम्}}\Afootnote{H; \skt{सं}\ A; \skt{ध्व्वन्नम्}\ K\textsuperscript{ac}; \skt{लध्व्वन्नम्}\ K\textsuperscript{pc}}}}{\lemma{\skt{अश्नीयात्तु}}\Afootnote{K H; \skt{अर्गमा}\ A}}\skt{ विरिक्तवत् । }\pend
\pstart\versepar \skt{न }\edtext{\skt{लंघ\-ये\-न्मा\-रु\-तजे}}{\Afootnote{H; \skt{लङ्घ\-ये\-न्मा\-रु\-तजे}\ A; \skt{लङ्घ\-ये\-त्मा\-रु\-तजे}\ K}}\skt{ क्षयजे मानसे तथा ॥}\edtext{}{\lemma{\skt{विरिक्तवत् । न}}\Afootnote{K H; \skt{चरेत् । न}\ A}}\pend

\pstart\versepar \ledsidenote{\passref{6.39.103}}\edtext{\edtext{\skt{अलंघ्याश्चैव}}{\lemma{\skt{अलंघ्याश्}}\Afootnote{K H; \skt{अलङ्घ्याश्}\ A}}\skt{ ये}}{\lemma{\skt{चैव ये}}\Afootnote{H; \skt{चापि ये}\ A K}}\skt{ पूर्व्वन्द्विव्रणीये प्रकीर्त्तिताः ॥ }\pend
\pstart\versepar \skt{तद्धि }\edtext{\skt{मारुतक्ष्त्तृष्णा}}{\Afootnote{H; \emph{om.}\ A K}}\skt{ }\edtext{\skt{मुख\-शो\-ष\-भ्र\-मा\-न्विते}}{\Afootnote{H; \emph{om.}\ A K}}\skt{ ॥ }\pend
\pstart\versepar \skt{}\edtext{\skt{कार्यान्न बा}}{\Afootnote{H; \emph{om.}\ A K}}\skt{ }\edtext{\skt{लेन}}{\Afootnote{H; \emph{om.}\ A K}}\skt{ }\edtext{\skt{वृद्धे न}}{\Afootnote{H; \emph{om.}\ A K}}\skt{ }\edtext{\skt{गर्भ्भिण्योन्न}}{\Afootnote{H; \emph{om.}\ A K}}\skt{ }\edtext{\skt{दुर्ब्बले}}{\Afootnote{H; \emph{om.}\ A K}}\skt{ । }\pend
\pstart\versepar \skt{}\edtext{\skt{शोफे\-षू\-त्थि\-त\-मा\-त्रेषु}}{\Afootnote{H; \skt{ताः ।}\ A K}}\skt{ }\edtext{\skt{व्रणे\-षू\-ग्र\-रु\-जेषु च}}{\Afootnote{H; \emph{om.}\ A K}}\skt{ ॥ }\pend
\pstart\versepar \skt{}\edtext{\skt{अनवस्थि }\edtext{\skt{र्ल्लघ\-न्न\-न्दो\-ष\-पा\-च\-नम्}}{\lemma{\skt{ल्लघन्नन्}}\Afootnote{\skt{तदो\-षा\-ग्ने\-र्ल\-ङ्घनं}\ A; \skt{तदो\-षा\-ग्ने\-र्ल\-न्घनं}\ K; \skt{तदोषाग्ने र्ल्लघन्नन्}\ H}}}{\lemma{\skt{अनवस्थि र्}}\Afootnote{\skt{अनवस्थि}\ A K H}}\skt{ ।}\edtext{}{\lemma{\skt{प्रकीर्त्तिताः ॥ तद्धि}}\Afootnote{\skt{प्रकीर्ति}\ A K; \skt{प्रकीर्त्तिताः… ॥ तद्धि}\ H}}\pend

\pstart\versepar \ledsidenote{\passref{6.39.104}}\skt{ज्वरघ्नन्दीपनं }\edtext{\skt{काङ्क्षा\-रु\-चि\-ला\-घ\-व\-का\-र\-कम्}}{\Afootnote{A H; \skt{काङ्क्षा\-रु\-चि\-ला\-घ\-व\-का\-रकᳵ ।}\ K}}\skt{ ॥ }\pend
\pstart\versepar \skt{}\edtext{\skt{सृष्ट\-मा\-रु\-त\-वि\-ण्मूत्रं}}{\Afootnote{A H; \skt{सृष्ट\-मा\-रु\-त\-वि\-त्मूत्रां}\ K}}\skt{ क्षुत्पिपासासहं }\edtext{\skt{लघम्}}{\Afootnote{H; \skt{लघुम् ॥}\ A; \skt{लघु ॥}\ K}}\skt{ ।}\pend

\pstart\versepar \ledsidenote{\passref{6.39.105}}\skt{प्रसन्नात्मेन्द्रियं }\edtext{\skt{क्षामन्नर द्यात्सुलंघितम्}}{\lemma{\skt{नर द्यात्सुलंघितम्}}\Afootnote{\skt{नरं विद्या\-त्सु\-ल\-ङ्घि\-तम् ।}\ A; \skt{नरं विद्यात्सुलंघितᳵ ।}\ K; \skt{नरम्वि …द्या\-त्सु\-लं\-घि\-तम् ॥}\ H}}\skt{ ॥ }\pend
\pstart\versepar \skt{}\edtext{\skt{बल\-क्ष\-य\-तृ\-षा\-मू\-र्च्छा\-त\-न्द्रा\-नि\-द्रा\-भ्र\-म\-क्लमाः}}{\Afootnote{H; \skt{बलक्षयस्तृषा शोष\-स्त\-न्द्रा\-नि\-द्रा\-भ्र\-म\-क्लमाः ॥}\ A; \skt{वलक्षयस्तृषा मूर्च्छा\-त\-न्द्रा\-नि\-द्रा\-भ्र\-म\-क्लमाः ।}\ K}}\skt{ ।}\pend

\pstart\versepar \ledsidenote{\passref{6.39.106}}\skt{उपद्रवाश्च श्वासाद्याः सम्भ\-व\-न्त्य\-ति\-लं\-घ\-नात् ॥ }\pend
\pstart\versepar \skt{दीपनं कफविच्छेदि }\edtext{\skt{वर्च्छो\-वा\-ता\-नु\-लो\-मनं}}{\Afootnote{H; \skt{पित्त\-वा\-ता\-नु\-लो\-म\-नम् ॥}\ A; \skt{वर्चो\-वा\-ता\-नु\-लो\-मनं ।}\ K}}\skt{ ।}\edtext{}{\lemma{\skt{सम्भ\-व\-न्त्य\-ति\-लं\-घ\-नात् ॥ दीपनं}}\Afootnote{\skt{संभ\-व\-न्त्य\-ति\-ल\-ङ्घ\-नात् । दीपनं}\ A; \skt{सांभ\-वं\-त्य\-ति\-लं\-घ\-नात् । दीपनं}\ K; \skt{सस\-म्भ\-व\-न्त्य\-ति\-लं\-घ\-नात् ॥ दीपनं}\ H}}\pend

\pstart\versepar \ledsidenote{\passref{6.39.107}}\edtext{\skt{कफ\-वा\-त\-ज्व\-रा\-र्त्तानां}}{\Afootnote{K H; \skt{कफ\-वा\-त\-ज्व\-रा\-र्तेभ्यो}\ A}}\skt{ }\edtext{\skt{हितमुष्णाम्बु}}{\lemma{\skt{उष्णाम्बु}}\Afootnote{A H; \skt{उष्णांवु}\ K}}\skt{ }\edtext{\skt{तृष्यतम्}}{\Afootnote{H; \skt{तृट्छिदम् ।}\ A; \skt{तृष्यतᳵ ।}\ K}}\skt{ ॥ }\pend
\pstart\versepar \skt{}\edtext{\skt{तद्धि}}{\Afootnote{K; \skt{तद्धि}\ A H}}\skt{ }\edtext{\skt{मार्द्दवकृद्दोष}}{\Afootnote{A H; \skt{मार्दवकृद्दोष}\ K}}\skt{ स्रोतसां शीतमन्यथा ।}\pend

\pstart\versepar \ledsidenote{\passref{6.39.108}}\edtext{\skt{सेव्यमानेत}}{\Afootnote{H; \skt{सेव्यमानेन}\ A K}}\skt{ तोयेन ज्वरः शीतेन वर्धते ॥ }\pend
\pstart\versepar \skt{}\edtext{\edtext{\skt{पित्तमद्यविषो}}{\lemma{\skt{पित्तम्}}\Afootnote{H; \skt{पित्त}\ A K}}\skt{ तिक्तकैः}}{\lemma{\skt{अद्यविषो तिक्तकैः}}\Afootnote{\skt{मद्यविषोत्थेषु शीतलं तिक्तकैः}\ A; \skt{मद्यविषोत्थेषु तिक्तकैः}\ K; \skt{अद्यविषोत्थेषु तिक्तकैः}\ H}}\skt{ }\edtext{\skt{शृत}}{\Afootnote{H; \skt{शृतम् ॥}\ A; \skt{श्रित}\ K}}\skt{ }\edtext{\skt{शीतलं}}{\Afootnote{K H; \emph{om.}\ A}}\skt{ ।}\pend

\pstart \ledsidenote{\passref{6.39.109}}\edtext{\skt{दीपनी}}{\Afootnote{K H; \skt{गाङ्गे\-य\-ना\-ग\-रो\-शी\-र\-प\-र्प\-टो\-दी\-च्य\-च\-न्दनैः । दीपनी}\ A}}\skt{ पाचनी लघ्वी }\edtext{\skt{ज्वरार्न्तानां}}{\Afootnote{\skt{ज्वरार्तानां}\ A H K}}\skt{ ज्वरापहा ॥}\pend

\pstart \ledsidenote{\passref{6.39.110}}\skt{अन्नकाले हिता पेया यथास्वं पाचनैः }\edtext{\skt{कृता}}{\Afootnote{H; \skt{कृता । बहुदोषस्य मन्दाग्नेः सप्तरात्रात्परं ज्वरे ॥}\ A; \skt{कृताः ।}\ K}}\skt{ ।}\pend

\pstart\versepar \ledsidenote{\passref{6.39.111}}\edtext{\skt{लंघ\-ना\-म्बो\-य\-वा\-गू\-भि\-र्यदा}}{\lemma{\skt{लंघ\-ना\-म्बो\-य\-वा\-गू\-भिर्}}\Afootnote{H; \skt{लङ्घ\-ना\-म्बु\-य\-वा\-गू\-भिर्}\ A; \skt{लांघ\-नां\-भो\-य\-वा\-गु\-भिर्}\ K}}\skt{ दोषो न पच्यते ॥ }\pend
\pstart\versepar \skt{तदा }\edtext{\skt{तत्मुखवैरस्य}}{\Afootnote{\skt{तं मुखवैरस्य}\ A; \skt{तन्मुखवैरस्य}\ K H}}\skt{ }\edtext{\skt{तृष्णा\-रो\-च\-क\-ना\-शनैः}}{\Afootnote{K H; \skt{त्ष्णारो\-च\-क\-ना\-शनैः ॥}\ A}}\skt{ ।}\pend

\pstart\versepar \ledsidenote{\passref{6.39.112}}\skt{कषायैः पाच\-नै\-र्हृ\-द्यै\-र्ज्व\-रघ्नैः समुपाचरेत् ॥ }\pend
\pstart\versepar \skt{पञ्च\-मू\-ली\-क\-षा\-यन्तु पाच\-न\-म्प\-व\-न\-ज्वरे ।}\pend

\pstart\versepar \ledsidenote{\passref{6.39.113}}\skt{सक्षौद्रम्पैत्तिके मुस्त\-क\-टु\-के\-न्द्र\-यवैः }\edtext{\skt{कृतम्}}{\Afootnote{A H; \skt{कृतः ।}\ K}}\skt{ ॥ }\pend
\pstart\versepar \skt{पिप्प\-ल्या\-दि\-क\-षा\-यन्तु }\edtext{\skt{पाचनं}}{\Afootnote{K H; \emph{om.}\ A}}\skt{ कफजे }\edtext{\skt{ज्वरे}}{\Afootnote{K H; \skt{परिपाचनम् ॥}\ A}}\skt{ ।}\pend

\pstart \ledsidenote{\passref{6.39.114}}\edtext{\skt{संषृ\-ष्ट\-दो\-षे\-षु\-हितं}}{\Afootnote{H; \skt{द्वन्द्वजेषु तु संसृष्टं द}\ A; \skt{संसृष्टदोषेषु हितं}\ K}}\skt{ }\edtext{\skt{संसृष्टमथ}}{\Afootnote{H; \skt{द्यादथ}\ A; \skt{संसृष्टमथ}\ K}}\skt{ }\edtext{\skt{पाचनम्}}{\Afootnote{K H; \skt{विवर्जयेत् । पीताम्बुर्लङ्घितो भुक्तोऽजीर्णी क्षीणः पिपासितः ॥}\ A}}\skt{ ॥}\pend

\pstart\versepar \ledsidenote{\passref{6.39.115}}\edtext{\skt{मृदौ}}{\Afootnote{H; \skt{(तीक्ष्णे ज्वरे गुरौ देहे विबद्धेषु मलेषु च । सामदोषं विजानीयाज्ज्वरं पक्वमतोऽयथा ॥) मृदौ}\ A; \skt{मृद्वौ}\ K}}\skt{ ज्वरे लघौ देहे प्रचलेषु मलेषु च । }\pend
\pstart\versepar \skt{}\edtext{\skt{पक्व}}{\Afootnote{K H; \skt{पक्वं}\ A}}\skt{ }\edtext{\skt{दोष\-म्वि\-जा\-नी\-या\-त्तदा}}{\lemma{\skt{विजानीयात्तदा}}\Afootnote{H; \skt{विजानीयाज्ज्वरे}\ A; \skt{विजानीया तदा}\ K}}\skt{ देयन्तदौषधं ॥}\pend

\pstart \ledsidenote{\passref{6.39.116}}\edtext{\edtext{\skt{दोष\-प्र\-कृ\-ति\-वै\-कृ\-त्या\-दे\-के\-षा\-म्प\-क्व\-ल\-क्ष\-णम्}}{\lemma{\skt{दोष\-प्र\-कृ\-ति\-वै\-कृ\-त्या\-दे\-के\-षाम्}}\Afootnote{A K\textsuperscript{pc} H; \skt{दोष\-प्र\-कृ\-ति\-वै\-कृ\-त्या\-दे\-षा\-के\-षाम्}\ K\textsuperscript{ac}}}}{\lemma{\skt{पक्वलक्षणम्}}\Afootnote{K H; \skt{पक्वलक्षणम् । हृदयोद्वेष्टनं तन्द्रा लाला\-स्रु\-ति\-र\-रो\-चकः ॥}\ A}}\skt{ ।}\pend

\pstart \ledsidenote{\passref{6.39.119}}\edtext{\skt{सप्तरात्परं}}{\Afootnote{H; \skt{लिङ्गै\-रे\-भि\-र्वि\-जा\-नी\-या\-ज्ज्व\-र\-मामं विचक्षणः । सप्तरात्रात्परं}\ A; \skt{सप्त\-रा\-त्रा\-त्प\-रङ्}\ K}}\skt{ }\edtext{\skt{केचिन्मन्यन्ते}}{\Afootnote{\skt{केचिन्मन्यन्ते}\ A H K}}\skt{ }\edtext{\skt{देयमौषधम्}}{\lemma{\skt{औषधम्}}\Afootnote{A H; \skt{औषधᳵ ।}\ K}}\skt{ ॥}\pend

\pstart\versepar \ledsidenote{\passref{6.39.120}}\skt{दशरात्रात्परं केचि\-द्दा\-त\-व्य\-मिति निश्चिताः । }\pend
\pstart\versepar \skt{पैत्तिके वा ज्वरे देय\-म\-ल्प\-का\-ल\-स\-मु\-त्थिते ॥}\pend

\pstart \ledsidenote{\passref{6.39.121ab}}\edtext{\edtext{\skt{सुचि\-र\-ज्व\-रि\-त\-स्यापि}}{\lemma{\skt{सुचिर्}}\Afootnote{H; \emph{om.}\ A K}}}{\lemma{\skt{अज्वरितस्यापि}}\Afootnote{H; \skt{अचि\-र\-ज्व\-रि\-त\-स्यापि}\ A K}}\skt{ देयं स्याद्दोषपाकतः ।}\pend

\pstart \ledsidenote{\passref{6.39.123cd}}\edtext{\skt{यस्य}}{\Afootnote{K H; \skt{यदा}\ A}}\skt{ }\edtext{\skt{दोषानुगाः}}{\Afootnote{K H; \skt{कोष्ठानुगाः}\ A}}\skt{ }\edtext{\skt{पक्वा\-वि\-शु\-द्ध\-स्रो\-त\-सो\-मलाः}}{\Afootnote{H; \skt{पक्वा विबद्धाः स्रोतसां मलाः ॥}\ A; \skt{पक्व\-वि\-श\-द्ध\-स्रो\-तसौ मलाः ।}\ K}}\skt{ ॥}\pend

\pstart \ledsidenote{\passref{6.39.124}}\edtext{\skt{उत्तर अचि\-र\-ज्व\-रि\-त\-स्यापि}}{\Afootnote{\skt{अचि\-र\-ज्व\-रि\-त\-स्यापि}\ A K; \skt{अचि\-र\-ज्व\-रि\-त\-स्यापि…}\ H}}\skt{ }\edtext{\skt{तस्य}}{\Afootnote{K H; \skt{तदा}\ A}}\skt{ }\edtext{\skt{दद्या\-द्वि\-रे\-च\-नम्}}{\Afootnote{H; \skt{दद्या\-द्वि\-रे\-च\-नम् । पक्वो ह्यनिर्हृतो दोषो देहे तिष्ठ\-न्म\-हा\-त्य\-यम् ॥}\ A; \skt{दद्याद्विरेचनᳵ ।}\ K}}\skt{ ।}\pend

\pstart \ledsidenote{\passref{6.39.121.add1}}\edtext{\skt{लाला\-प्र\-से\-क\-हृ\-ल्ला\-स\-त\-म\-का\-षु\-द्ध्य\-रो\-चकाः}}{\Afootnote{\emph{om.}\ K\textsuperscript{ac}; \skt{लाला\-प्र\-से\-क\-हृ\-ल्ला\-स\-त\-म\-का\-शु\-द्ध\-रो\-चकाः ।}\ K\textsuperscript{pc}; \skt{लाला\-प्र\-से\-क\-हृ\-ल्ला\-स\-त\-म\-का\-षु\-द्ध्य\-रो\-चकाः ॥}\ H}}\skt{ ॥}\pend

\pstart\versepar \ledsidenote{\passref{6.39.121.add2}}\edtext{\skt{तन्द्रा\-ल\-स्या\-वि\-पा\-कस्य}}{\Afootnote{H; \emph{om.}\ K\textsuperscript{ac}; \skt{तन्द्रा\-ल\-स्या\-वि\-पा\-कास्य}\ K\textsuperscript{pc}}}\skt{ }\edtext{\skt{वैरस्यं}}{\Afootnote{H; \emph{om.}\ K\textsuperscript{ac}; \skt{वैरस्याङ्}\ K\textsuperscript{pc}}}\skt{ }\edtext{\skt{गुरुगात्रता}}{\Afootnote{K\textsuperscript{pc} H; \emph{om.}\ K\textsuperscript{ac}}}\skt{ । }\pend
\pstart\versepar \skt{}\edtext{\skt{मूत्रत्वं}}{\Afootnote{\emph{om.}\ K\textsuperscript{ac}; \skt{क्षुन्नाशो वहुमूत्रत्वं}\ K\textsuperscript{pc}; \skt{क्ष्मन्ना\-शो\-व\-मू\-त्रत्वं}\ H}}\skt{ }\edtext{\skt{स्तब्धता}}{\Afootnote{H; \emph{om.}\ K\textsuperscript{ac}; \skt{स्तव्धता}\ K\textsuperscript{pc}}}\skt{ }\edtext{\skt{बलवान्ज्वरः}}{\Afootnote{H; \emph{om.}\ K\textsuperscript{ac}; \skt{वलवांज्वरः ।}\ K\textsuperscript{pc}}}\skt{ ॥}\pend

\pstart \ledsidenote{\passref{6.39.121.add3}}\edtext{\skt{आमज्वरस्य}}{\Afootnote{K\textsuperscript{pc} H; \emph{om.}\ K\textsuperscript{ac}}}\skt{ }\edtext{\skt{लिङ्गानि न}}{\Afootnote{K\textsuperscript{pc} H; \emph{om.}\ K\textsuperscript{ac}}}\skt{ }\edtext{\skt{दध्यात्तत्र}}{\Afootnote{H; \emph{om.}\ K\textsuperscript{ac}; \skt{दद्यात्तत्र}\ K\textsuperscript{pc}}}\skt{ }\edtext{\skt{भेषजं}}{\Afootnote{K\textsuperscript{pc} H; \emph{om.}\ K\textsuperscript{ac}}}\skt{ ।}\pend

\pstart \ledsidenote{\passref{6.39.121cd}}\edtext{\skt{भेषजं }\edtext{\skt{ह्यामदोषस्य}}{\lemma{\skt{आमदोषस्य}}\Afootnote{A K\textsuperscript{pc} H; \emph{om.}\ K\textsuperscript{ac}}}}{\lemma{\skt{भेषजं ह्य्}}\Afootnote{A K\textsuperscript{pc} H; \emph{om.}\ K\textsuperscript{ac}}}\skt{ }\edtext{\skt{भूयो}}{\Afootnote{A K\textsuperscript{pc} H; \emph{om.}\ K\textsuperscript{ac}}}\skt{ }\edtext{\skt{ज्वलयति}}{\Afootnote{A K\textsuperscript{pc} H; \emph{om.}\ K\textsuperscript{ac}}}\skt{ }\edtext{\skt{ज्वरम्}}{\Afootnote{A K\textsuperscript{pc} H; \emph{om.}\ K\textsuperscript{ac}}}\skt{ ॥}\pend

\pstart \ledsidenote{\passref{6.39.122}}\edtext{\skt{शमनीयञ्च}}{\Afootnote{\skt{शोधनं शमनीयं तु}\ A; \emph{om.}\ K\textsuperscript{ac}; \skt{शोधनं शमनीयञ्च}\ K\textsuperscript{pc}; \skt{…शमनीयञ्च}\ H}}\skt{ }\edtext{\skt{करोति}}{\Afootnote{A K\textsuperscript{pc} H; \emph{om.}\ K\textsuperscript{ac}}}\skt{ }\edtext{\skt{विषमज्वरम्}}{\Afootnote{H; \skt{विषमज्वरम् । च्यवमानं ज्वरो\-त्क्लि\-ष्ट\-मु\-पे\-क्षेत मलं सदा ॥}\ A; \emph{om.}\ K\textsuperscript{ac}; \skt{विषमज्वरᳵ ।}\ K\textsuperscript{pc}}}\skt{ ।}\pend

\pstart \ledsidenote{\passref{6.39.124cd}}\edtext{\skt{पक्वो}}{\Afootnote{H; \emph{om.}\ K\textsuperscript{ac}; \skt{पक्वो ह्य्}\ K\textsuperscript{pc}}}\skt{ }\edtext{\skt{व्यतिहृतो}}{\Afootnote{H; \emph{om.}\ K\textsuperscript{ac}; \skt{अनिर्हृतो}\ K\textsuperscript{pc}}}\skt{ }\edtext{\skt{दोषो}}{\Afootnote{K\textsuperscript{pc} H; \emph{om.}\ K\textsuperscript{ac}}}\skt{ }\edtext{\skt{देहे}}{\Afootnote{K\textsuperscript{pc} H; \emph{om.}\ K\textsuperscript{ac}}}\skt{ }\edtext{\skt{तिष्ठन्म}}{\Afootnote{H; \emph{om.}\ K\textsuperscript{ac}; \skt{तिष्ठम्म}\ K\textsuperscript{pc}}}\skt{ }\edtext{\skt{दात्ययः ॥ प्त}}{\Afootnote{\emph{om.}\ K\textsuperscript{ac}; \skt{हात्ययᳵ ।}\ K\textsuperscript{pc}; \skt{दात्ययः ॥}\ H}}\skt{ }\edtext{\skt{भ्र}}{\Afootnote{\emph{om.}\ K H}}\pend

\pstart \ledsidenote{\passref{6.39.125}}\edtext{\skt{विषमं}}{\Afootnote{A H; \skt{विषमम्}\ K}}\skt{ वा ज्वरं }\edtext{\skt{कुर्या\-त्म\-ल\-व्य\-प\-द\-मेव}}{\lemma{\skt{कुर्या\-त्म\-ल\-व्य\-प\-दम्}}\Afootnote{H; \skt{कुर्या\-द्ब\-ल\-व्या\-प\-दम्}\ A; \skt{कुर्या\-त्म\-ल\-व्या\-प\-दम्}\ K}}\skt{ }\edtext{\skt{वा}}{\Afootnote{K H; \skt{च । तस्मान्निर्हरणं कार्यं दोषाणां वमनादिभिः ॥}\ A}}\skt{ ।}\pend

\pstart\versepar \ledsidenote{\passref{6.39.126}}\edtext{\skt{प्राक्कर्म्म}}{\Afootnote{K H; \skt{प्रक्कर्म}\ A}}\skt{ वमनञ्चास्य कार्य\-मा\-स्था\-प\-न\-न्तथा ॥ }\pend
\pstart\versepar \skt{}\edtext{\edtext{\skt{विरे\-च\-न\-न्त\-त\-श्चापि}}{\lemma{\skt{ततश्}}\Afootnote{K H; \skt{तथा}\ A}}}{\lemma{\skt{चापि}}\Afootnote{H; \skt{कुर्या}\ A; \skt{यस्मा}\ K}}\skt{ }\edtext{\skt{शिरसश्च}}{\Afootnote{H; \skt{च्छिरसश्च}\ A K}}\skt{ विरेचनम् ।}\pend

\pstart \ledsidenote{\passref{6.39.127ab}}\edtext{\skt{क्रमशो}}{\Afootnote{A K\textsuperscript{pc} H; \skt{क्रशो}\ K\textsuperscript{ac}}}\skt{ }\edtext{\skt{बलिने}}{\Afootnote{A H; \skt{वलिने}\ K}}\skt{ देयं वमनं }\edtext{\skt{श्लेष्मिके}}{\Afootnote{K H; \skt{श्लैष्मिके}\ A}}\skt{ ज्वरे ॥}\pend

\pstart \ledsidenote{\passref{6.39.128ab}}\edtext{\skt{सरुजो}}{\Afootnote{\skt{सरुजे}\ A K; \skt{सरु}\ H}}\skt{ }\edtext{\skt{निलजे}}{\Afootnote{K; \skt{ऽनिलजे}\ A; \skt{जोनिलजे}\ H}}\skt{ कार्यं सोदावर्त्ते }\edtext{\skt{निरूहनम्}}{\Afootnote{H; \skt{निरूहणम्}\ A K}}\skt{ ।}\pend

\pstart \ledsidenote{\passref{6.39.127cd}}\edtext{\skt{पित्तप्राये}}{\Afootnote{A H; \skt{पत्तिप्राय}\ K}}\skt{ विरेकः }\edtext{\skt{स्यात्कार्यं}}{\Afootnote{K; \emph{om.}\ A; \skt{स्या…त्कार्यं}\ H}}\skt{ }\edtext{\skt{शीर्ष\-वि\-रे\-च\-नम्}}{\Afootnote{K H; \skt{तु कार्यः प्रशिथिलाशये ॥}\ A}}\skt{ ॥}\pend

\pstart \ledsidenote{\passref{6.39.128cd}}\skt{कटी\-पृ\-ष्ठ\-ग्र\-हा\-र्त्तस्य }\edtext{\edtext{\skt{दीप्ता\-ग्ने\-र\-नु\-वा\-स\-नम्}}{\lemma{\skt{दीप्ताग्नेर्}}\Afootnote{A; \skt{दीप्ताग्ने}\ K; \skt{…दीप्ताग्नेर्}\ H}}}{\lemma{\skt{अनुवासनम्}}\Afootnote{A H; \skt{रनुवासनं ।}\ K}}\skt{ ।}\pend

\pstart \ledsidenote{\passref{6.39.134ab}}\skt{कृशं चैवाल्पदोषञ्च शम\-नी\-यै\-रु\-पा\-च\-रेत् ॥}\pend

\pstart \ledsidenote{\passref{6.39.129cd}}\skt{कफाभिपन्ने शिरसि कार्यं }\edtext{\skt{शीर्ष\-वि\-रे\-च\-नम्}}{\Afootnote{H; \skt{मूर्ध\-वि\-रे\-च\-नम् ॥}\ A; \skt{शीर्षविरेचनᳵ ।}\ K}}\skt{ ।}\pend

\pstart \ledsidenote{\passref{6.39.129ab}}\edtext{\skt{शिरो\-गौ\-र\-व\-शू\-ल\-घ्न\-मि\-न्द्रि\-य\-प्र\-ति\-बो\-ध\-नम्}}{\Afootnote{A; \skt{शिरो\-गौ\-र\-व\-शू\-ल\-घ्न\-मि\-न्द्रि\-य\-प्र\-ति\-बो\-धनं ।}\ K; \skt{शिरो\-गौ\-र\-व\-शू\-ल\-घ्न…\-मि\-न्द्रि\-य\-प्र\-ति\-बो\-ध\-नम् ॥}\ H}}\skt{ ॥}\pend

\pstart \ledsidenote{\passref{6.39.134cd}}\edtext{\edtext{\skt{उप\-वा\-सै\-र्ब्ब\-ल\-स्थन्तु}}{\lemma{\skt{उपवासैर्}}\Afootnote{K H; \skt{उपवासै}\ A}}}{\lemma{\skt{ब्बलस्थन्तु}}\Afootnote{H; \skt{र्बलस्थं तु}\ A; \skt{वलस्थन्तु}\ K}}\skt{ ज्वरे }\edtext{\skt{सन्तर्प्पणान्थिते}}{\Afootnote{H; \skt{सन्तर्पणोत्थिते}\ A K}}\skt{ ॥}\pend

\pstart \ledsidenote{\passref{6.39.135}}\edtext{\skt{क्लिन्नां}}{\Afootnote{A H; \skt{क्लिन्ना}\ K}}\skt{ यवागूं मन्दाग्निं }\edtext{\skt{पिपासार्त्तञ्च}}{\Afootnote{K H; \skt{तृषार्तं}\ A}}\skt{ }\edtext{\skt{पाययेत्}}{\Afootnote{K H; \skt{पाययेन्नरम् । तृट्छ\-र्दि\-दा\-ह\-घ\-र्मार्तं मद्यपं लाजतर्पणम् ॥}\ A}}\skt{ ।}\pend

\pstart \ledsidenote{\passref{6.39.136a}}\skt{श्रमो\-प\-वा\-सा\-नि\-लजे }\edtext{\skt{हितोनित्यं}}{\Afootnote{\skt{हितो नित्यं}\ H}}\skt{ रसोदनः ॥}\pend

\pstart \ledsidenote{\passref{6.39.137ab}}\edtext{\edtext{\skt{मुद्ग\-यू\-षो\-द\-न\-ञ्चापि}}{\lemma{\skt{मुद्ग\-यू\-षो\-द\-नञ्}}\Afootnote{H; \skt{दीप्ताग्निं भोजयेत्प्राज्ञो नरं मांस\-र\-सौ\-द\-नम् ।}\ A; \skt{मुद्ग\-यू\-षो\-द\-नश्}\ K}}}{\lemma{\skt{चापि}}\Afootnote{K; \emph{om.}\ A; \skt{चापि …}\ H}}\skt{ }\edtext{\skt{देयः}}{\Afootnote{K H; \emph{om.}\ A}}\skt{ }\edtext{\skt{कफसमुत्थिते}}{\Afootnote{K H; \emph{om.}\ A}}\skt{ ।}\pend

\pstart \ledsidenote{\passref{6.39.138}}\skt{स एव सितया युक्तः }\edtext{\skt{शीतः}}{\Afootnote{A H; \skt{शीत}\ K}}\skt{ पित्तज्वरे }\edtext{\skt{हितः}}{\Afootnote{K H; \skt{हितः । दाडि\-मा\-म\-ल\-मु\-द्गानां यूष\-श्चा\-नि\-ल\-पै\-त्तिके ॥}\ A}}\skt{ ॥}\pend

\pstart \ledsidenote{\passref{6.39.137cd}}\edtext{\skt{मुद्गा\-म\-ल\-क\-यू\-षस्तु}}{\Afootnote{K; \skt{मुद्ग\-यू\-षौ\-द\-न\-श्चापि हितः कफसमुत्थिते ॥}\ A; \skt{मुद्गा\-म\-ल\-क\-यू\-षस्तु}\ H}}\skt{ }\edtext{\skt{वातपित्तात्मके}}{\Afootnote{H; \emph{om.}\ A; \skt{वात्तपित्तात्मके}\ K}}\skt{ }\edtext{\skt{ज्वरे}}{\Afootnote{K H; \emph{om.}\ A}}\skt{ ।}\pend

\pstart\versepar \ledsidenote{\passref{6.39.139}}\edtext{\skt{ह्रस्व\-मू\-ल\-क\-यू\-षस्तु}}{\Afootnote{A H; \skt{ह्रस्व\-मू\-ल\-क\-षा\-यस्तु}\ K}}\skt{ }\edtext{\skt{वातश्लेष्मातके}}{\Afootnote{H; \skt{वातश्लेष्माधिके}\ A; \skt{वातश्लेष्मात्मके}\ K}}\skt{ }\edtext{\skt{ज्वरे}}{\Afootnote{K H; \skt{हि}\ A}}\skt{ ॥ }\pend
\pstart\versepar \skt{}\edtext{\skt{निम्ब\-को\-ल\-क\-यू\-षस्तु}}{\Afootnote{H; \skt{तः । पटो\-ल\-नि\-म्ब\-यू\-षस्तु}\ A; \skt{निम्व\-को\-ल\-क\-यू\-षस्तु}\ K}}\skt{ }\edtext{\skt{हितः}}{\Afootnote{H; \skt{पथ्यः}\ A; \skt{हित}\ K}}\skt{ }\edtext{\skt{पित्तकफात्मके}}{\Afootnote{A H; \skt{ञ्पित्तकफात्मके ।}\ K}}\skt{ ।}\pend

\pstart\versepar \ledsidenote{\passref{6.39.140}}\edtext{\skt{दाहवम्यर्द्दितं}}{\Afootnote{H; \skt{दाहच्छर्दियुतं}\ A; \skt{दाहवर्म्म्यर्दितं}\ K}}\skt{ क्षामं निरन्नं }\edtext{\skt{तृष्ण\-आ\-या\-न्वि\-तम्}}{\Afootnote{H\textsuperscript{ac}; \skt{तृष्णयार्दितम् ।}\ A; \skt{तृष्णयान्वितᳵ ।}\ K; \skt{तृष्णयान्वितम् ॥}\ H\textsuperscript{pc}}}\skt{ ॥ }\pend
\pstart\versepar \skt{}\edtext{\skt{शर्क्क\-रा\-म\-धु\-सं\-युक्तं}}{\Afootnote{K H; \skt{सिताक्षौद्रयुतं}\ A}}\skt{ }\edtext{\skt{पाय\-ये\-ल्ला\-ज\-त\-र्पणं}}{\Afootnote{H; \skt{लाजतर्पणं पाययेत च ॥}\ A; \skt{पाय\-ये\-ल्ला\-ज\-त्र\-र्पणं ।}\ K}}\skt{ ॥}\pend

\pstart\versepar \ledsidenote{\passref{6.39.141}}\skt{कफ\-पि\-त्त\-प\-री\-तस्य }\edtext{\skt{ग्रीष्मेऽसृक्पित्त}}{\lemma{\skt{ऽसृक्पित्त}}\Afootnote{H; \skt{ऽसृक्पित्ति}\ A; \skt{सृक्पित्त}\ K}}\skt{ }\edtext{\skt{ऊर्द्धगे}}{\Afootnote{H; \skt{नस्तथा ।}\ A; \skt{ऊर्ध्वगे ॥}\ K}}\skt{ ॥ }\pend
\pstart\versepar \skt{}\edtext{\skt{मद्यनित्यस्य न}}{\Afootnote{K H; \skt{मद्यनित्यस्य न}\ A}}\skt{ हिता }\edtext{\edtext{\skt{यवा\-गू\-भि\-रु\-पा\-च\-रेत्}}{\lemma{\skt{यवागूभिर्}}\Afootnote{H; \skt{यवागूस्त}\ A; \skt{यवागूस्तम्}\ K}}}{\lemma{\skt{उपाचरेत्}}\Afootnote{K H; \skt{मुपाचरेत् ॥}\ A}}\skt{ ।}\pend

\pstart \ledsidenote{\passref{6.39.142}}\edtext{\skt{मद्यं वा}}{\Afootnote{\skt{यूषै\-र\-म्लै\-र\-न\-म्लैर्वा जाङ्गलैश्च रसैर्हितैः । मद्यं}\ A; \skt{यूषै\-र\-म्लै\-र\-न\-म्लैर्वा जाङ्गलैर्वा रसैर्हितं । मद्यं वा}\ K; \skt{यूषै\-र\-म्लै\-र\-न\-म्लैर्व्वा जाङ्गलैर्व्वा रसैर्हितैः ॥ मद्यं वा}\ H}}\skt{ }\edtext{\skt{मद्यसात्म्याय}}{\Afootnote{K H; \emph{om.}\ A}}\skt{ }\edtext{\skt{सामाहाय}}{\Afootnote{H; \emph{om.}\ A; \skt{सानाहाय}\ K}}\skt{ }\edtext{\skt{प्रदापयेत्}}{\Afootnote{K H; \skt{पुराणं मन्दा\-ग्ने\-र्य\-वा\-न्नो\-प\-हितं हितम् ॥}\ A}}\skt{ ।}\pend

\pstart \ledsidenote{\passref{6.39.143}}\edtext{\skt{सव्यो\-ष\-म्वि\-त\-रे\-त्तात्र कफा\-रो\-च\-क\-पी\-डिते}}{\lemma{\skt{वितरेत्तात्र कफा\-रो\-च\-क\-पी\-डिते}}\Afootnote{\skt{वितरेत्तक्रं कफा\-रो\-च\-क\-पी\-डिते । कृशोऽल्पदोषो दीनश्च नरो जीर्णज्वरार्दितः ॥}\ A; \skt{वितरे तक्र\-ङ्क\-फा\-रो\-च\-क\-पी\-डिते ॥}\ K; \skt{वितरेत्तात्र… कफा\-रो\-च\-क\-पी\-डिते ॥}\ H}}\skt{ ॥}\pend

\pstart\versepar \ledsidenote{\passref{6.39.144}}\edtext{\skt{वद्ध\-प्र\-द्यु\-त\-दो\-षस्तु}}{\Afootnote{K; \skt{विबद्धः सृष्टदोषश्}\ A; \skt{वद्धप्र र्युत दोषस्तु}\ H}}\skt{ }\edtext{\skt{नरो}}{\Afootnote{K H; \skt{च रू}\ A}}\skt{ }\edtext{\skt{जीर्ण्णज्वरान्वितः}}{\Afootnote{K H; \skt{क्षः पित्तानिल}\ A}}\skt{ । }\pend
\pstart\versepar \skt{}\edtext{\skt{पिपासार्तः}}{\Afootnote{K H; \skt{ज्वरी । पिपासार्तः}\ A}}\skt{ सदाहो वा पयसा स सुखी भवेत् ॥}\pend

\pstart\versepar \ledsidenote{\passref{6.39.145}}\edtext{\skt{पिपाव}}{\Afootnote{H; \skt{त}\ A; \skt{तदेव}\ K}}\skt{ }\edtext{\skt{तु}}{\Afootnote{K H; \skt{दे}\ A}}\skt{ }\edtext{\skt{पयः}}{\Afootnote{K H; \skt{व}\ A}}\skt{ }\edtext{\edtext{\skt{पीतन्तरुणे}}{\lemma{\skt{पीतन्}}\Afootnote{K H; \emph{om.}\ A}}}{\lemma{\skt{तरुणे}}\Afootnote{A H; \skt{तरुणेर्}\ K}}\skt{ }\edtext{\skt{हन्ति}}{\Afootnote{H; \skt{पीतं विषवद्धन्ति}\ A; \skt{हन्ति माणवᳵ । ज्व}\ K\textsuperscript{ac}; \skt{हन्ति माणवᳵ । ॥}\ K\textsuperscript{pc}}}\skt{ }\edtext{\skt{मानवम्}}{\Afootnote{A H; \skt{रे॥}\ K\textsuperscript{ac}; \emph{om.}\ K\textsuperscript{pc}}}\skt{ । }\pend
\pstart\versepar \skt{सर्व्वज्वरेषु }\edtext{\edtext{\skt{सुल\-घु\-र्म्मा\-त्रा\-व\-द्भो\-जनं}}{\lemma{\skt{सुलघुर्}}\Afootnote{K H; \skt{सुलघु}\ A}}}{\lemma{\skt{म्मात्रावद्भोजनं}}\Afootnote{K; \skt{मात्रावद्भोजनं}\ A H}}\skt{ }\edtext{\skt{हि तम्}}{\Afootnote{A H; \skt{हितᳵ ।}\ K}}\skt{ ॥}\pend

\pstart\versepar \ledsidenote{\passref{6.39.146}}\edtext{\skt{वेगापाये न्यथा}}{\Afootnote{K H; \skt{वेगापायेऽन्यथा}\ A}}\skt{ }\edtext{\skt{तद्धि}}{\Afootnote{K; \skt{तद्धि}\ A H}}\skt{ }\edtext{\skt{ज्वर\-वे\-गा\-भि\-व\-र्द्ध\-नम्}}{\Afootnote{A H; \skt{ज्वर\-वे\-गा\-भि\-व\-र्धनᳵ ।}\ K}}\skt{ । }\pend
\pstart\versepar \skt{ज्वरितो }\edtext{\edtext{\edtext{\edtext{\skt{हित\-म\-श्नी\-या\-द्य\-द्य\-प्य\-स्या\-रु\-चि\-र्भ्भ\-वेत्}}{\lemma{\skt{अश्नीयाद्यद्य्}}\Afootnote{A K; \skt{अश्नीयात्य}\ H}}}{\lemma{\skt{अप्य्}}\Afootnote{A K; \skt{दपि}\ H}}}{\lemma{\skt{अस्यारुचिर्}}\Afootnote{A K\textsuperscript{pc} H; \skt{अरुचिर्}\ K\textsuperscript{ac}}}}{\lemma{\skt{भ्भवेत्}}\Afootnote{\skt{भवेत्}\ A H; \skt{भवेत ।}\ K}}\skt{ ॥}\pend

\pstart \ledsidenote{\passref{6.39.147}}\skt{अन्नका ले ह्यभुञ्जानः क्षीयते }\edtext{\skt{म्रियतेऽपि वा}}{\lemma{\skt{ऽपि वा}}\Afootnote{H; \skt{ऽथवा । स क्षीणः कृच्छ्रतां याति यात्य\-सा\-ध्य\-त्व\-मेव च ॥}\ A; \skt{पि वा ।}\ K}}\skt{ ।}\pend

\pstart \ledsidenote{\passref{6.39.148}}\edtext{\skt{गुर्व\-भि\-ष्य\-न्दि\-काले च}}{\Afootnote{K H; \skt{तस्माद्रक्षेद्बलं पुंसां बले सति हि जीवितम् । गुर्व\-भि\-ष्य\-न्द्य\-काले च}\ A}}\skt{ ज्वरी }\edtext{\skt{नाद्यात्कथञ्चन}}{\Afootnote{A H; \skt{नाद्यात्कथञ्चनः ।}\ K}}\skt{ ॥}\pend

\pstart\versepar \ledsidenote{\passref{6.39.149}}\skt{न तु तस्याहितं }\edtext{\skt{भुक्तमायुषे वा}}{\Afootnote{\skt{भुक्तमायुषे वा}\ A K H}}\skt{ सुखाय वा । }\pend
\pstart\versepar \skt{}\edtext{\skt{सन्त\-त\-म्वि\-ष\-म\-ञ्चापि}}{\lemma{\skt{चापि}}\Afootnote{K H; \skt{वापि}\ A}}\skt{ क्षीणस्य सुचि रोत्थितम् ॥}\pend

\pstart\versepar \ledsidenote{\passref{6.39.150}}\edtext{\skt{ज्वरो}}{\Afootnote{H; \skt{ज्वरं}\ A; \skt{ज्वरे}\ K}}\skt{ }\edtext{\edtext{\skt{सभोजनम्पथ्यं}}{\lemma{\skt{सभोजनम्}}\Afootnote{H; \skt{संभोजनैः}\ A; \skt{संभोजनं}\ K}}}{\lemma{\skt{पथ्यं}}\Afootnote{K H; \skt{पथ्यै}\ A}}\skt{ }\edtext{\skt{लघुभिः}}{\Afootnote{K H; \skt{र्लघुभिः}\ A}}\skt{ }\edtext{\skt{समुपक्रमैः}}{\Afootnote{\skt{समुपाचरे}\ A; \skt{समुपक्रमैः ।}\ K H}}\skt{ । }\pend
\pstart\versepar \skt{}\edtext{\edtext{\skt{मुद्गा\-न्म\-सू\-रा\-न्श्च\-ण\-का\-न्कु\-ल\-त्था\-न्स\-म\-कु\-ष्ठ\-कान्}}{\lemma{\skt{मुद्गान्मसूरान्श्}}\Afootnote{H; \skt{त् । मुद्गान्मसूरांश्}\ A; \skt{मुद्गान्मसूराश्}\ K}}}{\lemma{\skt{चण\-का\-न्कु\-ल\-त्था\-न्स\-म\-कु\-ष्ठ\-कान्}}\Afootnote{A H; \skt{चणकां कुलत्थां समकुष्टकां ॥}\ K}}\skt{ । }\pend

\pstart\versepar \ledsidenote{\passref{6.39.151}}\skt{आहारकाले }\edtext{\skt{यूषार्थं}}{\Afootnote{A H; \skt{यूषार्थे}\ K}}\skt{ ज्वरिताय प्रदापयेत् ॥ }\pend
\pstart\versepar \skt{पटोलपत्रं }\edtext{\skt{वार्त्ताकु}}{\Afootnote{K H; \skt{वार्ताकं}\ A}}\skt{ }\edtext{\skt{कुक्कुलं}}{\Afootnote{H; \skt{कठिल्लं}\ A; \skt{कुकूलम्}\ K\textsuperscript{ac}; \skt{कुक्कूलम्}\ K\textsuperscript{pc}}}\skt{ }\edtext{\skt{पापचेलकं}}{\Afootnote{K H; \skt{पापचैलिकम् ॥}\ A}}\skt{ ॥}\edtext{}{\lemma{\skt{प्रदापयेत् ॥ पटोलपत्रं}}\Afootnote{A H; \skt{प्रदापयेत् । पटोलपत्र}\ K}}\pend

\pstart\versepar \ledsidenote{\passref{6.39.152}}\edtext{\skt{कर्का\-ट\-क\-म्प\-र्प्प\-टकं}}{\lemma{\skt{कर्काटकम्}}\Afootnote{A; \skt{कर्कोटकं}\ K H}}\skt{ }\edtext{\skt{गोजिह्वा}}{\Afootnote{H; \skt{गोजिह्वां}\ A K}}\skt{ }\edtext{\skt{बालमूलकम्}}{\Afootnote{A K; \skt{वालमूलकम् ।}\ H}}\skt{ । }\pend
\pstart\versepar \skt{}\edtext{\skt{पत्रं}}{\Afootnote{A H; \skt{पत्र}\ K}}\skt{ }\edtext{\skt{गुडूच्याः}}{\Afootnote{A H; \skt{ङगुडूच्याः}\ K}}\skt{ }\edtext{\skt{शाकार्थे}}{\Afootnote{A H; \skt{शकार्थे}\ K}}\skt{ }\edtext{\skt{ज्वरिताय}}{\Afootnote{K; \skt{ज्वरितानां}\ A; \skt{ज्वरितोय}\ H}}\skt{ }\edtext{\skt{प्रदापयेत्}}{\Afootnote{A H; \skt{प्रदापयेत ।}\ K}}\skt{ ॥}\pend

\pstart\versepar \ledsidenote{\passref{6.39.153}}\edtext{\skt{लावा\-न्क\-पि\-ञ्ज\-ला\-ने\-णा\-न्पृ\-ष\-ता\-न्क\-र\-भा\-ञ्छ\-शान्}}{\lemma{\skt{लावा\-न्क\-पि\-ञ्ज\-ला\-ने\-णा\-न्पृ\-ष\-ता\-न्क\-र\-भाञ्}}\Afootnote{H; \skt{लावा\-न्क\-पि\-ञ्ज\-ला\-ने\-णा\-न्पृ\-ष\-ता\-ञ्छ\-र\-भाञ्}\ A; \skt{लावा\-ङ्क\-पिं\-ज\-ला\-नेणां पृषतां कार}\ K}}\skt{ । }\pend
\pstart\versepar \skt{उत्त\-र\-का\-ल\-पु\-च्छा\-न्कु\-रङ्गा तथैव }\edtext{\skt{मृगमातृकान्}}{\Afootnote{A H; \skt{मृगमातृकां ॥}\ K}}\skt{ ॥}\edtext{}{\lemma{\skt{छशान् । उत्त\-र\-का\-ल\-पु\-च्छा\-न्कु\-रङ्गा}}\Afootnote{\skt{छशान् । काल\-पु\-च्छा\-न्कु\-र\-ङ्गांश्च}\ A; \skt{भाच्छसां । कालपुछां कुरङ्गाश्च}\ K; \skt{छशान् । काल\-पु\-च्छा\-न्कु\-र\-ङ्गाञ्श्च}\ H}}\pend

\pstart\versepar \ledsidenote{\passref{6.39.154}}\skt{मांसार्थे }\edtext{\skt{मांससात्म्याय}}{\Afootnote{K H; \skt{मांससात्म्यानां}\ A}}\skt{ }\edtext{\skt{ज्वरिताय}}{\Afootnote{K H; \skt{ज्वरितानां}\ A}}\skt{ प्रदापयेत् । }\pend
\pstart\versepar \skt{}\edtext{\skt{सार\-स\-क्रौ\-ञ्च\-शि\-खि\-न\-स्तथा}}{\lemma{\skt{तथा}}\Afootnote{K H; \emph{om.}\ A}}\skt{ }\edtext{\skt{कुक्कुटतित्तिराः}}{\Afootnote{K H; \skt{कुक्कु\-टां\-स्ति\-त्ति\-रां\-स्(\-ति\-त्ति\-रीं)\-तथा ॥}\ A}}\skt{ ॥}\edtext{}{\lemma{\skt{प्रदापयेत् । सार\-स\-क्रौ\-ञ्च\-शि\-खि\-नस्}}\Afootnote{A H; \skt{प्रदापयेत् । सार\-स\-क्रो\-ञ्च\-शि\-खि\-नस्}\ K}}\pend

\pstart\versepar \ledsidenote{\passref{6.39.155}}\skt{गुरूष्णत्वान्न }\edtext{\skt{शस्यन्त}}{\Afootnote{K H; \skt{शंसन्ति}\ A}}\skt{ }\edtext{\skt{इति}}{\Afootnote{K H; \skt{ज्वरे}\ A}}\skt{ }\edtext{\skt{केचिद्व्यवस्थिताः}}{\Afootnote{H; \skt{केचिच्चिकित्सकाः ।}\ A; \skt{केचित्व्यस्थिताः ।}\ K}}\skt{ । }\pend
\pstart\versepar \skt{ज्वरितानां }\edtext{\skt{प्रकोपन्तु}}{\Afootnote{A H; \skt{प्रकोपे तु}\ K}}\skt{ यदा याति समीरणः ॥}\pend

\pstart\versepar \ledsidenote{\passref{6.39.156}}\edtext{\skt{तदा}}{\Afootnote{K H; \skt{तदै}\ A}}\skt{ }\edtext{\skt{तेपि हि}}{\Afootnote{K H; \skt{तेऽपि हि}\ A}}\skt{ शस्यन्ते }\edtext{\skt{मात्रा\-का\-लो\-प\-पा\-दिता}}{\Afootnote{H; \skt{मात्रा\-का\-लो\-प\-पा\-दिताः ।}\ A K}}\skt{ । }\pend
\pstart\versepar \skt{}\edtext{\skt{परी\-षे\-का\-न्प्र\-देहा}}{\Afootnote{\skt{परि\-षे\-का\-न्प्र\-दे\-हांश्च}\ A; \skt{परीषेका प्रदेहांश्च}\ K; \skt{परी\-षे\-का\-न्प्र\-दे\-हाञ्श्च}\ H}}\skt{ }\edtext{\skt{स्नेहा\-न्सं\-शो\-ध\-नानि}}{\Afootnote{A H; \skt{स्नेहां संशोधनानि}\ K}}\skt{ }\edtext{\skt{च}}{\Afootnote{A K H\textsuperscript{pc}; \emph{om.}\ H\textsuperscript{ac}}}\skt{ ॥}\pend

\pstart \ledsidenote{\passref{6.39.156a}}\skt{दिवास्वप्नं }\edtext{\skt{व्यवा प्त ञ्यञ्च}}{\Afootnote{\skt{व्यवा यञ्च}\ H}}\skt{ व्यायामं शिशिरञ्जलम् ।}\pend

\pstart\versepar \ledsidenote{\passref{6.39.158}}\skt{सारवन्ति च भोज्यानि वर्ज्ज\-ये\-त्त\-रु\-ण\-ज्वरी ॥ }\pend
\pstart\versepar \skt{}\edtext{\skt{ईहादि}}{\Afootnote{K H; \skt{तथैव}\ A}}\skt{ }\edtext{\skt{नवधान्यानि}}{\Afootnote{H; \skt{नवधान्या}\ A; \skt{नवधान्यादी}\ K}}\skt{ }\edtext{\skt{क्रोधादींश्च}}{\Afootnote{H; \emph{om.}\ A; \skt{क्रोधादीश्च}\ K}}\skt{ }\edtext{\skt{विवर्ज्जयेत्}}{\Afootnote{K H; \skt{दिं वर्जयेच्च समासतः ॥}\ A}}\skt{ ॥}\pend

\pstart\versepar \ledsidenote{\passref{6.39.158.add}}\edtext{\edtext{\skt{शोष\-श्छ\-र्द्दि\-र्म्म\-दो\-मू\-र्च्छा\-मो\-ह\-स्तृ\-ष्णा\-ह्य\-रो\-चकः}}{\lemma{\skt{शोषश्}}\Afootnote{H; \emph{om.}\ K\textsuperscript{ac}; \skt{शोष}\ K\textsuperscript{pc}}}}{\lemma{\skt{छर्द्दि\-र्म्म\-दो\-मू\-र्च्छा\-मो\-ह\-स्तृ\-ष्णा\-ह्य\-रो\-चकः}}\Afootnote{\emph{om.}\ K\textsuperscript{ac}; \skt{र्च्छर्दिर्मदो मूर्च्छा शोफस्तृष्णा ह्यरोचकः ।}\ K\textsuperscript{pc}; \skt{छर्द्दि\-र्म्म\-दो\-मू\-र्च्छा\-मो\-ह\-स्तृष्णा ह्यरोचकः ।}\ H}}\skt{ । }\pend
\pstart\versepar \skt{}\edtext{\edtext{\skt{प्राप्नो\-त्यु\-प\-द्र\-वाने}}{\lemma{\skt{प्राप्नोत्य्}}\Afootnote{K\textsuperscript{pc} H; \emph{om.}\ K\textsuperscript{ac}}}}{\lemma{\skt{उपद्रवाने}}\Afootnote{K\textsuperscript{pc} H; \emph{om.}\ K\textsuperscript{ac}}}\skt{ }\edtext{\skt{तान्प\-रि\-षे\-का\-दि\-से\-व\-नात्}}{\Afootnote{H; \emph{om.}\ K\textsuperscript{ac}; \skt{तां परि\-षे\-का\-दि\-से\-व\-नात् ।}\ K\textsuperscript{pc}}}\skt{ ॥}\pend

\pstart\versepar \ledsidenote{\passref{6.39.159}}\edtext{\edtext{\skt{अन\-व\-स्थि\-त\-दो\-षा\-ग्ने\-रेभिः}}{\lemma{\skt{अन\-व\-स्थि\-त\-दो\-षा\-ग्नेर्}}\Afootnote{K; \skt{अन\-व\-स्थि\-त\-दो\-षाग्ने}\ A H}}}{\lemma{\skt{एभिः}}\Afootnote{K; \skt{रेभिः}\ A; \skt{…रेभिः}\ H}}\skt{ }\edtext{\skt{सन्धुक्षितो}}{\Afootnote{A H; \skt{साधुक्षितो}\ K}}\skt{ ज्वरः । }\pend
\pstart\versepar \skt{}\edtext{\skt{गम्भी\-र\-ती\-क्ष्ण\-वे\-गत्वं}}{\Afootnote{A H; \skt{गम्भी\-र\-स्ती\-क्ष्ण\-वे\-गत्वं}\ K}}\skt{ }\edtext{\edtext{\skt{यात्य\-सा\-ध्य\-त्व\-मेव}}{\lemma{\skt{यात्य्}}\Afootnote{K H; \skt{या}\ A}}}{\lemma{\skt{असाध्यत्वम्}}\Afootnote{H; \skt{त्यमाध्यत्वम्}\ A; \skt{अमाध्यत्वम्}\ K}}\skt{ }\edtext{\skt{वा}}{\Afootnote{K H; \skt{च ॥}\ A}}\skt{ ॥}\pend

\pstart\versepar \ledsidenote{\passref{6.39.161}}\skt{मुक्त }\edtext{\edtext{\skt{दुर्ब्ब\-ल\-स्या\-हि\-तै\-र्ज्ज्वरः}}{\lemma{\skt{दुर्ब्ब\-ल\-स्या\-हि\-तैर्}}\Afootnote{\skt{स्यापि ज्वरेणाशु दुर्बलस्याहितै}\ A; \skt{स्यापि ज्वरेणाह दुर्ब\-ल\-स्या\-हि\-तैर्}\ K; \skt{स्यापि ज्वरेणेह दुर्ब्ब\-ल\-स्या\-हि\-तैर्}\ H}}}{\lemma{\skt{ज्ज्वरः}}\Afootnote{K H; \skt{र्ज्वरः ।}\ A}}\skt{ । }\pend
\pstart\versepar \skt{}\edtext{\skt{प्रत्यासन्ने}}{\Afootnote{\skt{प्रत्यापन्नो}\ A; \skt{प्रत्यासनो}\ K; \skt{प्रत्यासन्ने}\ H}}\skt{ }\edtext{\skt{दहेद्देहं}}{\Afootnote{K H; \skt{दहेद्देहं}\ A}}\skt{ }\edtext{\edtext{\skt{शुष्क\-क\-क्ष\-मि\-वा\-नलः}}{\lemma{\skt{शुष्ककक्षम्}}\Afootnote{K H; \skt{शुष्कं वृक्ष}\ A}}}{\lemma{\skt{इवानलः}}\Afootnote{H; \skt{मिवानलः ॥}\ A; \skt{इवानतः ।}\ K}}\skt{ ॥}\pend

\pstart\versepar \ledsidenote{\passref{6.39.162}}\edtext{\skt{न}}{\Afootnote{H; \skt{त}\ A K}}\skt{ }\edtext{\skt{स्यात्कार्यः}}{\Afootnote{H; \skt{स्मात्कार्यः}\ A K}}\skt{ }\edtext{\skt{परीहारो}}{\Afootnote{A K; \skt{परीहारो…}\ H}}\skt{ }\edtext{\skt{ज्वरमुक्ते}}{\Afootnote{K; \skt{ज्वरमुक्तै}\ A; \skt{ज्वरमुक्ते}\ H}}\skt{ }\edtext{\skt{न}}{\Afootnote{K H; \skt{र्वि}\ A}}\skt{ }\edtext{\skt{देहिना}}{\Afootnote{K H; \skt{रिक्तव}\ A}}\skt{ । }\pend
\pstart\versepar \skt{}\edtext{\skt{यावन्न}}{\Afootnote{H; \skt{त् । यावन्न}\ A; \skt{यावान्न}\ K}}\skt{ प्रकृतिस्थः }\edtext{\skt{स्याद्दोषतः}}{\Afootnote{K; \skt{स्याद्दोषतः}\ A; \skt{स्या…द्दोषतः}\ H}}\skt{ प्राणतस्तथा ॥}\pend

\pstart\versepar \ledsidenote{\passref{6.39.163}}\skt{ज्वरे प्रमोहो भवति स्वल्पै\-र\-प्य\-व\-चे\-ष्टितैः । }\pend
\pstart\versepar \skt{निषण्णं }\edtext{\skt{भोज\-ये\-त्त\-स्मा\-न्मू\-त्रो\-च्चारौ च}}{\Afootnote{A; \skt{भोज\-ये\-त्त\-स्मा\-त्मू\-त्रो\-चारौ च}\ K; \skt{भोज\-ये…\-त्त\-स्मा\-न्मू\-त्रो\-च्चारौ च}\ H}}\skt{ कारयेत् ॥}\pend

\pstart\versepar \ledsidenote{\passref{6.39.164}}\edtext{\skt{अरोचके}}{\Afootnote{A H; \skt{आरोचके}\ K}}\skt{ }\edtext{\skt{गात्रसादे}}{\Afootnote{A K; \skt{गात्रसादे…}\ H}}\skt{ }\edtext{\skt{वैवर्ण्णाया ङ्गमलादिषु}}{\Afootnote{H; \skt{वैव\-र्ण्येऽ\-ङ्ग\-म\-ला\-दिषु ।}\ A; \skt{वैव\-र्णा\-ङ्ग\-म\-ला\-दिषु ।}\ K}}\skt{ । }\pend
\pstart\versepar \skt{}\edtext{\skt{शान्तज्वरो}}{\Afootnote{A H; \skt{शान्तज्व}\ K}}\skt{ }\edtext{\skt{विशोध्यः}}{\Afootnote{H; \skt{ऽपि शोध्यः}\ A; \skt{रे विशोध्यः}\ K}}\skt{ }\edtext{\skt{स्याद\-नु\-ब\-न्द्भ\-या\-नरः}}{\Afootnote{\skt{स्याद\-नु\-ब\-न्ध\-भ\-या\-न्नरः ॥}\ A; \skt{स्याद\-नु\-व\-न्ध\-भया नरः ।}\ K; \skt{स्यादनुब… ॥ …॥ ॥}\ H}}\skt{ ॥}\pend

\pstart\versepar \ledsidenote{\passref{6.39.165}}\skt{न जातु }\edtext{\skt{स्नपयेत्प्राज्ञः}}{\Afootnote{K; \skt{स्नापयेत्प्राज्ञः}\ A}}\skt{ सहसा }\edtext{\skt{ज्वरकर्षितः}}{\Afootnote{K; \skt{ज्वरकर्शितम् ।}\ A}}\skt{ । }\pend
\pstart\versepar \skt{}\edtext{\skt{तस्य}}{\Afootnote{K; \skt{तेन}\ A}}\skt{ }\edtext{\skt{संशमितो}}{\Afootnote{\skt{संदूषितो ह्य्}\ A; \skt{संशमितो ह्य्}\ K}}\skt{ }\edtext{\skt{ह्यस्य}}{\Afootnote{\skt{अस्य}\ A K}}\skt{ पुनरेव भवेज्ज्वरः ।}\pend

\pstart\versepar \ledsidenote{\passref{6.39.166}}\skt{चिकित्सेच्च }\edtext{\skt{ज्वरां}}{\Afootnote{K; \skt{ज्वरा}\ A}}\skt{ }\edtext{\skt{सर्वान्निमित्ताय}}{\Afootnote{K; \skt{न्सर्वान्निमित्तानां}\ A}}\skt{ विपर्ययैः । }\pend
\pstart\versepar \skt{श्रम\-क्ष\-या\-भि\-घा\-तोत्थे }\edtext{\skt{मूल\-व्या\-धि\-मु\-पा\-च\-रेत्}}{\Afootnote{\skt{मूल\-व्या\-धि\-मु\-पा\-च\-रेत् ॥}\ A K}}\skt{ ॥}\pend

\pstart\versepar \ledsidenote{\passref{6.39.168}}\skt{अतः संशमनीयानि }\edtext{\skt{कषायानि}}{\Afootnote{K; \skt{कषायाणि}\ A}}\skt{ }\edtext{\skt{निवोध मे}}{\Afootnote{K; \skt{निबोध मे ।}\ A}}\skt{ ॥ }\pend
\pstart\versepar \skt{सर्वज्वरेषु देयानि यानि वैद्येन जानता}\pend

\pstart\versepar \ledsidenote{\passref{6.39.169}}\edtext{\skt{पिप्प\-ली\-सा\-रि\-वा\-द्रा\-क्षा\-श\-त\-पु\-ष्पा\-ह\-रे\-णुभिः}}{\Afootnote{K; \skt{पिप्प\-ली\-सा\-रि\-वा\-द्रा\-क्षा\-श\-त\-पु\-षो\-आ\-ह\-रे\-णुभिः ।}\ A}}\skt{ । }\pend
\pstart\versepar \skt{}\edtext{\skt{कृतङ्कषाय}}{\Afootnote{\skt{कृतः कषायः}\ A; \skt{कृतङ्कषायं}\ K}}\skt{ }\edtext{\skt{सगु\-ड\-ङ्ह\-न्या\-च्छ्व\-स\-न\-ज\-ञ्ज्वरं}}{\Afootnote{\skt{सगुडो हन्याच्छ्वसनजं ज्वरम् ॥}\ A; \skt{सगु\-ड\-ङ्ह\-न्या\-च्छ्व\-स\-न\-ज\-ञ्ज्वरᳵ ।}\ K}}\skt{ ।}\pend

\pstart\versepar \ledsidenote{\passref{6.39.170}}\edtext{\skt{श्रितं}}{\Afootnote{K; \skt{शृतं}\ A}}\skt{ शीतकषायम्वा गुडूच्याः }\edtext{\skt{पेयमेव}}{\Afootnote{\skt{पेयमेव तु ।}\ A; \skt{पेयमेव}\ K}}\skt{ }\edtext{\skt{च}}{\Afootnote{K; \skt{बला}\ A}}\skt{ । }\pend
\pstart\versepar \skt{}\edtext{\skt{दर्भ\-म्व\-ला\-ङ्गो\-क्ष\-र\-का\-म्य\-चे\-न्या\-दा\-व\-शे\-षितं}}{\Afootnote{\skt{दर्भश्वदंष्ट्राणां कषायं पादशेषितम् ॥}\ A; \skt{दर्भा\-म्व\-ला\-ङ्गो\-क्ष\-र\-का\-म्य\-चे\-न्या\-दा\-व\-शे\-षितं ।}\ K}}\skt{ ।}\pend

\pstart \ledsidenote{\passref{6.39.171}}\edtext{\skt{शर्क\-रा\-घृ\-त\-सं\-यु\-क्त\-म्पि\-बे\-द्वा\-त\-ज्व\-रा\-पहं}}{\lemma{\skt{पिबे\-द्वा\-त\-ज्व\-रा\-पहं}}\Afootnote{\skt{पिबे\-द्वा\-त\-ज्व\-रा\-प\-हम् । शत\-पु\-ष्पा\-व\-चा\-कु\-ष्ठ\-दे\-व\-दा\-रु\-ह\-रे\-णुकाः ॥}\ A; \skt{पिवे\-द्वा\-त\-ज्व\-रा\-पहᳵ ।}\ K}}\skt{ ।}\pend

\pstart \ledsidenote{\passref{6.39.175ab}}\edtext{\skt{स्वेदालेपः}}{\Afootnote{K; \emph{om.}\ A}}\skt{ }\edtext{\skt{घृता\-भ्य\-ङ्गो\-न\-व\-स्थासु च}}{\Afootnote{K; \skt{घृता\-भ्य\-ङ्ग\-स्वे\-द\-ले\-पा\-न\-व\-स्थासु च}\ A}}\skt{ योजयेत् ॥}\pend

\pstart \ledsidenote{\passref{6.39.176ab}}\skt{शर्करा मधुरो हन्ति कषायः पैत्तिकं }\edtext{\skt{ज्वरं}}{\Afootnote{A; \skt{ज्वरᳵ ।}\ K}}\skt{ ।}\pend

\pstart \ledsidenote{\passref{6.39.175cd}}\edtext{\skt{चन्द\-नो\-शी\-र\-श्री\-प\-र्ण्णी\-प\-रू\-ष\-क\-म\-धू\-कजः}}{\Afootnote{\skt{श्रीप\-र्णी\-च\-न्द\-नो\-शी\-र\-प\-रू\-ष\-क\-म\-धू\-कजः ॥}\ A; \skt{चन्दनो शीर\-श्री\-प\-र्ण्णी\-प\-रू\-ष\-क\-म\-धू\-कजः ।}\ K}}\skt{ ।}\pend

\pstart \ledsidenote{\passref{6.39.176cd}}\edtext{\skt{पीत्तम्पित्तज्वरं}}{\Afootnote{K; \skt{पीतं पित्तज्वरं}\ A}}\skt{ हन्यात्सारिवाद्यं }\edtext{\skt{सशर्करं}}{\Afootnote{A; \skt{सशर्करᳵ ।}\ K}}\skt{ ।}\pend

\pstart\versepar \ledsidenote{\passref{6.39.177}}\edtext{\skt{हन्या\-त्स\-य\-ष्टी\-म\-धुकं}}{\Afootnote{K; \skt{सयष्टीमधुकं}\ A}}\skt{ }\edtext{\skt{तथै\-वो\-त्प\-ल\-पू\-र्वकं}}{\Afootnote{\skt{हन्या\-त्त\-थै\-वो\-त्प\-ल\-पू\-र्व\-कम् ।}\ A; \skt{तथै\-वो\-त्प\-ल\-पू\-र्वकᳵ ।}\ K}}\skt{ । }\pend
\pstart\versepar \skt{}\edtext{\skt{पाक्यं}}{\Afootnote{K; \skt{शृतं}\ A}}\skt{ }\edtext{\skt{शीत\-क\-षा\-या\-सो\-त्पलं}}{\Afootnote{K; \skt{शीतकषायं वा सोत्पलं}\ A}}\skt{ शर्करायुतं ।}\pend

\pstart\versepar \ledsidenote{\passref{6.39.178}}\edtext{\skt{गुञ्डू\-ची\-प\-द्म\-लो\-ध्राणां}}{\Afootnote{K; \skt{गुडू\-ची\-प\-द्म\-रो\-ध्राणां}\ A}}\skt{ सारि\-वो\-त्प\-ल\-यो\-स्तथा । }\pend
\pstart\versepar \skt{शर्करामधुरः क्वाथः शीतः पित्तज्वरापहः ।}\pend

\pstart\versepar \ledsidenote{\passref{6.39.179}}\edtext{\edtext{\skt{द्राक्षा\-र\-ग्व\-ध\-यो\-श्चा\-पि\-का\-श्म\-र्य\-स्या\-थवा}}{\lemma{\skt{द्राक्षा\-र\-ग्व\-ध\-योश्}}\Afootnote{A K\textsuperscript{pc}; \skt{द्राक्षारधयोश्}\ K\textsuperscript{ac}}}}{\lemma{\skt{चापि\-का\-श्म\-र्य\-स्या\-थवा}}\Afootnote{A K\textsuperscript{pc}; \skt{चापि\-का\-श्म\-र्य\-स्या\-वा\-थवा}\ K\textsuperscript{ac}}}\skt{ पुनः । }\pend
\pstart\versepar \skt{स्वादु\-ति\-क्त\-क\-षा\-याणां कषायैः शर्करायुतैः ।}\pend

\pstart\versepar \ledsidenote{\passref{6.39.180}}\edtext{\skt{सुशी\-तै\-स्स\-म\-ये\-त्तृष्णां}}{\Afootnote{K; \skt{सुशीतैः शमयेत्तृष्णां}\ A}}\skt{ प्रवृद्धां }\edtext{\skt{दाहमेव च}}{\Afootnote{\skt{दाहमेव च ।}\ A K}}\skt{ । }\pend
\pstart\versepar \skt{}\edtext{\skt{वारिशीतं}}{\Afootnote{K; \skt{शीतं}\ A}}\skt{ }\edtext{\skt{मधु\-यु\-त\-मा\-कं\-ठाद्वा}}{\Afootnote{\skt{मधुयुतं तोयमाकण्ठाद्वा}\ A; \skt{मधु\-यु\-त\-मा\-कं\-ठाद्वा}\ K}}\skt{ पिपासितं ।}\pend

\pstart\versepar \ledsidenote{\passref{6.39.181}}\edtext{\skt{पाय\-ये\-द्वा\-म\-ये\-द्वापि}}{\Afootnote{\skt{वाम\-ये\-त्पा\-य\-यित्वा तु}\ A; \skt{पाय\-ये\-द्वा\-म\-ये\-द्वापि}\ K}}\skt{ }\edtext{\skt{ततस्तृष्णा}}{\Afootnote{\skt{तेन तृष्णा}\ A; \skt{ततस्तृष्णा}\ K}}\skt{ प्रशाम्यति । }\pend
\pstart\versepar \skt{क्षीरैः }\edtext{\skt{शीतकषायैश्च}}{\Afootnote{K; \skt{क्षीरिकषायैश्च}\ A}}\skt{ }\edtext{\skt{सुर्शी\-ते\-श्चं\-द\-ना\-युतैः}}{\Afootnote{\skt{सुशी\-तै\-श्च\-न्द\-ना\-युतैः}\ A K}}\skt{ ।}\pend

\pstart\versepar \ledsidenote{\passref{6.39.182}}\edtext{\skt{अन्त\-र्दा\-ह\-म्प्र\-श\-म\-ये\-दे\-तै\-र\-न्यैश्च}}{\Afootnote{\skt{अन्तर्दाहे विधा\-त\-व्य\-मे\-भि\-श्चा\-न्यैश्च}\ A; \skt{अन्त\-र्दा\-ह\-म्प्र\-श\-म\-ये\-दे\-तै\-र\-न्यैश्च}\ K}}\skt{ शीतलैः । }\pend
\pstart\versepar \skt{पद्मकं मधुकं }\edtext{\skt{द्राक्षा}}{\Afootnote{\skt{द्राक्षां}\ A K}}\skt{ }\edtext{\skt{पुण्ड\-री\-क\-म\-थो\-त्पलं}}{\Afootnote{A; \skt{पुण्ड\-री\-क\-म\-थो\-त्पलᳵ ।}\ K}}\skt{ ।}\pend

\pstart\versepar \ledsidenote{\passref{6.39.183}}\edtext{\skt{यवां}}{\Afootnote{K; \skt{यवा}\ A}}\skt{ }\edtext{\skt{भृष्टानुशीराणि}}{\Afootnote{K; \skt{न्भृष्टानुशीराणि}\ A}}\skt{ }\edtext{\skt{समंगा}}{\Afootnote{K; \skt{समङ्गां}\ A}}\skt{ }\edtext{\skt{काश्मरीफलं}}{\Afootnote{K; \skt{कश्मरीफलम् ।}\ A}}\skt{ । }\pend
\pstart\versepar \skt{}\edtext{\skt{निद\-ध्या\-द\-प्सु\-था\-लोड्य}}{\Afootnote{\skt{निदध्यादप्सु चालोड्य}\ A; \skt{निद\-ध्या\-द\-प्सु\-था\-लोड्य}\ K}}\skt{ }\edtext{\skt{निशा\-प\-य\-सि\-तन्तु}}{\Afootnote{K; \skt{निशापर्युषितं त}\ A}}\skt{ }\edtext{\skt{तत्}}{\Afootnote{K; \skt{तः ॥}\ A}}\skt{ ।}\pend

\pstart\versepar \ledsidenote{\passref{6.39.184}}\skt{क्षौद्रेण }\edtext{\skt{युक्तम्पिवतो}}{\Afootnote{K; \skt{युक्तं पिबतो}\ A}}\skt{ ज्वरदाहौ प्रशाम्यतः । }\pend
\pstart\versepar \skt{जिह्वा\-ता\-लु\-ग\-ल\-क्लो\-म\-शोषे }\edtext{\skt{मूर्धनि}}{\Afootnote{K; \skt{मूर्ध्नि च}\ A}}\skt{ दापयेत् ।}\pend

\pstart\versepar \ledsidenote{\passref{6.39.185}}\skt{केश\-र\-म्मा\-तु\-लु\-ङ्गस्य }\edtext{\skt{घृत\-सै\-न्ध\-व\-सं\-युतं}}{\Afootnote{\skt{मधु\-सै\-न्ध\-व\-सं\-यु\-तम् ।}\ A; \skt{घृत\-सै\-न्ध\-व\-सं\-युतᳵ ।}\ K}}\skt{ । }\pend
\pstart\versepar \skt{}\edtext{\skt{शर्क\-रा\-दा\-डि\-मा\-भ्यान्तु}}{\Afootnote{K; \skt{शर्क\-रा\-दा\-डि\-माभ्यां वा}\ A}}\skt{ }\edtext{\skt{द्राक्षा\-दा\-डि\-म\-यो\-स्तथा}}{\lemma{\skt{द्राक्षा\-दा\-डि\-म\-योस्}}\Afootnote{K; \skt{द्राक्षा\-ख\-र्जू\-र\-योस्}\ A}}\skt{ ।}\pend

\pstart\versepar \ledsidenote{\passref{6.39.186}}\edtext{\skt{वैर\-स्या\-द्धा\-र\-ये\-त्क\-ल्क\-ङ\-ग\-ण्डू\-षञ्च}}{\Afootnote{\skt{वैरस्ये धारयेत्कल्कं गण्डूषं च}\ A; \skt{वैर\-स्या\-द्धा\-र\-ये\-त्क\-ल्क\-ङ्ग\-ण्डू\-षञ्च}\ K}}\skt{ }\edtext{\skt{यथाहितं}}{\Afootnote{\skt{तथा हितम् ।}\ A; \skt{यथाहितᳵ ।}\ K}}\skt{ । }\pend
\pstart\versepar \skt{}\edtext{\skt{सप्त\-च्छ\-द\-ङ्गु\-ञ्डू\-चीञ्च}}{\Afootnote{K; \skt{सप्तच्छदं गुडूचीं च}\ A}}\skt{ }\edtext{\skt{निम्व}}{\Afootnote{K; \skt{निम्बं}\ A}}\skt{ }\edtext{\skt{स्फूर्जकमेव}}{\Afootnote{\skt{स्फूर्जकमेव}\ A K}}\skt{ च ॥}\pend

\pstart\versepar \ledsidenote{\passref{6.39.187}}\skt{क्वाथयित्वा }\edtext{\skt{पिबेत्क्वाथं}}{\Afootnote{A; \skt{पिवेत्क्वाथं}\ K}}\skt{ सक्षौद्रं कफजे }\edtext{\skt{ज्वरं}}{\Afootnote{\skt{ज्वरे ।}\ A; \skt{ज्वरᳵ ।}\ K}}\skt{ । }\pend
\pstart\versepar \skt{}\edtext{\skt{त्रिक\-टु\-न्ना\-ग\-पु\-ष्पञ्च}}{\Afootnote{K; \skt{कटुत्रिकं नागपुष्पं}\ A}}\skt{ }\edtext{\skt{हरिद्रे}}{\Afootnote{K; \skt{हरिद्रा}\ A}}\skt{ }\edtext{\skt{कटुरोहिणीं}}{\Afootnote{\skt{कटुरोहिणी ॥}\ A; \skt{कटुरोहिणीᳵ ।}\ K}}\skt{ ।}\pend

\pstart\versepar \ledsidenote{\passref{6.39.188}}\skt{कौटजञ्च फलं हन्यात्सेव्यमानं }\edtext{\skt{कफज्वरे}}{\Afootnote{K; \skt{कफज्वरम् ।}\ A}}\skt{ । }\pend
\pstart\versepar \skt{}\edtext{\skt{कटुकञ्चित्रकं}}{\Afootnote{K; \skt{हरिद्रां चिअत्रकं}\ A}}\skt{ }\edtext{\skt{निम्वं}}{\Afootnote{K; \skt{निम्ब}\ A}}\skt{ }\edtext{\skt{हरि\-द्रा\-ति\-वि\-षो\-भयां}}{\Afootnote{\skt{मुशीरातिविषे वचाम् ॥}\ A; \skt{हरि\-द्रा\-ति\-वि\-षो\-भयाᳵ ।}\ K}}\skt{ ।}\pend

\pstart\versepar \ledsidenote{\passref{6.39.189}}\edtext{\skt{कुष्ठ\-म्मि\-न्द्र\-य\-वा\-म्मूर्वा}}{\Afootnote{\skt{कुष्ठ\-मि\-न्द्र\-य\-वा\-न्मूर्वां}\ A; \skt{कुष्ठ\-म्मि\-न्द्र\-य\-वा\-म्मूर्वा}\ K}}\skt{ पटोलञ्चापि }\edtext{\skt{साधितं}}{\Afootnote{A; \skt{साधितᳵ ।}\ K}}\skt{ । }\pend
\pstart\versepar \skt{}\edtext{\skt{पिवे\-न्म\-रि\-च\-सं\-युक्तं}}{\Afootnote{\skt{पिबे\-न्म\-रि\-च\-सं\-युक्तं}\ A; \skt{पिवे\-न्म\-रि\-च\-सं\-युक्तं}\ K}}\skt{ }\edtext{\skt{सगुडङ्कफजे}}{\Afootnote{K; \skt{सक्षौद्रं कफजे}\ A}}\skt{ ज्वरे ॥}\pend

\pstart\versepar \ledsidenote{\passref{6.39.192}}\skt{राज\-वृ\-क्षा\-दि\-व\-र्गस्य }\edtext{\skt{कषा\-य\-म्म\-धु\-सं\-युतं}}{\Afootnote{\skt{कषायो मधुसंयुतः ।}\ A; \skt{कषा\-य\-म्म\-धु\-सं\-युतᳵ ।}\ K}}\skt{ । }\pend
\pstart\versepar \skt{कफवातज्वरं हन्याच्छीघ्रं काले }\edtext{\skt{च}}{\Afootnote{K; \skt{ऽव}\ A}}\skt{ }\edtext{\skt{चारितं}}{\Afootnote{\skt{चारितः ॥}\ A; \skt{चारितᳵ ।}\ K}}\skt{ ।}\pend

\pstart\versepar \ledsidenote{\passref{6.39.193}}\edtext{\skt{देव\-दा\-र्व्य\-भ\-या\-भा\-र्गी\-व\-चा\-प\-र्प\-ट\-ना\-गरैः}}{\Afootnote{K; \skt{नागरं धान्यकं भार्गीमभयां सुरदारु च । वचां पर्पट}\ A}}\skt{ । }\pend
\pstart\versepar \skt{}\edtext{\skt{पान्यां}}{\Afootnote{K; \emph{om.}\ A}}\skt{ }\edtext{\skt{कट्फ\-ल\-भू\-ती\-क\-मु\-स्ता\-भा\-गै\-स्स\-मा\-सिकैः}}{\Afootnote{K; \skt{कं मुस्तं भूतीकमथ कट्फलम् ॥}\ A}}\skt{ ।}\pend

\pstart\versepar \ledsidenote{\passref{6.39.194}}\edtext{\skt{सिद्ध\-मे\-तः\-क\-षायां तु}}{\Afootnote{\emph{om.}\ A; \skt{सिद्ध\-मे\-तः\-क\-षायां तु}\ K}}\skt{ }\edtext{\skt{हिङ्गुमाक्षिकं}}{\Afootnote{\emph{om.}\ A; \skt{हिङ्गुमाक्षिकं}\ K}}\skt{ }\edtext{\skt{संयुतं}}{\Afootnote{K; \emph{om.}\ A}}\skt{ ॥ }\pend
\pstart\versepar \skt{}\edtext{\skt{पाययेज्वरितं}}{\Afootnote{K; \emph{om.}\ A}}\skt{ }\edtext{\skt{युक्ता}}{\Afootnote{K; \skt{नि}\ A}}\skt{ }\edtext{\skt{वातश्लेष्मकृते}}{\Afootnote{K; \skt{ष्क्वा}\ A}}\skt{ }\edtext{\skt{ज्वरे}}{\Afootnote{K; \skt{थ्य}\ A}}\skt{ । }\pend
\pstart\versepar \skt{}\edtext{\skt{कफ\-प्र\-से\-क\-कासे च}}{\Afootnote{K; \skt{कफवातोत्थे क्षौद्र\-हि\-ङ्गु\-स\-म\-न्वि\-तम् । दात\-व्यं(\-अ\-पा\-तव्यं)}\ A}}\skt{ }\edtext{\skt{हिक्का\-श्वा\-स\-ग\-ल\-ग्रहे}}{\Afootnote{K; \skt{श्वासकासघ्नं श्लेष्मोत्सेके गलग्रहे ॥}\ A}}\skt{ । }\pend
\pstart\versepar \skt{}\edtext{\skt{कण्ठे}}{\Afootnote{K; \emph{om.}\ A}}\skt{ }\edtext{\skt{विवद्धे}}{\Afootnote{K; \emph{om.}\ A}}\skt{ }\edtext{\skt{विमुखे}}{\Afootnote{K; \emph{om.}\ A}}\skt{ }\edtext{\skt{ज्वरितानां}}{\Afootnote{K; \emph{om.}\ A}}\skt{ }\edtext{\skt{सदाहितं}}{\Afootnote{K; \emph{om.}\ A}}\skt{ ।}\pend

\pstart \ledsidenote{\passref{6.39.198}}\edtext{\skt{सश\-र्क\-रा\-म\-क्ष\-मा\-त्रा\-ङ्क\-टु\-का\-मु\-ष्ण\-वा\-रिणा}}{\Afootnote{\skt{काश्मर्यनागरैः क्वाथः सक्षौद्रः श्लेष्मपित्तजे । सश\-र्क\-रा\-म\-क्ष\-मात्रां कटु\-का\-मु\-ष्ण\-वा\-रिणा ॥}\ A; \skt{सश\-र्क\-रा\-क्ष\-म\-क्ष\-मा\-त्रा\-ङ्क\-टु\-का\-मु\-ष्ण\-वा\-रिणा ।}\ K\textsuperscript{ac}; \skt{सश\-र्क\-रा\-म\-क्ष\-मा\-त्रा\-ङ्क\-टु\-का\-मु\-ष्ण\-वा\-रिणा ।}\ K\textsuperscript{pc}}}\skt{ ।}\pend

\pstart\versepar \ledsidenote{\passref{6.39.199}}\skt{पीत्वा }\edtext{\skt{ज्वरजयेज्जन्तु}}{\Afootnote{\skt{ज्वरं जयेज्जन्तुः}\ A; \skt{ज्वराज्जयेज्जन्तु}\ K}}\skt{ कफ\-पि\-त्त\-स\-मु\-द्भवं । }\pend
\pstart\versepar \skt{}\edtext{\edtext{\skt{किरा\-त\-ति\-क्त\-म\-मृतां}}{\lemma{\skt{किराततिक्तम्}}\Afootnote{K; \skt{किराततिक्त}\ A}}}{\lemma{\skt{अमृतां}}\Afootnote{K; \skt{ममृतां}\ A}}\skt{ }\edtext{\skt{द्राक्षामामलकं}}{\Afootnote{A; \skt{द्राक्षामामलकं}\ K}}\skt{ }\edtext{\skt{शटीं}}{\Afootnote{A; \skt{शटीᳵ ।}\ K}}\skt{ ।}\pend

\pstart \ledsidenote{\passref{6.39.200}}\edtext{\skt{नित्क्वा\-थ्य\-पि\-त्ता\-नि\-लजे}}{\Afootnote{\skt{निष्क्वाथ्य वातपित्तोत्थे}\ A; \skt{निष्क्वाथ्य पित्तानिलजे}\ K}}\skt{ }\edtext{\skt{तत्क्वाथं}}{\Afootnote{\skt{तं क्वाथं}\ A; \skt{तन्क्वाथं}\ K}}\skt{ सगुडं }\edtext{\skt{पिबेत्}}{\Afootnote{\skt{पिबेत् । रास्ना वृषोऽथ त्रिफला राजवृक्षफलैः सह ॥}\ A; \skt{पिवेत् ।}\ K}}\skt{ ।}\pend

\pstart \ledsidenote{\passref{6.39.201}}\edtext{\skt{सर्व\-दो\-ष\-स\-मुत्थे तु}}{\Afootnote{K; \skt{कषायः साधितः पीतो वातपित्तज्वरं जयेत् । सर्व\-दो\-ष\-स\-मुत्थे तु}\ A}}\skt{ संसृ\-ष्टा\-न्न\-व\-चा\-र\-येत् ।}\pend

\pstart\versepar \ledsidenote{\passref{6.39.201.add1}}\edtext{\edtext{\skt{वर्धनैर्वापि}}{\lemma{\skt{वर्धनैर्}}\Afootnote{K\textsuperscript{pc}; \emph{om.}\ K\textsuperscript{ac}}}}{\lemma{\skt{वापि}}\Afootnote{K\textsuperscript{pc}; \emph{om.}\ K\textsuperscript{ac}}}\skt{ }\edtext{\skt{हीनस्य}}{\Afootnote{K\textsuperscript{pc}; \emph{om.}\ K\textsuperscript{ac}}}\skt{ }\edtext{\skt{हानै}}{\Afootnote{K\textsuperscript{pc}; \emph{om.}\ K\textsuperscript{ac}}}\skt{ }\edtext{\skt{धर्मवूच्छ्रितस्य वा}}{\Afootnote{K\textsuperscript{pc}; \emph{om.}\ K\textsuperscript{ac}}}\skt{ }\pend
\pstart\versepar \skt{}\edtext{\skt{कफ\-स्ना\-ना\-नु\-पू\-र्व्याच}}{\Afootnote{K\textsuperscript{pc}; \emph{om.}\ K\textsuperscript{ac}}}\skt{ }\edtext{\skt{संनि\-पा\-त\-ज्व\-र\-क्रिया}}{\Afootnote{K\textsuperscript{pc}; \emph{om.}\ K\textsuperscript{ac}}}\skt{ ।}\pend

\pstart\versepar \ledsidenote{\passref{6.39.201.add2}}\edtext{\skt{हीनस्य}}{\Afootnote{K\textsuperscript{pc}; \emph{om.}\ K\textsuperscript{ac}}}\skt{ }\edtext{\skt{वर्धनं}}{\Afootnote{K\textsuperscript{pc}; \emph{om.}\ K\textsuperscript{ac}}}\skt{ }\edtext{\edtext{\skt{हानि\-र्वृ\-द्ध\-यो\-रिति}}{\lemma{\skt{हानिर्वृद्धयोर्}}\Afootnote{\emph{om.}\ K\textsuperscript{ac}; \skt{हानिर्वृद्धयोर्}\ K\textsuperscript{pc}}}}{\lemma{\skt{इति}}\Afootnote{K\textsuperscript{pc}; \emph{om.}\ K\textsuperscript{ac}}}\skt{ }\edtext{\skt{निश्चयः}}{\Afootnote{K\textsuperscript{pc}; \emph{om.}\ K\textsuperscript{ac}}}\skt{ । }\pend
\pstart\versepar \skt{}\edtext{\skt{हाप\-ना\-द\-ति\-वृ\-द्धस्य}}{\Afootnote{\emph{om.}\ K\textsuperscript{ac}; \skt{हाप\-ना\-द\-ति\-वृ\-द्धस्य}\ K\textsuperscript{pc}}}\skt{ }\edtext{\edtext{\skt{हीनयोर्वृद्धि}}{\lemma{\skt{हीनयोर्}}\Afootnote{K\textsuperscript{pc}; \emph{om.}\ K\textsuperscript{ac}}}}{\lemma{\skt{वृद्धि}}\Afootnote{K\textsuperscript{pc}; \emph{om.}\ K\textsuperscript{ac}}}\skt{ }\edtext{\skt{सम्भवः}}{\Afootnote{K\textsuperscript{pc}; \emph{om.}\ K\textsuperscript{ac}}}\skt{ ।}\pend

\pstart\versepar \ledsidenote{\passref{6.39.201.add3}}\edtext{\skt{ततस्समत्वं}}{\Afootnote{K\textsuperscript{pc}; \emph{om.}\ K\textsuperscript{ac}}}\skt{ }\edtext{\skt{दोषाणामायं}}{\Afootnote{K\textsuperscript{pc}; \emph{om.}\ K\textsuperscript{ac}}}\skt{ }\edtext{\skt{स्थानकफस्य च}}{\Afootnote{K\textsuperscript{pc}; \emph{om.}\ K\textsuperscript{ac}}}\skt{ । }\pend
\pstart\versepar \skt{}\edtext{\skt{तद्वत्स्थानां}}{\Afootnote{K\textsuperscript{pc}; \emph{om.}\ K\textsuperscript{ac}}}\skt{ }\edtext{\skt{क्रिया\-त्त\-द्व\-ति\-ज्व\-र\-वि\-नि\-र्जयः}}{\Afootnote{K\textsuperscript{pc}; \emph{om.}\ K\textsuperscript{ac}}}\skt{ ।}\pend

\pstart \ledsidenote{\passref{6.39.202}}\edtext{\skt{वृश्ची\-व\-वि\-ल्व\-व\-र्षाभ्व}}{\Afootnote{K; \skt{यथादोषोच्छ्रयं चापि ज्वरा\-न्स\-र्वा\-नु\-पा\-च\-रेत् । वृश्ची\-व\-बि\-ल्व\-व\-र्षाभ्वः}\ A}}\skt{ }\edtext{\skt{पय\-श्चो\-द\-क\-मेव}}{\lemma{\skt{पयश्चोदकम्}}\Afootnote{K\textsuperscript{pc}; \skt{पयश्चोदकम्}\ A; \skt{पयम्}\ K\textsuperscript{ac}}}\skt{ च ।}\pend

\pstart\versepar \ledsidenote{\passref{6.39.203}}\edtext{\skt{पचे\-त्क्षी\-रा\-व\-शि\-ष्ट\-न्त\-द्धन्ति}}{\Afootnote{K; \skt{पचे\-त्क्षी\-रा\-व\-शिष्टं तु तद्धि}\ A}}\skt{ }\edtext{\skt{सर्वज्वरातपि}}{\Afootnote{K; \skt{सर्व\-ज्व\-रा\-प\-हम् ।}\ A}}\skt{ । }\pend
\pstart\versepar \skt{}\edtext{\skt{उदकं}}{\Afootnote{K; \skt{उदकां}\ A}}\skt{ }\edtext{\skt{त्रिगुणं}}{\Afootnote{K; \skt{शास्त्रयः}\ A}}\skt{ क्षीरं }\edtext{\skt{शिंश\-पा\-सा\-र\-मेव}}{\Afootnote{\skt{शिंश\-पा\-सा\-र\-संयु}\ A; \skt{शिंश\-पा\-सा\-र\-मेव}\ K}}\skt{ }\edtext{\skt{च}}{\Afootnote{K; \skt{तम् ॥}\ A}}\skt{ ।}\pend

\pstart \ledsidenote{\passref{6.39.204ab}}\edtext{\skt{तत्क्षी\-र\-शे\-षा\-ङ्क\-थि\-त\-म्पेयं}}{\lemma{\skt{तत्क्षी\-र\-शे\-षा\-ङ्क\-थि\-तम्}}\Afootnote{K; \skt{तत्क्षीरशेषं क्वथितं}\ A}}\skt{ }\edtext{\skt{सर्वज्वरापहं}}{\Afootnote{A; \skt{सर्वज्वरापहᳵ ।}\ K}}\skt{ ।}\pend

\pstart\versepar \ledsidenote{\passref{6.39.208}}\skt{अनन्तां }\edtext{\skt{वाल\-क\-म्मु\-स्त\-न्ना\-गरं}}{\Afootnote{K; \skt{बालकं मुस्तां नागरं}\ A}}\skt{ }\edtext{\skt{कटुरोहिणीं}}{\Afootnote{A; \skt{कटुरोहिणीᳵ ।}\ K}}\skt{ । }\pend
\pstart\versepar \skt{सुखाम्बुना }\edtext{\skt{पाययीत}}{\Afootnote{K; \skt{प्रागु\-द\-या\-त्पा\-य\-येता}\ A}}\skt{ }\edtext{\skt{पिष्ट्वा कर्षसमन्नरः}}{\Afootnote{K; \skt{क्षसंमितम् ॥}\ A}}\skt{ ।}\pend

\pstart\versepar \ledsidenote{\passref{6.39.209}}\skt{एष }\edtext{\skt{सर्वज्वरां}}{\Afootnote{K; \skt{सर्वज्वरा}\ A}}\skt{ }\edtext{\skt{हन्ति}}{\Afootnote{K; \skt{न्हन्ति}\ A}}\skt{ }\edtext{\skt{दीपयेत्स}}{\Afootnote{K; \skt{दीपयत्या}\ A}}\skt{ }\edtext{\skt{हुताशनं}}{\Afootnote{\skt{शु चानलम् ।}\ A; \skt{हुताशनᳵ ।}\ K}}\skt{ । }\pend
\pstart\versepar \skt{द्रव्याणि दीपनीयानि तथा वैरेचनानि च ।}\pend

\pstart \ledsidenote{\passref{6.39.210ab}}\edtext{\skt{एकैकशो}}{\Afootnote{K; \skt{एकशो}\ A}}\skt{ }\edtext{\skt{द्वन्द्वत्यो}}{\Afootnote{\skt{वा द्विशो वाप्}\ A; \skt{द्वन्द्वसो}\ K}}\skt{ वा ज्वरघ्नानि }\edtext{\skt{विभावयेत्}}{\Afootnote{K; \skt{प्रयोजयेत् ।}\ A}}\skt{ ।}\pend

\pstart \ledsidenote{\passref{6.39.204cd}}\edtext{\skt{नल\-वे\-त\-स\-यो\-र्मूलं}}{\lemma{\skt{मूलं}}\Afootnote{K; \skt{मूले}\ A}}\skt{ }\edtext{\skt{मूर्वया}}{\Afootnote{K; \skt{मूर्वायां}\ A}}\skt{ }\edtext{\skt{सुरदारुभिः}}{\Afootnote{K; \skt{देवदारुणि ॥}\ A}}\skt{ ।}\pend

\pstart\versepar \ledsidenote{\passref{6.39.210.add}}\skt{तृष्णाङ्गापि वल\-द्रा\-क्षा\-च\-न्द\-नै\-र्वा\-म\-य\-ष्टिभिः । }\pend
\pstart\versepar \skt{श्वास\-का\-सा\-ङ्ग\-शू\-ला\-र्त्त\-म्प\-ञ्च\-मू\-लेन वा श्रितः ।}\pend

\pstart \ledsidenote{\passref{6.39.205}}\edtext{\skt{हरि\-द्रा\-भ\-द्र\-मुस्ता च}}{\Afootnote{K; \skt{कषायं विधिवत्कृत्वा पेय\-मे\-त\-ज्ज्व\-रा\-प\-हम् । हरिद्रा भद्रमुस्तं च}\ A}}\skt{ त्रिफला }\edtext{\skt{कटुरोहिणीं}}{\Afootnote{\skt{कटुरोहिणी ॥}\ A; \skt{कटुरोहिणीᳵ ।}\ K}}\skt{ ।}\pend

\pstart\versepar \ledsidenote{\passref{6.39.206}}\edtext{\skt{पिचु\-म\-र्द\-प\-टोलौ च}}{\Afootnote{K; \skt{पिचुमन्दः पटोली च}\ A}}\skt{ }\edtext{\skt{देव\-दा\-रु\-नि\-दि\-द्विका}}{\Afootnote{\skt{देवदारु निदिग्धिका ।}\ A K}}\skt{ । }\pend
\pstart\versepar \skt{एषाङ्कषायः पीतस्तु संनिपातज्वरं जयेत् ।}\pend

\pstart\versepar \ledsidenote{\passref{6.39.207}}\edtext{\skt{अवि\-प\-क्ति\-प्र\-से\-कञ्च}}{\Afootnote{K; \skt{अविपक्तिं प्रसेकं च}\ A}}\skt{ }\edtext{\skt{शोफ\-म्प\-क्ति\-म\-रो\-चकः}}{\lemma{\skt{पक्तिमरोचकः}}\Afootnote{\skt{कास\-म\-रो\-च\-कम् ।}\ A; \skt{पक्तिमरोचकᳵ ।}\ K}}\skt{ । }\pend
\pstart\versepar \skt{त्रैफलो वा }\edtext{\skt{समर्प्पिष्टः}}{\Afootnote{\skt{ससर्पिष्कः}\ A K}}\skt{ }\edtext{\skt{क्वाथ\-र्म्य\-य\-न्त्रि\-दो\-षजे}}{\Afootnote{\skt{क्वाथः पेयस्त्रिदोषजे ॥}\ A; \skt{क्वाथ\-स्र्पे\-य\-स्त्रि\-दो\-षजे ।}\ K}}\skt{ ।}\pend

\pstart \ledsidenote{\passref{6.39.210cd}}\edtext{\skt{सर्पि\-म\-ध्व\-भ\-या\-तैलं}}{\Afootnote{K; \skt{सर्पि\-र्म\-ध्व\-भ\-या\-तैल}\ A}}\skt{ }\edtext{\skt{लेहोयं}}{\Afootnote{K; \skt{लेहोऽयं}\ A}}\skt{ सर्वजं ज्वरं ।}\pend

\pstart \ledsidenote{\passref{6.39.211}}\edtext{\skt{शान्ति\-न्न\-ये\-त्तृ\-व\-द्वापि}}{\Afootnote{\skt{शान्तिं नयेत्त्रिवृच्चापि}\ A; \skt{शान्ति\-न्न\-ये\-त्तृ\-व\-द्वापि}\ K}}\skt{ }\edtext{\skt{सक्षौ\-द्रा\-वि\-ष\-म\-ज्वरं}}{\Afootnote{\skt{सक्षौद्रा प्रबलं ज्वरम् । ज्वरे तु विषमे कार्यमूर्ध्वं चाधश्च शोधनम् ॥}\ A; \skt{सक्षौद्रा विषमज्वरᳵ ।}\ K}}\skt{ ।}\pend

\pstart\versepar \ledsidenote{\passref{6.39.212}}\edtext{\edtext{\skt{घृत\-प्ली\-हो\-द\-रो\-क्तन्तु}}{\lemma{\skt{घृत\-प्ली\-हो\-द\-रो\-क्तन्}}\Afootnote{K; \skt{घृतं प्लीहोदरोक्तं वा}\ A}}}{\lemma{\skt{तु}}\Afootnote{K; \skt{नि}\ A}}\skt{ }\edtext{\skt{हन्याअ विषमज्वरां}}{\Afootnote{\skt{हन्या\-द्वि\-ष\-म\-ज्व\-रम् ।}\ A; \skt{हन्या\-त्स\-वि\-ष\-म\-ज्वरां ॥}\ K}}\skt{ ॥ }\pend
\pstart\versepar \skt{}\edtext{\skt{गुड\-प्र\-गा\-ढ\-न्त्रि\-फला}}{\Afootnote{K; \skt{गुडप्रगाढां त्रिफलां}\ A}}\skt{ }\edtext{\skt{पिबेद्वा}}{\Afootnote{\skt{पिबेद्वा}\ A; \skt{पिवेद्वा}\ K}}\skt{ विषमार्दितः ।}\pend

\pstart\versepar \ledsidenote{\passref{6.39.213}}\edtext{\skt{गुडू\-ची\-मु\-स्त\-धा\-त्री\-णा\-ङ्क\-षा\-यम्वा}}{\Afootnote{K; \skt{गुडू\-ची\-नि\-म्ब\-धा\-त्रीणां कषायं वा}\ A}}\skt{ }\edtext{\skt{समाक्षिकं}}{\Afootnote{A; \skt{समाक्षिकᳵ ।}\ K}}\skt{ । }\pend
\pstart\versepar \skt{प्रातः }\edtext{\skt{प्रातःससर्पिर्वा}}{\Afootnote{K; \skt{प्रातः ससर्पिष्कं}\ A}}\skt{ रसो\-न\-मु\-प\-यो\-ज\-येत् ।}\pend

\pstart \ledsidenote{\passref{6.39.215}}\edtext{\skt{सर्पि\-क्षी\-र\-सि\-ता\-क्षौ\-द्र\-मा\-ग\-धीर्वा}}{\Afootnote{K; \skt{हरीतक्याश्च सर्वोऽयं त्रिविधो योग इष्यते । सर्पिः\-क्षी\-र\-सि\-ता\-क्षौ\-द्र\-मा\-ग\-धीर्वा}\ A}}\skt{ }\edtext{\skt{यथाबलं}}{\Afootnote{A; \skt{यथावलं}\ K}}\pend

\pstart \ledsidenote{\passref{6.39.216}}\edtext{\skt{पिप्प\-ली\-व\-र्द्ध\-मानं वा}}{\Afootnote{K; \skt{दश\-मू\-ली\-क\-षा\-येण मागधीर्वा प्रयोजयेत् । पिप्प\-ली\-व\-र्ध\-मानं वा}\ A}}\skt{ }\edtext{\skt{पिबे\-त्क्षी\-र\-र\-सा\-सनः}}{\Afootnote{\skt{पिबे\-त्क्षी\-र\-र\-सा\-शनः ॥}\ A; \skt{पिवे\-त्क्षी\-र\-र\-सा\-सनः ।}\ K}}\skt{ ।}\pend

\pstart\versepar \ledsidenote{\passref{6.39.217}}\skt{ताम्रचूडस्य मांसेन }\edtext{\skt{पिबेद्वा}}{\Afootnote{\skt{पिबेद्वा}\ A; \skt{पिवेद्वा}\ K}}\skt{ }\edtext{\skt{मद्यमुत्तमं}}{\Afootnote{\skt{मद्यमुत्तमम् ।}\ A; \skt{मद्यमुत्तमᳵ ।}\ K}}\skt{ । }\pend
\pstart\versepar \skt{}\edtext{\skt{तृफ\-ला\-ग्नि\-म\-न्थ\-को\-ल\-क्वाथे}}{\Afootnote{K; \skt{कोला\-ग्नि\-म\-न्थ\-त्रि\-फ\-ला\-क्वाथे}\ A}}\skt{ दध्ना घृतम्पचेत् ।}\pend

\pstart\versepar \ledsidenote{\passref{6.39.218}}\edtext{\edtext{\skt{तिल्व\-का\-वा\-प\-मे\-तद्धि}}{\lemma{\skt{तिल्वकावापम्}}\Afootnote{K; \skt{तिल्वकावाप}\ A}}}{\lemma{\skt{एतद्धि}}\Afootnote{\skt{मेतद्धि}\ A; \skt{एतद्धि}\ K}}\skt{ विष\-म\-ज्व\-र\-ना\-शनं । }\pend
\pstart\versepar \skt{त्रिभिश्चतुर्भिः पिबेत्क्वाथं पञ्चभिर्वा समन्वितः । }\pend
\pstart\versepar \skt{मधुकस्य पटोलस्य रोहिण्या मुस्तकस्य च । }\pend
\pstart\versepar \skt{हरीतक्याश्च तत्सर्वन्त्रिविधो योग इष्यते ।}\edtext{}{\lemma{\skt{त्रिभिश्चतुर्भिः पिबेत्क्वाथं पञ्चभिर्वा समन्वितः । मधुकस्य पटोलस्य रोहिण्या मुस्तकस्य च । हरीतक्याश्च तत्सर्वन्त्रिविधो योग इष्यते}}\Afootnote{\skt{पिप्प\-ल्य\-ति\-वि\-षा\-द्रा\-क्षा\-रा\-रि\-वा\-बि\-ल्व\-च\-न्दनैः ॥}\ A; \emph{om.}\ K}}\pend

\pstart \ledsidenote{\passref{6.39.229}}\edtext{\skt{विडि\-ङ्ग\-मु\-स्त\-त्रि\-फ\-ला\-म\-ञ्जि\-ष्ठा\-दा\-डि\-मो\-त्पलैः}}{\Afootnote{K; \skt{हन्या\-न्न\-य\-न\-व\-द\-न\-श्र\-व\-ण\-घ्रा\-ण\-जा\-न्ग\-दान् । विड\-ङ्ग\-त्रि\-फ\-ला\-मु\-स्त\-म\-ञ्जि\-ष्ठा\-दा\-डि\-मो\-त्पलैः ॥}\ A}}\skt{ ।}\pend

\pstart\versepar \ledsidenote{\passref{6.39.230}}\edtext{\edtext{\skt{स्यामै\-ल\-वा\-लु\-कै\-ला\-भि\-श्च\-न्दना}}{\lemma{\skt{स्यामै\-ल\-वा\-लु\-कै\-ला\-भिश्}}\Afootnote{\skt{प्रिय\-ङ्ग्वे\-लै\-ल\-वा\-लूक}\ A; \skt{स्यामै\-ल\-वा\-लू\-कै\-ला\-भिश्}\ K}}\skt{ मरदारु च}}{\lemma{\skt{चन्दना मरदारु च}}\Afootnote{K; \skt{चन्द\-ना\-म\-र\-दा\-रुभिः ।}\ A}}\skt{ ॥ }\pend
\pstart\versepar \skt{}\edtext{\skt{बर्हि\-ष्ठ\-कु\-ष्ठ\-र\-जनी पर्ण्णिनी सारिवाद्वये}}{\Afootnote{\skt{बर्हि\-ष्ठ\-कु\-ष्ठ\-र\-ज\-नी\-प\-र्णि\-नी\-सा\-रि\-वा\-द्वयैः ॥}\ A; \skt{वर्हि\-ष्ठ\-कु\-ष्ठ\-र\-ज\-नी\-प\-र्ण्णि\-नी\-सा\-रि\-वा\-द्वये ।}\ K}}\skt{ ।}\pend

\pstart\versepar \ledsidenote{\passref{6.39.231}}\edtext{\skt{हरे\-णु\-का\-त्रि\-वृ\-द्द\-न्ती\-व\-चा\-ता\-ली\-श\-के\-सरैः}}{\Afootnote{A; \skt{हरे\-णु\-का\-तृ\-वृ\-द्द\-न्ती\-व\-चा\-ता\-ली\-श\-के\-सरैः ।}\ K}}\skt{ । }\pend
\pstart\versepar \skt{}\edtext{\skt{द्विःक्षीरं}}{\Afootnote{K; \skt{द्विक्षीरं}\ A}}\skt{ }\edtext{\skt{विप\-चे\-त्स\-र्पि\-र्मा\-ल\-ती\-कु\-सु\-मै\-स्सह}}{\lemma{\skt{सह}}\Afootnote{A; \skt{सहः ।}\ K}}\skt{ ।}\pend

\pstart \ledsidenote{\passref{6.39.232ab}}\edtext{\skt{विष\-म\-ज्व\-र\-श्वा\-स\-का\-स\-गु\-ल्मो\-न्मा\-द\-ज्व\-रा\-पहं}}{\Afootnote{\skt{जीर्ण\-ज्व\-र\-श्वा\-स\-का\-स\-गु\-ल्मो\-न्मा\-द\-ग\-रा\-प\-हम् ।}\ A; \skt{विष\-म\-ज्व\-र\-श्वा\-स\-का\-स\-गु\-ल्मो\-न्मा\-द\-ज्व\-रा\-पहᳵ ।}\ K}}\skt{ ।}\pend

\pstart\versepar \ledsidenote{\passref{6.39.232.add1}}\skt{एतदेव हविः }\edtext{\skt{पक्वा\-ञ्जी\-व\-नी\-यो\-प\-सं\-सृतः}}{\lemma{\skt{जीव\-नी\-यो\-प\-सं\-सृतः}}\Afootnote{\skt{जीवनीयो पसंसृतᳵ ।}\ K}}\skt{ । }\pend
\pstart\versepar \skt{चतुर्गुणेन पयसा }\edtext{\skt{महा\-क\-ल्पो\-श\-मु\-च्यते}}{\lemma{\skt{महाकल्पोशम्}}\Afootnote{\skt{महाकल्पेशम्}\ K}}\skt{ ।}\pend

\pstart\versepar \ledsidenote{\passref{6.39.232.add2}}\edtext{\skt{अप\-स्मा\-र\-गृ\-हा\-शोषं}}{\Afootnote{\skt{अपस्मारगृहां शोषं}\ K}}\skt{ }\edtext{\skt{क्लैव्य\-का\-र्श्य\-म\-की\-ङ्गीताः}}{\Afootnote{\skt{क्लैव्य\-का\-र्श्य\-म\-की\-ङ्गीताᳵ ।}\ K}}\skt{ । }\pend
\pstart\versepar \skt{घृत\-मे\-त\-न्नि\-ह\-न्याशु ये चापि विषमज्वराः ।}\pend

\pstart\versepar \ledsidenote{\passref{6.39.245}}\edtext{\skt{त्रिफ\-ला\-शी\-र\-स\-म्पा\-क\-क\-टु\-का\-ति\-वि\-षा\-घनैः}}{\Afootnote{K; \skt{त्रिफ\-लो\-शी\-र\-श\-म्पा\-क\-क\-टु\-का\-ति\-वि\-षा\-घनैः ।}\ A}}\skt{ । }\pend
\pstart\versepar \skt{}\edtext{\skt{शता\-व\-री\-स\-प्त\-प\-र्णी\-गु\-डू\-ची\-र\-ज\-नी\-द्वयैः}}{\Afootnote{K; \skt{शता\-व\-री\-स\-प्त\-प\-र्ण\-गु\-डू\-ची\-र\-ज\-नी\-द्वयैः ॥}\ A}}\skt{ ।}\pend

\pstart\versepar \ledsidenote{\passref{6.39.246}}\edtext{\skt{चित्र\-क\-त्रि\-वृ\-ता\-मू\-र्वा\-प\-टो\-ला\-रि\-ष्ट\-वा\-लकैः}}{\Afootnote{K; \skt{चित्र\-क\-त्रि\-वृ\-त\-क\-व\-चा\-वि\-शा\-ला\-प\-टो\-ला\-रि\-ष्ट\-वा\-लकैः ।}\ A}}\skt{ । }\pend
\pstart\versepar \skt{किरा\-त\-ति\-क्त\-क\-व\-चा\-वि\-शा\-ला\-प\-द्म\-को\-त्पलैः ।}\pend

\pstart\versepar \ledsidenote{\passref{6.39.247}}\skt{सारि\-वा\-द्व\-य\-य\-ष्ट्या\-ह्व\-च\-वि\-का\-र\-क्त\-च\-न्दनैः । }\pend
\pstart\versepar \skt{}\edtext{\skt{दुरा\-ल\-भा\-प\-र्प\-र्ट\-क\-त्रा\-य\-मा\-णा\-ट\-रू\-षकैः}}{\Afootnote{K; \skt{दुरा\-ल\-भा\-प\-र्प\-ट\-क\-त्रा\-य\-मा\-णा\-ट\-रू\-षकैः ॥}\ A}}\skt{ ॥}\pend

\pstart\versepar \ledsidenote{\passref{6.39.248}}\edtext{\skt{रास्ना\-कु\-ट\-ज\-म\-ञ्जि\-ष्ठा\-मा\-ग\-धी\-ना\-ग\-रै\-स्तथा}}{\lemma{\skt{रास्ना\-कु\-ट\-ज\-म\-ञ्जि\-ष्ठा\-मा\-ग\-धी\-ना\-ग\-रैस्}}\Afootnote{K; \skt{रास्ना\-कु\-ङ्कु\-म\-म\-ञ्जि\-ष्ठा\-मा\-ग\-धी\-ना\-ग\-रैस्}\ A}}\skt{ । }\pend
\pstart\versepar \skt{}\edtext{\skt{धात्रीफलरसे}}{\Afootnote{K; \skt{धात्रीफलरसैः}\ A}}\skt{ }\edtext{\skt{सम्यग्द्विगुणे}}{\Afootnote{K; \skt{सम्यग्द्विगुणैः}\ A}}\skt{ साधितं हविः ।}\pend

\pstart\versepar \ledsidenote{\passref{6.39.249}}\edtext{\skt{परि\-स\-र्प\-ज्व\-र\-श्वा\-स\-गु\-ल्म\-कु\-ष्ठ\-वि\-चा\-रणं}}{\Afootnote{\skt{परि\-स\-र्प\-ज्व\-र\-श्वा\-स\-गु\-ल्म\-कु\-ष्ठ\-नि\-वा\-र\-णम् ।}\ A; \skt{परि\-वा\-स\-र्प\-ज्व\-र\-श्वा\-स\-गु\-ल्म\-कु\-ष्ठ\-वि\-चा\-रणᳵ ।}\ K\textsuperscript{ac}; \skt{परि\-स\-र्प\-ज्व\-र\-श्वा\-स\-गु\-ल्म\-कु\-ष्ठ\-वि\-चा\-रणᳵ ।}\ K\textsuperscript{pc}}}\skt{ । }\pend
\pstart\versepar \skt{}\edtext{\skt{पाण्डु\-प्ली\-हा\-ग्नि\-सा\-देभ्य}}{\Afootnote{K\textsuperscript{pc}; \skt{पाण्डु\-प्ली\-हा\-ग्नि\-सा\-दिभ्य}\ A; \skt{पाण्डु\-प्ली\-हा\-ग्नि\-सा\-दे\-स्वभ्य}\ K\textsuperscript{ac}}}\skt{ एतदेव }\edtext{\skt{हितं}}{\Afootnote{K; \skt{परं हितम् ॥}\ A}}\skt{ }\edtext{\skt{परं}}{\Afootnote{\emph{om.}\ A; \skt{परᳵ ।}\ K}}\skt{ ।}\pend

\pstart\versepar \ledsidenote{\passref{6.39.250}}\edtext{\skt{दार्वी\-प\-टो\-ल\-ङ्क\-टु\-का\-न्नि\-म्व\-म्वा\-सा\-फ\-ल\-त्रिकं}}{\Afootnote{\skt{पटो\-ल\-क\-टु\-का\-दा\-र्वी\-नि\-म्ब\-वा\-सा\-फ\-ल\-त्रि\-कम् ।}\ A; \skt{दार्वीपटोलां कटु\-का\-न्नि\-म्व\-म्वा\-सा\-फ\-ल\-त्रिकं ।}\ K}}\skt{ । }\pend
\pstart\versepar \skt{}\edtext{\skt{दुरालभां}}{\Afootnote{K; \skt{दुरालभा}\ A}}\skt{ }\edtext{\skt{पर्पटकां}}{\Afootnote{\skt{पर्पटक}\ A; \skt{पर्पटकं}\ K}}\skt{ }\edtext{\skt{त्रायतीञ्च}}{\Afootnote{K; \skt{त्रायमाणाः}\ A}}\skt{ }\edtext{\skt{पलोन्मितां}}{\Afootnote{K; \skt{पलोन्मिताः ॥}\ A}}\skt{ ।}\pend

\pstart\versepar \ledsidenote{\passref{6.39.251}}\edtext{\skt{प्रस्थ\-मा\-म\-ल\-कानां च}}{\Afootnote{A; \skt{प्रस्थ\-मा\-म\-ल\-कानां च}\ K}}\skt{ }\edtext{\skt{क्वाथयेद्धि}}{\Afootnote{K; \skt{क्वाथयेत्स}\ A}}\skt{ }\edtext{\skt{मलस्तमि}}{\Afootnote{K; \skt{लिलार्मणे ।}\ A}}\skt{ । }\pend
\pstart\versepar \skt{तेन पादावशेषेण घृत\-प्र\-स्थ\-म्वि\-पा\-च\-येत् ॥}\pend

\pstart\versepar \ledsidenote{\passref{6.39.252}}\edtext{\skt{कल्कै\-र्भू\-नि\-म्व\-कु\-ट\-ज\-मु\-स्त\-य\-ष्ट्या\-ह्व\-च\-न्दनैः}}{\lemma{\skt{भूनि\-म्व\-कु\-ट\-ज\-मु\-स्त\-य\-ष्ट्या\-ह्व\-च\-न्दनैः}}\Afootnote{\skt{कुट\-ज\-भू\-नि\-म्ब\-घ\-न\-य\-ष्ट्या\-ह्व\-च\-न्दनैः ।}\ A; \skt{भूनि\-म्व\-कु\-ट\-ज\-मु\-स्त\-य\-ष्ट्या\-ह्व\-च\-न्दनैः ।}\ K}}\skt{ । }\pend
\pstart\versepar \skt{सपि\-प्प\-ली\-कै\-स्त\-त्सि\-द्ध\-ञ्च\-क्षुष्यं शुक्लयोर्हितं ॥}\pend

\pstart \ledsidenote{\passref{6.39.253}}\edtext{\skt{घ्राण\-क\-र्ण्णा\-ष्ठ\-व\-र्त्म\-तृ\-द्मु\-स्वे\-रो\-ग\-व्र\-णा\-पहं}}{\Afootnote{\skt{घ्राण\-क\-र्णा\-क्षि\-व\-द\-न\-व\-र्त्म\-रो\-ग\-व्र\-णा\-प\-हम् । रक्त्क\-पि\-त्त\-क\-फ\-स्वे\-द\-क्ले\-द\-पू\-यो\-प\-शो\-ष\-णम् ॥}\ A; \skt{घ्राण\-क\-र्ण्णा\-ष्ठ\-व\-र्त्म\-तृ\-द्मु\-स्वे\-रो\-ग\-व्र\-णा\-पहᳵ ।}\ K}}\skt{ ।}\pend

\pstart \ledsidenote{\passref{6.39.254ab}}\edtext{\skt{काम\-ला\-ज्व\-र\-वै\-स\-र्प\-ग\-ण्ड\-मा\-पि\-प्प\-ली\-सा\-रि\-वा\-मु\-स्त\-मु\-शीरं}}{\Afootnote{\skt{काम\-ला\-ज्व\-र\-वी\-स\-र्प\-ग\-ण्ड\-मा\-ला\-हरं परम् ।}\ A; \skt{काम\-ला\-ज्व\-र\-वै\-स\-र्प\-ग\-ण्ड\-मा\-हरं}\ K\textsuperscript{ac}; \skt{काम\-ला\-ज्व\-र\-वै\-स\-र्प\-ग\-ण्ड\-मा\-ला\-हरं परᳵ ।}\ K\textsuperscript{pc}}}\skt{ }\edtext{\skt{कटुरोहिणी}}{\Afootnote{\emph{om.}\ A K\textsuperscript{pc}; \skt{परᳵ ।}\ K\textsuperscript{ac}}}\skt{ ।}\pend

\pstart\versepar \ledsidenote{\passref{6.39.254.add1}}\edtext{\skt{पिप्प\-ली\-सा\-रि\-वा\-मु\-स्त\-मु\-शीरं}}{\Afootnote{K\textsuperscript{pc}; \emph{om.}\ K\textsuperscript{ac}}}\skt{ }\edtext{\skt{कटुरोहिणी}}{\Afootnote{K\textsuperscript{pc}; \emph{om.}\ K\textsuperscript{ac}}}\skt{ । }\pend
\pstart\versepar \skt{}\edtext{\skt{कलि\-ङ्ग\-का\-म्स्त्या\-म\-लकी}}{\Afootnote{K\textsuperscript{pc}; \emph{om.}\ K\textsuperscript{ac}}}\skt{ }\edtext{\skt{चन्दनातिविषै}}{\Afootnote{K\textsuperscript{pc}; \emph{om.}\ K\textsuperscript{ac}}}\skt{ }\edtext{\skt{स्थिरा}}{\Afootnote{K\textsuperscript{pc}; \emph{om.}\ K\textsuperscript{ac}}}\skt{ ।}\pend

\pstart\versepar \ledsidenote{\passref{6.39.254.add2}}\edtext{\skt{विण्ठ\-मा\-म\-ल\-के\-न्द्राक्षा}}{\Afootnote{\emph{om.}\ K\textsuperscript{ac}; \skt{विण्ठ\-मा\-म\-ल\-के\-न्द्राक्षा}\ K\textsuperscript{pc}}}\skt{ }\edtext{\skt{ओप\-मा\-ला\-नि\-दि\-ग्ठका}}{\Afootnote{K\textsuperscript{pc}; \emph{om.}\ K\textsuperscript{ac}}}\skt{ । }\pend
\pstart\versepar \skt{}\edtext{\skt{साधितं}}{\Afootnote{K\textsuperscript{pc}; \emph{om.}\ K\textsuperscript{ac}}}\skt{ }\edtext{\edtext{\skt{हविरेतैस्तु}}{\lemma{\skt{हविर्}}\Afootnote{K\textsuperscript{pc}; \emph{om.}\ K\textsuperscript{ac}}}}{\lemma{\skt{एतैस्तु}}\Afootnote{K\textsuperscript{pc}; \emph{om.}\ K\textsuperscript{ac}}}\skt{ }\edtext{\edtext{\skt{जीर्ण्ण\-ज्व\-र\-म\-पो\-हति}}{\lemma{\skt{जीर्ण्णज्वरम्}}\Afootnote{K\textsuperscript{pc}; \emph{om.}\ K\textsuperscript{ac}}}}{\lemma{\skt{अपोहति}}\Afootnote{K\textsuperscript{pc}; \emph{om.}\ K\textsuperscript{ac}}}\skt{ ॥}\pend

\pstart\versepar \ledsidenote{\passref{6.39.254.add3}}\edtext{\edtext{\skt{क्षय\-का\-स\-शि\-र\-शू\-ल\-पा\-र्श्व\-शू\-ल\-फ\-णी\-शकां।}}{\lemma{\skt{क्षयकासशिर्}}\Afootnote{K\textsuperscript{pc}; \emph{om.}\ K\textsuperscript{ac}}}}{\lemma{\skt{अशू\-ल\-पा\-र्श्व\-शू\-ल\-फ\-णी\-शकां}}\Afootnote{\emph{om.}\ K\textsuperscript{ac}; \skt{अशू\-ल\-पा\-र्श्व\-शू\-ल\-फ\-णी\-शकाᳵ।}\ K\textsuperscript{pc}}}\skt{ }\pend
\pstart\versepar \skt{}\edtext{\edtext{\skt{अंशा\-भि\-घा\-त\-म\-ग्नीञ्च}}{\lemma{\skt{अंशाभिघातम्}}\Afootnote{K\textsuperscript{pc}; \emph{om.}\ K\textsuperscript{ac}}}}{\lemma{\skt{अग्नीञ्च}}\Afootnote{K\textsuperscript{pc}; \emph{om.}\ K\textsuperscript{ac}}}\skt{ }\edtext{\skt{विषमं}}{\Afootnote{K\textsuperscript{pc}; \emph{om.}\ K\textsuperscript{ac}}}\skt{ }\edtext{\skt{संनियच्छति}}{\Afootnote{K\textsuperscript{pc}; \emph{om.}\ K\textsuperscript{ac}}}\skt{ ।}\pend

\pstart\versepar \ledsidenote{\passref{6.39.259}}\edtext{\skt{निर्वि\-षै\-र्भु\-ज\-गै\-र्ना\-गै\-र्वि\-नी\-तै\-र्नृ\-प\-त\-स्करैः}}{\Afootnote{\skt{निर्वि\-षै\-र्भु\-ज\-गै\-र्ना\-गै\-र्वि\-नीतैः कृततस्करैः ।}\ A; \skt{निर्वि\-षै\-र्भु\-ज\-गै\-र्ना\-गै\-र्वि\-नी\-तै\-र्नृ\-प\-त\-स्करैः ।}\ K}}\skt{ । }\pend
\pstart\versepar \skt{त्रासयेदागमे }\edtext{\skt{चैन\-न्त\-द\-ह\-र्भो\-ज\-येन्न}}{\lemma{\skt{चैनन्तदहर्}}\Afootnote{K; \skt{चैनं तदहर्}\ A}}\skt{ च ।}\pend

\pstart\versepar \ledsidenote{\passref{6.39.260}}\edtext{\skt{अत्य\-भि\-ष्य\-न्दि\-गु\-रु\-भि\-र्वा\-म\-येद्वा}}{\lemma{\skt{अत्य\-भि\-ष्य\-न्दि\-गु\-रु\-भिर्}}\Afootnote{A; \skt{अत्य\-भि\-ष्य\-न्दि\-गु\-रु\-भिर्}\ K}}\skt{ पुनः पुनः । }\pend
\pstart\versepar \skt{}\edtext{\skt{मद्यम्वा}}{\lemma{\skt{वा}}\Afootnote{K; \skt{तीक्ष्णं}\ A}}\skt{ }\edtext{\skt{पाययेत्तीक्ष्णं}}{\Afootnote{K; \skt{पाययेत}\ A}}\skt{ घृतम्वा }\edtext{\skt{ज्वरनाशनः}}{\Afootnote{\skt{ज्वरनाशनम् ॥}\ A; \skt{ज्वरनाशनᳵ ।}\ K}}\skt{ ।}\pend

\pstart\versepar \ledsidenote{\passref{6.39.261}}\skt{पुराणम्वा }\edtext{\skt{घृत\-ङ्का\-म\-मु\-दा\-रम्वा}}{\Afootnote{A; \skt{घृत\-ङ्का\-म\-मु\-दा\-रम्वा}\ K}}\skt{ विरेचनं । }\pend
\pstart\versepar \skt{}\edtext{\skt{निरूहयेद्वा}}{\Afootnote{\skt{निरूहयेद्वा}\ A K}}\skt{ }\edtext{\skt{मतिमां}}{\Afootnote{K; \skt{मतिमा}\ A}}\skt{ }\edtext{\skt{सुस्वि\-न्न\-न्त\-द\-ह\-न्नरः}}{\Afootnote{K; \skt{न्सुस्विन्नं तदहरि नरम् ॥}\ A}}\skt{ ।}\pend

\pstart\versepar \ledsidenote{\passref{6.39.262}}\edtext{\skt{अजा\-व्यो\-श्च\-र्म\-रो\-माणि}}{\Afootnote{\skt{अजा\-व्यो\-श्च\-र्म\-रो\-माणि}\ A K}}\skt{ वचा कुष्ठम्पलङ्कषा । }\edtext{\skt{जाथपीगन्ध}}{\Afootnote{K\textsuperscript{pc}; \emph{om.}\ A K\textsuperscript{ac}}}\skt{ }\pend
\pstart\versepar \skt{}\edtext{\skt{निम्वपत्राणि}}{\Afootnote{K; \skt{निम्बपत्रं म}\ A}}\skt{ }\edtext{\skt{मधु\-क\-धू\-प\-नार्थे}}{\Afootnote{K; \skt{धुयुतं धूपनं त}\ A}}\skt{ }\edtext{\skt{प्रदापयेत्}}{\Afootnote{K; \skt{स्य दापयेत् ॥}\ A}}\skt{ ।}\pend

\pstart\versepar \ledsidenote{\passref{6.39.263}}\skt{बैडालम्वा }\edtext{\skt{सकृद्योज्य}}{\Afootnote{K; \skt{शकृद्योज्यं}\ A}}\skt{ वेपमानस्य }\edtext{\skt{धूपने}}{\Afootnote{K; \skt{धूपनम् ।}\ A}}\skt{ । }\pend
\pstart\versepar \skt{पिप्पली सैन्धवन्तैलं }\edtext{\skt{मैना\-ह्वा\-चे\-क्ष\-ण\-ञ्जनं}}{\Afootnote{\skt{नेपाली चेक्षणाञ्जनम् ॥}\ A; \skt{मैना\-ह्वा\-चे\-क्ष\-ण\-ञ्जनᳵ ॥}\ K}}\skt{ ॥}\pend

\pstart\versepar \ledsidenote{\passref{6.39.265}}\edtext{\skt{भूत\-वि\-द्या\-स\-मु\-दि\-ष्टै\-र्व\-न्धा\-वे\-श\-न\-पू\-जनैः}}{\Afootnote{\skt{भूत\-वि\-द्या\-स\-मु\-द्दि\-ष्टै\-र्ब\-न्धा\-वे\-श\-न\-पू\-जनैः ।}\ A; \skt{भूत\-वि\-द्या\-स\-मु\-दि\-ष्टै\-र्व\-न्धा\-वे\-श\-न\-पू\-जनैः ।}\ K}}\skt{ । }\pend
\pstart\versepar \skt{}\edtext{\edtext{\skt{जये\-द्भू\-ता\-भि\-श\-ङ्को\-त्थ\-म्वि\-ज्ञा\-ना\-द्यैश्च}}{\lemma{\skt{जये\-द्भू\-ता\-भि\-श\-ङ्को\-त्थम्}}\Afootnote{K; \skt{जये\-द्भू\-ता\-भि\-ष\-ङ्गोत्थं}\ A}}}{\lemma{\skt{विज्ञानाद्यैश्च}}\Afootnote{K; \skt{विज्ञानाद्यैश्च}\ A}}\skt{ मानसं ।}\pend

\pstart \ledsidenote{\passref{6.39.266}}\edtext{\skt{अभि\-शा\-पा\-भि\-चा\-रोत्थौ}}{\Afootnote{K\textsuperscript{pc}; \skt{श्रमक्षयोत्थे भुञ्जीत घृताभ्यक्तो रसौदनम् । अभि\-शा\-पा\-भि\-चा\-रोत्थौ}\ A; \skt{अभि\-शा\-पा\-भि\-चात्थौ}\ K\textsuperscript{ac}}}\skt{ ज्वरौ होमादिना जयेत् ।}\pend

\pstart \ledsidenote{\passref{6.39.267}}\edtext{\skt{दान\-स्व\-स्त्य\-य\-ना\-ति\-थ्यैं\-रु\-त्पा\-त\-ग्र\-ह\-पी\-डितं}}{\Afootnote{\skt{दान\-स्व\-स्त्य\-य\-ना\-ति\-थ्यै\-रु\-त्पा\-त\-ग्र\-ह\-पी\-डि\-तम् । अभिघातज्वरे कुर्या\-त्क्रि\-या\-मु\-ष्ण\-वि\-व\-र्जि\-ताम् ॥}\ A; \skt{दान\-स्व\-स्त्य\-य\-ना\-ति\-थ्यैं\-रु\-त्पा\-त\-ग्र\-ह\-पी\-डितᳵ ॥}\ K}}\skt{ ॥}\pend

\pstart \ledsidenote{\passref{6.39.268}}\edtext{\skt{ओष\-धी\-ग\-न्ध\-वि\-षजौ}}{\Afootnote{K; \skt{कषायमधुरां स्निग्धां यथा\-दो\-ष\-म\-थापि वा । औष\-धी\-ग\-न्ध\-वि\-षजौ}\ A}}\skt{ }\edtext{\skt{विष\-पी\-त\-प्र\-सा\-धनैः}}{\Afootnote{K; \skt{विष\-पि\-त्त\-प्र\-सा\-धनैः ॥}\ A}}\skt{ ।}\pend

\pstart\versepar \ledsidenote{\passref{6.39.269}}\skt{जयेत्कषायञ्च हितं सर्व\-ग\-न्ध\-कृ\-त\-न्तथा । }\pend
\pstart\versepar \skt{}\edtext{\skt{निम्व\-दा\-रु\-क\-षा\-यम्वा}}{\Afootnote{K; \skt{निम्ब\-दा\-रु\-क\-षायं वा}\ A}}\skt{ हितं }\edtext{\skt{सौमनसन्तथा}}{\Afootnote{K; \skt{सौमनसं यथा ॥}\ A}}\skt{ ।}\pend

\pstart\versepar \ledsidenote{\passref{6.39.270}}\edtext{\skt{यवान्नविकृती}}{\Afootnote{K; \skt{यवान्नविकृतिः}\ A}}\skt{ }\edtext{\skt{सर्पर्मद्यञ्च}}{\Afootnote{\skt{सर्पिर्मद्यं च}\ A; \skt{सर्पर्मद्यञ्च}\ K}}\skt{ विषमे }\edtext{\skt{हितं}}{\Afootnote{A; \skt{हितᳵ ।}\ K}}\skt{ । }\pend
\pstart\versepar \skt{}\edtext{\skt{सम्पू\-ज\-ये\-द्द्वि\-जा\-ङ्गाश्च}}{\Afootnote{\skt{संपू\-ज\-ये\-द्वि\-जा\-न्गाश्च}\ A; \skt{सम्पू\-ज\-ये\-द्द्वि\-जा\-ङ्गाश्च}\ K}}\skt{ }\edtext{\skt{देव\-मी\-शा\-न\-म\-म्विकां}}{\Afootnote{\skt{देव\-मी\-शा\-न\-म\-म्बि\-काम् ॥}\ A; \skt{देव\-मी\-शा\-न\-म\-म्विकाᳵ ।}\ K}}\skt{ ।}\pend

\pstart\versepar \ledsidenote{\passref{6.39.271}}\edtext{\skt{कफ\-वा\-तो\-त्थ\-यो\-रे\-न\-ञ्ज्व\-रयोः}}{\lemma{\skt{एनञ्ज्वरयोः}}\Afootnote{K; \skt{चापि ज्वरयोः}\ A}}\skt{ }\edtext{\skt{शीतपीडिते}}{\Afootnote{K; \skt{शीतपीडितम् ।}\ A}}\skt{ । }\pend
\pstart\versepar \skt{}\edtext{\skt{दिहेदुष्णेण}}{\Afootnote{\skt{दिह्यादुष्णेन}\ A; \skt{दिहेदुष्णेण}\ K}}\skt{ वर्गेण }\edtext{\skt{परिश्चोष्णो}}{\Afootnote{\skt{परश्चोष्णो}\ A; \skt{परिश्चोष्णो}\ K}}\skt{ विधिर्हितः ।}\pend

\pstart\versepar \ledsidenote{\passref{6.39.272}}\edtext{\edtext{\skt{सिञ्चे\-त्सु\-स्यो\-ष्णै\-र\-थवा}}{\lemma{\skt{सिञ्चेत्सुस्योष्णैर्}}\Afootnote{K; \skt{सिञ्चेत्कोष्णैरा}\ A}}}{\lemma{\skt{अथवा}}\Afootnote{K; \skt{रनाल}\ A}}\skt{ }\edtext{\skt{शुक्त\-गो\-मू\-त्र\-म\-स्तुभिः}}{\Afootnote{A; \skt{शुक्त\-गो\-मू\-त्र\-म\-स्तुभिः ।}\ K}}\skt{ । }\pend
\pstart\versepar \skt{}\edtext{\skt{दिह्या\-त्प\-ला\-शै\-स्त\-रु\-णै\-स्सु\-र\-सा\-र्ज\-क\-शि\-ग्रुजैः}}{\lemma{\skt{तरुणैस्}}\Afootnote{K; \skt{पिष्टैर्वा}\ A}}\skt{ ॥}\pend

\pstart\versepar \ledsidenote{\passref{6.39.273}}\skt{क्षारतैलेन वाभ्य\-ङ्ग\-स्स\-शु\-क्तेन विधीयते। }\pend
\pstart\versepar \skt{}\edtext{\edtext{\skt{पान\-मा\-र\-ग्व\-धा\-देश्च}}{\lemma{\skt{पानम्}}\Afootnote{K; \skt{पान}\ A}}}{\lemma{\skt{आरग्वधादेश्च}}\Afootnote{K; \skt{मारग्वधादेश्च}\ A}}\skt{ क्वथितस्य विशेषतः ।}\pend

\pstart \ledsidenote{\passref{6.39.273a}}\skt{न्दितष्णसकृच्छ्र गुधि\-ता\-नि\-ल\-त्वात् ।}\pend

\pstart\versepar \ledsidenote{\passref{6.39.273b}}\skt{जीर्ण्णासने कुप्यति वात\-गु\-ल्म\-वि\-दा\-ह\-काले श्वसनस्य पित्तात् । }\pend
\pstart\versepar \skt{गुल्मः कफामूर्च्छति भुक्तमात्रे सच्वासुवेलासु च सन्निपाते ॥}\pend

\pstart\versepar \ledsidenote{\passref{6.39.273c}}\skt{स्त्रीणान्तु यः शोषि\-त\-को\-प\-गुल्मः स योनिमूलं तु मूलं करोति ॥ }\pend
\pstart\versepar \skt{योन्या परि\-स्रा\-व\-म\-नि\-ष्ट\-गंन्धं विशेषलिङ्गानि भवन्ति चैषां ॥}\pend

\pstart\versepar \ledsidenote{\passref{6.39.273d}}\skt{वातगुल्मार्दितं स्निग्धं युक्तां स्नेहविरेचनैः । }\pend
\pstart\versepar \skt{उपाचरेद्यथा कंलां निरूहैः सानुवासनैः ॥}\pend

\pstart\versepar \ledsidenote{\passref{6.39.273e}}\skt{पित्तगुल्मादितं पीतं काकोल्यादि कृतं घृतं }\pend
\pstart\versepar \skt{विरि\-क्त\-म\-पु\-रै\-यो\-गै\-र्नि\-रू\-है\-स्स\-मु\-पा\-च\-रेत् ॥}\pend

\pstart\versepar \ledsidenote{\passref{6.39.273f}}\skt{श्लेष्मगुर्ल्मार्दितं पिप्पल्यादि घृतेन तु । }\pend
\pstart\versepar \skt{तीक्ष्णैर्विरिक्तं तद्भावैर्निरूहैः समुपक्रमेत् ॥}\pend

\pstart\versepar \ledsidenote{\passref{6.39.273g}}\skt{सन्निपातोत्थिते गुल्मे त्रिदोषघ्ना विधिहितः । }\pend
\pstart\versepar \skt{पित्त\-व\-द्र\-क्त\-गु\-ल्मिन्याः}\pend

\pstart\versepar \ledsidenote{\passref{6.39.273h}}\skt{भ्यङ्गस्सशुक्तेन विधीयते । }\pend
\pstart\versepar \skt{पान\-मा\-र\-ग्व\-धा\-देश्च क्वथितस्य विशेषतः ।}\pend

\pstart\versepar \ledsidenote{\passref{6.39.274}}\edtext{\skt{अव\-गा\-ह\-सु\-खो\-ष्णश्च}}{\Afootnote{K; \skt{अवगाहः सुखोष्णश्च}\ A}}\skt{ }\edtext{\skt{वातघ्नं}}{\Afootnote{K; \skt{वातघ्न}\ A}}\skt{ }\edtext{\skt{क्वाथयोजिते}}{\Afootnote{K; \skt{क्वाथयोजितः ।}\ A}}\skt{ । }\pend
\pstart\versepar \skt{जित्वा शीतं }\edtext{\skt{क्रमैरेवं}}{\Afootnote{\skt{क्रमैरेभिः}\ A; \skt{क्रमैरेवं}\ K}}\skt{ }\edtext{\skt{सुखो\-ष्ण\-ज\-ल\-से\-चितं}}{\Afootnote{A; \skt{सुखो\-ष्ण\-ज\-ल\-से\-चितᳵ ।}\ K}}\skt{ ।}\pend

\pstart\versepar \ledsidenote{\passref{6.39.275}}\edtext{\skt{प्राचा\-रौ\-र्ण्णि\-क\-का\-र्पा\-स\-कौ\-षे\-यां\-व\-र\-सं\-भृतं}}{\Afootnote{K; \skt{प्रवे\-श्यौ\-र्णि\-क\-का\-र्पा\-स\-कौ\-शे\-या\-म्ब\-र\-सं\-वृ\-तम् ।}\ A}}\skt{ । }\pend
\pstart\versepar \skt{}\edtext{\skt{शाय\-ये\-ङ्ग्ला\-न\-दे\-हम्वा}}{\Afootnote{K; \skt{शाय\-ये\-त्क्षा\-म\-देहं च}\ A}}\skt{ }\edtext{\skt{काला\-गु\-रु\-वि\-भू\-षितं}}{\Afootnote{\skt{काल\-गु\-रु\-वि\-भू\-षि\-तम् ॥}\ A; \skt{काला\-गु\-रु\-वि\-भू\-षितᳵ ।}\ K}}\skt{ ।}\pend

\pstart\versepar \ledsidenote{\passref{6.39.276}}\edtext{\skt{पयो\-ध\-रा\-कु\-शला}}{\Afootnote{K; \skt{स्तनाढ्या रूपसंपन्नाः कुशला}\ A}}\skt{ नवयौवनाः । }\pend
\pstart\versepar \skt{}\edtext{\edtext{\skt{गात्रै\-र्भ\-जे\-र\-न्प्र\-मदाः}}{\lemma{\skt{गात्रैर्}}\Afootnote{K; \emph{om.}\ A}}}{\lemma{\skt{भजेरन्प्रमदाः}}\Afootnote{K; \skt{भजेयुः प्रमदा}\ A}}\skt{ }\edtext{\skt{शीतदैन्यापहाः}}{\Afootnote{K; \skt{गात्रैः शीतदैन्यापहाः}\ A}}\skt{ शुभाः ।}\pend

\pstart \ledsidenote{\passref{6.39.280}}\edtext{\skt{प्रह्लादेनञ्चास्य}}{\lemma{\skt{प्रह्लादेनञ्}}\Afootnote{\skt{गाढ\-मा\-लि\-ङ्ग\-ये\-युस्तं तरुं वनलता इव । प्रह्लादं}\ A; \skt{प्रह्लादनञ्}\ K}}\skt{ विज्ञाय }\edtext{\edtext{\skt{तास्त्री\-र\-प\-न\-ये\-त्पुनः}}{\lemma{\skt{त्रीर्}}\Afootnote{K; \skt{स्त्री}\ A}}}{\lemma{\skt{अपनयेत्पुनः}}\Afootnote{K; \skt{रपनयेत्पुनः ॥}\ A}}\skt{ ।}\pend

\pstart \ledsidenote{\passref{6.39.281}}\edtext{\skt{भोज\-ये\-द्धि\-त\-म\-न्नञ्च}}{\Afootnote{K; \skt{तासा\-म\-ङ्ग\-प\-रि\-ष्व\-ङ्ग\-नि\-वा\-रु\-त\-हि\-म\-ज्व\-रम् । भोज\-ये\-द्धि\-त\-मन्नं च}\ A}}\skt{ }\edtext{\skt{तथा}}{\Afootnote{K; \skt{यथा}\ A}}\skt{ }\edtext{\skt{सुख\-म\-वा\-प्नु\-यात्}}{\Afootnote{\skt{सुख\-म\-वा\-प्रु\-यात् ॥}\ A; \skt{सुख\-म\-वा\-प्नु\-यात् ।}\ K}}\skt{ ।}\pend

\pstart\versepar \ledsidenote{\passref{6.39.282}}\skt{दाहाभिभूते तु }\edtext{\skt{विधि\-ङ्कु\-र्या\-द्दा\-ह\-वि\-ना\-शनीं}}{\Afootnote{\skt{विधिं कुर्या\-द्दा\-ह\-वि\-ना\-श\-नम् ।}\ A; \skt{विधि\-ङ्कु\-र्या\-द्दा\-ह\-वि\-ना\-शनीᳵ ।}\ K}}\skt{ । }\pend
\pstart\versepar \skt{}\edtext{\skt{मधु\-फा\-नि\-त\-मि\-श्रेण}}{\Afootnote{K; \skt{मधु\-फ\-णि\-त\-यु\-क्तेन}\ A}}\skt{ }\edtext{\skt{निम्व\-प\-त्रा\-म्भ\-सापि}}{\Afootnote{K; \skt{निम्ब\-प\-त्रा\-म्भ\-सापि}\ A}}\skt{ वा ।}\pend

\pstart\versepar \ledsidenote{\passref{6.39.283}}\edtext{\skt{दाह\-ज्व\-रा\-र्त्त\-म्म\-तिमां}}{\Afootnote{K; \skt{दाहज्वरार्तं मतिमा}\ A}}\skt{ }\edtext{\skt{वाम\-ये\-त्क्षि\-प्र\-मेव तु}}{\Afootnote{\skt{न्वाम\-ये\-त्क्षि\-प्र\-मेव च ।}\ A; \skt{वाम\-ये\-त्क्षि\-प्र\-मेव तु ।}\ K}}\skt{ । }\pend
\pstart\versepar \skt{}\edtext{\skt{शत\-धौ\-त\-घृ\-ता\-भ्य\-क्त\-न्दि\-हे\-द्वा\-य\-व\-स\-क्तुभिः}}{\Afootnote{\skt{शत\-धौ\-त\-घृ\-ता\-भ्यक्तं दिह्याद्वा यवशक्तुभिः ॥}\ A; \skt{शत\-धौ\-त\-घृ\-ता\-भ्य\-क्त\-न्दि\-हेद्वा यवसक्तुभिः ।}\ K}}\skt{ ।}\pend

\pstart \ledsidenote{\passref{6.39.284}}\edtext{\skt{कोला\-म\-ल\-क\-सं\-यु\-क्तै\-र्धा\-न्या\-म्लेन}}{\Afootnote{\skt{कोला\-म\-ल\-क\-सं\-युक्तैः शुक्त\-धा\-न्या\-म्ल\-सं\-युतैः । अम्ल}\ A; \skt{कोला\-म\-ल\-क\-सं\-यु\-क्तै\-र्धा\-न्या\-म्लेन}\ K}}\skt{ }\edtext{\skt{च}}{\Afootnote{K; \skt{पिष्टैः सुशीतैश्च}\ A}}\skt{ }\edtext{\skt{वुद्धिमां}}{\Afootnote{\skt{फेनि\-ला\-प\-ल्ल\-वै\-स्तथा ॥}\ A; \skt{वुद्धिमाᳵ ।}\ K}}\skt{ ।}\pend

\pstart\versepar \ledsidenote{\passref{6.39.285}}\skt{अम्ल\-पि\-ष्टै\-स्सु\-शी\-तैर्वा पला\-श\-त\-रु\-जै\-र्दि\-हेत् । }\pend
\pstart\versepar \skt{वद\-री\-प\-ल्ल\-वो\-त्थेन फेने\-ना\-रि\-ष्ट\-कस्य }\edtext{\skt{वा}}{\Afootnote{K; \skt{च ॥}\ A}}\skt{ ।}\edtext{}{\lemma{\skt{दिहेत् । वद\-री\-प\-ल्ल\-वो\-त्थेन}}\Afootnote{K; \skt{दिहेत् । बद\-री\-प\-ल्ल\-वो\-त्थेन}\ A}}\pend

\pstart \ledsidenote{\passref{6.39.286}}\edtext{\skt{यवार्धकुडवं}}{\Afootnote{K; \skt{लिप्तेऽङ्गे दाहतृण्मूर्च्छाः प्रशाम्यन्ति च सर्वशः । यवार्धकुडवं}\ A}}\skt{ पिष्ट्वा मंजि\-ष्ठा\-र्ध\-प\-ल\-न्तथा ।}\pend

\pstart\versepar \ledsidenote{\passref{6.39.287}}\edtext{\skt{अंव्ल\-प्र\-स्थ\-श\-तो\-न्मिश्रं}}{\Afootnote{K; \skt{अम्ल\-प्र\-स्थ\-श\-तो\-न्मिश्रं}\ A}}\skt{ तैल\-प्र\-स्थ\-म्वि\-पा\-च\-येत् । }\pend
\pstart\versepar \skt{एत\-त्प्र\-ह्ला\-द\-न\-न्तैलं }\edtext{\skt{ज्वर\-दा\-ह\-वि\-ना\-शनं}}{\Afootnote{A; \skt{ज्वर\-दा\-ह\-वि\-ना\-शनᳵ ।}\ K}}\skt{ ।}\pend

\pstart\versepar \ledsidenote{\passref{6.39.288}}\edtext{\edtext{\skt{न्यग्रोधादिर्गणो}}{\lemma{\skt{न्यग्रोधादिर्}}\Afootnote{K; \skt{न्यग्रोधादि}\ A}}\skt{ यश्च}}{\lemma{\skt{गणो यश्च}}\Afootnote{K; \skt{र्गणो यस्तु}\ A}}\skt{ काकोल्यादिश्च यः }\edtext{\skt{स्मृतः}}{\Afootnote{K; \skt{गणः ।}\ A}}\skt{ । }\pend
\pstart\versepar \skt{}\edtext{\skt{उत्पलादिर्गणो यश्च}}{\Afootnote{K\textsuperscript{pc}; \skt{उत्पलादिर्गणो यस्तु}\ A; \skt{उत्प\-ला\-दि\-श्च\-र्गणो यश्च}\ K\textsuperscript{ac}}}\skt{ }\edtext{\edtext{\skt{पिष्टैर्वातैः}}{\lemma{\skt{पिष्टैर्}}\Afootnote{K; \skt{पिष्ट}\ A}}}{\lemma{\skt{वातैः}}\Afootnote{K; \skt{र्वा तैः}\ A}}\skt{ प्रलेपयेत् ।}\pend

\pstart\versepar \ledsidenote{\passref{6.39.289}}\edtext{\skt{तत्क\-षा\-या\-म्ल\-सं\-युक्ता}}{\Afootnote{K; \skt{तत्क\-षा\-या\-म्ल\-सं\-सिद्धाः}\ A}}\skt{ स्नेहाश्चाभ्यञ्जने हिताः । }\pend
\pstart\versepar \skt{तेषां शीतकषाये वा }\edtext{\skt{दाहा\-न्त\-म\-व\-गा\-ह\-येत्}}{\lemma{\skt{दाहान्तम्}}\Afootnote{K; \skt{दाहार्तम्}\ A}}\skt{ ।}\pend

\pstart\versepar \ledsidenote{\passref{6.39.290}}\skt{दाहवेगे त्वतिक्रान्ते }\edtext{\skt{तस्मा\-दु\-द्धृ\-त्य\-मा\-नवं}}{\Afootnote{\skt{तस्मादुद्धृत्य मानवम् ।}\ A; \skt{तस्मा\-दु\-द्धृ\-त्य\-मा\-नवᳵ ।}\ K}}\skt{ । }\pend
\pstart\versepar \skt{}\edtext{\skt{प्रदि\-ह्या\-त्छ\-त\-धौ\-तेन}}{\Afootnote{K; \skt{परिषिच्याम्बुधिः शीतैः प्रलि\-म्पे\-च्च\-न्द\-ना\-दिभिः ॥}\ A}}\skt{ }\edtext{\skt{सर्पिषा}}{\Afootnote{K; \emph{om.}\ A}}\skt{ }\edtext{\skt{चन्दनेन}}{\Afootnote{K; \emph{om.}\ A}}\skt{ }\edtext{\skt{वा}}{\Afootnote{K; \emph{om.}\ A}}\skt{ ।}\pend

\pstart\versepar \ledsidenote{\passref{6.39.291}}\skt{ग्लानम्वा }\edtext{\skt{दीन\-म\-न\-स\-मा\-श्लि\-षे\-यु\-र्व\-रा\-ङ्गनाः}}{\Afootnote{\skt{दीन\-म\-न\-स\-मा\-श्लि\-षे\-यु\-र्व\-रा\-ङ्गनाः ।}\ A; \skt{दीन\-म\-न\-स\-मा\-श्लि\-षे\-यु\-र्व\-रा\-ङ्गनाः ।}\ K}}\skt{ । }\pend
\pstart\versepar \skt{}\edtext{\skt{दुकू\-ल\-क्षौ\-म\-स\-म्वी\-ता\-श्च\-न्द\-ना\-र्द्र\-प\-यो\-धराः}}{\Afootnote{\skt{पेल\-व\-क्षौ\-म\-सं\-वी\-ता\-श्च\-न्द\-ना\-र्द्र\-प\-यो\-धराः ॥}\ A; \skt{दुकू\-ल\-क्षौ\-म\-स\-म्वी\-ता\-श्च\-न्द\-ना\-र्द्र\-प\-यो\-धराः ।}\ K}}\skt{ ।}\pend

\pstart\versepar \ledsidenote{\passref{6.39.292}}\edtext{\skt{विभ्रत्यः}}{\Afootnote{K; \skt{बिभ्रत्यो}\ A}}\skt{ }\edtext{\skt{स्रग्विचित्रास्तु}}{\Afootnote{K; \skt{ऽब्जस्रजश्चित्रा}\ A}}\skt{ }\edtext{\skt{मणि\-हा\-र\-वि\-भू\-षिताः}}{\Afootnote{K; \skt{मणि\-र\-त्न\-वि\-भू\-षिताः ।}\ A}}\skt{ । }\pend
\pstart\versepar \skt{}\edtext{\skt{भजेयुस्तं}}{\Afootnote{K; \skt{भजेयुस्ताः}\ A}}\skt{ }\edtext{\skt{तनैः}}{\Afootnote{K; \skt{स्तनैः}\ A}}\skt{ }\edtext{\skt{शीतैः\-स्पृ\-शं\-त्यां\-वु\-रु\-है\-स्सुखैः}}{\Afootnote{K; \skt{शीतैः स्पृशन्त्योऽम्बुरुहैः सुखैः ॥}\ A}}\skt{ ।}\pend

\pstart\versepar \ledsidenote{\passref{6.39.293}}\edtext{\skt{प्रह्लादनञ्च}}{\Afootnote{K; \skt{प्रह्लादं चास्य}\ A}}\skt{ }\edtext{\skt{विज्ञाय ता}}{\Afootnote{K; \skt{विज्ञाय ताः}\ A}}\skt{ }\edtext{\skt{स्त्रीर\-प\-न\-ये\-त्पुनः}}{\Afootnote{\skt{स्त्रीर\-प\-न\-ये\-त्पुनः}\ A K}}\skt{ । }\pend
\pstart\versepar \skt{}\edtext{\skt{भोज\-ये\-द्धि\-त\-मं\-नञ्च}}{\Afootnote{\skt{हितं च भोजयेदन्नं}\ A; \skt{भोज\-ये\-द्धि\-त\-मं\-नञ्च}\ K}}\skt{ तथा }\edtext{\skt{सुख\-म\-वा\-प्नु\-यात्}}{\Afootnote{\skt{प्नोति सुखं महत् ॥}\ A; \skt{सुख\-म\-वा\-प्नु\-यात् ॥}\ K}}\skt{ ॥}\pend

\pstart\versepar \ledsidenote{\passref{6.39.294}}\edtext{\skt{गम\-ये\-त्पि\-त्त\-मे\-वादौ}}{\lemma{\skt{गमयेत्पित्तम्}}\Afootnote{K; \skt{पित्तज्वरोक्तं शमनं विरेकोऽन्यद्धितं च यत् । निर्हरेत्पित्तम्}\ A}}\skt{ }\edtext{\skt{ज्वरेषु}}{\Afootnote{K; \skt{दोषेषु}\ A}}\skt{ समवायिषु ॥ }\pend
\pstart\versepar \skt{}\edtext{\edtext{\skt{दिर्नि\-म्वा\-र\-त\-र\-न्तद्धि}}{\lemma{\skt{दिर्नि\-म्वा\-र\-त\-रन्}}\Afootnote{K; \emph{om.}\ A}}}{\lemma{\skt{तद्धि}}\Afootnote{\emph{om.}\ A; \skt{तद्धि}\ K}}\skt{ }\edtext{\skt{ज्व}}{\Afootnote{K; \emph{om.}\ A}}\pend

\pstart\versepar \ledsidenote{\passref{6.39.294a}}\edtext{\skt{भार्गी\-शा\-रि\-वा\-द्व\-य\-चि\-त्रकैः}}{\Afootnote{K; \emph{om.}\ A}}\skt{ । }\pend
\pstart\versepar \skt{}\edtext{\skt{दन्ती च}}{\Afootnote{K; \emph{om.}\ A}}\skt{ }\edtext{\skt{श्रेयसी}}{\Afootnote{K; \emph{om.}\ A}}\skt{ }\edtext{\skt{पाठा}}{\Afootnote{K; \emph{om.}\ A}}\skt{ }\edtext{\skt{एतः}}{\Afootnote{K; \emph{om.}\ A}}\skt{ }\edtext{\edtext{\skt{प्रस्थ\-घृ\-त\-म्प\-चेत्}}{\lemma{\skt{प्रस्थघृतम्}}\Afootnote{K; \emph{om.}\ A}}}{\lemma{\skt{पचेत्}}\Afootnote{K; \emph{om.}\ A}}\skt{ ॥}\pend

\pstart\versepar \ledsidenote{\passref{6.39.294b}}\edtext{\skt{मूत्राशकृद्रसे}}{\Afootnote{K; \emph{om.}\ A}}\skt{ }\edtext{\skt{तुल्ये}}{\Afootnote{K; \emph{om.}\ A}}\skt{ }\edtext{\skt{दधिमूत्र}}{\Afootnote{K; \emph{om.}\ A}}\skt{ }\edtext{\skt{यथायुतं}}{\Afootnote{K; \emph{om.}\ A}}\skt{ । }\pend
\pstart\versepar \skt{}\edtext{\edtext{\skt{पञ्चगव्यमिति}}{\lemma{\skt{पञ्चगव्यम्}}\Afootnote{K; \emph{om.}\ A}}}{\lemma{\skt{इति}}\Afootnote{K; \emph{om.}\ A}}\skt{ }\edtext{\edtext{\edtext{\skt{ख्यात\-मे\-त\-त्स\-र्पि\-श्च\-तु\-र्थ\-कम्}}{\lemma{\skt{ख्यातम्}}\Afootnote{K; \emph{om.}\ A}}}{\lemma{\skt{एतत्सर्पिश्}}\Afootnote{K; \emph{om.}\ A}}}{\lemma{\skt{चतुर्थकम्}}\Afootnote{\emph{om.}\ A; \skt{चतुर्थकᳵ ।}\ K}}\skt{ ।}\pend

\pstart\versepar \ledsidenote{\passref{6.39.294c}}\edtext{\skt{गुल्म\-पां\-डा\-म\-ये\-कासं}}{\Afootnote{K; \emph{om.}\ A}}\skt{ }\edtext{\edtext{\skt{काम\-ला\-मु\-द\-राणि}}{\lemma{\skt{कामलाम्}}\Afootnote{K; \emph{om.}\ A}}\skt{ च}}{\lemma{\skt{उदराणि च}}\Afootnote{K; \emph{om.}\ A}}\skt{ । }\pend
\pstart\versepar \skt{}\edtext{\skt{अर्शांसि}}{\Afootnote{K; \emph{om.}\ A}}\skt{ }\edtext{\edtext{\edtext{\skt{शोष\-मु\-न्मा\-द\-म\-प\-स्मा\-र\-भ\-ग\-न्द\-रम्}}{\lemma{\skt{शोषम्}}\Afootnote{K; \emph{om.}\ A}}}{\lemma{\skt{उन्मादम्}}\Afootnote{K; \emph{om.}\ A}}}{\lemma{\skt{अप\-स्मा\-र\-भ\-ग\-न्द\-रम्}}\Afootnote{K; \emph{om.}\ A}}\skt{ ।}\pend

\pstart\versepar \ledsidenote{\passref{6.39.294d}}\edtext{\skt{ग्रहा न}}{\Afootnote{K; \emph{om.}\ A}}\skt{ }\edtext{\skt{लक्ष्मीपाध्मानं}}{\Afootnote{K; \emph{om.}\ A}}\skt{ }\edtext{\skt{पीतं}}{\Afootnote{K; \emph{om.}\ A}}\skt{ }\edtext{\skt{चैतद्व्यपोहति}}{\Afootnote{\emph{om.}\ A; \skt{चैतद्व्योहति ।}\ K}}\skt{ । }\pend
\pstart\versepar \skt{}\edtext{\skt{विना}}{\Afootnote{K; \emph{om.}\ A}}\skt{ }\edtext{\skt{वात्क्वाथकल्काभ्यां}}{\Afootnote{K; \emph{om.}\ A}}\skt{ }\edtext{\skt{पञ्चगाव्यं}}{\Afootnote{K; \emph{om.}\ A}}\skt{ }\edtext{\skt{घृतं}}{\Afootnote{K; \emph{om.}\ A}}\skt{ }\edtext{\skt{श्रितः}}{\Afootnote{\emph{om.}\ A; \skt{श्रितᳵ ।}\ K}}\skt{ ।}\pend

\pstart\versepar \ledsidenote{\passref{6.39.294e}}\edtext{\skt{तदपि}}{\Afootnote{\emph{om.}\ A; \skt{तदपि}\ K}}\skt{ }\edtext{\skt{ज्वर\-पा\-ण्डु\-त्व\-शो\-षा\-प\-स्मा\-र\-जि\-घृतं}}{\Afootnote{K; \emph{om.}\ A}}\skt{ । }\pend
\pstart\versepar \skt{}\edtext{\skt{त्रिफला}}{\Afootnote{K; \emph{om.}\ A}}\skt{ }\edtext{\skt{सप्त\-ला\-नि\-म्ब\-प\-टो\-ली\-वृ\-ष\-शा\-रिवा}}{\Afootnote{\emph{om.}\ A; \skt{सप्त\-ला\-नि\-म्व\-प\-टो\-ली\-वृ\-ष\-शा\-रिवा ।}\ K}}\skt{ ।}\pend

\pstart\versepar \ledsidenote{\passref{6.39.294f}}\edtext{\skt{कुट\-ज\-त्व\-ड्फ\-ला\-म्पाठा}}{\Afootnote{K; \emph{om.}\ A}}\skt{ }\edtext{\skt{सप्ता\-ह\-सु\-म\-ना\-वृता}}{\Afootnote{K; \emph{om.}\ A}}\skt{ । }\pend
\pstart\versepar \skt{}\edtext{\edtext{\skt{वेतसश्चन्दनं}}{\lemma{\skt{वेतसश्}}\Afootnote{K; \emph{om.}\ A}}}{\lemma{\skt{चन्दनं}}\Afootnote{K; \emph{om.}\ A}}\skt{ }\edtext{\skt{मूर्वा}}{\Afootnote{K; \emph{om.}\ A}}\skt{ }\edtext{\skt{कर\-ण्ड\-क\-टु\-रो\-हिणी}}{\Afootnote{K; \emph{om.}\ A}}\skt{ ।}\pend

\pstart\versepar \ledsidenote{\passref{6.39.294g}}\edtext{\skt{हरिद्रापिप्पलं}}{\Afootnote{K; \emph{om.}\ A}}\skt{ }\edtext{\skt{त्वैव}}{\Afootnote{K; \emph{om.}\ A}}\skt{ }\edtext{\skt{पिष्टैः}}{\Afootnote{K; \emph{om.}\ A}}\skt{ }\edtext{\skt{कर्षसमैः}}{\Afootnote{\emph{om.}\ A; \skt{कर्षसमै}\ K}}\skt{ पचेत् }\pend
\pstart\versepar \skt{जलद्रोणे }\edtext{\skt{घृतप्रस्थं}}{\Afootnote{K; \emph{om.}\ A}}\skt{ }\edtext{\skt{विष\-म\-ज्व\-र\-ना\-शनं}}{\Afootnote{\emph{om.}\ A; \skt{विषमं ज्वरनाशनᳵ ।}\ K}}\skt{ ।}\edtext{}{\lemma{\skt{पचेत्जलद्रोणे}}\Afootnote{K; \emph{om.}\ A}}\pend

\pstart\versepar \ledsidenote{\passref{6.39.294h}}\edtext{\skt{कासे}}{\Afootnote{K; \emph{om.}\ A}}\skt{ }\edtext{\skt{श्वासे}}{\Afootnote{K; \emph{om.}\ A}}\skt{ }\edtext{\skt{प्रतिश्याये}}{\Afootnote{K; \emph{om.}\ A}}\skt{ }\edtext{\skt{रक्तपित्ते}}{\Afootnote{K; \emph{om.}\ A}}\skt{ }\edtext{\skt{भगन्दरे}}{\Afootnote{K; \emph{om.}\ A}}\skt{ । }\pend
\pstart\versepar \skt{}\edtext{\skt{पाण्डु\-त्व\-मे\-ह\-कु\-ष्टेषु}}{\Afootnote{K; \emph{om.}\ A}}\skt{ }\edtext{\skt{हृद्रौगे}}{\Afootnote{\emph{om.}\ A; \skt{हृद्रओगे}\ K}}\skt{ }\edtext{\skt{वातशोणिते}}{\Afootnote{K; \emph{om.}\ A}}\skt{ ।}\pend

\pstart\versepar \ledsidenote{\passref{6.39.294i}}\edtext{\skt{पित्त\-गु\-ल्मै\-क\-चौ\-रैव}}{\Afootnote{K; \emph{om.}\ A}}\skt{ }\edtext{\skt{कामलायाञ्च}}{\Afootnote{K; \emph{om.}\ A}}\skt{ }\edtext{\skt{पूजितं}}{\Afootnote{\emph{om.}\ A; \skt{पूजितᳵ ।}\ K}}\skt{ । }\pend
\pstart\versepar \skt{}\edtext{\skt{शीलयेत्मगुडां}}{\Afootnote{\emph{om.}\ A; \skt{शीलयेत्सगुडां}\ K}}\skt{ }\edtext{\edtext{\edtext{\skt{शुञ्ठी\-म\-भ\-या\-ञ्चापि}}{\lemma{\skt{शुञ्ठीम्}}\Afootnote{K; \emph{om.}\ A}}}{\lemma{\skt{अभयाञ्}}\Afootnote{K; \emph{om.}\ A}}}{\lemma{\skt{चापि}}\Afootnote{K; \emph{om.}\ A}}\skt{ }\edtext{\skt{पूजितः}}{\Afootnote{K; \emph{om.}\ A}}\skt{ ।}\pend

\pstart\versepar \ledsidenote{\passref{6.39.294j}}\edtext{\edtext{\skt{ह्यत\-प्र\-धा\-न\-पि\-त्तन्तु}}{\lemma{\skt{ह्य्}}\Afootnote{\emph{om.}\ A K}}}{\lemma{\skt{अत\-प्र\-धा\-न\-पि\-त्तन्तु}}\Afootnote{\emph{om.}\ A; \skt{ह्यत\-प्र\-धा\-न\-पि\-त्तन्तु}\ K}}\skt{ }\edtext{\skt{स्वस्थो}}{\Afootnote{K; \emph{om.}\ A}}\skt{ }\edtext{\skt{जनयति}}{\Afootnote{K; \emph{om.}\ A}}\skt{ }\edtext{\skt{ज्वरं}}{\Afootnote{\emph{om.}\ A; \skt{ज्वरᳵ ।}\ K}}\skt{ । }\pend
\pstart\versepar \skt{}\edtext{\edtext{\skt{पिबे\-दि\-क्षु\-र\-सा\-न्तत्र}}{\lemma{\skt{पिबे\-दि\-क्षु\-र\-सान्}}\Afootnote{\emph{om.}\ A; \skt{पिवे\-दि\-क्षु\-र\-सान्}\ K}}}{\lemma{\skt{तत्र}}\Afootnote{K; \emph{om.}\ A}}\skt{ }\edtext{\skt{शीतम्वा}}{\Afootnote{K; \emph{om.}\ A}}\skt{ }\edtext{\skt{शर्करोदकं}}{\Afootnote{K; \emph{om.}\ A}}\skt{ ।}\pend

\pstart \ledsidenote{\passref{6.39.319}}\edtext{\skt{शालि\-ष\-ष्टि\-क\-यो\-श्चान्नं}}{\lemma{\skt{चान्नं}}\Afootnote{K; \skt{अन्नम्}\ A}}\skt{ }\edtext{\edtext{\skt{क्षीरे\-णा\-प्लु\-त\-म\-श्नु\-यात्}}{\lemma{\skt{क्षीरेणाप्लुतम्}}\Afootnote{K; \emph{om.}\ A}}}{\lemma{\skt{अश्नुयात्}}\Afootnote{K; \skt{अश्नी\-या\-त्क्षी\-र\-सं\-प्लु\-तम् । कफ\-वा\-तो\-त्थ\-यो\-रेव स्वेदाभ्यङ्गौ प्रयोजयेत् ॥}\ A}}\skt{ ।}\pend

\pstart\versepar \ledsidenote{\passref{6.39.320}}\edtext{\skt{घृत}}{\Afootnote{K; \skt{घृतं}\ A}}\skt{ द्वादशरात्रात्तु देयं }\edtext{\skt{सर्वज्वरेषु वै}}{\Afootnote{K; \skt{सर्वज्वरेषु च ।}\ A}}\skt{ । }\pend
\pstart\versepar \skt{}\edtext{\skt{तेना\-न्त\-रे\-ना\-शय स्व}}{\Afootnote{K; \skt{तेना\-न्त\-रे\-णा\-शयं स्वं}\ A}}\skt{ गता दोषा भवन्ति हि ।}\pend

\pstart\versepar \ledsidenote{\passref{6.39.320.add1}}\edtext{\skt{तदेव}}{\Afootnote{\skt{तदेव}\ K}}\skt{ शून्यशिरसां सर्पि\-र्न्न\-स्त\-म्प्र\-यो\-ज\-येत् । }\pend
\pstart\versepar \skt{श्लेष्मेणस्तु समुत्क्त्रेशे धूमां }\edtext{\skt{व्रायो\-शि\-क\-म्पि\-बेत्}}{\lemma{\skt{पिबेत्}}\Afootnote{\skt{पिवेत् ।}\ K}}\skt{ ।}\edtext{}{\lemma{\skt{प्रयोजयेत् । श्लेष्मेणस्तु}}\Afootnote{\skt{प्रयोजयेत् । श्लेष्मणस्तु}\ K}}\pend

\pstart \ledsidenote{\passref{6.39.320.add2}}\skt{मूर्च्छायामिति निद्रायां कुर्या\-च्छी\-त\-वि\-व\-र्जितः ।}\pend

\pstart\versepar \ledsidenote{\passref{6.39.321}}\edtext{\skt{धातूं}}{\Afootnote{K; \skt{धानू}\ A}}\skt{ }\edtext{\skt{संक्षोभयेद्दोषो}}{\Afootnote{\skt{न्प्रक्षोभयन्दोषो}\ A; \skt{संक्षोभयेद्दोषो}\ K}}\skt{ मोक्षकाले }\edtext{\skt{वलीयसि}}{\Afootnote{K; \skt{बलीयते ।}\ A}}\skt{ । }\pend
\pstart\versepar \skt{तेन }\edtext{\skt{व्याकुलचित्तस्य}}{\Afootnote{K; \skt{व्याकुलचित्तस्}\ A}}\skt{ }\edtext{\skt{विसंज्ञ}}{\Afootnote{K; \emph{om.}\ A}}\skt{ }\edtext{\skt{इव}}{\Afootnote{K; \emph{om.}\ A}}\skt{ }\edtext{\skt{लक्ष्यते}}{\Afootnote{K; \skt{तु म्रियमाण इवेहते ॥}\ A}}\skt{ ।}\pend

\pstart\versepar \ledsidenote{\passref{6.39.322}}\edtext{\skt{स्वेदो}}{\Afootnote{K; \emph{om.}\ A}}\skt{ लघुत्वं शिरसः }\edtext{\skt{कण्डूपाको}}{\Afootnote{K; \skt{स्वेदो}\ A}}\skt{ }\edtext{\skt{मुखस्य च}}{\Afootnote{K; \skt{मुखमापाण्डु पाकि च ।}\ A}}\skt{ । }\pend
\pstart\versepar \skt{}\edtext{\edtext{\skt{क्ष्ववथुश्चान्न}}{\lemma{\skt{क्ष्ववथुश्}}\Afootnote{K; \skt{क्षवथुश्}\ A}}\skt{ लिप्सा च}}{\lemma{\skt{चान्न लिप्सा च}}\Afootnote{K; \skt{चान्नकाङ्क्षा च}\ A}}\skt{ }\edtext{\skt{ज्वरमुक्तस्य}}{\Afootnote{K; \skt{ज्वरमुक्तस्य}\ A}}\skt{ }\edtext{\skt{लक्षणं}}{\Afootnote{A; \skt{लक्षणᳵ ।}\ K}}\skt{ ।}\pend

\pstart\versepar \ledsidenote{\passref{6.39.324}}\edtext{\skt{व्यापि\-त्वा\-त्स\-र्व\-सं\-स्प\-र्शा\-त्कृ\-च्छ्र\-त्वा\-द\-न्त\-स\-म्भवा}}{\Afootnote{K; \skt{व्यापि\-त्वा\-त्स\-र्व\-सं\-स्प\-र्शा\-त्कृ\-च्छ्र\-त्वा\-द\-न्त\-सं\-भवा}\ A}}\skt{ । }\pend
\pstart\versepar \skt{}\edtext{\skt{दन्तको ह्येष}}{\lemma{\skt{दन्तको ह्य्}}\Afootnote{K; \skt{त् । अन्तको ह्य्}\ A}}\skt{ }\edtext{\skt{मन्याना}}{\Afootnote{K; \skt{भूतानां}\ A}}\skt{ }\edtext{\skt{ज्वर}}{\Afootnote{A; \skt{न्ज्वर}\ K}}\skt{ इत्युपदिश्यते }\edtext{\skt{इति}}{\Afootnote{K; \emph{om.}\ A}}\skt{ ॥}\pend

\endnumbering
\beginnumbering
\eledsection*{Adhyāya 40}
\pstart \ledsidenote{\passref{6.40.1}}\edtext{\skt{अथा\-तोऽ\-ती\-सा\-र\-प्र\-ति\-षे\-ध\-म\-ध्यायं}}{\lemma{\skt{अध्यायं}}\Afootnote{K; \emph{om.}\ A}}\skt{ }\edtext{\skt{व्याख्यास्यामः}}{\Afootnote{A; \skt{व्या व्याख्यास्यामः ॥}\ K}}\skt{ ॥}\pend

\pstart\versepar \ledsidenote{\passref{6.40.3}}\edtext{\skt{गुर्वन्ति}}{\Afootnote{K; \skt{गुर्वति}\ A}}\skt{ }\edtext{\skt{ष्यन्दि\-रु\-क्ष्णो\-ष्ण\-द्र\-व\-स्थू\-ला\-ति\-शी\-तलैः}}{\Afootnote{K; \skt{स्निग्ध\-रू\-क्षो\-ष्ण\-द्र\-व\-स्थू\-ला\-ति\-शी\-तलैः ।}\ A}}\skt{ । }\pend
\pstart\versepar \skt{}\edtext{\edtext{\skt{विरु\-द्धा\-ध्य\-श\-ना\-जी\-र्णै\-र्वि\-ष\-मै\-श्चापि}}{\lemma{\skt{विरु\-द्धा\-ध्य\-श\-ना\-जी\-र्णैर्}}\Afootnote{K; \skt{विरु\-द्धा\-ध्य\-श\-ना\-जीर्णै}\ A}}}{\lemma{\skt{विषमैश्}}\Afootnote{K; \skt{रसात्म्यैश्}\ A}}\skt{ भोजनैः ।}\pend

\pstart\versepar \ledsidenote{\passref{6.40.4}}\edtext{\skt{युक्तै}}{\Afootnote{K; \skt{स्नेहा\-द्यै\-र\-ति\-यु\-क्तैश्च मिथ्यायुक्तै}\ A}}\skt{ । }\pend
\pstart\versepar \skt{}\edtext{\edtext{\skt{शोका\-दु\-ष्टा\-म्वु\-पा\-नैर्वा}}{\lemma{\skt{शोका\-दु\-ष्टा\-म्वु\-पा\-नैर्}}\Afootnote{K; \skt{र्विषाद्भयात् । शोका\-द्दु\-ष्टा\-म्बु\-मद्या}\ A}}}{\lemma{\skt{वा}}\Afootnote{K; \skt{ति}\ A}}\skt{ }\edtext{\skt{तथा}}{\Afootnote{K; \skt{पाना}\ A}}\skt{ }\edtext{\skt{च}}{\Afootnote{K; \skt{त्सा}\ A}}\skt{ }\edtext{\skt{तुर्विपययुयैः}}{\Afootnote{K; \skt{त्म्यर्तुपर्ययात् ॥}\ A}}\skt{ ।}\pend

\pstart\versepar \ledsidenote{\passref{6.40.5}}\edtext{\edtext{\skt{जला\-भि\-र\-म\-णै\-र्वे\-ग\-वि\-घातैः}}{\lemma{\skt{जला\-भि\-र\-म\-णैर्}}\Afootnote{K; \skt{जलातिरमणै}\ A}}}{\lemma{\skt{वेगविघातैः}}\Afootnote{K; \skt{र्वेगविघातैः}\ A}}\skt{ }\edtext{\skt{क्रिमिकोपतः}}{\Afootnote{K; \skt{कृमिदोषतः ।}\ A}}\skt{ । }\pend
\pstart\versepar \skt{नृणां भवत्यतीसारो लक्षणं }\edtext{\skt{संप्रवक्ष्यते}}{\Afootnote{K; \skt{तस्य वक्ष्यते ॥}\ A}}\skt{ ।}\pend

\pstart\versepar \ledsidenote{\passref{6.40.6}}\edtext{\edtext{\skt{संश\-म्या\-पा\-त्वा\-तु\-र\-ग्नि\-प्र\-वद्धः}}{\lemma{\skt{संश\-म्या\-पा\-त्वा\-तुर्}}\Afootnote{K; \skt{संशम्यापांधातु}\ A}}}{\lemma{\skt{अग्निप्रवद्धः}}\Afootnote{K; \skt{रन्तः}\ A}}\skt{ }\edtext{\skt{शकृभिग्रो}}{\Afootnote{K; \skt{कृशानुं वर्चोमिश्रो}\ A}}\skt{ }\edtext{\skt{वायुनाधः}}{\Afootnote{K; \skt{मारुतेन}\ A}}\skt{ प्रणुन्नः । }\pend
\pstart\versepar \skt{}\edtext{\edtext{\edtext{\skt{सर\-त्व\-ती\-दा\-ति\-सा\-र\-तु\-मा\-हु\-र्व्याधिं}}{\lemma{\skt{सर\-त्व\-ती\-दा\-ति\-सा\-र\-तुम्}}\Afootnote{K; \skt{वृद्धोऽतीवाधः सरत्येष}\ A}}}{\lemma{\skt{आहुर्}}\Afootnote{K; \skt{यस्मा}\ A}}}{\lemma{\skt{व्याधिं}}\Afootnote{K; \skt{द्व्याधिं}\ A}}\skt{ घोरं }\edtext{\skt{दत्ति}}{\Afootnote{K; \skt{तं त्वतीसारमाहुः ॥}\ A}}\pend

\pstart\versepar \ledsidenote{\passref{6.40.7}}\skt{एकैकशः सर्वशश्चापि दोषैः शोकेनान्यः षष्ठ आमेन चोक्तः । }\pend
\pstart\versepar \skt{}\edtext{\skt{केचिदाहुस्त्वे}}{\Afootnote{K; \skt{केचित्प्राहुर्नै}\ A}}\skt{ }\edtext{\skt{कविशन्प्रकारं}}{\Afootnote{K; \skt{करूपप्रकारं}\ A}}\skt{ }\edtext{\skt{नेत्वे}}{\Afootnote{K; \skt{नैवे}\ A}}\skt{ }\edtext{\skt{तस्मि\-न्का\-शि\-रा\-ज\-स्त्व\-वो\-चत्}}{\Afootnote{K; \skt{त्येवं काशि\-रा\-ज\-स्त्व\-वो\-चत् ॥}\ A}}\skt{ ।}\pend

\pstart\versepar \ledsidenote{\passref{6.40.8}}\skt{दोषावस्था }\edtext{\skt{तस्य}}{\Afootnote{K; \skt{स्तस्य}\ A}}\skt{ नैकप्रकाराः । काले काले }\edtext{\skt{व्याधितानां}}{\Afootnote{K; \skt{व्याधितस्यो}\ A}}\skt{ }\edtext{\skt{भवन्ति}}{\Afootnote{K; \skt{द्भवन्ति ।}\ A}}\skt{ ॥ }\pend
\pstart\versepar \skt{}\edtext{\skt{हृन्ना\-भि\-पा\-यू\-द\-र\-व\-स्ति\-तो\-द\-गा\-त्रा\-व\-सा\-दा\-नि\-ल\-वे\-ग\-रोधः}}{\Afootnote{K; \skt{हृन्ना\-भि\-पा\-यू\-द\-र\-कु\-क्षि\-तो\-द\-गा\-त्रा\-व\-सा\-दा\-नि\-ल\-स\-न्नि\-रोधाः ॥}\ A}}\skt{ ।}\pend

\pstart\versepar \ledsidenote{\passref{6.40.9}}\edtext{\edtext{\edtext{\skt{विट्स\-ङ्ग\-मा\-ध्मा\-न\-म\-था\-वि\-पाको}}{\lemma{\skt{विट्सङ्गम्}}\Afootnote{K; \skt{विट्सङ्ग}\ A}}}{\lemma{\skt{आध्मानम्}}\Afootnote{K; \skt{आध्मान}\ A}}}{\lemma{\skt{अथाविपाको}}\Afootnote{K; \skt{मथाविपाको}\ A}}\skt{ }\edtext{\skt{भविष्यतस्य}}{\Afootnote{K; \skt{भविष्यतस्तस्य}\ A}}\skt{ पुरस्सराणि । }\pend
\pstart\versepar \skt{}\edtext{\edtext{\edtext{\skt{शूला\-र्दि\-तो\-ल्प\-म्व\-द्ग\-श\-म्स\-फेन}}{\lemma{\skt{शूलार्दितोल्पम्}}\Afootnote{K; \skt{शूला}\ A}}}{\lemma{\skt{वद्गशम्}}\Afootnote{K; \emph{om.}\ A}}\skt{ ते}}{\lemma{\skt{सफेन ते}}\Afootnote{K; \emph{om.}\ A}}\skt{ }\edtext{\skt{सानि\-ल\-मा\-रु\-तेन}}{\Afootnote{K; \skt{विष्टः सक्त\-मू\-त्रोऽ\-न्त्र\-कूजी स्रस्तापानः सन्न\-क\-ट्यू\-रु\-गङ्घः ॥}\ A}}\skt{ ।}\pend

\pstart \ledsidenote{\passref{6.40.10}}\edtext{\edtext{\edtext{\skt{मांसा\-द\-का\-र\-श\-वे\-दा\-न्वि\-कृ\-ष्ण\-मु\-ष्ण\-न्द्र\-वाश्य}}{\lemma{\skt{मांसा\-द\-का\-र\-श\-वे\-दा\-न्वि\-कृ\-ष्णम्}}\Afootnote{K; \skt{वर्चो मुञ्चत्यल्पमल्पं सफेनं रूक्षं श्यावं सा}\ A}}}{\lemma{\skt{उष्णन्}}\Afootnote{K; \emph{om.}\ A}}}{\lemma{\skt{द्रवाश्य}}\Afootnote{K; \emph{om.}\ A}}\skt{ }\edtext{\edtext{\skt{नीलोपमम्वा}}{\lemma{\skt{नीलोपमम्}}\Afootnote{K; \skt{निलं}\ A}}}{\lemma{\skt{वा}}\Afootnote{K; \skt{मारुतेन । दुर्गन्ध्युष्णं वेग\-व\-न्मां\-स\-तो\-य\-प्रख्यं भिन्नं स्विन्न\-दे\-होऽ\-ति\-ती\-क्ष्णम् ॥}\ A}}\skt{ ।}\pend

\pstart \ledsidenote{\passref{6.40.11ab}}\edtext{\skt{पित्ता\-त्पी\-त\-ह\-रितं}}{\Afootnote{K; \skt{पित्तात्पीतं नीलमा}\ A}}\skt{ लोहितं वा }\edtext{\skt{तृष्णां}}{\Afootnote{K; \skt{तृष्णा}\ A}}\skt{ }\edtext{\skt{मूर्च्छा\-दा\-ह\-पा\-को\-प\-पन्नः}}{\Afootnote{K; \skt{मूर्च्छा\-दा\-ह\-पा\-क\-ज्व\-रार्तः ।}\ A}}\skt{ ।}\pend

\pstart \ledsidenote{\passref{6.40.12ab}}\skt{शुक्लं सान्द्रं }\edtext{\skt{सकफं}}{\Afootnote{K; \emph{om.}\ A}}\skt{ श्लेष्मणा }\edtext{\skt{तु}}{\Afootnote{K; \skt{श्लेष्मयुक्तं भक्तद्वेषी}\ A}}\skt{ }\edtext{\skt{विश्रां गो}}{\Afootnote{K; \skt{निःस्वनं}\ A}}\skt{ }\edtext{\skt{हृष्टरोभा}}{\Afootnote{K; \skt{हृष्टरोमा ।}\ A}}\skt{ }\edtext{\skt{मनुष्यः}}{\Afootnote{K; \emph{om.}\ A}}\skt{ ।}\pend

\pstart\versepar \ledsidenote{\passref{6.40.13}}\edtext{\skt{सर्वात्मक}}{\Afootnote{K; \skt{सर्वोद्भूते}\ A}}\skt{ सर्व\-लि\-ङ्गो\-प\-पत्तिः }\edtext{\edtext{\skt{कृच्छाम्चैषो}}{\lemma{\skt{कृच्छाम्}}\Afootnote{K; \skt{कृच्छ्रश्}\ A}}}{\lemma{\skt{चैषो}}\Afootnote{K; \skt{चायं}\ A}}\skt{ }\edtext{\skt{वाल\-वृ\-द्धे\-षे\-व\-साध्यः}}{\lemma{\skt{वालवृद्धेष्}}\Afootnote{K; \skt{बालवृद्धे}\ A}}\skt{ । }\pend
\pstart\versepar \skt{}\edtext{\edtext{\skt{तैस्तै\-र्भा\-वै\-शो\-च\-तो\-ल्पा\-स\-नस्य}}{\lemma{\skt{तैर्}}\Afootnote{K; \skt{तै}\ A}}}{\lemma{\skt{भावै\-शो\-च\-तो\-ल्पा\-स\-नस्य}}\Afootnote{K; \skt{र्भावैः}\ A}}\skt{ }\edtext{\skt{शुष्कं}}{\Afootnote{K; \emph{om.}\ A}}\skt{ }\edtext{\skt{शुक्तं}}{\Afootnote{K; \emph{om.}\ A}}\skt{ }\edtext{\skt{तेजे}}{\Afootnote{K; \emph{om.}\ A}}\skt{ }\edtext{\skt{ऊष्मो}}{\Afootnote{K; \emph{om.}\ A}}\skt{ }\edtext{\skt{गृहीत{\symfont ¦}}}{\Afootnote{K; \skt{पक्ति\-मा\-वि\-ध्य(\-आ\-वि\-श्य)\-जन्तोः ॥}\ A}}\skt{ ।}\edtext{}{\lemma{\skt{एवसाध्यः । तैस्}}\Afootnote{K; \skt{ष्वसाध्यः । तैस्}\ A}}\pend

\pstart\versepar \ledsidenote{\passref{6.40.14}}\skt{कोष्ठं गत्वा }\edtext{\edtext{\skt{विप्रमुष्ण}}{\lemma{\skt{विप्रम्}}\Afootnote{K; \skt{क्षोभ}\ A}}}{\lemma{\skt{उष्ण}}\Afootnote{K; \skt{यत्य}\ A}}\skt{ }\edtext{\edtext{\edtext{\skt{प्रणद्यस्तास्तं}}{\lemma{\skt{प्रणद्यस्}}\Afootnote{K; \skt{स्य रक्तं}\ A}}}{\lemma{\skt{तास्}}\Afootnote{K; \skt{त}\ A}}}{\lemma{\skt{तं}}\Afootnote{K; \skt{च्चा}\ A}}\skt{ }\edtext{\skt{रक्तात्काकणन्ति}}{\Afootnote{K; \skt{धस्तात्काकणन्ती}\ A}}\skt{ }\edtext{\skt{प्रकाश{\symfont ¦}}}{\Afootnote{K; \skt{प्रकाशम् ।}\ A}}\skt{ । }\pend
\pstart\versepar \skt{}\edtext{\skt{मयुरीशं}}{\Afootnote{K; \skt{वर्चोमिश्रं}\ A}}\skt{ }\edtext{\skt{निःपुरीगश्च}}{\Afootnote{K; \skt{निःपुरीषं}\ A}}\skt{ }\edtext{\skt{सग\-न्ध\-मा\-ग\-न्तश्च}}{\lemma{\skt{आगन्तश्च}}\Afootnote{K; \skt{निर्गन्धं वा}\ A}}\skt{ सार्यते }\edtext{\skt{योतिसारः}}{\Afootnote{K; \skt{तेन कृच्छ्रात् ॥}\ A}}\skt{ ।}\pend

\pstart\versepar \ledsidenote{\passref{6.40.15}}\skt{शोकोत्पन्नो }\edtext{\skt{दुश्चि\-कि\-त्स्यो\-ति\-सा\-मात्रं}}{\Afootnote{K; \skt{दुश्चि\-कि\-त्स्योऽ\-ति\-मात्रं}\ A}}\skt{ रोगो वैद्यैः कष्ट एष प्रदिष्टः । }\pend
\pstart\versepar \skt{}\edtext{\skt{आमा\-जी\-र्णो\-त्प्र\-द्रुता}}{\Afootnote{K; \skt{आमा\-जी\-र्णो\-प\-द्रुताः}\ A}}\skt{ }\edtext{\skt{क्षैभयन्तु}}{\Afootnote{K; \skt{क्षोभयन्तः}\ A}}\skt{ । कोष्ठं }\edtext{\edtext{\skt{जन्तार्पात्रमालाः}}{\lemma{\skt{जन्तार्}}\Afootnote{K; \emph{om.}\ A}}}{\lemma{\skt{पात्रमालाः}}\Afootnote{K; \emph{om.}\ A}}\skt{ }\edtext{\skt{सदोषाः}}{\Afootnote{K; \skt{दोषाः संप्रदुष्टाः सभक्तम् ॥}\ A}}\skt{ ।}\pend

\pstart \ledsidenote{\passref{6.40.16}}\edtext{\edtext{\skt{नाना\-व\-र्ण\-श्चा\-ति\-मात्रं}}{\lemma{\skt{नानावर्णश्}}\Afootnote{K; \skt{नानावर्णं}\ A}}\skt{ सु}}{\lemma{\skt{चातिमात्रं सु}}\Afootnote{K; \skt{नैकशः सारयन्ति कृच्छ्राज्जन्तोः}\ A}}\skt{ }\edtext{\edtext{\skt{षष्ठ\-मे\-त\-द्व\-दन्ति}}{\lemma{\skt{षष्ठम्}}\Afootnote{K; \skt{षष्ठ}\ A}}}{\lemma{\skt{एतद्वदन्ति}}\Afootnote{K; \skt{मेनं वदन्ति ।}\ A}}\skt{ ।}\pend

\pstart\versepar \ledsidenote{\passref{6.40.17}}\edtext{\edtext{\skt{मयि\-च्छ\-मे\-भि\-दो\-षै\-र्वा\-न्य\-स्त\-मय्नि}}{\lemma{\skt{मयि\-च्छ\-मे\-भि\-दो\-षैर्}}\Afootnote{K; \skt{संसृ\-ष्ट\-मे\-भि\-र्दोषै}\ A}}}{\lemma{\skt{वान्यस्तमय्नि}}\Afootnote{K; \skt{स्तु न्यस्तमप्स्व}\ A}}\skt{ }\edtext{\skt{ग ति}}{\Afootnote{K; \skt{वसीदति ।}\ A}}\skt{ । }\pend
\pstart\versepar \skt{पुरीषं }\edtext{\skt{भृगदुर्गन्धि}}{\Afootnote{K; \skt{भृशदुर्गन्धि}\ A}}\skt{ }\edtext{\edtext{\skt{विच्छि\-न्न\-श्चा\-म\-सं\-ज्ञि\-तम्}}{\lemma{\skt{विच्छिन्नश्}}\Afootnote{K; \skt{विच्छिन्नं}\ A}}\skt{ {\symfont ¦}}}{\lemma{\skt{चामसंज्ञितम् {\symfont ¦}}}\Afootnote{K; \skt{चामसंज्ञकम् ॥}\ A}}\skt{ ।}\pend

\pstart\versepar \ledsidenote{\passref{6.40.18}}\skt{एतान्ये व तु लिङ्गानि विपरीतानि }\edtext{\skt{यस्य वैः}}{\Afootnote{K; \skt{यस्य तु ।}\ A}}\skt{ । }\pend
\pstart\versepar \skt{}\edtext{\skt{लाघवश्च}}{\Afootnote{K; \skt{लाघवं च}\ A}}\skt{ मनुष्यस्य }\edtext{\edtext{\skt{तन्पक्क्वमिति}}{\lemma{\skt{तन्पक्क्वम्}}\Afootnote{K; \skt{तस्य प}\ A}}}{\lemma{\skt{इति}}\Afootnote{K; \skt{क्वं वि}\ A}}\skt{ }\edtext{\skt{निद्रिर्शते}}{\Afootnote{K; \skt{निर्दिशेत् ॥}\ A}}\skt{ ।}\pend

\pstart\versepar \ledsidenote{\passref{6.40.19}}\edtext{\skt{॥ मञ्ज\-य\-कृ\-द्वे\-श\-वा\-रा\-म्वु\-तैलं}}{\Afootnote{K; \skt{सर्पि\-र्भे\-दो\-वे\-स\-वा\-रा\-म्बु\-तैल}\ A}}\skt{ }\edtext{\skt{सर्पिः}}{\Afootnote{K; \skt{मज्जा}\ A}}\skt{ }\edtext{\skt{क्षौद्रं}}{\Afootnote{K; \skt{क्षीर}\ A}}\skt{ }\edtext{\skt{क्षीररूपं}}{\Afootnote{K; \skt{क्षौद्ररूपं}\ A}}\skt{ स्नवेत्{\symfont ¦} । }\pend
\pstart\versepar \skt{मञ्जि\-ष्ठा\-भ\-म्मा\-तु\-ल\-ङ्गो\-पम्वा विस्रं शीतं }\edtext{\skt{सर्व\-ग\-न्ध्य\-ञ्ज\-नाभ{\symfont ¦}}}{\Afootnote{K; \skt{प्रेत\-ग\-न्ध्य\-ञ्ज\-ना\-भम् ॥}\ A}}\skt{ ।}\edtext{}{\lemma{\skt{स्नवेत्{\symfont ¦} । मञ्जि\-ष्ठा\-भ\-म्मा\-तु\-ल\-ङ्गो\-पम्वा}}\Afootnote{K; \skt{स्रवेद्यत् । मञ्जिष्ठाभं मस्तुलुङ्गोपमं वा}\ A}}\pend

\pstart\versepar \ledsidenote{\passref{6.40.20}}\skt{राजीमद्वा }\edtext{\skt{चन्द्रिकैः}}{\Afootnote{K; \skt{चन्द्रकैः}\ A}}\skt{ सन्ततम्वा }\edtext{\edtext{\edtext{\skt{पूयद्रकश्चाश्च}}{\lemma{\skt{पूयद्रकश्}}\Afootnote{K; \skt{पूयप्रख्यं कर्दमाभं तथोष्ण्}\ A}}}{\lemma{\skt{चाश्}}\Afootnote{K; \skt{अम्}\ A}}}{\lemma{\skt{च}}\Afootnote{K; \skt{।}\ A}}\skt{ }\pend
\pstart\versepar \skt{}\edtext{\edtext{\skt{हन्या\-दे\-त\-द्वि\-प\-री\-तश्च}}{\lemma{\skt{हन्या\-दे\-त\-द्वि\-प\-री\-तश्}}\Afootnote{K; \skt{हन्या\-दे\-त\-द्य\-त्प्र\-तीपं}\ A}}}{\lemma{\skt{च}}\Afootnote{K; \skt{भ}\ A}}\skt{ }\edtext{\skt{यः}}{\Afootnote{K; \skt{वेच्च}\ A}}\skt{ क्षीणं }\edtext{\skt{हन्युर्वहव}}{\lemma{\skt{वहव}}\Afootnote{K; \skt{चो}\ A}}\skt{ }\edtext{\skt{पसर्गाः}}{\Afootnote{K; \skt{पसर्गाः प्रभूताः ॥}\ A}}\skt{ ।}\pend

\pstart\versepar \ledsidenote{\passref{6.40.21}}\skt{असंवृतगुदं क्षीणं दुरा\-ध्मा\-त\-मु\-प\-द्रुतं । }\pend
\pstart\versepar \skt{गुदे पक्वे गतो\-ष्मा\-ण\-म\-ती\-सा\-र\-किणं त्यजेत् ।}\pend

\pstart\versepar \ledsidenote{\passref{6.40.24}}\skt{आमपक्वक्रमं हित्वा }\edtext{\skt{नातीसारे}}{\Afootnote{K; \skt{नातिसारे}\ A}}\skt{ क्रिया }\edtext{\skt{मतः}}{\Afootnote{K; \skt{यतः ।}\ A}}\skt{ । }\pend
\pstart\versepar \skt{अतः }\edtext{\skt{सर्वातिसारेषु}}{\Afootnote{K; \skt{सर्वेऽतिसारास्तु}\ A}}\skt{ }\edtext{\skt{ज्ञेयं}}{\Afootnote{K; \skt{ज्ञेयाः}\ A}}\skt{ पक्वामलक्षणैः ।}\pend

\pstart\versepar \ledsidenote{\passref{6.40.25}}\edtext{\skt{हितं}}{\Afootnote{K; \skt{तत्रादौ}\ A}}\skt{ }\edtext{\skt{लङ्घानमेवादौ}}{\Afootnote{K; \skt{लङ्घनं कार्य}\ A}}\skt{ }\edtext{\skt{र्सरूपेषु}}{\Afootnote{K; \skt{मतिसारेषु}\ A}}\skt{ }\edtext{\skt{देहिना{\symfont ¦}}}{\Afootnote{K; \skt{देहिनाम् ।}\ A}}\skt{ । }\pend
\pstart\versepar \skt{}\edtext{\skt{तत्रा\-स\-च्छ\-र्द\-न\-द्कार्यं}}{\Afootnote{K; \skt{ततः}\ A}}\skt{ }\edtext{\skt{पिप्प\-ली\-ल\-व\-णाम्च}}{\lemma{\skt{पिप्प\-ली\-ल\-व\-णाम्}}\Afootnote{K; \skt{पा}\ A}}\skt{ }\edtext{\skt{ना}}{\Afootnote{K; \skt{नसंयुक्तो यवाग्वादिक्रमो हितः ॥}\ A}}\skt{ ।}\pend

\pstart \ledsidenote{\passref{6.40.27ab}}\edtext{\edtext{\skt{कार्यश्चानशनं}}{\lemma{\skt{कार्यश्}}\Afootnote{K; \skt{कार्यं}\ A}}}{\lemma{\skt{चानशनं}}\Afootnote{K; \skt{च वमन}\ A}}\skt{ स्यान्ते }\edtext{\skt{द्रवं}}{\Afootnote{K; \skt{प्रद्रवं}\ A}}\skt{ }\edtext{\skt{लघु च}}{\Afootnote{K; \skt{लघु}\ A}}\skt{ }\edtext{\skt{भोजन{\symfont ¦}}}{\Afootnote{K; \skt{भोजनम् ।}\ A}}\skt{ ।}\pend

\pstart\versepar \ledsidenote{\passref{6.40.29}}\edtext{\edtext{\skt{दोष\-स्त\-म्भ\-न\-मग्रे}}{\lemma{\skt{दोषस्तम्भनम्}}\Afootnote{K; \skt{आमातीसारिणां कार्यं नादौ संग्रहणं नृणाम् । तेषां दो}\ A}}\skt{ तु न}}{\lemma{\skt{अग्रे तु न}}\Afootnote{K; \emph{om.}\ A}}\skt{ }\edtext{\edtext{\skt{कर्त\-व्य\-म्वि\-जा\-नता}}{\lemma{\skt{कर्तव्यम्}}\Afootnote{K; \emph{om.}\ A}}}{\lemma{\skt{विजानता}}\Afootnote{K; \skt{षा}\ A}}\skt{ । }\pend
\pstart\versepar \skt{}\edtext{\skt{तदौ}}{\Afootnote{K; \skt{वि}\ A}}\skt{ }\edtext{\skt{वध्यमाणास्तु}}{\Afootnote{K; \skt{बद्धाः प्राग्जनयन्त्य्}\ A}}\skt{ }\edtext{\skt{वहुं}}{\Afootnote{K; \skt{आमया}\ A}}\skt{ }\edtext{\edtext{\skt{कुर्व\-त्यु\-प\-द्र\-वान्}}{\lemma{\skt{कुर्वत्य्}}\Afootnote{K; \skt{नि}\ A}}}{\lemma{\skt{उपद्रवान्}}\Afootnote{K; \skt{मान् ॥}\ A}}\pend

\pstart\versepar \ledsidenote{\passref{6.40.30.add}}\skt{नामे सङ्गुरुणं दद्यादतिसार केषञ्चनः । }\pend
\pstart\versepar \skt{अकाल स हीतो हि विकारां कुरुते वनं ।}\pend

\pstart\versepar \ledsidenote{\passref{6.40.30}}\edtext{\skt{दण्ड\-का\-ल\-स\-का\-त्मा\-न\-शो\-फो\-द\-र\-भ\-ग\-न्दराः}}{\Afootnote{K; \emph{om.}\ A}}\skt{ । }\pend
\pstart\versepar \skt{}\edtext{\skt{कुष्ठयां}}{\Afootnote{K; \emph{om.}\ A}}\skt{ }\edtext{\skt{मय प्लीह\-गु\-ल्मा\-या\-क\-रा\-दि\-कान्}}{\Afootnote{K; \skt{प्लीह\-पा\-ण्ड्वा\-म\-या\-ना\-ह\-मे\-ह\-कु\-ष्ठो\-द\-र\-ज्व\-रान् । शोफ\-गु\-ल्म\-ग्र\-ह\-ण्य\-र्शः\-शू\-ला\-ल\-स\-क\-हृ\-द्ग्र\-हान् ॥}\ A}}\skt{ ।}\pend

\pstart\versepar \ledsidenote{\passref{6.40.33}}\edtext{\skt{स्तोकस्तोकं}}{\Afootnote{K; \skt{स्तोकं स्तोकं}\ A}}\skt{ }\edtext{\skt{विवद्धम्वा}}{\Afootnote{K; \skt{विबद्धं वा}\ A}}\skt{ सशूलं }\edtext{\skt{योतिसार्यते}}{\Afootnote{K; \skt{योऽतिसार्यते ।}\ A}}\skt{ । }\pend
\pstart\versepar \skt{अभ\-या\-पि\-प्प\-ली\-क\-ल्कै\-स्सु\-खो\-ष्णै\-स्त\-म्वि\-रे\-च\-येत् ।}\pend

\pstart \ledsidenote{\passref{6.40.28cd}}\skt{हरि\-द्रा\-दी\-म्व\-चा\-दि\-म्वा\-त\-स्मा\-दादौ पिवेन्नरः ।}\pend

\pstart \ledsidenote{\passref{6.40.27cd}}\edtext{\skt{गल\-यु\-ष\-य\-य\-वा\-गूभिः}}{\Afootnote{K; \skt{खड\-यू\-ष\-य\-वा\-गूषु}\ A}}\skt{ }\edtext{\skt{पिप्पल्यादिश्च}}{\Afootnote{K; \skt{पिप्पल्याद्यं च}\ A}}\skt{ }\edtext{\skt{यो}}{\Afootnote{K; \skt{यो}\ A}}\skt{ }\edtext{\skt{जयेत}}{\Afootnote{K; \skt{जयेत् ॥}\ A}}\skt{ ।}\pend

\pstart \ledsidenote{\passref{6.40.35cd}}\edtext{\skt{देव\-दा\-रु\-व\-चा\-मु\-स्ता\-ना\-ग\-रा\-ति\-वि\-षा\-भया}}{\Afootnote{K; \skt{देव\-दा\-रु\-व\-चा\-मु\-स्ता\-ना\-ग\-रा\-ति\-वि\-षा\-भयाः ॥}\ A}}\skt{ ॥}\pend

\pstart\versepar \ledsidenote{\passref{6.40.36}}\skt{अभया }\edtext{\skt{धान्य\-क\-म्मु\-स्ता\-म्पि\-प्पली}}{\lemma{\skt{मुस्ताम्}}\Afootnote{K; \skt{मुस्तं}\ A}}\skt{ }\edtext{\skt{नागराम्वचा}}{\lemma{\skt{नागराम्}}\Afootnote{K; \skt{नागरं}\ A}}\skt{ । }\pend
\pstart\versepar \skt{}\edtext{\skt{पिपली}}{\Afootnote{K; \skt{नागरं}\ A}}\skt{ धान्यकं मुस्तं }\edtext{\skt{वाल\-क\-वि\-ल्व\-मेव}}{\lemma{\skt{वालकविल्वम्}}\Afootnote{K; \skt{वालकं बिल्वम्}\ A}}\skt{ च ।}\pend

\pstart\versepar \ledsidenote{\passref{6.40.37}}\skt{मुस्तं }\edtext{\skt{पर्पकं}}{\Afootnote{K; \skt{पर्पटकं}\ A}}\skt{ शुण्ठी वचा }\edtext{\skt{सातिविषाभया}}{\Afootnote{K; \skt{प्रतिविषाभया ।}\ A}}\skt{ । }\pend
\pstart\versepar \skt{}\edtext{\skt{वति}}{\Afootnote{K; \skt{अभयाति}\ A}}\skt{ विषा हिङ्गु वचा }\edtext{\skt{सौवर्चला}}{\Afootnote{K; \skt{सौवर्चलं}\ A}}\skt{ }\edtext{\skt{भया}}{\Afootnote{K; \skt{तथा ॥}\ A}}\skt{ ।}\pend

\pstart\versepar \ledsidenote{\passref{6.40.45}}\edtext{\skt{महौ\-ष\-ध\-म\-ति\-विषा}}{\lemma{\skt{अतिविषा}}\Afootnote{K; \skt{प्रतिविषा}\ A}}\skt{ }\edtext{\skt{मुस्ता}}{\Afootnote{K; \skt{मुस्तं}\ A}}\skt{ चेत्याम }\edtext{\skt{पाचनाः}}{\Afootnote{K; \skt{पाचनाः ।}\ A}}\skt{ । }\pend
\pstart\versepar \skt{}\edtext{\skt{एते}}{\Afootnote{K; \skt{प्रयो}\ A}}\skt{ }\edtext{\skt{योगास्मसारताः}}{\Afootnote{K; \skt{ज्या विंशतिर्योगाः}\ A}}\skt{ }\edtext{\edtext{\skt{श्लोका\-र्ध\-वि\-हि\-तास्तु}}{\lemma{\skt{श्लोका\-र्ध\-वि\-हि\-तास्}}\Afootnote{K; \skt{श्लोकार्धविहिता}\ A}}}{\lemma{\skt{तु}}\Afootnote{K; \skt{स्त्वि}\ A}}\skt{ मे ।}\pend

\pstart\versepar \ledsidenote{\passref{6.40.46}}\edtext{\skt{धान्या\-म्लो\-ष्णा\-म्बु\-म\-ध्यानां}}{\Afootnote{\skt{धान्या\-म्लो\-ष्णा\-म्बु\-म\-द्यानां}\ A; \skt{धान्या\-म्लो\-ष्णा\-म्वु\-म\-ध्यानां}\ K}}\skt{ }\edtext{\skt{पिबेद}}{\Afootnote{A; \skt{पिवेद}\ K}}\skt{ }\edtext{\skt{अन्यतमेन वा}}{\Afootnote{K; \skt{न्यतमेन वा ।}\ A}}\skt{ । }\pend
\pstart\versepar \skt{}\edtext{\skt{निष्क्वाथां वा}}{\Afootnote{K; \skt{निष्क्वाथान्वा}\ A}}\skt{ }\edtext{\skt{पिबेदेषां}}{\Afootnote{\skt{पिबेदेषां}\ A; \skt{पिवेदेषां}\ K}}\skt{ }\edtext{\skt{सुखोष्णां}}{\Afootnote{K; \skt{सुखोष्णा}\ A}}\skt{ }\edtext{\skt{साधु}}{\Afootnote{K; \skt{न्साधु}\ A}}\skt{ }\edtext{\skt{योजिता{\symfont ¦}}}{\Afootnote{K; \skt{साधितान् ॥}\ A}}\skt{ ।}\pend

\pstart \ledsidenote{\passref{6.40.63}}\skt{मुस्तं }\edtext{\skt{कुटजवीजानि}}{\Afootnote{K; \skt{कुटजबीजानि}\ A}}\skt{ }\edtext{\skt{भूनिम्च}}{\Afootnote{K; \skt{भूनिम्बं}\ A}}\skt{ }\edtext{\skt{सदुरालभ{\symfont ¦}}}{\Afootnote{K; \skt{सरसाञ्जनम् । दार्वी दुरालभा बिल्वं वालकं रक्तचन्दनम् ॥}\ A}}\skt{ ।}\pend

\pstart \ledsidenote{\passref{6.40.64}}\edtext{\skt{उशीरं}}{\Afootnote{K; \skt{चन्दनं वालकं मुस्तं भूनिम्बं सदुरालभम् । मृणालं}\ A}}\skt{ चन्दनं }\edtext{\skt{लोध्रां}}{\Afootnote{K; \skt{रोध्रं}\ A}}\skt{ नागरं }\edtext{\edtext{\skt{नीलमुत्पल{\symfont ¦}}}{\lemma{\skt{नीलम्}}\Afootnote{K; \skt{नील}\ A}}}{\lemma{\skt{उत्पल{\symfont ¦}}}\Afootnote{K; \skt{मुत्पलम् ॥}\ A}}\skt{ ।}\pend

\pstart\versepar \ledsidenote{\passref{6.40.65}}\skt{पाठा मुस्तं हरिद्रे द्वे पिप्पली कौटजं }\edtext{\skt{फल{\symfont ¦}}}{\Afootnote{K; \skt{फलम् ।}\ A}}\skt{ । }\pend
\pstart\versepar \skt{}\edtext{\skt{फलं}}{\Afootnote{K; \skt{फल}\ A}}\skt{ त्वचम्वत्सकस्य }\edtext{\skt{शृङ्वरे}}{\Afootnote{K; \skt{शृङ्गवेरं}\ A}}\skt{ }\edtext{\skt{घो}}{\Afootnote{K; \skt{घनं}\ A}}\skt{ वचा ।}\pend

\pstart\versepar \ledsidenote{\passref{6.40.71.add}}\skt{चत्वार एते योगास्स्युः }\edtext{\skt{पक्वा\-ती\-सा\-र\-ना\-श\-नामे}}{\Afootnote{K\textsuperscript{pc}; \skt{पक्वा\-ती\-सा\-र\-ना\-शना ॥}\ K\textsuperscript{ac}}}\skt{ ॥ }\pend
\pstart\versepar \skt{}\edtext{\skt{उक्ता\-म\-यो\-प\-यो\-ज्यास्ते}}{\Afootnote{\skt{उक्ता\-म\-यो\-प\-यो\-ज्यास्ते}\ K}}\skt{ }\edtext{\skt{सक्षौ\-द्रा\-स्त\-ण्डु\-लां\-बुना}}{\Afootnote{\skt{सक्षौ\-द्रा\-स्त\-ण्डु\-लां\-वुना ।}\ K}}\skt{ ।}\pend

\pstart\versepar \ledsidenote{\passref{6.40.39}}\skt{मूर्वा निर्दहनी पाठा }\edtext{\skt{त्रिकटु}}{\Afootnote{K; \skt{त्र्यूषणं}\ A}}\skt{ गजपिप्पली । }\pend
\pstart\versepar \skt{}\edtext{\skt{कटुराहिणो}}{\Afootnote{K; \skt{सिद्धार्थका भ}\ A}}\skt{ }\edtext{\skt{रुद्रदारु}}{\Afootnote{K; \skt{द्रदारु}\ A}}\skt{ शताह्वा }\edtext{\skt{गौरसर्षपा}}{\Afootnote{K; \skt{कटुरोहिणी ॥}\ A}}\skt{ ।}\pend

\pstart\versepar \ledsidenote{\passref{6.40.40}}\skt{एला }\edtext{\skt{शारिवकं}}{\Afootnote{K; \skt{सावरकं}\ A}}\skt{ कुष्ठं हरिद्रे }\edtext{\skt{शुक्रज}}{\Afootnote{K; \skt{कौटजा}\ A}}\skt{ यवा । }\pend
\pstart\versepar \skt{}\edtext{\skt{मेषशृण्ग}}{\Afootnote{K; \skt{मेषशृङ्गी}\ A}}\skt{ त्व }\edtext{\skt{कृमिघ्नं}}{\Afootnote{K; \skt{गेले च कृमिघ्नं}\ A}}\skt{ वृक्षकाणि च ।}\pend

\pstart\versepar \ledsidenote{\passref{6.40.41}}\skt{वृक्षादनी }\edtext{\skt{वीरतरां}}{\Afootnote{K; \skt{वीरतरु}\ A}}\skt{ }\edtext{\skt{वृहत्यौ द्वे}}{\Afootnote{K; \skt{र्बृहत्यौ द्वे}\ A}}\skt{ सहे तथा ॥ }\pend
\pstart\versepar \skt{}\edtext{\skt{अर\-लु\-त्व\-तै\-न्दुकी च}}{\Afootnote{K; \skt{अर\-लु\-त्व\-क्तै\-न्दुकी च}\ A}}\skt{ }\edtext{\edtext{\skt{दाडिमन्तर्व}}{\lemma{\skt{दाडिमन्तर्}}\Afootnote{K; \skt{दाडिमी कौटज्}\ A}}}{\lemma{\skt{व}}\Afootnote{K; \skt{ई शमी ॥}\ A}}\skt{ ।}\pend

\pstart\versepar \ledsidenote{\passref{6.40.41.add}}\skt{साल\-सा\-रा\-दि\-निर्हो लवणानि तथैव च । }\pend
\pstart\versepar \skt{यो गातयो स्यरा\-म\-ण\-र्श\-कराः स्मृताः ।}\pend

\pstart \ledsidenote{\passref{6.40.42ab}}\skt{पाठा }\edtext{\skt{तेजवती}}{\Afootnote{K; \skt{तेजोवती}\ A}}\skt{ }\edtext{\skt{मुस्तो}}{\Afootnote{K; \skt{मुस्तं}\ A}}\skt{ पिप्पलीकौटजं फलं ॥}\pend

\pstart\versepar \ledsidenote{\passref{6.40.46.add}}\skt{सुरापान्यम्ल सीनां }\edtext{\skt{पिबे\-द\-न्य\-त\-रेण तु}}{\Afootnote{\skt{पिवे\-द\-न्य\-त\-रेण तु ।}\ K}}\skt{ । }\pend
\pstart\versepar \skt{निस्क्वाथं वापि }\edtext{\skt{पिबेदेषां}}{\Afootnote{\skt{पिवेदेषां}\ K}}\skt{ सुखोष्णां साधु जोजितां ।}\pend

\pstart \ledsidenote{\passref{6.40.48cd}}\edtext{\edtext{\skt{हरी\-त\-की\-म\-ति\-विषां}}{\lemma{\skt{हरीतकीम्}}\Afootnote{K; \skt{हरीतकी}\ A}}}{\lemma{\skt{अतिविषां}}\Afootnote{K; \skt{मतिविषां}\ A}}\skt{ हिङ्गु }\edtext{\skt{सौवर्चालं}}{\Afootnote{K; \skt{सौवर्चलं}\ A}}\skt{ }\edtext{\skt{वचा{\symfont ¦}}}{\Afootnote{K; \skt{वचाम् ॥}\ A}}\skt{ ।}\pend

\pstart\versepar \ledsidenote{\passref{6.40.49}}\edtext{\skt{पिबेत्सुखांवुना}}{\Afootnote{\skt{पिबेत्सुखाम्बुना}\ A; \skt{पिवेत्सुखांवुना}\ K}}\skt{ }\edtext{\skt{जन्तु\-रा\-मा\-ती\-सा\-र\-पी\-डितः}}{\Afootnote{K; \skt{जन्तु\-रा\-मा\-ति\-सा\-र\-पी\-डितः ।}\ A}}\skt{ । }\pend
\pstart\versepar \skt{पटोलं दीप्यकं }\edtext{\skt{निम्व}}{\Afootnote{K; \skt{बिल्वं}\ A}}\skt{ }\edtext{\skt{वचा\-पि\-प्प\-लि\-ना\-गरां}}{\Afootnote{K; \skt{वचा\-पि\-प्प\-लि\-ना\-ग\-रम् ॥}\ A}}\skt{ ।}\pend

\pstart\versepar \ledsidenote{\passref{6.40.50}}\edtext{\skt{विडङ्गनाथ}}{\Afootnote{K; \skt{मुस्तं कुष्ठं विडङ्गं}\ A}}\skt{ }\edtext{\skt{कुष्ठश्च}}{\Afootnote{K; \skt{च}\ A}}\skt{ }\edtext{\skt{पिबेद्वापि}}{\Afootnote{\skt{पिबेद्वापि}\ A; \skt{पिवेद्वापि}\ K}}\skt{ }\edtext{\skt{सुखाम्वुना}}{\Afootnote{K; \skt{सुखाम्बुना ।}\ A}}\skt{ ॥ }\pend
\pstart\versepar \skt{शृङ्गवेरं }\edtext{\skt{गुडुचीश्च}}{\Afootnote{K; \skt{गुडूचीं च}\ A}}\skt{ }\edtext{\skt{पिवे}}{\Afootnote{K; \skt{पिबे}\ A}}\skt{ }\edtext{\skt{उष्णेन}}{\Afootnote{K; \skt{दुष्णेन}\ A}}\skt{ }\edtext{\skt{चांभसा}}{\Afootnote{K; \skt{वारिणा ॥}\ A}}\skt{ ।}\pend

\pstart \ledsidenote{\passref{6.40.51.add}}\skt{योगादारौ तेयो कराः स्मृताः ।}\pend

\pstart\versepar \ledsidenote{\passref{6.40.51}}\skt{लवणान्यथ पिप्पल्यो विडङ्गानि हरीतकी ॥ }\pend
\pstart\versepar \skt{चित्रकं शिंशपा पाठा }\edtext{\skt{शार्ंगेष्ठा}}{\Afootnote{K; \skt{शार्ङ्गेष्टा}\ A}}\skt{ }\edtext{\skt{लवनानि}}{\Afootnote{K; \skt{लवणानि}\ A}}\skt{ च ।}\pend

\pstart\versepar \ledsidenote{\passref{6.40.52}}\edtext{\skt{हिङ्गु\-वृ\-क्ष\-क\-वी\-जानि}}{\Afootnote{K; \skt{हिङ्गु\-र्वृ\-क्ष\-क\-बी\-जानि}\ A}}\skt{ लवणानि च भागशः । }\pend
\pstart\versepar \skt{}\edtext{\skt{पटोलं}}{\Afootnote{K; \skt{ह}\ A}}\skt{ ॥ }\edtext{\skt{हन्ति}}{\Afootnote{K; \skt{स्ति}\ A}}\skt{ दंत्यथ पिप्पल्यः कल्कावक्ष समौ स्मृतौ ॥}\pend

\pstart\versepar \ledsidenote{\passref{6.40.53}}\skt{वचा }\edtext{\skt{कुटकवीजानि}}{\Afootnote{K; \skt{गुडूचीकाण्डानि योग्}\ A}}\skt{ }\edtext{\edtext{\skt{योरोथं}}{\lemma{\skt{योर्}}\Afootnote{K; \skt{ओ}\ A}}}{\lemma{\skt{ओथं}}\Afootnote{K; \skt{ऽयं}\ A}}\skt{ }\edtext{\skt{पञ्चमो}}{\Afootnote{K; \skt{परमो}\ A}}\skt{ मतः । }\pend
\pstart\versepar \skt{एते }\edtext{\skt{सुखाम्वु}}{\Afootnote{K; \skt{सुखाम्बु}\ A}}\skt{ }\edtext{\skt{पेया}}{\Afootnote{K; \skt{ना यो}\ A}}\skt{ }\edtext{\skt{योगाः}}{\Afootnote{K; \skt{गा देयाः}\ A}}\skt{ }\edtext{\skt{पञ्चामना}}{\Afootnote{K; \skt{पश्च सतां}\ A}}\skt{ }\edtext{\skt{मनाः}}{\Afootnote{K; \skt{मताः ॥}\ A}}\skt{ ।}\pend

\pstart\versepar \ledsidenote{\passref{6.40.54.add}}\skt{पय\-स्यां\-न्का\-ष्म\-मु\-स्ता\-म्चाविं मति\-भ\-द्र\-सा\-ह्वया{\symfont ¦} । }\pend
\pstart\versepar \skt{क्षीरावगिष्टं पीत्वा हन्यामं शूलमेव च ।}\pend

\pstart\versepar \ledsidenote{\passref{6.40.54}}\edtext{\skt{निवृ\-त्ते\-ष्वा\-स\-शू\-लेषु}}{\Afootnote{K; \skt{निवृ\-त्ते\-ष्वा\-म\-शू\-लेषु}\ A}}\skt{ }\edtext{\skt{यस्या न}}{\Afootnote{K; \skt{यस्य न}\ A}}\skt{ }\edtext{\skt{प्रगुणोनिलः}}{\Afootnote{K; \skt{प्रगुणोऽनिलः ।}\ A}}\skt{ । }\pend
\pstart\versepar \skt{}\edtext{\skt{स्तोकस्तोकं}}{\Afootnote{K; \skt{स्तोकं स्तोकं}\ A}}\skt{ }\edtext{\skt{सशूलम्वा}}{\Afootnote{K; \skt{रुजामच्च सशूलं}\ A}}\skt{ }\edtext{\skt{रुजा\-वा\-न्या\-ति\-सा\-र्यते}}{\Afootnote{K; \skt{योऽतिसार्यते ॥}\ A}}\skt{ ।}\pend

\pstart\versepar \ledsidenote{\passref{6.40.55}}\edtext{\edtext{\skt{सक्षा\-र\-ल\-व\-णै\-यु\-क्त्म्म\-न्दाग्नि}}{\lemma{\skt{सक्षा\-र\-ल\-व\-णै\-युक्त्म्}}\Afootnote{K; \skt{सक्षा\-र\-ल\-व\-णै\-र्युक्तं}\ A}}}{\lemma{\skt{मन्दाग्नि}}\Afootnote{K; \skt{मन्दाग्निः}\ A}}\skt{ }\edtext{\skt{प्रपिबेद्घृत{\symfont ¦}}}{\Afootnote{\skt{स पिबेद्घृतम् ।}\ A; \skt{प्रपिवेद्घृत{\symfont ¦} ।}\ K}}\skt{ । }\pend
\pstart\versepar \skt{}\edtext{\skt{क्षार\-ना\-ग\-र\-चा\-ङ्गे\-री\-को\-ल\-द\-ध्य\-म्ल\-सा\-धित{\symfont ¦}}}{\Afootnote{K; \skt{क्षीर\-ना\-ग\-र\-चा\-ङ्ग\-री\-को\-ल\-द\-ध्य\-म्ल\-सा\-धि\-तम् ॥}\ A}}\skt{ ।}\pend

\pstart\versepar \ledsidenote{\passref{6.40.56}}\skt{सर्पिरच्छं }\edtext{\skt{पिबेद्वापि}}{\Afootnote{\skt{पिबेद्वापि}\ A; \skt{पिवेद्वापि}\ K}}\skt{ }\edtext{\skt{शूला\-तु\-सा\-र\-शा\-न्तये}}{\Afootnote{K; \skt{शूला\-ती\-सा\-र\-शा\-न्तये ।}\ A}}\skt{ । }\pend
\pstart\versepar \skt{}\edtext{\skt{दध्नाम्लघृतं}}{\Afootnote{K; \skt{दध्ना तैलघृतं}\ A}}\skt{ पक्वं }\edtext{\skt{सव्यो\-षा\-जा\-ति\-चि\-त्रक{\symfont ¦}}}{\Afootnote{K; \skt{सव्यो\-षा\-जा\-ति\-चि\-त्रकैः ॥}\ A}}\skt{ ।}\pend

\pstart\versepar \ledsidenote{\passref{6.40.57}}\edtext{\skt{सेचव्य}}{\Afootnote{K; \skt{सबिल्व}\ A}}\skt{ }\edtext{\skt{पिप्पलीमूलं}}{\Afootnote{K; \skt{पिप्पलीमूल}\ A}}\skt{ दाडिमैर्वा }\edtext{\skt{जगर्दितः}}{\Afootnote{K; \skt{रुगन्वितैः ।}\ A}}\skt{ । }\pend
\pstart\versepar \skt{}\edtext{\skt{निखिले}}{\Afootnote{K; \skt{निखि}\ A}}\skt{ }\edtext{\skt{नोपदिष्टायं}}{\Afootnote{K; \skt{लो}\ A}}\skt{ }\edtext{\skt{विधि\-रा\-मा\-य\-शा\-न्तये}}{\Afootnote{K; \skt{विधिरुक्तोऽयं वात\-श्ले\-ष्मो\-प\-शा\-न्तये ॥}\ A}}\skt{ ।}\pend

\pstart\versepar \ledsidenote{\passref{6.40.58}}\edtext{\skt{तूक्ष्णोष्ण}}{\Afootnote{K; \skt{तीक्ष्णोष्ण}\ A}}\skt{ }\edtext{\skt{वर्ज्यमेतन्तु}}{\Afootnote{K; \skt{वर्ज्यमेनं तु}\ A}}\skt{ विदध्यात्पित्तजे }\edtext{\skt{विधि{\symfont ¦}}}{\Afootnote{K; \skt{भिष}\ A}}\skt{ । }\pend
\pstart\versepar \skt{}\edtext{\skt{यथो\-क्त\-मु\-प\-वा\-सान्ते}}{\lemma{\skt{यथोक्तम्}}\Afootnote{K; \skt{क् । यथोक्तम्}\ A}}\skt{ }\edtext{\skt{यवा\-गू\-श्चा\-स्य\-प्र\-द्दते}}{\lemma{\skt{चास्यप्रद्दते}}\Afootnote{K; \skt{च प्रशस्यते ॥}\ A}}\skt{ ।}\pend

\pstart\versepar \ledsidenote{\passref{6.40.59}}\edtext{\skt{वलयोरं}}{\lemma{\skt{वलयोर्}}\Afootnote{K; \skt{बलयोर्}\ A}}\skt{ शुमत्याम्च }\edtext{\skt{श्वद्रं\-ष्टा\-वृ\-ह\-तीषु च}}{\Afootnote{K; \skt{श्वदं\-ष्ट्र\-बृ\-ह\-तीषु च ।}\ A}}\skt{ । }\pend
\pstart\versepar \skt{}\edtext{\skt{शतावर्याश्च}}{\Afootnote{K; \skt{शतावर्यां च}\ A}}\skt{ }\edtext{\skt{संसिद्धा}}{\Afootnote{K; \skt{संसिद्धाः}\ A}}\skt{ सुशीता मधुसंयुताः ।}\pend

\pstart\versepar \ledsidenote{\passref{6.40.60}}\skt{मुद्गादिषु च }\edtext{\skt{युषां}}{\Afootnote{K; \skt{यूषाः}\ A}}\skt{ }\edtext{\skt{स्युर्द्रव्यैरेतिः}}{\Afootnote{K; \skt{स्युर्द्रव्यैरेतैः}\ A}}\skt{ सुसंस्कृताः । }\pend
\pstart\versepar \skt{}\edtext{\skt{मृदु\-भि\-र्दी\-प\-नै\-स्ति\-क्तै\-र्द्र\-व्यै\-स्या\-ध्या\-म\-पा\-चन{\symfont ¦}}}{\lemma{\skt{दीप\-नै\-स्ति\-क्तै\-र्द्र\-व्यै\-स्या\-ध्या\-म\-पा\-चन{\symfont ¦}}}\Afootnote{K; \skt{दीप\-नै\-स्ति\-क्तै\-र्द्रव्यैः स्यादा\-म\-पा\-च\-नम् ॥}\ A}}\skt{ ।}\pend

\pstart\versepar \ledsidenote{\passref{6.40.68}}\skt{यदा }\edtext{\skt{पक्वोप्यतीसार}}{\Afootnote{K; \skt{पक्वोऽप्यतीसारः}\ A}}\skt{ }\edtext{\edtext{\skt{सरत्येवं}}{\lemma{\skt{सरत्य्}}\Afootnote{K; \skt{सर}\ A}}}{\lemma{\skt{एवं}}\Afootnote{K; \skt{स्येव}\ A}}\skt{ }\edtext{\skt{मुर्हुमुर्हुः}}{\Afootnote{K; \skt{मुहुर्मुहुः ।}\ A}}\skt{ । }\pend
\pstart\versepar \skt{ग्रहण्या }\edtext{\skt{मार्द\-वा\-त\-ज\-न्तो\-स्तत्र}}{\lemma{\skt{मार्द\-वा\-त\-ज\-न्तोस्}}\Afootnote{K; \skt{मार्दवाज्जन्तोस्}\ A}}\skt{ }\edtext{\skt{सस्तम्भनो}}{\Afootnote{K; \skt{संस्तम्भनं}\ A}}\skt{ }\edtext{\skt{हितः}}{\Afootnote{K; \skt{हितम् ॥}\ A}}\skt{ ।}\pend

\pstart\versepar \ledsidenote{\passref{6.40.69}}\skt{समङ्गा }\edtext{\edtext{\skt{धात\-की\-पु\-ष्पं\-स\-म्व\-ष्ठा\-लो\-ध्र\-मेव}}{\lemma{\skt{धात\-की\-पु\-ष्पं\-स\-म्व\-ष्ठा\-लो\-ध्रम्}}\Afootnote{K; \skt{धातकीपुष्पं मञ्जिष्ठा लोध्र}\ A}}\skt{ च}}{\lemma{\skt{एव च}}\Afootnote{K; \skt{मुस्तकम् ।}\ A}}\skt{ । }\pend
\pstart\versepar \skt{}\edtext{\skt{शाल्मलीवेष्टकं}}{\Afootnote{K; \skt{शाल्मलीवेष्टको}\ A}}\skt{ }\edtext{\skt{वापि}}{\Afootnote{K; \skt{रोध्रं}\ A}}\skt{ }\edtext{\edtext{\skt{विद\-ध्या\-ध्यो\-ग\-मु\-त्तम{\symfont ¦}}}{\lemma{\skt{विदध्याध्योगम्}}\Afootnote{K; \skt{वृक्षदाडि}\ A}}}{\lemma{\skt{उत्तम{\symfont ¦}}}\Afootnote{K; \skt{मयोस्त्वचौ ॥}\ A}}\skt{ ।}\pend

\pstart \ledsidenote{\passref{6.40.69.add}}\skt{समङ्गा धात\-की\-लो\-ध्र\-वृ\-क्ष\-दा\-डि\-म\-यो\-स्वचौ ।}\pend

\pstart\versepar \ledsidenote{\passref{6.40.70}}\edtext{\skt{आम्रा\-स्थि\-म\-ध्य\-लो\-ध्रञ्च}}{\Afootnote{K; \skt{आम्रास्थिमध्यं लोध्रं च}\ A}}\skt{ }\edtext{\skt{विल्व\-म\-ध्य\-प्रि\-य\-ङ्गवः}}{\Afootnote{K; \skt{बिल्वमध्यं प्रियङ्गवः ।}\ A}}\skt{ । }\pend
\pstart\versepar \skt{}\edtext{\skt{मधु\-क\-म्मे\-ष\-शृ\-ङ्गास्य}}{\lemma{\skt{मेषशृङ्गास्य}}\Afootnote{K; \skt{शृङ्गवेरं च}\ A}}\skt{ दीर्घ\-वृ\-न्त\-त्व\-गेव च ।}\pend

\pstart\versepar \ledsidenote{\passref{6.40.71}}\skt{चत्वार एते }\edtext{\skt{योगा स्युः}}{\Afootnote{K; \skt{योगाः स्युः}\ A}}\skt{ पक्वा\-ती\-सा\-र\-ना\-शनाः । }\pend
\pstart\versepar \skt{उक्ता }\edtext{\skt{मयोपयोज्यास्ते}}{\Afootnote{K; \skt{य उपयोज्यास्ते}\ A}}\skt{ }\edtext{\skt{क्षौर्द्रा\-ता\-ण्ड\-लां\-वुना}}{\Afootnote{K; \skt{सक्षौ\-द्रा\-स्त\-ण्डु\-ला\-म्बुना ॥}\ A}}\skt{ ।}\pend

\pstart\versepar \ledsidenote{\passref{6.40.72}}\edtext{\edtext{\skt{मुस्ता\-क\-षा\-य\-मेकं}}{\lemma{\skt{मुस्ताकषायम्}}\Afootnote{K; \skt{मौस्तं कषाय}\ A}}\skt{ वा}}{\lemma{\skt{एकं वा}}\Afootnote{K; \skt{मेकं वा}\ A}}\skt{ पेय\-म्म\-धु\-स\-मा\-युतं । }\pend
\pstart\versepar \skt{लोध्रां }\edtext{\edtext{\skt{वष्ठा\-प्रि\-य\-ङ्ग्ग्वा\-दी\-ङ्ग\-णा\-नेवं}}{\lemma{\skt{वष्ठा\-प्रि\-य\-ङ्ग्ग्वा\-दीङ्}}\Afootnote{K; \skt{बष्ठाप्रियङ्ग्वादी}\ A}}}{\lemma{\skt{गणानेवं}}\Afootnote{K; \skt{न्गणानेवं}\ A}}\skt{ प्रयोजयेत् ॥}\pend

\pstart\versepar \ledsidenote{\passref{6.40.73}}\edtext{\skt{पध्मां}}{\Afootnote{K; \skt{पद्मां}\ A}}\skt{ }\edtext{\skt{समाङ्गां}}{\Afootnote{K; \skt{समङ्गां}\ A}}\skt{ मधुकं }\edtext{\skt{विल्व\-जं\-वु\-श\-ला\-टुभिः}}{\Afootnote{K; \skt{बिल्व\-ज\-म्बू\-श\-लाटु च ।}\ A}}\skt{ । }\pend
\pstart\versepar \skt{}\edtext{\skt{पिबे\-त्त\-ण्डु\-ल\-तो\-येन}}{\Afootnote{A; \skt{पिवे\-त्त\-ण्डु\-ल\-तो\-येन}\ K}}\skt{ सक्षौद्रमगदं }\edtext{\skt{पर{\symfont ¦}}}{\Afootnote{K; \skt{करम् ॥}\ A}}\skt{ ।}\pend

\pstart\versepar \ledsidenote{\passref{6.40.74}}\edtext{\skt{कच्छ\-ग\-मू\-ल\-कल्कं}}{\Afootnote{K; \skt{कच्छु\-रा\-मू\-ल\-कल्कं}\ A}}\skt{ }\edtext{\skt{वाप्यु\-दु\-म्व\-र\-फ\-लो\-पम{\symfont ¦}}}{\Afootnote{K; \skt{वाप्यु\-दु\-म्ब\-र\-फ\-लो\-प\-मम् ।}\ A}}\skt{ । }\pend
\pstart\versepar \skt{पयस्या चन्दनं }\edtext{\skt{पद्मं}}{\Afootnote{K; \skt{पद्मा}\ A}}\skt{ }\edtext{\skt{शर्करा}}{\Afootnote{K; \skt{सितामु}\ A}}\skt{ }\edtext{\skt{नागकेशर{\symfont ¦}}}{\Afootnote{K; \skt{स्ताब्जकेशरम् ॥}\ A}}\skt{ ।}\pend

\pstart\versepar \ledsidenote{\passref{6.40.75}}\edtext{\skt{पक्वातीसार}}{\Afootnote{K; \skt{पक्वातिसारं}\ A}}\skt{ }\edtext{\skt{योगोयं}}{\Afootnote{K; \skt{योगोऽयं}\ A}}\skt{ }\edtext{\skt{जपेत्पीतं}}{\Afootnote{K; \skt{जयेत्पीतः}\ A}}\skt{ }\edtext{\skt{सशोणित{\symfont ¦}}}{\Afootnote{K; \skt{सशोणितम् ।}\ A}}\skt{ । }\pend
\pstart\versepar \skt{निरामरूपं }\edtext{\skt{शूला\-र्त्त\-ल\-ङ्घ\-ना\-ध्यैस्च}}{\Afootnote{K; \skt{शूलार्तं लङ्घनाद्यैश्च}\ A}}\skt{ }\edtext{\skt{कर्षित{\symfont ¦}}}{\Afootnote{K; \skt{कर्षितम् ॥}\ A}}\skt{ ।}\pend

\pstart \ledsidenote{\passref{6.40.76ab}}\skt{नरं }\edtext{\edtext{\skt{रूक्षमेवेक्षाग्निं}}{\lemma{\skt{रूक्षम्}}\Afootnote{K; \skt{रूक्ष}\ A}}}{\lemma{\skt{एवेक्षाग्निं}}\Afootnote{K; \skt{मवेक्ष्याग्निं}\ A}}\skt{ सक्षारं }\edtext{\skt{पाचयेद्धृत{\symfont ¦}}}{\Afootnote{K; \skt{पाययेद्घृतम् ।}\ A}}\skt{ ।}\pend

\pstart \ledsidenote{\passref{6.40.79cd}}\edtext{\skt{पयस्तैलश्च}}{\Afootnote{K; \skt{पयो घृतं च}\ A}}\skt{ मधु च }\edtext{\skt{कम्विल्वजंवु}}{\Afootnote{K; \emph{om.}\ A}}\skt{ }\edtext{\skt{पिवे\-च्छू\-लै\-रू\-प\-द्रुतः}}{\Afootnote{K; \skt{पिबे\-च्छू\-लै\-र\-भि\-द्रुतः ॥}\ A}}\skt{ ।}\pend

\pstart\versepar \ledsidenote{\passref{6.40.80}}\edtext{\skt{सिता\-ज\-मो\-द\-क\-ट्वा\-ङ्ग\-म\-धु\-कै\-र\-व\-चू\-र्ण्णितः}}{\Afootnote{K; \skt{सिता\-ज\-मो\-द\-क\-ट्व\-ङ्ग\-म\-धु\-कै\-र\-व\-चू\-र्णि\-तम् ।}\ A}}\skt{ । }\pend
\pstart\versepar \skt{अवेदनं सुसंपक्वं }\edtext{\skt{दीप्ताग्ने}}{\Afootnote{K; \skt{दीप्ताग्नेः}\ A}}\skt{ }\edtext{\skt{रचिरोत्थित{\symfont ¦}}}{\Afootnote{K; \skt{सुचिरोत्थितम् ॥}\ A}}\skt{ ॥}\pend

\pstart\versepar \ledsidenote{\passref{6.40.81}}\edtext{\edtext{\skt{नाना\-व\-र्ण\-म\-ती\-सारं}}{\lemma{\skt{नानावर्णम्}}\Afootnote{K; \skt{नानावर्ण}\ A}}}{\lemma{\skt{अतीसारं}}\Afootnote{K; \skt{मतीसारं}\ A}}\skt{ }\edtext{\skt{पुट\-पा\-कै\-रु\-पा\-च\-रेत}}{\lemma{\skt{उपाचरेत}}\Afootnote{K; \skt{उपाचरे}\ A}}\skt{ । }\pend
\pstart\versepar \skt{}\edtext{\skt{त्वकि\-ण्ड\-दी\-र्घ\-वृ\-न्तस्य}}{\Afootnote{K; \skt{त् । त्वक्पिण्डं दीर्घवृन्तस्य}\ A}}\skt{ }\edtext{\skt{काश्म\-री\-च्छ\-द\-ना\-वृतं}}{\Afootnote{K; \skt{पद्म\-के\-स\-र\-सं\-यु\-तम् ॥}\ A}}\skt{ ।}\pend

\pstart \ledsidenote{\passref{6.40.82cd}}\skt{मृदावलिप्तं }\edtext{\edtext{\edtext{\skt{सुकृ\-त\-म\-ङ्गा\-रे\-ष्व\-व\-कु\-ल\-येत}}{\lemma{\skt{सुकृतम्}}\Afootnote{K; \skt{सुकृत}\ A}}}{\lemma{\skt{अङ्गारेष्व्}}\Afootnote{K; \skt{मङ्गारेष्व्}\ A}}}{\lemma{\skt{अवकुलयेत}}\Afootnote{K; \skt{अवकूलयेत् ॥}\ A}}\skt{ ।}\pend

\pstart\versepar \ledsidenote{\passref{6.40.83}}\edtext{\edtext{\skt{स्विन्नमुधृत्य}}{\lemma{\skt{स्विन्नम्}}\Afootnote{K; \skt{स्विन्न}\ A}}}{\lemma{\skt{उधृत्य}}\Afootnote{K; \skt{मुद्धृत्य}\ A}}\skt{ निष्पीड्य }\edtext{\edtext{\skt{रसमादाय}}{\lemma{\skt{रसम्}}\Afootnote{K; \skt{रस}\ A}}\skt{ तन्ततः}}{\lemma{\skt{आदाय तन्}}\Afootnote{K; \skt{मादाय तं}\ A}}\skt{ । }\pend
\pstart\versepar \skt{}\edtext{\skt{शीतीकृत्वा}}{\Afootnote{K; \skt{शीतं}\ A}}\skt{ मधुयुतं }\edtext{\skt{पाय\-ये\-दु\-द\-रा\-मये}}{\Afootnote{K; \skt{कृत्वा पाय\-ये\-तो\-द\-रा\-मये ॥}\ A}}\skt{ ।}\pend

\pstart\versepar \ledsidenote{\passref{6.40.84}}\skt{जीव\-न्ती\-मे\-ष\-शृ\-ङ्ग्या\-दि\-ष्वेवं }\edtext{\skt{वृत्तेषु}}{\Afootnote{K; \skt{द्रव्येषु}\ A}}\skt{ कारयेत् । }\pend
\pstart\versepar \skt{तित्तिरिं लुञ्चितं }\edtext{\skt{गोधृ\-निः\-कृ\-ष्या\-न्त्रन्तु}}{\Afootnote{K; \skt{सम्यक्निष्कृष्टान्त्रं तु}\ A}}\skt{ }\edtext{\skt{पूरयेत}}{\Afootnote{K; \skt{पूरयेत् ॥}\ A}}\skt{ ।}\edtext{}{\lemma{\skt{कारयेत् । तित्तिरिं}}\Afootnote{K; \skt{साधयेत् । तित्तिरिं}\ A}}\pend

\pstart\versepar \ledsidenote{\passref{6.40.85}}\edtext{\skt{न्यग्रो\-धा\-धि\-त्व\-चा\-कल्कैः}}{\Afootnote{K; \skt{न्यग्रोधादित्वचां कल्कैः}\ A}}\skt{ }\edtext{\skt{पूर्व\-व\-च्चा\-व\-कु\-ल\-येत}}{\Afootnote{K; \skt{पूर्व\-व\-च्चा\-व\-कू\-ल\-येत् । र}\ A}}\skt{ । }\pend
\pstart\versepar \skt{}\edtext{\edtext{\skt{रसमादाय}}{\lemma{\skt{रसम्}}\Afootnote{K; \skt{स}\ A}}}{\lemma{\skt{आदाय}}\Afootnote{K; \skt{मादाय}\ A}}\skt{ तस्याथ सुस्विन्नस्य }\edtext{\skt{समाक्षिक{\symfont ¦}}}{\Afootnote{K; \skt{समाक्षिकम् ॥}\ A}}\skt{ ।}\pend

\pstart \ledsidenote{\passref{6.40.86ab}}\skt{शर्करोपहितं शीतं }\edtext{\skt{पाय\-ये\-दु\-द\-रा\-मये}}{\Afootnote{K; \skt{पाय\-ये\-तो\-द\-रा\-मये ।}\ A}}\skt{ ।}\pend

\pstart \ledsidenote{\passref{6.40.88cd}}\skt{एवं }\edtext{\skt{प्ररोह}}{\Afootnote{K; \skt{प्ररोहैः}\ A}}\skt{ }\edtext{\skt{कुर्वीतन्य}}{\Afootnote{K; \skt{कुर्वीत व}\ A}}\skt{ }\edtext{\edtext{\skt{योपा\-ध्यै\-र्वि\-च\-क्षणः}}{\lemma{\skt{योपाध्यैर्}}\Afootnote{K; \skt{टादीनं}\ A}}}{\lemma{\skt{विचक्षणः}}\Afootnote{K; \skt{विधानवत् ॥}\ A}}\skt{ ।}\pend

\pstart \ledsidenote{\passref{6.40.89ab}}\edtext{\skt{पुट\-पा\-का\-न्य\-था\-योगां}}{\Afootnote{K; \skt{पुट\-पा\-का\-न्य\-था\-योगं}\ A}}\skt{ }\edtext{\skt{जाङ्गलोपहितां}}{\Afootnote{K; \skt{जाङ्गलोपहिता}\ A}}\skt{ }\edtext{\skt{शुभा{\symfont ¦}}}{\Afootnote{K; \skt{न्शुभान् ।}\ A}}\skt{ ।}\pend

\pstart\versepar \ledsidenote{\passref{6.40.90}}\skt{कौटजं }\edtext{\edtext{\skt{फाणितश्चापि}}{\lemma{\skt{फाणितश्}}\Afootnote{K; \skt{फाणितं}\ A}}}{\lemma{\skt{चापि}}\Afootnote{K; \skt{वापि}\ A}}\skt{ }\edtext{\skt{हंन्त्य\-ती\-सा\-र\-मो\-जसा}}{\Afootnote{K; \skt{हन्त्य\-ती\-सा\-र\-मो\-जसा ।}\ A}}\skt{ । }\pend
\pstart\versepar \skt{}\edtext{\skt{अम्व\-ष्ठा\-दि\-म\-धु\-युतं}}{\Afootnote{K; \skt{अम्ब\-ष्ठा\-दि\-म\-धु\-युतं}\ A}}\skt{ पि अर\-लु\-त्व\-क्प्रि\-यङ्गु च मधु\-क\-दा\-डि\-मा\-ङ्कुरां । }\pend
\pstart\versepar \skt{आवाप्य }\edtext{\skt{पिष्ठ्वा विप\-चे\-द्य\-वा\-गू\-न्द\-ध्नि\-ता\-म्पि\-बेत्}}{\Afootnote{\emph{om.}\ A K}}\skt{ ।}\edtext{}{\lemma{\skt{पि अर\-लु\-त्व\-क्प्रि\-यङ्गु च मधु\-क\-दा\-डि\-मा\-ङ्कुरां । आवाप्य}}\Afootnote{\skt{पिप्प\-ल्या\-दि\-स\-म\-न्वि\-तम् ॥}\ A; \skt{पिल्या\-दि\-स\-मा\-युत{\symfont ¦} ।}\ K}}\pend

\pstart \ledsidenote{\passref{6.40.93ab}}\edtext{\skt{सर्वातीसारान्तां हि}}{\Afootnote{K; \skt{एषा सर्वानतीसारा}\ A}}\skt{ }\edtext{\skt{हन्ति}}{\Afootnote{K; \skt{न्हन्ति}\ A}}\skt{ }\edtext{\skt{पक्वानशंशयः}}{\Afootnote{K; \skt{पक्वा\-न\-सं\-श\-यम् ।}\ A}}\skt{ ।}\pend

\pstart \ledsidenote{\passref{6.40.98cd}}\skt{कृतं शाल्मलिवृन्तेषु कषायं }\edtext{\skt{हिमसंज्ञित{\symfont ¦}}}{\Afootnote{K; \skt{हिमसंज्ञितम् ॥}\ A}}\skt{ ।}\pend

\pstart\versepar \ledsidenote{\passref{6.40.99}}\edtext{\skt{निशा\-प\-र्यु\-षि\-त\-म्पे\-य\-म्सा\-क्षि\-के\-श\-स\-मा\-युत{\symfont ¦}}}{\lemma{\skt{साक्षि\-के\-श\-स\-मा\-युत{\symfont ¦}}}\Afootnote{K; \skt{सक्षौद्रं मधुकान्वितम् ।}\ A}}\skt{ । }\pend
\pstart\versepar \skt{}\edtext{\skt{विव\-ध\-वा\-त\-वि\-ट्शू\-ल\-प\-रीतः}}{\Afootnote{K; \skt{विब\-द्ध\-वा\-त\-वि\-ट्शू\-ल\-प\-रीतः}\ A}}\skt{ }\edtext{\skt{संप्रवाहिकः}}{\Afootnote{K; \skt{सप्रवाहिकः ॥}\ A}}\skt{ ।}\pend

\pstart\versepar \ledsidenote{\passref{6.40.100}}\skt{सरक्तपित्तश्च पयः }\edtext{\skt{पिबेत्तृष्णा}}{\Afootnote{A; \skt{पिवेत्तृष्णा}\ K}}\skt{ }\edtext{\skt{समान्वितः}}{\Afootnote{K; \skt{समन्वितः ।}\ A}}\skt{ । }\pend
\pstart\versepar \skt{}\edtext{\skt{यथामृतन्तथा}}{\lemma{\skt{तथा}}\Afootnote{K; \skt{यथा}\ A}}\skt{ }\edtext{\skt{क्षीर\-म\-ती\-सा\-रेषु}}{\Afootnote{A K\textsuperscript{pc}; \skt{क्षीर\-म\-ती\-सा\-रेत{\symfont ¦} ।}\ K\textsuperscript{ac}}}\skt{ }\edtext{\skt{त{\symfont ¦}}}{\Afootnote{K\textsuperscript{pc}; \skt{पूजितम् ॥}\ A; \emph{om.}\ K\textsuperscript{ac}}}\skt{ ।}\pend

\pstart\versepar \ledsidenote{\passref{6.40.101}}\skt{चिरोत्थितेषु तत्पेयमपां }\edtext{\skt{भारान}}{\Afootnote{K; \skt{भागैस्त्रि}\ A}}\skt{ }\edtext{\skt{साधित{\symfont ¦}}}{\Afootnote{K; \skt{भिः शृतम् ।}\ A}}\skt{ । }\pend
\pstart\versepar \skt{दोषशेषं हरेत्तद्धि तस्मात्पथ्यतमं }\edtext{\skt{स्मृत{\symfont ¦}}}{\Afootnote{K; \skt{स्मृतम् ॥}\ A}}\skt{ ॥}\pend

\pstart \ledsidenote{\passref{6.40.102ab}}\edtext{\skt{हितं}}{\Afootnote{K; \skt{हितः}\ A}}\skt{ स्नेहविरेको वा }\edtext{\skt{वस्तयः}}{\Afootnote{K; \skt{बस्तयः}\ A}}\skt{ पिच्छिलाश्च ये ।}\pend

\pstart\versepar \ledsidenote{\passref{6.40.103}}\skt{शकृता यस्तु संसृ\-ष्ट\-म\-ति\-सा\-र्येत }\edtext{\skt{शोणित{\symfont ¦}}}{\Afootnote{K; \skt{शोणितम् ।}\ A}}\skt{ । }\pend
\pstart\versepar \skt{}\edtext{\edtext{\skt{प्राक्पश्चाद्वा}}{\lemma{\skt{प्राक्पश्}}\Afootnote{K; \skt{प्राक्प}\ A}}}{\lemma{\skt{चाद्वा}}\Afootnote{K; \skt{श्चाद्वा}\ A}}\skt{ पुरीषस्य }\edtext{\skt{सरुक्षः}}{\Afootnote{K; \skt{सरुक्स}\ A}}\skt{ }\edtext{\skt{परिकीर्त्तिकं}}{\Afootnote{K; \skt{परिकर्तिकः ॥}\ A}}\skt{ ।}\pend

\pstart\versepar \ledsidenote{\passref{6.40.104}}\edtext{\skt{क्षीरशुक्लाश्रितं}}{\Afootnote{K; \skt{क्षीरिशुङ्गाशृतं}\ A}}\skt{ सर्पिः }\edtext{\skt{पिबे\-त्स\-क्षौ\-द्र\-श\-र्कर{\symfont ¦}}}{\Afootnote{\skt{पिबे\-त्स\-क्षौ\-द्र\-श\-र्क\-रम् ।}\ A; \skt{पिवे\-त्स\-क्षौ\-द्र\-श\-र्कर{\symfont ¦} ।}\ K}}\skt{ । }\pend
\pstart\versepar \skt{}\edtext{\skt{दार्वी\-त्व\-क्पि\-प्प\-ली\-शु\-ण्ठी\-ला\-क्षा\-श\-कु\-र्य\-र्वै\-घृत{\symfont ¦}}}{\Afootnote{K; \skt{दार्वी\-त्व\-क्पि\-प्प\-ली\-शु\-ण्ठी\-ला\-क्षा\-श\-क्र\-य\-वै\-र्घृ\-तम् ॥}\ A}}\skt{ ।}\pend

\pstart\versepar \ledsidenote{\passref{6.40.105}}\skt{संयु\-क्त\-म्भ\-द्र\-रो\-हिण्या }\edtext{\skt{पक्व\-म्पे\-या\-वि\-मि\-श्रित{\symfont ¦}}}{\lemma{\skt{पेयाविमिश्रित{\symfont ¦}}}\Afootnote{K; \skt{पेया\-दि\-मि\-श्रि\-तम् ।}\ A}}\skt{ । }\pend
\pstart\versepar \skt{}\edtext{\skt{त्रिदो\-ष\-म\-प्य\-ती\-सारं}}{\Afootnote{K; \skt{त्रिदो\-ष\-म\-प्य\-ती\-सारं}\ A}}\skt{ पीतं हन्ति }\edtext{\skt{सुदारण{\symfont ¦}}}{\Afootnote{K; \skt{सुदारुणम् ॥}\ A}}\skt{ ।}\pend

\pstart\versepar \ledsidenote{\passref{6.40.106}}\skt{गौरवे वमनम्पथ्यं यस्य }\edtext{\skt{स्यात्प्रवलः}}{\Afootnote{K; \skt{स्यात्प्रबलः}\ A}}\skt{ कफः । }\pend
\pstart\versepar \skt{ज्वरे }\edtext{\skt{दाघथविड्वन्ध}}{\Afootnote{K; \skt{दाहे सविङ्बन्धे}\ A}}\skt{ }\edtext{\skt{मारू\-ता\-द्र\-क्त\-पि\-त्त\-वत्}}{\Afootnote{K; \skt{मारु\-ता\-द्र\-क्त\-पि\-त्त\-वत् ॥}\ A}}\skt{ ।}\pend

\pstart \ledsidenote{\passref{6.40.107cd}}\edtext{\edtext{\skt{कार्यमास्थापनं}}{\lemma{\skt{कार्यम्}}\Afootnote{K; \skt{कार्य}\ A}}}{\lemma{\skt{आस्थापनं}}\Afootnote{K; \skt{मास्थापनं}\ A}}\skt{ क्षिप्रं }\edtext{\skt{सदा}}{\Afootnote{K; \skt{तथा}\ A}}\skt{ }\edtext{\edtext{\skt{चाप्यानुवामन{\symfont ¦}}}{\lemma{\skt{चाप्य्}}\Afootnote{K; \skt{चै}\ A}}}{\lemma{\skt{आनुवामन{\symfont ¦}}}\Afootnote{K; \skt{वानुवासनम् ॥}\ A}}\skt{ ।}\pend

\pstart\versepar \ledsidenote{\passref{6.40.109}}\skt{गुदपाकस्तु पित्तेन यस्य स्यादहिताशिनः । }\pend
\pstart\versepar \skt{तस्य पित्तहराः सेका\-स्त\-त्सि\-द्धा\-श्चा\-नु\-वा\-सनाः ।}\pend

\pstart\versepar \ledsidenote{\passref{6.40.110}}\edtext{\skt{दधि\-म\-ण्ड\-सु\-रा\-वि\-ल्व\-सिद्धं}}{\Afootnote{K; \skt{दधि\-म\-ण्ड\-सु\-रा\-बि\-ल्व\-सिद्धं}\ A}}\skt{ तैलं समारुते । }\pend
\pstart\versepar \skt{भोजने च हितं क्षीरं }\edtext{\skt{कच्छु\-रा\-मू\-ल\-सा\-धिर{\symfont ¦}}}{\Afootnote{K; \skt{कच्छु\-रा\-मू\-ल\-सा\-धि\-तम् ॥}\ A}}\skt{ ।}\pend

\pstart \ledsidenote{\passref{6.40.111cd}}\edtext{\edtext{\edtext{\skt{रुज\-या\-श्चा\-प्य\-शाम्यां}}{\lemma{\skt{रुजयाश्}}\Afootnote{K; \skt{यदा वायुर्विबद्धश्}\ A}}}{\lemma{\skt{चाप्य्}}\Afootnote{K; \skt{च}\ A}}}{\lemma{\skt{अशाम्यां}}\Afootnote{K; \emph{om.}\ A}}\skt{ }\edtext{\skt{त्वापि\-च्छा\-व\-स्ति\-धी\-यते}}{\Afootnote{K; \skt{पिच्छाबस्तिस्तदा हितः ॥}\ A}}\skt{ ।}\pend

\pstart\versepar \ledsidenote{\passref{6.40.112}}\edtext{\skt{प्रायेन}}{\Afootnote{K; \skt{प्रायेण}\ A}}\skt{ }\edtext{\skt{गुददौर्वल्यं}}{\Afootnote{K; \skt{गुददौर्बल्यं}\ A}}\skt{ }\edtext{\skt{दीर्घ\-का\-ला\-ति\-सा\-रिणा{\symfont ¦}}}{\Afootnote{K; \skt{दीर्घ\-का\-ला\-ति\-सा\-रि\-णाम् ।}\ A}}\skt{ । }\pend
\pstart\versepar \skt{}\edtext{\skt{भवेत्तस्याधितं}}{\Afootnote{K; \skt{भवेत्तस्माद्धितं}\ A}}\skt{ तेषाङ्गुदे }\edtext{\skt{तैलावचारण{\symfont ¦}}}{\Afootnote{K; \skt{तैला\-व\-चा\-र\-णम् ॥}\ A}}\skt{ ।}\pend

\pstart\versepar \ledsidenote{\passref{6.40.113}}\edtext{\skt{कपि\-त्थ\-शा\-ल्म\-ली\-फ\-ञ्जी\-व\-टा\-का\-र्पा\-सि\-दा\-डिमा}}{\Afootnote{K; \skt{कपि\-त्थ\-शा\-ल्म\-ली\-फ\-ञ्जी\-व\-ट\-का\-र्पा\-स\-दा\-डिमाः ।}\ A}}\skt{ । }\pend
\pstart\versepar \skt{यूथिका कच्छुरा }\edtext{\skt{शेलू}}{\Afootnote{K; \skt{शेलुः}\ A}}\skt{ }\edtext{\skt{सणामूर्वा च}}{\Afootnote{K; \skt{शणश्चुच्चूश्च}\ A}}\skt{ दाधिकाः ।}\pend

\pstart \ledsidenote{\passref{6.40.117cd}}\skt{यो रक्तं शकृतः }\edtext{\skt{पूर्वां}}{\Afootnote{K; \skt{पूर्वं}\ A}}\skt{ }\edtext{\edtext{\skt{पश्चाद्वा}}{\lemma{\skt{पश्}}\Afootnote{K; \skt{प}\ A}}}{\lemma{\skt{चाद्वा}}\Afootnote{K; \skt{श्चाद्वा}\ A}}\skt{ प्रतिसार्यते ॥}\pend

\pstart \ledsidenote{\passref{6.40.118}}\skt{स }\edtext{\skt{पल्ल\-वै\-र्व\-टा\-दीनां}}{\Afootnote{K; \skt{पल्ल\-वै\-र्व\-टा\-दीनं}\ A}}\skt{ सस\-र्पि\-स्सा\-धि\-त\-म्पयः }\edtext{\skt{पिबे\-त्स\-श\-र्क\-रा\-क्षौ\-द्र\-म\-थ\-वा\-प्य\-भि\-मथ्य}}{\Afootnote{A; \skt{पिवे\-त्स\-श\-र्क\-रा\-क्षौ\-द्र\-म\-थ\-वा\-प्य\-भि\-मथ्य}\ K}}\skt{ तत् ।}\pend

\pstart \ledsidenote{\passref{6.40.119ab}}\edtext{\edtext{\skt{नवनीतमतो}}{\lemma{\skt{नवनीतम्}}\Afootnote{K; \skt{नवनीत}\ A}}}{\lemma{\skt{अतो}}\Afootnote{K; \skt{मथो}\ A}}\skt{ }\edtext{\skt{लिह्यात्चक्रं}}{\Afootnote{K; \skt{लिह्यात्तक्रं}\ A}}\skt{ }\edtext{\skt{येन}}{\Afootnote{K; \skt{चानु}\ A}}\skt{ }\edtext{\edtext{\skt{पिवेन्नरः}}{\lemma{\skt{पिवेन्}}\Afootnote{K; \skt{पिबे}\ A}}}{\lemma{\skt{नरः}}\Afootnote{K; \skt{त्ततः ।}\ A}}\skt{ ।}\pend

\pstart \ledsidenote{\passref{6.40.115cd}}\skt{तिलकल्को हितश्चापि मोद्गौ मुद्गरसस्तथा ॥}\pend

\pstart \ledsidenote{\passref{6.40.119cd}}\edtext{\skt{पिया\-ल\-श\-ल्ल\-की\-ज\-म्वु\-शा\-ल्म\-ली\-ति\-नि\-श\-त्वचः}}{\Afootnote{K; \skt{प्रिया\-ल\-शा\-ल्म\-ली\-प्ल\-क्ष\-श\-ल्ल\-की\-ति\-नि\-श\-त्वचः ॥}\ A}}\skt{ ।}\pend

\pstart\versepar \ledsidenote{\passref{6.40.120}}\skt{क्षीरे }\edtext{\skt{शमृदिता}}{\Afootnote{K; \skt{विमृदिताः}\ A}}\skt{ }\edtext{\skt{पीतास्सक्षौर्द्रा}}{\lemma{\skt{सक्षौर्द्रा}}\Afootnote{K; \skt{सक्षौद्रा}\ A}}\skt{ रक्तनाशनाः । }\pend
\pstart\versepar \skt{मधुकं }\edtext{\skt{शर्करालोध्रं}}{\Afootnote{K; \skt{शर्करां लोध्रं}\ A}}\skt{ }\edtext{\edtext{\skt{पयस्यामथ}}{\lemma{\skt{पयस्याम्}}\Afootnote{K; \skt{पयस्या}\ A}}}{\lemma{\skt{अथ}}\Afootnote{K; \skt{मथ}\ A}}\skt{ }\edtext{\skt{शारिवा{\symfont ¦}}}{\Afootnote{K; \skt{सारिवाम् ॥}\ A}}\skt{ ।}\pend

\pstart\versepar \ledsidenote{\passref{6.40.122}}\edtext{\skt{पिबेत्पद्माश्च}}{\Afootnote{\skt{पिबेत्पद्मां च}\ A; \skt{पिवेत्पद्माश्च}\ K}}\skt{ }\edtext{\skt{दुग्धे\-न\-च्छा\-रो\-ना\-सृ\-क्प्र\-शा\-न्तये}}{\Afootnote{K; \skt{दुग्धेन छागे\-ना\-सृ\-क्प्र\-शा\-न्तये ।}\ A}}\skt{ ॥ }\pend
\pstart\versepar \skt{}\edtext{\skt{शर्क\-रो\-त्प\-ल\-लो\-ध्रानि}}{\Afootnote{K; \skt{शर्क\-रो\-त्प\-ल\-लो\-ध्राणि}\ A}}\skt{ समङ्गा मधुकन्तिलाः ।}\pend

\pstart \ledsidenote{\passref{6.40.123ab}}\skt{तिलाः }\edtext{\skt{कृष्णा\-स्स\-म\-प्र\-ष्ट्याह्वा}}{\lemma{\skt{समप्रष्ट्याह्वा}}\Afootnote{K; \skt{सयष्ट्याह्वाः}\ A}}\skt{ समङ्गा चोत्पलानि च ।}\pend

\pstart \ledsidenote{\passref{6.40.121cd}}\edtext{\skt{मंजिष्ठाशर्करा}}{\Afootnote{K; \skt{मञ्जिष्ठां सारिवां}\ A}}\skt{ लोध्रं पद्मकं }\edtext{\skt{कुमुदोत्पलैः}}{\Afootnote{K; \skt{कुमुदोत्पलम् ॥}\ A}}\skt{ ।}\pend

\pstart \ledsidenote{\passref{6.40.122.add}}\edtext{\skt{पिबेद्पद्माश्च}}{\Afootnote{\skt{पिवेद्पद्माश्च}\ K}}\skt{ दुग्धे\-न\-च्छा\-गे\-ना\-सृ\-क्प्र\-शा\-न्तये ।}\pend

\pstart \ledsidenote{\passref{6.40.123cd}}\edtext{\edtext{\skt{लोध्रन्तिलातमो}}{\lemma{\skt{लोध्रन्}}\Afootnote{K; \emph{om.}\ A}}}{\lemma{\skt{तिलातमो}}\Afootnote{K; \skt{तिला मो}\ A}}\skt{ च }\edtext{\edtext{\skt{रसान्तथैव}}{\lemma{\skt{रसान्}}\Afootnote{K; \skt{रसो}\ A}}}{\lemma{\skt{तथैव}}\Afootnote{K; \skt{लोध्रं तथैव}\ A}}\skt{ }\edtext{\skt{मधुकार्यणैः}}{\Afootnote{K; \skt{मधुकोत्पलम् ॥}\ A}}\skt{ ।}\pend

\pstart\versepar \ledsidenote{\passref{6.40.124}}\skt{कच्छुरा तिलकल्कश्च }\edtext{\skt{योगाम्चत्वार}}{\lemma{\skt{योगाम्}}\Afootnote{K; \skt{योगाश्}\ A}}\skt{ एव च । }\pend
\pstart\versepar \skt{आजेन पयसा }\edtext{\skt{पेया}}{\Afootnote{K; \skt{पेयाः}\ A}}\skt{ }\edtext{\skt{सरक्तो}}{\Afootnote{K; \skt{सरक्ते}\ A}}\skt{ }\edtext{\skt{मधुसंयुता}}{\Afootnote{K; \skt{मधुसंयुताः ॥}\ A}}\skt{ ।}\pend

\pstart\versepar \ledsidenote{\passref{6.40.126}}\edtext{\skt{वसवानां}}{\Afootnote{K; \skt{कोशकारं}\ A}}\skt{ }\edtext{\skt{घृत\-ला\-ज\-चू\-र्ण\-सि\-ता\-मधु}}{\Afootnote{K; \skt{घृते भृष्टं लाजचूर्णं सिता मधु ।}\ A}}\skt{ । }\pend
\pstart\versepar \skt{}\edtext{\skt{सशू\-ल\-र\-क्त\-पित्तं तु}}{\Afootnote{K; \skt{सशूलं रक्तपित्तोत्थं}\ A}}\skt{ लीढं }\edtext{\skt{हन्त्युदरामय{\symfont ¦}}}{\Afootnote{K; \skt{हन्त्यु\-द\-रा\-म\-यम् ॥}\ A}}\skt{ ॥}\pend

\pstart\versepar \ledsidenote{\passref{6.40.125}}\skt{द्रवे }\edtext{\skt{सरक्तो}}{\Afootnote{K; \skt{सरक्ते}\ A}}\skt{ }\edtext{\skt{सरति}}{\Afootnote{K; \skt{स्रवति}\ A}}\skt{ }\edtext{\skt{वालविल्वं}}{\Afootnote{K; \skt{बालबिल्वं}\ A}}\skt{ सफाणितं । }\pend
\pstart\versepar \skt{सक्षौद्रतैलं प्रागेव लिह्यादाशु }\edtext{\skt{मुदितं}}{\Afootnote{K; \skt{हितं}\ A}}\skt{ }\edtext{\skt{कर{\symfont ¦}}}{\Afootnote{K; \skt{हि तत् ॥}\ A}}\skt{ ।}\pend

\pstart\versepar \ledsidenote{\passref{6.40.130}}\edtext{\skt{दीप्ता\-ग्नि\-म्चा\-ष्य\-ती\-सारे}}{\Afootnote{K; \skt{सक्त\-वि\-ड्दो\-ष\-ब\-हुलं दीप्ताग्निर्योऽतिसा}\ A}}\skt{ }\edtext{\skt{वर्धमान}}{\Afootnote{K; \emph{om.}\ A}}\skt{ }\edtext{\skt{वलाम्विते}}{\Afootnote{K; \skt{र्य}\ A}}\skt{ । }\pend
\pstart\versepar \skt{}\edtext{\skt{कृष्णा\-वि\-डा\-ङ्ग\-तृ\-फ\-ला\-क\-षा\-येण}}{\Afootnote{K; \skt{ते । विड\-ङ्ग\-त्रि\-फ\-ला\-कृ\-ष्णा\-क\-षा\-यैस्तं}\ A}}\skt{ }\edtext{\skt{विरेचयेत}}{\Afootnote{K; \skt{विरेचयेत् ॥}\ A}}\skt{ ।}\pend

\pstart\versepar \ledsidenote{\passref{6.40.131}}\skt{अथ वैरण्डसिद्धेन पयसा केवलेन वा । }\pend
\pstart\versepar \skt{}\edtext{\edtext{\skt{यवागूम्चास्य}}{\lemma{\skt{यवागूम्}}\Afootnote{K; \skt{यवागू}\ A}}}{\lemma{\skt{चास्य}}\Afootnote{K; \emph{om.}\ A}}\skt{ }\edtext{\edtext{\skt{वित\-रे\-द्वा\-त\-घ्नै\-र्दी\-पनैः}}{\lemma{\skt{वित\-रे\-द्वा\-त\-घ्नैर्}}\Afootnote{K; \skt{र्वितरेच्चास्य वातघ्नै}\ A}}}{\lemma{\skt{दीपनैः}}\Afootnote{K; \skt{र्दीपनैः}\ A}}\skt{ }\edtext{\skt{कृतां}}{\Afootnote{K; \skt{कृताः ॥}\ A}}\skt{ ।}\pend

\pstart\versepar \ledsidenote{\passref{6.40.132}}\edtext{\skt{दीप्ताग्निर्ः}}{\Afootnote{K; \skt{दीप्ताग्निर्नि}\ A}}\skt{ }\edtext{\skt{पुरीषा यः}}{\Afootnote{K; \skt{ष्पुरीषो यः}\ A}}\skt{ }\edtext{\skt{सूवे\-च्छ\-न\-स\-वे\-दन{\symfont ¦}}}{\Afootnote{K; \skt{सार्यते फेनिलं शकृत् । स}\ A}}\skt{ । }\pend
\pstart\versepar \skt{}\edtext{\skt{पिबेत्फाणितं}}{\Afootnote{A; \skt{पिवेत्फाणितं}\ K}}\skt{ }\edtext{\skt{शुण्ठीदधितैलं}}{\Afootnote{K; \skt{शुण्ठीदधितैल}\ A}}\skt{ }\edtext{\skt{घृतपयः}}{\Afootnote{K; \skt{पयोघृतम् ॥}\ A}}\skt{ ।}\pend

\pstart\versepar \ledsidenote{\passref{6.40.133}}\edtext{\skt{स्विन्नाम्वा}}{\Afootnote{K; \skt{स्विन्नानि}\ A}}\skt{ }\edtext{\skt{गुड\-तै\-ल\-न\-भ\-क्ष\-ये\-द्व\-द\-राणि वा}}{\Afootnote{K; \skt{गुडतैलाभ्यां भक्ष\-ये\-द्ब\-द\-राणि च ।}\ A}}\skt{ । }\pend
\pstart\versepar \skt{}\edtext{\skt{सस्विन्नां}}{\Afootnote{K; \skt{स्विन्नानि}\ A}}\skt{ पिष्टवद्वापि }\edtext{\skt{समा}}{\Afootnote{K; \skt{समं}\ A}}\skt{ }\edtext{\skt{विल्वशलादुभिः}}{\Afootnote{K; \skt{बिल्वशलाटुभिः ॥}\ A}}\skt{ ।}\pend

\pstart\versepar \ledsidenote{\passref{6.40.134}}\skt{दध्नोपयुज्य }\edtext{\edtext{\skt{कुल्मा\-षा\-म्श्वे\-ता\-म\-नु\-पि\-बे\-त्सुरा{\symfont ¦}}}{\lemma{\skt{कुल्माषाम्}}\Afootnote{K; \skt{कुल्माषा}\ A}}}{\lemma{\skt{श्वेता\-म\-नु\-पि\-बे\-त्सुरा{\symfont ¦}}}\Afootnote{\skt{न्श्वेता\-म\-नु\-पि\-बे\-त्सु\-राम् ।}\ A; \skt{श्वेता\-म\-नु\-पि\-वे\-त्सुरा{\symfont ¦} ।}\ K}}\skt{ । }\pend
\pstart\versepar \skt{शशमांसं }\edtext{\skt{सरुधिरधिर{\symfont ¦}}}{\Afootnote{K; \skt{सरुधिरं}\ A}}\skt{ }\edtext{\skt{समांशं}}{\Afootnote{K; \skt{समङ्गां}\ A}}\skt{ सघृतं दधि ।}\pend

\pstart\versepar \ledsidenote{\passref{6.40.135}}\edtext{\skt{विपाच्य}}{\Afootnote{K; \skt{खादेद्विपाच्य}\ A}}\skt{ }\edtext{\skt{शोदको}}{\Afootnote{K; \skt{से}\ A}}\skt{ }\edtext{\skt{वेच्च}}{\Afootnote{K; \skt{वेत}\ A}}\skt{ मृद्वन्नं शकृतः क्षये ॥ }\pend
\pstart\versepar \skt{संस्कृतो यमके }\edtext{\skt{माषायवकोलः}}{\Afootnote{K; \skt{माषयवकोल}\ A}}\skt{ रसः शुभः ।}\pend

\pstart\versepar \ledsidenote{\passref{6.40.136}}\edtext{\skt{भोजनार्थे}}{\Afootnote{K; \skt{भोजनार्थं}\ A}}\skt{ प्रदातव्यो }\edtext{\skt{दधि\-डा\-दा\-म\-सा\-रिवः}}{\Afootnote{K; \skt{दधि\-दा\-डि\-म\-सा\-धितः ।}\ A}}\skt{ । }\pend
\pstart\versepar \skt{}\edtext{\skt{गुडा\-वि\-ल्व\-श\-ला\-दूनि}}{\Afootnote{K; \skt{विडं बिल्वशलाटूनि}\ A}}\skt{ }\edtext{\edtext{\skt{दाडिमश्चापि}}{\lemma{\skt{दाडिमश्}}\Afootnote{K; \skt{नागरं}\ A}}}{\lemma{\skt{चापि}}\Afootnote{K; \skt{चाम्ल}\ A}}\skt{ }\edtext{\skt{पेशिम{\symfont ¦}}}{\Afootnote{K; \skt{पेषितम् ॥}\ A}}\skt{ ।}\pend

\pstart \ledsidenote{\passref{6.40.137ab}}\edtext{\skt{वष्क}}{\Afootnote{K; \skt{दध्नः}\ A}}\skt{ सरश्च यमके भृष्टो वर्चः क्षये हितः ।}\pend

\pstart\versepar \ledsidenote{\passref{6.40.138}}\skt{वायुः प्रवृद्धो निचितं }\edtext{\skt{वलासं}}{\Afootnote{K; \skt{बलासं}\ A}}\skt{ नुद\-त्य\-ध\-स्ता\-द\-हि\-ता\-श\-नस्य । }\pend
\pstart\versepar \skt{प्रवाहमाणस्य }\edtext{\skt{मुहु\-र्मा\-ला\-क्त\-म्प्र\-वा\-हि\-का\-म्ता\-म्प्र\-व\-दन्ति}}{\lemma{\skt{मुहुर्मालाक्तम्}}\Afootnote{K; \skt{मुहुर्मलाक्तं}\ A}}\skt{ तज्ज्ञाः ।}\pend

\pstart \ledsidenote{\passref{6.40.139ab}}\skt{प्रवाहिका }\edtext{\edtext{\skt{वात\-कृ\-ता\-स्स\-शूलाः}}{\lemma{\skt{वातकृतास्}}\Afootnote{K; \skt{वातकृता}\ A}}}{\lemma{\skt{सशूलाः}}\Afootnote{K; \skt{सशूला}\ A}}\skt{ पित्तात्सदाहा सकफा कफाच्च ।}\pend

\pstart\versepar \ledsidenote{\passref{6.40.127}}\edtext{\skt{विल्वमध्यो}}{\Afootnote{K; \skt{बिल्वमध्यं}\ A}}\skt{ समधुकं }\edtext{\skt{शर्क\-रा\-चू\-र्ण्ण\-सं\-युत{\symfont ¦}}}{\Afootnote{K; \skt{शर्क\-रा\-क्षौ\-द्र\-सं\-यु\-तम् ।}\ A}}\skt{ । }\pend
\pstart\versepar \skt{तण्डुलाम्बुयुतो }\edtext{\skt{योगस्पित्तं}}{\Afootnote{K; \skt{योगः पित्त}\ A}}\skt{ }\edtext{\skt{रक्तोत्थितां}}{\Afootnote{K; \skt{रक्तोत्थितं}\ A}}\skt{ जयेत् ।}\pend

\pstart\versepar \ledsidenote{\passref{6.40.128}}\edtext{\skt{योगान्वरां}}{\Afootnote{K; \skt{योगान्सां}\ A}}\skt{ }\edtext{\edtext{\skt{ग्राहि\-क\-म्चा\-यि\-न्कु\-र्या\-त्क्षौ\-द्र\-श\-र्करां}}{\lemma{\skt{ग्राहिकम्}}\Afootnote{K; \skt{ग्राहिकांश्}\ A}}}{\lemma{\skt{चायि\-न्कु\-र्या\-त्क्षौ\-द्र\-श\-र्करां}}\Afootnote{K; \skt{चान्या\-न्पि\-बे\-त्स\-क्षौ\-द्र\-श\-र्क\-रान् ।}\ A}}\skt{ । }\pend
\pstart\versepar \skt{}\edtext{\skt{न्यग्राधादिषु}}{\Afootnote{K; \skt{न्यग्रोधादिषु}\ A}}\skt{ }\edtext{\skt{कुर्व्-इत}}{\Afootnote{K; \skt{कुर्याच्च}\ A}}\skt{ }\edtext{\skt{पुट\-पा\-का\-न्य\-ण\-रितां}}{\Afootnote{K; \skt{पुट\-पा\-का\-न्य\-थे\-रि\-तान् ॥}\ A}}\skt{ ॥}\pend

\pstart\versepar \ledsidenote{\passref{6.40.130.add}}\skt{दीप्ता\-शौ\-वा\-ष्प\-ती\-सार विवर्धति वलाग्निते । }\pend
\pstart\versepar \skt{त्रिफ\-ला\-वि\-ष\-ङ्ग\-कृ\-ष्णा\-क\-षा\-येण विरेचयेत ।}\pend

\pstart \ledsidenote{\passref{6.40.139cd}}\skt{सशोणिता }\edtext{\skt{शोणितसम्भवा}}{\Afootnote{K; \skt{शोणितसंभवा तु ताः}\ A}}\skt{ }\edtext{\skt{या}}{\Afootnote{K; \skt{स्नेह}\ A}}\skt{ रूक्षप्रभवा }\edtext{\skt{मता}}{\Afootnote{K; \skt{मतास्}\ A}}\skt{ तु ।}\pend

\pstart\versepar \ledsidenote{\passref{6.40.140}}\edtext{\edtext{\skt{तासा\-म\-ती\-सा\-र\-व\-दा\-दि\-शेच्च}}{\lemma{\skt{तासाम्}}\Afootnote{K; \skt{तासा}\ A}}}{\lemma{\skt{अती\-सा\-र\-व\-दा\-दि\-शेच्च}}\Afootnote{K; \skt{मती\-सा\-र\-व\-दा\-दि\-शेच्च}\ A}}\skt{ लिङ्गं }\edtext{\skt{क्रमे श्चा}}{\lemma{\skt{क्रमे श्}}\Afootnote{K; \skt{क्रमं}\ A}}\skt{ समुद्भवानां ॥ }\pend
\pstart\versepar \skt{न }\edtext{\skt{सिद्ध्मां}}{\Afootnote{K; \skt{शान्तिमा}\ A}}\skt{ }\edtext{\skt{यान्ति}}{\Afootnote{K; \skt{याति}\ A}}\skt{ }\edtext{\skt{विलंघ्न}}{\Afootnote{K; \skt{विलङ्घ}\ A}}\skt{ }\edtext{\skt{नाध्यौ\-र्यो\-गै\-रु\-दीर्णा}}{\lemma{\skt{नाध्यौर्}}\Afootnote{K; \skt{नैर्या}\ A}}\skt{ यदि }\edtext{\skt{पाचनैर्या}}{\lemma{\skt{या}}\Afootnote{K; \skt{वा ॥}\ A}}\skt{ ।}\edtext{}{\lemma{\skt{समुद्भवानां ॥ न}}\Afootnote{K; \skt{मविपक्वतां च । न}\ A}}\pend

\pstart \ledsidenote{\passref{6.40.141}}\edtext{\skt{ताः}}{\Afootnote{K; \skt{तां}\ A}}\skt{ }\edtext{\skt{क्षीरसमेवाशु}}{\Afootnote{K; \skt{क्षीरमेवाशु}\ A}}\skt{ शृतं निहन्ति तैलन्तिलाः }\edtext{\skt{पिच्छि\-ल\-व\-स्त\-स्त\-यम्च}}{\Afootnote{K; \skt{पिच्छि\-ल\-ब\-स्त\-यश्च । आर्द्रैः कुशैः संप\-रि\-वे\-ष्टि\-तानि वृन्तान्यथार्द्राणि हि शाल्मलीनाम् ॥}\ A}}\skt{ ।}\pend

\pstart\versepar \ledsidenote{\passref{6.40.143-144}}\edtext{\skt{क्षीरेण}}{\Afootnote{\skt{बस्तिं विद\-ध्या\-द्भि\-ष\-ग\-प्र\-मत्तः प्रवाहिकामूत्र}\ A}}\skt{ }\edtext{\edtext{\edtext{\skt{चास्था\-प\-न\-म\-ग्र्य\-मुक्त}}{\lemma{\skt{चास्थापनम्}}\Afootnote{\skt{पुरीषसङ्}\ A}}}{\lemma{\skt{अग्र्यम्}}\Afootnote{\skt{गे ।}\ A}}}{\lemma{\skt{उक्त}}\Afootnote{\emph{om.}\ A}}\skt{ }\edtext{\skt{द्विप\-ञ्च\-मू\-ली\-क्व\-थि\-तेन}}{\Afootnote{\skt{द्विप\-ञ्ज्च\-मू\-ली\-क्व\-थि\-तेन}\ A}}\skt{ शूले ॥ }\pend
\pstart\versepar \skt{}\edtext{\skt{प्रवाहिकायाश्च}}{\Afootnote{\skt{प्रवाहमाणस्य}\ A}}\skt{ }\edtext{\skt{समाक्षिकेन}}{\Afootnote{\skt{समाक्षिकेण ॥ क्षीरेण चास्थापनम}\ A}}\skt{ }\edtext{\skt{तैलेन}}{\Afootnote{\skt{ग्र्यमुक्तं तैलेन}\ A}}\skt{ }\edtext{\skt{युंज्या\-द\-नु\-वा\-स\-नञ्च}}{\Afootnote{\skt{युञ्ज्या\-द\-नु\-वा\-सनं च ।}\ A}}\skt{ । }\pend
\pstart\versepar \skt{}\edtext{\skt{वात\-घ्न\-व\-र्ग\-म\-धु\-रेषु}}{\Afootnote{\skt{वातघ्नवर्गे लवणेषु}\ A}}\skt{ चैव }\edtext{\skt{तैलश्च}}{\Afootnote{\skt{तैलं च}\ A}}\skt{ सिद्धं हितमन्नपाने ।}\pend

\pstart\versepar \ledsidenote{\passref{6.40.145}}\skt{लोध्रं विडं }\edtext{\skt{विल्वशलाटु}}{\Afootnote{K; \skt{बिल्वशलाटु}\ A}}\skt{ चैव }\edtext{\skt{लिह्यच्च}}{\Afootnote{K; \skt{लिह्याच्च}\ A}}\skt{ }\edtext{\skt{तैन}}{\Afootnote{K; \skt{तैलेन}\ A}}\skt{ कटुत्रिकाढ्यं । }\pend
\pstart\versepar \skt{दध्ना }\edtext{\skt{ससारेण}}{\Afootnote{K; \emph{om.}\ A}}\skt{ ससारेण समाक्षिकेण भुञ्जीत }\edtext{\skt{निश्चा\-र\-क\-पी\-डि\-तस्तु}}{\lemma{\skt{चार\-क\-पी\-डि\-तस्}}\Afootnote{K; \skt{चार\-क(\-निः\-सा\-र\-क))\-पी\-डि\-तस्}\ A}}\skt{ ।}\pend

\pstart\versepar \ledsidenote{\passref{6.40.146}}\edtext{\skt{सुत\-प्त\-कु\-प्य\-न्क्व\-थि\-तेन}}{\Afootnote{K; \skt{सुत\-प्त\-कु\-प्य\-क्व\-थि\-तेन}\ A}}\skt{ वापि क्षीरेण शीतेन }\edtext{\skt{मधुप्लुतेनं}}{\Afootnote{K; \skt{मधुप्लुतेन ।}\ A}}\skt{ । }\pend
\pstart\versepar \skt{}\edtext{\skt{क्षीराप्लुतौ}}{\Afootnote{K; \skt{शूलार्दितो}\ A}}\skt{ व्योष\-वि\-दा\-रि\-ग\-न्धा\-सि\-द्धेन दुग्धेन हिताय भोज्यः ।}\pend

\pstart\versepar \ledsidenote{\passref{6.40.147}}\skt{वात\-घ्न\-स\-ङ्ग्रा\-हि\-क\-दी\-प\-नीयैः }\edtext{\skt{कृता\-न्ष\-डा\-म्चा\-प्यु\-प\-यो\-ज\-यीत्}}{\lemma{\skt{कृतान्षडाम्}}\Afootnote{K; \skt{कृतान्षडांश्}\ A}}\skt{ । }\pend
\pstart\versepar \skt{खादेच्च }\edtext{\skt{मत्यांस}}{\Afootnote{K; \skt{मत्स्यान्र}\ A}}\skt{ }\edtext{\edtext{\skt{रसस्तु}}{\lemma{\skt{रसस्}}\Afootnote{K; \skt{समा}\ A}}}{\lemma{\skt{तु}}\Afootnote{K; \skt{प्नु}\ A}}\skt{ याच्च }\edtext{\skt{वातघ्नसिद्धां}}{\Afootnote{K; \skt{वातघ्नसिद्धं}\ A}}\skt{ सघृतं । }\edtext{\skt{सतैलां}}{\Afootnote{K; \skt{सतैलम् ॥}\ A}}\skt{ ।}\edtext{}{\lemma{\skt{उपयोजयीत् । खादेच्च}}\Afootnote{K; \skt{उपभोजयेच्च । खादेच्च}\ A}}\pend

\pstart\versepar \ledsidenote{\passref{6.40.148}}\edtext{\skt{एणाव्यजनाश्च}}{\Afootnote{K; \skt{एणाव्यजानां तु}\ A}}\skt{ वटप्रवालैः सिद्धानि सार्धम्पिशितानि }\edtext{\skt{खादेत}}{\Afootnote{K; \skt{खादेत् । मेध्यस्य}\ A}}\skt{ । }\pend
\pstart\versepar \skt{}\edtext{\skt{मेध्यं}}{\Afootnote{K; \skt{(मेद्य}\ A}}\skt{ }\edtext{\skt{सुसिद्धं}}{\Afootnote{K; \skt{स्य) सिद्धं}\ A}}\skt{ त्वथ वापि रक्तं }\edtext{\edtext{\skt{वस्तस्य}}{\lemma{\skt{वस्}}\Afootnote{K; \skt{ब}\ A}}}{\lemma{\skt{तस्य}}\Afootnote{K; \skt{स्तस्य}\ A}}\skt{ दध्ना }\edtext{\skt{घृततैलयुक्तः}}{\Afootnote{K; \skt{घृत\-तै\-ल\-यु\-क्तम् ॥}\ A}}\skt{ ।}\pend

\pstart\versepar \ledsidenote{\passref{6.40.149}}\edtext{\skt{खादेत}}{\Afootnote{K; \skt{खादे}\ A}}\skt{ }\edtext{\skt{प्रदेहैः}}{\Afootnote{K; \skt{त्प्रदेहैः}\ A}}\skt{ }\edtext{\skt{शिखिलावर्जर्वा}}{\Afootnote{K; \skt{शिखिलावजैर्वा}\ A}}\skt{ }\edtext{\skt{यूषैश्च}}{\Afootnote{K; \skt{भुञ्जीत यूषै}\ A}}\skt{ }\edtext{\edtext{\skt{मुख्यैर्दधिभिश्च}}{\lemma{\skt{मुख्यैर्}}\Afootnote{K; \emph{om.}\ A}}}{\lemma{\skt{दधिभिश्च}}\Afootnote{K; \skt{र्दधिभिश्च}\ A}}\skt{ }\edtext{\skt{भुंजेत}}{\Afootnote{K; \skt{मुख्यैः । मा}\ A}}\skt{ । }\pend
\pstart\versepar \skt{}\edtext{\skt{माषां}}{\Afootnote{K; \skt{षान्सु}\ A}}\skt{ }\edtext{\skt{सुद्धिं}}{\Afootnote{K; \skt{सिद्धा}\ A}}\skt{ }\edtext{\skt{घृतमण्डयुक्तां}}{\Afootnote{K; \skt{न्घृतमण्डयुक्ता}\ A}}\skt{ }\edtext{\skt{खादेच्च}}{\Afootnote{K; \skt{न्खादेच्च}\ A}}\skt{ दध्ना }\edtext{\skt{मरिचोपदंशां}}{\Afootnote{K; \skt{मरि\-चो\-प\-दं\-शान् ॥}\ A}}\skt{ ।}\pend

\pstart\versepar \ledsidenote{\passref{6.40.150.add1}}\skt{विडं दुडं विल्वसलाटु चैव लिह्याच्च तैल न कटुत्रिकाद्य । }\pend
\pstart\versepar \skt{राजसेनात्क दाडिमानां रसेषु सिद्धांरूघृतास्स तैलाः ।}\pend

\pstart\versepar \ledsidenote{\passref{6.40.151.add2}}\skt{प्रवाहमानाय हिता यवा\-गू\-रा\-प्नु\-यानं पयसा हितम्वा । }\pend
\pstart\versepar \skt{रुफा\-ति\-मू\-त्र\-कृ\-च्छ\-क्षी\-रेण सक्षौ\-द्र\-घृ\-तो\-त्प\-लन् ।}\pend

\pstart \ledsidenote{\passref{6.40.152.add3}}\skt{स्यष्टि\-नो\-च्छा\-प्य\-त\-तो\-नु\-वा\-स्य\-तै\-लेन काको\-लि\-ग\-ण\-श्रि\-तेन ।}\pend

\pstart \ledsidenote{\passref{6.40.152ab}}\skt{रात्रावहनि वा नित्यं रुजार्तो यो }\edtext{\skt{भवेन्नर{\symfont ¦}}}{\lemma{\skt{नर{\symfont ¦}}}\Afootnote{K; \skt{नरः ।}\ A}}\skt{ ।}\pend

\pstart \ledsidenote{\passref{6.40.152.add4}}\skt{तैलं सावा\-हा\-र्ये\-षु\-मे\-के\-ङ्ग\-नु\-वा\-सन ।}\pend

\pstart \ledsidenote{\passref{6.40.152cd}}\skt{यथा यथा }\edtext{\skt{तैलं}}{\Afootnote{K; \skt{सतैलः}\ A}}\skt{ }\edtext{\skt{स्यात्तत्रहा}}{\Afootnote{K; \skt{स्या}\ A}}\skt{ }\edtext{\skt{शाम्यति}}{\Afootnote{K; \skt{द्वात}\ A}}\skt{ }\edtext{\skt{मारुतः}}{\Afootnote{K; \skt{शान्तिस्तथा तथा ॥}\ A}}\skt{ ।}\pend

\pstart \ledsidenote{\passref{6.40.153ab}}\skt{प्रशान्ते मारुते }\edtext{\skt{चास्य}}{\Afootnote{K; \skt{चापि}\ A}}\skt{ }\edtext{\skt{शान्ति}}{\Afootnote{K\textsuperscript{pc}; \skt{शान्तिं}\ A; \skt{शा}\ K\textsuperscript{ac}}}\skt{ याति प्रवाहिका ।}\pend

\pstart\versepar \ledsidenote{\passref{6.40.154}}\edtext{\skt{पाठा\-ज\-मो\-दौ\-कु\-ट\-जस्य}}{\Afootnote{\skt{पाठा\-ज\-मो\-दा\-कु\-ट\-जोत्प}\ A; \skt{पाठा\-ज\-मो\-दौ\-कु\-त्̤\-अ\-जस्य}\ K}}\skt{ }\edtext{\skt{वीजं}}{\Afootnote{K; \skt{लं च}\ A}}\skt{ शुण्ठी समा मागधिकाश्च पिष्टाः । }\pend
\pstart\versepar \skt{}\edtext{\skt{सुखाम्वुपीताः}}{\Afootnote{K; \skt{सुखाम्बुपीताः}\ A}}\skt{ शमयन्ति रोगं }\edtext{\edtext{\skt{मेषा\-ण्ड\-सि\-द्ध\-म्प\-यसो}}{\lemma{\skt{मेषाण्डसिद्धम्}}\Afootnote{K; \skt{मेध्याण्डसिद्धं}\ A}}}{\lemma{\skt{पयसो}}\Afootnote{K; \skt{सघृतं}\ A}}\skt{ }\edtext{\skt{हितम्वा}}{\lemma{\skt{हितम्}}\Afootnote{K; \skt{पयो}\ A}}\skt{ ।}\pend

\pstart \ledsidenote{\passref{6.40.155.add}}\skt{पित्तान्शिकायां मधुकं सिताश्च मयिश्च व करसंयुतश्च ।}\pend

\pstart \ledsidenote{\passref{6.40.155}}\edtext{\skt{लीढ्वाथ}}{\Afootnote{K; \skt{शुण्ठीं घृतं सक्षवकं सतैलं विपाच्य लीढ्वामयमाशु हन्यात् । गजा\-श\-ना\-कु\-म्भि\-क\-दा\-डि\-मानां रसैः कृता तैल}\ A}}\skt{ }\edtext{\edtext{\skt{सक्षौ\-द्र\-स\-म\-म्पि\-बेद्वा}}{\lemma{\skt{सक्षौद्रसमम्}}\Afootnote{K; \skt{घृते सदध्नि ॥}\ A}}}{\lemma{\skt{पिबेद्वा}}\Afootnote{\emph{om.}\ A; \skt{पिवेद्वा}\ K}}\skt{ }\edtext{\skt{जम्वु\-रा\-दा\-म\-ल\-काच्च}}{\Afootnote{K; \emph{om.}\ A}}\skt{ }\edtext{\skt{कल्क{\symfont ¦}}}{\Afootnote{K; \emph{om.}\ A}}\skt{ ।}\pend

\pstart \ledsidenote{\passref{6.40.156.add}}\skt{सिता\-द\-वि\-क्षौ\-द्र\-यु\-त\-न्त\-थौव विल्वन्न तदाक दाडिमश्च ।}\pend

\pstart \ledsidenote{\passref{6.40.157.add}}\skt{निलां सयष्टीमधुकं सीताश्च क्षौद्रा तण्डु\-ल\-पा\-च\-नेन ।}\pend

\pstart \ledsidenote{\passref{6.40.158.add}}\skt{सर्पिर्वि वंधे भृशे दाह मूविपाच्यदुग्धं न लवते साभ्यां ।}\pend

\pstart \ledsidenote{\passref{6.40.159.add}}\skt{मूलैः पयसोत्पल यष्टि\-का\-मी\-वि\-दा\-र्य\-का\-श्र्मर्य वलासमतैः ।}\pend

\pstart \ledsidenote{\passref{6.40.160.add}}\skt{देया }\edtext{\skt{घृत\-क्षौ\-द्र\-यु\-तो\-नु\-वा\-स\-घृ\-तेन}}{\Afootnote{K\textsuperscript{pc}; \skt{घृत\-क्षौ\-यु\-तो\-नु\-वा\-स\-घृ\-तेन}\ K\textsuperscript{ac}}}\skt{ चान्ते मधुकश्रितेन ।}\pend

\pstart \ledsidenote{\passref{6.40.161.add}}\skt{पित्तान्मिकारसः सपयाथ विमिश्रितस्तु ।}\pend

\pstart \ledsidenote{\passref{6.40.162.add}}\skt{पद्मां समद्म्वुज धातकीभिः कल्की कृता\-भि\-र्म\-धु\-सं\-यु\-ताभिः ।}\pend

\pstart \ledsidenote{\passref{6.40.163.add}}\skt{शाम्यन्तथैवं न तस्य समङ्गा पद्मा पयः पादपयोत कीर्तिः ।}\pend

\pstart \ledsidenote{\passref{6.40.164.add}}\skt{पक्वा\-प\-यो\-वि\-ल्व\-क\-पी\-त\-ना\-म्भमो या सर्पिमयुनारु ।}\pend

\pstart \ledsidenote{\passref{6.40.165.add}}\skt{एतत्च्छिते नैव घृतेन सम्यशये चैवमुपैति शान्हि{\symfont ¦} ।}\pend

\pstart \ledsidenote{\passref{6.40.156cd}}\skt{लघूनि पथ्यान्यथ दीपनानि स्निग्धानि }\edtext{\skt{भोज्या\-न्यु\-द\-रा\-म\-यु\-एषु}}{\Afootnote{K; \skt{भोज्या\-न्यु\-द\-रा\-म\-येषु ॥}\ A}}\skt{ ॥}\pend

\pstart \ledsidenote{\passref{6.40.157}}\skt{हिताय नित्यं }\edtext{\skt{वितरेद्विभोक्षं}}{\Afootnote{K; \skt{वितरेद्विभज्य}\ A}}\skt{ }\edtext{\skt{योगांस्च}}{\Afootnote{K; \skt{योगांश्च}\ A}}\skt{ तां }\edtext{\skt{स्तां}}{\Afootnote{K; \skt{स्ता}\ A}}\skt{ }\edtext{\skt{भिषप्रमत्तः}}{\Afootnote{K; \skt{न्भिषगप्रमत्तः ॥}\ A}}\skt{ ।}\pend

\pstart\versepar \ledsidenote{\passref{6.40.158}}\skt{तृष्णापनयनी }\edtext{\skt{लध्वी}}{\Afootnote{K; \skt{लघ्वी}\ A}}\skt{ दीपनी }\edtext{\skt{वस्तिशोधनी}}{\Afootnote{K; \skt{बस्तिशोधनी ।}\ A}}\skt{ । }\pend
\pstart\versepar \skt{ज्वरे चैवातिसारे च }\edtext{\skt{यच्चाद्विस्रत्रं}}{\Afootnote{K; \skt{य}\ A}}\skt{ }\edtext{\skt{सम्यग्वायुम्च}}{\Afootnote{K; \emph{om.}\ A}}\skt{ }\edtext{\skt{च्छति}}{\Afootnote{K; \skt{वागूः सर्वदा हिता ॥}\ A}}\skt{ ।}\pend

\pstart \ledsidenote{\passref{6.40.169.add}}\skt{वागूसर्वदा हिता ।}\pend

\pstart \ledsidenote{\passref{6.40.158cd}}\skt{रौ क्षान्प्र\-वृ\-त्त\-ती\-सा\-र\-क्रि\-या\-स्निग्धा विधीयते ।}\pend

\pstart \ledsidenote{\passref{6.40.159cd}}\edtext{\skt{रौक्षा\-न्ग्र\-ह\-नि\-मित्ते}}{\Afootnote{K; \skt{भयजे}\ A}}\skt{ }\edtext{\skt{च शोकजे}}{\Afootnote{K; \skt{सान्त्वनापूर्वा शोकजे}\ A}}\skt{ शोक }\edtext{\skt{नाशनी}}{\Afootnote{K; \skt{नाशिनी ॥}\ A}}\skt{ ।}\pend

\pstart \ledsidenote{\passref{6.40.160ab}}\skt{विषार्शः }\edtext{\skt{क्रिमिसंभूते}}{\Afootnote{K; \skt{कृमिसंभूते}\ A}}\skt{ हिता चोभयशर्मदा ।}\pend

\pstart\versepar \ledsidenote{\passref{6.40.161}}\skt{समवाये तु दोषाणां }\edtext{\skt{पूर्व्व\-म्पि\-त्त\-मु\-पा\-च\-रेत्}}{\lemma{\skt{पित्तम्}}\Afootnote{K; \skt{पित्त}\ A}}\skt{ । }\pend
\pstart\versepar \skt{ज्वरे चैवातिसारे च सर्व्वत्रान्यत्र }\edtext{\skt{मारुत{\symfont ¦}}}{\Afootnote{K; \skt{मारुतम् ॥}\ A}}\skt{ ।}\edtext{}{\lemma{\skt{उपाचरेत् । ज्वरे}}\Afootnote{K; \skt{मुपाचरेत् । ज्वरे}\ A}}\pend

\pstart\versepar \ledsidenote{\passref{6.40.172.add}}\skt{यस्योच्चाराद्विना मूत्रं सम्यग्वायुश्च गच्छति । }\pend
\pstart\versepar \skt{दीप्ता\-ग्ने\-र्ल\-घु\-को\-ष्ठ\-स्य\-स्थि\-त\-स्त\-स्यो\-दरो मयः ।}\pend

\pstart\versepar \ledsidenote{\passref{6.40.167}}\skt{अतिसारे निवृत्ते }\edtext{\skt{पि}}{\Afootnote{K; \skt{ऽपि}\ A}}\skt{ }\edtext{\edtext{\skt{मन्दा\-ग्ने\-र\-हि\-ता\-शिनः}}{\lemma{\skt{मन्दाग्नेर्}}\Afootnote{K; \skt{मन्दाग्ने}\ A}}}{\lemma{\skt{अहिताशिनः}}\Afootnote{K; \skt{रहिताशिनः ।}\ A}}\skt{ । }\pend
\pstart\versepar \skt{भूयः संदूषितो }\edtext{\edtext{\skt{वन्हि\-र्ग्र\-ह\-णी\-मयि}}{\lemma{\skt{वन्हिर्}}\Afootnote{K; \skt{वह्नि}\ A}}}{\lemma{\skt{ग्रहणीमयि}}\Afootnote{K; \skt{र्ग्रहणीमभि}\ A}}\skt{ दूषयेत् ।}\pend

\pstart \ledsidenote{\passref{6.40.168}}\skt{तस्मात्कार्यः }\edtext{\skt{परी\-हा\-र\-स्मा\-ती\-सारे}}{\Afootnote{K; \skt{परी\-हा\-र\-स्त्व\-ती\-सारे}\ A}}\skt{ विरिक्तवत् । यावन्न }\edtext{\skt{प्रकृ\-ति\-स्थ\-स्स्यो\-द्दो\-षतः}}{\Afootnote{K; \skt{प्रकृतिस्थः स्याद्दोषतः}\ A}}\skt{ प्राणतस्तथा ॥}\pend

\pstart\versepar \ledsidenote{\passref{6.40.169}}\skt{षष्ठी }\edtext{\skt{कला}}{\Afootnote{K; \emph{om.}\ A}}\skt{ पित्तधरा }\edtext{\skt{या}}{\Afootnote{K; \skt{नाम या}\ A}}\skt{ }\edtext{\skt{पूर्व\-म्प\-रि\-की\-र्त्तिता}}{\lemma{\skt{पूर्वम्}}\Afootnote{K; \skt{कला}\ A}}\skt{ । }\pend
\pstart\versepar \skt{}\edtext{\skt{पक्क्वा\-मा\-श\-य\-म\-ध्यस्था}}{\Afootnote{K; \skt{पक्वा\-मा\-श\-य\-म\-ध्यस्था}\ A}}\skt{ ग्रहणी सा प्रकीर्त्तिता ।}\pend

\pstart\versepar \ledsidenote{\passref{6.40.170}}\skt{ग्रहण्या }\edtext{\edtext{\skt{वलमग्निर्हि}}{\lemma{\skt{वलम्}}\Afootnote{K; \skt{बल}\ A}}\skt{ स}}{\lemma{\skt{अग्निर्हि स}}\Afootnote{K; \skt{मग्निर्हि स}\ A}}\skt{ चापि }\edtext{\skt{ग्रहणी}}{\Afootnote{K; \skt{ग्रहणीं}\ A}}\skt{ श्रितः । }\pend
\pstart\versepar \skt{}\edtext{\skt{तस्मात्सदूषिते}}{\Afootnote{K; \skt{तस्मात्संदूषिते}\ A}}\skt{ }\edtext{\skt{वन्हौ}}{\Afootnote{K; \skt{वह्नौ}\ A}}\skt{ ग्रहणी सम्प्रदुष्यति ।}\pend

\pstart\versepar \ledsidenote{\passref{6.40.171}}\edtext{\edtext{\skt{एकै\-क\-श\-स्स\-र्वम्च}}{\lemma{\skt{एकैकशस्}}\Afootnote{K; \skt{एकशः}\ A}}}{\lemma{\skt{सर्वम्च}}\Afootnote{K; \skt{सर्वशश्}\ A}}\skt{ }\edtext{\edtext{\skt{चैथै\-र\-त्य\-र्थ\-मु\-च्छ्रितैः}}{\lemma{\skt{चैथैर्}}\Afootnote{K; \skt{चैव दोषै}\ A}}}{\lemma{\skt{अत्यर्थम्}}\Afootnote{K; \skt{रत्यर्थ}\ A}}\skt{ । }\pend
\pstart\versepar \skt{सा दुष्टा }\edtext{\skt{वहुशो}}{\Afootnote{K; \skt{बहुशो}\ A}}\skt{ }\edtext{\skt{भुक्तामाममेव}}{\lemma{\skt{भुक्तामामम्}}\Afootnote{K; \skt{भुक्तमामम्}\ A}}\skt{ विमुञ्चति ।}\edtext{}{\lemma{\skt{उच्छ्रितैः । सा}}\Afootnote{K; \skt{मुच्छ्रितैः । सा}\ A}}\pend

\pstart\versepar \ledsidenote{\passref{6.40.172}}\skt{पक्वम्वा }\edtext{\skt{सरू\-ज\-पू\-ति\-र्मु\-र्हु\-र्वद्धं}}{\Afootnote{K; \skt{सरुजं पूति मुहुर्बद्धं}\ A}}\skt{ }\edtext{\skt{मुर्हुर्द्रव{\symfont ¦}}}{\Afootnote{K; \skt{मुहुर्द्रवम् ।}\ A}}\skt{ । }\pend
\pstart\versepar \skt{}\edtext{\skt{ग्रह\-णी\-रो\-ग\-मा\-हुस्तं}}{\Afootnote{K; \skt{ग्रह\-णी\-रो\-ग\-मा\-हुस्त}\ A}}\skt{ }\edtext{\skt{मायु\-र्वे\-द\-वि\-शा\-रदाः}}{\Afootnote{K; \skt{मायुर्वेदविदो जनाः ॥}\ A}}\skt{ ।}\pend

\pstart\versepar \ledsidenote{\passref{6.40.173}}\skt{तस्योत्पत्तौ }\edtext{\skt{सदाहोन्न}}{\Afootnote{K; \skt{विदाहो+आने}\ A}}\skt{ }\edtext{\skt{सदानालस्य}}{\Afootnote{K; \skt{सदनालस्य}\ A}}\skt{ तृट्क्लमाः । }\pend
\pstart\versepar \skt{}\edtext{\edtext{\skt{वलक्षयोरुचिः}}{\lemma{\skt{वलक्षयोर्}}\Afootnote{K; \skt{बलक्षयो}\ A}}}{\lemma{\skt{उचिः}}\Afootnote{K; \skt{ऽरुचिः}\ A}}\skt{ कासः }\edtext{\skt{कर्ण्नक्ष्वेडो}}{\Afootnote{K; \skt{कर्णक्ष्वेडो}\ A}}\skt{ }\edtext{\skt{तु}}{\Afootnote{K; \skt{ऽन्त्र}\ A}}\skt{ कूजनं ।}\pend

\pstart\versepar \ledsidenote{\passref{6.40.174}}\edtext{\skt{अथा}}{\Afootnote{K; \skt{अथ}\ A}}\skt{ जाते }\edtext{\skt{भवेज्जन्तोः ः}}{\Afootnote{K; \skt{भवेज्जन्तुः}\ A}}\skt{ । }\pend
\pstart\versepar \skt{}\edtext{\skt{निछ\-र्दि\-ज्व\-रा\-रो\-च\-क\-दा\-ह\-वान्}}{\Afootnote{K; \skt{शूनपादकरः कृशः । पर्व\-रु\-ग्लौ\-ल्य\-तृ\-ट्छ\-र्दि\-ज्व\-रा\-रो\-च\-क\-दा\-ह\-वाम् ॥}\ A}}\skt{ ।}\pend

\pstart\versepar \ledsidenote{\passref{6.40.175}}\edtext{\skt{उद्गि\-रे\-च्छु\-क्त\-ति\-क्ता\-म्ल\-तस्य}}{\Afootnote{K; \skt{उद्गि\-रे\-च्छु\-क्त\-ति\-क्ता\-म्ल\-लोह}\ A}}\skt{ }\edtext{\skt{धूमामगन्धिमां}}{\Afootnote{K; \skt{धूमा\-म\-ग\-न्धि\-कम् ।}\ A}}\skt{ । }\pend
\pstart\versepar \skt{}\edtext{\skt{प्रसे\-क\-मु\-ख\-वै\-र\-स्य\-त\-म\-का\-रु\-चि\-पी\-डिताः}}{\Afootnote{K; \skt{प्रसे\-क\-मु\-ख\-वै\-र\-स्य\-त\-म\-का\-रु\-चि\-पी\-डितः ॥}\ A}}\skt{ ।}\pend

\pstart\versepar \ledsidenote{\passref{6.40.176}}\skt{वाताच्छूलाधिकैः }\edtext{\skt{पायु\-हृ\-त्पा\-र्शो\-द\-र\-म\-स्तकैः}}{\Afootnote{K\textsuperscript{pc}; \skt{पायु\-हृ\-त्पा\-र्श्वो\-द\-र\-म\-स्तकैः ।}\ A; \skt{पायु\-हृ\-त्पा\-र्शो\-द\-म\-स्तकैः ।}\ K\textsuperscript{ac}}}\skt{ । }\pend
\pstart\versepar \skt{}\edtext{\edtext{\skt{पित्ता\-त्स\-दा\-है\-र्गु\-रुभिः}}{\lemma{\skt{पित्तात्सदाहैर्}}\Afootnote{K; \skt{पित्तात्सदाहै}\ A}}}{\lemma{\skt{गुरुभिः}}\Afootnote{K; \skt{र्गुरुभिः}\ A}}\skt{ }\edtext{\skt{कफा }\edtext{\skt{त्रिभ्यस्त्रिलक्षणैः}}{\lemma{\skt{त्रिलक्षणैः}}\Afootnote{A; \skt{त्रिलक्षणौः ।}\ K}}}{\lemma{\skt{कफा त्रिभ्यस्}}\Afootnote{K; \skt{कफात्त्रिभ्यस्}\ A}}\skt{ ।}\pend

\endnumbering
\beginnumbering
\eledsection*{Adhyāya 41}
\pstart \ledsidenote{\passref{6.41.1}}\skt{अथातः }\edtext{\skt{शोषप्रतिषेधं}}{\Afootnote{A K; \skt{शोष\-चि\-कि\-त्सि\-तम्}\ H}}\skt{ }\edtext{\skt{व्याख्यास्यामः}}{\Afootnote{A H; \skt{व्या ॥}\ K}}\skt{ ॥}\pend

\pstart\versepar \ledsidenote{\passref{6.41.3}}\edtext{\skt{अने\-क\-दो\-षा\-नु\-गतो}}{\Afootnote{K H; \skt{अने\-क\-रो\-गा\-नु\-गतो}\ A}}\skt{ }\edtext{\skt{वहुरोगागमः}}{\Afootnote{K H; \skt{बहु\-रो\-ग\-पु\-रो\-गमः ।}\ A}}\skt{ । }\pend
\pstart\versepar \skt{दुर्विज्ञेयो }\edtext{\skt{दुर्निवारः}}{\Afootnote{A K; \skt{दुर्न्निवार्यः}\ H}}\skt{ शोषो }\edtext{\edtext{\skt{व्याधिर्महावलः}}{\lemma{\skt{व्याधिर्}}\Afootnote{K H; \skt{व्याधि}\ A}}}{\lemma{\skt{महावलः}}\Afootnote{K H; \skt{र्महाबलः ॥}\ A}}\skt{ ॥}\pend

\pstart\versepar \ledsidenote{\passref{6.41.4}}\skt{संशो\-ष\-णा\-द्र\-सा\-दीनां शोष इत्यभिधीयते । }\pend
\pstart\versepar \skt{क्रिया\-क्ष\-य\-क\-र\-त्वाच्च क्षय इत्युच्यते पुनः ।}\pend

\pstart\versepar \ledsidenote{\passref{6.41.5}}\skt{राज्ञस्चंद्रमसो }\edtext{\skt{यस्मादभूदेष}}{\Afootnote{K; \skt{यस्मादभूदेष}\ A; \skt{यस्मादभूदेष}\ H}}\skt{ किलामयः । }\pend
\pstart\versepar \skt{}\edtext{\skt{तस्मात्त}}{\Afootnote{K; \skt{तस्मात्तं}\ A H}}\skt{ राजयक्ष्मेति }\edtext{\skt{केचिदाहुः}}{\Afootnote{K H; \skt{केचिदाहुः}\ A}}\skt{ }\edtext{\edtext{\skt{पुनर्जनाः}}{\lemma{\skt{पुनर्}}\Afootnote{K; \skt{पुन}\ A H}}}{\lemma{\skt{जनाः}}\Afootnote{K; \skt{र्जनाः ॥}\ A; \skt{र्ज्जनाः ।}\ H}}\skt{ ।}\pend

\pstart\versepar \ledsidenote{\passref{6.41.6}}\skt{स }\edtext{\skt{व्यस्तैर्जयते}}{\lemma{\skt{जयते}}\Afootnote{K; \skt{जायते}\ A H}}\skt{ दोषैरिति केचिद्वदन्ति हि ॥ }\pend
\pstart\versepar \skt{}\edtext{\skt{एका\-द\-शा\-ना\-मे\-कस्मिं}}{\lemma{\skt{एकस्मिं}}\Afootnote{K; \skt{एकस्मि}\ A H}}\skt{ }\edtext{\skt{सांनि\-ध्या\-त्त\-न्त्र\-युक्तिः}}{\Afootnote{K; \skt{न्सान्नि\-ध्या\-त्त\-न्त्र\-यु\-क्तितः}\ A H}}\skt{ ।}\pend

\pstart\versepar \ledsidenote{\passref{6.41.7}}\edtext{\edtext{\skt{क्रिया\-णा\-श्चा\-वि\-भा\-गेन}}{\lemma{\skt{क्रियाणाश्}}\Afootnote{K; \skt{क्रियाणाम्}\ A; \skt{क्रियानां}\ H}}}{\lemma{\skt{चाविभागेन}}\Afootnote{K H; \skt{अविभागेन}\ A}}\skt{ }\edtext{\skt{प्रागे\-वो\-त्पा\-द\-नेन च}}{\Afootnote{K; \skt{प्रागे\-को\-त्पा\-द\-नेन च}\ A H}}\skt{ । }\pend
\pstart\versepar \skt{एक एव मतः }\edtext{\skt{शोष}}{\Afootnote{K H; \skt{शोषः}\ A}}\skt{ }\edtext{\skt{सन्निपातात्मको}}{\Afootnote{K H; \skt{सन्निपातात्मको ह्य्}\ A}}\skt{ }\edtext{\skt{ह्यतः}}{\Afootnote{K H; \skt{अतः ॥}\ A}}\skt{ ।}\pend

\pstart\versepar \ledsidenote{\passref{6.41.8}}\skt{उद्रे }\edtext{\skt{कान्तस्य}}{\Afootnote{H; \skt{कात्तत्र}\ A; \skt{कान्त…स्य}\ K}}\skt{ लिङ्गानि दोषाणां }\edtext{\skt{तु}}{\Afootnote{K H; \skt{नि}\ A}}\skt{ }\edtext{\skt{भवन्ति ह}}{\Afootnote{K H; \skt{पतन्ति हि ।}\ A}}\skt{ । }\pend
\pstart\versepar \skt{}\edtext{\skt{क्षया\-द्वे\-ग\-प्र\-ती\-घा\-ता\-दा\-घा\-ता\-द्वि\-ष\-मा\-श\-नात्}}{\Afootnote{K; \skt{क्सया\-द्वे\-ग\-प्र\-ती\-घा\-ता\-दा\-घा\-ता\-द्वि\-ष\-मा\-श\-नात् ॥}\ A; \skt{क्षया\-द्वे\-ग\-प्र\-ती\-घा\-ता\-दा\-घा\-ता\-द्वि\-ष\-मा\-श\-नात् ।}\ H}}\skt{ ।}\pend

\pstart\versepar \ledsidenote{\passref{6.41.9}}\edtext{\skt{जयते}}{\Afootnote{K; \skt{जायते}\ A; \skt{जायन्ते}\ H}}\skt{ }\edtext{\edtext{\edtext{\skt{कुपि\-तै\-र्दो\-षै\-र्व्या\-प्त\-दे\-हस्य}}{\lemma{\skt{कुपितैर्}}\Afootnote{K H; \skt{कुपितै}\ A}}}{\lemma{\skt{दोषैर्}}\Afootnote{K H; \skt{र्दोषै}\ A}}}{\lemma{\skt{व्याप्तदेहस्य}}\Afootnote{K H; \skt{र्व्याप्तदेहस्य}\ A}}\skt{ देहिनः । }\pend
\pstart\versepar \skt{}\edtext{\edtext{\skt{कफ\-प्र\-धा\-नै\-र्दो\-षैस्तु}}{\lemma{\skt{कफप्रधानैर्}}\Afootnote{K; \skt{कफप्रधानै}\ A; \skt{कफनैर्}\ H}}}{\lemma{\skt{दोषैस्तु}}\Afootnote{K H; \skt{र्दोषैर्हि}\ A}}\skt{ }\edtext{\skt{रुद्धेषु}}{\Afootnote{A K; \skt{रूद्धे}\ H}}\skt{ रसवर्त्मसु ।}\pend

\pstart\versepar \ledsidenote{\passref{6.41.10}}\skt{अतिव्यवायिनो वापि }\edtext{\skt{क्षीण\-रे\-त\-स्य\-ना\-न्तरं}}{\Afootnote{K; \skt{क्षीणे रेत\-स्य\-न\-न्त\-रम् ।}\ A; \skt{क्षीण\-रे\-त\-स्य\-न\-न्तरं ॥}\ H}}\skt{ ॥ }\pend
\pstart\versepar \skt{क्षीयन्ते धातवस्सर्वे ततः }\edtext{\skt{शुष्यति}}{\Afootnote{K H; \skt{शुष्यन्ति}\ A}}\skt{ मानवः ।}\pend

\pstart\versepar \ledsidenote{\passref{6.41.12}}\edtext{\skt{स्वर\-भे\-दो\-नि\-ला\-च्छू\-ल\-सं\-काच}}{\Afootnote{K; \skt{स्वर\-भे\-दोऽ\-नि\-ला\-ञ्छूलं संलोचश्}\ A; \skt{स्वर\-भे\-दोऽ\-नि\-ला\-च्छूलं संकाचं}\ H}}\skt{ }\edtext{\skt{स्वांसपार्श्व याः}}{\Afootnote{K; \skt{चांसपार्श्वयोः ।}\ A; \skt{श्वांसपार्श्वः ॥}\ H}}\skt{ । }\pend
\pstart\versepar \skt{}\edtext{\skt{ज्वरां}}{\Afootnote{K; \skt{ज्वरो}\ A H}}\skt{ }\edtext{\skt{दाहोतिसारश्च}}{\Afootnote{K H; \skt{दाहोऽतिसारश्च}\ A}}\skt{ }\edtext{\skt{पित्ताद्रक्तस्य}}{\Afootnote{A K; \skt{पित्ताद्रक्तस्य}\ H}}\skt{ }\edtext{\skt{चागामः}}{\Afootnote{K; \skt{चागमः ॥}\ A H}}\skt{ ।}\pend

\pstart\versepar \ledsidenote{\passref{6.41.13}}\edtext{\skt{शिरसः}}{\Afootnote{A K; \skt{शिरसः परिसः}\ H}}\skt{ }\edtext{\skt{परि\-पू\-र्ण\-त्वं\-म\-भ\-क्त\-च्छन्द}}{\Afootnote{K; \skt{परि\-पू\-र्ण\-त्व\-म\-भ\-क्त\-च्छन्द}\ A; \skt{परि\-पू\-र्ण\-त्व\-म\-भ\-क्त\-च्छन्द}\ H}}\skt{ एव च । }\pend
\pstart\versepar \skt{कासः कण्ठस्य }\edtext{\skt{चोध्वंसो}}{\Afootnote{A K; \skt{चोद्वंसो}\ H}}\skt{ विज्ञेयः }\edtext{\skt{कफकोपतः}}{\Afootnote{A K; \skt{ककफकोपतः ।}\ H}}\skt{ ।}\pend

\pstart \ledsidenote{\passref{6.41.14ab}}\edtext{\skt{एकादसभिः}}{\Afootnote{K; \skt{एकादशभिर्}\ A H}}\skt{ }\edtext{\skt{षड्भिर्वा}}{\lemma{\skt{षड्भिर्}}\Afootnote{K H; \skt{एभिर्वा षड्भिर्}\ A}}\skt{ }\edtext{\skt{पञ्चभिर्वा}}{\Afootnote{K; \skt{पि}\ A; \skt{पञ्चभि}\ H}}\skt{ }\edtext{\skt{समन्वितः}}{\Afootnote{K; \skt{समन्वितम्}\ A H}}\skt{ ।}\pend

\pstart \ledsidenote{\passref{6.41.15cd}}\edtext{\skt{जह्या\-च्छो\-षा\-र्दि\-त\-ज\-न्तु\-मिच्छं}}{\Afootnote{K; \skt{) जह्याच्छोषार्दितं जन्तुमिच्छ}\ A; \skt{जह्याच्छोषार्द्दितं जन्तुमिच्छ}\ H}}\skt{ }\edtext{\skt{सुविपुलं}}{\Afootnote{K; \skt{न्सुविपुलं}\ A H}}\skt{ यशः ।}\pend

\pstart\versepar \ledsidenote{\passref{6.41.16}}\edtext{\skt{व्यवा\-य\-शो\-क\-स्था\-विर्या}}{\Afootnote{K; \skt{व्याय\-शो\-क\-स्था\-विर्य}\ A; \skt{व्यवा\-य\-शो\-क\-स्था\-विर्य}\ H}}\skt{ }\edtext{\skt{व्याया\-मा\-ध्व\-प्र\-शा\-षिणां}}{\Afootnote{K H; \skt{व्याया\-म\-ध्वो\-प\-वा\-सतः ।}\ A}}\skt{ । }\pend
\pstart\versepar \skt{}\edtext{\skt{क्षरति}}{\Afootnote{K; \skt{व्रणोरःक्षत}\ A; \skt{व्रणक्षरति}\ H}}\skt{ }\edtext{\skt{यद्वागुंशोषाणां}}{\Afootnote{K; \skt{पीडाभ्यां शोषा}\ A; \skt{यन्धातुं शोषाणां}\ H}}\skt{ }\edtext{\skt{लक्षणं}}{\Afootnote{K H; \emph{om.}\ A}}\skt{ }\edtext{\skt{शृणु}}{\Afootnote{K H; \skt{नन्ये वदन्ति हि ॥}\ A}}\skt{ ॥}\pend

\pstart\versepar \ledsidenote{\passref{6.41.17}}\skt{व्यवा\-य\-शो\-षी\-शु\-क्रस्य }\edtext{\skt{क्षय\-लि\-ङ्गै\-रू\-प\-द्रुतः}}{\Afootnote{K; \skt{क्षय\-लि\-ङ्गै\-रु\-प\-द्रुतः ।}\ A; \skt{क्षय\-लि\-ङ्गै\-रू\-प\-द्रुतः ।}\ H}}\skt{ }\pend
\pstart\versepar \skt{पाण्डुदेहो यथापूर्वं }\edtext{\skt{क्षीयन्ते}}{\Afootnote{A K; \skt{क्षायन्ते}\ H}}\skt{ चास्य धातवः ।}\pend

\pstart \ledsidenote{\passref{6.41.18ab}}\edtext{\skt{प्रधानशीलः}}{\Afootnote{K H; \skt{प्रध्यानशीलः}\ A}}\skt{ }\edtext{\skt{श्रतांगः}}{\Afootnote{K; \skt{स्रस्ताङ्गः}\ A H}}\skt{ }\edtext{\skt{शोकशोष्यति}}{\Afootnote{K H; \skt{शोकशोष्यपि}\ A}}\skt{ तादृशः ।}\pend

\pstart\versepar \ledsidenote{\passref{6.41.19}}\skt{जराशोषी कृशो }\edtext{\skt{मन्दो}}{\Afootnote{K H; \skt{मन्द}\ A}}\skt{ }\edtext{\skt{हत\-वु\-द्धि\-व\-ले\-न्द्रियः}}{\Afootnote{K; \skt{वीर्य\-बु\-द्धि\-ब\-ले\-न्दिर्यः ।}\ A; \skt{हत{\symfont ¦}वुद्धि बलेन्द्रियः ॥}\ H\textsuperscript{ac}; \skt{हत{\symfont ¦}वु\-द्धि\-श्ब\-ले\-न्द्रियः ॥}\ H\textsuperscript{pc}}}\skt{ । }\pend
\pstart\versepar \skt{}\edtext{\skt{कम्प\-नो\-रू\-चि\-मं\-भि\-न्न\-कां\-स\-पा\-त्र\-स\-म\-स्वरः}}{\Afootnote{K; \skt{कम्प\-नोऽ\-रु\-चि\-मा\-न्भि\-न्न\-कां\-स्य\-पा\-त्र\-ह\-त\-स्व\-रः(नः) ॥}\ A; \skt{कम्प\-नो\-रू\-चि\-मा\-न्भि\-न्न\-कां\-स\-पा\-त्र\-स\-म\-स्वरः ।}\ H}}\skt{ ।}\pend

\pstart \ledsidenote{\passref{6.41.20cd}}\edtext{\skt{सप्र\-स्रु\-ता\-स्य\-ना\-साक्षि}}{\Afootnote{K; \skt{संप्र\-स्रु\-ता\-स्य\-ना\-साक्षः}\ A; \skt{संप्र\-स्रु\-ता\-स्य\-ना\-साक्षि}\ H}}\skt{ }\edtext{\skt{सुप्त\-रू\-क्ष\-म\-ल\-च्छविः}}{\Afootnote{A K; \skt{सुप्तरूक्ष{\symfont ¦}समच्चविः ॥}\ H}}\skt{ ।}\pend

\pstart\versepar \ledsidenote{\passref{6.41.21}}\skt{अध्वप्रशोषी }\edtext{\skt{श्रस्तांगं}}{\Afootnote{K; \skt{स्रस्ताङ्गः}\ A H}}\skt{ }\edtext{\skt{सम्भृ\-ष्ट\-प\-रु\-ष\-च्छविः}}{\Afootnote{K; \skt{संभृ\-ष्ट\-प\-रु\-ष\-च्छविः ।}\ A; \skt{संभृ\-ष्त\-प\-रु\-ष\-च्छविः ।}\ H}}\skt{ }\pend
\pstart\versepar \skt{प्रसु\-प्त\-गा\-त्रा\-व\-यवः शुष्क\-क्लो\-म\-ग\-ला\-ननः ।}\pend

\pstart \ledsidenote{\passref{6.41.20ab}}\skt{ष्ठीवति }\edtext{\skt{श्लेष्मणा}}{\Afootnote{A K; \skt{श्लेष्मणां}\ H}}\skt{ }\edtext{\skt{हीन}}{\Afootnote{K; \skt{हीनं}\ A H}}\skt{ }\edtext{\skt{गौर्वा\-रू\-चि\-पी\-डितः}}{\Afootnote{K; \skt{गौर\-वा\-रु\-चि\-पी\-डितः ।}\ A; \skt{गौर\-वा\-रू\-चि\-पी\-डितः ।}\ H}}\skt{ ।}\pend

\pstart\versepar \ledsidenote{\passref{6.41.22}}\edtext{\skt{व्याया\-म\-शो\-षी\-स्भू\-यि\-ष्ठ\-मे\-ति\-रेव}}{\Afootnote{K; \skt{व्यायामशोषी भूयि\-ष्ठ\-मे\-भि\-रेव}\ A; \skt{व्यवा{\symfont ¦}यशोषी भूयि\-ष्ठ\-मे\-भि\-रेव}\ H}}\skt{ }\edtext{\skt{समन्वितः}}{\Afootnote{K H; \skt{समन्वितः ।}\ A}}\skt{ । }\pend
\pstart\versepar \skt{}\edtext{\skt{लिङ्ग\-रु\-र\-क्ष\-त\-कृतः}}{\Afootnote{K; \skt{उरः\-क्ष\-त\-कृ\-तै\-र्लिङ्गैः}\ A; \skt{लिङ्गै रूक्षयकृतैः}\ H}}\skt{ }\edtext{\skt{संयुक्तश्च}}{\Afootnote{A K; \skt{संयुक्तस्य}\ H}}\skt{ क्षताद्विना ।}\pend

\pstart\versepar \ledsidenote{\passref{6.41.23}}\edtext{\skt{रक्त\-क्ष\-या\-द्वे\-द\-ना\-भि\-स्त\-थै\-वा\-हा\-र\-य\-न्त्र\-णात्}}{\lemma{\skt{रक्त\-क्ष\-या\-द्वे\-द\-ना\-भिस्}}\Afootnote{K; \skt{रक्त\-क्ष\-या\-द्वे\-द\-ना\-भिस्}\ A H}}\skt{ । }\pend
\pstart\versepar \skt{व्रणितस्य भवेच्छोषो }\edtext{\skt{दोषः}}{\Afootnote{K; \skt{स}\ A; \skt{दोषै}\ H}}\skt{ }\edtext{\skt{सामश्मः}}{\Afootnote{K; \skt{चासाध्यतमः स्मृतः ॥}\ A; \skt{ह्̤ साध्यतमस्य च ।}\ H}}\skt{ ।}\pend

\pstart\versepar \ledsidenote{\passref{6.41.24}}\edtext{\skt{व्याया\-म\-भा\-ष्या\-ध्य\-य\-ने\-रै\-भि\-घा\-ता\-ति\-मै\-थुनैः}}{\Afootnote{\skt{व्याया\-म\-भा\-रा\-ध्य\-य\-नै\-र\-भि\-घा\-ता\-ति\-मै\-थुनैः ।}\ A}}\skt{ । }\pend
\pstart\versepar \skt{कर्मणा }\edtext{\skt{वाप्युरस्येन}}{\Afootnote{\skt{चाप्युरस्येन}\ A}}\skt{ }\edtext{\skt{वक्ष्यो}}{\Afootnote{\skt{वक्षो}\ A}}\skt{ यस्य विदारितं ।}\pend

\pstart \ledsidenote{\passref{6.41.ef}}\skt{तस्योपरिक्षते रक्तुयश्लेष्मा च गच्छति ।}\pend

\pstart\versepar \ledsidenote{\passref{6.41.25}}\edtext{\skt{काशं\-प्र\-र्च्छ\-र्द\-ये\-च्चापि}}{\Afootnote{K; \skt{कास\-मा\-न\-श्छ\-र्द\-येच्च}\ A; \skt{कांस\-प्र\-च्छ\-द\-ये\-च्चापि}\ H}}\skt{ }\edtext{\skt{पीत\-र\-क्ता\-सि\-ता\-रूणं}}{\Afootnote{K H; \skt{पीत\-र\-क्ता\-सि\-ता\-रु\-णम् ।}\ A}}\skt{ । }\pend
\pstart\versepar \skt{}\edtext{\skt{सन्त\-प्त\-व\-क्ष\-शो\-त्यर्थं}}{\Afootnote{K; \skt{संतप्तवक्षाः सोऽत्यर्थं}\ A; \skt{सन्त\-प्त\-व\-क्ष\-सोऽ\-त्यर्थं}\ H}}\skt{ }\edtext{\skt{दूयमानेन}}{\Afootnote{K H; \skt{दूयनात्परि}\ A}}\skt{ }\edtext{\skt{ता}}{\Afootnote{K; \skt{ताम्यति}\ A H}}\skt{ ।}\pend

\pstart\versepar \ledsidenote{\passref{6.41.26}}\edtext{\skt{वदनो}}{\Afootnote{K; \skt{दुर्ग\-न्ध\-व\-द\-नो\-च्छ्वासो}\ A; \skt{दुर्ग\-न्धो\-च्च्वा\-स\-व\-दनो}\ H}}\skt{ }\edtext{\skt{भिन्नवर्णस्वरो}}{\Afootnote{A K; \skt{भिन्नवर्ण्णः स्वरो}\ H}}\skt{ नरः । }\pend
\pstart\versepar \skt{}\edtext{\skt{केषा\-ञ्चि\-दे\-व\-शो\-षोहि}}{\Afootnote{K; \skt{केषांचिदेवं शोषो हि}\ A; \skt{केषाञ्चिदेव शोषो हि}\ H}}\skt{ }\edtext{\edtext{\skt{कार\-णै\-भे\-द\-मा\-गतः}}{\lemma{\skt{कारणैभेदम्}}\Afootnote{K; \skt{कारणैर्भेद}\ A; \skt{कारणो भेदम्}\ H}}}{\lemma{\skt{आगतः}}\Afootnote{K H; \skt{मागतः ॥}\ A}}\skt{ ।}\pend

\pstart\versepar \ledsidenote{\passref{6.41.27}}\skt{न तत्र }\edtext{\skt{दोष\-लि\-ङ्गा\-ना\-स\-म\-स्तानां}}{\Afootnote{K; \skt{दोषलिङ्गानां समस्तानां}\ A H}}\skt{ }\edtext{\skt{निपातनं}}{\Afootnote{A K; \skt{निवातनं ।}\ H}}\skt{ । }\pend
\pstart\versepar \skt{क्षया एव हि }\edtext{\skt{विज्ञेया}}{\Afootnote{K H; \skt{ते ज्ञेयाः}\ A}}\skt{ }\edtext{\skt{प्रत्ये\-क\-धा\-तु\-स\-ञ्ज्ञिताः}}{\Afootnote{K; \skt{प्रत्येकं धातुसंज्ञिताः ॥}\ A; \skt{प्रत्ये\-क\-न्धा\-तु\-सं\-ज्ञिताः ॥}\ H}}\skt{ ।}\pend

\pstart \ledsidenote{\passref{6.41.28}}\skt{चिकित्सितन्तु तेषां हि प्रागुक्तं धातुसंक्षये ॥}\pend

\pstart\versepar \ledsidenote{\passref{6.41.29}}\edtext{\skt{श्वासां\-ग\-सा\-द\-क\-फ\-सं\-श्र\-य\-ता\-लु\-शो\-षा\-च्छ\-र्द्य\-ग्नि\-सा\-द\-प\-रि\-पी\-न\-स\-पा\-ण्डु\-निद्राः}}{\Afootnote{K; \skt{श्वासा\-ङ्ग\-सा\-द\-क\-फ\-सं\-स्र\-व\-ता\-लु\-शो\-ष\-च्छ\-र्द्य\-ग्नि\-सा\-द\-म\-द\-पी\-न\-स\-पा\-ण्डु\-निद्राः ।}\ A; \skt{श्वासाङ्गसादक{\symfont ¦}फ\-सं\-श्र\-य\-ता\-लु\-शो\-अ\-छ\-र्द्य\-ग्नि\-सा\-द\-प\-रि\-पी\-न\-स\-पा\-ण्डु\-निद्रा ।}\ H}}\skt{ । }\pend
\pstart\versepar \skt{शोषे भविष्यति भवन्ति स चापि जंतुः }\edtext{\skt{शुक्लक्षणो}}{\Afootnote{K; \skt{शुक्लेक्षणे}\ A; \skt{शुक्लक्षण}\ H}}\skt{ भवति मांसपरो }\edtext{\skt{रिरंसुः}}{\Afootnote{K H; \skt{रिरिंसुः ॥}\ A}}\skt{ ।}\pend

\pstart\versepar \ledsidenote{\passref{6.41.30}}\skt{स्वप्नेषु }\edtext{\skt{काक\-शु\-क\-श\-ल्य\-क\-नी\-ल\-क\-ण्ठा\-गृ\-ध्रा\-न्त\-थैव}}{\lemma{\skt{काक\-शु\-क\-श\-ल्य\-क\-नी\-ल\-क\-ण्ठा\-गृ\-ध्रान्}}\Afootnote{K; \skt{काक\-शु\-क\-श\-ल्ल\-कि\-नी\-ल\-क\-ण्ठ\-गृ\-ध्रास्}\ A; \skt{काक\-शु\-क\-श\-ल्य\-क\-नी\-ल\-क\-ण्ठा\-न्गृ\-ध्रास्}\ H}}\skt{ }\edtext{\skt{च}}{\Afootnote{K H; \skt{क}\ A}}\skt{ }\edtext{\skt{कपीं}}{\Afootnote{K; \skt{पयः}\ A; \skt{कपी}\ H}}\skt{ }\edtext{\skt{कृकलासकांश्च}}{\Afootnote{K; \skt{कृकलासकाश्च ।}\ A; \skt{स्कृकलासेकांश्च ।}\ H}}\skt{ । }\pend
\pstart\versepar \skt{}\edtext{\skt{पश्येत्सरांसि}}{\Afootnote{K; \skt{तं}\ A; \skt{पश्चेत्सरांसि}\ H}}\skt{ }\edtext{\skt{स}}{\Afootnote{K; \skt{वाहयन्ति स}\ A; \skt{विमलआः स}\ H}}\skt{ }\edtext{\edtext{\skt{नदीम्विमलांल्स्च}}{\lemma{\skt{नदीम्}}\Afootnote{K; \skt{नदी}\ A; \skt{नदीश्}\ H}}}{\lemma{\skt{विमलांल्स्च}}\Afootnote{K; \skt{र्विजलाश्च}\ A; \skt{च}\ H}}\skt{ }\edtext{\skt{वापीं}}{\Afootnote{K; \skt{पश्ये}\ A; \skt{वाइः}\ H}}\skt{ }\edtext{\skt{शुष्कां\-म्स्त\-रू\-न\-रु\-ण\-धू\-म\-वा\-र्दि\-ताश्च}}{\lemma{\skt{शुष्कां\-म्स्त\-रू\-न\-रु\-ण\-धू\-म\-वा\-र्दि\-ताश्}}\Afootnote{K; \skt{च्छुष्कां\-स्त\-रू\-न्प\-व\-न\-धू\-म\-द\-वा\-र्दि\-तांश्}\ A; \skt{शुष्कां\-स्त\-रू\-न\-रू\-ण\-धू\-म\-द\-वा\-र्द्दि\-तांश्}\ H}}\skt{ ॥}\pend

\pstart\versepar \ledsidenote{\passref{6.41.31}}\edtext{\skt{महासनं}}{\Afootnote{K; \skt{महाशनं}\ A; \skt{अहाशनं}\ H}}\skt{ }\edtext{\skt{क्षीयमाणां}}{\Afootnote{K; \skt{क्षीयमाण}\ A; \skt{परिक्षीणं}\ H}}\skt{ }\edtext{\skt{सोषो\-ति\-सा\-र\-पी\-डितं}}{\Afootnote{K; \skt{मती\-सा\-र\-नि\-पी\-डि\-तम् ।}\ A; \skt{शोषा\-ती\-सा\-र\-पी\-डितं ।}\ H}}\skt{ । }\pend
\pstart\versepar \skt{}\edtext{\edtext{\skt{शून\-मू\-स्को\-द\-र\-श्चापि}}{\lemma{\skt{शून\-मू\-स्को\-द\-रश्}}\Afootnote{K; \skt{शूनमुष्कोदरं}\ A; \skt{शून\-मू\-ष्को\-द\-रञ्}\ H}}}{\lemma{\skt{चापि}}\Afootnote{K H; \skt{चैव}\ A}}\skt{ यक्ष्मि\-ण\-म्प\-रि\-व\-र्ज\-येत् ।}\pend

\pstart\versepar \ledsidenote{\passref{6.41.32}}\edtext{\skt{उपा\-च\-रे\-दा\-त्म\-वन्तं}}{\Afootnote{K H; \skt{उपा\-च\-रे\-दा\-त्म\-वन्तं}\ A}}\skt{ }\edtext{\edtext{\edtext{\skt{दीप्ता\-ग्नि\-म\-कृ\-श\-न्नर{\symfont ¦}}}{\lemma{\skt{दीप्ताग्निम्}}\Afootnote{K; \skt{दीप्ताग्नि}\ A H}}}{\lemma{\skt{अकृशन्}}\Afootnote{K; \skt{मकृशं}\ A H}}}{\lemma{\skt{नर{\symfont ¦}}}\Afootnote{K; \skt{नवम् ।}\ A; \skt{नरं ।}\ H}}\skt{ । }\pend
\pstart\versepar \skt{स्थिरा\-दि\-व\-र्ग\-सि\-द्धेन }\edtext{\skt{घृतेना}}{\Afootnote{K; \skt{धृतेना}\ A; \skt{पयसा}\ H}}\skt{ जाविकेन }\edtext{\skt{वा}}{\Afootnote{K; \skt{च ॥}\ A; \skt{वै ॥}\ H}}\skt{ ।}\pend

\pstart\versepar \ledsidenote{\passref{6.41.33}}\edtext{\skt{स्निग्धस्य}}{\Afootnote{A K H\textsuperscript{pc}; \skt{स्निग्ध}\ H\textsuperscript{ac}}}\skt{ मृदु }\edtext{\skt{कर्त\-व्य\-मू\-र्ध्व\-श्चा\-धश्च}}{\lemma{\skt{कर्तव्यमूर्ध्वश्}}\Afootnote{K; \skt{कर्तव्यमूर्ध्वं}\ A; \skt{कर्तव्यमूर्ध्वंञ्}\ H}}\skt{ शोधनं । }\pend
\pstart\versepar \skt{}\edtext{\skt{आस्थापनन्ततः}}{\lemma{\skt{ततः}}\Afootnote{K; \skt{तथा}\ A; \skt{त्थाः}\ H}}\skt{ }\edtext{\skt{पश्चाच्छिरस}}{\Afootnote{K; \skt{कार्यं शिरसश्च}\ A; \skt{पश्चाच्छिरसश्च}\ H}}\skt{ }\edtext{\skt{विरेचन{\symfont ¦}}}{\Afootnote{K; \skt{विरेचनम्}\ A H}}\skt{ ।}\pend

\pstart \ledsidenote{\passref{6.41.34ab}}\edtext{\edtext{\skt{रसैर्भुञ्जीत}}{\lemma{\skt{रसैर्}}\Afootnote{K H; \emph{om.}\ A}}}{\lemma{\skt{भुञ्जीत}}\Afootnote{K H; \emph{om.}\ A}}\skt{ }\edtext{\skt{शुद्धश्च}}{\Afootnote{K H; \emph{om.}\ A}}\skt{ }\edtext{\skt{यव\-गो\-धू\-म\-शा\-लयः}}{\Afootnote{K H; \skt{यव\-गो\-धू\-म\-शा\-लींश्च रसैर्भुञ्जीत शोधितः ।}\ A}}\skt{ ।}\pend

\pstart\versepar \ledsidenote{\passref{6.41.35}}\skt{व्यवायशोषिणं प्रायो भजन्ते वातजा गदाः । }\pend
\pstart\versepar \skt{}\edtext{\skt{वृंहनीलो}}{\Afootnote{K; \skt{बृंहणीयो}\ A; \skt{वृंहणीयो}\ H}}\skt{ }\edtext{\edtext{\skt{हितस्तस्य}}{\lemma{\skt{हितस्}}\Afootnote{K H; \skt{विधि}\ A}}}{\lemma{\skt{तस्य}}\Afootnote{K H; \skt{स्तस्मै}\ A}}\skt{ }\edtext{\skt{हित\-स्नि\-ग्धा\-नि\-ला\-पहः}}{\Afootnote{K H; \skt{हितः स्निघ्दोऽनिलापहः ॥}\ A}}\skt{ ।}\pend

\pstart\versepar \ledsidenote{\passref{6.41.36}}\edtext{\skt{गृध्रातुलूकां}}{\Afootnote{K; \skt{काकानुलूकान्}\ A; \skt{गृध्रानुलूका}\ H}}\skt{ }\edtext{\skt{सकु\-ना\-म्त\-था\-ग\-ण्द\-पदां}}{\Afootnote{K; \skt{नकु\-ला\-न्बि\-डा\-ला\-न्ग\-ण्डू\-पदा}\ A; \skt{न्शकु\-नां\-स्त\-था\-खू\-न्ग\-ण्ड\-पदा}\ H}}\skt{ }\edtext{\edtext{\skt{व्याल\-वि\-ले\-श\-याश्च}}{\lemma{\skt{व्याल\-वि\-ले\-श\-याश्}}\Afootnote{K; \skt{न्व्यालबिलेशया}\ A; \skt{न्व्यालविलेश}\ H}}}{\lemma{\skt{च}}\Afootnote{K; \skt{खून् । गृध्रांश्च}\ A; \skt{याञ्श्च ॥}\ H}}\skt{ । }\pend
\pstart\versepar \skt{}\edtext{\skt{दद्यात्प्रकारै}}{\Afootnote{K; \skt{दद्या}\ A; \skt{दद्यात्प्रकारैर्}\ H}}\skt{ }\edtext{\skt{द्विविधैर्विधिज्ञः ।सर्सतृवां}}{\lemma{\skt{विधिज्ञः ।सर्सतृवां}}\Afootnote{K; \skt{प्रवादैः ससैन्धवा}\ A; \skt{व्विधिज्ञः ससैन्धवा}\ H}}\skt{ }\edtext{\skt{शर्ष\-प\-तै\-ल\-भृष्टां}}{\Afootnote{K; \skt{न्सर्ष\-प\-तै\-ल\-भृ\-ष्टान् ॥}\ A H}}\skt{ ।}\pend

\pstart\versepar \ledsidenote{\passref{6.41.37}}\skt{देयानि मांसानि }\edtext{\skt{वजाङ्गलानि}}{\Afootnote{K H; \skt{च जाङ्गलानि}\ A}}\skt{ }\edtext{\skt{मुद्गा\-ढ\-की\-यू\-ष\-राश्च}}{\Afootnote{K; \skt{मुद्गा\-ढ\-की\-सू\-प\-र\-साश्च}\ A; \skt{मुद्गा\-ढ\-की\-यू\-ष\-र\-साश्च}\ H}}\skt{ }\edtext{\skt{पथ्याः}}{\Afootnote{K H; \skt{हृद्याः ।}\ A}}\skt{ । }\pend
\pstart\versepar \skt{}\edtext{\edtext{\skt{खरोष्ट्रनागाश्च}}{\lemma{\skt{खरोष्ट्रनागाश्}}\Afootnote{K H; \skt{खरोष्ट्रनागा}\ A}}}{\lemma{\skt{च}}\Afootnote{K H; \skt{श्व}\ A}}\skt{ }\edtext{\skt{तराशि}}{\Afootnote{K H; \skt{तराश्व}\ A}}\skt{ }\edtext{\skt{चैव}}{\Afootnote{K; \skt{जा}\ A; \emph{om.}\ H}}\skt{ }\edtext{\skt{खादेत्मुभृष्ठानि}}{\Afootnote{K H; \skt{नि देयानि मांसानि}\ A}}\skt{ सुकल्पितानि ।}\pend

\pstart\versepar \ledsidenote{\passref{6.41.38}}\edtext{\skt{मांसोपदंशाञ्श्च}}{\Afootnote{A H; \skt{मांसोपदंशां}\ K}}\skt{ }\edtext{\skt{पिवे\-द\-रि\-ष्टा\-न्मा\-ध्वी\-क\-सेवी}}{\Afootnote{H; \skt{पिबे\-द\-रि\-ष्टा\-न्मा\-र्द्वी\-क\-युक्ता}\ A; \emph{om.}\ K}}\skt{ }\edtext{\skt{च भवेत्मनुष्यः}}{\Afootnote{H; \skt{न्मदिराश्च सेव्याः ।}\ A; \emph{om.}\ K}}\skt{ । }\pend
\pstart\versepar \skt{}\edtext{\skt{अर्क्का\-मृ\-ता\-क्षा\-र\-ज\-लो\-षि\-तेभ्यः}}{\Afootnote{A H; \emph{om.}\ K}}\skt{ }\edtext{\skt{कृत्वा}}{\Afootnote{A H; \emph{om.}\ K}}\skt{ }\edtext{\skt{यवेभ्यो}}{\Afootnote{A H; \emph{om.}\ K}}\skt{ }\edtext{\skt{विविधाञ्श्च}}{\Afootnote{A H; \emph{om.}\ K}}\skt{ }\edtext{\skt{भक्ष्यान्}}{\Afootnote{A H; \emph{om.}\ K}}\skt{ ॥}\pend

\pstart\versepar \ledsidenote{\passref{6.41.39}}\edtext{\skt{खादेत्पिवेत्चापि}}{\Afootnote{H; \skt{खादेत्पिबेत्सर्पि}\ A}}\skt{ }\edtext{\skt{अजाविकं वा}}{\Afootnote{H; \skt{रजाविकं वा}\ A}}\skt{ कृशो यवाग्वा सह भक्तकाले । }\pend
\pstart\versepar \skt{}\edtext{\skt{व्योषं}}{\Afootnote{H; \emph{om.}\ A}}\skt{ }\edtext{\skt{घृत\-क्षौ\-द्र\-यु\-तञ्च लिह्याच्च}}{\Afootnote{H; \skt{सर्पिर्मधुभ्यां त्रिकटु प्रलिह्याच्च}\ A}}\skt{ व्यावि\-ड\-ङ्गो\-प\-हितं क्षयार्त्तः ॥}\pend

\pstart\versepar \ledsidenote{\passref{6.41.40}}\skt{मांसादमांसेषु घृतञ्च सिद्धं }\edtext{\skt{शोषापहं}}{\Afootnote{H; \skt{शोसापहं}\ A}}\skt{ }\edtext{\skt{क्षौद्रयुतं}}{\Afootnote{H; \skt{क्षौद्रकणा}\ A}}\skt{ }\edtext{\skt{वदन्ति}}{\Afootnote{H; \skt{समेतम् ।}\ A}}\skt{ । }\pend
\pstart\versepar \skt{}\edtext{\edtext{\skt{द्राक्षा\-सि\-ता\-मा\-ग\-धि\-काश्च}}{\lemma{\skt{द्राक्षा\-सि\-ता\-मा\-ग\-धि\-काश्}}\Afootnote{H; \skt{द्राक्षा\-सि\-ता\-मा\-ग\-धिका}\ A}}}{\lemma{\skt{च}}\Afootnote{H; \skt{व}\ A}}\skt{ }\edtext{\skt{लेहाः}}{\Afootnote{H; \skt{लेहः}\ A}}\skt{ }\edtext{\skt{सक्षौद्रतैलं}}{\Afootnote{H; \skt{सक्षौद्रतैलः}\ A}}\skt{ }\edtext{\skt{क्षयकासहन्ता}}{\Afootnote{H; \skt{क्षयरोगघाती ॥}\ A}}\skt{ ॥}\pend

\pstart\versepar \ledsidenote{\passref{6.41.41}}\skt{घृतेन चाजेन }\edtext{\skt{समाक्षिकेन}}{\Afootnote{H; \skt{समाक्षिकेण}\ A}}\skt{ तुर\-ङ्ग\-ग\-न्धा\-ति\-ल\-मा\-ष\-चूर्ण्णं । }\pend
\pstart\versepar \skt{सिता\-श्व\-ग\-न्धा\-म\-ग\-धो\-द्भ\-वानां }\edtext{\skt{चूर्ण\-घृ\-त\-क्षौ\-द्र\-युतं}}{\Afootnote{H; \skt{चूर्णं घृतक्षौद्रयुतं}\ A}}\skt{ }\edtext{\skt{च}}{\Afootnote{H; \skt{प्र}\ A}}\skt{ }\edtext{\skt{लिह्या}}{\Afootnote{H; \skt{लिह्यात् ॥}\ A}}\skt{ ॥}\pend

\pstart\versepar \ledsidenote{\passref{6.41.42}}\skt{क्षीरं }\edtext{\skt{पिवेद्वाप्यथ}}{\Afootnote{H; \skt{पिबेद्वाप्यथ}\ A}}\skt{ }\edtext{\edtext{\skt{वाजि\-ग\-न्धा\-वि\-प\-क्व\-मेवं}}{\lemma{\skt{वाजि\-ग\-न्धा\-वि\-प\-क्वम्}}\Afootnote{H; \skt{वाजि\-ग\-न्धा\-वि\-पक्व}\ A}}}{\lemma{\skt{एवं}}\Afootnote{H; \skt{मेवं}\ A}}\skt{ }\edtext{\skt{लभतेङ्गमुष्ठिं}}{\Afootnote{H; \skt{लभ\-तेऽ\-ङ्ग\-पु\-ष्टिम् ।}\ A}}\skt{ । }\pend
\pstart\versepar \skt{तदुत्थितं क्षीरघृतं सिताढ्यं प्रातः }\edtext{\skt{पिवेद्वापि}}{\Afootnote{H; \skt{पिबेद्वापि}\ A}}\skt{ पयोनुपानम् ॥}\pend

\pstart\versepar \ledsidenote{\passref{6.41.43}}\edtext{\skt{सेवते}}{\Afootnote{H; \emph{om.}\ A}}\skt{ }\edtext{\skt{चाजेन}}{\Afootnote{H; \emph{om.}\ A}}\skt{ }\edtext{\skt{समाक्षिकेण}}{\Afootnote{H; \emph{om.}\ A}}\skt{ }\edtext{\skt{तुर\-ङ्ग\-ग\-न्धा\-ति\-ल\-मा\-ष\-चूर्णं}}{\Afootnote{H; \emph{om.}\ A}}\skt{ । }\pend
\pstart\versepar \skt{उत्सादने चापि तुरंगगन्धा योज्या }\edtext{\edtext{\skt{यवाश्चैव}}{\lemma{\skt{यवाश्}}\Afootnote{A; \skt{यवा}\ H}}}{\lemma{\skt{चैव}}\Afootnote{A; \skt{श्चैव}\ H}}\skt{ }\edtext{\skt{पुनर्नवो च}}{\Afootnote{H; \skt{पुनर्नवे च ।}\ A}}\skt{ ॥ }\pend
\pstart\versepar \skt{कृत्स्ने वृषे तत्कुसुमैश्च सिद्धं सर्प्पिः }\edtext{\skt{पिवेत्क्षौद्रहितं}}{\Afootnote{H; \skt{पिबेत्क्षौद्रयुतं}\ A}}\skt{ हिताशी ।}\pend

\pstart\versepar \ledsidenote{\passref{6.41.44}}\edtext{\skt{यक्ष्मा\-ण\-मे\-त\-त्प्र\-व\-लञ्च}}{\Afootnote{H; \skt{यक्ष्मा\-ण\-मे\-त\-त्प्र\-बलं च}\ A}}\skt{ कासं }\edtext{\skt{श्वास\-न्नि\-ह\-न्या\-दपि}}{\lemma{\skt{निहन्यादपि}}\Afootnote{H; \skt{च हन्यादपि}\ A}}\skt{ पाण्डुताञ्च ॥ }\pend
\pstart\versepar \skt{}\edtext{\edtext{\skt{शकृ\-द्र\-सा\-श्चा\-श्व\-ग\-वा\-ज\-माव्याः}}{\lemma{\skt{शकृद्रसाश्}}\Afootnote{H; \skt{शकृद्रसा}\ A}}}{\lemma{\skt{चाश्व\-ग\-वा\-ज\-माव्याः}}\Afootnote{H; \skt{गोश्व\-ग\-जा\-व्य\-जानां}\ A}}\skt{ क्वाथा }\edtext{\skt{मिताश्चैव}}{\lemma{\skt{चैव}}\Afootnote{H; \skt{चापि}\ A}}\skt{ तत्थैव भागैः ।}\pend

\pstart\versepar \ledsidenote{\passref{6.41.45}}\skt{मूर्व्वा\-ह\-रि\-द्रा\-ख\-दि\-र\-द्रु\-माणां क्षीरस्य }\edtext{\edtext{\skt{भागस्त्वपरो}}{\lemma{\skt{भागस्}}\Afootnote{H; \skt{भाग}\ A}}}{\lemma{\skt{त्वपरो}}\Afootnote{H; \skt{स्त्वपरो}\ A}}\skt{ घृतस्य ॥ }\pend
\pstart\versepar \skt{}\edtext{\skt{भागान्नवैता}}{\lemma{\skt{नवैता}}\Afootnote{H; \skt{दशौता}\ A}}\skt{ }\edtext{\skt{विपचेद्विदिज्ञो}}{\Afootnote{H; \skt{न्विपचेद्विधिज्ञो}\ A}}\skt{ दत्वा त्रिव\-र्ग\-म्म\-धु\-रञ्च कृत्स्नं ।}\pend

\pstart\versepar \ledsidenote{\passref{6.41.46}}\skt{कटुत्रिकञ्चैव }\edtext{\skt{सभ\-द्र\-दा\-रू\-घृ\-तो\-तमं}}{\Afootnote{H; \skt{सभद्रदारु घृतोत्तमं}\ A}}\skt{ यक्ष्म\-नि\-व्-\-अ\-र\-णीयं ॥ }\pend
\pstart\versepar \skt{द्वे }\edtext{\skt{पञ्चमूले}}{\Afootnote{H; \skt{पञ्चमूल्यौ}\ A}}\skt{ }\edtext{\skt{वरुणः}}{\Afootnote{H; \skt{वरुणं}\ A}}\skt{ }\edtext{\skt{करञ्ज}}{\Afootnote{H; \skt{करञ्जं}\ A}}\skt{ }\edtext{\skt{भल्ला\-त\-का\-वि\-ल्व\-पु\-न\-र्नवो}}{\Afootnote{H; \skt{भल्लातकं बिल्वपुनर्नवे च ॥}\ A}}\skt{ }\edtext{\skt{च}}{\Afootnote{H; \emph{om.}\ A}}\skt{ ॥}\edtext{}{\lemma{\skt{यक्ष्म\-नि\-व्-\-अ\-र\-णीयं ॥ द्वे}}\Afootnote{H; \skt{यक्ष्म\-नि\-वा\-र\-णाय । द्वे}\ A}}\pend

\pstart\versepar \ledsidenote{\passref{6.41.47}}\edtext{\edtext{\skt{यवान्कुलत्थाञ्श्च}}{\lemma{\skt{यवान्कुलत्थाञ्श्}}\Afootnote{H; \skt{यवान्कुलत्था}\ A}}}{\lemma{\skt{च}}\Afootnote{H; \skt{न्ब}\ A}}\skt{ }\edtext{\skt{वराणि}}{\Afootnote{H; \skt{दराणि}\ A}}\skt{ }\edtext{\skt{चैव}}{\Afootnote{H; \skt{भार्गीं पाठा}\ A}}\skt{ }\edtext{\skt{स}}{\Afootnote{H; \skt{ं}\ A}}\skt{ }\edtext{\edtext{\skt{चित्तकञ्चापि}}{\lemma{\skt{चित्तकञ्}}\Afootnote{H; \skt{हुताशं}\ A}}}{\lemma{\skt{चापि}}\Afootnote{H; \skt{स}\ A}}\skt{ }\edtext{\skt{महीकदम्वं}}{\Afootnote{H; \skt{महीकदम्बम् ।}\ A}}\skt{ ॥ }\pend
\pstart\versepar \skt{कृत्वा कषा\-य\-म्वि\-प\-चेधि तस्य }\edtext{\skt{षड्भिर्हि}}{\Afootnote{H; \skt{षड्भिर्हि}\ A}}\skt{ }\edtext{\skt{पात्रे}}{\Afootnote{H; \skt{पात्रै}\ A}}\skt{ }\edtext{\edtext{\skt{घृतपात्रमेकं}}{\lemma{\skt{घृतपात्रम्}}\Afootnote{H; \skt{र्घृतपात्र}\ A}}}{\lemma{\skt{एकं}}\Afootnote{H; \skt{मेकम् ॥}\ A}}\skt{ ।}\pend

\pstart\versepar \ledsidenote{\passref{6.41.48}}\edtext{\skt{महा\-वृ\-क्ष\-प\-यो\-भ\-याञ्च}}{\Afootnote{\skt{व्योषं महा\-वृ\-क्ष\-प\-योऽ\-भयां च}\ A; \emph{om.}\ K; \skt{दत्वा महा\-वृ\-क्ष\-प\-यो\-भ\-याञ्च}\ H}}\skt{ }\edtext{\skt{व्यांशु\-का.\-ल\-व\-णो\-त्त\-मञ्च}}{\Afootnote{H; \skt{चव्यं सुराख्यं लवणोत्तमं च ।}\ A; \emph{om.}\ K}}\skt{ ॥ }\pend
\pstart\versepar \skt{}\edtext{\skt{तद्धि}}{\Afootnote{K; \skt{एतद्धि}\ A; \skt{एतधि}\ H}}\skt{ शोषं जठराणि चैव हन्यात्प्रमेहांश्च सहानिलेन ।}\pend

\pstart\versepar \ledsidenote{\passref{6.41.55}}\skt{प्लीहोदरोक्तं }\edtext{\skt{तु}}{\Afootnote{K H; \skt{वि}\ A}}\skt{ }\edtext{\skt{हितश्च}}{\Afootnote{K; \skt{हितं च}\ A H}}\skt{ }\edtext{\skt{सर्पिस्त्रीण्येव}}{\Afootnote{A K; \skt{सर्प्पिस्त्रीण्येव}\ H}}\skt{ चान्यानि हितानि चात्र । }\pend
\pstart\versepar \skt{}\edtext{\skt{उप\-द्र\-का\-स्व\-र\-वै\-कृ\-ता\-दी\-ञ्ज\-ये\-द\-वे\-क्ष\-शि\-खि\-लेन}}{\Afootnote{K; \skt{उपद्रवांश्च स्वर\-वै\-कृ\-ता\-दी\-न्ज\-ये\-द्य\-थास्वं प्रसमीक्ष्य शास्त्रम् ॥}\ A; \skt{उपद्रवाञ्श्च स्वर\-वै\-कृ\-ता\-दीन् ।ये\-द\-व।\-नि\-खि\-लेन}\ H}}\skt{ }\edtext{\skt{शांतं}}{\Afootnote{K; \emph{om.}\ A; \skt{शान्तां ।}\ H}}\skt{ ।}\pend

\pstart\versepar \ledsidenote{\passref{6.41.50}}\edtext{\edtext{\skt{एला\-ज\-मो\-दा\-म\-ले\-का\-भ\-योश्च}}{\lemma{\skt{एला\-ज\-मो\-दा\-म\-ले\-का\-भ\-योश्}}\Afootnote{K; \skt{एला\-ज\-मो\-दा\-म\-ल\-का\-भया}\ A; \skt{एला\-ज\-मो\-दा\-म\-ल\-का\-भ\-याश्}\ H}}}{\lemma{\skt{च}}\Afootnote{K H; \skt{क्ष}\ A}}\skt{ }\edtext{\skt{गाय\-त्रि\-नि\-म्वा\-स\-न\-सा\-ल\-सारा}}{\Afootnote{K; \skt{गाय\-त्र्य\-रि\-ष्टा\-स\-न\-शा\-ल\-सारा}\ A; \skt{गाय\-त्रि\-नि\-म्वा\-स\-न\-सा\-ल\-सारां ॥}\ H}}\skt{ ॥ }\pend
\pstart\versepar \skt{}\edtext{\edtext{\skt{विड\-ङ्ग\-भ\-ल्ला\-त\-क\-चि\-त्र\-काश्च}}{\lemma{\skt{विड\-ङ्ग\-भ\-ल्ला\-त\-क\-चि\-त्र\-काश्}}\Afootnote{K; \skt{न् । विड\-ङ्ग\-भ\-ल्ला\-त\-क\-चि\-त्रको}\ A; \skt{विड\-ङ्ग\-भ\-ल्ला\-त\-क\-चि\-त्र\-काञ्श्}\ H}}}{\lemma{\skt{च}}\Afootnote{K H; \skt{ग्रा}\ A}}\skt{ }\edtext{\edtext{\skt{कटुत्रिकश्चापि}}{\lemma{\skt{कटुत्रिकश्}}\Afootnote{K; \skt{कटुत्रिका}\ A; \skt{कटुत्रिक..}\ H}}}{\lemma{\skt{चापि}}\Afootnote{K; \skt{म्भोद}\ A; \emph{om.}\ H}}\skt{ }\edtext{\edtext{\skt{सुराष्ट्रजाश्च}}{\lemma{\skt{सुराष्ट्रजाश्}}\Afootnote{K; \skt{सुराष्ट्रजांश्}\ A; \skt{सुराष्ट्रजञ्श् ।}\ H}}}{\lemma{\skt{च}}\Afootnote{A K; \emph{om.}\ H}}\skt{ ।}\pend

\pstart\versepar \ledsidenote{\passref{6.41.51}}\edtext{\skt{पक्वा}}{\Afootnote{K H; \skt{पक्त्वा}\ A}}\skt{ जले }\edtext{\skt{तेन}}{\Afootnote{A K; \skt{येन}\ H}}\skt{ }\edtext{\edtext{\skt{पचेत्सर्पिस्तस्मिं}}{\lemma{\skt{पचेत्सर्पिस्}}\Afootnote{K H; \skt{पचेद्धि सर्पि}\ A}}}{\lemma{\skt{तस्मिं}}\Afootnote{K; \skt{स्तस्मि}\ A; \skt{तस्मि}\ H}}\skt{ }\edtext{\skt{सुसिद्धे}}{\Afootnote{K; \skt{न्सुसिद्धे}\ A H}}\skt{ त्ववतारिते च । }\pend
\pstart\versepar \skt{}\edtext{\skt{त्रिंशत्पलान्यत्र}}{\Afootnote{A K; \skt{त्रिंशत्पलान्यत्र}\ H}}\skt{ सितोपलाया }\edtext{\edtext{\skt{दद्या\-न्तु\-गा\-क्षी\-रि\-प\-लानि}}{\lemma{\skt{दद्यान्}}\Afootnote{K; \skt{दत्त्वा}\ A; \skt{दद्या}\ H}}}{\lemma{\skt{तुगा\-क्षी\-रि\-प\-लानि}}\Afootnote{A K; \skt{त्तुगा\-क्षी\-रि\-प\-लानि}\ H}}\skt{ }\edtext{\skt{षट्वा}}{\Afootnote{K; \skt{षट्च ॥}\ A; \skt{पञ्च ।}\ H}}\skt{ ॥}\pend

\pstart\versepar \ledsidenote{\passref{6.41.52}}\skt{प्रस्थे }\edtext{\skt{घ्र्̥तस्य}}{\Afootnote{K; \skt{घृतस्य}\ A H}}\skt{ }\edtext{\skt{द्विगुणश्च}}{\Afootnote{K; \skt{द्विगुणं च}\ A H}}\skt{ }\edtext{\skt{दद्यात्क्षौद्रं}}{\Afootnote{A K; \skt{दद्या}\ H}}\skt{ }\edtext{\skt{ततो}}{\Afootnote{A K; \emph{om.}\ H}}\skt{ }\edtext{\skt{मन्थहृतं}}{\Afootnote{K; \skt{मन्थहतं}\ A; \emph{om.}\ H}}\skt{ विदद्यात् । }\pend
\pstart\versepar \skt{पलं }\edtext{\skt{पलं}}{\Afootnote{A K; \skt{य}\ H}}\skt{ }\edtext{\skt{प्रातरतो}}{\Afootnote{A K; \skt{प्रातरेतो{\symfont ¦}}\ H}}\skt{ }\edtext{\skt{लिहेच्च}}{\Afootnote{K H; \skt{प्रलिह्य}\ A}}\skt{ }\edtext{\skt{लीह्वा}}{\Afootnote{K H; \skt{पश्चा}\ A}}\skt{ }\edtext{\edtext{\skt{पिबे\-त्क्षी\-र\-म\-त\-न्द्रि\-तस्तु}}{\lemma{\skt{पिबे\-त्क्षी\-र\-म\-त\-न्द्रि\-तस्}}\Afootnote{\skt{त्पिबे\-त्क्षी\-र\-म\-त\-न्द्रि\-तश्}\ A; \skt{पिवे\-त्क्षी\-र\-म\-त\-न्द्रि\-तस्}\ K; \skt{पिवे\-न्क्षी\-र\-म\-त\-न्द्रि\-तस्}\ H}}}{\lemma{\skt{तु}}\Afootnote{K H; \skt{च ॥}\ A}}\skt{ ।}\edtext{}{\lemma{\skt{विदद्यात् । पलं}}\Afootnote{K; \skt{विदध्यात् । पलं}\ A; \skt{त्पलम्}\ H}}\pend

\pstart\versepar \ledsidenote{\passref{6.41.53}}\skt{एतद्धि मेध्य\-म्प\-र\-म\-म्प\-वित्रं }\edtext{\edtext{\skt{चक्षु\-स्य\-मा\-यु\-त\-म\-न्त\-थैव}}{\lemma{\skt{चक्षु\-स्य\-मा\-यु\-त\-मन्}}\Afootnote{K; \skt{चक्षु\-ष्य\-मा\-यु\-ष्य\-मथो}\ A; \skt{चक्षु\-ष्य\-मा\-यु\-ष\-त\-मन्}\ H}}}{\lemma{\skt{तथैव}}\Afootnote{K; \skt{यशस्यम् ।}\ A; \skt{तथैअव ।}\ H}}\skt{ । }\pend
\pstart\versepar \skt{}\edtext{\skt{यक्षाणमाशु}}{\Afootnote{K; \skt{यक्ष्माणमाशु}\ A; \skt{यक्ष्माणमायुर्}\ H}}\skt{ }\edtext{\skt{व्यवहन्ति}}{\Afootnote{K; \skt{व्यपहन्ति}\ A H}}\skt{ }\edtext{\skt{चैव}}{\Afootnote{K H; \skt{चैत}\ A}}\skt{ }\edtext{\edtext{\skt{पाण्ड्वामयश्चापि}}{\lemma{\skt{पाण्ड्वामयश्}}\Afootnote{K; \skt{त्पाण्ड्वामयं}\ A; \skt{पाण्ड्वामयञ्}\ H}}}{\lemma{\skt{चापि}}\Afootnote{K H; \skt{चैव}\ A}}\skt{ }\edtext{\skt{भगन्दरश्च}}{\lemma{\skt{भगन्दरश्}}\Afootnote{K; \skt{भगन्दरं}\ A H}}\skt{ ।}\pend

\pstart \ledsidenote{\passref{6.41.54cd}}\skt{न चात्र किञ्चि\-त्प\-रि\-व\-र्ज\-नीयं रसायनं }\edtext{\skt{चैत\-दु\-पा\-स्य\-माणं}}{\Afootnote{K; \skt{चैत\-दु\-पा\-स्य\-मा\-नम्}\ A H}}\skt{ ।}\pend

\pstart\versepar \ledsidenote{\passref{6.41.58}}\edtext{\skt{स्त्रियो}}{\Afootnote{K H; \skt{शोकं स्त्रियं}\ A}}\skt{ }\edtext{\skt{रति\-क्रो\-ध\-म\-सू\-य\-नोश्च}}{\Afootnote{K; \skt{क्रोधमसूयनं च}\ A; \skt{रति\-क्रो\-ध\-म\-सू\-य\-ताञ्च}\ H}}\skt{ }\edtext{\skt{त्यजेदुदारां}}{\Afootnote{K; \skt{त्यजेदुदारा}\ A; \skt{त्यहेत्द्दुरां}\ H}}\skt{ }\edtext{\skt{विषयं}}{\Afootnote{K; \skt{न्विषया}\ A; \skt{विषया}\ H}}\skt{ भजेत् । }\pend
\pstart\versepar \skt{द्विजागुरूं }\edtext{\edtext{\edtext{\skt{तान्स्त्रिदशांश्च}}{\lemma{\skt{तान्स्}}\Afootnote{K; \skt{तीं}\ A; \skt{ता}\ H}}}{\lemma{\skt{त्रिदशांश्}}\Afootnote{K; \skt{स्त्रिदशा}\ A; \skt{ञ्स्त्रीदशाञ्श्}\ H}}}{\lemma{\skt{च}}\Afootnote{K H; \skt{न्गु}\ A}}\skt{ }\edtext{\skt{पूजेद्वाचश्च}}{\Afootnote{K H; \skt{रूश्च वाचश्च}\ A}}\skt{ पुण्याः }\edtext{\skt{शृणुयाद्द्विजेभ्य}}{\Afootnote{K H; \skt{शृणुयाद्बिजेभ्यः ॥}\ A}}\skt{ }\edtext{\skt{इति}}{\Afootnote{K H; \emph{om.}\ A}}\skt{ ॥}\edtext{}{\lemma{\skt{भजेत् । द्विजागुरूं}}\Afootnote{K; \skt{न्भजेत । वैद्यान्द्विजा}\ A; \skt{न्भजेत् । द्विजान्गुरूञ्स्}\ H}}\pend

\pstart \ledsidenote{\passref{6.41.trailer}}\edtext{\skt{काय}}{\Afootnote{K; \skt{इति सुश्रु\-त\-सं\-हि\-ता\-या\-मु\-त्त\-र\-त\-न्त्रा\-न्त\-र्गते काय}\ A; \skt{उत्तरतन्त्रे शोषविकित्सा अष्टत्रिंसत्तमः काय}\ H}}\skt{ }\edtext{\skt{चि}}{\Afootnote{K; \skt{चिकित्सातन्त्रे (तृतीयोऽध्यायः, आदित) एक\-च\-त्वा\-रिं\-शोऽ\-ध्यायः ॥४१ ॥}\ A; \skt{चिकित्सायां तृतीयोध्यायः ॥}\ H}}\skt{ ३८ ॥}\pend

\endnumbering
\beginnumbering
\eledsection*{Adhyāya 42}
\pstart \ledsidenote{\passref{6.42.1}}\skt{अथातो }\edtext{\skt{गुल्म\-रो\-ग\-प्र\-ति\-षेधं}}{\Afootnote{K; \skt{गुल्मप्रतिषेधं}\ A H}}\skt{ व्याख्यास्यामः ॥}\pend

\pstart\versepar \ledsidenote{\passref{6.42.3}}\skt{यथोक्तैः }\edtext{\skt{कोपनैदोषः}}{\Afootnote{K; \skt{कोपनैर्दोषाः}\ A; \skt{कोपनैर्दोषाः}\ H}}\skt{ कुपिताः कोष्ठमागताः । }\pend
\pstart\versepar \skt{जनयन्ति }\edtext{\skt{नृणाङ्गुल्मं स}}{\Afootnote{A K; \skt{नृणाञ्गुल्मं स}\ H}}\skt{ पञ्चविध उच्यते ।}\pend

\pstart\versepar \ledsidenote{\passref{6.42.4}}\edtext{\skt{हृन्नास्योरन्तरे}}{\lemma{\skt{हृन्नास्योर्}}\Afootnote{K H; \skt{हृद्बस्त्योर्}\ A}}\skt{ ग्रन्थिः }\edtext{\skt{साञ्चारी}}{\Afootnote{K; \skt{संचारी}\ A H}}\skt{ }\edtext{\skt{यदु}}{\Afootnote{K; \skt{यदि}\ A H}}\skt{ वाचलः । }\pend
\pstart\versepar \skt{}\edtext{\skt{वृत्तश्च}}{\Afootnote{K; \emph{om.}\ A; \skt{वृत्तश्}\ H}}\skt{ }\edtext{\skt{चयः}}{\Afootnote{K H; \skt{चया}\ A}}\skt{ }\edtext{\skt{सञ्चय}}{\Afootnote{K H; \skt{पचय}\ A}}\skt{ }\edtext{\skt{वां}}{\Afootnote{K; \skt{वा}\ A H}}\skt{ }\edtext{\skt{स}}{\Afootnote{K; \skt{न्वृत्तः स}\ A; \skt{न्स}\ H}}\skt{ गुल्म इति }\edtext{\skt{सञ्जितः}}{\Afootnote{K; \skt{कीर्तितः ॥}\ A; \skt{संज्ञितः ॥}\ H}}\skt{ ॥}\pend

\pstart \ledsidenote{\passref{6.42.5cd}}\edtext{\skt{गुपि\-ता\-नि\-ल\-मू\-ल\-त्वा\-द्गु\-ढ\-गु\-ल्मो\-द\-या\-दपि}}{\Afootnote{K H; \skt{गुपि\-ता\-नि\-ल\-मू\-ल\-त्वा\-द्गू\-ढ\-मू\-लो\-द\-या\-दपि ॥}\ A}}\skt{ ॥}\pend

\pstart\versepar \ledsidenote{\passref{6.42.6}}\edtext{\skt{गुल्मवद्धा}}{\Afootnote{K; \skt{गुल्मवद्वा}\ A H}}\skt{ विशालत्वाद्गुल्म }\edtext{\skt{इत्यभिधीयते}}{\lemma{\skt{इत्य्}}\Afootnote{A K; \skt{इ}\ H}}\skt{ ॥ }\pend
\pstart\versepar \skt{स यस्मादात्मनि चयं }\edtext{\edtext{\edtext{\skt{गच्छ\-त्य\-प्स्व\-वु\-द्धुदः}}{\lemma{\skt{गच्छत्य्}}\Afootnote{A K; \skt{गच्छ}\ H}}}{\lemma{\skt{अप्स्व्}}\Afootnote{A K; \skt{त्य}\ H}}}{\lemma{\skt{अवुद्धुदः}}\Afootnote{K; \skt{इव बुद्बुदः ॥}\ A; \skt{प्स्विव वुद्वुदः ।}\ H}}\skt{ ॥}\edtext{}{\lemma{\skt{अभिधीयते ॥ स}}\Afootnote{A K; \skt{त्यभिधीयते ॥ स}\ H}}\pend

\pstart\versepar \ledsidenote{\passref{6.42.7}}\edtext{\skt{अन्तश्चरति}}{\lemma{\skt{चरति}}\Afootnote{K H; \skt{सरति}\ A}}\skt{ }\edtext{\skt{यस्माच्च}}{\Afootnote{K; \skt{यस्माच्च न}\ A H}}\skt{ }\edtext{\skt{पाक\-मु\-प\-या\-त्यतः}}{\lemma{\skt{उपयात्य्}}\Afootnote{K; \skt{उपया}\ A; \skt{वुपया}\ H}}\skt{ । }\pend
\pstart\versepar \skt{स व्यस्तैर्जायते दोषैः समस्तैरपि चोच्छ्रितैः ।}\edtext{}{\lemma{\skt{अतः । स}}\Afootnote{K; \skt{त्यतः स}\ A H}}\pend

\pstart\versepar \ledsidenote{\passref{6.42.8}}\skt{पुरुषाणान्तथा स्त्रीणां ज्ञेयो रक्तेन }\edtext{\skt{चापरः}}{\Afootnote{A K; \skt{यो परः ॥}\ H}}\skt{ । }\pend
\pstart\versepar \skt{}\edtext{\skt{कृच्छ्र\-वि\-ण्मू\-त्र\-वा\-त\-त्व\-माना}}{\Afootnote{K H; \skt{सदनं मदन्ता वह्नेराटोपो}\ A}}\skt{ }\edtext{\skt{होत्रनिकूजनं /}}{\Afootnote{K H; \skt{ऽन्त्र\-वि\-कू\-ज\-नम् ॥}\ A}}\pend

\pstart \ledsidenote{\passref{6.42.9cd}}\edtext{\skt{द्वेसोन्ने च}}{\Afootnote{K; \skt{द्वेषोऽन्ने}\ A; \skt{द्वेषोन्ने}\ H}}\skt{ }\edtext{\skt{वायुरूर्ध्वश्च}}{\Afootnote{K; \skt{वायुरूर्ध्वं च}\ A; \skt{वायूरूर्ध्वश्च}\ H}}\skt{ पूर्वरूपेषु गुल्मिनां ।}\pend

\pstart\versepar \ledsidenote{\passref{6.42.10}}\edtext{\skt{हृत्कुक्षिशुलं}}{\Afootnote{K; \skt{हृत्कुक्षिशूलं}\ A; \skt{हृत्कूक्षिशूलं}\ H}}\skt{ }\edtext{\skt{मुखकण्ठशोषो}}{\Afootnote{A K; \skt{मूखख ण्ठशोषो}\ H}}\skt{ वायो\-र्नि\-रो\-णा\-वि\-ष\-मा\-ग्निता च । }\pend
\pstart\versepar \skt{ते ते विकाराः }\edtext{\skt{पवनात्मिकाश्च}}{\Afootnote{K H; \skt{पवनात्मकाश्च}\ A}}\skt{ भवन्ति गुल्मे }\edtext{\skt{निलसंभवे}}{\Afootnote{K; \skt{ऽनिलसंभवे}\ A; \skt{निलसम्भवे}\ H}}\skt{ तु ।}\edtext{}{\lemma{\skt{वायो\-र्नि\-रो\-णा\-वि\-ष\-मा\-ग्निता च । ते ते}}\Afootnote{K; \skt{वायोर्निरोधो विषमाग्निता च । ते ते}\ A; \skt{वायोर्न्निरोधो विषमाग्निता च । ते ते}\ H}}\pend

\pstart\versepar \ledsidenote{\passref{6.42.11}}\skt{स्वेद\-ज्व\-रा\-हा\-र\-वि\-दा\-ह\-दा\-हा\-स्तृ\-ष्णा\-ङ्ग\-रागः }\edtext{\skt{कटुवक्रता च}}{\Afootnote{K; \skt{कटुवक्त्रता च ।}\ A; \skt{क टूवक्रता च ।}\ H}}\skt{ । }\pend
\pstart\versepar \skt{पित्तस्य }\edtext{\skt{लिङ्गानि च}}{\Afootnote{K H; \skt{लिङ्गान्यखिलानि}\ A}}\skt{ यानि }\edtext{\skt{तानि}}{\Afootnote{K H; \skt{पित्तात्मके तानि}\ A}}\skt{ भवन्ति }\edtext{\skt{पित्तप्रभवे तु}}{\Afootnote{K H; \emph{om.}\ A}}\skt{ गुल्मे ॥}\pend

\pstart\versepar \ledsidenote{\passref{6.42.12}}\edtext{\skt{अन्ने}}{\Afootnote{K H; \skt{स्तैमित्यमन्ने}\ A}}\skt{ }\edtext{\skt{न}}{\Afootnote{K H; \skt{ऽ}\ A}}\skt{ }\edtext{\skt{लिप्सा}}{\Afootnote{K H; \emph{om.}\ A}}\skt{ }\edtext{\skt{गुरूगात्रता च}}{\Afootnote{K H; \skt{रुचि\-र\-ङ्ग\-सा\-दश्}\ A}}\skt{ }\edtext{\skt{छर्दिप्रसेको}}{\Afootnote{K H; \skt{छर्दिः प्रसेको}\ A}}\skt{ मधुरास्य ता च । }\pend
\pstart\versepar \skt{}\edtext{\skt{कफस्य}}{\Afootnote{A K; \skt{कफ प्र}\ H}}\skt{ लिङ्गानि च यानि तानि भवन्ति गुल्मे }\edtext{\skt{कफसम्भवे}}{\Afootnote{K H; \skt{कफसंभवे}\ A}}\skt{ }\edtext{\skt{च}}{\Afootnote{K H; \skt{तु ॥}\ A}}\skt{ ।}\pend

\pstart\versepar \ledsidenote{\passref{6.42.13}}\edtext{\skt{सर्वात्माकः}}{\Afootnote{K; \skt{सर्वात्मकः}\ A H}}\skt{ }\edtext{\edtext{\skt{सर्व\-रू\-जो\-प\-प\-न्न\-स्त\-श्चाप्य}}{\lemma{\skt{सर्व\-रू\-जो\-प\-प\-न्न\-स्तश्}}\Afootnote{K; \skt{सर्व\-वि\-का\-र\-युक्तः सो}\ A; \skt{सर्व्व\-रू\-जो\-प\-प\-न्नस्तं}\ H}}}{\lemma{\skt{चाप्य}}\Afootnote{K; \skt{ऽसाध्य}\ A; \skt{याप्य}\ H}}\skt{ }\edtext{\edtext{\skt{सध्यम्प्रवदन्ति}}{\lemma{\skt{सध्यम्}}\Afootnote{K; \skt{उक्तः}\ A; \skt{साध्याम्}\ H}}}{\lemma{\skt{प्रवदन्ति}}\Afootnote{K H; \skt{क्षतजं प्र}\ A}}\skt{ }\edtext{\skt{तज्ज्ञाः}}{\Afootnote{K H; \skt{वक्ष्ये ।}\ A}}\skt{ । }\pend
\pstart\versepar \skt{}\edtext{\skt{नव\-प्र\-सु\-ता\-ता\-हि\-त\-भो\-जना या या}}{\Afootnote{K; \skt{नव\-प्र\-सू\-ता\-हि\-त\-भो\-जना या या}\ A H}}\skt{ }\edtext{\edtext{\skt{चाम\-ग\-र्भ\-म्वि\-सृ\-जो\-दृतै}}{\lemma{\skt{चामगर्भम्}}\Afootnote{A K; \skt{चाम लि ग र्भम्}\ H}}}{\lemma{\skt{विसृजोदृतै}}\Afootnote{K; \skt{विसृजेदृतौ}\ A H}}\skt{ वा ॥}\pend

\pstart\versepar \ledsidenote{\passref{6.42.14}}\edtext{\edtext{\skt{वायुर्हितस्यो}}{\lemma{\skt{वायुर्}}\Afootnote{A K; \skt{वायु}\ H}}}{\lemma{\skt{हितस्यो}}\Afootnote{K; \skt{हि तस्याः}\ A; \skt{र्हि तस्याः}\ H}}\skt{ }\edtext{\skt{विगृह्य}}{\Afootnote{K; \skt{परिगृह्य}\ A; \skt{प्रविगृह्य}\ H}}\skt{ }\edtext{\skt{रक्तां}}{\Afootnote{K; \skt{रक्तं}\ A H}}\skt{ करोति गुल्मं }\edtext{\skt{सरूजां}}{\Afootnote{K; \skt{सरुजं}\ A; \skt{सरूजं}\ H}}\skt{ सदाहं । }\pend
\pstart\versepar \skt{}\edtext{\skt{पित्तस्य}}{\Afootnote{K H; \skt{पैत्तस्य}\ A}}\skt{ लिङ्गेन समानलिङ्गं विशेषणं }\edtext{\edtext{\skt{चाप्य\-प\-र\-न्नि\-वोधः}}{\lemma{\skt{चाप्यपरन्}}\Afootnote{K; \skt{चाप्यपरं}\ A; \skt{चापरमं}\ H}}}{\lemma{\skt{निवोधः}}\Afootnote{K; \skt{निबोध ॥}\ A; \skt{निवोध ।}\ H}}\skt{ ।}\pend

\pstart\versepar \ledsidenote{\passref{6.42.15}}\skt{न स्पन्दते नोदरमेति वृद्धिम्भवन्ति लिङ्गानि च गर्भिनीनां । }\pend
\pstart\versepar \skt{तं }\edtext{\skt{रक्तगुल्मं}}{\Afootnote{H; \skt{गर्भकाला}\ A; \emph{om.}\ K}}\skt{ }\edtext{\skt{प्रवदन्ति}}{\Afootnote{H; \skt{तिगमे चैकि\-त्स्य\-म\-सृ\-ग्भवं गुल्ममुशन्ति}\ A; \skt{रक्तगुल्मं प्रवद}\ K}}\skt{ }\edtext{\skt{तज्ज्ञाः स}}{\Afootnote{H; \skt{तज्ज्ञाः ॥}\ A; \emph{om.}\ K}}\skt{ }\edtext{\skt{चाप्य}}{\Afootnote{H; \emph{om.}\ A K}}\skt{ }\edtext{\skt{साध्या तु}}{\Afootnote{H; \emph{om.}\ A K}}\skt{ }\edtext{\skt{पितानिलत्वात्}}{\Afootnote{H; \emph{om.}\ A K}}\skt{ ।}\edtext{}{\lemma{\skt{गर्भिनीनां । तं}}\Afootnote{K; \skt{गर्भिणीनाम्तं}\ A H}}\pend

\pstart\versepar \ledsidenote{\passref{6.42.16}}\skt{वातगुल्मार्द्दितं स्निग्धं युक्तं स्नेहविरेचनैः । }\pend
\pstart\versepar \skt{उपा\-च\-रे\-द्य\-था\-कालं निरूहैः सानुवासनैः ॥}\pend

\pstart\versepar \ledsidenote{\passref{6.42.16.add1}}\skt{जीर्णासने कुप्यति वातगुल्मे विदा\-ह\-का\-ले\-त्वे\-श\-नस्य पित्तात् ॥ }\pend
\pstart\versepar \skt{गुल्मः कफात्मूर्छति रुकमात्रे सर्वासु वलासु च सन्निपाते ।}\pend

\pstart\versepar \ledsidenote{\passref{6.42.16.add2}}\skt{स्त्रीणा तु यः शोणि\-त\-को\-ष\-गुल्मः स योनिशूलं तु मूलं करोति ॥ }\pend
\pstart\versepar \skt{योन्याः परि\-स्रा\-व\-म\-नि\-ष्त\-ग\-न्ध\-म्वि\-शे\-ष\-लि\-ङ्गानि भवन्ति चैषां ॥}\pend

\pstart\versepar \ledsidenote{\passref{6.42.17}}\skt{पित्तगुल्मार्द्दितं }\edtext{\skt{पीतं}}{\Afootnote{H; \skt{स्निग्धं}\ A}}\skt{ }\edtext{\skt{काकोल्यादिकृतं}}{\Afootnote{H; \skt{काकोल्यादिघृते}\ A}}\skt{ }\edtext{\skt{घृतं}}{\Afootnote{H; \skt{न तु ।}\ A}}\skt{ ॥ }\pend
\pstart\versepar \skt{}\edtext{\edtext{\skt{विरि\-क्त\-म्म\-धु\-रै\-यो\-है\-र्न्नि\-रूहैः}}{\lemma{\skt{मधुरैयोहैर्}}\Afootnote{H; \skt{मधुरैर्योगै}\ A}}}{\lemma{\skt{न्निरूहैः}}\Afootnote{H; \skt{र्निरूहैः}\ A}}\skt{ समुपाचरेत् ॥}\pend

\pstart\versepar \ledsidenote{\passref{6.42.18}}\skt{श्लेष्मगुल्मार्द्दितं स्निग्धं }\edtext{\skt{पिप्पल्यादिघृतं न तु}}{\Afootnote{H; \skt{पिप्प\-ल्या\-दि\-घृ\-तेन तु ।}\ A}}\skt{ । }\pend
\pstart\versepar \skt{}\edtext{\edtext{\edtext{\edtext{\skt{तीक्षै\-र्व्वि\-रि\-क्त\-न्त\-द्रू\-पै\-र्न्नि\-रूहैः}}{\lemma{\skt{तीक्षैर्}}\Afootnote{H; \skt{तीक्ष्णै}\ A}}}{\lemma{\skt{व्विरिक्तन्}}\Afootnote{H; \skt{र्विरिक्तं}\ A}}}{\lemma{\skt{तद्रूपैर्}}\Afootnote{H; \skt{तद्रूपै}\ A}}}{\lemma{\skt{न्निरूहैः}}\Afootnote{H; \skt{र्निरूहैः}\ A}}\skt{ समुपाचरेत् ॥}\pend

\pstart\versepar \ledsidenote{\passref{6.42.19}}\edtext{\skt{सन्निपातोत्थिते}}{\Afootnote{A H; \emph{om.}\ K}}\skt{ }\edtext{\skt{गुल्मे}}{\Afootnote{A H; \emph{om.}\ K}}\skt{ }\edtext{\skt{त्रिदोषघ्नो}}{\Afootnote{A H; \emph{om.}\ K}}\skt{ }\edtext{\skt{विधिर्हितः}}{\Afootnote{A H; \emph{om.}\ K}}\skt{ । }\pend
\pstart\versepar \skt{}\edtext{\skt{पित्त\-व\-द्र\-क्त\-गु\-ल्मिन्या}}{\Afootnote{A H; \emph{om.}\ K}}\skt{ }\edtext{\skt{नार्याः}}{\Afootnote{A H; \emph{om.}\ K}}\skt{ }\edtext{\skt{कार्यः}}{\Afootnote{A H; \skt{नार्याः}\ K}}\skt{ क्रियाविधिः ॥}\pend

\pstart\versepar \ledsidenote{\passref{6.42.20}}\edtext{\edtext{\skt{विशेषमपरं}}{\lemma{\skt{विशेषम्}}\Afootnote{K; \skt{विशेष}\ A H}}}{\lemma{\skt{अपरं}}\Afootnote{K; \skt{मपरं}\ A H}}\skt{ }\edtext{\skt{चायिः}}{\Afootnote{K; \skt{चास्याः}\ A H}}\skt{ शृणु }\edtext{\skt{रक्तप्रभेदनं}}{\Afootnote{K H; \skt{रक्त\-वि\-भे\-द\-नम् ।}\ A}}\skt{ । }\pend
\pstart\versepar \skt{}\edtext{\skt{पला\-श\-क्षा\-र\-तो\-येन}}{\Afootnote{A K; \skt{पलाशक्षोर तोयेन}\ H}}\skt{ सिद्धं सर्पिः प्रयोजयेत् ।}\pend

\pstart \ledsidenote{\passref{6.42.21cd}}\skt{उष्णैर्वा भेदयेत्भिन्ने }\edtext{\edtext{\skt{विधिरा}}{\lemma{\skt{विधिर्}}\Afootnote{H; \skt{विधि}\ A K}}\skt{ मृग्दरो}}{\lemma{\skt{आ मृग्दरो}}\Afootnote{\skt{रासृग्दरो}\ A; \skt{रामृग्दरो}\ K; \skt{आसृग्दरो}\ H}}\skt{ हितः ।}\pend

\pstart\versepar \ledsidenote{\passref{6.42.22}}\skt{आनू\-पौ\-द\-क\-म\-ज्जानो }\edtext{\skt{वसा\-तै\-ल\-घृ\-तानि च}}{\Afootnote{K; \skt{वसा तैलं घृतं दधि ।}\ A; \skt{वसातैलं घृतानि च ॥}\ H}}\skt{ । }\pend
\pstart\versepar \skt{}\edtext{\skt{विपचेदेकतः}}{\Afootnote{K H; \skt{विपक्वमेकतः}\ A}}\skt{ शस्तं }\edtext{\skt{वात\-गु\-ल्मेऽ\-नु\-वा\-सनं}}{\lemma{\skt{ऽनुवासनं}}\Afootnote{A H; \skt{नुवासनं ।}\ K}}\skt{ ।}\pend

\pstart\versepar \ledsidenote{\passref{6.42.23}}\edtext{\skt{जाङ्ग\-लै\-क\-श\-फा\-नांञ्च}}{\Afootnote{K; \skt{जाङ्ग\-लै\-क\-श\-फानां तु}\ A; \skt{जाङ्ग\-लै\-क\-श\-फा\-नाञ्च}\ H}}\skt{ }\edtext{\skt{वसासर्पिश्च}}{\Afootnote{A K; \skt{वसा सर्प्पिञ्च}\ H}}\skt{ }\edtext{\skt{पित्तजा}}{\Afootnote{K; \skt{पैत्तिके ।}\ A; \skt{पित्तजे ॥}\ H}}\skt{ । }\pend
\pstart\versepar \skt{तैलं }\edtext{\edtext{\skt{जाङ्ग\-ल\-म\-ज्जा\-न\-मेवं}}{\lemma{\skt{जाङ्ग\-ल\-म\-ज्जा\-नम्}}\Afootnote{K H; \skt{जाङ्गलमज्जान}\ A}}}{\lemma{\skt{एवं}}\Afootnote{A K; \skt{एव}\ H}}\skt{ }\edtext{\skt{गुल्मैकफोत्थिते}}{\Afootnote{K H; \skt{गुल्मे कफोत्थिते ॥}\ A}}\skt{ ॥}\pend

\pstart\versepar \ledsidenote{\passref{6.42.24}}\edtext{\skt{धात्री\-फ\-ला\-ना\-स्व\-रसे}}{\Afootnote{K; \skt{धात्रीफलानां स्वरसे}\ A H}}\skt{ }\edtext{\skt{षडङ्गे}}{\Afootnote{K; \skt{षडङ्गं}\ A; \skt{षडङ्ग}\ H}}\skt{ }\edtext{\skt{विपचेद्धृतं}}{\Afootnote{K; \skt{विपचेद्घृतम्}\ A H}}\skt{ । }\pend
\pstart\versepar \skt{}\edtext{\skt{शर्क\-रा\-सै\-न्ध\-वो\-पेतं}}{\Afootnote{A K; \skt{शर्क्क\-रा\-स\-न्धै\-वो\-पेतं}\ H}}\skt{ }\edtext{\skt{तधितं}}{\Afootnote{K H; \skt{तद्धितं}\ A}}\skt{ }\edtext{\skt{वातगुल्मिनां}}{\Afootnote{K H; \skt{वातगुल्मिने ॥}\ A}}\skt{ ॥}\pend

\pstart\versepar \ledsidenote{\passref{6.42.34}}\edtext{\skt{कुल\-त्थ\-मू\-ल\-क\-क्वा\-थ\-द\-धि\-का\-ञ्जि\-क\-स\-र्पिषां}}{\Afootnote{K H; \skt{दधि सौवीरकं सर्पिः क्वाथौ मुद्गकुलत्थजौ ।}\ A}}\skt{ । }\pend
\pstart\versepar \skt{पञ्चाढकानि विपचेदावाप्य }\edtext{\skt{द्विपलान्य}}{\Afootnote{K; \skt{द्विपलान्यथ ॥}\ A; \skt{द्विपलानि तु ॥}\ H}}\skt{ ।}\pend

\pstart\versepar \ledsidenote{\passref{6.42.35}}\edtext{\skt{सौर्वर्चलं}}{\Afootnote{K; \skt{सौवर्चलं}\ A H}}\skt{ }\edtext{\skt{वितं}}{\Afootnote{K H; \skt{सर्जि}\ A}}\skt{ }\edtext{\skt{हिंगुदेवदाथ}}{\Afootnote{K; \skt{कां च देवदार्वथ}\ A; \skt{हिंगु\-दे\-व\-दा\-र्व्वथ}\ H}}\skt{ सैन्धवं । }\pend
\pstart\versepar \skt{वातगुल्मापहं }\edtext{\skt{सर्पि\-रे\-त\-द्दी\-प\-न\-मेव}}{\lemma{\skt{एतद्दीपनम्}}\Afootnote{A K; \skt{एतद्वीपणम्}\ H}}\skt{ च ॥}\pend

\pstart \ledsidenote{\passref{6.42.37cd}}\edtext{\skt{आगग्वधादौ}}{\Afootnote{K; \skt{आरग्वधादौ}\ A H}}\skt{ }\edtext{\skt{विप\-चे\-द्दी\-प\-नी\-य\-युतं}}{\Afootnote{A K; \skt{विप\-चे\-द्दी\-प\-नी\-य\-युतं}\ H}}\skt{ घृतं ।}\pend

\pstart\versepar \ledsidenote{\passref{6.42.38}}\skt{क्षारवर्गे }\edtext{\skt{पचेद्चान्य{\symfont ¦}}}{\Afootnote{K; \skt{पचेच्चान्य}\ A; \skt{पचेद्चान्यं}\ H}}\skt{ }\edtext{\skt{पचे\-न्मू\-त्र\-ग\-णे\-परं}}{\Afootnote{K H; \skt{त्पचे\-न्मू\-त्र\-ग\-णेऽ\-प\-रम् ।}\ A}}\skt{ । }\pend
\pstart\versepar \skt{घ्नन्ति गुल्मं कफोद्भूतं }\edtext{\edtext{\edtext{\skt{घृता\-न्ये\-ता\-न्य\-सं\-शयं}}{\lemma{\skt{घृतान्य्}}\Afootnote{K H; \skt{घृता}\ A}}}{\lemma{\skt{एतान्य्}}\Afootnote{K H; \skt{न्येता}\ A}}}{\lemma{\skt{असंशयं}}\Afootnote{K H; \skt{न्यसंशयम् ॥}\ A}}\skt{ ।}\pend

\pstart\versepar \ledsidenote{\passref{6.42.36}}\edtext{\skt{तृण\-मू\-ली\-क\-षाये तु}}{\Afootnote{K H; \skt{तृण\-मू\-ल\-क\-षाये तु}\ A}}\skt{ जीवनीयैः }\edtext{\skt{पचेद्घृतं}}{\Afootnote{K H; \skt{पचेद्घृतम् ।}\ A}}\skt{ । }\pend
\pstart\versepar \skt{}\edtext{\skt{न्यग्रो\-धा\-दि\-क\-षा\-येण}}{\Afootnote{K H; \skt{न्यग्रोधादिगणे वापि}\ A}}\skt{ }\edtext{\skt{घृतं}}{\Afootnote{K; \skt{गणे}\ A; \skt{घृते}\ H}}\skt{ }\edtext{\skt{वाप्यपलादिके}}{\Afootnote{K H; \skt{वाप्युत्पलादिके ॥}\ A}}\skt{ ।}\pend

\pstart \ledsidenote{\passref{6.42.37ab}}\skt{रक्तपित्तोत्थितं }\edtext{\skt{हन्ति}}{\Afootnote{K H; \skt{घ्नन्ति}\ A}}\skt{ }\edtext{\edtext{\edtext{\skt{घृता\-न्ये\-ता\-न्य\-सं\-शयं}}{\lemma{\skt{घृतान्य्}}\Afootnote{A K; \skt{घृ}\ H}}}{\lemma{\skt{एतान्य्}}\Afootnote{A K; \skt{ता}\ H}}}{\lemma{\skt{असंशयं}}\Afootnote{A K; \skt{न्यसंशयं ।}\ H}}\skt{ ।}\pend

\pstart\versepar \ledsidenote{\passref{6.42.39}}\skt{यथा }\edtext{\skt{दोषाच्छ्रायञ्चापि}}{\lemma{\skt{दोषाच्छ्रायञ्}}\Afootnote{K H; \skt{दोषोच्छ्रयं}\ A}}\skt{ चिकि\-त्से\-त्सा\-न्नि\-पा\-तिकं । }\pend
\pstart\versepar \skt{चूर्णं }\edtext{\edtext{\skt{हिग्वादिकञ्चापि}}{\lemma{\skt{हिग्वादिकञ्}}\Afootnote{K H; \skt{हिङ्ग्वादिकं}\ A}}}{\lemma{\skt{चापि}}\Afootnote{K H; \skt{वापि}\ A}}\skt{ घृतम्वा प्लीहनाशनं ।}\pend

\pstart \ledsidenote{\passref{6.42.40ab}}\edtext{\skt{पिबेद्गुल्मापहं}}{\Afootnote{A; \skt{पिवेद्गुल्मापहं}\ K H}}\skt{ काले }\edtext{\edtext{\skt{सर्पि\-श्तै\-ल्व\-क\-मेव}}{\lemma{\skt{सर्पिश्}}\Afootnote{K H; \skt{सर्पि}\ A}}}{\lemma{\skt{तैल्वकम्}}\Afootnote{K H; \skt{स्तैल्वकम्}\ A}}\skt{ वा ॥}\pend

\pstart\versepar \ledsidenote{\passref{6.42.40.add1}}\skt{पाठा स च व्यावृ\-ह\-ती\-द्व\-यञ्च }\edtext{\skt{रमूलेन}}{\Afootnote{K; \skt{एव शुमूलेन}\ H}}\skt{ }\edtext{\skt{समानि }\edtext{\skt{श्च}}{\lemma{\skt{च}}\Afootnote{K; \skt{पञ्च ।}\ H}}}{\lemma{\skt{समानि श्}}\Afootnote{K; \skt{समानि}\ H}}\skt{ । }\pend
\pstart\versepar \skt{संच्छि\-न्द्य\-सं\-चू\-र्ध्याथ धारयित्वा जालादकं तैः सह पाचनीयं ।}\pend

\pstart\versepar \ledsidenote{\passref{6.42.40.add2}}\edtext{\skt{अथा\-व\-ता\-वा\-व\-शेष}}{\Afootnote{K; \skt{अथावताव्यं वावशेष}\ H}}\skt{ }\edtext{\skt{मत्राक्षमात्रं}}{\Afootnote{K; \skt{मात्राक्षमात्रं}\ H}}\skt{ तृवृतं विय । }\pend
\pstart\versepar \skt{स नागरं गुल्मवधाययें }\edtext{\skt{तद्वातगुल्मेषु}}{\Afootnote{K; \skt{तद्वातगुल्मेषु}\ H}}\skt{ विरेचनं स्यात् ।}\edtext{}{\lemma{\skt{विय । स}}\Afootnote{K; \skt{विधाय । स}\ H}}\pend

\pstart\versepar \ledsidenote{\passref{6.42.40.add3}}\skt{दन्ती\-व\-चा\-हिं\-ग्व\-ज\-मो\-द\-पाण क्षारं यवानां }\edtext{\skt{लवणत्रिकश्च}}{\Afootnote{K; \skt{लवणत्रिकञ्च ।}\ H}}\skt{ । }\pend
\pstart\versepar \skt{सचित्रकं }\edtext{\skt{क्षीर\-क\-शृ\-द्ग\-वेरं}}{\Afootnote{K; \skt{जीरकवेरं}\ H}}\skt{ }\edtext{\skt{स्यान्मातुलं}}{\Afootnote{K; \skt{स्यान्मातुलुङ्}\ H}}\skt{ }\edtext{\skt{गम्य}}{\Afootnote{K; \skt{गस्य}\ H}}\skt{ }\edtext{\skt{भवस्य}}{\Afootnote{K; \skt{च}\ H}}\skt{ }\edtext{\skt{मूलं}}{\Afootnote{K; \skt{मूलकञ्च ॥}\ H}}\skt{ ।}\pend

\pstart\versepar \ledsidenote{\passref{6.42.40.add4}}\skt{एषां समानं }\edtext{\skt{स्वश\-ये\-षि\-मा\-नां\-सु\-खा\-म्वुना}}{\Afootnote{K; \skt{स्वशुपेषितानां सुखाम्वुना}\ H}}\skt{ । }\pend
\pstart\versepar \skt{}\edtext{\skt{नरो}}{\Afootnote{K; \skt{प्रातरूपोस्य मात्रा । नरो}\ H}}\skt{ }\edtext{\skt{भृषा\-म\-यो\-द\-र\-प्र\-वृद्वो}}{\Afootnote{K; \skt{भृषा\-म\-प्य\-द\-र\-प्र\-वृद्वां}\ H}}\skt{ सद्योरूजां वातकृतां निहन्ति ।}\pend

\pstart\versepar \ledsidenote{\passref{6.42.40.add5}}\skt{द्राक्षापयस्या मधुकञ्च }\edtext{\skt{मुख्यं}}{\Afootnote{K; \skt{मूख्यं}\ H}}\skt{ सचन्दनं }\edtext{\edtext{\skt{पद्म\-क\-मु\-त्प\-लञ्च}}{\lemma{\skt{पद्मकम्}}\Afootnote{K; \skt{पद्मक}\ H}}}{\lemma{\skt{उत्पलञ्च}}\Afootnote{K; \skt{मूत्पलञ्च ।}\ H}}\skt{ । }\pend
\pstart\versepar \skt{क्षौद्रा\-न्वि\-त\-न्डु\-ल\-धा\-व\-नेन पीत्वा नुदेत्पित्तकृतं तु गुल्मं ।}\pend

\pstart\versepar \ledsidenote{\passref{6.42.40.add6}}\edtext{\skt{भागीपोटोलं}}{\Afootnote{K; \skt{भाग्नीपटोलीं}\ H}}\skt{ मधुकञ्च पाठां }\edtext{\skt{रोहि\-ल्य\-त्थै\-ला\-सह}}{\Afootnote{K; \skt{रोहिल्यत्थौला सह}\ H}}\skt{ पिप्पलीभिः । }\pend
\pstart\versepar \skt{एतानि पिष्ठ्वा मधुना }\edtext{\skt{वलिह}}{\Afootnote{K; \skt{वलि}\ H}}\skt{ पीतं }\edtext{\skt{समुति}}{\Afootnote{K; \skt{शमं नाशति}\ H}}\skt{ पित्तगुल्मं ।}\pend

\pstart\versepar \ledsidenote{\passref{6.42.40.add7}}\edtext{\skt{हरी\-त\-का\-त्या\-म\-लकं}}{\Afootnote{K; \skt{हरी\-त\-का\-न्या\-म\-लम्}\ H}}\skt{ }\edtext{\skt{वितंगं}}{\Afootnote{K; \skt{विडङ्गं}\ H}}\skt{ विश्वासधं चेति समानि कृत्वा । }\pend
\pstart\versepar \skt{क्षौद्रान्वितं स त्रिफला कषायं यः श्लेष्मगुल्मं प्रपिवन्ति हन्यात् ॥}\pend

\pstart\versepar \ledsidenote{\passref{6.42.40.add8}}\skt{च व्या घृतं }\edtext{\skt{तत्वृता}}{\Afootnote{K; \skt{त्रिव्टा}\ H}}\skt{ }\edtext{\skt{त्रितीयं}}{\Afootnote{K; \skt{त्रिईयन्}\ H}}\skt{ तथा }\edtext{\skt{चतुर्थी}}{\Afootnote{K; \skt{चतुर्था}\ H}}\skt{ गजपिप्पलीञ्च । }\pend
\pstart\versepar \skt{फलत्रयाणां सह पिप्पलीभिः पचेत्समुस्ता सलिलाढकस्थं ॥}\pend

\pstart\versepar \ledsidenote{\passref{6.42.40.add9}}\skt{मृद्वग्निसिद्धं क्षतवावशेषं पूतं }\edtext{\skt{पिबेद्क्षौद्रयूतं}}{\Afootnote{\skt{पिवेद्क्षौद्रयूतं}\ K; \skt{पिवेत्क्षौद्रयूतं}\ H}}\skt{ हिताय । }\pend
\pstart\versepar \skt{}\edtext{\skt{तत्श्लेष्मं}}{\Afootnote{K; \skt{तत्श्लेष्मगुल्मं}\ H}}\skt{ त्वथ वातगुल्मं }\edtext{\edtext{\edtext{\skt{निर्वा\-प\-ये\-थ्या\-ग्नि\-मि\-वा\-म्वु\-वेगः}}{\lemma{\skt{निर्वापयेथ्य्}}\Afootnote{K; \skt{निर्वापयेथ्य}\ H}}}{\lemma{\skt{आग्निम्}}\Afootnote{K; \skt{ग्नि}\ H}}}{\lemma{\skt{इवाम्वुवेगः}}\Afootnote{K; \skt{मिवाम्वुवेगः ॥}\ H}}\skt{ ।}\pend

\pstart\versepar \ledsidenote{\passref{6.42.40.add10}}\skt{सदेवकाष्ठं चिरिविल्वजं त्वक्सपिप्पली }\edtext{\skt{पिप्प\-लि\-मू\-ल\-भागौ}}{\Afootnote{K; \skt{पिप्प\-लि\-मू\-ल\-भाग्नी ।}\ H}}\skt{ । }\pend
\pstart\versepar \skt{पीत्वा }\edtext{\skt{नलक्वाथयुतं तु}}{\Afootnote{K H\textsuperscript{pc}; \skt{नलक्वाथयुतं}\ H\textsuperscript{ac}}}\skt{ वारि प्रमुच्यते }\edtext{\skt{शोणि\-त\-गु\-ल्म\-मू\-लात्}}{\Afootnote{K; \skt{शोणि\-अ\-पि\-त्त\-गु\-ल्मात् ॥}\ H}}\skt{ ।}\pend

\pstart\versepar \ledsidenote{\passref{6.42.40.add11}}\edtext{\skt{घोषञ्चभाईं च}}{\lemma{\skt{चभाईं च}}\Afootnote{K; \skt{च भाग्जीं च}\ H}}\skt{ }\edtext{\skt{गुरूप्रगाढां}}{\Afootnote{K; \skt{गुडप्रगाढां}\ H}}\skt{ }\edtext{\skt{नारीं}}{\Afootnote{K; \skt{नारी}\ H}}\skt{ }\edtext{\skt{पिबेत्मामृदितां}}{\Afootnote{\skt{पिवेत्मामृदितां}\ K H}}\skt{ घृतेन । }\pend
\pstart\versepar \skt{सुखाम्वुना }\edtext{\edtext{\skt{हन्ह्यथ}}{\lemma{\skt{हन्ह्य्}}\Afootnote{K; \skt{हन्ह्य}\ H}}}{\lemma{\skt{अथ}}\Afootnote{K; \skt{थ}\ H}}\skt{ रक्तगुल्मं नष्टञ्च पुष्पं }\edtext{\edtext{\skt{पुनरादधाति}}{\lemma{\skt{पुनर्}}\Afootnote{K; \skt{पुन}\ H}}}{\lemma{\skt{आदधाति}}\Afootnote{K; \skt{रादधाति ॥}\ H}}\skt{ ॥}\pend

\pstart\versepar \ledsidenote{\passref{6.42.40.add12}}\edtext{\skt{समु\-ष्क\-क\-ञ्चे\-क्ष\-रकं}}{\lemma{\skt{चेक्षरकं}}\Afootnote{K; \skt{चेक्षरसं}\ H}}\skt{ }\edtext{\skt{स्वदंष्ट्रांशोकं}}{\Afootnote{K; \skt{स्वदंट्रांशोकम्}\ H}}\skt{ }\edtext{\skt{पलासं}}{\Afootnote{K; \skt{पलाशेन्}\ H}}\skt{ }\edtext{\edtext{\skt{तिलनालमेव}}{\lemma{\skt{तिलनालम्}}\Afootnote{K; \skt{तिलनाल}\ H}}}{\lemma{\skt{एव}}\Afootnote{K; \skt{मेव ।}\ H}}\skt{ । }\pend
\pstart\versepar \skt{नीपक्वदम्वञ्च }\edtext{\skt{समातुलुंगं}}{\Afootnote{K; \skt{समातुलुङ्ग}\ H}}\skt{ पूतीकरञ्च }\edtext{\skt{सकपित्थ विल्वं}}{\Afootnote{H; \skt{सक\-पि\-त्थ\-वि\-विल्वं ।}\ K}}\skt{ ।}\pend

\pstart\versepar \ledsidenote{\passref{6.42.40.add13}}\edtext{\skt{वल्य\-न्ते\-मा\-म\-प्य\-मृतां}}{\lemma{\skt{वल्यन्तेमाम्}}\Afootnote{K; \skt{वल्लूत्तमाम्}\ H}}\skt{ च वल्लीं वृक्षादनीं }\edtext{\skt{जातिमथैंगुञ्च}}{\lemma{\skt{अथैंगुञ्च}}\Afootnote{K; \skt{अत्थैङ्गुदीञ्च ।}\ H}}\skt{ । }\pend
\pstart\versepar \skt{एताः प्रताना }\edtext{\skt{क्वथयाश्व}}{\Afootnote{K; \skt{क्वथयादश्व}\ H}}\skt{ प्रदा\-ह\-य\-न्क्षा\-र\-म\-था\-द\-दीत ।}\pend

\pstart\versepar \ledsidenote{\passref{6.42.40.add14}}\skt{पृथ\-क्पृ\-थ\-क्का\-ढ\-क\-मा\-न\-युक्तं }\edtext{\skt{तत्स्रु\-त\-मा\-द\-दीत}}{\lemma{\skt{तत्स्रुतम्}}\Afootnote{K; \skt{तत्मपुकृत्वः स्रुतम्}\ H}}\skt{ । }\pend
\pstart\versepar \skt{सुखा\-म्वु\-ना\-वे\-श\-मनि सन्निपाते }\edtext{\skt{कृतोयवस्त्राणि}}{\Afootnote{K; \skt{कृतोयवसानि}\ H}}\skt{ यतः क्षमावान् ।}\pend

\pstart\versepar \ledsidenote{\passref{6.42.40.add15}}\skt{अष्टाढकं }\edtext{\skt{तन्प्रसितं}}{\Afootnote{K; \skt{तत्प्रसितम्}\ H}}\skt{ }\edtext{\skt{विदध्यान्मूलन्तु}}{\Afootnote{K; \skt{विदध्यात्मूलं तु}\ H}}\skt{ }\edtext{\skt{माम्वाप्यथ}}{\Afootnote{K; \skt{लाम्वाप्यथ}\ H}}\skt{ चित्रकस्य । }\pend
\pstart\versepar \skt{सुकल्पितां तां }\edtext{\skt{विपचेद्विधिक्षाः}}{\Afootnote{K; \skt{विपचेद्विधिक्षाः}\ H}}\skt{ }\edtext{\skt{पूत\-वि\-द\-ध्या\-द\-थ\-या\-द\-लेष}}{\Afootnote{K; \skt{पूत\-म्वि\-ध्या\-द\-थ\-या\-द\-शेषं ॥}\ H}}\skt{ ।}\pend

\pstart\versepar \ledsidenote{\passref{6.42.40.add16}}\edtext{\skt{पलारि}}{\Afootnote{K; \skt{पलानि चत्वारि}\ H}}\skt{ }\edtext{\skt{पूनस्तथास्य}}{\Afootnote{\skt{पूनस्तथास्य}\ K; \skt{पुनस्तथास्य नि}\ H}}\skt{ }\edtext{\skt{कृत्वानि}}{\Afootnote{K; \skt{कृत्वा}\ H}}\skt{ }\edtext{\skt{क्वव\-था\-चि\-त्र\-कस्य}}{\Afootnote{K\textsuperscript{pc}; \skt{क्वथाचित्रकस्य ।}\ K\textsuperscript{ac}; \skt{क्वथचित्रकस्य ।}\ H}}\skt{ । }\pend
\pstart\versepar \skt{}\edtext{\skt{तद्वारके}}{\Afootnote{\skt{तद्वारके}\ K; \skt{तद्वाढके}\ H}}\skt{ चापि }\edtext{\skt{जलस्य}}{\Afootnote{K; \skt{जलास्य}\ H}}\skt{ सिद्धं }\edtext{\skt{पादा\-व\-शे\-षा\-द्व\-षे\-गा\-ल\-नीयं}}{\Afootnote{K; \skt{पादावशेषं क्वथगालनीयं ॥}\ H}}\skt{ ।}\pend

\pstart\versepar \ledsidenote{\passref{6.42.40.add17}}\skt{तेन }\edtext{\skt{प्रपिष्येत्व}}{\Afootnote{K; \skt{प्रपिष्येव}\ H}}\skt{ }\edtext{\skt{पलो\-न्मि\-ता\-ना\-मि\-सं\-व\-रिष्ठं}}{\Afootnote{K; \skt{पलोन्मितानां मिसम्वरिष्ठं}\ H}}\skt{ गणमौषधानां । }\pend
\pstart\versepar \skt{कटुत्रिवर्गं }\edtext{\edtext{\skt{लवणान्यथाष्टौ}}{\lemma{\skt{लवणान्य्}}\Afootnote{K; \skt{लवणा}\ H}}\skt{ द्वे}}{\lemma{\skt{अथाष्टौ द्वे}}\Afootnote{K; \skt{न्यहाष्ठौ द्वे}\ H}}\skt{ }\edtext{\skt{रोहि\-णी\-ञ्चा\-म्ल\-क\-वे\-त\-मञ्च}}{\lemma{\skt{चाम्ल\-क\-वे\-त\-मञ्}}\Afootnote{K; \skt{चाम्ल\-क\-वे\-त\-सञ्}\ H}}\skt{ ।}\pend

\pstart\versepar \ledsidenote{\passref{6.42.40.add18}}\edtext{\skt{ञ्च}}{\lemma{\skt{ञ्}}\Afootnote{K; \emph{om.}\ H}}\skt{ }\edtext{\skt{भागीञ्च}}{\Afootnote{K; \skt{भाग्नीं च}\ H}}\skt{ वयाञ्च पाठं }\edtext{\skt{सजीवकं}}{\Afootnote{K; \skt{सजीरकं}\ H}}\skt{ }\edtext{\edtext{\skt{मूल\-म\-था\-ज\-मोदां}}{\lemma{\skt{मूलम्}}\Afootnote{K; \skt{मूल}\ H}}}{\lemma{\skt{अथाजमोदां}}\Afootnote{K; \skt{महाजमोदां ।}\ H}}\skt{ । }\pend
\pstart\versepar \skt{}\edtext{\skt{हिंगुन्तु }\edtext{\skt{माञ्चातिविषञ्च}}{\lemma{\skt{चातिविषञ्च}}\Afootnote{K; \skt{चातिविषञ्च्q}\ H}}}{\lemma{\skt{हिंगुन्तु माञ्}}\Afootnote{K; \skt{हिङ्गुं तु सं}\ H}}\skt{ }\edtext{\edtext{\skt{मुस्तमेञ्च}}{\lemma{\skt{मुस्तमेञ्}}\Afootnote{\skt{मुस्तमेञ्च}\ K; \skt{मुस्तमेलां स च}\ H}}\skt{ गजपिप्पलीञ्च}}{\lemma{\skt{च गजपिप्पलीञ्}}\Afootnote{\skt{गजपिप्पलीञ्}\ K; \skt{व्याग\-ज\-पि\-प्प\-लीञ्}\ H}}\skt{ ।}\pend

\pstart\versepar \ledsidenote{\passref{6.42.40.add19}}\edtext{\skt{आलो\-ग\-र्भ\-म\-म\-त्र\-पिष्ठं}}{\Afootnote{K; \skt{आलो गर्भं कृत्वि\-म\-म\-त्र\-मिष्टं}\ H}}\skt{ }\edtext{\skt{पूर्वश्रूतं}}{\Afootnote{K; \skt{पूर्व्वस्रुतं}\ H}}\skt{ }\edtext{\edtext{\skt{क्षार\-म\-थो\-प\-चेत्तु}}{\lemma{\skt{क्षारम्}}\Afootnote{K; \skt{क्षार}\ H}}}{\lemma{\skt{अथोपचेत्तु}}\Afootnote{K; \skt{मथापचेत्तु ।}\ H}}\skt{ । }\pend
\pstart\versepar \skt{मृद्वग्निना }\edtext{\skt{वाक्क्वथितं}}{\Afootnote{K; \skt{वा क्वथितं}\ H}}\skt{ }\edtext{\skt{सनैस्तु}}{\Afootnote{K; \skt{शनैस्तु}\ H}}\skt{ }\edtext{\skt{मोद\-वि\-ले\-पा\-त्प\-रि\-शो\-षि\-तस्तु}}{\lemma{\skt{मोद\-वि\-ले\-पा\-त्प\-रि\-शो\-षि\-तस्}}\Afootnote{K; \skt{मोद\-र्व्वि\-ले\-पा\-त्प\-रि\-शो\-षि\-तस्}\ H}}\skt{ ।}\pend

\pstart\versepar \ledsidenote{\passref{6.42.40.add20}}\edtext{\skt{कोलप्रमाणां}}{\Afootnote{K; \skt{कोलप्रमाणं}\ H}}\skt{ गुटिकां विद\-ध्या\-द्द्वे\-भ\-क्ष\-यि\-त्वैव मतोन्नमद्यात् । }\pend
\pstart\versepar \skt{ततः पदानां }\edtext{\edtext{\skt{शामेव}}{\lemma{\skt{शाम्}}\Afootnote{K; \skt{श}\ H}}}{\lemma{\skt{एव}}\Afootnote{K; \skt{तमेव}\ H}}\skt{ गत्वा }\edtext{\skt{यतप्रकांक्षः}}{\Afootnote{K; \skt{यतः प्रकांक्षः}\ H}}\skt{ }\edtext{\edtext{\skt{पुनरेव}}{\lemma{\skt{पुनर्}}\Afootnote{K; \skt{पुन}\ H}}}{\lemma{\skt{एव}}\Afootnote{K; \skt{रेव ॥}\ H}}\skt{ ।}\edtext{}{\lemma{\skt{मतोन्नमद्यात् । ततः}}\Afootnote{K; \skt{मओन्नमद्यात् । ततः}\ H}}\pend

\pstart\versepar \ledsidenote{\passref{6.42.40.add21}}\skt{द्वे भक्षयित्वौव मतो\-न्न\-म\-द्या\-द\-पू\-र्व\-व\-द्जा\-त\-म\-हा\-प्र\-कांक्षः । }\pend
\pstart\versepar \skt{स एव गुटि\-का\-प्र\-यो\-गा\-द\-हह्̤ सुखेनैव सहस्रसापि ।}\pend

\pstart\versepar \ledsidenote{\passref{6.42.40.add22}}\skt{तिस्रस्तु गुल्मीगुटिकां प्रयुक्त्वा मूलात्ततः पञ्च निहन्ति गुल्मान् । }\pend
\pstart\versepar \skt{}\edtext{\skt{चत्वारम्चार्मांमि}}{\lemma{\skt{चार्मांमि}}\Afootnote{K; \skt{अशांसि}\ H}}\skt{ निहनि }\edtext{\skt{पीताः}}{\Afootnote{K; \skt{पीताः}\ H}}\skt{ }\edtext{\skt{त था}}{\Afootnote{K; \skt{ध्योरे तथा}\ H}}\skt{ }\edtext{\skt{वातकहादरे}}{\Afootnote{K; \skt{वातकफादरे}\ H}}\skt{ च ।}\edtext{}{\lemma{\skt{गुल्मान् । चत्वारम्}}\Afootnote{K; \skt{गुल्मान् । चत्वारि}\ H}}\pend

\pstart\versepar \ledsidenote{\passref{6.42.40.add23}}\skt{येवा परे }\edtext{\skt{वाअ\-क\-फ\-प्र\-वृ\-त्ता\-रा\-नृणां}}{\Afootnote{K; \skt{वात\-क\-फ\-प्र\-वृत्ता द्योरानृणां}\ H}}\skt{ कुक्षिगता विकाराः । }\pend
\pstart\versepar \skt{}\edtext{\skt{तां}}{\Afootnote{K; \skt{ता}\ H}}\skt{ }\edtext{\skt{हन्ति}}{\Afootnote{K; \skt{न्हनि}\ H}}\skt{ }\edtext{\skt{सर्वांश्च}}{\Afootnote{K; \skt{सर्व्वान्श्च}\ H}}\skt{ }\edtext{\skt{यथंञ्च}}{\Afootnote{K; \skt{यहुंश्च}\ H}}\skt{ पीताः घृतेन चोक्तेन जलेन चापि ।}\pend

\pstart\versepar \ledsidenote{\passref{6.42.40.add24}}\skt{कायग्निसैन्यं }\edtext{\edtext{\skt{ग्रह\-णी\-प्र\-दो\-ष\-म\-स्त\-वि\-दा\-म\-वाप्य}}{\lemma{\skt{ग्रह\-णी\-प्र\-दो\-ष\-म\-स्त\-वि\-दाम्}}\Afootnote{K; \skt{ग्रह\-णी\-प्र\-दो\-ष\-म\-स्रं\-विदा}\ H}}}{\lemma{\skt{अवाप्य}}\Afootnote{K; \skt{हित्वमवाप्य}\ H}}\skt{ जीर्णं । }\pend
\pstart\versepar \skt{निहन्ति }\edtext{\skt{सद्यो}}{\Afootnote{K; \skt{सद्यो}\ H}}\skt{ जरयन्ति चान्नं }\edtext{\skt{पासा\-ण\-चू\-र्णै\-रपि}}{\lemma{\skt{पासाणचूर्णैर्}}\Afootnote{K; \skt{पाषाणचूर्णैर्}\ H}}\skt{ मिश्रितं यत् ॥}\pend

\pstart\versepar \ledsidenote{\passref{6.42.40.add25}}\edtext{\skt{कुरिणी}}{\Afootnote{K; \skt{कुम्वुरूणि}\ H}}\skt{ }\edtext{\skt{त्रिप\-ह्ला\-म\-रि\-चानि}}{\Afootnote{K; \skt{त्रिफला मरिचानि}\ H}}\skt{ }\edtext{\skt{विडंगंकं}}{\Afootnote{K; \skt{विडङ्गकं ।}\ H}}\skt{ । }\pend
\pstart\versepar \skt{}\edtext{\skt{पिप्प\-ली\-चा\-द\-मो\-दञ्च}}{\Afootnote{K; \skt{पिप्प\-ल्य\-श्चा\-ज\-मोदं च}\ H}}\skt{ }\edtext{\skt{कुटाजस्य}}{\Afootnote{K; \skt{कूटजस्य}\ H}}\skt{ फलानि च ।}\pend

\pstart\versepar \ledsidenote{\passref{6.42.40.add26}}\skt{चित्रकं कालिकाजाजी एकैकं पलिकं }\edtext{\skt{मितं}}{\Afootnote{K; \skt{मिअं ।}\ H}}\skt{ । }\pend
\pstart\versepar \skt{अष्टौपलानि त्रिवृतः }\edtext{\skt{पञ्चाशः}}{\Afootnote{K; \skt{पञ्चाश}\ H}}\skt{ }\edtext{\skt{स्याद्गुडस्य}}{\Afootnote{K; \skt{त्स्याद्गुडस्य}\ H}}\skt{ तु ।}\pend

\pstart\versepar \ledsidenote{\passref{6.42.40.add27}}\skt{तैलस्य }\edtext{\edtext{\skt{रस\-मा\-म\-ल\-कस्य}}{\lemma{\skt{रसम्}}\Afootnote{K; \skt{कुडद्वौ तु रस}\ H}}\skt{ तु}}{\lemma{\skt{आमलकस्य तु}}\Afootnote{K; \skt{मामलकस्य च ।}\ H}}\skt{ । }\pend
\pstart\versepar \skt{त्रयः प्रस्ठाच्च }\edtext{\skt{तत्सर्वां}}{\Afootnote{K; \skt{तत्सर्व्वं}\ H}}\skt{ }\edtext{\skt{शनैमृद्वग्निना}}{\Afootnote{K; \skt{शनैर्म्मृद्वग्निना}\ H}}\skt{ पचेत् ।}\pend

\pstart\versepar \ledsidenote{\passref{6.42.40.add28}}\skt{घनीभूतं फलैः कृत्वा }\edtext{\edtext{\skt{सम\-मा\-म\-ल\-कस्य}}{\lemma{\skt{समम्}}\Afootnote{K; \skt{सम}\ H}}\skt{ तु}}{\lemma{\skt{आमलकस्य तु}}\Afootnote{K; \skt{मामलकस्य तु ।}\ H}}\skt{ । }\pend
\pstart\versepar \skt{अग्निप्रमाणं विज्ञाय }\edtext{\skt{भक्ष\-ये\-त्प्रा\-त\-रु\-त्थितः}}{\Afootnote{K; \skt{भक्ष\-ये\-त्प्रा\-त\-रू\-त्थितः ॥}\ H}}\pend

\pstart\versepar \ledsidenote{\passref{6.42.40.add29}}\edtext{\skt{पञ्चगुल्मां}}{\Afootnote{K; \skt{प}\ H}}\skt{ }\edtext{\skt{प्रमेहांञ्च}}{\Afootnote{K; \skt{प्रमेहाश्च}\ H}}\skt{ }\edtext{\skt{पाण्डुरोगं}}{\Afootnote{K; \skt{पाण्डुरोगा}\ H}}\skt{ }\edtext{\skt{सकामलां}}{\Afootnote{K; \skt{न्सकामलां ।}\ H}}\skt{ । }\pend
\pstart\versepar \skt{शोषं }\edtext{\skt{कुष्ठोदरामीम}}{\Afootnote{K; \skt{कुष्थोदराणांसि}\ H}}\skt{ }\edtext{\skt{समयेच्च}}{\Afootnote{K; \skt{शमयेच्च}\ H}}\skt{ }\edtext{\skt{भगन्दरां}}{\Afootnote{K; \skt{भगन्दरान् ॥}\ H}}\skt{ ।}\pend

\pstart\versepar \ledsidenote{\passref{6.42.40.add30}}\edtext{\skt{योनिदोषहरो}}{\Afootnote{K; \skt{योनिरोगहरो}\ H}}\skt{ वृष्यो मेधा }\edtext{\skt{वलवपुः}}{\Afootnote{K; \skt{वलवपु}\ H}}\skt{ }\edtext{\skt{करः}}{\Afootnote{K; \skt{क्कराः ।}\ H}}\skt{ । }\pend
\pstart\versepar \skt{गुड\-क\-ल्या\-ण\-सं\-क्षस्तु मतश्चापि रसायनैः ॥}\pend

\pstart \ledsidenote{\passref{6.42.40cd}}\skt{तिले क्षुरकपालाश सार्षपं }\edtext{\skt{यवनालजं}}{\Afootnote{K H; \skt{यावनालजम् ॥}\ A}}\skt{ ।}\pend

\pstart\versepar \ledsidenote{\passref{6.42.41}}\edtext{\skt{भस्मा\-मू\-क\-ज\-ञ्चापि}}{\lemma{\skt{भस्मामूकजञ्}}\Afootnote{K; \skt{भस्म मूलकजं}\ A H}}\skt{ }\edtext{\skt{गोज\-वि\-ख\-र\-ह\-स्तिनां}}{\Afootnote{K H; \skt{गोजा\-वि\-ख\-र\-ह\-न्ति\-नाम् ।}\ A}}\skt{ । }\pend
\pstart\versepar \skt{मूत्रेण महिषीणाञ्च }\edtext{\edtext{\skt{पलिकैश्चापि}}{\lemma{\skt{पलिकैश्}}\Afootnote{K H; \skt{पालिकैश्}\ A}}}{\lemma{\skt{चापि}}\Afootnote{K H; \skt{चाव}\ A}}\skt{ चूर्ण्णितैः ।}\pend

\pstart\versepar \ledsidenote{\passref{6.42.42}}\edtext{\skt{कुष्ठ\-सै\-न्ध\-व\-य\-ष्ठी\-क\-ना\-ग\-र\-क्रि\-मि\-ह\-न्तिभिः}}{\Afootnote{K; \skt{कुष्ठ\-सै\-न्ध\-व\-य\-ष्ट्या\-ह्व\-ना\-ग\-र\-कृ\-मि\-घा\-तिभिः ।}\ A; \skt{कुष्ठ\-सै\-न्ध\-व\-य\-ष्ठीक नाग\-र\-क्रि\-मि\-हं\-त्रिभिः ॥}\ H}}\skt{ । }\pend
\pstart\versepar \skt{}\edtext{\skt{अजमोदाच्च}}{\Afootnote{K H; \skt{साजमोदैश्च}\ A}}\skt{ दशभिः सामुद्राच्च }\edtext{\skt{पलैयुतं}}{\Afootnote{K; \skt{पलैर्युतम् ॥}\ A; \skt{पलैर्युतं ।}\ H}}\skt{ ।}\pend

\pstart\versepar \ledsidenote{\passref{6.42.43}}\edtext{\skt{अय\-स्पा\-त्रे\-ग्नि\-क\-ल्पेन}}{\lemma{\skt{पात्रेग्निकल्पेन}}\Afootnote{K; \skt{पात्रेऽग्निनाल्पेन}\ A; \skt{पात्रेऽ\-ग्नि\-क\-क\-स्येन}\ H}}\skt{ पक्त्वा }\edtext{\skt{लेह्य\-म\-थो\-द्ध\-रेत्}}{\lemma{\skt{लेह्यम्}}\Afootnote{K; \skt{लेह्यम्(लहम्)}\ A; \skt{लेह्य म}\ H}}\skt{ । }\pend
\pstart\versepar \skt{तस्य मात्रां }\edtext{\skt{पिबेद्दध्ना}}{\Afootnote{A; \skt{पिवेद्दध्ना}\ K H}}\skt{ }\edtext{\skt{सुरयाषा पि}}{\Afootnote{\skt{सुरया सर्पिषापि}\ A H; \skt{सुरयास्̤आ पि}\ K}}\skt{ वा ।}\edtext{}{\lemma{\skt{अथोद्धरेत् । तस्य}}\Afootnote{A K; \skt{त्थोद्धरेत् । तस्य}\ H}}\pend

\pstart\versepar \ledsidenote{\passref{6.42.44}}\skt{धान्या\-म्ले\-नो\-ष्ण\-तो\-येन कौलत्थेन रसेन वा । }\pend
\pstart\versepar \skt{}\edtext{\skt{गुल्मं}}{\Afootnote{K; \skt{गुल्मा}\ A; \skt{गुल्म}\ H}}\skt{ }\edtext{\skt{वातविकाराश्च}}{\Afootnote{K H; \skt{न्वातविकारांश्च}\ A}}\skt{ }\edtext{\skt{क्षारोयं}}{\Afootnote{K H; \skt{क्षारोऽ}\ A}}\skt{ }\edtext{\edtext{\skt{हन्त्यसंशयं}}{\lemma{\skt{हन्त्य्}}\Afootnote{K; \skt{यं हन्}\ A; \skt{हन्}\ H}}}{\lemma{\skt{असंशयं}}\Afootnote{K; \skt{त्यसंशयम् ॥}\ A H}}\skt{ ॥}\pend

\pstart\versepar \ledsidenote{\passref{6.42.45}}\edtext{\skt{स्वर्जि\-का\-कु\-ष्ट\-ह\-सहैः}}{\Afootnote{K; \skt{स्वर्जि\-का\-कु\-ष्ठ\-स\-हितः}\ A H}}\skt{ क्षारः }\edtext{\skt{केतकजो}}{\Afootnote{K; \skt{केतकिजो}\ A; \skt{केतकजो}\ H}}\skt{ }\edtext{\skt{पि वा}}{\Afootnote{A K; \skt{पि वा ।}\ H}}\skt{ । }\pend
\pstart\versepar \skt{}\edtext{\skt{पीततैलेन}}{\Afootnote{K H; \skt{तैलेन}\ A}}\skt{ }\edtext{\skt{शमयेद्दुल्मं}}{\Afootnote{K; \skt{शमयेत्पीतो गुल्मं}\ A; \skt{श मयेत्गुल्मं}\ H}}\skt{ पवनसंभवं ।}\pend

\pstart\versepar \ledsidenote{\passref{6.42.46}}\skt{पीतं }\edtext{\skt{सुखांवुनावापि}}{\Afootnote{K H; \skt{सुखाम्बुना वापि}\ A}}\skt{ स्वर्जि\-का\-कु\-ष्ठ\-सै\-न्धवं ॥ }\pend
\pstart\versepar \skt{}\edtext{\skt{वृश्ची\-क\-मु\-रू\-वू\-कश्च}}{\Afootnote{K; \skt{वृश्ची\-व\-मु\-रु\-बू\-कं(\-वृ\-श्ची\-र\-मु\-रु\-बूकं) च}\ A; \skt{वृश्चि\-क\-मु\-रू\-वू\-कश्च}\ H}}\skt{ }\edtext{\skt{वर्षाभूवृहती}}{\Afootnote{K H; \skt{वर्षाभूर्बृहती}\ A}}\skt{ द्वयं ।}\pend

\pstart\versepar \ledsidenote{\passref{6.42.47}}\skt{चित्रकञ्च }\edtext{\skt{जलद्रोणै}}{\Afootnote{K; \skt{जलद्रोणे}\ A H}}\skt{ }\edtext{\skt{पचे\-त्पा\-दा\-व\-शे\-षितं}}{\Afootnote{K H; \skt{पक्त्वा पादा\-व\-शे\-षि\-तम् ।}\ A}}\skt{ । }\pend
\pstart\versepar \skt{}\edtext{\skt{मागधाचिकुम्भे}}{\Afootnote{K; \skt{माग\-धी\-चि\-क्त्र\-क\-क्षौ\-द्र\-लिप्ते कुम्भे}\ A; \skt{मागधचित्रक लिप्ते कुम्भे}\ H}}\skt{ निधापयेत् ।}\pend

\pstart\versepar \ledsidenote{\passref{6.42.48}}\skt{मधुनः प्रस्थमावाप्य पथ्या\-चू\-र्ण्णा\-र्ध\-सं\-युतं । }\pend
\pstart\versepar \skt{}\edtext{\skt{विशोषितं}}{\Afootnote{K H; \skt{बुसोषितं}\ A}}\skt{ दशाहन्तु }\edtext{\skt{जीर्ण्णभकः}}{\Afootnote{K; \skt{जीर्णभक्तः}\ A H}}\skt{ }\edtext{\skt{पिवेन्नरः}}{\lemma{\skt{पिवेन्}}\Afootnote{K H; \skt{पिबेन्}\ A}}\skt{ ।}\pend

\pstart\versepar \ledsidenote{\passref{6.42.49}}\edtext{\skt{अरिष्थोयं}}{\Afootnote{K; \skt{अरिष्टोऽयं}\ A; \skt{अरिष्ठोयं}\ H}}\skt{ }\edtext{\skt{जये\-द्गु\-ल्म\-वि\-पाकं}}{\Afootnote{K; \skt{जये\-द्गु\-ल्म\-म\-वि\-पाक}\ A; \skt{जये\-द्गु\-ल्म\-म\-वि\-पा\-कञ्}\ H}}\skt{ }\edtext{\skt{च}}{\Afootnote{K H; \skt{म}\ A}}\skt{ }\edtext{\skt{दुस्तरं}}{\Afootnote{K H; \skt{रोचकम् ।}\ A}}\skt{ ॥ }\pend
\pstart\versepar \skt{}\edtext{\skt{पाठा\-नि\-कु\-म्भ\-र\-ज\-नी\-त्रि\-क\-टु\-त्रि\-फ\-ल्-\-अ\-ग्निका}}{\Afootnote{K; \skt{पाठा\-नि\-कु\-म्भ\-र\-ज\-नी\-त्रि\-क\-टु\-त्रि\-फ\-ला\-ग्नि\-कम् ॥}\ A; \skt{पाठा\-नि\-कु\-म्भ\-र\-जनी त्रिक\-टु\-त्रि\-फ\-ला\-ग्निका ।}\ H}}\skt{ ।}\pend

\pstart \ledsidenote{\passref{6.42.50ab}}\skt{लवणं }\edtext{\skt{वृक्षकवीजानि}}{\Afootnote{K H; \skt{वृक्ष}\ A}}\skt{ }\edtext{\skt{चूर्ण्ण\-गो\-मू\-त्र\-सा\-धितं}}{\Afootnote{K H; \skt{बीजं च तुल्यं स्यादनवो गुडः ।}\ A}}\skt{ ।}\pend

\pstart\versepar \ledsidenote{\passref{6.42.51}}\edtext{\skt{घनीभूतं तु}}{\Afootnote{K H; \skt{गुटि\-का\-स्त\-द्घृ\-नी\-भूतं}\ A}}\skt{ }\edtext{\skt{गुटिकां}}{\Afootnote{K H; \emph{om.}\ A}}\skt{ कृत्वा खादे\-द\-भु\-क्त\-वान् । }\pend
\pstart\versepar \skt{गुल्म\-प्ली\-हा\-ग्नि\-सा\-दांस्तां }\edtext{\edtext{\skt{नाश\-ये\-यु\-र\-शे\-षतः}}{\lemma{\skt{नाशयेयुर्}}\Afootnote{K H; \skt{नाशयेयु}\ A}}}{\lemma{\skt{अशेषतः}}\Afootnote{K H; \skt{रशेषतः ॥}\ A}}\skt{ ॥}\edtext{}{\lemma{\skt{खादे\-द\-भु\-क्त\-वान् । गुल्म\-प्ली\-हा\-ग्नि\-सा\-दांस्तां}}\Afootnote{K H; \skt{खादे\-द\-भु\-क्त\-वान् । गुल्म\-प्ली\-हा\-ग्नि\-सा\-दांता}\ A}}\pend

\pstart \ledsidenote{\passref{6.42.53cd}}\edtext{\skt{सुखो\-ष्ठा\-जां\-ग\-ल\-रसाः}}{\Afootnote{K H; \skt{सुखोष्णा जाङ्गलरसाः}\ A}}\skt{ सुस्निग्धा व्यक्तसैन्धवाः ।}\pend

\pstart \ledsidenote{\passref{6.42.54cd}}\skt{पेया }\edtext{\skt{वातहरैः}}{\Afootnote{A K; \skt{वा अहरैः}\ H}}\skt{ }\edtext{\skt{सिद्धाः}}{\Afootnote{A K; \skt{सिद्धा}\ H}}\skt{ कौलत्थाः }\edtext{\skt{संस्कृताः}}{\Afootnote{A K; \skt{संकृता}\ H}}\skt{ रसाः ॥}\pend

\pstart \ledsidenote{\passref{6.42.55ab}}\edtext{\skt{खलाः}}{\Afootnote{A K; \skt{खला}\ H}}\skt{ सपञ्चमूलाश्च }\edtext{\skt{गुल्मिनो}}{\Afootnote{K H; \skt{गुल्मिनां}\ A}}\skt{ भोजने हिताः ।}\pend

\pstart\versepar \ledsidenote{\passref{6.42.56}}\edtext{\skt{कुम्भी\-पि\-ण्डे\-ष्ट\-क\-स्वेदां}}{\Afootnote{K; \skt{कुम्भी\-पि\-ण्डे\-ष्ट\-का\-स्वेदा}\ A; \skt{कुम्भी\-पि\-ण्डे\-ष्ट\-क\-स्वेदा}\ H}}\skt{ }\edtext{\skt{कुर्याच्चवकुशो}}{\Afootnote{K; \skt{न्कारयेत्कुशलो}\ A; \skt{न कुर्याच्चवकुशो}\ H}}\skt{ भिषक् । }\pend
\pstart\versepar \skt{गुल्मिनः सर्व एवोक्ता }\edtext{\skt{दुर्विरेच्यतमा}}{\Afootnote{A K; \skt{दुर्व्विरेच्य तमा}\ H}}\skt{ }\edtext{\skt{भृशाः}}{\Afootnote{K H; \skt{भृशम् ॥}\ A}}\skt{ ।}\pend

\pstart\versepar \ledsidenote{\passref{6.42.57}}\edtext{\edtext{\skt{अतश्चैनं}}{\lemma{\skt{अतश्}}\Afootnote{K H; \skt{अत}\ A}}\skt{ तु}}{\lemma{\skt{चैनं तु}}\Afootnote{K; \skt{शैतांस्तु}\ A; \skt{चना -}\ H}}\skt{ }\edtext{\skt{सुस्विन्नां}}{\Afootnote{K H; \skt{सुस्विन्ना}\ A}}\skt{ स्रंस\-ने\-नो\-प\-पा\-द\-येत् । }\pend
\pstart\versepar \skt{विम्लापनं }\edtext{\skt{भेदनानि}}{\Afootnote{K H; \skt{भ्यञ्जनानि}\ A}}\skt{ तथैव }\edtext{\skt{स्यञ्जनानि}}{\Afootnote{K; \skt{दहनानि}\ A; \skt{स्यञ्जनानि}\ H}}\skt{ च ।}\edtext{}{\lemma{\skt{स्रंस\-ने\-नो\-प\-पा\-द\-येत् । विम्लापनं}}\Afootnote{K; \skt{न्स्रंस\-ने\-नो\-प\-पा\-द\-येत् । विम्लापना}\ A; \skt{स्रंस\-ने\-नो\-प\-पा\-द\-येत ॥ विम्लापनं}\ H}}\pend

\pstart\versepar \ledsidenote{\passref{6.42.58}}\edtext{\skt{उपनाहाश्च}}{\Afootnote{A K; \skt{उपनाहोश्च}\ H}}\skt{ कर्तव्याः }\edtext{\skt{सुखोष्ठाः}}{\Afootnote{K H; \skt{सुखोष्णाः}\ A}}\skt{ }\edtext{\skt{साल्वलादयः}}{\Afootnote{K H; \skt{साल्वणादयः ।}\ A}}\skt{ । }\pend
\pstart\versepar \skt{उदरोक्तानि }\edtext{\skt{सर्पी\-षि\-मू\-त्र\-वर्तिः}}{\Afootnote{K; \skt{सर्पींषि मूत्रवर्ति}\ A; \skt{सर्पि\-षि\-मू\-त्र\-वर्तिः}\ H}}\skt{ }\edtext{\skt{क्रियास्तथा}}{\Afootnote{K; \skt{क्रियास्तथा ॥}\ A; \skt{क्रिया था ।}\ H}}\skt{ ।}\pend

\pstart \ledsidenote{\passref{6.42.59ab}}\edtext{\skt{लवणानि च}}{\Afootnote{A K; \skt{लवणाणि च}\ H}}\skt{ योज्यानि }\edtext{\skt{यान्यु\-क्ता\-न्य\-नि\-ला\-च्कुर्ये}}{\Afootnote{K H; \skt{यान्यु\-क्ता\-न्य\-नि\-ला\-मये ।}\ A}}\skt{ ।}\pend

\pstart \ledsidenote{\passref{6.42.60cd}}\edtext{\skt{दन्ती\-द्र\-व\-न्ती\-मू\-लेषु}}{\Afootnote{K H; \skt{दन्ती\-चि\-त्र\-क\-मू\-लेषु}\ A}}\skt{ तथा वातहरेषु च ।}\pend

\pstart \ledsidenote{\passref{6.42.61ab}}\edtext{\skt{कार\-ये\-दा\-स\-वा\-रि\-ष्ठा\-श्लो\-क\-स्थाने}}{\Afootnote{K; \skt{कुर्या\-द\-रि\-ष्टा\-न्स\-र्वांश्च श्लोकस्थाने}\ A; \skt{कार\-ये\-दा\-स\-वा\-रि\-ष्ठा\-न्श्लो\-क\-स्थाने}\ H}}\skt{ }\edtext{\skt{यथेरितां}}{\Afootnote{K; \skt{यथेरितान्}\ A H}}\skt{ ।}\pend

\pstart \ledsidenote{\passref{6.42.59cd}}\edtext{\skt{वात\-व\-र्चो\-नि\-रो\-धैषु}}{\Afootnote{K; \skt{वात\-व\-र्चो\-नि\-रोधे तु}\ A; \skt{वातवर्च्चानि रोधे वा}\ H}}\skt{ सामु\-द्रा\-र्द्र\-क\-स\-र्षपैः}\pend

\pstart \ledsidenote{\passref{6.42.60ab}}\skt{कृत्वा }\edtext{\skt{पायौ}}{\Afootnote{A K; \skt{पायो}\ H}}\skt{ }\edtext{\skt{निधातव्या}}{\Afootnote{K H; \skt{विधातव्या}\ A}}\skt{ वर्त्तयो मरिचोत्तराः ।}\pend

\pstart \ledsidenote{\passref{6.42.61cd}}\edtext{\skt{खादेद्वान्यांकुरां}}{\Afootnote{K; \skt{खादेद्वाप्यङ्कुरा}\ A; \skt{खादेद्वान्याङ्कुरा}\ H}}\skt{ }\edtext{\skt{भृष्तां}}{\Afootnote{K; \skt{न्भृष्टा}\ A; \skt{भृष्टां}\ H}}\skt{ }\edtext{\skt{पूती\-क\-नृ\-प\-वृ\-क्षयोः}}{\Afootnote{K H; \skt{न्पूती\-क\-नृ\-प\-वृ\-क्षयोः ॥}\ A}}\skt{ ।}\pend

\pstart\versepar \ledsidenote{\passref{6.42.62}}\edtext{\skt{ऊर्ध्ववातन्तु}}{\Afootnote{K H; \skt{ऊर्ध्ववातं}\ A}}\skt{ }\edtext{\skt{मनुजं}}{\Afootnote{K H; \skt{मनुष्यं च}\ A}}\skt{ गुल्मिनं न निरूहये ॥ }\pend
\pstart\versepar \skt{}\edtext{\skt{पिबे\-त्त्रि\-वृ\-न्ना\-ग\-रम्वा}}{\Afootnote{\skt{त् । पिबे\-त्त्रि\-वृ\-न्ना\-गरं वा}\ A; \skt{पिवे\-त्त्रि\-वृ\-न्ना\-ग\-रम्वा}\ K; \skt{त् ॥ पिवे\-त्त्रि\-वृ\-न्ना\-ग\-रम्वा}\ H}}\skt{ }\edtext{\skt{सगुडा वा}}{\Afootnote{K; \skt{सगुडां वा}\ A; \skt{सगुडम्वा}\ H}}\skt{ हरीतकीं ।}\pend

\pstart\versepar \ledsidenote{\passref{6.42.63}}\edtext{\skt{गुग्गुलं}}{\Afootnote{K; \skt{गुग्गुलुं}\ A; \skt{गुग्गुल}\ H}}\skt{ }\edtext{\skt{वृवृतां}}{\Afootnote{K; \skt{त्रिवृतां}\ A; \skt{न्भृवृतां}\ H}}\skt{ }\edtext{\skt{दनीं}}{\Afootnote{K H; \skt{दन्तीं}\ A}}\skt{ द्रवन्तीं }\edtext{\skt{सैन्धवं}}{\Afootnote{A K; \skt{सैन्धवाम्}\ H}}\skt{ वचां । }\pend
\pstart\versepar \skt{}\edtext{\edtext{\skt{मूत्र\-म\-द्य\-प\-यो\-द्रा\-क्षा\-र\-सै\-र्वीक्ष्य}}{\lemma{\skt{मूत्र\-म\-द्य\-प\-यो\-द्रा\-क्षा\-र\-सैर्}}\Afootnote{K H; \skt{मूत्र\-म\-द्य\-प\-यो\-द्रा\-क्षा\-रसै}\ A}}}{\lemma{\skt{वीक्ष्य}}\Afootnote{K; \skt{र्वीक्ष्य}\ A; \skt{व्वीक्ष}\ H}}\skt{ }\edtext{\skt{वलावलं}}{\Afootnote{K H; \skt{बलाबलम् ॥}\ A}}\skt{ ।}\pend

\pstart\versepar \ledsidenote{\passref{6.42.64}}\skt{एवं पीलूनि }\edtext{\skt{पिष्ठानि}}{\Afootnote{K H; \skt{भृष्टानि}\ A}}\skt{ }\edtext{\skt{पिबे\-त्स\-ल\-व\-णानि}}{\Afootnote{\skt{पिबे\-त्स\-ल\-व\-णानि तु ।}\ A; \skt{पिवे\-त्स\-ल\-व\-णानि}\ K; \skt{पिवे\-त्स\-ल\-व\-णानि तु ॥}\ H}}\skt{ त्य् ॥ }\pend
\pstart\versepar \skt{पिप्पै\-पि\-प्प\-ली\-मूलं }\edtext{\skt{चित्र\-का\-जा\-जि\-सै\-न्धवैः}}{\Afootnote{K; \skt{चव्य\-चि\-त्र\-क\-सै\-न्धवैः ॥}\ A; \skt{उल\-चि\-त्र\-का\-जा\-जि\-सै\-न्धवैः ।}\ H}}\skt{ ।}\edtext{}{\lemma{\skt{त्य् ॥ पिप्पै\-पि\-प्प\-ली\-मूलं}}\Afootnote{K; \skt{पिप्प\-ली\-पि\-प्प\-ली\-मूल}\ A; \skt{पिप्पैपिप्पै}\ H}}\pend

\pstart\versepar \ledsidenote{\passref{6.42.65}}\skt{युक्ता }\edtext{\skt{सुरा}}{\Afootnote{K H; \skt{हन्ति सुरा}\ A}}\skt{ }\edtext{\skt{घ्नन्ति}}{\Afootnote{K H; \emph{om.}\ A}}\skt{ गुल्मं }\edtext{\skt{शीघ्रकाले}}{\Afootnote{K H; \skt{शीघ्रं काले}\ A}}\skt{ प्रयोजिता । }\pend
\pstart\versepar \skt{}\edtext{\skt{वद्धविन्मारूतो}}{\Afootnote{K H; \skt{बद्धविण्मारुतो}\ A}}\skt{ गुल्मी भुञ्जीत }\edtext{\skt{पये}}{\Afootnote{K; \skt{पय}\ A H}}\skt{ }\edtext{\skt{यावां}}{\Afootnote{K; \skt{सा यवान्}\ A H}}\skt{ ।}\pend

\pstart \ledsidenote{\passref{6.42.66ab}}\edtext{\skt{कुल्मांषाश्चा}}{\Afootnote{K; \skt{कुल्माषान्वा}\ A H}}\skt{ }\edtext{\skt{वहुस्नेहां}}{\Afootnote{K; \skt{बहुस्नेहा}\ A; \skt{वहुस्नेहा}\ H}}\skt{ }\edtext{\skt{भक्ष\-ये\-ल्ल\-व\-णै\-त्तरां}}{\Afootnote{K; \skt{न्भक्ष्ये\-ल्ल\-व\-णो\-त्त\-रान् ।}\ A; \skt{न्भक्ष\-ये\-ल्ल\-व\-णा\-त्त\-रान् ॥}\ H}}\skt{ ।}\pend

\pstart \ledsidenote{\passref{6.42.75cd}}\skt{वल्लूरम्मूलकं }\edtext{\skt{मत्स्यां}}{\Afootnote{K; \skt{मत्स्या}\ A; \skt{मत्स्यं}\ H}}\skt{ }\edtext{\skt{शुष्कशाकानि}}{\Afootnote{K H; \skt{न्शुष्कशाकानि}\ A}}\skt{ वैदलं ।}\pend

\pstart \ledsidenote{\passref{6.42.76ab}}\skt{न }\edtext{\skt{खादेदालुकं}}{\Afootnote{A K; \skt{खादेच्चालूकं}\ H}}\skt{ गुल्मी मधुराणि फलानि च ।}\pend

\pstart\versepar \ledsidenote{\passref{6.42.76.add1}}\edtext{\skt{लङ्घनं}}{\Afootnote{K; \skt{लघ्नननन्}\ H}}\skt{ दीपनं }\edtext{\skt{स्निग्धं}}{\Afootnote{K; \skt{स्निग्ध}\ H}}\skt{ उष्णं वातानुलोमनं । }\pend
\pstart\versepar \skt{वृंहणञ्च भवेद्यस्तु तधितं वातगुल्मिनां ।}\pend

\pstart\versepar \ledsidenote{\passref{6.42.76.add2}}\skt{स्थानावसेको रक्तस्य वाहुमध्ये सिरांव्यधन् । }\pend
\pstart\versepar \skt{स्वेदोनुलोमनं चैव तधितं वातगुल्मिनां ।}\edtext{}{\lemma{\skt{सिरांव्यधन् । स्वेदोनुलोमनं}}\Afootnote{K; \skt{सिराव्यधं । स्वेदो\-नु\-लो\-म\-नञ्}\ H}}\pend

\pstart\versepar \ledsidenote{\passref{6.42.76.add3}}\skt{तृष्णा }\edtext{\skt{याभिपरीतो यो}}{\Afootnote{K; \skt{याभिपरी तोयः}\ H}}\skt{ }\edtext{\skt{आध्या\-तो\-र\-क्त\-पी\-डितः}}{\Afootnote{K; \skt{आध्यातो भक्तपीडितः ।}\ H}}\skt{ । }\pend
\pstart\versepar \skt{}\edtext{\skt{प्राणा\-क्षि\-शो\-ति\-सारी च}}{\Afootnote{K; \skt{प्राण\-क्षी\-शोऽ\-ति\-सारी च यश्}\ H}}\skt{ }\edtext{\skt{यशूलैः}}{\Afootnote{K; \skt{च शूलैः}\ H}}\skt{ }\edtext{\skt{प्रदीद्यते}}{\Afootnote{K; \skt{प्रपीद्यते ॥}\ H}}\skt{ ।}\pend

\pstart \ledsidenote{\passref{6.42.76.add4}}\skt{वर्जनीया भवेन्ह्येते }\edtext{\skt{विरि\-क्के\-यो\-भि\-पू\-ज्यते}}{\Afootnote{K; \skt{विरि\-क्के\-योऽ\-भि\-पू\-ज्यते}\ H}}\skt{ इति ॥}\pend

\pstart \ledsidenote{\passref{6.42.trailer}}\edtext{\skt{कायचि}}{\Afootnote{K; \skt{इति सुश्रु\-त\-सं\-हि\-ता\-या\-मु\-त्त\-र\-त\-न्त्रा\-न्त\-र्गते कायचि}\ A; \skt{उत्तर}\ H}}\skt{ ॥ }\edtext{\skt{थ्व}}{\Afootnote{K; \skt{कित्सातन्त्रे गुल्मप्रतिषेधो नाम (चतुर्थोऽध्यायः, आदितः) द्विच\-त्वा\-रिं\-शोऽ\-ध्यायः ॥४२ ॥}\ A; \skt{तन्त्रे गुल्मप्रतिषेध न चत्वारिंशत्तमः कायचिकित्सायां चतुर्थोऽद्यायः ॥ ॥}\ H}}\skt{ ॥ ० ॥}\pend

\endnumbering
\beginnumbering
\eledsection*{Adhyāya 43}
\pstart \ledsidenote{\passref{6.43.1}}\skt{अथातो हृद्रो\-ग\-प्र\-ति\-षेधं }\edtext{\skt{व्याख्यास्यामः}}{\Afootnote{A H; \skt{व्या ॥}\ K}}\skt{ ॥}\pend

\pstart\versepar \ledsidenote{\passref{6.43.3}}\skt{वेगा\-घा\-तो\-ष्ण\-रू\-क्षा\-न्नै\-र\-ति\-मा\-त्रो\-प\-से\-वितैः । }\pend
\pstart\versepar \skt{}\edtext{\skt{विरु\-द्धा\-ध्य\-श\-ना\-जी\-र्णै\-र्र\-सा\-त्म्यै\-श्चापि}}{\lemma{\skt{विरु\-द्धा\-ध्य\-श\-ना\-जी\-र्णैर्}}\Afootnote{K\textsuperscript{pc}; \skt{विरु\-द्धा\-ध्य\-श\-ना\-जीर्णै}\ A; \skt{विरु\-द्धा\-ध्य\-ना\-जी\-र्णैर्}\ K\textsuperscript{ac}; \skt{विरू\-द्धा\-द्य\-श\-ना\-ज्जीर्णै}\ H}}\skt{ }\edtext{\skt{भोजनैः}}{\Afootnote{K H; \skt{(चाति) भोजनैः ॥}\ A}}\skt{ ।}\pend

\pstart\versepar \ledsidenote{\passref{6.43.4}}\skt{दूषयित्वा }\edtext{\skt{रसां}}{\Afootnote{K; \skt{रसं}\ A H}}\skt{ दोषा विगुणा }\edtext{\skt{हृदं यं}}{\Afootnote{K; \skt{हृदयं}\ A H}}\skt{ }\edtext{\skt{गताः}}{\Afootnote{K H; \skt{गताः ।}\ A}}\skt{ । }\pend
\pstart\versepar \skt{}\edtext{\skt{हृदिवाधां}}{\Afootnote{K; \emph{om.}\ A; \skt{हृदि बाधाम्}\ H}}\skt{ }\edtext{\skt{प्रकुव्यन्ति}}{\Afootnote{K; \skt{कुर्वन्ति}\ A; \skt{प्रकुर्व्वन्ति}\ H}}\skt{ }\edtext{\skt{हृद्रोगां तं}}{\Afootnote{K; \skt{हृदये बाधां हृद्रोगं तं}\ A; \skt{हृद्रोगन्तम्}\ H}}\skt{ प्रचक्षते ।}\pend

\pstart\versepar \ledsidenote{\passref{6.43.6}}\skt{आयम्यते मा॥॥॥ }\pend
\pstart\versepar \skt{}\edtext{\skt{निर्मथ्यते}}{\Afootnote{K; \skt{रुतजे हृदयं तुद्यते तथा । निर्मथ्यते}\ A; \skt{रुतेन हृदयन्तुद्यते तथा । निर्म्मथ्यते}\ H}}\skt{ दीर्यते च }\edtext{\skt{योद्यतेपि}}{\Afootnote{K; \skt{स्फोट्यते}\ A; \skt{…}\ H}}\skt{ }\edtext{\skt{पाट्यतेपि}}{\Afootnote{K; \skt{पाट्यतेऽपि}\ A; \skt{पोट्यतेपि}\ H}}\skt{ च ।}\pend

\pstart\versepar \ledsidenote{\passref{6.43.7}}\edtext{\skt{तृष्णो\-ष्णा\-दा\-ह\-चो\-षास्युः}}{\Afootnote{K; \skt{तृष्णो\-षा\-दा\-ह\-चोषाः स्युः}\ A; \skt{तृष्णो\-ष्णा\-दा\-ह\-चोषाः स्युः}\ H}}\skt{ पैत्तिके }\edtext{\skt{हृदयक्लमः}}{\Afootnote{A K; \skt{हृदय क्लमाः ।}\ H}}\skt{ । }\pend
\pstart\versepar \skt{धूमायनं च मूर्छा च स्वेदः शोषो मुखस्य च ।}\pend

\pstart\versepar \ledsidenote{\passref{6.43.8}}\skt{गौरवं }\edtext{\skt{कफ\-सं\-स्रा\-वो\-रू\-चि\-स्त\-स्त\-म्भो\-ग्नि\-मा\-र्दवं}}{\Afootnote{K; \skt{कफ\-सं\-स्रा\-वोऽ\-रुचिः स्तम्भोऽ\-ग्नि\-मा\-र्द\-वम् ।}\ A; \skt{कफसंस्रावो रुचि स्तम्भोऽ\-ग्नि\-मा\-र्द्द\-वम् ।}\ H}}\skt{ । }\pend
\pstart\versepar \skt{माधुर्यमपि चा }\edtext{\skt{सस्य}}{\Afootnote{K H; \skt{स्यस्य}\ A}}\skt{ }\edtext{\skt{वलासापहते}}{\Afootnote{K; \skt{बलासावतते}\ A; \skt{बलासापहते}\ H}}\skt{ हृदि ॥}\pend

\pstart\versepar \ledsidenote{\passref{6.43.9}}\skt{उत्क्लेशः ष्ठीवनं }\edtext{\skt{तोदः}}{\Afootnote{A K; \skt{तो दः}\ H}}\skt{ }\edtext{\skt{शूलं}}{\Afootnote{K H; \skt{शूलो}\ A}}\skt{ }\edtext{\skt{हृल्लासकोवमी}}{\Afootnote{K; \skt{हृल्लासकस्तमः ।}\ A; \skt{हृल्लासकोचमी ।}\ H}}\skt{ । }\pend
\pstart\versepar \skt{अरुचिः }\edtext{\skt{श्यावनेत्रत्वं}}{\Afootnote{K H; \skt{श्याव्नेत्रत्वं}\ A}}\skt{ }\edtext{\skt{शोषश्च}}{\Afootnote{A K; \skt{शोथश्च}\ H}}\skt{ }\edtext{\skt{क्रिमिजो}}{\Afootnote{K; \skt{कृमिजे}\ A; \skt{कृमिजो}\ H}}\skt{ भवेत् ।}\pend

\pstart\versepar \ledsidenote{\passref{6.43.10}}\edtext{\skt{भ्रमक्लमौ}}{\Afootnote{A K; \skt{द्रमक्लमो}\ H}}\skt{ }\edtext{\skt{मोदशोषौ}}{\Afootnote{K; \skt{सादशोषौ}\ A; \skt{सादशोथौ}\ H}}\skt{ }\edtext{\edtext{\skt{ज्ञेया\-स्ते\-षा\-मु\-प\-द्रवाः}}{\lemma{\skt{ज्ञेयास्}}\Afootnote{A K; \skt{ज्ञेया}\ H}}}{\lemma{\skt{तेषाम्}}\Afootnote{A K; \skt{स्तेषाम्}\ H}}\skt{ । }\pend
\pstart\versepar \skt{}\edtext{\skt{क्रिमिजो}}{\Afootnote{K; \skt{कृमिजे}\ A; \skt{क्रिमिजे}\ H}}\skt{ }\edtext{\skt{क्रिमिजातीनां}}{\Afootnote{K H; \skt{कृमिजातीनां}\ A}}\skt{ }\edtext{\skt{श्लेमिकाणाश्च यो}}{\Afootnote{K; \skt{श्लैष्मिकाणां च ये}\ A; \skt{श्लेष्मिकाणां च ये}\ H}}\skt{ मताः ॥}\pend

\pstart\versepar \ledsidenote{\passref{6.43.11}}\skt{वातोपसृष्टे हृदये वाम\-ये\-त्स्नि\-ग्ध\-मा\-तुरं । }\pend
\pstart\versepar \skt{}\edtext{\skt{द्विप\-ञ्च\-मू\-ली\-क्वा\-थेन}}{\Afootnote{K H; \skt{द्विप\-ञ्च\-मू\-ल\-क्वा\-थेन}\ A}}\skt{ सस्ने\-ह\-ल\-व\-णेन तु ॥}\pend

\pstart\versepar \ledsidenote{\passref{6.43.12}}\edtext{\skt{पिप्प\-य़े\-ला\-व\-चा\-हिं\-गु\-य\-व\-भ\-स्मानि}}{\Afootnote{K; \skt{पिप्प\-ल्ये\-ला\-व\-चा\-हि\-ङ्गु\-य\-व\-भ\-स्मानि}\ A H}}\skt{ }\edtext{\skt{सैन्धवं}}{\Afootnote{A K; \skt{सैन्धवः ।}\ H}}\skt{ । }\pend
\pstart\versepar \skt{}\edtext{\skt{सौर्वर्चलं}}{\Afootnote{K; \skt{सौवर्चल}\ A; \skt{सौवर्च्चलं}\ H}}\skt{ }\edtext{\skt{तथा}}{\Afootnote{K H; \skt{मथो}\ A}}\skt{ }\edtext{\skt{शुन्ठी}}{\Afootnote{K; \skt{शुण्ठी}\ A H}}\skt{ }\edtext{\skt{मजोमोदश्च}}{\Afootnote{K; \skt{मजमोदां च}\ A; \skt{मजमोदं च}\ H}}\skt{ }\edtext{\skt{चूर्णितं}}{\Afootnote{A K; \skt{चूर्ण्णितः ॥}\ H}}\skt{ ।}\pend

\pstart\versepar \ledsidenote{\passref{6.43.13}}\skt{फल\-धा\-न्या\-म्ल\-कौ\-ल\-त्थ\-द\-धि\-म\-द्या\-स\-वा\-दिभिः । }\pend
\pstart\versepar \skt{}\edtext{\skt{पाययेत्च}}{\Afootnote{K; \skt{पाययेत}\ A; \skt{पायये}\ H}}\skt{ }\edtext{\skt{शुद्धदेहश्च}}{\Afootnote{K; \skt{विशुद्धं च}\ A; \skt{च्छुद्धदेहञ्च}\ H}}\skt{ }\edtext{\skt{स्नेह\-ना\-न्य\-त\-मेन}}{\Afootnote{K; \skt{स्नेहे\-ना\-न्य\-त\-मेन}\ A H}}\skt{ वा ।}\pend

\pstart\versepar \ledsidenote{\passref{6.43.14}}\edtext{\skt{भोज\-ये\-ज्जी\-र्ण\-भै\-षज्यां}}{\Afootnote{K; \skt{भोज\-ये\-ज्जी\-र्ण\-शा\-ल्यन्नं}\ A; \skt{भोज\-ये\-ज्जी\-र्ण्ण\-भै\-षज्यं}\ H}}\skt{ }\edtext{\skt{जंगलैसघृतै}}{\Afootnote{K; \skt{जाङ्गलैः सघृतै}\ A H}}\skt{ रसैः । }\pend
\pstart\versepar \skt{वातघ्नसिद्धं तैलञ्च }\edtext{\skt{वस्तिं}}{\Afootnote{K; \skt{दद्या}\ A H}}\skt{ }\edtext{\skt{दद्याद्विचक्षणाः}}{\Afootnote{K; \skt{द्बस्तिं प्रमाणतः ॥}\ A; \skt{द्बस्ति विचकषणः ॥}\ H}}\skt{ ।}\pend

\pstart\versepar \ledsidenote{\passref{6.43.15}}\edtext{\skt{श्रीप\-णि\-म\-धु\-क\-क्षौ\-द्र\-सि\-ता\-गु\-ड\-ज\-लै\-र्व\-मेत्}}{\lemma{\skt{श्रीप\-णि\-म\-धु\-क\-क्षौ\-द्र\-सि\-ता\-गु\-ड\-ज\-लैर्}}\Afootnote{K; \skt{श्रीप\-र्णी\-म\-धु\-क\-क्षौ\-द्र\-सि\-तो\-त्प\-ल\-जलै}\ A; \skt{श्रीप\-र्ण्णी\-म\-धु\-क\-क्षौद्र सितागुडजलै}\ H}}\skt{ । }\pend
\pstart\versepar \skt{पित्तोपसृष्टे }\edtext{\skt{हृदयेश्च}}{\Afootnote{K; \skt{हृदये}\ A H}}\skt{ }\edtext{\skt{सेवेत}}{\Afootnote{A K; \skt{सेवते}\ H}}\skt{ }\edtext{\skt{मधुरश्रितं}}{\Afootnote{K H; \skt{मधुरैः शृतम् ॥}\ A}}\skt{ ।}\edtext{}{\lemma{\skt{वमेत् । पित्तोपसृष्टे}}\Afootnote{K; \skt{र्वमेत् । पित्तोपसृष्टे}\ A; \skt{र्व्वमेत् । पित्तोपसृष्ट}\ H}}\pend

\pstart\versepar \ledsidenote{\passref{6.43.16}}\skt{घृतं }\edtext{\edtext{\skt{कषायश्चोद्दिष्टं}}{\lemma{\skt{कषायश्}}\Afootnote{K; \skt{कषायांश्}\ A; \skt{क षायां}\ H}}}{\lemma{\skt{चोद्दिष्टं}}\Afootnote{K H; \skt{चोद्दिष्टा}\ A}}\skt{ }\edtext{\skt{पित्त\-ज्वा\-र\-वि\-ना\-शनं}}{\Afootnote{K; \skt{न्पित्त\-ज्व\-र\-वि\-ना\-श\-नान् ।}\ A; \skt{पित्त\-ज्व\-र\-वि\-ना\-श\-नाम् ।}\ H}}\skt{ }\pend
\pstart\versepar \skt{तृप्तस्य च }\edtext{\edtext{\skt{रसौ\-र्मु\-ख्यै\-र्जं\-गलैः}}{\lemma{\skt{रसौर्मुख्यैर्}}\Afootnote{K; \skt{रसैर्मुख्यै}\ A; \skt{रसैर्म्मुख्यैर्}\ H}}}{\lemma{\skt{जंगलैः}}\Afootnote{K; \skt{र्मधुरैः}\ A; \skt{जांगलैः}\ H}}\skt{ }\edtext{\skt{सघृतैभिषक्}}{\Afootnote{K; \skt{सघृतैर्भिषक् ॥}\ A; \skt{सघृतैर्भ्भिषक् ॥}\ H}}\skt{ ।}\pend

\pstart\versepar \ledsidenote{\passref{6.43.17}}\skt{सक्षौद्रं }\edtext{\skt{वित\-रे\-द्व\-स्ति\-तैलं}}{\Afootnote{K; \skt{वितरेद्बस्तौ तैलं}\ A; \skt{वित\-रे\-ब\-स्ति\-म्तैलं}\ H}}\skt{ }\edtext{\skt{मधुकसाधितं}}{\Afootnote{A K; \skt{मधु\-क\-फा\-धि\-कम् ।}\ H}}\skt{ । }\pend
\pstart\versepar \skt{}\edtext{\skt{वचा\-नि\-म्व\-क\-षा\-याभ्यां}}{\Afootnote{K; \skt{वचा\-नि\-म्ब\-क\-षा\-याभ्यां}\ A; \skt{वचा\-नि\-म्ब\-के\-षा\-याभ्यां}\ H}}\skt{ वान्तं हृदि }\edtext{\skt{कफोत्थिते}}{\Afootnote{K H; \skt{कफात्मके ॥}\ A}}\skt{ ।}\pend

\pstart\versepar \ledsidenote{\passref{6.43.18}}\edtext{\skt{पाययेत्चूर्णं}}{\Afootnote{K; \skt{चूर्णं तु पाययेतो}\ A; \skt{पाययेच्चूर्ण्ण}\ H}}\skt{ }\edtext{\skt{उद्दिष्टं}}{\Afootnote{K; \skt{क्तं}\ A; \skt{मुद्दिष्टम्}\ H}}\skt{ }\edtext{\skt{वातके}}{\Afootnote{K; \skt{वातजे}\ A; \skt{वातजो}\ H}}\skt{ }\edtext{\skt{भोजयेत्च्तम्}}{\Afootnote{K; \skt{भोजयेच्च तम् ।}\ A H}}\skt{ । }\pend
\pstart\versepar \skt{}\edtext{\skt{वलातैलं}}{\Afootnote{K; \skt{फला}\ A; \skt{वत्मा तैलम्}\ H}}\skt{ }\edtext{\skt{विदद्याच्च}}{\Afootnote{K; \emph{om.}\ A; \skt{विद्ध्याच्च}\ H}}\skt{ }\edtext{\skt{वस्ति\-व\-स्ति\-वि\-शा\-रदः}}{\Afootnote{K; \skt{दिमथ मुस्तादिं त्रिफलां वा पिबेन्नरः ॥}\ A; \skt{बस्तिं बस्तिविशारदाः ॥}\ H}}\skt{ ।}\pend

\pstart\versepar \ledsidenote{\passref{6.43.19}}\edtext{\skt{क्रिमिहृद्रोगिणां}}{\Afootnote{K; \emph{om.}\ A}}\skt{ }\edtext{\skt{स्निग्धं}}{\Afootnote{K; \emph{om.}\ A}}\skt{ }\edtext{\skt{भोज\-त्पि\-शि\-तो\-दनं}}{\Afootnote{K; \skt{श्यामात्रि}\ A}}\skt{ । }\pend
\pstart\versepar \skt{}\edtext{\skt{दघ्ना च}}{\Afootnote{K; \emph{om.}\ A}}\skt{ }\edtext{\skt{फललोपेतं}}{\Afootnote{K; \emph{om.}\ A}}\skt{ }\edtext{\skt{त्र्यहं}}{\Afootnote{K; \emph{om.}\ A}}\skt{ }\edtext{\skt{पश्चा\-द्वि\-रे\-च\-येत्}}{\Afootnote{K; \skt{वृत्कल्कयुतं घृतं वापि विरेचनम् । बला\-तै\-लै\-र्वि\-द\-ध्याच्च बस्तिं बस्तिविशारदः ॥}\ A}}\skt{ ॥}\pend

\pstart\versepar \ledsidenote{\passref{6.43.20}}\edtext{\skt{सुगन्धिभिः}}{\Afootnote{K; \emph{om.}\ A H}}\skt{ }\edtext{\edtext{\skt{शलवणैर्योगैः}}{\lemma{\skt{शलवणैर्}}\Afootnote{K; \emph{om.}\ A H}}}{\lemma{\skt{योगैः}}\Afootnote{K; \emph{om.}\ A H}}\skt{ }\edtext{\skt{साजाजिशर्करैः}}{\Afootnote{K; \emph{om.}\ A H}}\skt{ । }\pend
\pstart\versepar \skt{}\edtext{\skt{विडं\-ग\-गा\-ढै\-धा\-न्याम्लं}}{\Afootnote{K; \emph{om.}\ A H}}\skt{ }\edtext{\edtext{\skt{पाय\-ये\-धि\-त\-मु\-त्तमः}}{\lemma{\skt{पाययेधितम्}}\Afootnote{K; \emph{om.}\ A H}}}{\lemma{\skt{उत्तमः}}\Afootnote{K; \skt{कृमिहृद्रोगिणं स्निग्धं भोज\-ये\-त्पि\-शि\-तौ\-द\-नम् । दध्ना च पललोपेतं त्र्यहं पश्चा\-द्वि\-रे\-च\-येत् ॥}\ A; \skt{कृमिहृद्रोगिणां स्निग्ध\-म्भो\-ज\-ये\-त्पि\-शि\-तो\-द\-नम् । दध्ना च पल\-लो\-पे\-त\-म्त्र्यहं पश्चा\-द्वि\-रे\-च\-येत् ॥}\ H}}\skt{ ।}\pend

\pstart\versepar \ledsidenote{\passref{6.43.21}}\edtext{\skt{हृदिस्थिताः}}{\Afootnote{K; \skt{सुगन्धिभिः}\ A H}}\skt{ }\edtext{\skt{पतन्त्येव}}{\Afootnote{K; \skt{सलव}\ A H}}\skt{ }\edtext{\skt{मधस्तात्क्रिमयो}}{\Afootnote{K; \skt{णैर्योगैः साजा}\ A H}}\skt{ }\edtext{\skt{नृणां}}{\Afootnote{K; \emph{om.}\ A H}}\skt{ । }\pend
\pstart\versepar \skt{}\edtext{\skt{यवान्नं}}{\Afootnote{K; \emph{om.}\ A H}}\skt{ }\edtext{\skt{वित\-रे\-च्चा\-स्मै\-स\-वि\-डं\-ग\-मतः}}{\Afootnote{K; \skt{जिशर्करैः । विडङ्गगाढं}\ A; \skt{जिशर्क्करैः । विडंगगाढैर्}\ H}}\skt{ }\edtext{\edtext{\skt{परमिति}}{\lemma{\skt{परम्}}\Afootnote{K; \skt{धान्याम्लं पाय\-ये\-ता\-प्य\-न\-न्त\-रम् ॥}\ A; \skt{धान्या\-म्ल\-म्पा\-य\-ये\-द्धि\-अ\-मु\-लम् ॥}\ H}}}{\lemma{\skt{इति}}\Afootnote{K; \emph{om.}\ A H}}\skt{ ॥ }\edtext{\skt{कायचि}}{\Afootnote{K; \emph{om.}\ A H}}\skt{ ॥ }\edtext{\skt{कृ}}{\Afootnote{K; \emph{om.}\ A H}}\skt{ ॥ ॥}\pend

\endnumbering
\beginnumbering
\eledsection*{Adhyāya 44}
\pstart \ledsidenote{\passref{6.44.1}}\skt{अथातः }\edtext{\skt{पाण्डु\-रो\-ग\-प्र\-ति\-षेधं}}{\Afootnote{A K; \skt{पाण्डु\-र्प्ग\-प्र\-ति\-षेधं}\ H}}\skt{ }\edtext{\skt{व्याख्यास्यामः}}{\Afootnote{A H; \skt{व्या ॥}\ K}}\skt{ ॥}\pend

\pstart\versepar \ledsidenote{\passref{6.44.3}}\edtext{\skt{व्याया\-म\-म\-म्ल\-म्ल\-व\-णानि}}{\lemma{\skt{व्यायाममम्लम्}}\Afootnote{K H; \skt{व्यवायमम्लं}\ A}}\skt{ मद्यं मृदं दिवा\-स्व\-प्न\-म\-ती\-व\-तीक्ष्णं । }\pend
\pstart\versepar \skt{}\edtext{\skt{निषेव्यमाणस्य}}{\Afootnote{K H; \skt{निषेवमाणस्य}\ A}}\skt{ विदूष्य रक्तं कुर्वन्ति }\edtext{\edtext{\skt{दोषास्त्वचि}}{\lemma{\skt{दोषास्}}\Afootnote{K; \skt{दोषा}\ A H}}}{\lemma{\skt{त्वचि}}\Afootnote{K; \skt{स्त्वचि}\ A H}}\skt{ पाण्डुभावं ।}\pend

\pstart\versepar \ledsidenote{\passref{6.44.4}}\edtext{\skt{पाण्ड्वा\-म\-यो\-ष्टा\-र्ध\-विधः}}{\Afootnote{K; \skt{पाण्ड्वा\-म\-योऽ\-ष्टा\-र्ध\-विधः}\ A; \skt{पाण्द्वा\-म\-यो\-ष्टा\-र्द्ध\-विधः}\ H}}\skt{ प्रदिष्टः }\edtext{\skt{पृथ\-क्स\-म\-स्तै\-यु\-ग\-पच्च}}{\Afootnote{K; \skt{पृथ\-क्स\-म\-स्तै\-र्यु\-ग\-पच्च}\ A; \skt{पृथ\-क्स\-म\-स्तै\-र्यु\-ग\-पच्च}\ H}}\skt{ }\edtext{\skt{दोसैः}}{\Afootnote{K; \skt{दोषैः ।}\ A H}}\skt{ । }\pend
\pstart\versepar \skt{सर्वेषु }\edtext{\skt{चैव\-ष्वि\-ह\-पा\-ण्ड\-भावो}}{\Afootnote{K H; \skt{चैतेष्विह पाण्डुभावो}\ A}}\skt{ }\edtext{\skt{यतोधिकोतः}}{\Afootnote{K; \skt{यतोऽधिकोऽतः}\ A; \skt{यतोऽधिकोतः}\ H}}\skt{ खलु पाण्डुरोगः ॥}\pend

\pstart\versepar \ledsidenote{\passref{6.44.5}}\skt{त्वक्स्फो\-ट\-न\-ष्ठी\-व\-न\-गात्र }\pend
\pstart\versepar \skt{विण्मू\-त्र\-पी\-त\-त्व\-म\-था\-वि\-पाको }\edtext{\edtext{\skt{भविष्यतस्तस्य}}{\lemma{\skt{भविष्यतस्}}\Afootnote{A K; \skt{भविष्यत}\ H}}}{\lemma{\skt{तस्य}}\Afootnote{A K; \skt{स्तस्य}\ H}}\skt{ पुरः सराणि ।}\edtext{}{\lemma{\skt{त्वक्स्फो\-ट\-न\-ष्ठी\-व\-न\-गात्र विण्मू\-त्र\-पी\-त\-त्व\-म\-था\-वि\-पाको}}\Afootnote{\skt{त्वक्स्फोटनं ष्ठीव\-न\-गा\-त्र\-सादौ मृद्भक्षणं प्रेक्ष\-ण\-कू\-ट\-शोथः । विण्मू\-त्र\-पी\-त\-त्व\-म\-था\-वि\-पाको}\ A; \skt{त्वक्स्फो\-ट\-न\-ष्ठी\-व\-न\-गात्र… विण्मू\-त्र\-पी\-त\-त्व\-म\-था\-वि\-पाको}\ K; \skt{त्वक्स्फो\-ट\-ष्ठी\-व\-न\-गा\-त्र\-साद मृद्भ\-क्ष\-ण\-प्रे\-क्ष\-ण\-कू\-ट\-शोथाः । विण्मू\-त्र\-पी\-त\-त्व\-म\-था\-वि\-पाको}\ H}}\pend

\pstart\versepar \ledsidenote{\passref{6.44.6}}\edtext{\skt{मला\-पा\-न्न\-कि\-पा\-ण्डु\-रोगः}}{\Afootnote{\skt{स काम\-ला\-पा\-न\-कि\-पा\-ण्डु\-रोगः}\ A; \skt{…म\-ला\-पा\-न्न\-कि\-पा\-ण्डु\-रोगः ।}\ K; \skt{स काम\-ला\-पा\-न्न\-कि\-पा\-ण्डु\-रोगः}\ H}}\skt{ । }\pend
\pstart\versepar \skt{कुम्भाह्वयो }\edtext{\skt{लाक\-र\-को\-ल\-साख्यः}}{\Afootnote{K H; \skt{लाघ\-र\-को(\-ला\-घ\-व\-को)+\-अ\-ल\-साख्यः ।}\ A}}\skt{ । }\pend
\pstart\versepar \skt{}\edtext{\skt{विभास्यते}}{\Afootnote{K; \skt{विभाष्यते}\ A H}}\skt{ लक्षणमस्य कृत्स्नं }\edtext{\skt{निवोध्य}}{\Afootnote{K H; \skt{निबोध}\ A}}\skt{ }\edtext{\skt{वक्ष्या\-म्नु\-नु\-पू\-र्व\-शस्ते}}{\Afootnote{K; \skt{वक्ष्या\-म्य\-नु\-पू\-र्व\-श\-स्तत् ॥}\ A; \skt{वक्ष्या\-म्नु\-पू\-र्व्व\-शस्ते ॥}\ H}}\skt{ ।}\pend

\pstart\versepar \ledsidenote{\passref{6.44.7}}\edtext{\skt{कृष्नेक्षणं}}{\Afootnote{K; \skt{कृष्णेक्षणं}\ A H}}\skt{ कृष्ण\-सि\-रा\-व\-नद्धं }\edtext{\skt{तद्व\-र्ण\-वि\-ण्मू\-त्र\-न\-खा\-न\-नञ्च}}{\Afootnote{K H; \skt{तद्व\-र्ण\-वि\-ण्मू\-त्र\-न\-खा\-ननं च ।}\ A}}\skt{ । }\pend
\pstart\versepar \skt{वातेन पाण्डुं मनुजं }\edtext{\edtext{\skt{व्यव\-स्ये\-द्यु\-क्त\-त\-था\-न्यै\-स्त\-दू\-प\-द्र\-वैश्च}}{\lemma{\skt{व्यव\-स्ये\-द्यु\-क्त\-त\-था\-न्यैस्}}\Afootnote{K; \skt{व्यवस्येद्युक्तं तथान्यैस्}\ A; \skt{व्यवस्ये द्युक्तन्तथान्यैस्}\ H}}}{\lemma{\skt{तदूपद्रवैश्}}\Afootnote{K; \skt{तदुपद्रवैश्}\ A; \skt{तदुपद्रवैश्}\ H}}\skt{ ।}\pend

\pstart\versepar \ledsidenote{\passref{6.44.8}}\skt{पीतेक्षणं पीत\-सि\-रा\-व\-नद्धं }\edtext{\skt{तद्व\-र्ण\-वि\-ण्मू\-त्र\-न\-खा\-न\-नञ्च}}{\Afootnote{K; \skt{तद्व\-र्ण\-वि\-ण्मू\-त्र\-न\-खा\-ननं च ।}\ A; \skt{तद्व\-र्ण्ण\-वि\-ण्मू\-त्र\-न\-खा\-नन्ञ्च ।}\ H}}\skt{ । }\pend
\pstart\versepar \skt{पित्तेन पाण्डुम्मनुजं }\edtext{\edtext{\skt{व्यव\-स्ये\-द्यु\-क्त\-त\-था\-न्यै\-स्त\-दु\-प\-द्रैश्च}}{\lemma{\skt{व्यव\-स्ये\-द्यु\-क्त\-त\-था\-न्यैस्}}\Afootnote{K; \skt{व्यवस्येद्युक्तं तथान्यैस्}\ A; \skt{व्यवस्ये द्युक्तंन्तथान्यैस्}\ H}}}{\lemma{\skt{तदुपद्रैश्}}\Afootnote{K; \skt{तदुपद्रवैश्}\ A; \skt{तदूपद्रवैश्}\ H}}\skt{ ।}\pend

\pstart\versepar \ledsidenote{\passref{6.44.9}}\skt{शुक्लेक्षणं शुक्ल\-सि\-रा\-व\-नद्धं }\edtext{\skt{तद्व\-र्ण\-वि\-ण्मू\-त्र\-न\-खा\-न\-नञ्च}}{\Afootnote{K H; \skt{तद्व\-र्ण\-वि\-ण्मू\-त्र\-न\-खा\-ननं च ।}\ A}}\skt{ । }\pend
\pstart\versepar \skt{कफेन पाण्डुं मनुजं }\edtext{\skt{व्यवस्येद्युक्तं}}{\Afootnote{A K; \skt{व्यवस्ये ध्युक्तन्}\ H}}\skt{ }\edtext{\edtext{\skt{तेषा\-न्यै\-स्त\-दु\-प\-द्र\-वैश्च}}{\lemma{\skt{तेषान्यैस्}}\Afootnote{K; \skt{तथान्यैस्}\ A H}}}{\lemma{\skt{तदुपद्रवैश्}}\Afootnote{K H; \skt{तदुपद्रवैश्}\ A}}\skt{ ।}\pend

\pstart\versepar \ledsidenote{\passref{6.44.10}}\edtext{\edtext{\skt{अभ\-क्त\-रू\-ग्नै\-र\-व\-गात्र}}{\lemma{\skt{अभक्तरूग्नैर्}}\Afootnote{K H; \emph{om.}\ A}}}{\lemma{\skt{अवगात्र}}\Afootnote{K H; \emph{om.}\ A}}\skt{ }\edtext{\skt{सादाः}}{\Afootnote{K H; \emph{om.}\ A}}\skt{ }\edtext{\skt{सर्वेषु}}{\Afootnote{K H; \emph{om.}\ A}}\skt{ }\edtext{\skt{तन्द्री च}}{\Afootnote{K H; \emph{om.}\ A}}\skt{ }\edtext{\skt{वलक्षयञ्च}}{\Afootnote{K; \emph{om.}\ A; \skt{वलक्सयञ्च ।}\ H}}\skt{ । }\pend
\pstart\versepar \skt{सर्वात्मके }\edtext{\skt{सर्वमिदं}}{\Afootnote{K H; \skt{सर्वमिदं}\ A}}\skt{ }\edtext{\skt{व्यवस्येल्लिङ्गानि}}{\Afootnote{K H; \skt{व्यवस्येद्वक्ष्यामि}\ A}}\skt{ }\edtext{\edtext{\skt{वक्स्̤याम्यथ}}{\lemma{\skt{वक्स्̤याम्य्}}\Afootnote{K; \skt{लिङ्गा}\ A; \skt{वक्ष्या}\ H}}}{\lemma{\skt{अथ}}\Afootnote{K; \skt{न्यथ}\ A; \skt{म्यथ}\ H}}\skt{ }\edtext{\skt{कामलायाः}}{\Afootnote{K H; \skt{कामलायाः । यो ह्यामयान्ते सह\-सा\-न्न\-म\-म्ल\-म\-द्या\-द\-प\-थ्यानि च तस्य पित्तम् ॥}\ A}}\skt{ ।}\pend

\pstart\versepar \ledsidenote{\passref{6.44.11}}\edtext{\skt{यो}}{\Afootnote{K; \emph{om.}\ A H}}\skt{ }\edtext{\skt{ह्यामयान्ते}}{\Afootnote{K; \emph{om.}\ A H}}\skt{ }\edtext{\skt{सह\-सा\-न्न\-म\-म्ल\-म\-द्या\-द\-प\-थ्यानि च}}{\Afootnote{K; \emph{om.}\ A; \skt{सहसान्नमम्ल मद्यादपथ्यानि च}\ H}}\skt{ }\edtext{\skt{तस्य}}{\Afootnote{K H; \emph{om.}\ A}}\skt{ }\edtext{\skt{पित्तं}}{\Afootnote{K H; \emph{om.}\ A}}\skt{ }\pend
\pstart\versepar \skt{करोति }\edtext{\skt{पाण्डुवदनं}}{\Afootnote{K; \skt{पाण्डुं वदनं}\ A H}}\skt{ }\edtext{\skt{विशे\-षा\-त्पू\-र्व\-रितौ}}{\Afootnote{K H; \skt{विशे\-षा\-त्पू\-र्वे\-रितौ}\ A}}\skt{ }\edtext{\skt{तन्द्रिवलक्षयौ}}{\Afootnote{K H; \skt{तन्द्रिबलक्षयौ}\ A}}\skt{ }\edtext{\skt{च}}{\Afootnote{K H; \skt{च । भेदस्तु तस्याः खलु कुम्भसाह्वः शोफो महांस्तत्र च पर्वभेदः ॥}\ A}}\skt{ ।}\pend

\pstart\versepar \ledsidenote{\passref{6.44.12}}\edtext{\skt{भेदस्तु}}{\Afootnote{K H; \emph{om.}\ A}}\skt{ }\edtext{\skt{तस्ये\-त\-र\-कु\-म्भ\-साह्यः}}{\Afootnote{K; \emph{om.}\ A; \skt{तस्येतर कुम्भसाह्वः}\ H}}\skt{ }\edtext{\skt{सोफं}}{\Afootnote{K; \emph{om.}\ A; \skt{शोफो}\ H}}\skt{ }\edtext{\skt{महांस्तत्र च}}{\Afootnote{K; \emph{om.}\ A; \skt{महान्तत्र च}\ H}}\skt{ }\edtext{\skt{पर्वभेदः}}{\Afootnote{K H; \emph{om.}\ A}}\skt{ । }\pend
\pstart\versepar \skt{}\edtext{\skt{ज्वरां\-ग\-म\-र्द\-भ्र\-म\-दा\-ह\-तृष्णा}}{\Afootnote{K H; \skt{ज्वरा\-ङ्ग\-म\-र्द\-भ्र\-म\-सा\-द\-तन्द्रा}\ A}}\skt{ }\edtext{\skt{क्षयान्वितो}}{\Afootnote{A K; \skt{क्षयोन्वितो}\ H}}\skt{ }\edtext{\skt{लाक\-र\-का\-ल\-सेतु}}{\Afootnote{K; \skt{लाघ\-र\-को(\-ला\-घ\-व\-को)+\-अ\-ल\-साख्यः । तं वात\-पि\-त्ता\-द्ध\-रि\-पी\-त\-नीलं हलीमकं नाम वदन्ति तज्ज्ञाः ॥}\ A; \skt{लाक\-र\-का\-ल\-सेत्तु ॥}\ H}}\skt{ ।}\pend

\pstart\versepar \ledsidenote{\passref{6.44.13}}\edtext{\skt{तं}}{\Afootnote{K H; \emph{om.}\ A}}\skt{ }\edtext{\skt{वात\-पि\-त्ता\-ध\-रि\-पी\-त\-नीलं}}{\Afootnote{K H; \emph{om.}\ A}}\skt{ }\edtext{\skt{हलीमकेत्तु}}{\Afootnote{K; \emph{om.}\ A; \skt{हलीमकेति}\ H}}\skt{ }\edtext{\skt{प्रवदन्ति}}{\Afootnote{K H; \emph{om.}\ A}}\skt{ केचित् । }\pend
\pstart\versepar \skt{उप\-द्र\-वा\-स्स्ते\-ष्व\-र\-रूचिः पिपासा }\edtext{\skt{छर्दिर्ज्वराः}}{\Afootnote{K H; \skt{छर्दिर्ज्वरो}\ A}}\skt{ }\edtext{\skt{शीर्ष\-रु\-ज\-ग्नि\-सादः}}{\Afootnote{K; \skt{मूर्ध\-रु\-जा\-ग्नि\-सादः । शोफस्तथा कण्ठ\-ग\-तोऽ\-ब\-लत्वं मूर्च्छा क्लमो हृद्यवपीडनं च ॥}\ A; \skt{शीर्षौजाग्नि सादः ।}\ H}}\skt{ ।}\edtext{}{\lemma{\skt{केचित् । उप\-द्र\-वा\-स्स्ते\-ष्व\-र\-रूचिः}}\Afootnote{K; \skt{उप\-द्र\-वा\-स्ते\-ष्व\-रुचिः}\ A; \skt{केचित् ॥ उप\-द्र\-वा\-स्ते\-ष्व\-रूचिः}\ H}}\pend

\pstart\versepar \ledsidenote{\passref{6.44.14}}\edtext{\skt{सोफ\-स्थि\-रु\-क्क\-ण्ठ\-ग\-तो\-व\-लत्वं}}{\Afootnote{K; \emph{om.}\ A; \skt{शोषो\-स्थि\-रू\-क्क\-ण्ठ\-ग\-तो\-व\-लत्वं}\ H}}\skt{ }\edtext{\skt{मूर्च्छा}}{\Afootnote{K H; \emph{om.}\ A}}\skt{ }\edtext{\skt{कफे}}{\Afootnote{K H; \emph{om.}\ A}}\skt{ }\edtext{\skt{हृद्य\-व\-पी\-ड\-नञ्च}}{\Afootnote{K H; \emph{om.}\ A}}\skt{ । }\pend
\pstart\versepar \skt{}\edtext{\skt{साध्ये तु}}{\Afootnote{K H; \skt{साध्यं तु}\ A}}\skt{ पाण्ड्वामयिनं समीक्ष्य स्निग्धं }\edtext{\skt{घृते\-नो\-र्ध\-म\-धश्च}}{\Afootnote{K; \skt{घृते\-नो\-र्ध्व\-म\-धश्च}\ A; \skt{घृते तूर्द्ध्वमधश्च}\ H}}\skt{ }\edtext{\skt{शुद्धं}}{\Afootnote{K H; \skt{शुद्धम् । संपा\-द\-ये\-त्क्षौ\-द्र\-घृ\-त\-प्र\-गा\-ढै\-र्ह\-री\-त\-की\-चू\-र्ण\-युतैः प्रयोगैः ॥}\ A}}\skt{ ।}\pend

\pstart\versepar \ledsidenote{\passref{6.44.15}}\edtext{\edtext{\skt{सम्पा\-द\-ये\-त्क्षौ\-द्र\-घृ\-त\-प्र\-गा\-ढै\-र्ह\-री\-त\-की\-चू\-र्ण\-युतैः}}{\lemma{\skt{सम्पा\-द\-ये\-त्क्षौ\-द्र\-घृ\-त\-प्र\-गा\-ढैर्}}\Afootnote{K H; \emph{om.}\ A}}}{\lemma{\skt{हरी\-त\-की\-चू\-र्ण\-युतैः}}\Afootnote{K H; \emph{om.}\ A}}\skt{ }\edtext{\skt{प्रयोगैः}}{\Afootnote{K H; \emph{om.}\ A}}\skt{ । }\pend
\pstart\versepar \skt{}\edtext{\skt{पिबेद्घृतं वा}}{\Afootnote{A; \skt{पिवेद्घृतं वा}\ K; \skt{पिवेद्घृतं वा}\ H}}\skt{ }\edtext{\skt{रजनादि}}{\Afootnote{K H; \skt{रजनीवि}\ A}}\skt{ }\edtext{\skt{पक्वंयत्रैफलं}}{\Afootnote{K; \skt{पक्वं यन्त्रैफलं}\ A; \skt{पक्क्वं यत्रैफलन्}\ H}}\skt{ }\edtext{\edtext{\skt{तैल्वकमेव}}{\lemma{\skt{तैल्वकम्}}\Afootnote{A K; \skt{तैल्वक}\ H}}}{\lemma{\skt{एव}}\Afootnote{A K; \skt{मेव}\ H}}\skt{ }\edtext{\skt{चापि}}{\Afootnote{K H; \skt{वापि । विरे\-च\-न\-द्र\-व्य\-कृतं पिबेद्वा योगांश्च वैरे\-च\-नि\-दा\-न्घृ\-तेन ॥}\ A}}\skt{ ।}\pend

\pstart\versepar \ledsidenote{\passref{6.44.16}}\edtext{\skt{विरे\-च\-ना\-द्र\-व्य\-वि\-पा\-चितं वा}}{\Afootnote{K; \emph{om.}\ A; \skt{वैरे\-च\-न\-द्र\-व्य\-वि\-पा\-चित्म्वा}\ H}}\skt{ }\edtext{\skt{योगाश्च}}{\Afootnote{K H; \emph{om.}\ A}}\skt{ }\edtext{\skt{वैरेचनिकां}}{\Afootnote{K; \emph{om.}\ A; \skt{वैरेचनिका}\ H}}\skt{ }\edtext{\skt{घृतेन}}{\Afootnote{K H; \emph{om.}\ A}}\skt{ । }\pend
\pstart\versepar \skt{मूत्रे }\edtext{\skt{निकुम्भार्धनं}}{\Afootnote{K; \skt{निकुम्भार्धपलं}\ A; \skt{निकु\-म्भा\-र्द्ध\-ओ\-अ\-लम्}\ H}}\skt{ विपाच्य }\edtext{\skt{पिबेदभीक्ष्णं}}{\Afootnote{A; \skt{पिवेदभीक्ष्णं}\ K H}}\skt{ }\edtext{\skt{कुडवार्धमात्रं}}{\Afootnote{K H; \skt{कुड\-वा\-र्ध\-मा\-त्रम् । खादेद्गुडं वाप्य\-भ\-या\-वि\-प\-क्व\-मा\-र\-ग्व\-धा\-दि\-क्व\-थितं पिबेद्वा ॥}\ A}}\skt{ ।}\pend

\pstart\versepar \ledsidenote{\passref{6.44.17}}\edtext{\edtext{\skt{खादेद्गुडम्वाप्य}}{\lemma{\skt{खादेद्गुडम्}}\Afootnote{K H; \emph{om.}\ A}}}{\lemma{\skt{वाप्य}}\Afootnote{K H; \emph{om.}\ A}}\skt{ }\edtext{\skt{भयाविपक्व}}{\Afootnote{K H; \emph{om.}\ A}}\skt{ }\edtext{\skt{मार\-ग्व\-धा\-दि\-क्व\-थितं}}{\Afootnote{K; \emph{om.}\ A; \skt{मार\-ग्व\-धा\-धि\-क्व\-थि\-तम्}\ H}}\skt{ }\edtext{\skt{पिबेद्वा}}{\Afootnote{\emph{om.}\ A; \skt{पिवेद्वा ॥}\ K H}}\skt{ ॥ }\pend
\pstart\versepar \skt{}\edtext{\skt{अयो\-र\-ज\-व्यो\-ष\-वि\-द\-ङ्ग\-चूर्णं}}{\Afootnote{K; \skt{अयो\-र\-जो\-व्यो\-ष\-वि\-ड\-ङ्ग\-चूर्णं}\ A; \skt{अयोएअजः व्योष\-वि\-ड\-ङ्ग\-चूर्ण्णं}\ H}}\skt{ }\edtext{\skt{लिह्याधरिद्रां}}{\Afootnote{K H; \skt{लिह्याद्धरिद्रां}\ A}}\skt{ }\edtext{\skt{त्रिफलान्विताम्वा}}{\lemma{\skt{वा}}\Afootnote{K H; \skt{वा । सर्पिर्मधुभ्यां विदधीत वापि शास्त्र\-प्र\-दे\-शा\-भि\-हि\-तांश्च योगान् ॥}\ A}}\skt{ ।}\pend

\pstart\versepar \ledsidenote{\passref{6.44.18}}\edtext{\skt{सर्पिमधुभ्यां}}{\Afootnote{K; \emph{om.}\ A; \skt{सर्प्पिर्म्मधुभ्याम्}\ H}}\skt{ }\edtext{\skt{विदधीत}}{\Afootnote{K H; \emph{om.}\ A}}\skt{ }\edtext{\skt{वापि}}{\Afootnote{K H; \emph{om.}\ A}}\skt{ }\edtext{\skt{शास्त्र\-प्र\-दे\-शा\-भि\-हितां}}{\Afootnote{K; \emph{om.}\ A; \skt{शास्त्र\-प\-दे\-शा\-भि\-हि\-तांश्}\ H}}\skt{ }\edtext{\skt{प्रयोगां}}{\Afootnote{K; \emph{om.}\ A; \skt{च योगाम् ॥}\ H}}\skt{ । }\pend
\pstart\versepar \skt{}\edtext{\skt{हरेच्च}}{\Afootnote{K H; \skt{अह्रेच्च}\ A}}\skt{ }\edtext{\skt{दोषां}}{\Afootnote{K; \skt{दोषा}\ A H}}\skt{ }\edtext{\skt{वहुशोल्पमात्रां}}{\Afootnote{K; \skt{न्बहुशोऽल्पमात्रा}\ A; \skt{न्वकु\-शोऽ\-ल्प\-मा\-त्राञ्}\ H}}\skt{ }\edtext{\skt{श्वयेद्धि}}{\Afootnote{K; \skt{न्श्वयेद्धि}\ A; \skt{च येधि}\ H}}\skt{ }\edtext{\skt{दोषे\-ष्व\-तो\-नि\-र्हृ\-तेषु}}{\Afootnote{K; \skt{दोषे\-ष्व\-ति\-नि\-र्हृ\-तेषु । धात्रीफलानां रसमिक्षुजं च मन्थं पिबेत्क्षौद्रयुतं हिताशी ॥}\ A; \skt{दोषेषु विनिर्हतेषु ।}\ H}}\skt{ ।}\pend

\pstart\versepar \ledsidenote{\passref{6.44.19}}\edtext{\skt{धात्रीफलानां}}{\Afootnote{K H; \emph{om.}\ A}}\skt{ }\edtext{\skt{रसमिक्षुजञ्च}}{\Afootnote{K H; \emph{om.}\ A}}\skt{ }\edtext{\edtext{\skt{मन्थ\-म्पि\-बे\-त्क्षौ\-द्र\-युतं}}{\lemma{\skt{मन्थम्}}\Afootnote{K H; \emph{om.}\ A}}}{\lemma{\skt{पिबेत्क्षौद्रयुतं}}\Afootnote{\emph{om.}\ A; \skt{पिवेत्क्षौद्रयुतं}\ K H}}\skt{ }\edtext{\skt{समांशं}}{\Afootnote{K H; \emph{om.}\ A}}\skt{ । }\pend
\pstart\versepar \skt{उभे }\edtext{\skt{वृहत्यौ}}{\Afootnote{K H; \skt{बृहत्यौ}\ A}}\skt{ }\edtext{\skt{रजाने}}{\Afootnote{K; \skt{रजनीं}\ A; \skt{रजने}\ H}}\skt{ शुकाख्यां }\edtext{\skt{काकादनीञ्चापि}}{\lemma{\skt{काकादनीञ्}}\Afootnote{K; \skt{काकादनीं}\ A H}}\skt{ }\edtext{\skt{सकाकमाची}}{\Afootnote{K H; \skt{सकाकमाचीम् । आदारिबिम्बीं सकदम्बपुष्पीं विपाच्य सर्पि\-र्वि\-प\-चे\-त्क\-षाये ॥}\ A}}\skt{ ।}\pend

\pstart\versepar \ledsidenote{\passref{6.44.20}}\edtext{\skt{आदारि}}{\Afootnote{K H; \emph{om.}\ A}}\skt{ }\edtext{\skt{विम्ख्यो}}{\Afootnote{K; \emph{om.}\ A; \skt{विंख्यो}\ H}}\skt{ }\edtext{\skt{सकदम्वपुष्पौ}}{\Afootnote{K; \emph{om.}\ A; \skt{सकदम्वपुष्पो}\ H}}\skt{ }\edtext{\skt{विपाच्य}}{\Afootnote{K; \emph{om.}\ A; \skt{विपच्य}\ H}}\skt{ }\edtext{\skt{सर्पि\-र्वि\-प\-चे\-त्क\-षाये}}{\Afootnote{K H; \emph{om.}\ A}}\skt{ । }\pend
\pstart\versepar \skt{तत्पाण्डुतां }\edtext{\skt{हन्त्यु\-प\-से\-व्य\-मानं}}{\Afootnote{K H; \skt{हन्त्यु\-प\-यु\-ज्य\-मानं}\ A}}\skt{ }\edtext{\skt{क्षीरेण वा}}{\Afootnote{K H; \skt{क्षीएण वा}\ A}}\skt{ मागधिका }\edtext{\skt{यथाग्निं}}{\Afootnote{K H; \skt{यथाग्नि । हितं च यष्टीमधुजं कषायं चूर्णं समं वा मधु\-ना\-व\-लि\-ह्यात् ॥}\ A}}\skt{ ।}\pend

\pstart\versepar \ledsidenote{\passref{6.44.21}}\edtext{\skt{पिवेच्च}}{\Afootnote{K H; \emph{om.}\ A}}\skt{ }\edtext{\skt{यष्ठीमधुकं}}{\Afootnote{K H; \emph{om.}\ A}}\skt{ }\edtext{\skt{कषायं}}{\Afootnote{K H; \emph{om.}\ A}}\skt{ }\edtext{\skt{चूर्णीकृतम्वा}}{\Afootnote{K; \emph{om.}\ A; \skt{चूर्ण्णकृतम्वा}\ H}}\skt{ }\edtext{\skt{मधुनासमांशं}}{\Afootnote{K; \emph{om.}\ A; \skt{मधुनाशं ॥}\ H}}\skt{ । }\pend
\pstart\versepar \skt{गोमूत्रयुक्तं }\edtext{\skt{त्रिपलादलानां}}{\Afootnote{K; \skt{त्रिफलादलानां}\ A H}}\skt{ दत्वायसं }\edtext{\skt{चूर्ण\-म\-न\-ल्प\-कालं}}{\Afootnote{K H; \skt{चूर्ण\-म\-न\-ल्प\-का\-लम् । प्रवा\-ल\-मु\-क्ता\-ञ्ज\-न\-श\-ङ्ख\-चूर्णं लिह्यात्तथा काञ्च\-न\-गै\-रि\-को\-त्थम् ॥}\ A}}\skt{ ।}\pend

\pstart\versepar \ledsidenote{\passref{6.44.22}}\edtext{\skt{प्रवा\-ल\-मु\-क्ता\-ञ्ज\-न\-शं\-ख\-चूर्णां}}{\Afootnote{K; \emph{om.}\ A; \skt{प्रवा\-ल\-मु\-क्तां\-ज\-न\-शं\-ख\-चूर्ण्णं}\ H}}\skt{ }\edtext{\skt{लिह्यात्तथा}}{\Afootnote{K H; \emph{om.}\ A}}\skt{ }\edtext{\skt{गैरिकजश्च}}{\Afootnote{K; \emph{om.}\ A; \skt{गरिकजं च}\ H}}\skt{ }\edtext{\skt{चूर्णं}}{\Afootnote{K H; \emph{om.}\ A}}\skt{ । }\pend
\pstart\versepar \skt{}\edtext{\skt{आजं}}{\Afootnote{A K; \skt{आद्यं}\ H}}\skt{ }\edtext{\skt{शकृ\-त्स्या\-त्कु\-ड\-व\-प्र\-माणं}}{\Afootnote{A K; \skt{श . त्स्यात्कु\-ड\-व\-प्र\-माणं}\ H}}\skt{ विडं }\edtext{\skt{हरिद्रां}}{\Afootnote{K H; \skt{हरिद्रा}\ A}}\skt{ लवणोत्तमञ्च । पृथक्पलांशानि }\edtext{\edtext{\skt{सम\-ग्र\-मे\-त\-न्चूर्णं}}{\lemma{\skt{समग्रम्}}\Afootnote{K H; \skt{समग्र}\ A}}}{\lemma{\skt{एतन्चूर्णं}}\Afootnote{K; \skt{मेतच्चूर्णं}\ A; \skt{एतत्चूर्ण्णं}\ H}}\skt{ हिताशी मधु\-ना\-व\-लि\-ह्यात् ।}\pend

\pstart\versepar \ledsidenote{\passref{6.44.25}}\skt{सौवर्चलं }\edtext{\skt{हिंगु}}{\Afootnote{A K; \skt{लिङ्गु}\ H}}\skt{ किराततिक्तं कलायमात्राणि }\edtext{\skt{सुखाम्वुना वा}}{\Afootnote{K H; \skt{सुखाम्बुना वा ।}\ A}}\skt{ । }\pend
\pstart\versepar \skt{मूर्वाहरिद्रा मलकं }\edtext{\skt{पिबेद्वा}}{\Afootnote{\skt{च लिह्या}\ A; \skt{पिवेद्वा}\ K H}}\skt{ }\edtext{\skt{स्थितं}}{\Afootnote{K H; \skt{त्स्थितं}\ A}}\skt{ गवां सप्तदिनानि मूत्रे ।}\pend

\pstart\versepar \ledsidenote{\passref{6.44.26}}\edtext{\skt{मूल\-म्व\-ला\-चि\-त्र\-कयोः}}{\Afootnote{K H; \skt{मूलं बलाचित्रकयोः}\ A}}\skt{ }\edtext{\skt{पिबेद्वा}}{\Afootnote{\skt{पिबेद्वा}\ A; \skt{पिवेद्वा}\ K H}}\skt{ }\edtext{\skt{पाण्डामयार्तोक्ष}}{\Afootnote{K H; \skt{पाण्ड्वामयार्तोऽक्ष}\ A}}\skt{ समं हिताशी । }\pend
\pstart\versepar \skt{}\edtext{\skt{सुखाम्वुना वा}}{\Afootnote{K H; \skt{सुखाम्बुना वा}\ A}}\skt{ लवणेन तुल्यं }\edtext{\skt{शिग्राः}}{\Afootnote{K H; \skt{शिग्रोः}\ A}}\skt{ फलं क्षीर\-भु\-जो\-प\-योज्यं ।}\pend

\pstart\versepar \ledsidenote{\passref{6.44.27}}\edtext{\skt{न्यग्रोधवर्गस्य}}{\Afootnote{A K; \skt{त्न्यग्रोधवर्गस्य}\ H}}\skt{ }\edtext{\skt{पिबेत्कषायं}}{\Afootnote{A; \skt{पिवेत्कषायं}\ K H}}\skt{ शीतं सिताक्षौद्रयुतं हिताशी । }\pend
\pstart\versepar \skt{}\edtext{\edtext{\skt{शाला\-दि\-क\-म्चा\-प्यथ}}{\lemma{\skt{शालादिकम्}}\Afootnote{A K; \skt{सालादिकं}\ H}}}{\lemma{\skt{चाप्यथ}}\Afootnote{K H; \skt{चाप्यथ}\ A}}\skt{ }\edtext{\skt{सालचूर्णं}}{\Afootnote{K H; \skt{सारचूर्णं}\ A}}\skt{ }\edtext{\skt{क्षौद्राप्लुतम्वा}}{\Afootnote{K H; \skt{धात्रीफलं वा}\ A}}\skt{ }\edtext{\skt{मलकं}}{\Afootnote{K H; \skt{म}\ A}}\skt{ }\edtext{\skt{सुखार्थी}}{\Afootnote{K H; \skt{धुनावलिह्यात् ॥}\ A}}\skt{ ॥}\pend

\pstart\versepar \ledsidenote{\passref{6.44.28}}\edtext{\skt{विडं\-ग\-मु\-स्ता\-र\-जा\-ना\-ज\-मो\-द\-प\-रू\-ष\-क\-व्यो\-ष\-वि\-नि\-र्द\-हन्यः}}{\Afootnote{K; \skt{विड\-ङ्ग\-मु\-स्त\-त्रि\-फ\-ला\-ज\-मो\-द\-प\-रू\-ष\-क\-व्यो\-ष\-वि\-नि\-र्द\-हन्यः ।}\ A; \skt{विडङ्गमुस्तो रजनाजमोद परू\-ष\-का\-व्यो\-ष\-वि\-नि\-र्द्दि\-हन्यः ।}\ H}}\skt{ । }\pend
\pstart\versepar \skt{चूर्णानि कृत्वा }\edtext{\skt{गुडर्शर्करे च}}{\Afootnote{K; \skt{गुडशर्करे च}\ A H}}\skt{ तथैव सर्पिर्मधुनी }\edtext{\skt{शुभे}}{\Afootnote{K H; \skt{शुमे}\ A}}\skt{ च ।}\pend

\pstart\versepar \ledsidenote{\passref{6.44.29}}\edtext{\edtext{\skt{सम्भा\-र\-मे\-त\-दि\-प\-चे\-थि\-ताय}}{\lemma{\skt{सम्भारम्}}\Afootnote{K; \skt{संभार}\ A H}}}{\lemma{\skt{एत\-दि\-प\-चे\-थि\-ताय}}\Afootnote{K; \skt{मेत\-द्वि\-प\-चे\-न्नि\-धाय}\ A; \skt{येत\-द्वि\-प\-चे\-थि\-ताय}\ H}}\skt{ }\edtext{\skt{सारवतो}}{\Afootnote{K; \skt{सारोदके सारवतो}\ A H}}\skt{ गणस्य । }\pend
\pstart\versepar \skt{}\edtext{\skt{दातञ्च}}{\Afootnote{K; \skt{जातं च}\ A H}}\skt{ }\edtext{\skt{लेह्य॥।विदित्वा}}{\Afootnote{K; \skt{लेह्यं मतिमान्विदित्वा}\ A H}}\skt{ }\edtext{\skt{निधा\-प\-ये\-न्मु\-स्कजे}}{\Afootnote{K; \skt{निधा\-प\-ये\-न्मो\-क्ष\-कजे}\ A; \skt{निधापयेन्मुष्क कजे}\ H}}\skt{ समुद्गे ।}\pend

\pstart\versepar \ledsidenote{\passref{6.44.30}}\edtext{\edtext{\skt{हंत्येष}}{\lemma{\skt{हंत्य्}}\Afootnote{K; \skt{हन्}\ A H}}}{\lemma{\skt{एष}}\Afootnote{K; \skt{त्येष}\ A H}}\skt{ }\edtext{\skt{योगः}}{\Afootnote{K H; \skt{लेहः}\ A}}\skt{ }\edtext{\skt{लु}}{\Afootnote{\skt{खलु}\ A H; \skt{…लु}\ K}}\skt{ पाण्डुरोगं }\edtext{\edtext{\skt{सशो\-थ\-मु\-ग्रा\-मपि}}{\lemma{\skt{सशोथमुग्राम्}}\Afootnote{A K; \skt{सशोथमुग्रा}\ H}}}{\lemma{\skt{अपि}}\Afootnote{A K; \skt{मपि}\ H}}\skt{ }\edtext{\skt{कामलञ्च}}{\Afootnote{K; \skt{कामलां च}\ A H}}\skt{ । }\pend
\pstart\versepar \skt{सश\-र्क\-रा\-का\-म\-लिनां त्रिभण्डी हिता गवाक्षी सगुडा च शुण्ठी ॥}\pend

\pstart\versepar \ledsidenote{\passref{6.44.31}}\skt{कालेयके चापि घृतं विपक्वं हितञ्च }\edtext{\skt{तस्या\-द्र\-ज\-ना\-दि\-मि\-श्रम्}}{\Afootnote{K; \skt{तत्स्या\-द्र\-ज\-नी\-वि\-मि\-श्रम् ।}\ A; \skt{तस्यार्द्रजनादि मित्रं ।}\ H}}\skt{ । }\pend
\pstart\versepar \skt{}\edtext{\skt{शिला\-ह्व\-य\-ञ्वा\-प्यथ}}{\Afootnote{K; \skt{धातुं नदीजं जतु शैल}\ A; \skt{शालाह्वयं वाप्यथ}\ H}}\skt{ }\edtext{\skt{माक्षिकम्वा}}{\Afootnote{K H; \skt{जं वा}\ A}}\skt{ कुम्भाह्वये }\edtext{\skt{मूत्र\-यु\-त\-म्पि\-बेद्वा}}{\lemma{\skt{पिबेद्वा}}\Afootnote{\skt{पिबेद्वा ॥}\ A; \skt{पिवेद्वा}\ K H}}\skt{ ।}\pend

\pstart\versepar \ledsidenote{\passref{6.44.32}}\edtext{\skt{मूत्रस्थितं}}{\Afootnote{K H; \skt{मूत्रे स्थितं}\ A}}\skt{ सैन्ध\-व\-सं\-प्र\-युक्तं }\edtext{\skt{समां}}{\Afootnote{K; \skt{मासं}\ A H}}\skt{ }\edtext{\skt{पिबेत्वापि}}{\Afootnote{\skt{पिबेद्वापि}\ A; \skt{पिवेत्वापि}\ K H}}\skt{ }\edtext{\skt{हिरण्यकिट्टं}}{\Afootnote{K; \skt{हि लोहकिट्टम् ।}\ A; \skt{हिरण्ये किट्टम् ।}\ H}}\skt{ । }\pend
\pstart\versepar \skt{}\edtext{\skt{द्राक्षा\-गु\-डू\-च्या\-म\-ल\-की\-र\-सेन}}{\Afootnote{K H; \skt{दग्ध्वा\-क्ष\-का\-ष्ठै\-र्मल}\ A}}\skt{ }\edtext{\skt{सिद्धं}}{\Afootnote{K H; \emph{om.}\ A}}\skt{ }\edtext{\skt{घृतं}}{\Afootnote{K H; \emph{om.}\ A}}\skt{ }\edtext{\skt{लाकरके}}{\Afootnote{K H; \emph{om.}\ A}}\skt{ }\edtext{\skt{हिताये}}{\Afootnote{K H; \skt{मायसं वा गोमू\-त्र\-नि\-र्वा\-पि\-त\-म\-ष्ट\-वा\-रान् ॥}\ A}}\skt{ ।}\pend

\pstart\versepar \ledsidenote{\passref{6.44.33}}\edtext{\skt{सिद्धास्य}}{\Afootnote{K H; \skt{विचूर्ण्य}\ A}}\skt{ }\edtext{\skt{वस्तिमधुरो}}{\Afootnote{K; \skt{लीढं मधु}\ A; \skt{वस्तिम्मधुरो}\ H}}\skt{ }\edtext{\skt{हलीमे}}{\Afootnote{K H; \emph{om.}\ A}}\skt{ }\edtext{\skt{हन्यस्त्वरिष्टा}}{\Afootnote{K H; \emph{om.}\ A}}\skt{ }\edtext{\skt{मृदवश्च}}{\Afootnote{K; \emph{om.}\ A; \skt{मृदवश}\ H}}\skt{ }\edtext{\skt{पथ्या}}{\Afootnote{K H; \skt{ना}\ A}}\skt{ । }\pend
\pstart\versepar \skt{}\edtext{\skt{गौडानरिष्टां}}{\Afootnote{K; \skt{चिरेण कुम्भाह्वयं पाण्डुगदं निहन्यात् । सिन्धूद्भवं}\ A; \skt{गौडा\-न\-रि\-ष्टा\-न्मधु}\ H}}\skt{ }\edtext{\skt{मधुशर्करांञ्च}}{\Afootnote{K; \skt{वाग्निसमं च कृत्वा क्षिप्त्वा च}\ A; \skt{शर्क्करांञ्च}\ H}}\skt{ }\edtext{\skt{मूत्रासवां}}{\Afootnote{K; \skt{मूत्रे सकृ}\ A; \skt{मूत्रासवा}\ H}}\skt{ }\edtext{\skt{क्षार\-घृ\-तां\-स्त\-थैव}}{\Afootnote{K; \skt{देव तप्तम् ॥}\ A; \skt{न्क्षार\-घृ\-तां\-स्त\-थैव ॥}\ H}}\skt{ ।}\pend

\pstart\versepar \ledsidenote{\passref{6.44.34}}\edtext{\edtext{\skt{स्निग्धा\-न्र\-सा\-ना\-म\-ल\-कै\-रू\-पेतां}}{\lemma{\skt{स्निग्धा\-न्र\-सा\-ना\-म\-ल\-कैर्}}\Afootnote{K H; \emph{om.}\ A}}}{\lemma{\skt{ऊपेतां}}\Afootnote{K H; \emph{om.}\ A}}\skt{ }\edtext{\skt{कोलान्वितां}}{\Afootnote{K H; \emph{om.}\ A}}\skt{ }\edtext{\skt{वापि हि}}{\Afootnote{K H; \emph{om.}\ A}}\skt{ }\edtext{\skt{जंगलानां}}{\Afootnote{K; \emph{om.}\ A; \skt{जाङ्गलानां ।}\ H}}\skt{ । }\pend
\pstart\versepar \skt{}\edtext{\skt{सेवेत}}{\Afootnote{K H; \emph{om.}\ A}}\skt{ }\edtext{\skt{शोथाभिहितांञ्च}}{\Afootnote{K; \emph{om.}\ A; \skt{शोथाभिहितां श्च}\ H}}\skt{ }\edtext{\skt{योगां}}{\Afootnote{K; \emph{om.}\ A; \skt{योगा}\ H}}\skt{ }\edtext{\skt{पाण्ड्वामयी}}{\Afootnote{K; \emph{om.}\ A; \skt{न्पाण्ड्वामयी}\ H}}\skt{ }\edtext{\skt{शालियवाञ्च}}{\Afootnote{K; \skt{लौहं च}\ A; \skt{शालियवञ्च}\ H}}\skt{ }\edtext{\skt{नित्यं}}{\Afootnote{K H; \skt{किट्टं बहुशश्च तप्त्वा निर्वाप्य मूत्रे बहुशस्तथैव । एकीकृतं गोज\-ल\-पि\-ष्ट\-मे\-त\-दै\-क\-ध्य\-मा\-वाप्य पचेदुखायाम् ॥}\ A}}\skt{ ।}\pend

\pstart\versepar \ledsidenote{\passref{6.44.35}}\edtext{\skt{ज्वरं}}{\Afootnote{K H; \skt{य}\ A}}\skt{ }\edtext{\skt{तृषां}}{\Afootnote{K H; \skt{था न}\ A}}\skt{ }\edtext{\edtext{\skt{छर्दि\-म\-था\-ति\-सारं}}{\lemma{\skt{छर्दिम्}}\Afootnote{K; \skt{दह्ये}\ A; \skt{छर्द्दि}\ H}}}{\lemma{\skt{अथातिसारं}}\Afootnote{K; \skt{त तथा विशुष्कं}\ A; \skt{मथातिसारं}\ H}}\skt{ }\edtext{\skt{मूर्च्छा\-क्ष\-यं\-श्वा\-स\-म\-न्न\-लिप्सां}}{\Afootnote{K H; \skt{चूर्णीकृतं पेयमुदश्विता}\ A}}\skt{ । }\pend
\pstart\versepar \skt{}\edtext{\skt{तथा}}{\Afootnote{K; \skt{तत् । तक्रौदनाशी विजयेत रोगं}\ A; \skt{तथावि}\ H}}\skt{ }\edtext{\skt{विपा\-क\-स्व\-र\-भे\-द\-सादां}}{\Afootnote{K; \skt{पाण्डुं तथा दीपयतेऽनलं}\ A; \skt{पाकं\-स्व\-र\-भे\-द\-सा\-दाम्}\ H}}\skt{ }\edtext{\skt{जायेद्यथास्वां}}{\Afootnote{K; \skt{च ॥}\ A; \skt{जयेद्यथास्वां}\ H}}\skt{ }\edtext{\skt{प्रसमीक्षदोषां}}{\Afootnote{K; \emph{om.}\ A; \skt{प्रस\-मी\-क्ष्य\-दो\-षान् ॥}\ H}}\skt{ ॥}\pend

\pstart\versepar \ledsidenote{\passref{6.44.36}}\edtext{\skt{अन्तेषु}}{\Afootnote{K H; \emph{om.}\ A}}\skt{ }\edtext{\skt{शूनं}}{\Afootnote{K H; \emph{om.}\ A}}\skt{ }\edtext{\skt{परिहीनमध्यं}}{\Afootnote{K; \emph{om.}\ A; \skt{परुहीनमध्यं}\ H}}\skt{ }\edtext{\skt{म्लानन्तथान्तेषु च}}{\Afootnote{K H; \emph{om.}\ A}}\skt{ }\edtext{\skt{मध्यशूनं}}{\Afootnote{K H; \emph{om.}\ A}}\skt{ । }\pend
\pstart\versepar \skt{}\edtext{\skt{गुदेषु}}{\Afootnote{K H; \emph{om.}\ A}}\skt{ }\edtext{\skt{शोफ\-स्त्व\-थ\-मु\-ष्क\-शूनं}}{\Afootnote{K; \emph{om.}\ A; \skt{शोफा\-स्त्व\-थ\-मु\-ष्क\-शूनं}\ H}}\skt{ }\edtext{\skt{प्रातम्यमानञ्च}}{\Afootnote{K; \emph{om.}\ A; \skt{प्राताम्यमानञ्च}\ H}}\skt{ }\edtext{\skt{विसञ्ज्ञकल्पं}}{\Afootnote{K H; \skt{द्राक्षा\-गु\-डू\-च्या\-म\-ल\-की\-र\-सैश्च सिद्धं घृतं लाघ\-र\-के(\-ला\-घ\-वके) हितं च । गौडा\-न\-रि\-ष्टा\-न्म\-धु\-श\-र्क\-राश्च मूत्रा\-स\-वा\-न्क्षा\-र\-कृ\-तां\-स्त\-थैव ॥}\ A}}\skt{ ।}\pend

\pstart \ledsidenote{\passref{6.44.37}}\edtext{\skt{विवि\-र्ज\-ये\-त्पा\-ण्डु\-किनं}}{\Afootnote{K; \skt{स्निग्धा\-न्र\-सा\-ना\-म\-ल\-कै\-रु\-पे\-ता\-न्को\-ला\-न्वि\-ता\-न्वापि हि जाङ्गलानाम् । सेवेत शोफाभिहितांश्च योगान्पाण्ड्वामयी शालि}\ A; \skt{विव\-र्ज्ज\-ये\-त्पा\-ण्डु\-किनं}\ H}}\skt{ }\edtext{\skt{यशोर्थी}}{\Afootnote{\skt{यवांश्च}\ A; \skt{यशो}\ K; \skt{यशोऽ}\ H}}\skt{ }\edtext{\edtext{\skt{तथा\-ति\-सा\-र\-ज्व\-र\-पी\-डि\-त\-ञ्चेति}}{\lemma{\skt{तथा\-ति\-सा\-र\-ज्व\-र\-पी\-डि\-तञ्}}\Afootnote{\skt{नित्यम् ॥}\ A; \skt{र्थी तथा\-ति\-सा\-र\-ज्व\-र\-पी\-डि\-तञ्}\ K; \skt{र्थो तथा\-ति\-सा\-र\-ज्व\-र\-पी\-डितं}\ H}}}{\lemma{\skt{चेति}}\Afootnote{K H; \emph{om.}\ A}}\skt{ ॥}\pend

\pstart \ledsidenote{\passref{6.44.trailer}}\edtext{\skt{कायचिकित्सा}}{\Afootnote{\skt{इति सुश्रु\-त\-सं\-हि\-ता\-या\-मु\-त्त\-र\-त\-न्त्रा\-न्त\-र्गते काय\-चि\-कि\-त्सा\-तन्त्रे पाण्डु\-रो\-ग\-प्र\-ति\-षेधो नाम (षष्ठोऽध्यायः, आदितः) चतु\-श्च\-त्वा\-रिं\-श\-त्त\-मोऽ\-ध्यायः ॥४४ ॥}\ A; \skt{कायचि कायचिकित्सा ॥ ॥}\ K; \skt{उत्तरतन्रे एक\-च\-त्वा\-रिं\-श\-त्तमः कायचिकित्सायां षष्ठोध्यायः ॥ ॥}\ H}}\skt{ ॥ ॥}\pend

\endnumbering
\beginnumbering
\eledsection*{Adhyāya 45}
\pstart \ledsidenote{\passref{6.45.1}}\skt{अथातो रक्त\-पि\-त्त\-प्र\-ति\-षेधं }\edtext{\skt{व्याख्यास्यामः}}{\Afootnote{A H; \skt{व्या ॥}\ K}}\skt{ ॥}\pend

\pstart\versepar \ledsidenote{\passref{6.45.3x}}\edtext{\skt{व्याया\-म\-भा\-रा\-ध्य\-य\-न\-व्य\-वो\-या\-त\-प\-से\-वनैः}}{\Afootnote{K; \skt{व्याया\-म\-भा\-रा\-ध्य\-य\-न\-व्यवो यातपसेवनैः ।}\ H}}\skt{ । }\pend
\pstart\versepar \skt{}\edtext{\skt{तीक्ष्णो\-ष्ण\-क्षा\-र\-ल\-व\-णै\-अम्लैः}}{\Afootnote{K; \skt{तीक्ष्णो\-ष्ण\-क्षा\-र\-ल\-व\-णै\-रम्लैः}\ H}}\skt{ }\edtext{\skt{कट्वाभिरेव}}{\lemma{\skt{कट्वाभिर्}}\Afootnote{K; \skt{कट्वभिर्}\ H}}\skt{ च ।}\pend

\pstart\versepar \ledsidenote{\passref{6.45.4x}}\edtext{\skt{अभिक्ष्णं}}{\Afootnote{K; \skt{अभीक्ष्णं}\ H}}\skt{ सेवतो }\edtext{\skt{तूष्णो}}{\Afootnote{K; \skt{त्युष्णो}\ H}}\skt{ रसः पित्तं प्रदूषयेत् । }\pend
\pstart\versepar \skt{पित्तं विदग्धं }\edtext{\edtext{\skt{स्वगु\-णै\-वि\-द\-ह\-त्याशु}}{\lemma{\skt{स्वगु\-णै\-वि\-द\-हत्य्}}\Afootnote{K; \skt{स्वगुणैर्व्विदह}\ H}}}{\lemma{\skt{आशु}}\Afootnote{K; \skt{त्याशु}\ H}}\skt{ शोणितं ।}\pend

\pstart\versepar \ledsidenote{\passref{6.45.5}}\skt{ततः }\edtext{\skt{प्रवर्तते}}{\Afootnote{A K; \skt{प्रवत्तते}\ H}}\skt{ }\edtext{\skt{रक्तमूर्ध्वं}}{\Afootnote{K; \skt{रक्तमूर्ध्वं}\ A; \skt{रक्तमूर्द्द्धञ्}\ H}}\skt{ चाधो द्विधापि वा । }\pend
\pstart\versepar \skt{}\edtext{\skt{आमा\-श\-या\-द्व्र\-जे\-दू\-र्ध्व\-मधः}}{\lemma{\skt{आमा\-श\-या\-द्व्र\-जे\-दू\-र्ध्वम्}}\Afootnote{A K; \skt{आमा\-स\-या\-द्व्र\-जे\-दू\-र्ध्वम्}\ H}}\skt{ }\edtext{\skt{पक्वा\-श\-या\-द्व्र\-जेत्}}{\Afootnote{A K; \skt{पक्वा\-स\-या\-द्र\-हेत् ॥}\ H}}\skt{ ।}\pend

\pstart \ledsidenote{\passref{6.45.6ab}}\edtext{\edtext{\skt{विद\-ग्ध\-यो\-श्चा\-प्यु\-भ\-यो\-र्द्वि\-धा\-भागं}}{\lemma{\skt{विदग्धयोश्}}\Afootnote{K H; \skt{विद\-ग्ध\-यो\-र्द्व\-योश्}\ A}}}{\lemma{\skt{उभयोर्}}\Afootnote{K H; \emph{om.}\ A}}\skt{ }\edtext{\skt{प्रवर्तते}}{\Afootnote{A K; \skt{प्रवर्त्तेन ।}\ H}}\skt{ ।}\pend

\pstart \ledsidenote{\passref{6.45.7ab}}\edtext{\skt{ऊर्ध्वं}}{\Afootnote{A K; \skt{ऊर्ध्व}\ H}}\skt{ }\edtext{\skt{साध्य\-म\-धो\-या\-प्य\-म\-साध्यं}}{\Afootnote{K; \skt{साध्यमधो याप्यमसाध्यं}\ A; \skt{न्साध्य\-म\-धो\-या\-प्य\-म\-साध्यं}\ H}}\skt{ }\edtext{\skt{युगपद्गतं}}{\Afootnote{K; \skt{युगपद्गतम्}\ A H}}\skt{ ।}\pend

\pstart \ledsidenote{\passref{6.45.6cd}}\edtext{\edtext{\skt{यकृ\-न्प्ली\-हा\-स्या\-मि\-छन्ति}}{\lemma{\skt{यकृन्प्लीहास्याम्}}\Afootnote{K H; \emph{om.}\ A}}}{\lemma{\skt{इछन्ति}}\Afootnote{K H; \emph{om.}\ A}}\skt{ }\edtext{\skt{केचिद्रक्तगतिं}}{\Afootnote{K H; \skt{केचित्सयकृतः प्लीह्नः प्रवदन्त्यसृजो गतिम् ॥}\ A}}\skt{ }\edtext{\skt{जानाः}}{\Afootnote{K; \emph{om.}\ A; \skt{जनाः ।}\ H}}\skt{ ॥}\pend

\pstart \ledsidenote{\passref{6.45.7cd}}\skt{सदनं शीतकामित्वं }\edtext{\skt{कण्ठ\-धू\-मा\-य\-न\-म्वमी}}{\lemma{\skt{वमी}}\Afootnote{K; \skt{वमिः ॥}\ A H}}\skt{ ।}\pend

\pstart\versepar \ledsidenote{\passref{6.45.8}}\skt{लोहगन्धिश्च }\edtext{\skt{निश्वासो}}{\Afootnote{K H; \skt{निःश्वासो}\ A}}\skt{ }\edtext{\skt{भवत्यस्मिं}}{\lemma{\skt{अस्मिं}}\Afootnote{K; \skt{अस्मि}\ A H}}\skt{ }\edtext{\skt{भविष्यति}}{\Afootnote{K; \skt{न्भविष्यति}\ A H}}\skt{ । }\pend
\pstart\versepar \skt{}\edtext{\skt{वाह्या\-सृ\-ग्ल\-क्ष\-णै\-स्तस्य}}{\lemma{\skt{वाह्या\-सृ\-ग्ल\-क्ष\-णैस्}}\Afootnote{K H; \skt{बाह्या\-सृ\-ग्ल\-क्ष\-णैस्}\ A}}\skt{ }\edtext{\skt{संख्या}}{\Afootnote{K; \skt{सङ्ख्या}\ A; \skt{स ख्या}\ H}}\skt{ }\edtext{\skt{दोषोच्छ्रितैर्विदुः}}{\lemma{\skt{दोषोच्छ्रितैर्}}\Afootnote{K H; \skt{दोषोच्छ्रितीर्}\ A}}\skt{ ।}\pend

\pstart\versepar \ledsidenote{\passref{6.45.9}}\edtext{\skt{दौर्व\-ल्य\-श्वा\-स\-का\-स\-ज्व\-र\-व\-थु\-मदा}}{\Afootnote{K; \skt{दौर्ब\-ल्य\-श्वा\-स\-का\-स\-ज्व\-र\-व\-म\-थु\-मदा}\ A; \skt{दौर्व्व\-ल्यं\-श्वा\-स\-का\-स\-ज्व\-र\-व\-म\-थु\-मदा}\ H}}\skt{ }\edtext{\skt{पाण्डु\-ता\-दा\-ह\-मूर्च्छा}}{\Afootnote{K H; \skt{स्तन्द्रि\-ता\-दा\-ह\-मूर्च्छा}\ A}}\skt{ }\pend
\pstart\versepar \skt{भुक्ते चान्ने विदा\-ह\-स्त्व\-धृ\-ति\-रपि सदा हृद्य तुल्या }\edtext{\skt{च}}{\Afootnote{A K; \skt{व}\ H}}\skt{ पीडा । }\pend
\pstart\versepar \skt{तृष्णाकण्ठस्य भेदः शिरसि च }\edtext{\skt{पवनं}}{\Afootnote{K H; \skt{दवनं}\ A}}\skt{ }\edtext{\skt{पूतिनिः}}{\Afootnote{K; \skt{पूतिनि}\ A H}}\skt{ ष्ठीवनञ्च }\pend
\pstart\versepar \skt{द्वेषो }\edtext{\skt{भक्ते}}{\Afootnote{A K; \skt{बक्ते}\ H}}\skt{ }\edtext{\skt{विपाको}}{\Afootnote{K H; \skt{ऽविपाको}\ A}}\skt{ विरतिरपि }\edtext{\skt{भवे\-द्र\-क्त\-पि\-त्तो\-प\-स\-र्गात्}}{\Afootnote{K; \skt{रते रक्त\-पि\-त्तो\-प\-सर्गाः ॥}\ A; \skt{भव\-द्र\-क्त\-पि\-त्तो\-प\-स\-र्गात् ॥}\ H}}\skt{ ॥}\pend

\pstart\versepar \ledsidenote{\passref{6.45.10}}\skt{मांस\-प्र\-क्षा\-ल\-नाभं }\edtext{\skt{क्वथितमपि}}{\lemma{\skt{अपि}}\Afootnote{K H; \skt{इव}\ A}}\skt{ }\edtext{\skt{यकृ\-त्क\-र्द\-मा\-ञ्भो\-नि\-भम्वा}}{\Afootnote{K; \skt{च यत्क\-र्द\-मा\-म्भो\-निभं वा}\ A; \skt{यकृ\-त्क\-र्द्द\-मा\-म्भो\-नि\-भम्वा}\ H}}\skt{ }\pend
\pstart\versepar \skt{मेदः }\edtext{\skt{पूयास्तिकल्पं}}{\Afootnote{K; \skt{पूयास्रकल्पं}\ A; \skt{पूयास्थिकल्पं}\ H}}\skt{ यकृदिव यदि वा पक्व\-ज\-म्बू\-फ\-ला\-भम् । }\pend
\pstart\versepar \skt{यत्कृष्णं यच्च }\edtext{\skt{नीलं}}{\Afootnote{A K; \skt{नील}\ H}}\skt{ }\edtext{\edtext{\skt{भृश\-म\-ति\-कु\-णपं}}{\lemma{\skt{भृशम्}}\Afootnote{A K; \skt{भृश}\ H}}}{\lemma{\skt{अतिकुणपं}}\Afootnote{A K; \skt{मतिकुणपं}\ H}}\skt{ यत्र चोक्ता विकाराः }\pend
\pstart\versepar \skt{तद्वर्ज्यं रक्तपित्तं सुर\-प\-ति\-ध\-नुषा यच्च तुल्यं विभाति ॥}\pend

\pstart\versepar \ledsidenote{\passref{6.45.11}}\edtext{\edtext{\skt{नोदृक्तमादौ}}{\lemma{\skt{नोदृक्तम्}}\Afootnote{K H; \emph{om.}\ A}}}{\lemma{\skt{आदौ}}\Afootnote{K H; \skt{नादौ}\ A}}\skt{ }\edtext{\skt{संग्राह्यं}}{\Afootnote{A K; \skt{सं}\ H}}\skt{ }\edtext{\skt{वलिनो}}{\Afootnote{K; \skt{उद्रिक्तं यदसृग्बलिनो}\ A; \skt{ग्र्ह्यां वलिनो}\ H}}\skt{ }\edtext{\skt{भिषज}}{\Afootnote{K; \skt{ऽश्नतः ।}\ A; \skt{भिषजा}\ H}}\skt{ }\edtext{\skt{सदा}}{\Afootnote{K H; \skt{त}\ A}}\skt{ । }\pend
\pstart\versepar \skt{}\edtext{\skt{हृत्पा\-ण्डु\-ग्र\-ह\-णी\-रो\-ग\-प्ली\-ह\-गु\-ल्म\-ज्व\-रा\-वहं}}{\Afootnote{K H; \skt{त्पाण्डु\-ग्र\-ह\-णी\-कु\-ष्ठ\-प्ली\-ह\-गु\-ल्म\-ज्व\-रा\-व\-हम् ॥}\ A}}\skt{ ।}\pend

\pstart\versepar \ledsidenote{\passref{6.45.12}}\skt{अधः प्रवृत्तं }\edtext{\skt{वम\-नै\-रू\-र्ध्व\-गन्तु}}{\lemma{\skt{ऊर्ध्वगन्तु}}\Afootnote{K; \skt{ऊर्ध्वगं च}\ A; \skt{ऊर्द्धगन्तु}\ H}}\skt{ विरेचनैः । }\pend
\pstart\versepar \skt{}\edtext{\skt{जयेदन्यतरं}}{\Afootnote{K H; \skt{जयेदन्यतर}\ A}}\skt{ }\edtext{\skt{चापि}}{\Afootnote{K H; \skt{द्वापि}\ A}}\skt{ क्षीणस्य शमनैरसृक् ।}\pend

\pstart\versepar \ledsidenote{\passref{6.45.15}}\edtext{\skt{द्राक्षा\-का\-र्श्म\-य\-म\-धुकं}}{\Afootnote{K; \skt{द्राक्षामधुकका}\ A; \skt{द्राक्षा\-का\-श्म\-र्य\-म\-धुक}\ H}}\skt{ }\edtext{\skt{सितायेतं}}{\Afootnote{K H; \skt{श्मर्यसितायुक्तं}\ A}}\skt{ विरेचनं । }\pend
\pstart\versepar \skt{}\edtext{\skt{यष्ठी\-म\-धु\-क\-सि\-द्धञ्च}}{\Afootnote{K H; \skt{यष्टी\-म\-धु\-क\-युक्तं च}\ A}}\skt{ सक्षौद्रं वमनं हितं ।}\pend

\pstart\versepar \ledsidenote{\passref{6.45.16}}\skt{पयांसि शीतानि रसाश्च जाङ्गलाः }\edtext{\skt{सती\-न\-मू\-द्गा\-य\-व\-शा\-लि\-ष\-ष्टिका}}{\Afootnote{K H; \skt{सतीनयूषाश्च सशालिषष्टिकाः ।}\ A}}\skt{ । }\pend
\pstart\versepar \skt{}\edtext{\skt{जीवन्ति}}{\Afootnote{K H; \skt{पटोल}\ A}}\skt{ }\edtext{\edtext{\skt{शैलू\-सु\-नि\-ष\-ण्ण\-यू\-थि\-कां\-स्त\-था\-ति\-मुक्तां}}{\lemma{\skt{शैलू\-सु\-नि\-ष\-ण्ण\-यू\-थि\-कांस्}}\Afootnote{K H; \skt{शेलू\-सु\-नि\-ष\-ण्ण\-यू\-थिका}\ A}}}{\lemma{\skt{तथातिमुक्तां}}\Afootnote{K; \skt{वटातिमुक्ताङ्}\ A; \skt{तथातिमुक्तांङ्}\ H}}\skt{ }\edtext{\skt{कुररजाः}}{\Afootnote{K; \skt{कुर\-सि\-न्दु\-वा\-र\-जम् ॥}\ A; \skt{कुर\-सि\-न्दु\-वा\-रजाः ॥}\ H}}\skt{ ।}\pend

\pstart\versepar \ledsidenote{\passref{6.45.17}}\edtext{\skt{शाकं}}{\Afootnote{K H; \skt{हितं च शाकं}\ A}}\skt{ }\edtext{\skt{हितं}}{\Afootnote{K; \skt{घृत}\ A; \skt{हित}\ H}}\skt{ }\edtext{\skt{सर्पिषिसंस्कृतं}}{\Afootnote{K; \skt{संस्कृतं}\ A; \skt{न्सर्प्पिषिसंस्कृतं}\ H}}\skt{ सदा }\edtext{\skt{पथ्याञ्च}}{\Afootnote{K H; \skt{तथैव}\ A}}\skt{ धात्री\-फ\-ल\-दा\-डि\-मा\-न्वितं । }\pend
\pstart\versepar \skt{रसाश्च }\edtext{\skt{पारा\-व\-त\-हं\-स\-कु\-र्म\-ज\-पेया}}{\Afootnote{K; \skt{पारा\-व\-त\-श\-ङ्ख\-कूर्म}\ A; \skt{पारावतहं सकुर्म्मजजाः}\ H}}\skt{ }\edtext{\skt{पेयास्तथा}}{\lemma{\skt{पेयास्}}\Afootnote{K H; \skt{जास्}\ A}}\skt{ }\edtext{\skt{चापि}}{\Afootnote{K H; \skt{यवाग्वो विहिता}\ A}}\skt{ }\edtext{\skt{घृतान्तरा}}{\Afootnote{K H; \skt{घृतोत्तराः ॥}\ A}}\skt{ }\edtext{\skt{हिताः}}{\Afootnote{K H; \emph{om.}\ A}}\skt{ ।}\pend

\pstart\versepar \ledsidenote{\passref{6.45.18}}\edtext{\edtext{\skt{सन्ता\-नि\-का\-श्चो\-त्प\-ल\-पू\-र्व\-केषु}}{\lemma{\skt{सन्तानिकाश्}}\Afootnote{A K; \skt{सान्तानिकाश्}\ H}}}{\lemma{\skt{चोत्प\-ल\-पू\-र्व\-केषु}}\Afootnote{K H; \skt{चोत्पलवर्गसा}\ A}}\skt{ }\edtext{\skt{घृतेषु}}{\Afootnote{K H; \skt{धिते क्षी}\ A}}\skt{ }\edtext{\skt{सिद्धाः}}{\Afootnote{K H; \skt{रे}\ A}}\skt{ }\edtext{\skt{पयसः}}{\Afootnote{K H; \skt{प्रशस्ता}\ A}}\skt{ }\edtext{\skt{सिताह्वा}}{\Afootnote{K; \skt{मधु\-श\-र्क\-रो\-त्तराः ।}\ A; \skt{सिता .ः ।}\ H}}\skt{ । }\pend
\pstart\versepar \skt{}\edtext{\skt{हिमाः}}{\Afootnote{A K; \skt{हिमा}\ H}}\skt{ प्रदेहा मधुरा }\edtext{\skt{गणाश्च}}{\Afootnote{K H; \skt{गणाश्च ये}\ A}}\skt{ घृतानि पथ्यानि च }\edtext{\skt{रक्तपित्ते}}{\Afootnote{K H; \skt{रक्तपित्तिनाम् ॥}\ A}}\skt{ ।}\pend

\pstart\versepar \ledsidenote{\passref{6.45.19}}\edtext{\skt{मधू\-क\-शो\-भा\-ञ्ज\-न\-को\-दि\-रजैः}}{\Afootnote{K; \skt{मधू\-क\-शो\-भा\-ञ्ज\-न\-को\-वि\-दा\-रजैः}\ A; \skt{मधू\-क\-सौ\-भा\-ञ्ज\-न\-को\-दि\-रजैः}\ H}}\skt{ }\edtext{\skt{प्रियंगुकानां}}{\Afootnote{K H; \skt{प्रियङ्गुकायाः}\ A}}\skt{ कुसुमैश्च }\edtext{\skt{चूर्णनितः}}{\Afootnote{K; \skt{चूर्णितैः ।}\ A; \skt{चूर्ण्णितः ।}\ H}}\skt{ । }\pend
\pstart\versepar \skt{भिष\-ग्वि\-द\-ध्या\-च्च\-तुरः }\edtext{\skt{समाक्षिकां}}{\Afootnote{K H; \skt{समाक्षिका}\ A}}\skt{ }\edtext{\skt{हिताय}}{\Afootnote{K H; \skt{न्हिताय}\ A}}\skt{ }\edtext{\skt{लेहानसृजाः}}{\Afootnote{K; \skt{लेहानसृजः}\ A H}}\skt{ प्रशान्तये ।}\pend

\pstart\versepar \ledsidenote{\passref{6.45.20}}\skt{लिह्याच्च }\edtext{\edtext{\skt{दुग्ध\-द्रु\-म\-ज\-न्न\-वां\-कुरां}}{\lemma{\skt{दुग्धद्रुमजन्}}\Afootnote{K; \skt{दूर्वावटजांश्}\ A; \skt{दुग्धद्रुमजान्}\ H}}}{\lemma{\skt{नवांकुरां}}\Afootnote{K; \skt{च पल्लवा}\ A; \skt{न वा .रा}\ H}}\skt{ }\edtext{\skt{मधुद्वितीयां}}{\Afootnote{K; \skt{न्मधुद्वितीया}\ A; \skt{न्मधुद्वितीयां}\ H}}\skt{ }\edtext{\skt{सितकर्णिकस्य}}{\Afootnote{K H; \skt{न्सितकर्णिकस्य}\ A}}\skt{ । }\pend
\pstart\versepar \skt{}\edtext{\skt{पथ्याञ्च}}{\Afootnote{K H; \skt{च । हितं च}\ A}}\skt{ खर्जूरफलं समाक्षिकं फलानि }\edtext{\skt{चान्या\-न्य\-पि\-त\-ण्गु\-णा\-न्यथ}}{\Afootnote{K; \skt{चान्यान्यपि तद्गुणान्यथ ॥}\ A; \skt{चान्या\-न्य\-पि\-त\-ङ्गु\-णा\-न्यथा ॥}\ H}}\skt{ ॥}\pend

\pstart \ledsidenote{\passref{6.45.21ab}}\edtext{\skt{रक्ता\-नि\-सा\-रा\-भि\-हि\-ता\-न्यो\-गा\-न\-त्रापि}}{\Afootnote{K; \skt{रक्ता\-ति\-सा\-र\-प्रो\-क्तांश्च योगानत्रापि}\ A}}\skt{ योजयेत् ।}\pend

\pstart \ledsidenote{\passref{6.45.21x}}\edtext{\skt{नीलो\-त्प\-ला\-ना\-म्म\-धुना}}{\lemma{\skt{नीलोत्पलानाम्}}\Afootnote{H; \skt{॥नि\-लो\-त्प\-ला\-नाम्}\ K}}\skt{ भस्म }\edtext{\skt{वापि}}{\Afootnote{K; \skt{वारि}\ H}}\skt{ }\edtext{\skt{परिस्रुतं}}{\Afootnote{K; \skt{परिप्लुतं ॥}\ H}}\skt{ ।}\pend

\pstart \ledsidenote{\passref{6.45.21cd}}\edtext{\edtext{\skt{शुद्धे\-क्षु\-ग\-ण्ड\-मा\-पोष्य}}{\lemma{\skt{शुद्धेक्षुगण्डम्}}\Afootnote{\skt{शुद्धेक्षुकाण्डम्}\ A; \skt{शुद्धेक्षुगण्ड}\ K H}}}{\lemma{\skt{आपोष्य}}\Afootnote{\skt{आपोथ्य}\ A; \skt{मापोष्य}\ K; \skt{मायोष्ये}\ H}}\skt{ नवे कुम्भे }\edtext{\skt{हितांभसा}}{\Afootnote{K; \skt{हिमाम्भसा}\ A H}}\skt{ ।}\pend

\pstart\versepar \ledsidenote{\passref{6.45.22}}\skt{योजयित्वा }\edtext{\skt{स्थितं}}{\Afootnote{K H; \skt{क्षिपे}\ A}}\skt{ }\edtext{\skt{रात्रावाकाशे}}{\Afootnote{K H; \skt{द्रात्रावाकाशे}\ A}}\skt{ सोत्पलं तु तत् । }\pend
\pstart\versepar \skt{प्रातः }\edtext{\skt{श्रुतं}}{\Afootnote{K; \skt{स्रुतं}\ A H}}\skt{ क्षौद्रयुतं }\edtext{\skt{पिवे\-च्छो\-णि\-तै\-पि\-त्तकः}}{\Afootnote{K; \skt{पिबे\-च्छो\-णि\-त\-पि\-त्त\-वान् ॥}\ A; \skt{पिवे\-च्छो\-णि\-त\-यै\-पि\-त्तकः ।}\ H}}\skt{ ।}\pend

\pstart\versepar \ledsidenote{\passref{6.45.23}}\edtext{\skt{पिवे\-च्छि\-त\-क\-षायं वा}}{\Afootnote{K; \skt{पिबे\-च्छी\-त\-क\-षायं वा}\ A; \skt{पीवे\-च्छी\-त\-क\-षा\-यम्वा}\ H}}\skt{ }\edtext{\skt{जम्वा\-म्रा\-र्जु\-न\-स\-म्भवं}}{\Afootnote{K; \skt{जम्ब्वा\-म्रा\-र्जु\-न\-सं\-भ\-वम् ।}\ A; \skt{जम्व्वा\-म्रा\-र्जु\-न\-स\-म्भवं ।}\ H}}\skt{ । }\pend
\pstart\versepar \skt{}\edtext{\skt{आमो\-दु\-म्व\-र\-नि\-र्यासं}}{\Afootnote{K; \skt{उदुम्बर}\ A; \skt{आमो\-दु\-म्व\-र\-नि\-र्य\-सम्}\ H}}\skt{ }\edtext{\skt{पिबे\-त्स\-क्षौ\-द्र\-श\-र्करं}}{\Afootnote{\skt{फलं पिष्ट्वा पिबे\-त्त\-द्र\-स\-मेव वा ॥}\ A; \skt{पिवे\-त्स\-क्षौ\-द्र\-श\-र्करं}\ K H}}\skt{ ।}\pend

\pstart\versepar \ledsidenote{\passref{6.45.24}}\edtext{\skt{त्रपु\-शी\-मू\-ल\-क\-ल्कम्वा}}{\Afootnote{K; \skt{त्रपु\-षी\-मू\-ल\-कल्कं वा}\ A; \skt{त्रय\-सी\-मू\-ल\-क\-ल्कम्वा}\ H}}\skt{ सक्षौद्रं }\edtext{\skt{तण्डुतालाम्वुना}}{\Afootnote{K; \skt{तण्डुलाम्बुना ।}\ A; \skt{तण्डुलाम्वुना ॥}\ H}}\skt{ । }\pend
\pstart\versepar \skt{}\edtext{\skt{पिबेदक्षसमं}}{\Afootnote{A; \skt{पिवेदक्षसमं}\ K; \skt{पेवे\-द\-क्ष\-स\-मङ्}\ H}}\skt{ कल्कं }\edtext{\edtext{\skt{मधु\-क\-ल्पै\-व\-मेव}}{\lemma{\skt{मधुकल्पैवम्}}\Afootnote{K; \skt{यष्टीमधुकम्}\ A; \skt{मधुकस्यैव}\ H}}}{\lemma{\skt{एव}}\Afootnote{A K; \skt{मेव}\ H}}\skt{ वा ।}\pend

\pstart\versepar \ledsidenote{\passref{6.45.25}}\skt{चन्दनं }\edtext{\skt{पद्मं}}{\Afootnote{K; \skt{मधु}\ A; \skt{पद्म}\ H}}\skt{ }\edtext{\edtext{\skt{कंलोध्रमेवं}}{\lemma{\skt{कंलोध्रम्}}\Afootnote{K; \skt{कं रोध्रम्}\ A; \skt{कं लोध्र}\ H}}}{\lemma{\skt{एवं}}\Afootnote{A K; \skt{मेव}\ H}}\skt{ }\edtext{\skt{मेवं}}{\Afootnote{K H; \skt{एव}\ A}}\skt{ समं पिबेत् । }\pend
\pstart\versepar \skt{}\edtext{\skt{कर\-ञ्ज\-वी\-ज\-मे\-वम्वा}}{\lemma{\skt{एवम्}}\Afootnote{A K; \skt{ए}\ H}}\skt{ }\edtext{\skt{सिताक्षौद्रयुतं}}{\Afootnote{A K; \skt{त्सिता\-क्षौ\-द्र\-यु\-तम्}\ H}}\skt{ }\edtext{\skt{पिबेत्}}{\Afootnote{A; \skt{पिवेत् ।}\ K H}}\skt{ ।}\edtext{}{\lemma{\skt{पिबेत् । करञ्जवीजम्}}\Afootnote{\skt{पिबेत् । करञ्जबीजम्}\ A; \skt{पिवे\-त्क\-र\-ञ्ज\-वी\-जम्}\ K H}}\pend

\pstart\versepar \ledsidenote{\passref{6.45.26}}\edtext{\edtext{\skt{मज्जा\-न\-मि\-द्गु\-द\-स्यै\-व\-म्पेयो}}{\lemma{\skt{मज्जा\-न\-मि\-द्गु\-द\-स्यै\-वम्}}\Afootnote{K H; \skt{मज्जा\-न\-मि\-ङ्गु\-द\-स्यैवं}\ A}}}{\lemma{\skt{पेयो}}\Afootnote{K H; \skt{पिबे}\ A}}\skt{ }\edtext{\skt{मधुकसंयुताः}}{\Afootnote{K H; \skt{न्मधु\-क\-सं\-यु\-तम् ।}\ A}}\skt{ । }\pend
\pstart\versepar \skt{}\edtext{\skt{सुखोष्णलवणं}}{\Afootnote{K H; \skt{सुखोष्णं लवणं}\ A}}\skt{ }\edtext{\skt{वीजं}}{\Afootnote{K H; \skt{बीजं}\ A}}\skt{ }\edtext{\skt{कर\-ञ्ज\-द\-धि\-म\-स्तुना}}{\Afootnote{K H; \skt{कारञ्जं दधिमस्तुना ॥}\ A}}\skt{ ।}\pend

\pstart\versepar \ledsidenote{\passref{6.45.27}}\edtext{\skt{पिबेद्वापि}}{\Afootnote{A; \skt{पिवेद्वापि}\ K H}}\skt{ }\edtext{\skt{त्र्यहं}}{\Afootnote{A K; \skt{त्र्यअहं}\ H}}\skt{ मर्त्यो रक्त\-पि\-त्ता\-भि\-पी\-डितः । }\pend
\pstart\versepar \skt{रक्तपित्तहराः शस्ताः षडेते }\edtext{\edtext{\skt{योगमुत्तमाः}}{\lemma{\skt{योगम्}}\Afootnote{K; \skt{योग}\ A H}}}{\lemma{\skt{उत्तमाः}}\Afootnote{K; \skt{सत्तमाः ॥}\ A; \skt{मुलमाः ।}\ H}}\skt{ ।}\pend

\pstart\versepar \ledsidenote{\passref{6.45.28}}\edtext{\skt{पथ्या}}{\Afootnote{K; \skt{पथ्याश्}\ A}}\skt{ चैवावपीडेषु }\edtext{\skt{ग्राणतः}}{\Afootnote{K; \skt{घ्राणतः}\ A}}\skt{ }\edtext{\skt{प्रमृते}}{\Afootnote{K; \skt{प्रस्रुते}\ A}}\skt{ }\edtext{\skt{सृजि}}{\Afootnote{K; \skt{ऽसृजि ।}\ A}}\skt{ । }\pend
\pstart\versepar \skt{}\edtext{\skt{अति\-नि\-स्स्रु\-त\-रक्तो वा}}{\Afootnote{K; \skt{अति\-नि\-स्रु\-त\-रक्तो वा}\ A}}\skt{ क्षौद्रयुक्तं पिबेदसृक् । }\pend
\pstart\versepar \skt{यकृद्वा }\edtext{\skt{भक्ष\-ये\-दा\-जा\-मांसं}}{\Afootnote{K; \skt{भक्ष\-ये\-दा\-ज\-मामं}\ A}}\skt{ पित्तसमायुतं ।}\edtext{}{\lemma{\skt{पिबेदसृक् । यकृद्वा}}\Afootnote{A; \skt{पिवेदसृक् । यकृद्वा}\ K}}\pend

\pstart\versepar \ledsidenote{\passref{6.45.29}}\edtext{\skt{पला\-श\-वृ\-न्त\-स्व\-रसे}}{\Afootnote{K H; \skt{पला\-श\-वृ\-क्ष\-स्व\-रसे}\ A}}\skt{ विपक्वं }\pend
\pstart\versepar \skt{}\edtext{\skt{सर्पि}}{\Afootnote{K; \skt{सर्पिः}\ A H}}\skt{ }\edtext{\skt{पिबेत्क्षौद्रयुतं}}{\Afootnote{\skt{पिबेत्क्षौद्रयु}\ A; \skt{पिवेत्क्षौद्रयुतं}\ K H}}\skt{ }\edtext{\skt{हिताशी}}{\Afootnote{K H; \emph{om.}\ A}}\skt{ । }\pend
\pstart\versepar \skt{}\edtext{\skt{सशर्करं}}{\Afootnote{K H; \emph{om.}\ A}}\skt{ }\edtext{\skt{क्षीरघृतं}}{\Afootnote{K H; \skt{तं}\ A}}\skt{ }\edtext{\skt{पिबेद्वा}}{\Afootnote{\skt{सुशीतम् ।}\ A; \skt{पिवेद्वा}\ K H}}\skt{ }\pend
\pstart\versepar \skt{}\edtext{\skt{वनस्पतीं}}{\Afootnote{K; \skt{वनस्पतीनां}\ A H}}\skt{ स्वरसैः }\edtext{\skt{कृतम्वा}}{\lemma{\skt{वा}}\Afootnote{K H; \skt{वा सशर्करं क्षीरघृतं पिबेद्वा ॥}\ A}}\skt{ ।}\pend

\pstart\versepar \ledsidenote{\passref{6.45.30}}\skt{द्राक्षा\-मु\-शी\-रा\-ण्यथ पद्मकं सिता }\pend
\pstart\versepar \skt{}\edtext{\skt{पृथ\-क्प\-लां\-श्यु\-दके}}{\Afootnote{K; \skt{पृथ\-क्प\-लां\-शा\-न्यु\-दके}\ A; \skt{पृथ\-क्प\-लां\-शो\-न्यू\-दके}\ H}}\skt{ समावयेत् । }\pend
\pstart\versepar \skt{स्थितां }\edtext{\skt{निशा\-न्त\-द्रु\-धि\-रा\-म\-य\-ञ्ज\-येत्}}{\lemma{\skt{तद्रु\-धि\-रा\-म\-यञ्}}\Afootnote{K H; \skt{तद्रुधिरामयं}\ A}}\skt{ }\pend
\pstart\versepar \skt{पीतं पयो }\edtext{\skt{वाम्वुसमं}}{\Afootnote{K H; \skt{वाम्बुसमं}\ A}}\skt{ }\edtext{\skt{हिताशिना}}{\Afootnote{K H; \skt{हिताशिनः ॥}\ A}}\skt{ ।}\edtext{}{\lemma{\skt{समावयेत् । स्थितां}}\Afootnote{K; \skt{समावपेत् । स्थितं}\ A; \skt{समावयेत् । स्थितं}\ H}}\pend

\pstart\versepar \ledsidenote{\passref{6.45.31}}\edtext{\skt{समाक्षिकं}}{\Afootnote{K H; \skt{तुर\-ङ्ग\-व\-र्चः\-स्व\-रसं समाक्षिकं}\ A}}\skt{ }\edtext{\skt{वाजिशकृद्रसं वा}}{\Afootnote{K H; \emph{om.}\ A}}\skt{ }\pend
\pstart\versepar \skt{}\edtext{\skt{वासाकषायं}}{\Afootnote{K H; \emph{om.}\ A}}\skt{ }\edtext{\skt{सहितं}}{\Afootnote{K H; \emph{om.}\ A}}\skt{ }\edtext{\skt{पिबेद्वा}}{\Afootnote{\skt{पिबे\-त्सि\-ता\-क्षौ\-द्र\-युतं}\ A; \skt{पिवेद्वा ।}\ K; \skt{पि}\ H}}\skt{ । }\pend
\pstart\versepar \skt{}\edtext{\skt{लिह्यात्तथा}}{\Afootnote{K; \skt{वृषस्य वा । लिहेत्तथा}\ A; \skt{वेद्वा । लिह्यात्तथा}\ H}}\skt{ }\edtext{\skt{वास्तु\-क\-वी\-ज\-चूर्णं}}{\Afootnote{K H; \skt{वास्तु\-क\-बी\-ज\-चूर्णं}\ A}}\skt{ }\pend
\pstart\versepar \skt{}\edtext{\skt{क्षौद्रप्लुतं}}{\Afootnote{K; \skt{क्षौद्रान्वितं}\ A; \skt{क्षौद्राप्लुतं}\ H}}\skt{ तण्डु\-ल\-सा\-ह्व\-यम्वा ।}\pend

\pstart\versepar \ledsidenote{\passref{6.45.32}}\edtext{\skt{लिहेत}}{\Afootnote{K H; \skt{लिह्याच्च}\ A}}\skt{ }\edtext{\edtext{\skt{काला\-ञ्ज\-न\-चू\-र्णा\-मेव}}{\lemma{\skt{काला\-ञ्ज\-न\-चू\-र्णाम्}}\Afootnote{K; \skt{लाजाञ्जनचूर्ण}\ A; \skt{काला\-ञ्ज\-न\-चू\-र्ण्णम्}\ H}}}{\lemma{\skt{एव}}\Afootnote{K H; \skt{मेक}\ A}}\skt{ }\pend
\pstart\versepar \skt{}\edtext{\skt{लिहेत}}{\Afootnote{K H; \skt{मेवं सि}\ A}}\skt{ }\edtext{\skt{वा}}{\Afootnote{K H; \skt{ता}\ A}}\skt{ क्षौद्रयुतां तु }\edtext{\skt{काख्यां}}{\Afootnote{K H; \skt{गाख्याम् ।}\ A}}\skt{ । }\pend
\pstart\versepar \skt{}\edtext{\skt{द्राषा\-सि\-ता\-ति\-क्त\-क\-रो\-हि\-णीञ्च}}{\Afootnote{K; \skt{द्राक्षां सितां तिक्तकरोहिणीं च}\ A; \skt{द्राक्षा\-सि\-ता\-न्ति\-का\-क\-रो\-हि\-णीञ्च}\ H}}\skt{ }\pend
\pstart\versepar \skt{}\edtext{\skt{हिमाम्वुना}}{\Afootnote{K H; \skt{हिमाम्बुना}\ A}}\skt{ वा }\edtext{\edtext{\skt{मधुकञ्जयेयं}}{\lemma{\skt{मधुकञ्}}\Afootnote{K H; \skt{मधुके}\ A}}}{\lemma{\skt{जयेयं}}\Afootnote{K H; \skt{न युक्ताम् ॥}\ A}}\skt{ ।}\pend

\pstart\versepar \ledsidenote{\passref{6.45.33ab}}\edtext{\skt{पथ्यामहींस्रां}}{\Afootnote{K; \skt{पथ्यामहिंस्रां}\ A H}}\skt{ }\edtext{\skt{रजनीं}}{\Afootnote{A K; \skt{रजनी}\ H}}\skt{ }\edtext{\skt{घ्टञ्च}}{\Afootnote{K; \skt{घृतं च}\ A H}}\skt{ }\pend
\pstart\versepar \skt{}\edtext{\skt{लिहेत्तथा}}{\Afootnote{K H; \skt{लिह्यात्तथा}\ A}}\skt{ शोणि\-त\-पि\-त्त\-रोगी ।}\pend

\pstart\versepar \ledsidenote{\passref{6.45.36cd}}\skt{मूलानि पुष्पाणि च मातुलुङ्ग्याः }\pend
\pstart\versepar \skt{पिष्ट्वा }\edtext{\skt{पिबे\-त्त\-ण्डु\-ल\-धा\-व\-नेन}}{\Afootnote{A; \skt{पिवे\-त्त\-ण्डु\-ल\-धा\-व\-नेन}\ K H}}\skt{ ।}\pend

\pstart\versepar \ledsidenote{\passref{6.45.37}}\skt{घ्राणप्रवृत्ते }\edtext{\skt{सृजि}}{\Afootnote{K; \skt{ज}\ A; \skt{ऽसृजि}\ H}}\skt{ }\edtext{\skt{नस्य}}{\Afootnote{K H; \skt{लमा}\ A}}\skt{ }\edtext{\skt{मुक्तं}}{\Afootnote{K; \skt{शु दे}\ A; \skt{मुक्तां}\ H}}\skt{ }\pend
\pstart\versepar \skt{}\edtext{\skt{सशरकरां}}{\Afootnote{K; \skt{यं सशर्करं}\ A; \skt{सशर्क्करान्}\ H}}\skt{ }\edtext{\skt{नासिकयोः}}{\Afootnote{K H; \skt{नासिकया}\ A}}\skt{ पयो वा । }\pend
\pstart\versepar \skt{द्राक्षारसं क्षीरघृतं }\edtext{\skt{पिबेद्वा}}{\Afootnote{\skt{पिबेद्वा}\ A; \skt{पिवेद्वा}\ K; \skt{पिवेद्वा}\ H}}\skt{ }\pend
\pstart\versepar \skt{सशर्करं चेक्षुरसं }\edtext{\skt{हितम्वा}}{\lemma{\skt{हितम्}}\Afootnote{K H; \skt{हिमं}\ A}}\skt{ ।}\pend

\pstart\versepar \ledsidenote{\passref{6.45.38}}\edtext{\skt{शितोपचारं}}{\Afootnote{K; \skt{शीतोपचारं}\ A H}}\skt{ मधुरञ्च कुर्याद् }\pend
\pstart\versepar \skt{विशेषतः }\edtext{\skt{शोणि\-त\-पि\-त्त\-रोगै}}{\Afootnote{K; \skt{शोणि\-त\-पि\-त्त\-रोगे ।}\ A H}}\skt{ । }\pend
\pstart\versepar \skt{}\edtext{\skt{शिता\-घृ\-त\-ञ्क्षौ\-द्र\-यु\-तेन}}{\Afootnote{K; \skt{द्राक्षा\-घृ\-त\-क्षौ\-द्र\-सि\-ता\-यु\-तेन}\ A; \skt{सिताघृतं क्षौद्रयुतेन}\ H}}\skt{ }\edtext{\skt{चापि}}{\Afootnote{K H; \emph{om.}\ A}}\skt{ }\pend
\pstart\versepar \skt{}\edtext{\skt{विदा\-रि\-ग\-न्धा\-दि\-ग\-ण\-श्रि\-तेन}}{\Afootnote{K H; \skt{विदा\-रि\-ग\-न्धा\-दि\-वि\-पा\-चि\-तेन ॥}\ A}}\skt{ ।}\pend

\pstart\versepar \ledsidenote{\passref{6.45.39}}\skt{क्षीरेण }\edtext{\skt{चास्था\-प\-न\-म\-ग्र्य\-मुक्तं}}{\Afootnote{K H; \skt{चास्था\-प\-न\-म\-ग्र्य\-मुक्तं}\ A}}\skt{ }\pend
\pstart\versepar \skt{हितं }\edtext{\skt{घ्तं}}{\Afootnote{K; \skt{घृतं}\ A H}}\skt{ चाप्य\-नु\-वा\-स\-नार्थं ॥ }\pend
\pstart\versepar \skt{}\edtext{\skt{मञ्जि\-ष्ठ\-रो\-ध्रा\-ञ्ज\-न\-गै\-रि\-को\-त्पलैः}}{\Afootnote{K; \skt{प्रिय\-ङ्गु\-रो\-ध्रा\-ञ्ज\-न\-गै\-रि\-को\-त्पलैः}\ A; \skt{मञ्जि\-ष्ठ\-लो\-ध्रा\-ञ्ज\-न\-गै\-रि\-को\-त्पलैः}\ H}}\skt{ }\pend
\pstart\versepar \skt{}\edtext{\skt{सुव\-र्ण\-का\-ली\-य\-क\-शं\-ख\-च\-न्दनैः}}{\Afootnote{K H; \skt{सुव\-र्ण\-का\-ली\-य\-क\-र\-क्त\-च\-न्दनैः ॥}\ A}}\skt{ ।}\pend

\pstart\versepar \ledsidenote{\passref{6.45.40}}\edtext{\skt{सिता\-श्व\-ग\-न्धा\-म्बु\-ज\-य\-ष्ठि\-सा\-ह्वयैः}}{\Afootnote{K H; \skt{सिता\-श्व\-ग\-न्धा\-म्बु\-द\-य\-ष्टि\-का\-ह्व\-यैर्}\ A}}\skt{ }\pend
\pstart\versepar \skt{}\edtext{\skt{मृणा\-ल\-सौ\-ग\-न्धि\-क\-तु\-ल्य\-पेषिः}}{\Afootnote{K H; \skt{मृणा\-ल\-सौ\-ग\-न्धि\-क\-तु\-ल्य\-पेषि}\ A}}\skt{ । }\pend
\pstart\versepar \skt{}\edtext{\skt{घृताप्लुतैः}}{\Afootnote{K; \skt{तैः । निरूह्य चैनं पयसा समा\-क्षि\-कै\-र्घृ\-त\-प्लुतैः}\ A; \skt{आस्था\-प्य\-ये\-न\-म्प\-यसा समा\-क्षि\-कै\-र्घृ\-ता\-प्लुतैः}\ H}}\skt{ }\edtext{\skt{शित\-ज\-ला\-व\-से\-वितं}}{\Afootnote{K; \skt{शीत\-ज\-ला\-व\-से\-चि\-तम् ॥}\ A; \skt{शीत\-ज\-ला\-व\-से\-वितं ॥}\ H}}\skt{ ।}\pend

\pstart\versepar \ledsidenote{\passref{6.45.41}}\edtext{\skt{क्षीरोदनं}}{\Afootnote{K; \skt{क्षीरौदनं}\ A; \skt{क्षीरोदन.}\ H}}\skt{ }\edtext{\skt{भ्क्तम\-म\-था\-नु\-वा\-स\-येद्}}{\lemma{\skt{भ्क्तमम्}}\Afootnote{K; \skt{भुक्तम्}\ A; \skt{उक्त}\ H}}\skt{ }\pend
\pstart\versepar \skt{घृतेन }\edtext{\skt{यष्ठी\-म\-धु\-क\-श्रि\-तेन}}{\Afootnote{K H; \skt{यष्टी\-म\-धु\-सा\-धि\-तेन च ।}\ A}}\skt{ । }\pend
\pstart\versepar \skt{अधोवहं }\edtext{\edtext{\skt{शोणितमेष}}{\lemma{\skt{शोणितम्}}\Afootnote{K; \skt{शोणित}\ A H}}}{\lemma{\skt{एष}}\Afootnote{K; \skt{मेष}\ A H}}\skt{ नाशयेत् }\pend
\pstart\versepar \skt{तथातिसारं }\edtext{\skt{रुधिरस्य}}{\Afootnote{K; \skt{रुधिरस्य दुस्तरम् ॥}\ A; \skt{रूधिरस्य दुस्तरं ॥}\ H}}\skt{ ।}\edtext{}{\lemma{\skt{अथा\-नु\-वा\-स\-ये\-द्घृ\-तेन}}\Afootnote{A K; \skt{मथानुवासये घृतेन}\ H}}\pend

\pstart\versepar \ledsidenote{\passref{6.45.42}}\edtext{\skt{कयोगे}}{\Afootnote{K; \skt{विरेकयोगे}\ A H}}\skt{ }\edtext{\skt{ष्वति}}{\Afootnote{K H; \skt{त्वति}\ A}}\skt{ चैव शस्यते }\pend
\pstart\versepar \skt{}\edtext{\skt{वाम्यत्सूरश्च}}{\Afootnote{K; \skt{वाम्यश्च}\ A; \skt{वाम्यत्सू}\ H}}\skt{ }\edtext{\skt{रक्ते}}{\Afootnote{A K; \skt{रक्तम्}\ H}}\skt{ विजिते }\edtext{\skt{वलान्वितः}}{\Afootnote{K H; \skt{बलान्वितः ॥}\ A}}\skt{ ।}\pend

\pstart \ledsidenote{\passref{6.45.43}}\skt{एवं विधा }\edtext{\skt{त्तरवस्तयञ्च}}{\Afootnote{K; \skt{उत्त\-र\-बा\-स्त\-यश्च}\ A; \skt{त्तरवस्तुयश्च}\ H}}\skt{ }\edtext{\skt{मूत्राशयस्थे}}{\Afootnote{A K; \skt{मूत्रासयस्थे}\ H}}\skt{ }\edtext{\skt{रुधिरे}}{\Afootnote{A K; \skt{रूधिरे}\ H}}\skt{ }\edtext{\skt{विधेयाः}}{\Afootnote{K H; \skt{विधेयाः । प्रवृत्तरक्तेषु च पायुजेषु कुर्याद्विधानं खलु रक्तपैत्तम् ॥}\ A}}\skt{ ।}\pend

\pstart\versepar \ledsidenote{\passref{6.45.44}}\edtext{\skt{असृग्दरे}}{\Afootnote{K H; \skt{विधिश्चासृग्दरे}\ A}}\skt{ }\edtext{\edtext{\skt{प्येष}}{\lemma{\skt{प्य्}}\Afootnote{K; \skt{ऽ}\ A; \emph{om.}\ H}}}{\lemma{\skt{एष}}\Afootnote{K; \skt{प्येष}\ A H}}\skt{ }\edtext{\skt{विधिः}}{\Afootnote{K H; \emph{om.}\ A}}\skt{ }\edtext{\skt{स्त्रीणाञ्कार्यो}}{\Afootnote{K; \skt{स्त्रीणां कार्यो}\ A; \skt{स्त्रीणाञ्कार्यो}\ H}}\skt{ विजानता । }\pend
\pstart\versepar \skt{शस्त्रकर्मणि }\edtext{\skt{रक्तञ्यस्यातीव}}{\Afootnote{K H; \skt{रक्तं च यस्यातीव}\ A}}\skt{ प्रवर्त्तते ।}\pend

\pstart\versepar \ledsidenote{\passref{6.45.45}}\skt{त्रयाणामपि दोषाणां }\edtext{\skt{शोणितस्य च}}{\Afootnote{K H; \skt{शोणितेऽपि च}\ A}}\skt{ सर्वशः । }\pend
\pstart\versepar \skt{लिङ्गान्यालोक्य }\edtext{\edtext{\skt{वित\-रे\-च्चि\-कि\-त्सि\-त\-म\-न\-न्त\-र\-मिति}}{\lemma{\skt{वित\-रे\-च्चि\-कि\-त्सि\-तम्}}\Afootnote{K H; \skt{कर्तव्यं चिकित्सितम्}\ A}}}{\lemma{\skt{इति}}\Afootnote{K H; \emph{om.}\ A}}\skt{ ॥}\pend

\pstart \ledsidenote{\passref{6.45.trailer}}\edtext{\skt{कायचि}}{\Afootnote{K; \skt{इति सुश्रु\-त\-सं\-हि\-ता\-या\-मु\-त्त\-र\-त\-न्त्रा\-न्त\-र्गते काय\-चि\-कि\-त्सा\-तन्त्रे रक्त\-पि\-त्त\-प्र\-ति\-षेधो नाम (सप्तमोऽध्यायः, आदितः) पञ्च\-च\-त्वा\-रिं\-श\-त्त\-मोऽ\-ध्यायः ॥४५ ॥}\ A; \skt{उत्तरतन्त्रे रक्त\-पि\-त्त\-चि\-कित्सा द्वाच\-त्वा\-रिं\-स\-त्तमः कायचिकित्सायां सप्तमः ॥ ॥}\ H}}\skt{ ॥ ॥ ॥}\pend

\endnumbering
\beginnumbering
\eledsection*{Adhyāya 46}
\pstart \ledsidenote{\passref{6.46.1}}\skt{अथातो मूर्च्छा प्रतिषेधं }\edtext{\skt{व्याख्यास्यामः}}{\Afootnote{A H; \emph{om.}\ K}}\skt{ ॥}\pend

\pstart\versepar \ledsidenote{\passref{6.46.3}}\skt{क्षीणस्य बहुदोषस्य विरु\-द्धा\-हा\-र\-से\-विनः । }\pend
\pstart\versepar \skt{}\edtext{\skt{वेगा\-घा\-ता\-द\-भी\-घा\-ता\-द्धी\-न\-स\-त्वस्य वा}}{\Afootnote{A K; \skt{वेगा\-घा\-ता\-द\-ती\-पाता द्धीनसत्वस्य वा}\ H}}\skt{ पुनः ।}\pend

\pstart\versepar \ledsidenote{\passref{6.46.4}}\edtext{\edtext{\skt{कर\-णा\-य\-त\-ने\-षू\-ग्रा\-बा\-ह्ये\-ष्व\-भ्य\-न्त\-रेषु}}{\lemma{\skt{कर\-णा\-य\-त\-ने\-षू\-ग्रा\-बा\-ह्येष्व्}}\Afootnote{A K; \skt{कर\-णा\-य\-त\-ने\-षूग्रा बाह्येष्वा}\ H}}\skt{ च}}{\lemma{\skt{अभ्यन्तरेषु च}}\Afootnote{K; \skt{आभ्यन्तरेषु च ।}\ A; \skt{भ्यन्तरेषु च ।}\ H}}\skt{ । }\pend
\pstart\versepar \skt{निविशन्ते यदा }\edtext{\edtext{\skt{दोषास्तदा}}{\lemma{\skt{दोषास्}}\Afootnote{A K; \skt{दोषा}\ H}}}{\lemma{\skt{तदा}}\Afootnote{A K; \skt{स्तथा}\ H}}\skt{ }\edtext{\skt{मूर्च्छति}}{\Afootnote{K H; \skt{मूर्च्छन्ति}\ A}}\skt{ }\edtext{\skt{मानवः}}{\Afootnote{K H; \skt{मानवाः ॥}\ A}}\skt{ ।}\pend

\pstart\versepar \ledsidenote{\passref{6.46.6}}\skt{संज्ञावहासु नाडीषु }\edtext{\skt{पिथि\-ता\-स्व\-नि\-ला\-दिषु}}{\Afootnote{K H; \skt{पिहि\-ता\-स्व\-नि\-ला\-दिभिः ।}\ A}}\skt{ । }\pend
\pstart\versepar \skt{}\edtext{\skt{तमोभ्युपैति}}{\Afootnote{K; \skt{तमोऽभ्युपैति}\ A; \skt{तमोऽभ्युदेति}\ H}}\skt{ सहसा }\edtext{\skt{सुख\-दः\-ख\-व्य\-पो\-ह\-कृत्}}{\Afootnote{K; \skt{सुख\-दुः\-ख\-व्य\-पो\-ह\-कृत् ॥}\ A H}}\skt{ ।}\pend

\pstart\versepar \ledsidenote{\passref{6.46.7}}\skt{सुख\-दुः\-ख\-व्य\-पो\-हाच्च }\edtext{\skt{नरः}}{\Afootnote{K; \emph{om.}\ A H}}\skt{ नरः पतति काष्ठवत् । }\pend
\pstart\versepar \skt{मोहो }\edtext{\skt{मूर्च्छति}}{\Afootnote{K; \skt{मूर्च्छेति}\ A H}}\skt{ }\edtext{\edtext{\skt{तामाहुः}}{\lemma{\skt{ताम्}}\Afootnote{A K; \skt{ता}\ H}}}{\lemma{\skt{आहुः}}\Afootnote{K; \skt{प्राहुः}\ A; \skt{माहुः}\ H}}\skt{ षड्विधा सा }\edtext{\skt{च}}{\Afootnote{K H; \skt{प्र}\ A}}\skt{ }\edtext{\skt{कीर्त्यते}}{\Afootnote{K H; \skt{कीर्तिता ॥}\ A}}\skt{ ।}\pend

\pstart\versepar \ledsidenote{\passref{6.46.8}}\skt{वातादिभिः शोणितेन मद्येन च }\edtext{\skt{विशेण च}}{\Afootnote{K; \skt{विषेण च ।}\ A H}}\skt{ । }\pend
\pstart\versepar \skt{}\edtext{\skt{षट्स्वेतासु च}}{\Afootnote{K; \skt{षट्स्वप्येतासु}\ A; \skt{षड्स्वेतासु च}\ H}}\skt{ }\edtext{\skt{पैत्तीन्तु}}{\Afootnote{K; \skt{पित्तं हि}\ A; \skt{पित्तीन्तु}\ H}}\skt{ प्रभु\-त्वे\-ना\-व\-ति\-ष्ठते ।}\pend

\pstart\versepar \ledsidenote{\passref{6.46.8.1}}\skt{हृत्पीडा जृम्भणं ग्लानिः संज्ञा\-दौ\-र्ब\-ल्य\-मेव च । }\pend
\pstart\versepar \skt{सर्वासां पूर्वरूपाणि यथास्वञ्च विभावयेत् ।}\pend

\pstart\versepar \ledsidenote{\passref{6.46.9}}\edtext{\skt{अप\-स्मा\-र\-ण\-लि\-ङ्गानि}}{\Afootnote{K; \skt{अप\-स्मा\-रो\-क्त\-लि\-ङ्गानि}\ A; \skt{अपस्मारेण लिङ्गानि}\ H}}\skt{ }\edtext{\skt{तासामुक्तानि}}{\Afootnote{K; \skt{तासामुक्तानि}\ A; \skt{तासां मुक्तानि}\ H}}\skt{ तत्वतः । }\pend
\pstart\versepar \skt{}\edtext{\skt{पृथि\-व्या\-प\-स्त\-मो\-रूपं}}{\Afootnote{K; \skt{पृथि\-व्य\-म्भ\-स्त\-मो\-रूपं}\ A; \skt{पृथि\-व्या\-प\-स्त\-नौ\-रूपं}\ H}}\skt{ }\edtext{\skt{रक्तं}}{\Afootnote{K H; \skt{रक्त}\ A}}\skt{ गन्धश्च }\edtext{\skt{तत्र}}{\Afootnote{K H; \skt{तन्म}\ A}}\skt{ यः ।}\pend

\pstart\versepar \ledsidenote{\passref{6.46.10}}\skt{तस्माद्रक्तस्य गन्धेन मूर्च्छन्ति भुवि मानवाः । }\pend
\pstart\versepar \skt{}\edtext{\skt{द्रव्य\-स्व\-भा\-व\-मि\-त्येके}}{\lemma{\skt{द्रव्यस्वभावम्}}\Afootnote{K H; \skt{द्रव्यस्वभाव}\ A}}\skt{ }\edtext{\skt{दृष्ट्वा}}{\Afootnote{K H; \emph{om.}\ A}}\skt{ }\edtext{\skt{यदपि}}{\Afootnote{K H; \emph{om.}\ A}}\skt{ }\edtext{\skt{मुह्यति}}{\Afootnote{K H; \skt{स्थितास्तु विषमद्ययोः ॥}\ A}}\skt{ ।}\pend

\pstart\versepar \ledsidenote{\passref{6.46.11}}\edtext{\skt{गुण\-स्ती\-व्र\-त\-र\-त्वेन}}{\lemma{\skt{गुणस्}}\Afootnote{K; \skt{गुणास्}\ A H}}\skt{ स्थितास्तु विषमद्ययोः । }\pend
\pstart\versepar \skt{त एव }\edtext{\skt{तस्मादाभ्यां तु}}{\Afootnote{K H; \skt{तस्माज्जायेत}\ A}}\skt{ मोहो }\edtext{\skt{जायेद्यथेरिता}}{\Afootnote{K H; \skt{ताभ्यां यथेरितः ॥}\ A}}\skt{ ।}\pend

\pstart \ledsidenote{\passref{6.46.11ef}}\edtext{\edtext{\skt{स्तब्धा\-ङ्ग\-दृ\-ष्टि\-र\-सृजा}}{\lemma{\skt{स्तब्धाङ्गदृष्टिर्}}\Afootnote{H; \emph{om.}\ A K\textsuperscript{ac}; \skt{तप्धाङ्गदृष्टि}\ K\textsuperscript{pc}}}}{\lemma{\skt{असृजा}}\Afootnote{\emph{om.}\ A K\textsuperscript{ac}; \skt{रसृजा}\ K\textsuperscript{pc}; \skt{असृजाङ्}\ H}}\skt{ }\edtext{\skt{गूढोच्छ्वासश्च}}{\Afootnote{K\textsuperscript{pc} H; \emph{om.}\ A K\textsuperscript{ac}}}\skt{ }\edtext{\skt{मूर्च्छितः}}{\Afootnote{K\textsuperscript{pc} H; \emph{om.}\ A K\textsuperscript{ac}}}\skt{ ।}\pend

\pstart\versepar \ledsidenote{\passref{6.46.12}}\skt{मद्येन }\edtext{\skt{विलपं}}{\Afootnote{K H; \skt{विलपंश्}\ A}}\skt{ च्छेते }\edtext{\skt{निष्टनं}}{\Afootnote{K H; \skt{नष्टवि}\ A}}\skt{ }\edtext{\skt{भ्रान्तचेतसः}}{\Afootnote{K H; \skt{भ्रान्तमानसः ।}\ A}}\skt{ । }\pend
\pstart\versepar \skt{गात्राणि }\edtext{\skt{विक्षिपं}}{\Afootnote{K; \skt{विक्षिप}\ A H}}\skt{ }\edtext{\skt{भूयो}}{\Afootnote{K; \skt{न्भूमौ जरां}\ A; \skt{न्भूयो}\ H}}\skt{ }\edtext{\skt{यावत्पक्तिं न}}{\Afootnote{K H; \skt{यावन्न}\ A}}\skt{ याति तत् ।}\pend

\pstart\versepar \ledsidenote{\passref{6.46.13}}\edtext{\skt{वेपथुः}}{\Afootnote{K H; \skt{वेपथु}\ A}}\skt{ }\edtext{\skt{स्वप्नतृष्णास्यु}}{\Afootnote{K H; \skt{स्वप्नतृष्णाः स्युः}\ A}}\skt{ }\edtext{\skt{स्तम्भश्च}}{\Afootnote{A K; \skt{स्तमश्च}\ H}}\skt{ }\edtext{\skt{विशमूर्च्छिते}}{\Afootnote{K; \skt{विषमूर्च्छिते}\ A H}}\skt{ । }\pend
\pstart\versepar \skt{वेदितव्यं तीव्रतरं यथास्वं विषलक्षणैः ॥}\pend

\pstart\versepar \ledsidenote{\passref{6.46.14}}\skt{सेकावगाहौ मणयः सहाराः शीताः प्रदेहाः व्यजनानिलाश्च । }\pend
\pstart\versepar \skt{शीतानि }\edtext{\skt{पानानि च}}{\Afootnote{A K H\textsuperscript{pc}; \skt{पानानि}\ H\textsuperscript{ac}}}\skt{ गन्धवन्ति सर्वासु मूर्च्छा\-स्व\-नि\-वा\-रि\-तानि ।}\pend

\pstart \ledsidenote{\passref{6.46.15}}\edtext{\skt{शीतेन}}{\Afootnote{K; \emph{om.}\ A H}}\skt{ }\edtext{\skt{तोयेन}}{\Afootnote{K; \emph{om.}\ A H}}\skt{ }\edtext{\skt{तथा}}{\Afootnote{K; \emph{om.}\ A H}}\skt{ }\edtext{\skt{सलेह्यात् । क्षौद्रेण}}{\Afootnote{K; \emph{om.}\ A H}}\skt{ }\edtext{\skt{कृष्णां ससितां}}{\Afootnote{K; \skt{सिता\-प्रि\-या\-ले\-क्षु\-र\-स\-प्लु\-तानि द्राक्षा\-म\-धू\-क\-स्व\-र\-सान्}\ A; \skt{सश\-र्क्क\-रा\-णी\-क्षु\-र\-स\-म्पि\-यालं द्राक्षा\-म\-धू\-क\-स्व\-र\-सान्}\ H}}\skt{ }\edtext{\skt{हिताय}}{\Afootnote{K; \skt{वितानि । खर्जू\-र\-का\-श्म\-र्य\-रसैः शृतानि पानानि सर्पींषि च जीवनानि ॥}\ A; \skt{वितानि । खर्ज्जू\-र\-का\-श्म\-र्य\-र\-सा\-न्वि\-तानि तानि सर्प्पींषि च जीवनानि ॥}\ H}}\skt{ ।}\pend

\pstart\versepar \ledsidenote{\passref{6.46.18}}\skt{कुर्याच्च नासा\-व\-द\-ना\-व\-रोधं क्षीरं }\edtext{\edtext{\skt{पिबेद्वाप्यथ}}{\lemma{\skt{पिबेद्वाप्य्}}\Afootnote{A K; \skt{पिबेद्वा}\ H}}}{\lemma{\skt{अथ}}\Afootnote{A K; \skt{प्यथ}\ H}}\skt{ }\edtext{\skt{मानुषीणां}}{\Afootnote{A K; \skt{मानुषाणाम् ।}\ H}}\skt{ । }\pend
\pstart\versepar \skt{}\edtext{\skt{मूर्च्छाप्रसक्तान्तु}}{\Afootnote{K H; \skt{मूर्च्छां प्रसक्तां तु}\ A}}\skt{ }\edtext{\edtext{\skt{शिरो\-व\-रे\-कै\-र्ज\-ये\-द\-भीक्ष्णं}}{\lemma{\skt{शिरोवरेकैर्}}\Afootnote{K; \skt{शिरोविरेकै}\ A H}}}{\lemma{\skt{जयेदभीक्ष्णं}}\Afootnote{K; \skt{र्जयेदभीक्ष्णं}\ A H}}\skt{ वमनैश्च तीक्ष्णैः ।}\pend

\pstart\versepar \ledsidenote{\passref{6.46.19}}\edtext{\skt{हरी\-त\-की\-क्वा\-थ\-कृतं}}{\Afootnote{K H; \skt{हरी\-त\-की\-क्वा\-थ\-शृतं}\ A}}\skt{ }\edtext{\skt{पिबेद्वा}}{\Afootnote{K H; \skt{घृतं वा}\ A}}\skt{ }\edtext{\skt{धात्री\-फ\-ला\-नां\-स्व\-रसे}}{\Afootnote{K H; \skt{धात्रीफलानां स्वरसैः}\ A}}\skt{ कृतं वा । }\pend
\pstart\versepar \skt{द्राक्षा\-सि\-ता\-दा\-डि\-म\-ला\-ज\-वन्ति शीतानि }\edtext{\skt{नीलो\-त्प\-ल\-प\-द्म\-वन्ति}}{\Afootnote{A K; \skt{नीलो\-त्प\-ल\-प\-द्म\-यन्ति ॥}\ H}}\skt{ ।}\pend

\pstart\versepar \ledsidenote{\passref{6.46.20ab}}\skt{पिबत्कषायाणि च गन्धवन्ति । पित्तज्वरं याति शमन्न यन्ति । }\pend
\pstart\versepar \skt{}\edtext{\skt{प्रभू\-त\-दो\-ष\-स्तम}}{\lemma{\skt{तम}}\Afootnote{\skt{तमसो तिरे\-का\-त्सं\-मू\-र्च्छितो नैव विबुध्यते यः ।}\ K}}\pend

\pstart \ledsidenote{\passref{6.46.21}}\edtext{\skt{सव्यस्तसंज्ञो}}{\Afootnote{K H; \skt{संन्यस्तसंज्ञो}\ A}}\skt{ }\edtext{\skt{भृशदुश्चिकित्स्यो}}{\Afootnote{K H; \skt{भृशदुश्चिकित्स्यो}\ A}}\skt{ }\edtext{\skt{ज्ञेयस्तथा}}{\lemma{\skt{तथा}}\Afootnote{K H; \skt{तदा}\ A}}\skt{ बुद्धिमता }\edtext{\skt{मनुष्यः}}{\Afootnote{K H; \skt{मनुष्याः । यथा मलोष्टं सलिले निषिक्तं समु\-द्धे\-रे\-दा\-श्व\-वि\-ली\-न\-मेव ॥}\ A}}\skt{ ।}\pend

\pstart \ledsidenote{\passref{6.46.21a}}\edtext{\skt{मृतोपमन्तु}}{\Afootnote{K H; \emph{om.}\ A}}\skt{ }\edtext{\skt{त्वरया}}{\Afootnote{K H; \emph{om.}\ A}}\skt{ }\edtext{\skt{चिकि\-त्से\-दा\-श्वे\-वतं}}{\Afootnote{K H; \emph{om.}\ A}}\skt{ }\edtext{\skt{मृत्युवशं}}{\Afootnote{K H; \emph{om.}\ A}}\skt{ }\edtext{\skt{प्रयान्तं}}{\Afootnote{K H; \emph{om.}\ A}}\skt{ ।}\pend

\pstart \ledsidenote{\passref{6.46.22cd}}\edtext{\skt{तीक्ष्णा\-ञ्ज\-ना\-भ्य\-ञ्ज\-न\-धू\-म\-यो\-गै\-स्तथा}}{\lemma{\skt{तीक्ष्णा\-ञ्ज\-ना\-भ्य\-ञ्ज\-न\-धू\-म\-यो\-गैस्}}\Afootnote{A K; \skt{तीक्ष्णा…\-भ्य\-ञ्ज\-न\-धू\-म\-यो\-गैस्}\ H}}\skt{ }\edtext{\skt{नखा\-भ्य\-न्त\-र\-सू\-ति\-पातैः}}{\Afootnote{K H; \skt{नखाभ्यन्तरतो त्रपातैः ॥}\ A}}\skt{ ।}\pend

\pstart \ledsidenote{\passref{6.46.23ab}}\edtext{\edtext{\skt{वादि\-त्र\-गी\-ता\-नु\-र\-यै\-र\-पू\-र्व्वै\-र्वि\-स्मा\-पनै}}{\lemma{\skt{वादि\-त्र\-गी\-ता\-नु\-र\-यै\-र\-पू\-र्व्वैर्}}\Afootnote{K H; \skt{वादि\-त्र\-गी\-ता\-नु\-न\-यै\-र\-पू\-र्वैर्}\ A}}}{\lemma{\skt{विस्मापनै}}\Afootnote{K; \skt{विघट्टनैर्}\ A; \skt{व्विस्मापनैर्}\ H}}\skt{ }\edtext{\skt{गुप्त\-फ\-ल\-प्र\-कर्षैः}}{\Afootnote{K; \skt{गुप्त\-फ\-ला\-व\-घर्षैः ।}\ A; \skt{गुप्तफलैः प्रकर्षैः ॥}\ H}}\skt{ ।}\pend

\pstart \ledsidenote{\passref{6.46.23ef}}\edtext{\edtext{\edtext{\edtext{\skt{तद्व\-त्प्र\-स\-क्तै\-र्व\-म\-ना\-नु\-लो\-म्यै\-स्ती\-क्ष्नै\-र्वि\-शुद्धं}}{\lemma{\skt{तद्वत्प्रसक्तैर्}}\Afootnote{K; \emph{om.}\ A; \skt{एतत्प्रसक्तैर्}\ H}}}{\lemma{\skt{वम\-ना\-नु\-लो\-म्यैस्}}\Afootnote{K H; \emph{om.}\ A}}}{\lemma{\skt{तीक्ष्नैर्}}\Afootnote{K H; \emph{om.}\ A}}}{\lemma{\skt{विशुद्धं}}\Afootnote{K H; \emph{om.}\ A}}\skt{ }\edtext{\skt{लघुपथ्यभोज्यैः}}{\Afootnote{K H; \emph{om.}\ A}}\skt{ ।}\pend

\pstart\versepar \ledsidenote{\passref{6.46.24cdef}}\skt{फलत्रिकैः }\edtext{\skt{पिप्प\-लि\-चि\-त्र\-काढ्यैः}}{\Afootnote{K; \skt{चित्रकना}\ A; \skt{पिप्प\-लि\-चि\-त्र\-काद्यैः}\ H}}\skt{ }\edtext{\edtext{\skt{शिलाजतोश्चापि}}{\lemma{\skt{शिलाजतोश्}}\Afootnote{K H; \skt{गराढ्यैस्}\ A}}}{\lemma{\skt{चापि}}\Afootnote{K H; \emph{om.}\ A}}\skt{ तथा }\edtext{\skt{प्रयोगैः}}{\Afootnote{K H; \skt{श्मजाताज्जतुनः प्रयोगैः ।}\ A}}\skt{ । }\pend
\pstart\versepar \skt{}\edtext{\edtext{\skt{सश\-र्क\-रौ\-र्मा\-स\-मु\-पा\-च\-रेत}}{\lemma{\skt{सशर्करौर्}}\Afootnote{K; \skt{सशर्करैर्}\ A H}}}{\lemma{\skt{उपाचरेत}}\Afootnote{K; \skt{उपक्रमेत}\ A; \skt{उपचरेत}\ H}}\skt{ विशेषतो जीर्ण्णघृतं }\edtext{\skt{पिबेच्च}}{\Afootnote{K H; \skt{स पाय्यः ॥}\ A}}\skt{ ।}\pend

\pstart\versepar \ledsidenote{\passref{6.46.25}}\skt{यथास्वञ्च ज्वरघ्नानि }\edtext{\skt{कषायाणि}}{\Afootnote{K; \skt{कषायाण्यु}\ A; \skt{कषायानि}\ H}}\skt{ प्रयोजयेत् । }\pend
\pstart\versepar \skt{सर्व\-मू\-र्च्छा\-प\-री\-तानां विषजायां }\edtext{\edtext{\skt{विषापहमिति}}{\lemma{\skt{विषापहम्}}\Afootnote{A K; \skt{विषापह}\ H}}}{\lemma{\skt{इति}}\Afootnote{K; \emph{om.}\ A; \skt{मिति ॥}\ H}}\skt{ ॥}\edtext{}{\lemma{\skt{प्रयोजयेत् । सर्व\-मू\-र्च्छा\-प\-री\-तानां}}\Afootnote{K H; \skt{पयोजयेत् । सर्व\-मू\-र्च्छा\-प\-री\-तानां}\ A}}\pend

\pstart \ledsidenote{\passref{6.46.trailer}}\edtext{\skt{कायचि}}{\Afootnote{K; \skt{इति सुश्रु\-त\-सं\-हि\-ता\-या\-मु\-त्त\-र\-त\-न्त्रा\-न्त\-र्गते काय\-चि\-कि\-त्सा\-तन्त्रे मूर्च्छाप्रतिषेधो नाम (अष्टमोऽध्यायः, आदितः) षट्च\-त्वा\-रिं\-श\-त्त\-मोऽ\-ध्यायः ॥४६ ॥}\ A; \skt{उत्तरतन्त्रे मूर्च्छाचिकित्सा त्रिच\-त्वा\-रिं\-श\-त्तमः काय\-चि\-कि\-त्सा\-या\-म\-ष्ट\-मो\-ध्यायः ॥}\ H}}\skt{ ॥ ८ ॥ ० ॥}\pend

\endnumbering
\beginnumbering
\eledsection*{Adhyāya 47}
\pstart \ledsidenote{\passref{6.47.1}}\skt{अथातः पाना\-त्य\-य\-प्र\-ति\-षेधं }\edtext{\skt{व्याख्यास्यामः}}{\Afootnote{A H; \skt{व्या व्याख्यास्यामः ॥}\ K}}\skt{ ॥}\pend

\pstart\versepar \ledsidenote{\passref{6.47.3}}\edtext{\edtext{\skt{मद्यमुष्णन्तथा}}{\lemma{\skt{मद्यम्}}\Afootnote{K; \skt{मद्य}\ A H}}}{\lemma{\skt{उष्णन्}}\Afootnote{K; \skt{मुष्णं}\ A H}}\skt{ तीक्ष्णं सूक्ष्मं }\edtext{\skt{विशदमेव}}{\lemma{\skt{विशदम्}}\Afootnote{A K; \skt{विषदम्}\ H}}\skt{ च । }\pend
\pstart\versepar \skt{}\edtext{\edtext{\skt{रूक्ष\-मा\-शु\-क\-र\-ञ्चैव}}{\lemma{\skt{रूक्षम्}}\Afootnote{K; \skt{रूक्ष}\ A H}}}{\lemma{\skt{आशुकरञ्}}\Afootnote{K; \skt{माशुकरं}\ A H}}\skt{ व्यवायि च विकाशि च ।}\pend

\pstart\versepar \ledsidenote{\passref{6.47.4}}\edtext{\skt{औष्ण्या\-च्छी\-तो\-प\-चा\-रञ्च}}{\lemma{\skt{च}}\Afootnote{K H; \skt{त}\ A}}\skt{ }\edtext{\skt{तैक्ष्ण्याद्धन्ति}}{\Afootnote{K; \skt{त्तैक्ष्ण्याद्धन्ति}\ A; \skt{तैक्ष्णा हन्ति}\ H}}\skt{ }\edtext{\skt{मनो}}{\Afootnote{K; \skt{मनो}\ A H}}\skt{ गतिं । }\pend
\pstart\versepar \skt{}\edtext{\skt{पुंस्त्वं}}{\Afootnote{K H; \emph{om.}\ A}}\skt{ }\edtext{\skt{कफञ्च}}{\Afootnote{K H; \emph{om.}\ A}}\skt{ सूक्ष्म\-त्वा\-द्धि\-म\-त्य\-व\-य\-वा\-न्व\-हृत् । }\pend
\pstart\versepar \skt{वैषम्यात्तु }\edtext{\skt{कफं}}{\Afootnote{K H; \skt{त्क}\ A}}\skt{ }\edtext{\skt{हन्याद्रूचिवापि}}{\Afootnote{K; \skt{फ}\ A; \skt{हन्याद्रुचिम्वापि}\ H}}\skt{ }\edtext{\skt{प्रवर्त्तयेत्}}{\Afootnote{K H; \skt{शुक्रनुत् ॥}\ A}}\skt{ ।}\edtext{}{\lemma{\skt{सूक्ष्म\-त्वा\-द्धि\-म\-त्य\-व\-य\-वा\-न्व\-हृत् । वैषम्यात्तु}}\Afootnote{K; \skt{विश\-त्य\-व\-य\-वा\-न्सौ\-क्ष्म्या\-द्वै\-शद्या}\ A; \skt{सूक्ष्म\-त्वा\-द्वि\-श\-त्य\-व\-य\-वा\-न्ब\-हून् ॥ वैशद्यान्तु}\ H}}\pend

\pstart\versepar \ledsidenote{\passref{6.47.5}}\skt{मारुतं कोप\-ये\-द्रौ\-क्ष्या\-दा\-शु\-त्वा\-दा\-शु\-क\-र्म\-कृत् । }\pend
\pstart\versepar \skt{हर्षदञ्च }\edtext{\skt{व्यवा\-यि\-त्वा\-द्वि\-का\-शि\-त्वा\-न्वि\-स\-र्पिणं}}{\Afootnote{K; \skt{व्यवा\-यि\-त्वा\-द्वि\-का\-शि\-त्वा\-द्वि\-स\-र्पति ॥}\ A; \skt{व्यवा\-यि\-त्वा\-द्वि\-का\-सि\-त्वा\-द्वि\-सर्प्प्णं ।}\ H}}\skt{ ।}\edtext{}{\lemma{\skt{कोप\-ये\-द्रौ\-क्ष्या\-दा\-शु\-त्वा\-दा\-शु\-क\-र्म\-कृत् । हर्षदञ्च}}\Afootnote{K; \skt{कोप\-ये\-द्रौ\-क्ष्या\-दा\-शु\-त्वा\-च्चा\-शु\-क\-र्म\-कृत् । हर्षदं च}\ A; \skt{कोप\-ये\-द्रौ\-क्ष्या\-दा\-शु\-त्वा\-दा\-शु\-क\-र्म\-कृत् ॥ हर्षदं च}\ H}}\pend

\pstart\versepar \ledsidenote{\passref{6.47.6}}\edtext{\edtext{\skt{तद\-म्व्ल\-र\-स\-त\-स्तूक्तं}}{\lemma{\skt{तद\-म्व्ल\-र\-स\-तस्}}\Afootnote{K; \skt{तदम्लं रसतः}\ A; \skt{तदम्ल रसतस्}\ H}}}{\lemma{\skt{तूक्तं}}\Afootnote{K H; \skt{प्रोक्तं}\ A}}\skt{ }\edtext{\edtext{\skt{लघु\-दी\-प\-न\-मेव}}{\lemma{\skt{लघुदीपनम्}}\Afootnote{K H; \skt{लघु रोच\-न\-दी\-प\-नम् ।}\ A}}\skt{ च}}{\lemma{\skt{एव च}}\Afootnote{K H; \emph{om.}\ A}}\skt{ । }\pend
\pstart\versepar \skt{}\edtext{\skt{केचि\-ल्ल\-व\-ण\-व\-र्ज्यांस्तु}}{\Afootnote{A K; \skt{केचि\-ल्ल\-व\-ण\-व\-र्ज्ज्यास्तु}\ H}}\skt{ रसा\-न\-त्रा\-दि\-शन्ति }\edtext{\skt{ह}}{\Afootnote{K H; \skt{हि ॥}\ A}}\skt{ ।}\pend

\pstart\versepar \ledsidenote{\passref{6.47.7}}\edtext{\edtext{\skt{स्निग्धै\-स्त\-द\-म्व्लै\-र्मां\-सैश्च}}{\lemma{\skt{तदम्व्लैर्}}\Afootnote{K; \skt{तदन्नै}\ A; \skt{तदम्लैर्}\ H}}}{\lemma{\skt{मांसैश्च}}\Afootnote{K H; \skt{र्मांसैश्च}\ A}}\skt{ भक्ष्यैश्च }\edtext{\skt{सह}}{\Afootnote{A K H\textsuperscript{pc}; \skt{मह}\ H\textsuperscript{ac}}}\skt{ सेवितं । }\pend
\pstart\versepar \skt{}\edtext{\skt{भवेदायुः प्रकर्षाय}}{\Afootnote{K H; \skt{भवे\-दा\-युः\-प्र\-क\-र्षाय}\ A}}\skt{ }\edtext{\skt{वला\-यो\-प\-च\-याय}}{\Afootnote{K; \skt{बला\-यो\-प\-च\-याय}\ A H}}\skt{ च ।}\pend

\pstart\versepar \ledsidenote{\passref{6.47.8}}\skt{काम्यता }\edtext{\edtext{\skt{मन\-स\-स्तु\-ष्टि\-स्तेजो}}{\lemma{\skt{मनसस्}}\Afootnote{K; \skt{मनस}\ A H}}}{\lemma{\skt{तुष्टिस्}}\Afootnote{K; \skt{स्तुष्टिर्धैर्यं}\ A; \skt{सि तुष्टि स्}\ H}}\skt{ }\edtext{\skt{विक्रम}}{\Afootnote{K H; \skt{ऽतिविक्रमः ।}\ A}}\skt{ }\edtext{\skt{एव च}}{\Afootnote{K H; \emph{om.}\ A}}\skt{ । }\pend
\pstart\versepar \skt{विधि\-व\-त्से\-व्य\-माने तु मद्ये सन्निहिता गुणाः ।}\pend

\pstart\versepar \ledsidenote{\passref{6.47.9}}\edtext{\edtext{\skt{तदे\-वा\-न\-न्न\-म\-ज्ञेन}}{\lemma{\skt{तदेवानन्नम्}}\Afootnote{K H; \skt{तदेवानन्न}\ A}}}{\lemma{\skt{अज्ञेन}}\Afootnote{K H; \skt{मज्ञेन}\ A}}\skt{ }\edtext{\skt{सेव्य\-मा\-न\-म\-मा\-त्रया}}{\Afootnote{K; \skt{सेव्य\-मा\-न\-म\-मा\-त्रया ।}\ A; \skt{से वामा\-न\-म\-मा\-त्रया ॥}\ H\textsuperscript{ac}; \skt{से वाव्य\-मा\-न\-म\-मा\-त्रया ॥}\ H\textsuperscript{pc}}}\skt{ । }\pend
\pstart\versepar \skt{}\edtext{\skt{कायाग्निना ह्यग्निसमं}}{\lemma{\skt{कायाग्निना ह्य्}}\Afootnote{K; \skt{कयाग्निना ह्य्}\ A H}}\skt{ समेत्य कुरुते मदं ।}\pend

\pstart\versepar \ledsidenote{\passref{6.47.10}}\skt{मदेन करणानान्तु भावान्यत्वे कृते सति । }\pend
\pstart\versepar \skt{निगूढमपि }\edtext{\skt{भावं स्वं}}{\Afootnote{A K; \skt{भाव स्वं}\ H}}\skt{ प्रकाशीकुरुते }\edtext{\skt{वशः}}{\Afootnote{K H; \skt{ऽवशः ॥}\ A}}\skt{ ।}\pend

\pstart\versepar \ledsidenote{\passref{6.47.13}}\edtext{\skt{श्लैष्मि\-का\-न\-ल्प\-पि\-त्तांस्तु}}{\Afootnote{K; \skt{श्लैष्मि\-का\-न\-ल्प\-पि\-त्तांश्च}\ A; \skt{श्लेष्मि\-का\-न्य\-ल्प\-पि\-त्तांस्तु}\ H}}\skt{ स्निग्धां मात्रोपसेविनः । }\pend
\pstart\versepar \skt{पानं न }\edtext{\edtext{\skt{वाधतेत्यर्थं}}{\lemma{\skt{वाधतेत्य्}}\Afootnote{K; \skt{बाधतेऽ}\ A; \skt{बाधने}\ H}}}{\lemma{\skt{अर्थं}}\Afootnote{K; \skt{त्यर्थं}\ A H}}\skt{ विपरीतांस्तु बाधते ॥}\pend

\pstart\versepar \ledsidenote{\passref{6.47.13-1}}\edtext{\skt{बुद्धि\-स्मृ\-ति\-प्री\-ति\-करः}}{\Afootnote{K; \skt{बुद्धि\-स्मृ\-तिः\-प्री\-ति\-करः}\ H}}\skt{ }\edtext{\skt{सुखम्च}}{\Afootnote{K; \skt{सुखश्च}\ H}}\skt{ पानान्न निद्रारति }\edtext{\skt{च}}{\Afootnote{K; \skt{व}\ H}}\skt{ }\edtext{\skt{वनम्च}}{\Afootnote{K; \skt{र्द्धनश्च ।}\ H}}\skt{ । }\pend
\pstart\versepar \skt{}\edtext{\skt{संपा\-ठ\-गी\-त\-स्व\-र\-व\-व\-नम्च}}{\lemma{\skt{संपा\-ठ\-गी\-त\-स्व\-र\-व\-व\-नम्}}\Afootnote{K; \skt{सम्पा य गीत\-स्व\-र\-व\-र्द्ध\-नश्}\ H}}\skt{ }\edtext{\skt{प्रोक्तौतिरम्यः}}{\Afootnote{K; \skt{प्रोक्तोतिरम्यः}\ H}}\skt{ प्रथमो मदो हि ॥}\pend

\pstart\versepar \ledsidenote{\passref{6.47.13-2}}\edtext{\skt{अव्य\-क्त\-वु\-द्धि\-स्मृ\-ति\-वा\-ग्वि\-चेष्टः}}{\Afootnote{\skt{अव्य\-क्त\-बु\-द्धि\-स्मृ\-ति\-वा\-ग्वि\-चेष्टः}\ K H}}\skt{ }\edtext{\skt{सोन्म\-त्त\-ली\-ला\-कृति}}{\Afootnote{K; \skt{सोन्म\-त्त\-ली\-ला\-कृ\-तिर्}\ H}}\skt{ }\edtext{\skt{च}}{\Afootnote{K; \skt{अ}\ H}}\skt{ प्रशान्तः । }\pend
\pstart\versepar \skt{आल\-स्य\-नि\-द्रा\-भि\-हतो मुहुश्च मध्येन मत्तः पुरुषो मदेन ॥}\pend

\pstart\versepar \ledsidenote{\passref{6.47.13-3}}\edtext{\skt{गच्छेदगम्यान्न}}{\Afootnote{K; \skt{गच्छेदगम्यान्न}\ H}}\skt{ }\edtext{\skt{गुरुश्च}}{\Afootnote{K; \skt{गुरूंश्च}\ H}}\skt{ }\edtext{\skt{मन्ये}}{\Afootnote{K; \skt{मने}\ H}}\skt{ }\edtext{\skt{त्वादेदभक्ष्याणि च}}{\Afootnote{\skt{खादेदभक्ष्याणि च}\ K; \skt{य्त्रवा\-दे\-द\-भ\-क्ष्याणि च}\ H}}\skt{ नष्टसंज्ञः । }\pend
\pstart\versepar \skt{ब्रूयाच्च गुह्यानि हृदि स्थितानि मदे }\edtext{\skt{तृतीये}}{\Afootnote{K\textsuperscript{pc} H; \skt{तृतीयो}\ K\textsuperscript{ac}}}\skt{ पुरुषो स्वतन्त्रः ।}\pend

\pstart\versepar \ledsidenote{\passref{6.47.13-4}}\skt{चतुर्थे तु }\edtext{\skt{मदै}}{\Afootnote{K; \skt{मदे}\ H}}\skt{ }\edtext{\skt{मूढो}}{\Afootnote{K; \skt{मूढा}\ H}}\skt{ }\edtext{\skt{भग्नदार्विव}}{\Afootnote{K; \skt{भग्नदार्विवे}\ H}}\skt{ }\edtext{\skt{निःक्रियाः}}{\Afootnote{K; \skt{निष्क्रिययः ।}\ H}}\skt{ । }\pend
\pstart\versepar \skt{कार्याकार्या }\edtext{\skt{विभा\-ग\-ज्ञा\-मृ\-ता\-द\-ध्य\-परो}}{\Afootnote{K; \skt{वभागज्ञो मृतादप्यपरो}\ H}}\skt{ }\edtext{\skt{मतः}}{\Afootnote{K; \skt{मृतः ॥}\ H}}\skt{ ।}\pend

\pstart\versepar \ledsidenote{\passref{6.47.13-5}}\skt{कोमृतं तादृशं गच्छे\-दु\-न्मा\-द\-मिव चापरं । }\pend
\pstart\versepar \skt{}\edtext{\skt{बहु\-दो\-ष\-मि\-वा\-मूढः}}{\lemma{\skt{इवामूढः}}\Afootnote{K; \skt{इव अमुढः}\ H}}\skt{ कान्तारेष्ववशं कृती ।}\pend

\pstart\versepar \ledsidenote{\passref{6.47.14}}\edtext{\skt{निर्भ\-क्त\-मे\-का\-न्तत}}{\Afootnote{K; \skt{निर्भ\-क्त\-मे\-का\-न्तत}\ A; \skt{निर्भ्भ\-क्त\-मे\-का\-न्तत}\ H}}\skt{ एव मद्यं निषेव्यमाणं }\edtext{\skt{मनुजेन}}{\Afootnote{A K H\textsuperscript{pc}; \skt{मनुजे}\ H\textsuperscript{ac}}}\skt{ नित्यं । }\pend
\pstart\versepar \skt{}\edtext{\skt{उत्पा\-द\-ये\-त्क\-ष्ट\-तरां विका\-रा\-ना\-पा\-द\-ये\-च्चापि}}{\Afootnote{K; \skt{उत्पा\-द\-ये\-त्क\-ष्ट\-त\-मा\-न्वि\-का\-रा\-ना\-पा\-द\-ये\-च्चापि}\ A; \skt{आयादये क्षषु\-त\-मा\-न्वि\-का\-ना\-पा\-द\-ये\-च्चापि}\ H}}\skt{ }\edtext{\skt{शरीरमेदं}}{\Afootnote{K; \skt{शरीरभेदम्}\ A H}}\skt{ ।}\pend

\pstart\versepar \ledsidenote{\passref{6.47.15}}\skt{क्रुद्धेन भीतेन }\edtext{\skt{पिपासितेन}}{\Afootnote{A K; \skt{पिपासिते}\ H}}\skt{ शोकाभितप्तेन }\edtext{\skt{वुभुक्षितेन}}{\Afootnote{K; \skt{बुभुक्षितेन}\ A H}}\skt{ । }\pend
\pstart\versepar \skt{}\edtext{\skt{व्याया\-म\-भा\-रा\-ध्व\-प\-रि\-क्ष\-तेन}}{\Afootnote{A K; \skt{व्याया\-म\-भा\-री\-ध्व\-प\-रि\-र्क्ष\-तेन}\ H\textsuperscript{ac}; \skt{व्याया\-म\-भा\-रा\-ध्व\-प\-रि\-र्क्ष\-तेन}\ H\textsuperscript{pc}}}\skt{ वेगा\-व\-रो\-धा\-भि\-ह\-तेन चापि ।}\pend

\pstart\versepar \ledsidenote{\passref{6.47.16}}\edtext{\skt{अत्य\-म्ल\-भ\-क्ष्या\-व\-त\-तो\-द\-रेण}}{\Afootnote{A K; \skt{अत्य\-म्बु\-र\-क्षा\-व\-तो\-द\-रेण}\ H}}\skt{ }\edtext{\skt{साजीर्ण्णभक्तेन}}{\Afootnote{K H; \skt{साजीर्णभुक्तेन}\ A}}\skt{ }\edtext{\skt{तथावलेन}}{\Afootnote{K; \skt{तथाबलेन ।}\ A H}}\skt{ । }\pend
\pstart\versepar \skt{उष्णाभितप्तेन च सेव्यमानं करोति मद्यं }\edtext{\skt{विविधां}}{\Afootnote{K; \skt{विविधा}\ A H}}\skt{ }\edtext{\skt{विकारां}}{\Afootnote{K; \skt{न्विकारान् ॥}\ A H\textsuperscript{pc}; \skt{विकारान् ॥}\ H\textsuperscript{ac}}}\skt{ ।}\pend

\pstart\versepar \ledsidenote{\passref{6.47.17}}\skt{पानात्ययं }\edtext{\skt{सुपरमं}}{\Afootnote{K H; \skt{परमदं}\ A}}\skt{ }\edtext{\skt{पाना\-जी\-र्ण्ण\-म\-थापि च}}{\Afootnote{K; \skt{पाना\-जी\-र्ण\-म\-थापि वा ।}\ A; \skt{पाना\-जी\-र्ण्ण\-म\-थापि च ।}\ H}}\skt{ । }\pend
\pstart\versepar \skt{पान\-वि\-भ्र\-म\-मु\-ग्रञ्च तेषां वक्ष्यामि लक्षणं ॥}\pend

\pstart\versepar \ledsidenote{\passref{6.47.18}}\edtext{\skt{स्तम्भा\-ङ्ग\-म\-र्द\-हृ\-द\-य\-ग्र\-ह\-तो\-द\-कम्पा}}{\Afootnote{K; \skt{स्तम्भा\-ङ्ग\-म\-र्धृ\-द\-य\-ग्र\-ह\-तो\-द\-कम्पाः}\ A; \skt{स्तब्धा\-ङ्ग\-म\-र्द्द\-हृ\-द\-य\-ग्र\-ह\-तो\-द\-कम्पा}\ H}}\skt{ }\edtext{\skt{पाना\-त्य\-ये\-नि\-ल\-कृते}}{\Afootnote{K; \skt{पाना\-त्य\-येऽ\-नि\-ल\-कृते}\ A H}}\skt{ शिरसो रुजश्च । }\pend
\pstart\versepar \skt{}\edtext{\skt{स्वेद\-प्र\-ला\-प\-मु\-ख\-शो\-ष\-ण\-दा\-ह\-मूर्च्छा}}{\Afootnote{K H; \skt{स्वेद\-प्र\-ला\-प\-मु\-ख\-शो\-ष\-ण\-दा\-ह\-मूर्च्छाः}\ A}}\skt{ }\edtext{\skt{पित्तात्यये}}{\Afootnote{K H; \skt{पित्तात्मके}\ A}}\skt{ }\edtext{\skt{वद\-न\-लो\-हि\-त\-पी\-तता}}{\Afootnote{K H; \skt{वद\-न\-लो\-च\-न\-पी\-तता}\ A}}\skt{ च ।}\pend

\pstart\versepar \ledsidenote{\passref{6.47.19}}\edtext{\skt{श्लेष्मात्यये}}{\Afootnote{K H; \skt{श्लेष्मात्मके}\ A}}\skt{ }\edtext{\skt{वम\-थु\-शी\-त\-क\-फ\-प्र\-सक्ताः}}{\Afootnote{K; \skt{वम\-थु\-शी\-त\-क\-फ\-प्र\-सेकाः}\ A H}}\skt{ । सर्वात्मके भवति सर्व\-वि\-का\-र\-स\-म्पत् । }\pend
\pstart\versepar \skt{ततन्द्रां }\edtext{\skt{शरीरगुरुतां}}{\Afootnote{K H; \skt{णमङ्गगुरुतां}\ A}}\skt{ विरसाननत्वं }\edtext{\edtext{\skt{श्लेष्मा\-धि\-क\-त्व\-म\-रु\-चि\-न्म\-ल\-मू\-त्र\-संगां}}{\lemma{\skt{श्लेष्माधिकत्वम्}}\Afootnote{K; \skt{श्लेष्माधिकत्व}\ A H}}}{\lemma{\skt{अरु\-चि\-न्म\-ल\-मू\-त्र\-संगां}}\Afootnote{K; \skt{मरुचिं मल\-मू\-त्र\-स\-ङ्गम् ॥}\ A; \skt{मरु\-चि\-र्म्म\-ल\-मू\-त्र\-सं\-गान् ॥}\ H}}\skt{ ।}\edtext{}{\lemma{\skt{सर्व\-वि\-का\-र\-स\-म्पत् । ततन्द्रां}}\Afootnote{K; \skt{सर्व\-वि\-का\-र\-सं\-पत् । ऊष्मा}\ A; \skt{कार विकारसंपत् । तन्द्रा}\ H}}\pend

\pstart\versepar \ledsidenote{\passref{6.47.20}}\skt{लिंगं }\edtext{\skt{मदस्य}}{\Afootnote{K H; \skt{परस्य}\ A}}\skt{ }\edtext{\skt{परमस्य}}{\Afootnote{K H; \skt{तु मदस्य}\ A}}\skt{ वदन्ति }\edtext{\skt{तज्ज्ञास्तृष्णां}}{\Afootnote{A K; \skt{तज्ञाःस्तृष्णा}\ H}}\skt{ रुजां शिरसि सन्धिषु चापि भेदं । }\pend
\pstart\versepar \skt{}\edtext{\skt{आध्मा\-न\-मु\-द्गि\-र\-ण\-म\-म्व्ल\-रसो}}{\Afootnote{K; \skt{आध्मा\-न\-मु\-द्गि\-र\-ण\-म\-म्ल\-रसो}\ A; \skt{आध्मा\-न\-मु\-द्गि\-र\-ण\-म\-म्ल\-रसो}\ H}}\skt{ }\edtext{\skt{विदाहः}}{\Afootnote{A K; \skt{विदाहेः}\ H}}\skt{ । }\edtext{\skt{पाने}}{\Afootnote{K H; \emph{om.}\ A}}\skt{ }\edtext{\edtext{\skt{त्वजी\-र्ण्ण\-मु\-प\-ग\-च्छति}}{\lemma{\skt{त्वजीर्ण्णम्}}\Afootnote{K H; \skt{ऽजीर्ण}\ A}}}{\lemma{\skt{उपगच्छति}}\Afootnote{K H; \skt{स्य पानजनि}\ A}}\skt{ }\edtext{\skt{लक्षणानि}}{\Afootnote{K H; \skt{तस्य वदन्ति लिङ्गम् ॥}\ A}}\skt{ ।}\pend

\pstart \ledsidenote{\passref{6.47.21}}\edtext{\skt{हृद्गा\-त्र\-तो\-द\-व\-म\-थु\-र्ज्व\-र\-क\-ण्ठ\-धू\-म\-मू\-र्च्छा\-क\-फ\-स्र\-व\-ण\-मू\-ध्व\-रुजो}}{\Afootnote{K; \skt{ज्ञेयानि तत्र भिषजा सुविनिश्चितानि पित्त\-प्र\-को\-प\-ज\-नि\-तानि च कारणानि । हृद्गा\-त्र\-तो\-द\-व\-म\-थु\-ज्व\-र\-क\-ण्ठ\-धू\-म\-मू\-र्च्छा\-क\-फ\-स्र\-व\-ण\-मू\-र्ध\-रुजो}\ A; \skt{हृद्गा\-त्र\-तो\-द\-व\-म\-थु\-ज्व\-र\-क\-ण्ठ\-धू\-म\-मू\-र्च्छा\-क\-फ\-श्र\-व\-ण\-मू\-र्द्ध्व\-रुजा}\ H}}\skt{ }\edtext{\skt{विदाहाः}}{\Afootnote{K H; \skt{विदाहः ॥}\ A}}\skt{ ।}\pend

\pstart \ledsidenote{\passref{6.47.22}}\skt{द्वेषः सुरान्न विकृतेषु च तेषु तेषु तं }\edtext{\skt{पान\-वि\-भ्र\-म\-मु\-श\-न्त्य\-खि\-लेन}}{\Afootnote{K; \skt{पान\-वि\-भ्र\-म\-मु\-श\-न्त्य\-खि\-लेन}\ A; \skt{पान\-वि\-भ्र\-म\-मु\-श\-न्त्य\-खि\-लेन}\ H}}\skt{ }\edtext{\skt{धीराः}}{\Afootnote{K H; \skt{धीराः । हीनो\-त्त\-रौ\-ष्ठ\-म\-ति\-शी\-त\-म\-म\-न्द\-दाहं तैल\-प्र\-भा\-स्य\-म\-ति\-पा\-न\-हतं (तै\-ल\-प्र\-भा\-स्य\-म\-पि\-पा\-न\-हतं) विजह्यात् ॥}\ A}}\skt{ ।}\pend

\pstart \ledsidenote{\passref{6.47.23}}\edtext{\skt{हिक्काज्वरौ}}{\Afootnote{K; \skt{जिह्वौ\-ष्ठ\-द\-न्त\-म\-सितं त्वथवापि नीलं पीते च यस्य नयने रुधिरप्रमे च । हिक्काज्वरौ}\ A; \skt{हिक्काज्वरो}\ H}}\skt{ वमथु वेपथु पार्श्वशूलाः }\edtext{\edtext{\skt{कासभ्रमावपि}}{\lemma{\skt{कासभ्रमाव्}}\Afootnote{K; \skt{कासभ्रमा}\ A H}}\skt{ च}}{\lemma{\skt{अपि च}}\Afootnote{K; \skt{वपि च}\ A H}}\skt{ }\edtext{\skt{पानहतां}}{\Afootnote{K; \skt{पानहतं}\ A; \skt{पानहता}\ H}}\skt{ }\edtext{\skt{भजन्ते}}{\Afootnote{A K; \skt{न्भजन्ते ।}\ H}}\skt{ ।}\pend

\pstart\versepar \ledsidenote{\passref{6.47.24}}\edtext{\skt{तेषा\-न्नि\-वा\-र\-ण\-विधिं हि}}{\lemma{\skt{निवारणविधिं हि}}\Afootnote{K H; \skt{निवा\-र\-ण\-मि\-दंहि}\ A}}\skt{ मयोच्यमानं }\edtext{\edtext{\skt{व्यक्ता\-भि\-धा\-ण\-म\-खि\-लेन}}{\lemma{\skt{व्यक्ताभिधाणम्}}\Afootnote{K; \skt{व्यक्ताभिधान}\ A H}}}{\lemma{\skt{अखिलेन}}\Afootnote{K; \skt{मखिलेन}\ A H}}\skt{ विधिं }\edtext{\skt{शृणुश्व}}{\Afootnote{K; \skt{निबोध}\ A H}}\skt{ । }\pend
\pstart\versepar \skt{मद्यन्तु }\edtext{\skt{चुक्र\-म\-रि\-चा\-र्ज\-क\-दी\-प्य\-काढ्यां}}{\Afootnote{K; \skt{चुक्र\-म\-रि\-चा\-र्द्र\-क\-दी\-प्य\-कुष्ठ}\ A; \skt{रचु\-क्र\-म\-रि\-चा\-र्ज\-क\-दी\-प्य\-काद्यं}\ H}}\skt{ सौव\-र्च\-ला\-यु\-त\-मलं पवनस्य }\edtext{\skt{शान्त्यै}}{\Afootnote{A K; \skt{शान्त्यैः ।}\ H}}\skt{ ।}\pend

\pstart\versepar \ledsidenote{\passref{6.47.25}}\skt{पृथ्वी\-क\-दी\-प्य\-क\-म\-हौ\-ष\-ध\-हिं\-गु\-भिर्वा सौवर्चलेन च युतं }\edtext{\skt{वितरेत्सुखाय}}{\Afootnote{A K; \skt{वितरे सुखाय ॥}\ H}}\skt{ । }\pend
\pstart\versepar \skt{}\edtext{\skt{आम्रा\-त\-का\-म्र\-फ\-ल\-दा\-डि\-से\-वेत वा फल\-र\-सो\-प\-हि\-ताम्च}}{\Afootnote{\skt{आम्रा\-त\-का\-म्र\-फ\-ल\-दा\-डि\-म\-मा\-तु\-लुङ्गैः कुर्याच्छुभान्यपि}\ A; \skt{आम्रा\-त\-का\-म्र\-फ\-ल\-दा\-डि\-मा\-तु\-लुङ्गैः । कुर्याच्छुभाण्यपि च}\ K\textsuperscript{ac}; \skt{आम्रा\-त\-का\-म्र\-फ\-ल\-दा\-डि\-म\-मा\-तु\-लुङ्गैः । कुर्याच्छुभाण्यपि च}\ K\textsuperscript{pc}; \skt{आम्रा\-त\-का\-म्र\-फ\-ल\-दा\-डि\-म\-मा\-तु\-लुङ्गैः कुर्यात्सुभा}\ H}}\skt{ }\edtext{\skt{हृद्या\-ना\-नू\-प\-व\-र्ग\-पि\-शि\-तानि च}}{\Afootnote{\skt{च षड\-व\-पा\-न\-कानि ॥}\ A; \skt{षाड\-व\-पा\-न\-कानि ।}\ K; \skt{न्यपि च षड\-व\-पा\-न\-कानि ।}\ H}}\skt{ }\edtext{\skt{गन्धवन्ति}}{\Afootnote{\emph{om.}\ A K H}}\skt{ ॥}\pend

\pstart \ledsidenote{\passref{6.47.26}}\edtext{\skt{पित्तात्यये}}{\Afootnote{\skt{सेवेत वा फल\-र\-सो\-प\-हि\-ता\-न्र\-सा\-दी\-ना\-नू\-प\-व\-र्ग\-पि\-शि\-ता\-न्यपि गन्धवन्ति । पित्तात्मके}\ A; \skt{सेवेत वा फल\-र\-सो\-प\-हि\-ताम्च हृद्या\-ना\-नू\-प\-व\-र्ग\-पि\-शि\-तानि च गन्धवन्ति ॥ पित्तात्यये}\ K; \skt{सेवेत वा फल\-र\-सो\-प\-ति\-ता\-न्सु\-हृ\-द्या\-ना\-नू\-प\-व\-र्ग्ग\-पि\-शि\-तानि च गन्धव न्ति ॥ पित्तात्यये}\ H}}\skt{ मधु\-र\-व\-र्ग\-क\-षा\-य\-मिश्रं मद्यं }\edtext{\skt{हितां}}{\Afootnote{K; \skt{हितं}\ A H}}\skt{ सम\-धु\-श\-र्क\-र\-मि\-ष्ट\-गन्धं ।}\pend

\pstart\versepar \ledsidenote{\passref{6.47.27}}\edtext{\skt{पीत्वा च}}{\Afootnote{A K\textsuperscript{pc} H; \skt{पीत्वा}\ K\textsuperscript{ac}}}\skt{ }\edtext{\edtext{\skt{मद्यमपि}}{\lemma{\skt{मद्यम्}}\Afootnote{A K; \skt{मद्य}\ H}}}{\lemma{\skt{अपि}}\Afootnote{A K; \skt{मपि}\ H}}\skt{ }\edtext{\skt{चेक्षु\-र\-स\-प्र\-गाढं}}{\Afootnote{A K; \skt{चेक्षुरसं प्रगाढं}\ H}}\skt{ । }\edtext{\skt{निःशेषतः}}{\Afootnote{A K; \skt{निःशेषत}\ H}}\skt{ }\edtext{\skt{क्षण\-म\-व\-स्थि\-त\-मु\-ल्लि\-खेच्च}}{\Afootnote{A K; \skt{त्क्षण\-म\-व\-स्थि\-त\-मु\-ल्लि\-खेच्च ॥}\ H}}\skt{ । }\pend
\pstart\versepar \skt{}\edtext{\skt{लावै\-ण\-ति\-त्ति\-र\-र\-सांश्च}}{\Afootnote{K\textsuperscript{pc}; \skt{लावै\-ण\-ति\-त्ति\-रि\-र\-सांश्च}\ A; \skt{लावै\-ण\-ति\-त्ति\-र\-सांश्च}\ K\textsuperscript{ac}; \skt{लावेन तित्तिरिरसांश्च}\ H}}\skt{ }\edtext{\skt{पिबेदनम्व्लां}}{\Afootnote{\skt{पिबेदनम्ला}\ A; \skt{पिवेदनम्व्लां}\ K; \skt{विबेदनम्ला}\ H\textsuperscript{ac}; \skt{विपिबेदनम्ला}\ H\textsuperscript{pc}}}\skt{ }\edtext{\skt{मौद्गं}}{\Afootnote{K; \skt{न्मौद्गा}\ A H}}\skt{ }\edtext{\edtext{\skt{सश\-र्क\-र\-घृ\-तांश्च}}{\lemma{\skt{सश\-र्क\-र\-घृ\-तांश्}}\Afootnote{K; \skt{न्सुखाय सघृता}\ A; \skt{न्सशर्क्कर घृतांश्}\ H}}}{\lemma{\skt{च}}\Afootnote{K H; \skt{न्स}\ A}}\skt{ }\edtext{\skt{हिताय}}{\Afootnote{K H; \skt{सितांश्च}\ A}}\skt{ यूषान् ॥}\pend

\pstart\versepar \ledsidenote{\passref{6.47.28}}\skt{पानात्यये कफकृते कफमुल्लिखेच्च । }\edtext{\skt{मद्येन}}{\Afootnote{A K; \skt{सद्येन}\ H}}\skt{ }\edtext{\skt{विम्वि\-वि\-दु\-लो\-द\-क\-सं\-यु\-तेन}}{\Afootnote{K; \skt{बिम्बि\-वि\-दु\-लो\-द\-क\-सं\-यु\-तेन ।}\ A; \skt{बिम्बि\-वि\-द\-ला\-द\-क\-सं\-यु\-तेन ॥}\ H}}\skt{ । }\pend
\pstart\versepar \skt{सेवेत तिक्तकटुकांश्च }\edtext{\skt{रसा\-नु\-दा\-रा\-न्यो\-गांश्च}}{\Afootnote{K H; \skt{रसा\-नु\-दा\-रा\-न्यू\-षांश्च}\ A}}\skt{ }\edtext{\skt{तीक्त\-क\-टु\-को\-प\-हितां}}{\Afootnote{K; \skt{तिक्त\-क\-टु\-को\-प\-हिता}\ A; \skt{तीक्ष्ण\-क\-टु\-को\-प\-हिता}\ H}}\skt{ }\edtext{\skt{हिताय}}{\Afootnote{K; \skt{न्हिताय}\ A H}}\skt{ ।}\pend

\pstart \ledsidenote{\passref{6.47.29}}\edtext{\skt{पथ्यां}}{\Afootnote{K; \skt{पथ्यं}\ A H}}\skt{ यवां न विकृतानि च जांगलानि । श्लेष्म\-घ्न\-म\-न्य\-दपि यच्च निरत्ययं }\edtext{\skt{स्यात्}}{\Afootnote{K H; \skt{स्यात् । कुर्याच्च सर्वमथ सर्वभवे विधानं द्द्वन्द्वोद्भवे द्वयमवेक्ष्य यथाप्रधानम् ॥}\ A}}\skt{ ।}\pend

\pstart \ledsidenote{\passref{6.47.30}}\edtext{\skt{त्वङ्गा\-ग\-पु\-ष्प\-म\-ग\-धै\-ल\-म\-धू\-क\-धान्यैः}}{\Afootnote{K; \skt{सामा\-न्य\-म\-न्य\-दपि यच्च समग्रमग्र्यं वक्ष्यामि यच्च मनसो मदकृत्सुखं च । त्वङ्गा\-ग\-पु\-ष्प\-म\-ग\-धै\-ल\-म\-धू\-क\-धान्यैः}\ A; \skt{त्वग\-पु\-ष्प\-म\-ग\-धै\-ल\-म\-धू\-क\-धान्यैः}\ H}}\skt{ }\edtext{\skt{श्लक्ष्णै\-र\-जो\-जि\-म\-रि\-चैश्च}}{\Afootnote{K; \skt{श्लक्ष्णै\-र\-जा\-जि\-म\-रि\-चैश्च}\ A H}}\skt{ }\edtext{\skt{सुराह्वयांशैः}}{\Afootnote{K; \skt{कृतं समांशैः ॥}\ A; \skt{सुराद्क यांशैः ।}\ H}}\skt{ ।}\pend

\pstart\versepar \ledsidenote{\passref{6.47.31}}\edtext{\skt{पेयं}}{\Afootnote{K H; \skt{पानं}\ A}}\skt{ }\edtext{\skt{कपि\-त्थ\-र\-स\-वा\-रि\-प\-रू\-काढ्यं}}{\Afootnote{K; \skt{कपि\-त्थ\-र\-स\-वा\-रि\-प\-रू\-ष\-काढ्यं}\ A; \skt{कपि\-त्थ\-र\-स\-वा\-रि\-प\-रू\-ष\-काद्यं}\ H}}\skt{ }\edtext{\skt{पानात्ययेषु}}{\Afootnote{A K; \skt{वारात्ययेषु}\ H}}\skt{ }\edtext{\skt{विधि\-व\-त्स्रु\-त\-म\-म्व\-रान्ते}}{\Afootnote{K; \skt{विधि\-व\-त्स्रु\-त\-म\-म्ब\-रान्ते}\ A H}}\skt{ । }\pend
\pstart\versepar \skt{}\edtext{\skt{ह्रीवे\-र\-प\-द्म\-प\-रि\-पे\-ल\-व\-मु\-स्त\-काढ्यैः}}{\Afootnote{K; \skt{ह्रीवे\-र\-प\-द्म\-प\-रि\-पे\-ल\-व\-सं\-प्र\-युक्तैः}\ A; \skt{ह्रीवे\-र\-प\-द्म\-प\-रि\-पे\-ल\-व\-मु\-स्त\-काद्यैः}\ H}}\skt{ पुष्पै\-र्वि\-लि\-प्य\-क\-र\-वी\-र\-ज\-लो\-द्भ\-वैश्च ।}\pend

\pstart\versepar \ledsidenote{\passref{6.47.32}}\skt{पिष्टैः }\edtext{\edtext{\skt{सप\-द्म\-क\-यु\-तै\-रपि}}{\lemma{\skt{सप\-द्म\-क\-यु\-तैर्}}\Afootnote{A K; \skt{सपद्मकयुतै}\ H}}}{\lemma{\skt{अपि}}\Afootnote{A K; \skt{रपि}\ H}}\skt{ सारिवाद्यैः । सेकं जलैश्च }\edtext{\skt{वितरेदमलैः}}{\Afootnote{A K; \skt{वितरे…मलैः}\ H}}\skt{ सुशीतैः । }\pend
\pstart\versepar \skt{}\edtext{\skt{त्वक्प\-त्र\-चो\-च\-म\-रि\-चै\-ल\-भु\-ज\-ङ्ग\-पु\-ष्प\-श्ले\-ष्मा\-त\-क\-प्र\-स\-व\-क\-ल्क\-गु\-डै\-रु\-पेतं}}{\lemma{\skt{त्वक्प\-त्र\-चो\-च\-म\-रि\-चै\-ल\-भु\-ज\-ङ्ग\-पु\-ष्प\-श्ले\-ष्मा\-त\-क\-प्र\-स\-व\-क\-ल्क\-गु\-डैर्}}\Afootnote{K; \skt{त्वक्प\-त्र\-चो\-च\-म\-रि\-चै\-ल\-भु\-ज\-ङ्ग\-पु\-ष्प\-श्ले\-ष्मा\-त\-क\-प्र\-स\-व\-व\-ल्क\-गु\-डैर्}\ A; \skt{त्वक्प\-त्र\-चो\-च\-प\-रि\-चै\-ल\-भु\-ज\-ङ्ग\-पुष्प श्लेष्मा\-न्त\-क\-प्र\-स\-व\-क\-ल्क\-गु\-डैर्}\ H}}\skt{ ।}\pend

\pstart\versepar \ledsidenote{\passref{6.47.33}}\skt{द्राक्षायुतं हृतमलं }\edtext{\skt{मदिरामयार्त्तै}}{\Afootnote{K; \skt{मदि\-रा\-म\-या\-र्तैस्}\ A H}}\skt{ । }\edtext{\skt{स्तन्पानकं}}{\Afootnote{K; \skt{तत्पानकं}\ A H}}\skt{ शुचि सुगन्धि नरैर्निषेव्यं । }\pend
\pstart\versepar \skt{}\edtext{\skt{पिष्टं}}{\Afootnote{K H; \skt{पिष्ट्वा}\ A}}\skt{ }\edtext{\skt{पिवेच्च}}{\Afootnote{K; \skt{पिबेच्च}\ A H}}\skt{ मधुकं कटुरोहिणीं च । }\edtext{\skt{मूलञ्च}}{\Afootnote{K; \skt{द्राक्षां च}\ A; \skt{मूलं तु}\ H}}\skt{ }\edtext{\skt{तुल्य\-म\-स\-कृ\-त्त्र\-पुशी}}{\Afootnote{K; \skt{मूल\-म\-स\-कृ\-त्त्र\-पुषी}\ A; \skt{तुल्य\-म\-स\-कृ\-त्त्र\-पुषी}\ H}}\skt{ भवं यत् ।}\pend

\pstart\versepar \ledsidenote{\passref{6.47.34}}\edtext{\skt{कार्पा\-सि\-मे\-व\-म\-थ\-ना\-ग\-वलां च}}{\Afootnote{K; \skt{कार्पासिनीमथ च नागबलां च}\ A; \skt{कर्या\-सि\-मे\-व\-मथ नागबलाञ्च}\ H}}\skt{ तुल्यां पीत्वा सुखी भवति साधु }\edtext{\skt{सुवर्चलाञ्च}}{\Afootnote{A K; \skt{सुवर्च्चला च ॥}\ H}}\skt{ । }\pend
\pstart\versepar \skt{}\edtext{\skt{कार्श्म\-र्य\-दा\-रु\-वि\-ड\-पि\-प्प\-लि\-दा\-डि\-मेषु}}{\Afootnote{K; \skt{काश्म\-र्य\-दा\-रु\-वि\-ड\-दा\-डि\-म\-पि\-प्प\-लीषु}\ A; \skt{काश्म\-र्य\-दा\-रु\-वि\-ड\-पि\-प्प\-ली\-दा\-डि\-मेषु}\ H}}\skt{ }\edtext{\skt{द्राक्षान्वितेषु}}{\Afootnote{K; \skt{द्राक्षान्वितासु}\ A; \skt{द्राक्षान्वितेसु}\ H}}\skt{ }\edtext{\skt{कृतमाम्वुनि}}{\Afootnote{K; \skt{कृतमम्बुनि}\ A H}}\skt{ पानकं यत् ।}\pend

\pstart\versepar \ledsidenote{\passref{6.47.35}}\edtext{\skt{तद्वी\-ज\-पू\-र\-क\-र\-सा\-यु\-त\-माशु}}{\Afootnote{K; \skt{तद्बी\-ज\-पू\-र\-क\-र\-सा\-यु\-त\-माशु}\ A; \skt{तद्बी\-ज\-पू\-र\-क\-र\-सा\-यु\-त\-माशु}\ H}}\skt{ पीतं । शान्तिं }\edtext{\skt{परां}}{\Afootnote{A K; \skt{परा}\ H}}\skt{ }\edtext{\edtext{\skt{मद\-ग\-दे\-ष्व\-चि\-रा\-त्क\-रोति}}{\lemma{\skt{मदगदेष्व्}}\Afootnote{K; \skt{परमदे}\ A; \skt{मदगदैष्व्}\ H}}}{\lemma{\skt{अचिरात्करोति}}\Afootnote{K H; \skt{त्वचिरात्करोति ।}\ A}}\skt{ । }\pend
\pstart\versepar \skt{द्राक्षा\-सि\-ता\-म\-धु\-क\-जी\-र\-क\-धा\-न्य\-कृ\-ष्णा\-स्वेवं कृतं }\edtext{\skt{सतृवृतासु}}{\Afootnote{K H; \skt{त्रिवृतया च}\ A}}\skt{ }\edtext{\skt{पिबेदथापि}}{\Afootnote{H; \skt{पिबेत्तथैव ॥}\ A; \skt{पिवेदथापि ।}\ K}}\skt{ ।}\pend

\pstart\versepar \ledsidenote{\passref{6.47.36}}\edtext{\skt{सौव\-र्च\-ला\-यु\-त\-मु\-दा\-र\-रसां}}{\Afootnote{K; \skt{सौव\-र्च\-ला\-यु\-त\-मु\-दा\-र\-रसं}\ A; \skt{सौव\-र्च्च\-ला\-यु\-त\-मु\-दा\-र\-र\-सम्}\ H}}\skt{ }\edtext{\skt{फलाम्व्ल}}{\Afootnote{K; \skt{फलाम्लं}\ A H}}\skt{ । भार्गी }\edtext{\skt{श्रितेन च}}{\Afootnote{K; \skt{शृतेन च}\ A; \skt{स्रितेन च}\ H}}\skt{ जलेन }\edtext{\skt{हितोवसेकः}}{\Afootnote{\skt{हितोऽवसेकः ॥}\ A; \skt{हितोवसेकः ।}\ K; \skt{हितो वसेकः ॥}\ H}}\skt{ । }\pend
\pstart\versepar \skt{}\edtext{\skt{इक्ष्वा\-कु\-धा\-मा\-र्ग\-व\-वृ\-क्ष\-कानि काका\-ह्व\-यो\-दु\-म्व\-रिकां च दुग्धे}}{\Afootnote{\emph{om.}\ A K H}}\pend

\pstart \ledsidenote{\passref{6.47.37}}\edtext{\skt{विपाच्य}}{\Afootnote{\skt{इक्ष्वा\-कु\-धा\-मा\-र्ग\-व\-वृ\-क्ष\-काणि काका\-ह्व\-यो\-दु\-म्ब\-रि\-काश्च दुग्धे । विपाच्य}\ A; \skt{इक्ष्वा\-कु\-धा\-मा\-र्ग\-व\-वृ\-क्ष\-कानि काका\-ह्व\-यो\-दु\-म्व\-रिकां च दुग्धे विपाच्य}\ K; \skt{इक्ष्वा\-कु\-धा\-मा\-र्ग्ग\-व\-वृ\-क्ष\-कानि काका\-ह्व\-यो\-दु\-म्ब\-रि\-कांश्च दुग्धे । विपाच्य}\ H}}\skt{ तस्यांजलिना वमेद्धि मद्यं }\edtext{\skt{पिवेच्चाह्निगते}}{\Afootnote{K; \skt{पिबेच्चाह्नि गते}\ A; \skt{पिबेच्चापि गते}\ H}}\skt{ त्वजीर्ण्णे ।}\pend

\pstart\versepar \ledsidenote{\passref{6.47.38}}\edtext{\skt{त्वङ्ना\-ग\-पु\-ष्प\-वि\-ड\-जी\-र\-क\-हिं\-गु\-कृष्णाः}}{\Afootnote{K; \skt{त्वक्पि\-प्प\-ली\-भु\-ज\-ग\-पु\-ष्प\-वि\-डै\-रु\-पेतं}\ A; \skt{त्वग\-पु\-ष्प\-वि\-ड\-जी\-रक हिङ्गुकृष्णाः}\ H}}\skt{ । सेवेत }\edtext{\skt{चापि}}{\Afootnote{K H; \skt{हिङ्गु}\ A}}\skt{ मरिचैलयुतं }\edtext{\skt{फलाम्व्लं}}{\Afootnote{K; \skt{फलाम्लाम्}\ A H}}\skt{ । }\pend
\pstart\versepar \skt{}\edtext{\skt{उष्णाम्वु}}{\Afootnote{K; \skt{उष्णाम्बु}\ A H}}\skt{ }\edtext{\edtext{\skt{सैन्ध\-व\-यु\-ता\-स्त्व\-थवा}}{\lemma{\skt{सैन्धवयुतास्}}\Afootnote{K H; \skt{सैन्धवयुता}\ A}}}{\lemma{\skt{त्वथवा}}\Afootnote{K H; \skt{स्त्वथवा}\ A}}\skt{ }\edtext{\skt{यथोक्तं}}{\Afootnote{K H; \skt{विडत्व}\ A}}\skt{ । }\edtext{\skt{चव्यै\-ल\-हिं\-गु\-म\-ग\-धा\-फ\-ल\-मू\-ल\-शुण्ठी}}{\Afootnote{K H; \skt{क्चव्यै\-ल\-हि\-ङ्गु\-म\-ग\-धा\-फ\-ल\-मू\-ल\-शुण्ठीः ॥}\ A}}\skt{ ।}\pend

\pstart\versepar \ledsidenote{\passref{6.47.39}}\skt{हृद्यैः }\edtext{\skt{खलैरपि}}{\lemma{\skt{खलैर्}}\Afootnote{K H; \skt{खडैर्}\ A}}\skt{ च }\edtext{\edtext{\skt{भोजनमत्र}}{\lemma{\skt{भोजनम्}}\Afootnote{K; \skt{भोजन}\ A H}}}{\lemma{\skt{अत्र}}\Afootnote{K; \skt{मत्र}\ A H}}\skt{ शस्तं । }\edtext{\skt{द्राक्षा\-क\-पि\-त्थ\-फ\-ल\-दा\-डि\-म\-पा\-नकं}}{\Afootnote{A K; \skt{द्राक्षा\-क\-पि\-त्थ\-वि\-ड\-दा\-डि\-म\-पा\-नकं}\ H}}\skt{ यत् । }\pend
\pstart\versepar \skt{तत्पा\-न\-वि\-भ्र\-म\-हरं }\edtext{\skt{मधु\-श\-र्क\-रा\-ढ्य\-मा\-म्रा\-त\-को\-ल\-र\-स\-पा\-न\-क\-मेव}}{\lemma{\skt{मधु\-श\-र्क\-रा\-ढ्य\-मा\-म्रा\-त\-को\-ल\-र\-स\-पा\-न\-कम्}}\Afootnote{A K; \skt{मधु\-श\-र्क्क\-रा\-ढ्य\-मा\-म्रा\-त\-को\-ल\-र\-स\-पा\-न\-कम्}\ H}}\skt{ चापि ।}\pend

\pstart\versepar \ledsidenote{\passref{6.47.40}}\edtext{\skt{खर्जू\-र\-वे\-त्र\-क\-र\-वी\-श\-प\-रु\-ष\-केषु}}{\Afootnote{K; \skt{खर्जू\-र\-वे\-त्र\-क\-क\-री\-र\-प\-रू\-ष\-केषु}\ A; \skt{खर्ज्जू\-र\-वे\-त्र\-क\-र\-वी\-र\-प\-रू\-ष\-केषु}\ H}}\skt{ }\edtext{\skt{द्राक्षातृवृत्सु च}}{\Afootnote{K H; \skt{द्राक्षात्रिवृत्सु च}\ A}}\skt{ कृतं । }\edtext{\skt{सघृतं}}{\Afootnote{K H; \skt{ससितं}\ A}}\skt{ }\edtext{\skt{हितं वा}}{\Afootnote{K H; \skt{हिमं वा ।}\ A}}\skt{ । }\pend
\pstart\versepar \skt{श्रीप\-र्ण्णि\-यु\-क्त\-म\-थवा तु }\edtext{\skt{पिबेदिमानि}}{\Afootnote{\skt{पिबेदिमानि}\ A; \skt{पिवेदिमानि}\ K H}}\skt{ }\edtext{\skt{यष्टी\-क\-सो\-त्प\-ल\-हि\-मांवु}}{\Afootnote{K; \skt{यष्ट्या\-ह्व\-यो\-त्प\-ल\-हि\-माम्बु}\ A; \skt{यष्टी\-क\-सो\-त्प\-ल\-हि\-माम्बु}\ H}}\skt{ विमिश्रितानि ।}\pend

\pstart\versepar \ledsidenote{\passref{6.47.41}}\edtext{\skt{क्षीरी\-प्र\-वा\-ल\-वि\-स\-जी\-र\-क\-ना\-ग\-पु\-ष्प\-प\-त्रे\-ल\-वा\-लु\-सि\-त\-सा\-रि\-व\-प\-द्म\-कानि}}{\Afootnote{K; \skt{क्षीरि\-प्र\-वा\-ल\-बि\-स\-जी\-र\-क\-ना\-ग\-पु\-ष्प\-प\-त्रै\-ल\-वा\-लु\-सि\-त\-सा\-रि\-व\-प\-द्म\-कानि ।}\ A; \skt{क्षीरी\-प्र\-वा\-ल\-वि\-स\-जी\-र\-क\-ना\-ग\-पु\-ष्प\-प\-त्रै\-ल\-वा\-लु\-सि\-त\-शा\-रि\-व\-प\-द्म\-कानि ।}\ H}}\skt{ । }\pend
\pstart\versepar \skt{}\edtext{\skt{आम्रा\-त\-भ\-व्य\-क\-र\-म\-र्द\-क\-पि\-त्थ\-को\-ल\-वृ\-क्षा\-म्व्ल\-वे\-त्र\-सि\-त\-बी\-ज\-क\-दा\-डि\-मानि}}{\Afootnote{K; \skt{आम्रा\-त\-भ\-व्य\-क\-र\-म\-र्द\-क\-पि\-त्थ\-को\-ल\-वृ\-क्षा\-म्ल\-वे\-त्र\-फ\-ल\-जी\-र\-क\-दा\-डि\-मानि ॥}\ A; \skt{आम्रा\-त\-भ\-व्य\-क\-र\-म\-र्द्द\-क\-पि\-त्थ\-को\-ल\-वृ\-क्षा\-म्ल\-वे\-त्र\-सि\-त\-बी\-ज\-क\-दा\-डि\-मानि ॥}\ H}}\skt{ ।}\pend

\pstart\versepar \ledsidenote{\passref{6.47.42}}\skt{सेवेत वा }\edtext{\skt{मरि\-च\-जी\-र\-क\-ना\-ग\-पु\-ष्प\-त्व\-क्प\-त्र\-शु\-ण्ठि\-च\-वि\-कै\-ल\-यु\-ता\-न्र\-सांस्तु}}{\Afootnote{K H; \skt{मरि\-च\-जी\-र\-क\-ना\-ग\-पु\-ष्प\-त्व\-क्प\-त्र\-वि\-श्व\-च\-वि\-कै\-ल\-यु\-ता\-न्र\-सांश्च ।}\ A}}\skt{ । }\pend
\pstart\versepar \skt{}\edtext{\skt{सूक्ष्मा\-म्व\-र\-श्रु\-त\-हि\-मां\-शु\-सु\-ग\-न्धि\-गन्धां}}{\Afootnote{K; \skt{सूक्ष्मा\-म्ब\-र\-स्रु\-त\-हि\-मांच सुगन्धिग}\ A; \skt{सूक्ष्मा\-म्ब\-र\-स्रु\-त\-हि\-मासु सुगन्धिग}\ H}}\skt{ }\edtext{\skt{पञ्चे\-न्द्रि\-या\-र्थ\-वि\-धयो मृदुपानयोगाः हृद्याः सुखाश्च मनसः सततन्निषेव्या}}{\Afootnote{\skt{न्धान्पा\-नो\-द्भ\-वा\-न्नु\-दति सप्त\-ग\-दा\-न\-शे\-षान् ॥}\ A; \skt{नोद्भवां नुदति सप्त\-ग\-दा\-न\-शे\-षान् ।}\ K\textsuperscript{ac}; \skt{मानोद्भवां नुदति सप्त\-ग\-दा\-न\-शे\-षान् ।}\ K\textsuperscript{pc}; \skt{त्वां मानोद्भवां नुदति सप्त\-ग\-दा\-न\-शे\-षान् ॥}\ H}}\skt{ ॥}\pend

\pstart \ledsidenote{\passref{6.47.43}}\edtext{\skt{पानात्ययेषु}}{\Afootnote{\skt{पञ्चे\-न्द्रि\-या\-र्थ\-वि\-षया मृदुपानयोगा हृद्याः सुखाश्च मनसः सततं निषेव्याः । पानात्ययेषु}\ A; \skt{पञ्चे\-न्द्रि\-या\-र्थ\-वि\-धयो मृदुपानयोगाः हृद्याः सुखाश्च मनसः सततन्निषेव्या ॥ पानात्ययेरषु}\ K\textsuperscript{ac}; \skt{पञ्चे\-न्द्रि\-या\-र्थ\-वि\-धयो मृदुपानयोगाः हृद्याः सुखाश्च मनसः सततन्निषेव्या ॥ पानात्ययेषु}\ K\textsuperscript{pc}; \skt{पञ्चे\-न्द्रि\-या\-र्थे\-वि\-धयो मृदुशेषयोगाः हृद्याः सुखाश्च मनसः सततं निषेव्याः । पानात्ययेषु … ॥}\ H\textsuperscript{ac}; \skt{पञ्चे\-न्द्रि\-या\-र्थ\-वि\-धयो मृदुशेषयोगा हृद्याः सुखाश्च मनसः सततं निषेव्याः । पानात्ययेषु … ॥}\ H\textsuperscript{pc}}}\skt{ }\edtext{\skt{नव\-यौ\-व\-न\-पी\-न\-गात्र्यः}}{\Afootnote{K; \skt{विक\-टो\-रु\-नि\-त\-म्ब\-वत्यः}\ A; \emph{om.}\ H}}\skt{ }\edtext{\edtext{\skt{सेव्याश्च}}{\lemma{\skt{सेव्याश्}}\Afootnote{K; \skt{पीनोन्}\ A; \emph{om.}\ H}}}{\lemma{\skt{च}}\Afootnote{K; \skt{न}\ A; \emph{om.}\ H}}\skt{ }\edtext{\skt{पञ्च\-वि\-ष\-या\-ति\-श\-या\-यु\-वत्यः}}{\Afootnote{K; \skt{तस्त\-न\-भ\-रा\-न\-त\-म\-ध्य\-देशाः ॥}\ A; \emph{om.}\ H}}\skt{ ।}\pend

\pstart\versepar \ledsidenote{\passref{6.47.45}}\edtext{\skt{पिबेद्रसं}}{\Afootnote{\skt{पिबेद्रसं}\ A H; \skt{पिवेद्रसं}\ K}}\skt{ पुष्प\-फ\-लो\-द्भ\-वम्वा }\edtext{\skt{सिता\-म\-धू\-क\-त्रि\-सु\-ग\-न्ध\-युक्तं}}{\Afootnote{K; \skt{सिता\-म\-धू\-क\-त्रि\-सु\-ग\-न्धि\-यु\-क्तम् ।}\ A; \skt{सितामधूकं त्रिसुगन्धयुक्तं ।}\ H}}\skt{ । }\pend
\pstart\versepar \skt{}\edtext{\skt{सञ्चूर्ण्य}}{\Afootnote{A K; \skt{संचूर्ण्ण्या}\ H}}\skt{ संयोज्य च }\edtext{\skt{नाग\-पु\-ष्पै\-र\-जा\-जि\-कृ\-ष्णा\-म\-रि\-चैश्च}}{\Afootnote{A K; \skt{नागपुष्पैः सजा\-जि\-कृ\-ष्णा\-म\-रि\-चैश्च}\ H}}\skt{ तुल्यैः ।}\pend

\pstart\versepar \ledsidenote{\passref{6.47.46}}\edtext{\skt{वर्षा\-भू\-य\-ष्ट्या\-ह्व\-म\-धू\-क\-युक्ताः}}{\Afootnote{K; \skt{वर्षा\-भू\-य\-ष्ट्या\-ह्व\-म\-धू\-क\-लाक्षा}\ A; \skt{वर्षा\-भू\-य\-ष्ट्या\-ह्व\-म\-धू\-क\-युक्ता}\ H}}\skt{ }\edtext{\skt{शस्या\-न्यु\-दा\-रां\-कु\-र\-जी\-र\-कानि}}{\Afootnote{K; \skt{त्वक्क\-र्बु\-दा\-रा\-ङ्कु\-र\-जी\-र\-काणि ।}\ A; \skt{शम्या\-नु\-दा\-रां\-कु\-र\-जी\-र\-कानि ।}\ H}}\skt{ । }\pend
\pstart\versepar \skt{}\edtext{\skt{द्राक्षा\-म\-कृ\-ष्णा\-म\-धुकं}}{\Afootnote{K; \skt{द्राक्षां च कृष्णामथ के}\ A; \skt{द्राक्षा स कृष्णामधुकं}\ H}}\skt{ }\edtext{\skt{रसं च}}{\Afootnote{K H; \skt{शरं च}\ A}}\skt{ }\edtext{\skt{क्षीरं}}{\Afootnote{K H; \skt{क्षीरे}\ A}}\skt{ समालोड्य }\edtext{\skt{पिबेत्श्रुतन्तत्}}{\Afootnote{\skt{पिबेत्सुखेप्सुः ॥}\ A; \skt{पिवेत्श्रुतन्तत् ।}\ K; \skt{पिबे त्सुतन्तत् ॥}\ H}}\skt{ ।}\pend

\pstart\versepar \ledsidenote{\passref{6.47.47}}\skt{भवेच्च मद्येन तु येन पातितः । }\edtext{\skt{सुरा\-स\-वा\-म्वा\-म\-धु\-ना\-थ\-वा\-नरः।}}{\lemma{\skt{सुरा\-स\-वा\-म्वा\-म\-धु\-ना\-थ\-वा\-नरः}}\Afootnote{K; \skt{प्रकामपीतेन सुरासवादिना ।}\ A; \skt{सुरा\-स\-वा\-भ्या\-म्म\-धु\-ना\-थवा नरः ।}\ H}}\skt{ }\pend
\pstart\versepar \skt{}\edtext{\skt{तदेव}}{\Afootnote{A H; \skt{तदेव}\ K}}\skt{ तस्मै }\edtext{\skt{विधि\-व\-त्प्र\-धा\-प\-ये\-द्वि\-प\-र्यये}}{\Afootnote{K; \skt{विधि\-व\-त्प्र\-दा\-प\-ये\-द्वि\-प\-र्यये}\ A H}}\skt{ }\edtext{\edtext{\skt{भ्रंशमतो}}{\lemma{\skt{भ्रंशम्}}\Afootnote{K; \skt{भ्रंश}\ A; \skt{भ्रंस}\ H}}}{\lemma{\skt{अतो}}\Afootnote{K; \skt{मव}\ A; \skt{गतो}\ H}}\skt{ }\edtext{\skt{न्यथार्च्छति}}{\Afootnote{K; \skt{श्यमृच्छति ॥}\ A; \skt{न्य्थाच्छति ॥}\ H}}\pend

\pstart\versepar \ledsidenote{\passref{6.47.48}}\skt{यथा नरे\-न्द्रो\-प\-ह\-तस्य कस्य\-चि\-द्भ\-वे\-त्प्र\-सा\-द\-स्तत एव नान्यतः । }\pend
\pstart\versepar \skt{ध्रुवं तथा मद्यहतस्य देहिनो भवे\-त्प्र\-सा\-द\-स्तत एव नान्यतः ॥}\pend

\pstart\versepar \ledsidenote{\passref{6.47.49}}\skt{विच्छिन्नमद्यः }\edtext{\skt{सहसा }\edtext{\skt{यस्तु}}{\lemma{\skt{तु}}\Afootnote{K; \skt{ऽति}\ A; \skt{तस्तु}\ H}}}{\lemma{\skt{सहसा यस्}}\Afootnote{A K; \skt{सहसा य}\ H}}\skt{ }\edtext{\skt{मद्यन्निषेवते}}{\Afootnote{K H; \skt{मद्यं विषेवते ।}\ A}}\skt{ । }\pend
\pstart\versepar \skt{तस्य पानात्ययोद्दिष्टा }\edtext{\skt{विका\-रा\-सं\-भ\-वन्ति}}{\Afootnote{K; \skt{विकाराः संभवन्ति}\ A; \skt{विकाराः सम्भवन्ति}\ H}}\skt{ हि ।}\pend

\pstart\versepar \ledsidenote{\passref{6.47.50}}\skt{मद्य\-स्या\-ग्ने\-य\-वा\-यव्यौ }\edtext{\skt{गुणा\-व\-म्वु\-व\-हानि तु}}{\Afootnote{K; \skt{गुणा\-व\-म्बु\-व\-हा\-नि(\-गु\-णा\-व\-म्भो\-व\-हानि) तु ।}\ A; \skt{गुणा\-व\-म्बु\-व\-हानि तु ।}\ H}}\skt{ । }\pend
\pstart\versepar \skt{स्रोतांसि }\edtext{\skt{शोषयेतां हि}}{\Afootnote{K; \skt{शोषयेयातां}\ A; \skt{शोषयेत्तां हि}\ H}}\skt{ }\edtext{\skt{तत\-स्तृ\-ष्णो\-प\-जा\-यते}}{\Afootnote{K; \skt{तेन तृष्णोपजायते ॥}\ A; \skt{तत\-स्तृ\-ष्णो\-प\-जा\-यते ॥}\ H}}\skt{ ।}\pend

\pstart\versepar \ledsidenote{\passref{6.47.51}}\edtext{\skt{पाट\-ल्यु\-त्प\-ल\-द\-ण्डेषु}}{\Afootnote{K; \skt{पाट\-लो\-त्प\-ल\-क\-न्देषु}\ A; \skt{पाट\-ल्यु\-त्प\-ल\-द\-ण्डेन}\ H}}\skt{ }\edtext{\skt{मुद्गपर्ण्या च}}{\Afootnote{K H; \skt{मुद्गपर्ण्यां च}\ A}}\skt{ साधितं । }\pend
\pstart\versepar \skt{}\edtext{\skt{पिबे\-त्पि\-प्प\-लि\-सं\-मिश्रं}}{\Afootnote{H; \skt{पिबे\-न्मा\-ग\-धि\-को\-न्मिश्रं}\ A; \skt{पिवे\-त्पि\-प्प\-लि\-सं\-मिश्रं}\ K}}\skt{ तत्राम्भो हिमशीतलं ।}\pend

\pstart\versepar \ledsidenote{\passref{6.47.52}}\edtext{\skt{सर्पि\-स्तै\-ल\-व\-सा\-दु\-ग्ध\-द\-धि\-भृ\-ङ्ग\-र\-सै\-र्युतं}}{\Afootnote{K H; \skt{सर्पि\-स्तै\-ल\-व\-सा\-म\-ज्ज\-द\-धि\-भृ\-ङ्ग\-र\-सै\-र्यु\-तम् ।}\ A}}\skt{ । }\pend
\pstart\versepar \skt{क्वाथेन }\edtext{\skt{विल्वयवयोः}}{\Afootnote{K; \skt{बिल्वयवयोः}\ A H}}\skt{ सर्वगन्धैश्च }\edtext{\skt{पेषिमैः}}{\Afootnote{K H; \skt{पेषितैः ॥}\ A}}\skt{ ।}\pend

\pstart\versepar \ledsidenote{\passref{6.47.53}}\edtext{\skt{पक्वमभ्यञ्जने}}{\Afootnote{K; \skt{पक्वमभ्यञ्जने}\ A; \skt{पक्वमभ्यञ्जनं}\ H}}\skt{ श्रेष्ठं सेके क्वाथश्च शीतलः। }\pend
\pstart\versepar \skt{रसवन्ति च }\edtext{\skt{भोज्या\-नि\-य\-था\-स्व\-म\-व\-चा\-र\-येत्}}{\Afootnote{K; \skt{भोज्यानि यथा\-स्व\-म\-व\-चा\-र\-येत्}\ A H}}\skt{ ।}\pend

\pstart\versepar \ledsidenote{\passref{6.47.54}}\edtext{\skt{त्वचं}}{\Afootnote{K H; \skt{पानकानि सुशीतानि हृद्यानि सुरभीणि च । त्वचं}\ A}}\skt{ प्राप्तस्तु पानोष्मा }\edtext{\skt{पित्त\-र\-क्ता\-ति\-मू\-र्च्छितः}}{\Afootnote{K; \skt{पित्त\-र\-क्ता\-भि\-मू\-र्च्छितः ।}\ A H}}\skt{ । }\pend
\pstart\versepar \skt{दाहं प्रकुरुते घोरं पित्तवत्तत्त्र भेषजं ॥}\pend

\pstart\versepar \ledsidenote{\passref{6.47.55}}\skt{शीतं }\edtext{\skt{विधाणमत}}{\Afootnote{K; \skt{विधानमत}\ A H}}\skt{ }\edtext{\skt{ऊर्ध्व\-म\-नु\-क्र\-मिष्ये}}{\Afootnote{K; \skt{ऊर्ध्वमहं प्रवक्ष्ये}\ A; \skt{ऊर्द्ध्व\-म\-नु\-र\-मिष्ये}\ H}}\skt{ }\edtext{\skt{दाह\-प्र\-शां\-ति\-क\-र\-मृ\-द्धि\-व\-ता\-न्न\-न\-राणां}}{\Afootnote{K; \skt{दाह\-प्र\-शा\-न्ति\-क\-र\-मृ\-द्धि\-मतां नराणाम् ।}\ A; \skt{दाह\-प्र\-शा\-न्ति\-क\-र\-मृ\-द्धि\-व\-ता\-न्न\-राणां ।}\ H}}\skt{ । }\pend
\pstart\versepar \skt{तत्रादितो मलयजेन हितः }\edtext{\skt{प्रदे\-ह\-श्च\-न्द्रां\-शु\-हा\-र\-तु\-हि\-नो\-द\-क\-शी\-त\-लेन}}{\Afootnote{K; \skt{प्रदे\-ह\-श्च\-न्द्रां\-शु\-हा\-र\-तु\-हि\-नो\-द\-क\-शी\-त\-लेन ॥}\ A H}}\skt{ ।}\pend

\pstart\versepar \ledsidenote{\passref{6.47.56}}\edtext{\skt{शीता\-म्व\-शी\-त\-ल\-त\-रैश्च}}{\Afootnote{K; \skt{शीता\-म्बु\-शी\-त\-ल\-त\-रैश्च}\ A; \skt{शीता\-म्बु\-शी\-त\-ल\-न\-रैश्च}\ H}}\skt{ }\edtext{\skt{शयानयेनं}}{\Afootnote{K; \skt{शयानमेनं}\ A H}}\skt{ }\edtext{\skt{हारै\-र्मृ\-णा\-ल\-व\-ल\-यै\-र\-वलाः}}{\Afootnote{K; \skt{हारै\-र्मृ\-णा\-ल\-व\-ल\-यै\-र\-बलाः}\ A; \skt{हारै\-र्म्मृ\-णा\-ल\-व\-ल\-यै\-र\-बला}\ H}}\skt{ स्पृशेयुः । }\pend
\pstart\versepar \skt{}\edtext{\skt{भिन्नो\-त्प\-लो\-ज्व\-ल\-हिमे}}{\Afootnote{K; \skt{भिन्नो\-त्प\-लो\-ज्ज्व\-ल\-हिमे}\ A H}}\skt{ }\edtext{\skt{शयने}}{\Afootnote{A K; \skt{सघने}\ H}}\skt{ शयीत । पत्रेषु वा }\edtext{\skt{सजलविन्दुषु}}{\Afootnote{K; \skt{सजलबिन्दुषु}\ A H}}\skt{ पद्मिनीनां ।}\pend

\pstart\versepar \ledsidenote{\passref{6.47.57}}\edtext{\skt{सम्वी\-ज\-ये\-त्प\-व\-न\-मा\-हृ\-त\-म\-ङ्ग\-नाभिः}}{\Afootnote{K; \skt{आसा\-द\-य\-न्प\-व\-न\-मा\-हृ\-त\-म\-ङ्ग\-नाभिः}\ A; \skt{सम्बी\-ज\-ये\-त्प\-व\-न\-मा\-हृ\-त\-म\-ङ्ग\-नाभिः}\ H}}\skt{ । }\edtext{\skt{कल्हा\-र\-प\-द्म\-द\-ल\-शै\-व\-ल\-सं\-च\-येभ्यः}}{\Afootnote{K; \skt{कह्ला\-र\-प\-द्म\-द\-ल\-शै\-व\-ल\-सं\-च\-येषु ।}\ A; \skt{कह्ला\-र\-प\-द्म\-द\-ल\-शै\-व\-ल\-स\-ञ्च\-येभ्यः ।}\ H}}\skt{ । }\pend
\pstart\versepar \skt{}\edtext{\edtext{\skt{शीतै\-र्व\-ना\-न्त\-प\-व\-नै\-रु\-प\-वी\-ज्य\-मानः}}{\lemma{\skt{शीतै\-र्व\-ना\-न्त\-प\-व\-नैर्}}\Afootnote{K; \skt{कान्तै\-र्व\-ना\-न्त\-प\-वनैः}\ A; \skt{शीतै\-र्व्व\-ना\-न्त\-प\-व\-नैर्}\ H}}}{\lemma{\skt{उपवीज्यमानः}}\Afootnote{K H; \skt{परिमृश्यमानः}\ A}}\skt{ । }\edtext{\skt{प्रीत\-श्च\-रे\-द्भ\-व\-न\-का\-न\-न\-दी\-र्घि\-कासु}}{\Afootnote{K; \skt{शक्त\-श्च\-रे\-द्भ\-व\-न\-का\-न\-न\-दी\-र्घि\-कासु ॥}\ A; \skt{प्रीत\-श्च\-रे\-द्भ\-व\-न\-का\-न\-न\-दी\-र्घि\-कासु ॥}\ H}}\skt{ ।}\pend

\pstart\versepar \ledsidenote{\passref{6.47.58}}\skt{दाहा\-भि\-भू\-त\-म\-थ\-वा\-प\-रि\-षे\-च\-येत्तु }\edtext{\skt{लाम\-ज्ज\-का\-म्वु\-रु\-ह\-तो\-य\-मयैः}}{\Afootnote{K; \skt{लाम\-ज्ज\-का\-म्बु\-रु\-ह\-च\-न्द\-न\-तो\-य\-तोयैः ।}\ A; \skt{लाम\-ज्ज\-म्बु\-रु\-ह\-तोय समैः}\ H}}\skt{ }\edtext{\skt{सुशीतैः}}{\Afootnote{K H; \emph{om.}\ A}}\skt{ ॥ }\pend
\pstart\versepar \skt{विस्रावितां हृत\-म\-ला\-न्न\-व\-वा\-रि\-पूर्ण्णां }\edtext{\skt{पद्मो\-त्प\-ला\-कु\-ल\-ज\-ला\-म\-धि\-वा\-सि\-तोदां}}{\Afootnote{K H; \skt{पद्मो\-त्प\-ला\-कु\-ल\-ज\-ला\-म\-धि\-वा\-सि\-ता\-म्बुम् ॥}\ A}}\skt{ ।}\pend

\pstart\versepar \ledsidenote{\passref{6.47.59}}\edtext{\skt{वापी\-म्भ\-जे\-द्रु\-चि\-र\-च\-न्द\-न\-भू\-षि\-ताङ्गः}}{\Afootnote{K; \skt{वापीं भजेत हरि\-च\-न्द\-न\-भू\-षि\-ताङ्गः}\ A; \skt{वावी\-म्भ\-जे\-द्रु\-चि\-र\-च\-न्द\-न\-भू\-षि\-ताङ्गः}\ H}}\skt{ । }\edtext{\skt{कान्ता\-क\-र\-ग्र\-ह\-ण\-ह\-र्षि\-त\-रो\-म\-कूपः}}{\Afootnote{K; \skt{कान्ता\-क\-र\-स्पृ\-श\-न\-क\-र्क\-श\-रो\-म\-कूपः ।}\ A; \skt{कान्ता\-क\-र\-स्प\-ऋ\-श\-न\-ह\-र्षि\-त\-रो\-म\-कूपः ।}\ H}}\skt{ । }\pend
\pstart\versepar \skt{}\edtext{\skt{तत्रै\-न\-म\-म्वु\-रु\-ह\-प\-त्र\-समैः}}{\Afootnote{K\textsuperscript{pc}; \skt{तत्रै\-न\-म\-म्बु\-रु\-ह\-प\-त्र\-समैः}\ A H; \skt{तत्रै\-न\-म\-म्वु\-रु\-प\-ह\-प\-त्र\-समैः}\ K\textsuperscript{ac}}}\skt{ }\edtext{\skt{स्पृशत्यो}}{\Afootnote{K; \skt{स्पृशन्त्यः}\ A H}}\skt{ }\edtext{\skt{हस्तै\-र्नि\-म\-ग्न\-व\-दनैः}}{\Afootnote{K; \skt{शीतैः करोरुवदनैः}\ A; \skt{हस्तै\-र्न्नि\-म\-ग्न\-व\-दनैः}\ H}}\skt{ }\edtext{\skt{कठिनै}}{\Afootnote{K; \skt{कठिनैः}\ A H}}\skt{ स्तनैश्च ।}\pend

\pstart\versepar \ledsidenote{\passref{6.47.60}}\edtext{\skt{तोया\-व\-गा\-ह\-कु\-शला}}{\Afootnote{A K; \skt{तोया\-व\-गा\-ढ\-कु\-शला}\ H}}\skt{ }\edtext{\edtext{\skt{मधु\-र\-प्र\-ला\-पां\-स्सं\-ह\-र्ष\-ये\-यु\-र\-व\-लाश्च}}{\lemma{\skt{मधु\-र\-प्र\-ला\-पां\-स्सं\-ह\-र्ष\-ये\-यु\-र\-व\-लाश्}}\Afootnote{K; \skt{मधुरस्वभावाः संह\-र्ष\-ये\-यु\-र\-बलाः}\ A; \skt{मधुरप्रलापाः संह\-र्ष\-ये\-यु\-र\-ब\-लाश्}\ H}}}{\lemma{\skt{च}}\Afootnote{K H; \skt{सु}\ A}}\skt{ }\edtext{\skt{दुलैः}}{\Afootnote{K H; \skt{कलैः}\ A}}\skt{ }\edtext{\skt{स्वभावैः}}{\Afootnote{K H; \skt{प्रलापैः ।}\ A}}\skt{ ॥ }\pend
\pstart\versepar \skt{धारागृहे }\edtext{\skt{प्रच\-लि\-तो\-द\-र\-दु\-र्दि\-नान्ते}}{\Afootnote{K; \skt{प्रग\-लि\-तो\-द\-क\-दु\-र्दि\-नाभे}\ A; \skt{प्रग\-लि\-तो\-द\-र\-दु\-र्द्दि\-नान्ते}\ H}}\skt{ क्लान्तः शयीत सलि\-ला\-नि\-ल\-शी\-त\-कुक्षौ।}\pend

\pstart\versepar \ledsidenote{\passref{6.47.61}}\skt{गन्धोदकैः }\edtext{\skt{सकु\-सु\-मै\-रु\-प\-सि\-क्त\-भूके}}{\lemma{\skt{उपसिक्तभूके}}\Afootnote{K H; \skt{उपसिक्तभूमौ}\ A}}\skt{ }\edtext{\edtext{\skt{पत्रा\-म्वु\-च\-न्द\-न\-र\-सै\-रु\-प\-दि\-ग्ध\-कुड्ये}}{\lemma{\skt{पत्रा\-म्वु\-च\-न्द\-न\-र\-सैर्}}\Afootnote{K; \skt{पत्रा\-म्बु\-च\-न्द\-न\-र\-सैर्}\ A; \skt{तत्रा\-म्बु\-च\-न्द\-न\-र\-सैर्}\ H}}}{\lemma{\skt{उपदिग्धकुड्ये}}\Afootnote{K; \skt{उपलिप्तकुड्ये ।}\ A; \skt{उप\-दि\-ग\-ध\-कुड्ये ।}\ H}}\skt{ । }\pend
\pstart\versepar \skt{}\edtext{\skt{जात्यु\-त्प\-ल\-प्रि\-य\-क\-के\-श\-र\-पु\-ण्ड\-रीक}}{\Afootnote{A K; \skt{जात्यु\-त्प\-ल\-प्रि\-य\-त\-के\-स\-र\-पु\-ण्ड\-रीक}\ H}}\skt{ }\edtext{\skt{पुन्ना\-ग\-ना\-ग\-क\-र\-वी\-र\-कृ\-तो\-प\-कारै}}{\Afootnote{K; \skt{पुन्ना\-ग\-ना\-ग\-क\-र\-वी\-र\-कृ\-तो\-प\-चारे ॥}\ A; \skt{पुन्ना\-ग\-ना\-ग\-क\-र\-वी\-र\-कृ\-तो\-प\-कारे ॥}\ H}}\skt{ ।}\pend

\pstart\versepar \ledsidenote{\passref{6.47.62}}\edtext{\skt{तस्मिंगृहे}}{\Afootnote{K; \skt{तस्मिन्गृहे}\ A H}}\skt{ कम\-ल\-रे\-ण्व\-रुणे शयीत }\edtext{\skt{यत्ना\-हृ\-ता\-वि\-कं\-पि\-त\-पु\-ष्प\-दाम्नि}}{\Afootnote{K; \skt{यत्ना\-हृ\-ता\-नि\-ल\-वि\-क\-म्पि\-त\-पु\-ष्प\-दाम्नि ।}\ A; \skt{यन्ना\-द्वृ\-ता\-नि\-ल\-वि\-क\-म्पि\-त\-पु\-ष्प\-दाम्नि ।}\ H}}\skt{ । }\pend
\pstart\versepar \skt{}\edtext{\skt{पार्या\-त्र\-वि\-न्ध्य\-हि\-म\-व\-न्म\-ल\-या\-च\-लानां}}{\Afootnote{K; \skt{हेम\-न्त\-वि\-न्ध्य\-हि\-म\-व\-न्म\-ल\-या\-च\-लानां}\ A; \skt{पार्या\-त्र\-वि\-न्ध्य\-हि\-म\-व\-त्म\-ल\-या\-च\-लानां}\ H}}\skt{ शीताम्भसां सक\-द\-ली\-ह\-रि\-त\-द्रु\-माणां ।}\pend

\pstart\versepar \ledsidenote{\passref{6.47.63}}\edtext{\skt{उद्भि\-न्न\-नी\-ल\-न\-लि\-ना\-म्वु\-रु\-हां\-कु\-राणां}}{\Afootnote{K; \skt{उद्भि\-न्न\-नी\-ल\-न\-लि\-ना\-म्बु\-रु\-हा\-क\-राणां}\ A; \skt{उद्भि\-न्न\-नी\-ल\-न\-लि\-ना\-म्बु\-रु\-हा\-ङ्कु\-रा\-णाञ्}\ H}}\skt{ }\edtext{\skt{चन्द्रोदयस्य}}{\Afootnote{K; \skt{चन्द्रोदयस्य च}\ A H}}\skt{ }\edtext{\skt{शृणुयाच्चकथां}}{\Afootnote{K; \skt{कथाः शृणुयान्}\ A; \skt{शृणुयाच्च कथाम्}\ H}}\skt{ }\edtext{\skt{मनोज्ञां}}{\Afootnote{K H; \skt{मनोज्ञाः ।}\ A}}\skt{ । }\pend
\pstart\versepar \skt{}\edtext{\skt{म्लानं}}{\Afootnote{A K; \skt{ग्लानं}\ H}}\skt{ }\edtext{\skt{मुदीनमनसं}}{\Afootnote{K; \skt{प्रतान्तमनसं}\ A; \skt{सुदीनमनसं}\ H}}\skt{ }\edtext{\skt{मनसोनुकूलाः}}{\Afootnote{K; \skt{मनसोऽनुकूलाः}\ A H}}\skt{ । }\edtext{\skt{पीन\-स्त\-नो\-रु\-ज\-घ\-ना\-स्त\-रु\-सा\-र\-दिग्धाः}}{\Afootnote{K H; \skt{पीन\-स्त\-नो\-रु\-ज\-घना हरिचन्दनाङ्ग्यः ॥}\ A}}\skt{ ॥}\pend

\pstart \ledsidenote{\passref{6.47.64}}\edtext{\skt{तास्वे\-न\-मा\-र्द्र\-व\-सनाः}}{\Afootnote{K; \skt{ता एन\-मा\-र्द्र\-व\-सनाः}\ A; \skt{ता स्वेव\-मा\-र्द्द\-व\-स\-नाम्}\ H}}\skt{ }\edtext{\skt{सह}}{\Afootnote{A K; \skt{मह}\ H}}\skt{ सम्विशेयुः }\edtext{\skt{श्लिष्यावलाः}}{\Afootnote{K; \skt{श्लिष्ट्वावलाः}\ A; \skt{श्लिष्यावलोः}\ H}}\skt{ }\edtext{\skt{शिथि\-ल\-मे\-ख\-ल\-हा\-र\-यष्ट्यः}}{\Afootnote{A K; \skt{शिथि\-ल\-मे\-ख\-ल\-हा\-र\-यष्ट्याः ॥}\ H}}\skt{ ।}\pend

\pstart\versepar \ledsidenote{\passref{6.47.65}}\edtext{\skt{हर्षयेयुर्नरं}}{\Afootnote{A K H\textsuperscript{pc}; \skt{हर्ष\-ये\-यु\-र्न्न\-रंन्}\ H\textsuperscript{ac}}}\skt{ नार्यः }\edtext{\skt{स्वगुणैरम्भसि}}{\Afootnote{K H; \skt{स्वगुणै रहसि}\ A}}\skt{ }\edtext{\skt{स्थिताः}}{\Afootnote{K H; \skt{स्थिताः ।}\ A}}\skt{ । }\pend
\pstart\versepar \skt{}\edtext{\skt{हन्युःशैत्यान्विताः}}{\Afootnote{K H; \skt{ताः}\ A}}\skt{ }\edtext{\skt{पित्तं}}{\Afootnote{K H; \skt{शैत्या}\ A}}\skt{ }\edtext{\skt{हर्षयेयुरतः}}{\Afootnote{K H; \skt{च्छमयेयुश्च पित्त\-पा\-ना\-त्य\-या\-न्त\-रम् ॥}\ A}}\skt{ }\edtext{\skt{स्त्रियः}}{\Afootnote{K H; \emph{om.}\ A}}\skt{ ।}\pend

\pstart\versepar \ledsidenote{\passref{6.47.66}}\edtext{\skt{रक्त\-पि\-त्त\-स्तृ\-षा\-दाहे}}{\Afootnote{K; \skt{तृड्दा\-ह\-र\-क्त\-पि\-त्तेषु कार्यो}\ A; \skt{रक्त\-पि\-त्त\-स्तृ\-षा\-दाहे}\ H}}\skt{ }\edtext{\edtext{\skt{स्वयमेव}}{\lemma{\skt{स्वयम्}}\Afootnote{K H; \skt{ऽयं}\ A}}}{\lemma{\skt{एव}}\Afootnote{K H; \skt{भेष}\ A}}\skt{ }\edtext{\skt{विधि}}{\Afootnote{K; \skt{ज}\ A; \skt{विधिः}\ H}}\skt{ }\edtext{\skt{स्मृतः}}{\Afootnote{K H; \skt{क्रमः ।}\ A}}\skt{ । }\pend
\pstart\versepar \skt{सामान्यतो }\edtext{\skt{विशेषस्तु}}{\Afootnote{K H; \skt{विशेषं तु}\ A}}\skt{ शृणु }\edtext{\skt{दाहेश्वशेषतः}}{\Afootnote{K; \skt{दाहेष्वशेषतः ॥}\ A H}}\skt{ ।}\pend

\pstart\versepar \ledsidenote{\passref{6.47.67}}\skt{कृत्स्नदेहानुगं }\edtext{\skt{रक्तमुद्वृन्तः}}{\Afootnote{K; \skt{रक्तमुद्रिक्तं}\ A; \skt{रक्तंमुद्वृत्तः}\ H}}\skt{ }\edtext{\skt{प्रदहेन्नरं}}{\Afootnote{K H; \skt{दहति ह्यति । सं}\ A}}\skt{ । }\pend
\pstart\versepar \skt{}\edtext{\skt{दूष्यते}}{\Afootnote{K; \skt{चूष्यते}\ A; \skt{ऊष्यते}\ H}}\skt{ }\edtext{\skt{चूष्यते}}{\Afootnote{K H; \skt{दह्यते}\ A}}\skt{ }\edtext{\skt{चापि}}{\Afootnote{K H; \skt{च}\ A}}\skt{ ताम्रा\-भ\-स्ता\-म्र\-लो\-चनः ।}\pend

\pstart\versepar \ledsidenote{\passref{6.47.68}}\skt{लोह\-ग\-न्धा\-ङ्ग\-व\-दनो }\edtext{\skt{वह्निरेव}}{\Afootnote{K; \skt{वह्निनेवा}\ A; \skt{वह्निनेव}\ H}}\skt{ }\edtext{\skt{सदह्यते}}{\Afootnote{K H; \skt{वकीर्यते ।}\ A}}\skt{ । }\pend
\pstart\versepar \skt{}\edtext{\skt{कषा\-य\-ति\-क्त\-म\-धुरैः}}{\Afootnote{K H; \skt{तं विलङ्घ्य विधानेन}\ A}}\skt{ }\edtext{\skt{संस\-र्गा\-हा\-र\-मा\-दि\-शेत्}}{\Afootnote{K H; \skt{संसृ\-ष्टा\-हा\-र\-मा\-च\-रेत् ॥}\ A}}\skt{ ।}\pend

\pstart\versepar \ledsidenote{\passref{6.47.69}}\edtext{\skt{अशाम्यत्यथवा}}{\Afootnote{K H; \skt{अप्रशाम्यति}\ A}}\skt{ }\edtext{\skt{दाहे}}{\Afootnote{K; \skt{दाहे च}\ A; \skt{दाहो}\ H}}\skt{ }\edtext{\skt{रसैस्तृप्तस्य}}{\Afootnote{K; \skt{रसैस्तृप्तस्य}\ A; \skt{रसैस्तप्तस्य}\ H}}\skt{ जाङ्गलैः । }\pend
\pstart\versepar \skt{}\edtext{\skt{शाखाश्रयां}}{\Afootnote{K H; \skt{शाखाश्रया}\ A}}\skt{ यथान्यायं }\edtext{\skt{रोहिण्यां}}{\Afootnote{K; \skt{रोहिणी}\ A; \skt{रोगाणां}\ H}}\skt{ }\edtext{\skt{वेधयेत्शिरां}}{\Afootnote{K; \skt{र्व्यधयेत्सिराः ॥}\ A; \skt{वेधयेत्सिरां ॥}\ H}}\skt{ ।}\pend

\pstart\versepar \ledsidenote{\passref{6.47.70}}\edtext{\skt{पित्तज्वरहरो}}{\Afootnote{K; \skt{पित्तज्वरसमः पित्ता}\ A; \skt{पित्तज्वरंह}\ H}}\skt{ }\edtext{\edtext{\skt{यस्तु}}{\lemma{\skt{यस्}}\Afootnote{K; \skt{त्}\ A; \emph{om.}\ H}}}{\lemma{\skt{तु}}\Afootnote{K; \skt{स}\ A; \emph{om.}\ H}}\skt{ }\edtext{\skt{विधिःसोप्यत्र}}{\Afootnote{\skt{चाप्यस्य विधिर्}\ A; \skt{विधिःसो प्यत्र}\ K; \skt{रोयस्तु विधिः सो प्यत्र}\ H}}\skt{ }\edtext{\skt{पूजितः}}{\Afootnote{K H; \skt{हितः ।}\ A}}\skt{ । }\pend
\pstart\versepar \skt{}\edtext{\skt{तृष्णा\-नि\-रो\-धा\-द\-ब्धातौ}}{\Afootnote{K; \skt{तृष्णा\-नि\-रो\-धा\-द\-ब्धातौ}\ A; \skt{तृष्णा\-नि\-रो\-धा\-द\-ग्धादौ}\ H}}\skt{ क्षीणे तेजः }\edtext{\skt{समुद्धतः}}{\Afootnote{K H; \skt{समुद्धतम् ॥}\ A}}\skt{ ।}\pend

\pstart\versepar \ledsidenote{\passref{6.47.71}}\skt{स बाह्याभ्यन्तरं देहं }\edtext{\skt{प्रद\-हे\-त्म\-न्द\-चे\-तसः}}{\Afootnote{K; \skt{दहेद्वै मन्दचेतसः ।}\ A; \skt{प्रदेहे मन्दचेतसः ।}\ H}}\skt{ ॥ }\pend
\pstart\versepar \skt{}\edtext{\skt{स}}{\Afootnote{K H; \skt{सं}\ A}}\skt{ }\edtext{\skt{शुष्क\-ग\-ल\-ता\-ल्वौष्ठो}}{\Afootnote{K H; \skt{शुष्क\-ग\-ल\-ता\-ल्वोष्ठो}\ A}}\skt{ जिह्वां }\edtext{\skt{निःक्रिष्य}}{\Afootnote{K; \skt{निष्कृष्य}\ A H}}\skt{ }\edtext{\skt{चेष्टते}}{\Afootnote{A K; \skt{वेष्टने ॥}\ H}}\skt{ ।}\pend

\pstart\versepar \ledsidenote{\passref{6.47.72}}\edtext{\skt{तेज\-स्त\-त्रो\-प\-श\-म\-ये\-द\-व्धा\-तुश्च}}{\Afootnote{K; \skt{तत्रो\-प\-श\-म\-ये\-त्ते\-ज\-स्त्व\-ब्धातुं च}\ A; \skt{तेज\-स्त\-त्रो\-प\-श\-म\-ये\-द\-र्ध्वातुं च}\ H}}\skt{ विवर्धयेत् । }\pend
\pstart\versepar \skt{पाय\-ये\-त्का\-म\-म\-म्भश्च }\edtext{\skt{शर्करान्नः}}{\Afootnote{K; \skt{शर्कराढ्यं}\ A; \skt{शर्क्कराम्भः}\ H}}\skt{ पयो पि वा ।}\edtext{}{\lemma{\skt{विवर्धयेत् । पाय\-ये\-त्का\-म\-म\-म्भश्च}}\Afootnote{K; \skt{विवर्धयेत् । पाय\-ये\-त्का\-म\-म\-म्भश्च}\ A; \skt{विर्द्धयेत् । पाय\-ये\-त्का\-म\-म\-म्भश्च}\ H}}\pend

\pstart \ledsidenote{\passref{6.47.73ab}}\edtext{\skt{शीतमिक्षुरसं}}{\Afootnote{K; \skt{शीतमिक्षुरसं}\ A H}}\skt{ मन्थं }\edtext{\skt{वितरेच्चेरितं}}{\Afootnote{K; \skt{वितरेच्चेरितं}\ A; \skt{वित\-रे\-च्छो\-षे\-रितं}\ H}}\skt{ विधिं ।}\pend

\pstart \ledsidenote{\passref{6.47.73a}}\edtext{\skt{अप्राप्तम्वा}}{\Afootnote{K\textsuperscript{pc} H; \emph{om.}\ K\textsuperscript{ac}}}\skt{ }\edtext{\skt{प्रार्थयता}}{\Afootnote{K\textsuperscript{pc} H; \emph{om.}\ K\textsuperscript{ac}}}\skt{ }\edtext{\skt{तथा}}{\Afootnote{K\textsuperscript{pc} H; \emph{om.}\ K\textsuperscript{ac}}}\skt{ }\edtext{\edtext{\skt{चिन्तयतोऽपि}}{\lemma{\skt{चिन्तयतो}}\Afootnote{H; \emph{om.}\ K\textsuperscript{ac}; \skt{चिन्तयतो}\ K\textsuperscript{pc}}}}{\lemma{\skt{ऽपि}}\Afootnote{H; \emph{om.}\ K\textsuperscript{ac}; \skt{पि}\ K\textsuperscript{pc}}}\skt{ }\edtext{\skt{वा}}{\Afootnote{K\textsuperscript{pc} H; \emph{om.}\ K\textsuperscript{ac}}}\skt{ ।}\pend

\pstart \ledsidenote{\passref{6.47.73b}}\edtext{\skt{आहारस्य}}{\Afootnote{K\textsuperscript{pc} H\textsuperscript{pc}; \emph{om.}\ K\textsuperscript{ac}; \skt{आहारआस्य}\ H\textsuperscript{ac}}}\skt{ }\edtext{\skt{विरोधाच्च}}{\Afootnote{K\textsuperscript{pc} H; \emph{om.}\ K\textsuperscript{ac}}}\skt{ }\edtext{\edtext{\skt{शरी\-र\-मु\-प\-न\-श्यते}}{\lemma{\skt{शरीरम्}}\Afootnote{K\textsuperscript{pc} H; \emph{om.}\ K\textsuperscript{ac}}}}{\lemma{\skt{उपनश्यते}}\Afootnote{K\textsuperscript{pc}; \emph{om.}\ K\textsuperscript{ac}; \skt{अपनश्यति ।}\ H}}\skt{ ।}\pend

\pstart \ledsidenote{\passref{6.47.73c}}\edtext{\skt{ततः}}{\Afootnote{K\textsuperscript{pc} H; \emph{om.}\ K\textsuperscript{ac}}}\skt{ }\edtext{\skt{क्षयमवाप्नोति}}{\Afootnote{K\textsuperscript{pc} H; \emph{om.}\ K\textsuperscript{ac}}}\skt{ }\edtext{\edtext{\skt{दाह\-श्चा\-स्यो\-प\-जा\-यते}}{\lemma{\skt{दाहश्}}\Afootnote{K\textsuperscript{pc} H; \emph{om.}\ K\textsuperscript{ac}}}}{\lemma{\skt{चास्योपजायते}}\Afootnote{K\textsuperscript{pc} H; \emph{om.}\ K\textsuperscript{ac}}}\skt{ ।}\pend

\pstart \ledsidenote{\passref{6.47.73d}}\edtext{\skt{प्रमोहश्चश्च}}{\Afootnote{K\textsuperscript{pc}; \emph{om.}\ K\textsuperscript{ac}; \skt{प्रमोहश्च}\ H}}\skt{ }\edtext{\skt{प्रलापश्च}}{\Afootnote{K\textsuperscript{pc} H; \emph{om.}\ K\textsuperscript{ac}}}\skt{ }\edtext{\skt{मूर्च्छा\-शी\-ता\-भि\-न\-न्दनं}}{\Afootnote{K\textsuperscript{pc} H; \emph{om.}\ K\textsuperscript{ac}}}\skt{ ।}\pend

\pstart \ledsidenote{\passref{6.47.73e}}\edtext{\skt{रुजाभिः}}{\Afootnote{K\textsuperscript{pc} H; \emph{om.}\ K\textsuperscript{ac}}}\skt{ }\edtext{\edtext{\skt{पीड्यतेत्यर्थं}}{\lemma{\skt{पीड्यतेत्य्}}\Afootnote{K\textsuperscript{pc}; \emph{om.}\ K\textsuperscript{ac}; \skt{पीड्यते}\ H}}}{\lemma{\skt{अर्थं}}\Afootnote{K\textsuperscript{pc}; \emph{om.}\ K\textsuperscript{ac}; \skt{त्यर्थञ्}\ H}}\skt{ }\edtext{\skt{जन्तु\-र्दा\-ह\-क्ष\-या\-त्मके}}{\Afootnote{K\textsuperscript{pc}; \emph{om.}\ K\textsuperscript{ac}; \skt{जन्तुर्द्दाहे क्षयात्मके ॥}\ H}}\skt{ ।}\pend

\pstart \ledsidenote{\passref{6.47.73f}}\edtext{\skt{इष्टाः}}{\Afootnote{K\textsuperscript{pc} H; \emph{om.}\ K\textsuperscript{ac}}}\skt{ }\edtext{\edtext{\skt{शब्दादयस्तत्र}}{\lemma{\skt{शब्दादयस्}}\Afootnote{K\textsuperscript{pc} H; \emph{om.}\ K\textsuperscript{ac}}}}{\lemma{\skt{तत्र}}\Afootnote{K\textsuperscript{pc} H; \emph{om.}\ K\textsuperscript{ac}}}\skt{ }\edtext{\skt{पित्तघ्नश्च}}{\Afootnote{K\textsuperscript{pc} H; \emph{om.}\ K\textsuperscript{ac}}}\skt{ }\edtext{\skt{विधिः}}{\Afootnote{K\textsuperscript{pc} H; \emph{om.}\ K\textsuperscript{ac}}}\skt{ }\edtext{\skt{स्मृतः}}{\Afootnote{K\textsuperscript{pc} H; \emph{om.}\ K\textsuperscript{ac}}}\skt{ ।}\pend

\pstart \ledsidenote{\passref{6.47.73cd}}\edtext{\skt{असृजः}}{\Afootnote{K H; \skt{असृजा}\ A}}\skt{ पूर्ण्णकोष्ठस्य दाहो भवति }\edtext{\skt{दुस्तरः}}{\Afootnote{K H; \skt{दुःसहः ॥}\ A}}\skt{ ।}\pend

\pstart\versepar \ledsidenote{\passref{6.47.74}}\skt{विधिः }\edtext{\skt{सद्यो}}{\Afootnote{K; \skt{सद्यो}\ A H}}\skt{ }\edtext{\skt{व्रणीयोक्तं}}{\Afootnote{K; \skt{व्रणीयोक्तस्}\ A H}}\skt{ तस्य लक्षणमेव च । }\pend
\pstart\versepar \skt{}\edtext{\skt{धातुक्षयोक्तो यो}}{\Afootnote{A K; \skt{धातुक्षयो}\ H}}\skt{ }\edtext{\edtext{\skt{दाहस्तेन}}{\lemma{\skt{दाहस्}}\Afootnote{A K; \skt{दाह}\ H}}}{\lemma{\skt{तेन}}\Afootnote{A K; \skt{स्… तेन}\ H}}\skt{ मूर्च्छातृषान्वितः ।}\pend

\pstart \ledsidenote{\passref{6.47.75ab}}\skt{क्षामस्वरः क्रियाहीनः }\edtext{\skt{सीदते}}{\Afootnote{K H; \skt{भृशं सीदति}\ A}}\skt{ }\edtext{\skt{भृशपीडितः}}{\Afootnote{K H; \skt{पीडितः ।}\ A}}\skt{ ॥}\pend

\pstart\versepar \ledsidenote{\passref{6.47.77}}\edtext{\edtext{\skt{तमि\-ष्ट\-वि\-ष\-यो\-पेतं}}{\lemma{\skt{तम्}}\Afootnote{K; \skt{त}\ A H}}}{\lemma{\skt{इष्ट\-वि\-ष\-यो\-पेतं}}\Afootnote{K; \skt{मिष्ट\-वि\-ष\-यो\-पेतं}\ A H}}\skt{ सुहृ\-द्भि\-र\-भि\-सं\-वृतं । }\pend
\pstart\versepar \skt{क्षीर\-मां\-स\-र\-सा\-हारं विधिनोक्तेन साधयेत् ।}\pend

\pstart\versepar \ledsidenote{\passref{6.47.78}}\edtext{\skt{मर्माभिघातजो}}{\Afootnote{A K; \skt{ममाभिघातजो}\ H}}\skt{ प्यस्ति }\edtext{\skt{सोसाध्यः}}{\Afootnote{K; \skt{स चासाध्य}\ A; \skt{सोऽसाध्यः}\ H}}\skt{ }\edtext{\skt{सप्तमो}}{\Afootnote{K H; \skt{तमः}\ A}}\skt{ }\edtext{\skt{मदः}}{\Afootnote{K H; \skt{स्मृतः ।}\ A}}\skt{ । }\pend
\pstart\versepar \skt{सर्व }\edtext{\skt{एव तु}}{\Afootnote{K; \skt{एव च}\ A; \skt{एव तु च}\ H}}\skt{ }\edtext{\skt{वर्ज्या स्युः}}{\Afootnote{K; \skt{वर्ज्याः स्युः}\ A H}}\skt{ शीतगात्रेषु देहिषु ।}\pend

\pstart \ledsidenote{\passref{6.47.79}}\edtext{\skt{प्रशा\-न्तो\-प\-द्र\-व\-श्चापि}}{\lemma{\skt{प्रशा\-न्तो\-प\-द्र\-वश्}}\Afootnote{K; \skt{(एवंविधो भवेद्यस्तु मदि\-रा\-म\-य\-पी\-डितः ।) प्रशान्तोपद्रवे}\ A; \skt{प्रशा\-न्तो\-प\-द्र\-वाश्}\ H}}\skt{ }\edtext{\skt{शोधणं}}{\Afootnote{K; \skt{शोधनं}\ A H}}\skt{ }\edtext{\edtext{\skt{प्राप्तमाचरेत्}}{\lemma{\skt{प्राप्तम्}}\Afootnote{A K; \skt{प्राक्स}\ H}}}{\lemma{\skt{आचरेत्}}\Afootnote{A K; \skt{माचरेत् ॥}\ H}}\skt{ ।}\pend

\pstart\versepar \ledsidenote{\passref{6.47.80}}\edtext{\skt{सजी\-र\-क\-ण्या\-र्द्र\-क\-शृं\-ग\-वे\-र\-सौ\-व\-र्च\-ला\-र्न्य\-व\-ज\-ला\-प्लु\-तानि}}{\Afootnote{K; \skt{सजी\-र\-का\-ण्या\-र्द्र\-क\-शृ\-ङ्ग\-वे\-र\-सौ\-व\-र्च\-ला\-न्य\-र्ध\-ज\-ल\-प्लु\-तानि ।}\ A; \skt{सजी\-र\-का\-ण्या\-र्द्द्र\-क\-शृं\-ग\-वेरं सौव\-र्च्च\-ला\-ण्य\-र्ध\-ज\-ला\-प्लु\-तानि ।}\ H}}\skt{ । }\pend
\pstart\versepar \skt{मद्यानि }\edtext{\skt{हृद्यानि च}}{\Afootnote{K H; \skt{हृद्यान्यथ}\ A}}\skt{ }\edtext{\skt{गन्धवन्ति}}{\Afootnote{A K; \skt{सारवन्ति}\ H}}\skt{ पीतानि सद्यः शमयन्ति तृष्णां ॥}\pend

\pstart\versepar \ledsidenote{\passref{6.47.81}}\edtext{\edtext{\skt{जला\-प्लु\-त\-श्च\-न्द\-न\-भू\-षि\-तांगः}}{\lemma{\skt{जलाप्लुतश्}}\Afootnote{K; \skt{जलप्लुतश्}\ A; \skt{जलाप्लुतं}\ H}}}{\lemma{\skt{चन्द\-न\-भू\-षि\-तांगः}}\Afootnote{A K; \skt{चन्द\-न\-रू\-षि\-ताङ्गः}\ H}}\skt{ स्रग्वी }\edtext{\skt{सभक्तां}}{\Afootnote{A K; \skt{सभक्ता}\ H}}\skt{ }\edtext{\skt{पिशितोपदंशां}}{\Afootnote{A K; \skt{न्पिशितोपदंशा}\ H}}\skt{ । }\pend
\pstart\versepar \skt{}\edtext{\skt{पिबेत्सुरां}}{\Afootnote{\skt{पिबेन्सुरां}\ A; \skt{पिवेत्सुरां}\ K; \skt{न् । पिबेत्सुरां}\ H}}\skt{ नैव लभेत }\edtext{\skt{रोगां}}{\Afootnote{K; \skt{रोगा}\ A H}}\skt{ }\edtext{\skt{मनो}}{\Afootnote{K; \skt{न्मनो}\ A H}}\skt{ }\edtext{\skt{मतिघ्नञ्च}}{\Afootnote{K; \skt{नुविघ्नं च}\ A; \skt{गतिघ्नञ्च}\ H}}\skt{ मदन्न }\edtext{\skt{यातीति}}{\Afootnote{K; \skt{याति ॥}\ A; \skt{याति इति ॥ ॥}\ H}}\skt{ ॥०॥}\pend

\pstart \ledsidenote{\passref{6.47.trailer}}\edtext{\skt{कायचि}}{\Afootnote{K; \skt{इति सुश्रु\-त\-स्ं\-हि\-ता\-या\-मु\-त्त\-र\-त\-न्त्रा\-न्त\-र्गते काय\-चि\-कि\-त्सा\-तन्त्रे मदा\-त्य\-य\-प्र\-ति\-षेधो नाम (नवमोऽध्यायः, आदितः) सप्त\-च\-त्वा\-रिं\-श\-त्त\-मोऽ\-ध्यायः ॥४७ ॥}\ A; \skt{उत्तरतन्त्रे पान\-त्य\-य\-चि\-कित्सा चतु\-श्च\-त्वा\-रिं\-श\-त्तमः काय\-चि\-कि\-त्सा\-या\-न्न\-व\-मो\-ध्यायः ॥}\ H}}\skt{ ॥९}\pend

\endnumbering
\beginnumbering
\eledsection*{Adhyāya 48}
\pstart \ledsidenote{\passref{6.48.1}}\edtext{\skt{अथात}}{\Afootnote{K H; \skt{अथातस्}\ A}}\skt{ }\edtext{\skt{तृष्णाप्रतिषेधं}}{\Afootnote{A K; \skt{स्तृष्णाप्रतिषेध}\ H}}\skt{ }\edtext{\skt{व्याख्यास्यामः}}{\Afootnote{\skt{अध्यायं व्याख्यास्यामः ॥}\ A; \emph{om.}\ K; \skt{मध्यायं व्याख्यास्यामः ॥}\ H}}\skt{ ॥}\pend

\pstart \ledsidenote{\passref{6.48.4a}}\edtext{\skt{भयभ्रमाभ्यां}}{\Afootnote{K; \skt{भयश्रमाभ्यां}\ H}}\skt{ रससंक्षयाद्वा ऊर्ध्वं }\edtext{\skt{चित\-म्पि\-त्त\-वि\-व\-व\-नैश्च}}{\lemma{\skt{पित्त\-वि\-व\-व\-नैश्}}\Afootnote{K; \skt{पित्त\-वि\-व\-र्द्ध\-नैश्}\ H}}\skt{ ।}\pend

\pstart \ledsidenote{\passref{6.48.4b}}\skt{पित्तं सवातं कुपितं नराणां तालुप्रपन्नं जनयेत्पिपासां ।}\pend

\pstart \ledsidenote{\passref{6.48.5}}\edtext{\skt{स्रोतः}}{\Afootnote{K H; \skt{स्रोतांसि संदूषयतः समेतौ यान्यम्बुवाहीनि शरीरिणां हि । स्रोतः}\ A}}\skt{ }\edtext{\skt{स्वपा\-म्व\-स्वु\-षु\-दू\-षि\-तेषु}}{\Afootnote{K; \skt{स्वपांवाहिषु दूषितेषु}\ A; \skt{स्वबावाम्बुषु दूषितेषु}\ H}}\skt{ }\edtext{\edtext{\skt{दोषैस्तृषां}}{\lemma{\skt{दोषैस्}}\Afootnote{K H; \skt{जायेत तृष्णातिबला ततस्}\ A}}}{\lemma{\skt{तृषां}}\Afootnote{K H; \skt{तु ॥}\ A}}\skt{ }\edtext{\skt{सम्भवतीह}}{\Afootnote{K; \emph{om.}\ A; \skt{सन्सवतीह}\ H}}\skt{ }\edtext{\skt{जन्तोः}}{\Afootnote{K H; \emph{om.}\ A}}\skt{ ।}\pend

\pstart\versepar \ledsidenote{\passref{6.48.6}}\skt{तिस्रः स्मृता }\edtext{\skt{लोहितजा}}{\Afootnote{K H; \skt{स्ताः क्षतजा}\ A}}\skt{ चतुर्थी }\edtext{\edtext{\skt{क्षया\-त्त\-था\-न्या\-म\-स\-मु\-द्भवा}}{\lemma{\skt{क्षयात्तथान्याम्}}\Afootnote{A K; \skt{क्षयात्तथान्या}\ H}}\skt{ च}}{\lemma{\skt{असमुद्भवा च}}\Afootnote{A K; \skt{मसमुद्भवा च ॥}\ H}}\skt{ । }\pend
\pstart\versepar \skt{स्यात्सप्तमी }\edtext{\skt{भक्तसमुद्भवा}}{\Afootnote{K H; \skt{भक्तनिमित्तजा तु}\ A}}\skt{ }\edtext{\skt{च}}{\Afootnote{K H; \skt{निबोध}\ A}}\skt{ लिङ्गानि }\edtext{\skt{तासां}}{\Afootnote{K H; \skt{अनुपूर्वशस्तु ॥}\ A}}\skt{ }\edtext{\skt{शृणु}}{\Afootnote{K H; \emph{om.}\ A}}\skt{ }\edtext{\skt{सौषधानि}}{\Afootnote{K H; \emph{om.}\ A}}\skt{ ॥}\pend

\pstart\versepar \ledsidenote{\passref{6.48.8}}\edtext{\skt{शुष्काम्यता}}{\Afootnote{K; \skt{शुष्कास्यता}\ A H}}\skt{ }\edtext{\skt{मारु\-त\-स\-म्भ\-वाया}}{\Afootnote{K H; \skt{मारु\-त\-सं\-भ\-वायां}\ A}}\skt{ तोदस्तथा }\edtext{\skt{कण्ठशिरस्सु}}{\Afootnote{K H; \skt{शङ्खशिरःसु}\ A}}\skt{ चापि । }\pend
\pstart\versepar \skt{}\edtext{\skt{स्रोतो}}{\Afootnote{K; \skt{स्रोतो}\ A H}}\skt{ }\edtext{\skt{निरोधा}}{\Afootnote{K; \skt{निरोधो}\ A; \skt{विरोधा}\ H}}\skt{ }\edtext{\skt{विरसश्च}}{\Afootnote{K H; \skt{विरसं च}\ A}}\skt{ वक्त्रं }\edtext{\skt{सीताभिरद्भिश्च}}{\lemma{\skt{सीताभिर्}}\Afootnote{K; \skt{शीतभिर्}\ A H}}\skt{ }\edtext{\skt{विवृद्धिमेति}}{\Afootnote{K H; \skt{विवृद्धिमेति ॥}\ A}}\skt{ ।}\pend

\pstart\versepar \ledsidenote{\passref{6.48.9}}\skt{मूर्च्छा }\edtext{\skt{प्रलापो}}{\Afootnote{K H; \skt{प्रलापा}\ A}}\skt{ }\edtext{\skt{रुचिवक्त्रशोषो}}{\Afootnote{K; \skt{रुचिवक्त्रशोषाः}\ A; \skt{रुचितालुशोषौ}\ H}}\skt{ }\edtext{\skt{रक्ते\-क्ष\-ण\-त्वं\-प्र\-त\-तश्च}}{\Afootnote{K H; \skt{पीतेक्षणत्वं प्रततश्च}\ A}}\skt{ }\edtext{\skt{दोष}}{\Afootnote{K; \skt{दाहः ।}\ A; \skt{दोष्ः ॥}\ H}}\skt{ । }\pend
\pstart\versepar \skt{}\edtext{\skt{सीताभिकांक्षा}}{\Afootnote{K; \skt{शीताभिकाङ्क्षा}\ A H}}\skt{ मुखतिक्तता च पित्तात्मिकायां }\edtext{\skt{परिधूपणञ्च}}{\lemma{\skt{परिधूपणञ्}}\Afootnote{K; \skt{परिधूपनं}\ A H}}\skt{ ।}\pend

\pstart \ledsidenote{\passref{6.48.9a}}\skt{बाष्पावरोधः }\edtext{\skt{कफसम्भृतेग्नौ}}{\Afootnote{K; \skt{कफसंभवाग्नौ}\ H}}\skt{ तृष्णा }\edtext{\skt{बलासेन}}{\Afootnote{K; \skt{वलासेन}\ H}}\skt{ भवेत्तथा तु ।}\pend

\pstart \ledsidenote{\passref{6.48.10}}\edtext{\skt{निद्रागुरुत्वं}}{\Afootnote{K H; \skt{कफा\-वृ\-ता\-भ्या\-म\-नि\-ला\-न\-लाभ्यां कफोऽपि शुष्कः प्रकरोति तृष्णाम् । निद्रा गुरुत्वं}\ A}}\skt{ मधुरास्यता च }\edtext{\skt{भयार्दितः}}{\Afootnote{K H; \skt{तयार्दितः}\ A}}\skt{ शुष्यति चातिमात्रं ।}\pend

\pstart \ledsidenote{\passref{6.48.12}}\skt{क्षतस्य रुक्छो\-णि\-त\-नि\-र्ग\-माभ्यां तृष्णा चतुर्थी क्षतजा }\edtext{\skt{भवेत्सा}}{\Afootnote{K H; \skt{मता तु । तयाभिभूतस्य निशादिनानि गच्छन्ति दुःखं पिबतोऽपि तोयम् ॥}\ A}}\skt{ ।}\pend

\pstart\versepar \ledsidenote{\passref{6.48.13}}\edtext{\skt{रक्तक्षयाद्या}}{\Afootnote{K H; \skt{रसक्षयाद्या}\ A}}\skt{ क्षयजा }\edtext{\skt{मतासौ}}{\Afootnote{K H; \skt{मता सा}\ A}}\skt{ तयार्दितः शुष्यति दह्यते च । }\pend
\pstart\versepar \skt{}\edtext{\skt{अभ्य\-र्थ\-मा\-कां\-क्षति}}{\Afootnote{K H; \skt{अत्य\-र्थ\-मा\-का\-ङ्क्षति}\ A}}\skt{ चापि }\edtext{\skt{तोयं तां}}{\Afootnote{A K; \skt{नोयं तां}\ H}}\skt{ }\edtext{\skt{सन्निपातादिति}}{\Afootnote{K; \skt{सन्निपातादिति}\ A H}}\skt{ }\edtext{\skt{केचिदाहुः}}{\Afootnote{A K; \skt{केचिदाहुः ॥}\ H}}\skt{ ।}\pend

\pstart\versepar \ledsidenote{\passref{6.48.14}}\edtext{\skt{रक्तक्षयोक्तानि च}}{\Afootnote{K H; \skt{रसक्षयोक्तानि च}\ A}}\skt{ लक्षणानि तस्यामशेषेण भिषग्व्यवस्येत् । }\pend
\pstart\versepar \skt{त्रिदो\-ष\-लि\-ङ्गा\-म\-स\-मु\-द्भवा तु }\edtext{\skt{हृच्छू\-ल\-नि\-ष्ठी\-व\-न\-सं\-प्र\-युक्ता}}{\Afootnote{K H; \skt{हृच्छू\-ल\-नि\-ष्ठी\-व\-न\-सा\-द\-युक्ता ॥}\ A}}\skt{ ।}\edtext{}{\lemma{\skt{भिषग्व्यवस्येत् । त्रिदो\-ष\-लि\-ङ्गा\-म\-स\-मु\-द्भवा तु}}\Afootnote{K H; \skt{भिषग्व्यवस्येत् । त्रिदो\-ष\-लि\-ङ्गा\-म\-स\-मु\-द्भवा च}\ A}}\pend

\pstart \ledsidenote{\passref{6.48.15}}\skt{स्निग्धं तथाम्लं }\edtext{\skt{लवणञ्च}}{\Afootnote{A K; \skt{लवण च}\ H}}\skt{ भुक्तं }\edtext{\skt{गुर्वन्नमेवापि}}{\Afootnote{K; \skt{गुर्वन्नमेवाति}\ A; \skt{गुर्वन्नमेवापि}\ H}}\skt{ तृषां }\edtext{\skt{करोति}}{\Afootnote{K H; \skt{करोति । क्षीणं विचित्तं बधिरं तृषार्तं विव\-र्ज\-ये\-न्नि\-र्ग\-त\-जि\-ह्व\-माशु ॥}\ A}}\skt{ ।}\pend

\pstart\versepar \ledsidenote{\passref{6.48.16}}\edtext{\skt{तृष्णा\-भि\-वृ\-द्धा\-वु\-दरे च}}{\Afootnote{A K; \skt{तृष्णा\-भि\-वृ\-द्धा\-मु\-दरे च}\ H}}\skt{ पूर्ण्णे तं }\edtext{\skt{वाम\-ये\-त्मा\-ग\-धि\-को\-द\-केन}}{\Afootnote{K H; \skt{वाम\-ये\-न्मा\-ग\-धि\-को\-द\-केन ।}\ A}}\skt{ । }\pend
\pstart\versepar \skt{}\edtext{\skt{विलेखनञ्चात्र}}{\Afootnote{K H; \skt{विलोभनं चात्र}\ A}}\skt{ हितं }\edtext{\skt{वदन्ति}}{\Afootnote{K H; \skt{विधेयं}\ A}}\skt{ स्याद्दाडि माम्रा\-त\-क\-मा\-तु\-लुंगैः ।}\pend

\pstart\versepar \ledsidenote{\passref{6.48.17}}\edtext{\skt{तृष्णाप्रयोगैः}}{\Afootnote{K H; \skt{तिस्रः प्रयोगै}\ A}}\skt{ }\edtext{\skt{प्रतिवारणीया}}{\Afootnote{K H; \skt{रिह सन्निवार्याः}\ A}}\skt{ शीतैश्च सम्य\-ग्र\-स\-वी\-र्य\-जातैः । }\pend
\pstart\versepar \skt{}\edtext{\edtext{\skt{गण्डू\-ष\-म\-म्लै\-र्वि\-रसे}}{\lemma{\skt{गण्डूषमम्लैर्}}\Afootnote{K; \skt{गण्डूषमम्लै}\ A; \skt{गण्डूषमम्लैर्}\ H}}\skt{ च}}{\lemma{\skt{विरसे च}}\Afootnote{K H; \skt{र्विरसे च}\ A}}\skt{ वक्त्रे }\edtext{\edtext{\skt{कुर्या\-च्छु\-भै\-रा\-म\-ल\-कस्य}}{\lemma{\skt{कुर्याच्छुभैर्}}\Afootnote{K H; \skt{कुर्याच्छुभै}\ A}}}{\lemma{\skt{आमलकस्य}}\Afootnote{K H; \skt{रामलकस्य}\ A}}\skt{ चूर्ण्णैः ।}\pend

\pstart\versepar \ledsidenote{\passref{6.48.18}}\edtext{\edtext{\edtext{\skt{सुव\-र्ण्ण\-रू\-प्या\-दि\-भि\-र\-ग्नि\-त\-प्तै\-र्लोष्ट्रैः}}{\lemma{\skt{सुव\-र्ण्ण\-रू\-प्या\-दि\-भिर्}}\Afootnote{A K; \skt{सुव\-र्ण्ण\-रु\-ध्या\-दि\-भिर्}\ H\textsuperscript{ac}; \skt{सुव\-र्ण्ण\-रु\-ध्या\-रू\-प्या\-दि\-भिर्}\ H\textsuperscript{pc}}}}{\lemma{\skt{अग्नितप्तैर्}}\Afootnote{K H; \skt{अग्नितप्तै}\ A}}}{\lemma{\skt{लोष्ट्रैः}}\Afootnote{K; \skt{र्लोष्टैः}\ A; \skt{ल्लोष्टैः}\ H}}\skt{ }\edtext{\skt{कृतैर्वा}}{\Afootnote{K; \skt{कृतं वा}\ A H}}\skt{ }\edtext{\skt{सिकतासु}}{\Afootnote{K H; \skt{सिकतादि}\ A}}\skt{ }\edtext{\skt{चापि}}{\Afootnote{K H; \skt{भिर्वा ।}\ A}}\skt{ । }\pend
\pstart\versepar \skt{}\edtext{\skt{जलैः}}{\Afootnote{K; \skt{जलं}\ A H}}\skt{ }\edtext{\skt{सुखोष्णैः}}{\Afootnote{K H; \skt{सुखोष्णं}\ A}}\skt{ शमयेत्तु तृष्णां }\edtext{\skt{सशर्करैः}}{\Afootnote{K; \skt{सशर्करं}\ A H}}\skt{ क्षौद्रयुतं }\edtext{\skt{हितम्वा}}{\lemma{\skt{हितम्}}\Afootnote{K H; \skt{हिमं}\ A}}\skt{ ॥}\pend

\pstart\versepar \ledsidenote{\passref{6.48.19}}\skt{पञ्चांगिकाः पञ्चगणा य }\edtext{\edtext{\skt{उक्ता\-स्ते\-ष्व\-म्बु\-सिद्धं}}{\lemma{\skt{उक्तास्}}\Afootnote{K; \skt{उक्ता}\ A H}}}{\lemma{\skt{तेष्व्}}\Afootnote{K; \skt{स्तेष्व्}\ A H}}\skt{ प्रथमे गणे वा । }\pend
\pstart\versepar \skt{}\edtext{\skt{पिबं}}{\Afootnote{K; \skt{पिबे}\ A; \skt{पिब}\ H}}\skt{ }\edtext{\skt{सुखोष्णं}}{\Afootnote{K; \skt{त्सुखोष्णं}\ A; \skt{न्सुखोष्णं}\ H}}\skt{ मनुजो चिरेण तृषो }\edtext{\skt{विमुच्येत}}{\Afootnote{A K; \skt{विमुच्ये हि}\ H}}\skt{ }\edtext{\skt{हि}}{\Afootnote{A K; \skt{त}\ H}}\skt{ वातजायाः ॥}\pend

\pstart \ledsidenote{\passref{6.48.20a}}\skt{पित्तोत्थितां पित्त\-ह\-रै\-र्वि\-पक्वं निहन्ति तोयं पय एव वापि ॥}\pend

\pstart\versepar \ledsidenote{\passref{6.48.21}}\edtext{\skt{बिल्वा\-ढ\-की\-क\-न्य\-सि\-प\-ञ्च\-मूलं}}{\Afootnote{K H; \skt{बिल्वा\-ढ\-की\-क\-न्य\-क\-प\-ञ्च\-मूली}\ A}}\skt{ }\edtext{\skt{दर्भेषु}}{\Afootnote{A K; \skt{गर्भ्भेषु}\ H}}\skt{ सिद्धं }\edtext{\skt{कफजं}}{\Afootnote{K H; \skt{कफजां}\ A}}\skt{ निहन्ति । }\pend
\pstart\versepar \skt{हित\-म्भ\-वे\-च्छ\-र्द\-न\-मेव चात्र तप्ते\-न\-नि\-म्ब\-प्र\-स\-वो\-द\-केन ।}\pend

\pstart\versepar \ledsidenote{\passref{6.48.22}}\skt{सर्वासु }\edtext{\skt{तृष्णास्वथवापि}}{\Afootnote{A K; \skt{तृष्णास्त्वथवापि}\ H}}\skt{ }\edtext{\skt{पैत्त्ं}}{\Afootnote{K; \skt{पैत्तं}\ A H}}\skt{ कुर्याद्विधिं तेन }\edtext{\skt{शमंव्रजन्ति}}{\Afootnote{K H; \skt{हि ता न सन्ति ।}\ A}}\skt{ । }\pend
\pstart\versepar \skt{}\edtext{\skt{प्रया\-ग\-तो\-दु\-म्ब\-र\-जो\-र\-सस्तु}}{\Afootnote{K; \skt{पर्या\-ग\-तो\-दु\-म्ब\-रजो रसस्तु}\ A; \skt{पर्यागतो दुम्बरजो रसस्तु}\ H}}\skt{ }\edtext{\skt{सश\-र्क\-रा\-स्त्व\-क्क्व\-थि\-तो\-द\-कम्वा}}{\lemma{\skt{सश\-र्क\-रा\-स्त्व\-क्क्व\-थि\-तो\-द\-कम्}}\Afootnote{K; \skt{सश\-र्क\-र\-त\-त्क्व\-थि\-तो\-दकं}\ A; \skt{सश\-र्क्क\-रा\-स्त्व\-क्व\-थि\-तो\-द\-कम्}\ H}}\skt{ ।}\pend

\pstart\versepar \ledsidenote{\passref{6.48.23}}\skt{वर्गस्य सिद्धस्य च }\edtext{\skt{शारिवादेः}}{\Afootnote{K H; \skt{सारिवादेः}\ A}}\skt{ }\edtext{\skt{पेयं}}{\Afootnote{K H; \skt{पात}\ A}}\skt{ }\edtext{\skt{जलं}}{\Afootnote{K H; \skt{व्यमम्भः शिशिरं}\ A}}\skt{ }\edtext{\skt{स्यात्तु}}{\Afootnote{K H; \emph{om.}\ A}}\skt{ }\edtext{\skt{तृषाभिभूते}}{\Afootnote{K; \skt{तृषार्तैः ।}\ A; \skt{तृ षाभिभूतेः ।}\ H}}\skt{ । }\pend
\pstart\versepar \skt{}\edtext{\skt{कशे\-रु\-शृं\-गा\-ट\-क\-प\-द्म\-कैश्च}}{\lemma{\skt{कशे\-रु\-शृं\-गा\-ट\-क\-प\-द्म\-कैश्}}\Afootnote{K H; \skt{कशे\-रु\-शृ\-ङ्गा\-ट\-क\-प\-द्ममो}\ A}}\skt{ }\edtext{\skt{विशेषु}}{\Afootnote{K; \skt{बिसेक्षु}\ A; \skt{बिशेषु}\ H}}\skt{ सिद्धं }\edtext{\skt{त्वथवा}}{\Afootnote{K H; \skt{वां निहन्}\ A}}\skt{ }\edtext{\skt{पिबेत}}{\Afootnote{K H; \skt{ति ॥}\ A}}\pend

\pstart\versepar \ledsidenote{\passref{6.48.24}}\edtext{\skt{जलो\-त्प\-लो\-शी\-र\-कु\-च\-न्द\-नानि}}{\Afootnote{K; \skt{लाजो\-त्प\-लो\-शी\-र\-कु\-च\-न्द\-नानि}\ A H}}\skt{ दत्वा प्रवाते निशि }\edtext{\skt{वामयेत्तु}}{\Afootnote{K; \skt{वासयेत्तु ।}\ A; \skt{वासये तु ।}\ H}}\skt{ }\pend
\pstart\versepar \skt{}\edtext{\edtext{\skt{तदु\-त्त\-म\-न्तो\-य\-मु\-दा\-र\-गन्धि}}{\lemma{\skt{तदुत्तमन्}}\Afootnote{A K; \skt{तदुत्तमन्}\ H}}}{\lemma{\skt{तोय\-मु\-दा\-र\-गन्धि}}\Afootnote{A K; \skt{तोय\-मु\-दा\-र\-गन्धि}\ H}}\skt{ }\edtext{\skt{सितायुतं}}{\Afootnote{A K; \skt{सितायुत}\ H}}\skt{ }\edtext{\skt{क्षौद्रयुतं}}{\Afootnote{A K; \skt{…ञ्}\ H}}\skt{ }\edtext{\skt{च}}{\Afootnote{K H; \skt{व}\ A}}\skt{ }\edtext{\skt{पेयं}}{\Afootnote{K H; \skt{दन्ति ॥}\ A}}\skt{ ।}\pend

\pstart\versepar \ledsidenote{\passref{6.48.25}}\skt{द्राक्षाप्रगाढञ्च हिताय वैद्य\-स्तृ\-ष्णा\-र्दि\-तेभ्यो वितरेन्नरेभ्यः । }\pend
\pstart\versepar \skt{}\edtext{\skt{सशारिवादौ}}{\Afootnote{K H; \skt{ससारिवादौ}\ A}}\skt{ }\edtext{\skt{तृणपञ्चमूली}}{\Afootnote{K H; \skt{तृणपञ्चमूले}\ A}}\skt{ तथोत्पलादौ प्रथमे गणे च ।}\pend

\pstart\versepar \ledsidenote{\passref{6.48.26}}\edtext{\skt{कुर्यात्कषायाणि}}{\Afootnote{K H; \skt{कुर्यात्कषायं च}\ A}}\skt{ }\edtext{\skt{यथैतदुक्तं}}{\Afootnote{K H; \skt{यथेरितेन}\ A}}\skt{ }\edtext{\skt{मधू\-क\-पु\-ष्पा\-दिषु}}{\Afootnote{A K; \skt{मधूकपुष्पादि}\ H}}\skt{ }\edtext{\skt{चापरेषु}}{\Afootnote{A K; \skt{…परेषु ।}\ H}}\skt{ । }\pend
\pstart\versepar \skt{राजा\-द\-न\-क्षी\-रि\-क\-पी\-त\-नेषु षट्पानकान्यत्र हितानि च }\edtext{\skt{स्युः}}{\Afootnote{A K; \skt{स्युः ॥ राजा\-द\-न\-क्षी\-रि\-क\-पी\-त\-ने\-षु\-ष\-ट्पा\-न\-का\-न्यत्र हितानि च स्युः ॥}\ H}}\skt{ ॥}\pend

\pstart\versepar \ledsidenote{\passref{6.48.27}}\edtext{\skt{सतु\-ण्डि\-के\-रा\-ण्य\-थवा}}{\Afootnote{A K; \skt{सतु\-ण्डि\-के\-रा\-ण्य\-थवा}\ H}}\skt{ पिबेत्तु पिष्टानि }\edtext{\skt{कार्यानि}}{\Afootnote{K; \skt{कार्पास}\ A; \skt{कार्यास}\ H}}\skt{ समुद्भवानि । }\pend
\pstart\versepar \skt{क्षतोद्भवां रुग्वि\-नि\-वा\-र\-णेन }\edtext{\skt{जाये\-द्र\-सा\-ना\-म\-सृ\-जश्च}}{\Afootnote{K; \skt{जये\-द्र\-सा\-ना\-म\-सृ\-जश्च}\ A H}}\skt{ }\edtext{\skt{पानैः}}{\Afootnote{A K; \skt{मानैः ॥}\ H}}\skt{ ।}\pend

\pstart\versepar \ledsidenote{\passref{6.48.28}}\skt{क्षयोत्थितां }\edtext{\skt{क्षीरजलं}}{\Afootnote{K; \skt{क्षीरघृतं}\ A; \skt{क्षीरजलां}\ H}}\skt{ }\edtext{\skt{निह\-न्या\-न्मां\-सो\-द\-कम्वा}}{\Afootnote{K H; \skt{निह\-न्या\-न्मां\-सो\-दकं वा}\ A}}\skt{ }\edtext{\skt{मधुरोदकम्वा}}{\Afootnote{K H; \skt{मधुकोदकं वा ।}\ A}}\skt{ । }\pend
\pstart\versepar \skt{}\edtext{\skt{आमोत्थितां}}{\Afootnote{K; \skt{आमोद्भवां}\ A; \skt{आमोत्थिआं}\ H}}\skt{ }\edtext{\skt{बिल्व\-व\-चा\-यु\-तानां}}{\Afootnote{K H; \skt{बिल्व\-व\-चा\-यु\-तैस्तु}\ A}}\skt{ }\edtext{\skt{जाये\-त्क\-षा\-यै\-रथ}}{\Afootnote{K; \skt{जये\-त्क\-षा\-यै\-रथ}\ A H}}\skt{ }\edtext{\skt{दीपनानां}}{\Afootnote{K H; \skt{दीपनीयैः ॥}\ A}}\skt{ ।}\pend

\pstart \ledsidenote{\passref{6.48.29}}\edtext{\edtext{\skt{गुर्व\-न्न\-दा\-मु\-ल्लि\-ख\-नै\-र्ज\-येच्च}}{\lemma{\skt{गुर्व\-न्न\-दा\-मु\-ल्लि\-ख\-नैर्}}\Afootnote{K; \skt{आम्रा\-त\-भ\-ल्ला\-त\-ब\-ला\-यु\-तानि पिबे\-त्क\-षा\-या\-ण्यथ दीपनानि । गुर्वन्नजातां वमनै}\ A; \skt{गुर्व\-न्न\-जा\-मु\-ल्लि\-खनै}\ H}}}{\lemma{\skt{जयेच्च}}\Afootnote{K; \skt{र्जयेच्च}\ A H}}\skt{ क्षयादृते }\edtext{\skt{सर्वकृताश्च}}{\Afootnote{K H; \skt{सर्वकृतां च}\ A}}\skt{ }\edtext{\skt{तृष्णाः}}{\Afootnote{K; \skt{तृष्णाम् ॥}\ A; \skt{तृष्णा ॥}\ H}}\skt{ ।}\pend

\pstart\versepar \ledsidenote{\passref{6.48.33}}\edtext{\skt{लेपावगाहौ}}{\Afootnote{A K; \skt{लेपावगादौ}\ H}}\skt{ परिषेचनानि कुर्यात्तथा }\edtext{\skt{शीतगृहाणि}}{\Afootnote{A K; \skt{शीतगुहानि}\ H}}\skt{ चापि । }\pend
\pstart\versepar \skt{संशोधनं }\edtext{\skt{क्षीर\-र\-सा\-घृ\-तानि}}{\Afootnote{K H; \skt{क्षीररसौ घृतानि}\ A}}\skt{ सर्वासु }\edtext{\skt{लेहा}}{\Afootnote{K; \skt{लेहान्}\ A H}}\skt{ }\edtext{\skt{मधुरां}}{\Afootnote{K; \skt{मधुरा}\ A; \skt{मधुना}\ H}}\skt{ }\edtext{\edtext{\skt{हिमांश्चेति}}{\lemma{\skt{हिमांश्}}\Afootnote{K; \skt{न्हिमांश्}\ A; \skt{हिमाश्}\ H}}\skt{ ॥ {\symfont ❈}}}{\lemma{\skt{चेति ॥ {\symfont ❈}}}\Afootnote{K; \skt{च ॥}\ A; \skt{चेति ॥}\ H}}\skt{ ॥}\pend

\pstart \ledsidenote{\passref{6.48.33a}}\edtext{\skt{ज्वरोतिसारः}}{\Afootnote{K; \skt{ज्वरोऽतिसारः}\ H}}\skt{ शोषश्च गुल्म\-हृ\-त्पा\-ण्डु\-लो\-हितः}\pend

\pstart \ledsidenote{\passref{6.48.33b}}\skt{मूर्च्छा पानात्ययश्चैव तृष्णया पूर्यते }\edtext{\skt{दशः}}{\Afootnote{K; \skt{दश… ॥}\ H}}\skt{ ॥}\pend

\pstart \ledsidenote{\passref{6.48.trailer}}\edtext{\skt{कायचिकित्सायां}}{\Afootnote{K; \skt{इति सुश्रु\-त\-सं\-हि\-ता\-या\-मु\-त्त\-र\-त\-न्त्रा\-न्त\-र्गते काय\-चि\-कि\-त्सा\-तन्}\ A; \skt{उत्तरतन्त्रे तृष्ण्णा\-चि\-कि\-त्सा\-प\-ञ्च\-च\-त्वा\-रिं\-श\-त्तमः कायचिकित्सायां}\ H}}\skt{ }\edtext{\skt{प्रथमो}}{\Afootnote{K; \skt{त्रे (दशमोऽध्यायः आदितः)}\ A; \skt{दशमोऽध्यायः ॥}\ H}}\skt{ }\edtext{\skt{दश ॥ {\symfont ❈}}}{\Afootnote{K; \skt{अष्ट\-च\-त्वा\-रिं\-शोऽ\-ध्यायः ॥४८ ॥}\ A; \emph{om.}\ H}}\skt{ ॥}\pend

\endnumbering
\beginnumbering
\eledsection*{Adhyāya 49}
\pstart \ledsidenote{\passref{6.49.1}}\edtext{\edtext{\skt{अथा\-त\-श्छ\-र्द्दि\-प्र\-ति\-षेधं}}{\lemma{\skt{अथातश्}}\Afootnote{A K; \skt{अथात}\ H}}}{\lemma{\skt{छर्द्दिप्रतिषेधं}}\Afootnote{A K; \skt{श्छर्द्दिप्रतिषेध}\ H}}\skt{ }\edtext{\skt{व्याख्यास्यामः}}{\Afootnote{\skt{अध्यायं व्याख्यास्यामः ॥}\ A; \emph{om.}\ K; \skt{मध्यायं व्याख्यास्यामः ॥}\ H}}\skt{ ॥}\pend

\pstart\versepar \ledsidenote{\passref{6.49.6}}\skt{च्छादयन्नाननं }\edtext{\edtext{\skt{वेगै\-र\-र्द\-य\-न्नं\-ग\-भ\-ञ्जनैः}}{\lemma{\skt{वेगैर्}}\Afootnote{K H; \skt{वेगै}\ A}}}{\lemma{\skt{अर्द\-य\-न्नं\-ग\-भ\-ञ्जनैः}}\Afootnote{K; \skt{रर्द\-य\-न्न\-ङ्ग\-भ\-ञ्जनैः ।}\ A; \skt{अर्द्द\-य\-न्न\-ङ्ग\-भ\-ञ्जनैः ।}\ H}}\skt{ । }\pend
\pstart\versepar \skt{निरुच्यते }\edtext{\edtext{\skt{च्छर्दिरिति}}{\lemma{\skt{च्छर्दिर्}}\Afootnote{A K; \skt{च्छर्द्दि}\ H}}}{\lemma{\skt{इति}}\Afootnote{A K; \skt{रिति}\ H}}\skt{ दोषो }\edtext{\skt{वक्त्रात्प्रधावितः}}{\Afootnote{K H; \skt{वक्त्रा\-द्वि\-नि\-श्च\-रन् ॥}\ A}}\skt{ ।}\pend

\pstart \ledsidenote{\passref{6.49.6a}}\edtext{\skt{ईर\-य\-ञ्च्छ्ले\-ष्म\-पित्तौ तु}}{\Afootnote{K; \skt{ईर॥ …यन्श्लेष्मपित्तौ तु}\ H}}\skt{ }\edtext{\skt{उदानो}}{\Afootnote{K; \skt{उदानौ}\ H}}\skt{ व्यापदं गतः ।}\pend

\pstart \ledsidenote{\passref{6.49.7}}\edtext{\edtext{\skt{ऊर्ध्वमागच्छति}}{\lemma{\skt{ऊर्ध्वम्}}\Afootnote{K H; \skt{दोषा\-नु\-दी\-र\-य\-न्वृ\-द्धा\-नु\-दानो व्यानसङ्गतः । ऊर्ध्व}\ A}}}{\lemma{\skt{आगच्छति}}\Afootnote{K H; \skt{मागच्छति}\ A}}\skt{ भृशं }\edtext{\skt{विरु\-द्धा\-हा\-र\-से\-विनः}}{\Afootnote{K H; \skt{विरु\-द्धा\-हा\-र\-से\-व\-नात् ॥}\ A}}\skt{ ।}\pend

\pstart\versepar \ledsidenote{\passref{6.49.8}}\skt{प्रसेको }\edtext{\skt{हृद\-यो\-त्क्ले\-दो\-भ\-क्त\-स्या\-न\-भि\-न\-न्दनं}}{\Afootnote{K; \skt{हृदयोत्क्लेशो भक्त\-स्या\-न\-भि\-न\-न्द\-नम् ।}\ A; \skt{हृदयोत्क्लेदो भक्त\-स्या\-ना\-भि\-न\-न्दनं ।}\ H}}\skt{ । }\pend
\pstart\versepar \skt{}\edtext{\skt{पूर्वरूपां}}{\Afootnote{K; \skt{पूर्वरूपं}\ A H}}\skt{ मतं च्छर्द्या यथास्वं }\edtext{\skt{चापि}}{\Afootnote{K H; \skt{च वि}\ A}}\skt{ }\edtext{\skt{निर्दिशेत्}}{\Afootnote{K H; \skt{भावयेत् ॥}\ A}}\skt{ ॥}\pend

\pstart\versepar \ledsidenote{\passref{6.49.9}}\edtext{\skt{यः}}{\Afootnote{K H; \skt{प्रच्छर्दये}\ A}}\skt{ }\edtext{\skt{फेनिलं}}{\Afootnote{K H; \skt{त्फेनि}\ A}}\skt{ }\edtext{\skt{च्छर्दयते}}{\Afootnote{K H; \skt{लम}\ A}}\skt{ ल्पमल्पं }\edtext{\skt{शूला\-र्दि\-तो\-भ्य\-र्दि\-त\-पा\-र्श्व\-पृष्ठः}}{\Afootnote{K; \skt{शूला\-र्दि\-तोऽ\-भ्य\-र्दि\-त\-पा\-र्श्व\-पृष्ठः}\ A H}}\skt{ । }\pend
\pstart\versepar \skt{श्रान्तः }\edtext{\skt{सणोषं}}{\Afootnote{K; \skt{सघोषं}\ A; \skt{सशोषं}\ H}}\skt{ बहुशः }\edtext{\skt{सलीनं}}{\Afootnote{K; \skt{कषायं}\ A; \skt{स}\ H}}\skt{ }\edtext{\skt{सावा\-त\-को\-प\-प्र\-भ\-वातु}}{\Afootnote{K; \skt{जीर्णेऽधिकं सानिलजा वमिस्तु ॥}\ A; \skt{नीलं सावा\-त\-को\-प\-प्र\-भ\-वातु}\ H}}\skt{ }\edtext{\skt{च्छर्दिः}}{\Afootnote{K H; \emph{om.}\ A}}\skt{ ।}\pend

\pstart\versepar \ledsidenote{\passref{6.49.10}}\edtext{\skt{व्योम्व्लं}}{\Afootnote{K; \skt{योऽम्लं}\ A; \skt{अम्लं}\ H}}\skt{ भृशं वा कटु\-ति\-क्त\-व\-क्त्रः\-पीतं सरक्तं हरितं वमेद्वा । }\pend
\pstart\versepar \skt{}\edtext{\skt{सदा\-ह\-तो\-द\-ज्व\-र\-व\-क्त्र\-शोषं}}{\Afootnote{K H; \skt{सदा\-ह\-चो\-ष\-ज्व\-र\-व\-क्त्र\-शोषो}\ A}}\skt{ }\edtext{\skt{सापि\-त्त\-को\-प\-प्र\-भवा}}{\Afootnote{K H; \skt{मूर्च्छान्वितः पित्तनिमित्तजा सा ॥}\ A}}\skt{ }\edtext{\skt{मता}}{\Afootnote{K H; \emph{om.}\ A}}\skt{ }\edtext{\skt{तु}}{\Afootnote{K H; \emph{om.}\ A}}\skt{ ।}\pend

\pstart\versepar \ledsidenote{\passref{6.49.11}}\skt{यो हृष्टरोमा मधुरं प्रभूतं शुक्लं हिमं सांद्र\-क\-फा\-नु\-विद्धं । }\pend
\pstart\versepar \skt{}\edtext{\skt{अभ\-क्त\-रु\-ग्गौ\-र\-व\-सा\-द\-युक्तं}}{\Afootnote{K H; \skt{अभ\-क्त\-रु\-ग्गौ\-र\-व\-सा\-द\-युक्तो}\ A}}\skt{ वमेद्वमी सा कफकोपजा }\edtext{\skt{तु}}{\Afootnote{K H; \skt{स्यात् ॥}\ A}}\skt{ ॥}\pend

\pstart\versepar \ledsidenote{\passref{6.49.12}}\skt{सर्वाणि }\edtext{\skt{लिङ्गानि}}{\Afootnote{K H; \skt{रूपाणि}\ A}}\skt{ भवन्ति }\edtext{\skt{यस्याः सा}}{\Afootnote{K H; \skt{यस्यां सा}\ A}}\skt{ सर्व\-दो\-ष\-प्र\-भवा मता तु ॥ }\pend
\pstart\versepar \skt{बीभत्सजा }\edtext{\skt{दौहृ\-द\-ज\-न्म\-जात्व}}{\Afootnote{K H; \skt{दौहृदजामजा च}\ A}}\skt{ सात्म्य\-प्र\-को\-पा\-त्क्रि\-मिजा च या हि । }\pend
\pstart\versepar \skt{सा पञ्चमी ताञ्च }\edtext{\skt{विभावयीत}}{\Afootnote{K; \skt{विभावयेत्तु}\ A; \skt{निभावयेत}\ H}}\skt{ दोषोच्छ्रयेणैव यथोक्तमादौ ॥}\edtext{}{\lemma{\skt{सात्म्य\-प्र\-को\-पा\-त्क्रि\-मिजा च या हि । सा}}\Afootnote{K; \skt{सात्म्य\-प्र\-को\-पा\-त्कृ\-मिजा च या हि । सा}\ A; \skt{सात्म्यु प्रकोपात्क्रिमिजा च या हि ॥ सा}\ H}}\pend

\pstart\versepar \ledsidenote{\passref{6.49.13}}\skt{शूल\-हृ\-ल्ला\-स\-ब\-हुला }\edtext{\skt{क्रिमिजा तु}}{\Afootnote{K H; \skt{कृमिजा च}\ A}}\skt{ विशेषतः । }\pend
\pstart\versepar \skt{}\edtext{\skt{क्रिमि\-हृ\-द्रो\-ग\-तु\-ल्येन}}{\Afootnote{K; \skt{कृमि\-हृ\-द्रो\-ग\-तु\-ल्येन}\ A H}}\skt{ }\edtext{\skt{लक्षणेन च}}{\Afootnote{A K; \skt{लक्षणैन च}\ H}}\skt{ }\edtext{\skt{लक्षिता}}{\Afootnote{A K; \skt{लक्षिताः ॥}\ H}}\skt{ ॥}\pend

\pstart\versepar \ledsidenote{\passref{6.49.14}}\edtext{\skt{क्षीण\-स्यो\-प\-द्र\-व\-वतीं}}{\Afootnote{K; \skt{क्षीण\-स्यो\-प\-द्र\-वै\-र्युक्तां}\ A; \skt{क्षीण\-स्यो\-प\-द्र\-वतीं}\ H}}\skt{ }\edtext{\skt{सासृक्प्रायां}}{\Afootnote{K H; \skt{सासृक्पूयां}\ A}}\skt{ सचन्द्रिकां । }\pend
\pstart\versepar \skt{}\edtext{\skt{च्छर्दिप्रसक्तां}}{\Afootnote{K; \skt{छर्दिं प्रसक्तां}\ A H}}\skt{ कुशलो नारभेत चिकित्सितुं ॥}\pend

\pstart\versepar \ledsidenote{\passref{6.49.16}}\skt{वमीषु बहुदोषासु च्छर्दनं }\edtext{\skt{हितमुच्यते}}{\Afootnote{K H; \skt{हितमुच्यते ।}\ A}}\skt{ ॥ }\pend
\pstart\versepar \skt{विरेचनम्वा }\edtext{\skt{युञ्जीत}}{\Afootnote{K H; \skt{कुर्वीत}\ A}}\skt{ }\edtext{\skt{यथोदोषोच्छ्रयं}}{\Afootnote{K; \skt{यथादोषोच्छ्रयं}\ A H}}\skt{ भिषक् ।}\pend

\pstart\versepar \ledsidenote{\passref{6.49.17}}\edtext{\skt{संस\-र्गा\-श्चा\-नु\-पू\-र्वेण}}{\lemma{\skt{संसर्गाश्}}\Afootnote{K; \skt{संसर्गश्}\ A H}}\skt{ यथास्वं }\edtext{\skt{भेषजायुतं}}{\Afootnote{K H; \skt{भेषजायुतः ।}\ A}}\skt{ । }\pend
\pstart\versepar \skt{}\edtext{\skt{लघूणि}}{\Afootnote{K H; \skt{लघूनि}\ A}}\skt{ परिशुष्काणि }\edtext{\skt{सात्म्यान्यन्यानि}}{\Afootnote{K H; \skt{सात्म्यान्यन्नानि}\ A}}\skt{ }\edtext{\skt{चाचरेत्}}{\Afootnote{A K; \skt{चापरेत् ।}\ H}}\skt{ ।}\pend

\pstart\versepar \ledsidenote{\passref{6.49.18}}\skt{यथास्वञ्च कषायाणि ज्वरघ्नानि प्रयोजयेत् । }\pend
\pstart\versepar \skt{हन्यात्क्षीरोदकं पीतं च्छर्दिं }\edtext{\skt{पवनसंभ}}{\Afootnote{A K; \skt{पवनसम्भ}\ H}}\skt{ वां ।}\edtext{}{\lemma{\skt{प्रयोजयेत् । हन्यात्क्षीरोदकं}}\Afootnote{K H; \skt{प्रयोजयेत् । हन्यात्क्षीरघृतं}\ A}}\pend

\pstart\versepar \ledsidenote{\passref{6.49.19}}\edtext{\skt{मुद्गा\-म\-ल\-क\-यूषो वा}}{\Afootnote{K; \skt{ससैन्धवं पिबे\-त्स\-र्पि\-र्वा\-त\-च्छ\-र्दि\-नि\-वा\-र\-णम् । मुद्गा\-म\-ल\-क\-यूषो वा}\ A; \skt{…मूलकयूषो वा}\ H}}\skt{ ससर्पिष्कः ससैन्धवः । }\pend
\pstart\versepar \skt{यवागूं मधुमिश्राम्वा }\edtext{\skt{पञ्चमूलकृतां}}{\Afootnote{K H; \skt{पञ्चमूलीकृतां}\ A}}\skt{ पिबेत् ।}\pend

\pstart \ledsidenote{\passref{6.49.20}}\edtext{\skt{फलाम्ब्लं}}{\Afootnote{K; \emph{om.}\ A; \skt{फलाम्लं}\ H}}\skt{ }\edtext{\skt{विष्किररसं}}{\Afootnote{K H; \emph{om.}\ A}}\skt{ पिबेद्वा }\edtext{\skt{व्यक्तसैन्धवं}}{\Afootnote{K H; \skt{व्यक्तसिन्धूत्थं फलाम्लं वैष्किरं रसम् । सुखोष्णलवणं चात्र हितं स्नेह\-वि\-रे\-च\-नम् ॥}\ A}}\skt{ ।}\pend

\pstart\versepar \ledsidenote{\passref{6.49.21}}\skt{पित्तो\-प\-श\-म\-नी\-यानि पाक्यानि }\edtext{\skt{च}}{\Afootnote{K H; \skt{शि}\ A}}\skt{ }\edtext{\skt{हिमानि च}}{\Afootnote{K H; \skt{शिराणि च ।}\ A}}\skt{ । }\pend
\pstart\versepar \skt{}\edtext{\skt{कषा\-या\-न्यु\-प\-यु\-क्तानि}}{\Afootnote{K; \skt{कषा\-या\-ण्यु\-प\-यु\-क्तानि}\ A H}}\skt{ घ्नन्ति पित्तकृतां वमीं ।}\pend

\pstart\versepar \ledsidenote{\passref{6.49.22}}\edtext{\skt{शोधणं}}{\Afootnote{K; \skt{शोधनं}\ A H}}\skt{ }\edtext{\skt{मधुरञ्चात्र}}{\Afootnote{K H; \skt{मधुरैश्चात्र}\ A}}\skt{ }\edtext{\skt{द्राक्षा\-र\-स\-स\-मा\-युतं}}{\Afootnote{K H; \skt{द्राक्षा\-र\-स\-स\-मा\-युतैः ।}\ A}}\skt{ । }\pend
\pstart\versepar \skt{बलवत्यां प्रशंसन्ति }\edtext{\skt{सर्पि\-स्तै\-ल्व\-क\-मेव}}{\lemma{\skt{सर्पिस्तैल्वकम्}}\Afootnote{K; \skt{सर्पिस्तैल्वकम्}\ A; \skt{सर्प्पिःमैल्वकम्}\ H\textsuperscript{ac}; \skt{सर्प्पि\-स्तैः\-मै\-ल्व\-कम्}\ H\textsuperscript{pc}}}\skt{ }\edtext{\skt{वा}}{\Afootnote{K H; \skt{च ॥}\ A}}\skt{ ।}\pend

\pstart\versepar \ledsidenote{\passref{6.49.23}}\edtext{\skt{आर\-ग्व\-धा\-दि\-नि\-र्यूहं}}{\Afootnote{A K; \skt{आर\-ग्व\-धा\-दि\-नि\-र्युहं}\ H}}\skt{ }\edtext{\skt{दशांगं}}{\Afootnote{K; \skt{दशाङ्ग}\ A H}}\skt{ योगमेव वा । }\pend
\pstart\versepar \skt{}\edtext{\skt{पाय\-ये\-त्म\-धु\-सं\-युक्तं}}{\Afootnote{K H; \skt{पाययेताथ सक्षौद्रं}\ A}}\skt{ }\edtext{\skt{कफजायां}}{\Afootnote{A K; \skt{कफजायं}\ H}}\skt{ चिकित्सकः ।}\pend

\pstart\versepar \ledsidenote{\passref{6.49.24}}\skt{कृतं गुडूच्या विधिवत्कषायं हिमसञ्ज्ञितं । }\pend
\pstart\versepar \skt{तिसृष्वपि भवेत्पथ्यं }\edtext{\skt{माक्षिकेण}}{\Afootnote{A K; \skt{माक्षिकेन}\ H}}\skt{ }\edtext{\skt{समायुतं}}{\Afootnote{K H; \skt{समन्वितम् ॥}\ A}}\skt{ ।}\pend

\pstart\versepar \ledsidenote{\passref{6.49.25}}\edtext{\skt{बीभत्सजां}}{\Afootnote{A K; \skt{बीभत्स्यजां}\ H}}\skt{ }\edtext{\skt{हृद्य\-त\-मै\-र्दो\-ह्यदी}}{\Afootnote{K; \skt{हृद्य\-त\-मै\-र्दौ\-हृदीं}\ A; \skt{हृद्य\-त\-मैः\-र्द्दौ\-हृदी}\ H}}\skt{ }\edtext{\skt{कांक्षितैः}}{\Afootnote{A K; \skt{काक्षितैः}\ H}}\skt{ फलैः । }\pend
\pstart\versepar \skt{}\edtext{\edtext{\edtext{\skt{लंघ\-नै\-र्व\-म\-नै\-श्चासां}}{\lemma{\skt{लंघनैर्}}\Afootnote{K H; \skt{लङ्घनै}\ A}}}{\lemma{\skt{वमनैश्}}\Afootnote{K H; \skt{र्वमनैश्}\ A}}}{\lemma{\skt{चासां}}\Afootnote{K H; \skt{चामां}\ A}}\skt{ }\edtext{\skt{सात्म्यैर्वा}}{\lemma{\skt{वा}}\Afootnote{K H; \skt{सा}\ A}}\skt{ }\edtext{\skt{सात्म्यकोपजां}}{\Afootnote{K H; \skt{त्म्यप्रकोपजाम् ॥}\ A}}\skt{ ।}\pend

\pstart\versepar \ledsidenote{\passref{6.49.26}}\edtext{\skt{क्रिमि\-हृ\-द्रो\-ग\-व\-च्चापि}}{\Afootnote{K; \skt{कृमि\-हृ\-द्रो\-ग\-व\-च्चापि}\ A H}}\skt{ }\edtext{\skt{क्रिमिजं}}{\Afootnote{K; \skt{कृमिजां}\ A H}}\skt{ }\edtext{\skt{शोधयेद्वमीं}}{\Afootnote{K; \skt{साधयेद्वमीम्}\ A H}}\skt{ । }\pend
\pstart\versepar \skt{}\edtext{\skt{यथादोषाञ्च}}{\Afootnote{K; \skt{वितरेच्च यथादोषं}\ A; \skt{यथादोषञ्च}\ H}}\skt{ }\edtext{\skt{वितरेच्छस्तं}}{\Afootnote{K H; \skt{शस्तं}\ A}}\skt{ }\edtext{\skt{विधिमनन्तरं}}{\Afootnote{A K; \skt{विधिमनन्तरं ॥}\ H}}\skt{ ।}\pend

\pstart\versepar \ledsidenote{\passref{6.49.27}}\edtext{\skt{दधि\-त्थ\-र\-स\-सं\-युक्ताः}}{\Afootnote{K H; \skt{दधि\-त्थ\-र\-स\-स्ं\-सक्तां}\ A}}\skt{ }\edtext{\skt{पिप्प\-ली\-मा\-क्षि\-का\-न्विताः}}{\Afootnote{K; \skt{पिप्पलीं माक्षिकान्विताम् ।}\ A; \skt{पिप्पलीं मा क्षिकान्विता ।}\ H}}\skt{ । }\pend
\pstart\versepar \skt{}\edtext{\skt{मुहुर्मुहुर्नरो}}{\Afootnote{A K; \skt{मुहुर्म्मुहुर्न्नरो}\ H}}\skt{ लीढ्वा च्छर्दिभ्यः }\edtext{\skt{परिमुच्यते}}{\Afootnote{K H; \skt{प्रवि\-मु\-च्य\-ते(\-प्र\-ति\-मु\-च्यते) ॥}\ A}}\skt{ ॥}\pend

\pstart\versepar \ledsidenote{\passref{6.49.28}}\skt{समाक्षिका }\edtext{\skt{मधुरसा}}{\Afootnote{A K; \skt{मुस्तरसा}\ H}}\skt{ पीता वा तण्डुलाम्बुना । }\pend
\pstart\versepar \skt{}\edtext{\skt{तर्प्पणा वा}}{\Afootnote{K H; \skt{तर्पणं वा}\ A}}\skt{ }\edtext{\skt{मधुयुता}}{\Afootnote{K; \skt{मधुयुतं}\ A; \skt{मधुयुतास्}\ H}}\skt{ तिसृणामपि }\edtext{\skt{भेषजां}}{\Afootnote{K; \skt{भेषजम्}\ A H}}\skt{ ।}\pend

\pstart\versepar \ledsidenote{\passref{6.49.29}}\skt{स्वयं गुप्तां सयष्ट्याह्वां }\edtext{\skt{ताण्डू\-ला\-म्बु\-म\-धु\-द्रवां}}{\Afootnote{K; \skt{तण्डु\-ला\-म्बु\-म\-धु\-द्र\-वाम्}\ A H}}\skt{ । }\pend
\pstart\versepar \skt{पिबे\-द्य\-वा\-गू\-म\-थवा }\edtext{\skt{सिद्धां}}{\Afootnote{A K; \skt{सूतं}\ H}}\skt{ पत्रैः }\edtext{\skt{करञ्जजैः}}{\Afootnote{A K; \skt{करञ्जकैः ॥}\ H}}\skt{ ।}\pend

\pstart\versepar \ledsidenote{\passref{6.49.30}}\edtext{\skt{युक्ताम्व्लवणाः}}{\Afootnote{K; \skt{युक्ताम्ललवणाः}\ A H}}\skt{ पिष्टाः }\edtext{\skt{कुस्तुम्बुर्योथवा}}{\Afootnote{K; \skt{कुस्तुम्बुर्योऽथवा}\ A; \skt{कुस्तुम्बरो थवा}\ H}}\skt{ हिताः । }\pend
\pstart\versepar \skt{तण्डुलाम्बुयुतं }\edtext{\skt{खादेत्कपित्थं}}{\Afootnote{A K; \skt{खादेत्कपित्थ}\ H}}\skt{ }\edtext{\skt{त्र्यूषणेन}}{\Afootnote{A K; \skt{षणेन}\ H}}\skt{ वा ।}\pend

\pstart \ledsidenote{\passref{6.49.31}}\skt{सिता }\edtext{\skt{चन्दनमध्वक्तां}}{\Afootnote{K; \skt{चन्दनमध्वाक्तं}\ A H}}\skt{ }\edtext{\skt{लिह्याद्वा}}{\Afootnote{K H; \skt{लिह्याद्वा}\ A}}\skt{ }\edtext{\skt{मक्षिकाशकृत्}}{\Afootnote{K H; \skt{मक्षिकाशकृत् । पिबे\-त्प\-योऽ\-ग्नि\-तप्तं च निर्वाप्य गृहगोधिकाम् ॥}\ A}}\skt{ ।}\pend

\pstart \ledsidenote{\passref{6.49.32}}\edtext{\skt{सर्पिक्षौद्रसिता}}{\Afootnote{K; \skt{सर्पिःक्षौद्रयुता}\ A; \skt{सर्प्पिःक्षौद्रसिता}\ H}}\skt{ }\edtext{\skt{लाजा}}{\Afootnote{K; \skt{न्वापि लाज}\ A; \skt{लाजाः}\ H}}\skt{ }\edtext{\skt{शक्तुं}}{\Afootnote{K; \skt{शक्तू}\ A H}}\skt{ }\edtext{\skt{लिह्यात्तथापि}}{\Afootnote{K; \skt{न्पिबेत्तथा । स}\ A; \skt{न्लिह्यात्तथा}\ H}}\skt{ }\edtext{\skt{वा}}{\Afootnote{K; \skt{र्पिःक्षौ\-द्र\-सि\-तो\-पेतां मागधीं वा लिहेत्तथा ॥}\ A; \skt{मिवा ।}\ H}}\skt{ ।}\pend

\pstart\versepar \ledsidenote{\passref{6.49.33}}\edtext{\skt{धात्री\-र\-सै\-श्च\-न्द\-नम्वा}}{\Afootnote{K; \skt{धात्रीरसे चन्दनं वा}\ A; \skt{धात्री\-र\-सै\-श्च\-न्द\-नम्वा}\ H}}\skt{ }\edtext{\skt{श्रितां}}{\Afootnote{K; \skt{घृष्टं}\ A; \skt{श्रित}\ H}}\skt{ }\edtext{\skt{मुद्गदलाम्बुना}}{\Afootnote{A K; \skt{मुद्गदलाम्बुवा ॥}\ H}}\skt{ । }\pend
\pstart\versepar \skt{}\edtext{\skt{कोला\-म\-ल\-क\-म\-ज्जानो}}{\Afootnote{K; \skt{कोला\-म\-ल\-क\-म\-ज्जानं}\ A; \skt{कोम\-ला\-म\-ल\-क\-म\-ज्जानो}\ H}}\skt{ }\edtext{\skt{लिहेद्वापि}}{\Afootnote{K H; \skt{लिह्याद्वा त्रि}\ A}}\skt{ }\edtext{\skt{त्रिवर्ण्णकं}}{\Afootnote{K; \skt{सुगन्धिकम् ॥}\ A; \skt{त्रिवर्ण्णक्ं ।}\ H}}\skt{ ।}\pend

\pstart\versepar \ledsidenote{\passref{6.49.34}}\edtext{\skt{सक्षौद्रं}}{\Afootnote{K H; \skt{सक्षौद्रां}\ A}}\skt{ शालिलाजानां }\edtext{\skt{यवागू वा}}{\Afootnote{K; \skt{यवगूं वा}\ A; \skt{यवागूम्वा}\ H}}\skt{ पिबेन्नरः । }\pend
\pstart\versepar \skt{}\edtext{\skt{घ्रेया\-ण्यु\-प\-क्षि\-पे\-च्चापि}}{\Afootnote{K; \skt{घ्रेया\-ण्यु\-प\-ह\-रे\-च्चापि}\ A; \skt{घ्रेया\-ण्यु\-प\-क्षि\-पे\-च्चापि}\ H}}\skt{ }\edtext{\skt{सुगन्धीनि}}{\Afootnote{K H; \skt{मनो\-घ्रा\-ण\-सु\-खानि}\ A}}\skt{ }\edtext{\skt{चिकित्सकः}}{\Afootnote{K H; \skt{च ॥}\ A}}\skt{ ।}\pend

\pstart\versepar \ledsidenote{\passref{6.49.35}}\skt{जांगलानि च }\edtext{\skt{शूल्यानि}}{\Afootnote{K H; \skt{मांसानि}\ A}}\skt{ }\edtext{\skt{षाड\-वा\-ले\-ह्य\-पा\-नकाः}}{\Afootnote{K; \skt{शुभानि पानकानि च ।}\ A; \skt{षाड\-वा\-ने\-ह्य\-पा\-नकाः ॥}\ H}}\skt{ । }\pend
\pstart\versepar \skt{भोजनानि }\edtext{\skt{च}}{\Afootnote{K H; \skt{वि}\ A}}\skt{ चित्राणि }\edtext{\skt{कुर्या\-त्स\-र्व\-स्व\-त\-न्द्रितः}}{\Afootnote{K; \skt{कुर्या\-त्स\-र्वा\-स्व\-त\-न्द्रितः ॥}\ A; \skt{कुर्या\-त्स\-र्व्वा\-स्व\-त\-न्द्रितः ।}\ H}}\skt{ ।}\pend

\pstart \ledsidenote{\passref{6.49.35a}}\skt{अप्रियाण्यपि सेवेत जले\-नो\-द्वे\-ज\-ये\-द\-पीति ॥}\pend

\pstart \ledsidenote{\passref{6.49.trailer}}\edtext{\skt{कायचि}}{\Afootnote{K; \skt{इति सुश्रु\-त\-स्ं\-हि\-ता\-या\-मु\-त्त\-र\-त\-न्त्रा\-न्त\-र्गते कायचि}\ A; \skt{उत्तरतन्त्रे च्छर्द्दिचिकित्सा षट्चत्वारिं त्तमः कायचि}\ H}}\skt{ ॥ }\edtext{\skt{ल}}{\Afootnote{K; \skt{कित्सातन्त्रे छर्दिप्रतिषेधो नाम (ए\-का\-द\-शोऽ\-ध्यायः, आदितः) एको\-न\-प\-ञ्चा\-श\-त्त\-मोऽ\-ध्यायः ॥४९ ॥}\ A; \skt{कित्सा\-या\-मे\-क\-दशो ध्यायः ॥}\ H}}\pend

\endnumbering
\beginnumbering
\eledsection*{Adhyāya 50}
\pstart \ledsidenote{\passref{6.50.1}}\skt{अथातो }\edtext{\skt{हिक्काप्रतिषेधं}}{\Afootnote{A K; \skt{हिक्का\-प्र\-ति\-षे\-ध\-म\-ध्यायं}\ H}}\skt{ }\edtext{\skt{व्याख्यास्यामः}}{\Afootnote{A H; \emph{om.}\ K}}\skt{ ॥}\pend

\pstart\versepar \ledsidenote{\passref{6.50.3}}\edtext{\skt{विदा\-हि\-गु\-रु\-वि\-ष\-म्भि\-रू\-क्षा\-भि\-ष्य\-न्दि\-भो\-जनैः}}{\lemma{\skt{विदा\-हि\-गु\-रु\-वि\-ष\-म्भिर्}}\Afootnote{K; \skt{विदा\-हि\-गु\-रु\-वि\-ष्ट\-म्भिर्}\ A H}}\skt{ । }\pend
\pstart\versepar \skt{}\edtext{\skt{शीत\-या\-ना\-स\-न\-स्था\-न\-ऱजो}}{\Afootnote{K; \skt{शीत\-पा\-ना\-स\-न\-स्था\-न\-रजो}\ A; \skt{शीतया नासनस्थानां रजो}\ H}}\skt{ }\edtext{\skt{धूमानलानिलैः}}{\Afootnote{K H; \skt{धूमानिलानलैः ॥}\ A}}\skt{ ।}\pend

\pstart \ledsidenote{\passref{6.50.4}}\edtext{\skt{व्याया\-म\-क\-र्म\-भा\-रा\-ध्व\-वे\-गा\-घा\-ता\-प\-त\-र्पणैः}}{\Afootnote{K H; \skt{व्याया\-म\-क\-र्म\-भा\-रा\-ध्व\-वे\-गा\-घा\-ता\-प\-त\-र्पणैः । आम\-दो\-षा\-भि\-घा\-त\-स्त्री\-क्ष\-य\-दो\-ष\-प्र\-पी\-डनैः ॥}\ A}}\skt{ ।}\pend

\pstart \ledsidenote{\passref{6.50.5}}\edtext{\skt{हिक्काश्वासश्च}}{\Afootnote{K H; \skt{विष\-मा\-श\-ना\-ध्य\-न\-श\-नै\-स्तथा समशनैरपि । हिक्का श्वासश्च}\ A}}\skt{ कासश्च }\edtext{\skt{नृणां}}{\Afootnote{A K; \skt{नॄणा}\ H}}\skt{ }\edtext{\skt{समुपजायते}}{\Afootnote{A K; \skt{समुपजायते ॥}\ H}}\skt{ ।}\pend

\pstart\versepar \ledsidenote{\passref{6.50.6}}\edtext{\edtext{\skt{मुहु\-र्मु\-हु\-र्वा\-यु\-रु\-देति}}{\lemma{\skt{मुहुर्मुहुर्}}\Afootnote{K H; \skt{मुहुर्मुहु}\ A}}}{\lemma{\skt{वायुरुदेति}}\Afootnote{K; \skt{र्वायुरुदेति}\ A; \skt{वायुरुदेति}\ H}}\skt{ सस्वनो यकृत्प्लिहान्त्राणि }\edtext{\skt{मुखे}}{\Afootnote{K H; \skt{मुखा}\ A}}\skt{ समुत्क्षिपन् । }\pend
\pstart\versepar \skt{स }\edtext{\skt{घोषवानासु}}{\Afootnote{K; \skt{घोषवानाशु}\ A H}}\skt{ }\edtext{\skt{हिन}}{\Afootnote{K; \skt{हिन\-स्त्य\-सू\-न्य\-त\-स्त\-तस्तु हिक्केति भिषग्भिरुच्यते ॥}\ A; \skt{हिनस्ति यस्मात्ततस्तु हिक्केत्य भिधीयते बुधैः}\ H\textsuperscript{ac} H\textsuperscript{pc}}}\edtext{}{\lemma{\skt{समुत्क्षिपन् । स}}\Afootnote{K H; \skt{दिवाक्षिपन् । स}\ A}}\pend

\pstart\versepar \ledsidenote{\passref{6.50.7}}\edtext{\skt{अन्नजां}}{\Afootnote{A K; \skt{अनुजां}\ H}}\skt{ }\edtext{\skt{यमलां}}{\Afootnote{A K; \skt{यमला}\ H}}\skt{ क्षुद्रां }\edtext{\skt{गम्भीरां}}{\Afootnote{A K; \skt{गम्भीरा}\ H}}\skt{ महतीन्तथा । }\pend
\pstart\versepar \skt{वायुः कफेनानुगतः }\edtext{\skt{पञ्चहिक्का}}{\Afootnote{K H; \skt{पञ्च हिक्काः}\ A}}\skt{ करोति }\edtext{\skt{ह}}{\Afootnote{K H; \skt{हि ॥}\ A}}\skt{ ।}\pend

\pstart \ledsidenote{\passref{6.50.9}}\edtext{\edtext{\skt{वायु\-र\-न्नै\-र\-व\-स्त्तीर्णः}}{\lemma{\skt{वायुरन्नैर्}}\Afootnote{K; \skt{त्वरमाणस्य चाहारं भुञ्जानस्याथवा घनम् । वायुरन्नै}\ A; \skt{वायुरन्नै}\ H}}}{\lemma{\skt{अवस्त्तीर्णः}}\Afootnote{K; \skt{रवस्तीर्णः}\ A; \skt{र…जीर्ण्णः}\ H}}\skt{ }\edtext{\edtext{\skt{कटुकैरर्दितो}}{\lemma{\skt{कटुकैर्}}\Afootnote{K; \skt{कटुकै}\ A H}}}{\lemma{\skt{अर्दितो}}\Afootnote{K; \skt{रर्दितो}\ A H}}\skt{ भृशं ।}\pend

\pstart\versepar \ledsidenote{\passref{6.50.10}}\edtext{\skt{हिक्कयत्यूर्ध्वगो}}{\Afootnote{A K; \skt{हिक्कायत्यूर्ध्वगो}\ H}}\skt{ भूत्वा तां }\edtext{\skt{विद्यादन्नजां}}{\Afootnote{K; \skt{विद्यादन्नजां}\ A H}}\skt{ भिषक् । }\pend
\pstart\versepar \skt{चिरेण }\edtext{\edtext{\skt{यमलैर्वेगैर्या}}{\lemma{\skt{यमलैर्}}\Afootnote{K; \skt{यमलै}\ A H}}}{\lemma{\skt{वेगैर्या}}\Afootnote{K; \skt{र्वेगैर्या}\ A H}}\skt{ हिक्का सम्प्रवर्त्तते ।}\pend

\pstart\versepar \ledsidenote{\passref{6.50.11}}\skt{कम्पयन्ती शिरोग्रीवं यमलां तां विनिर्दिशेत् ॥ }\pend
\pstart\versepar \skt{विकृष्टकालैर्या }\edtext{\edtext{\skt{वेगैर्महद्भिः}}{\lemma{\skt{वेगैर्}}\Afootnote{K; \skt{वेगै}\ A H}}}{\lemma{\skt{महद्भिः}}\Afootnote{K; \skt{र्मन्दैः स}\ A; \skt{र्म्महद्भिः}\ H}}\skt{ }\edtext{\skt{सम्प्रवर्त्तते}}{\Afootnote{K H; \skt{मभिवर्तते ॥}\ A}}\skt{ ।}\edtext{}{\lemma{\skt{विनिर्दिशेत् ॥ विकृष्टकालैर्या}}\Afootnote{A K; \skt{विविर्द्दिशेत् ॥ विकृष्णकालैर्या}\ H}}\pend

\pstart\versepar \ledsidenote{\passref{6.50.12}}\skt{क्षुद्रिका नाम सा }\edtext{\skt{हिक्का\-ज\-त्रु\-मू\-ला\-त्प्र\-धा\-विता}}{\Afootnote{A K; \skt{हिक्का}\ H}}\skt{ ॥ }\pend
\pstart\versepar \skt{}\edtext{\skt{नाभिप्रवृत्ता या}}{\Afootnote{A K; \emph{om.}\ H}}\skt{ }\edtext{\skt{हिक्का}}{\Afootnote{A K; \skt{… । …}\ H}}\skt{ }\edtext{\skt{घोरा\-ग\-म्भी\-र\-ना\-दिनी}}{\Afootnote{A K; \skt{घोरा गंभीरनादिनी ॥}\ H}}\skt{ ।}\pend

\pstart \ledsidenote{\passref{6.50.13}}\edtext{\skt{अने\-को\-प\-द्र\-व\-वती}}{\Afootnote{K H\textsuperscript{pc}; \skt{शुष्कौ\-ष्ठ\-क\-ण्ठ\-जि\-ह्वा\-स्य\-श्वा\-स\-पा\-र्श्व\-रु\-जा\-करी । अने\-को\-प\-द्र\-व\-युता}\ A; \skt{अने\-को\-प\-द्र\-व\-व\-आती}\ H\textsuperscript{ac}}}\skt{ गम्भीरा नाम सा स्मृता ॥}\pend

\pstart\versepar \ledsidenote{\passref{6.50.14}}\edtext{\skt{मर्मा\-ण्या\-पी\-ड\-य\-न्तीव}}{\Afootnote{A K; \skt{मर्म्मा\-ण्या\-पी\-ड\-यन्ती च}\ H}}\skt{ सततं या प्रवर्त्तते । }\pend
\pstart\versepar \skt{}\edtext{\skt{महाहिक्केति सा}}{\Afootnote{K H; \skt{देहमायम्य वेगेन घोष\-य\-त्य\-ति\-तृ\-ष्यतः । महाहिक्केति सा}\ A}}\skt{ ज्ञेया }\edtext{\skt{सर्व\-गा\-त्र\-वि\-क\-म्पनी}}{\Afootnote{K H; \skt{सर्व\-गा\-त्र\-प्र\-क\-म्पिनी ॥}\ A}}\skt{ ॥}\pend

\pstart\versepar \ledsidenote{\passref{6.50.15}}\skt{आयम्यते }\edtext{\skt{हिक्कमानस्य}}{\Afootnote{K; \skt{हिक्कतोऽङ्गानि य}\ A; \skt{हिक्कामानस्य}\ H}}\skt{ }\edtext{\skt{देहे}}{\Afootnote{K; \skt{स्य}\ A; \skt{देहो}\ H}}\skt{ }\edtext{\skt{दृष्टिश्चोर्ध्वं}}{\Afootnote{K; \skt{दृष्टिश्चोर्ध्वं}\ A H}}\skt{ }\edtext{\skt{नाम्यते}}{\Afootnote{K H; \skt{ताम्यते}\ A}}\skt{ यस्य }\edtext{\skt{चापि}}{\Afootnote{K H; \skt{गाढम् ।}\ A}}\skt{ । }\pend
\pstart\versepar \skt{}\edtext{\skt{क्षीणो\-न्न\-द्वि\-द्वि\-क्षिपं यश्च}}{\Afootnote{K; \skt{क्षीणोऽ\-न्न\-द्वि\-ट्का\-सते यश्च}\ A; \skt{क्षीणो\-न्न\-द्वि\-द्वि\-क्षि\-प\-न्यश्च}\ H}}\skt{ }\edtext{\skt{देहं तौ}}{\Afootnote{K; \skt{हिक्की तौ}\ A; \skt{देहं तौ द्व्}\ H}}\skt{ }\edtext{\skt{द्वौ}}{\Afootnote{K; \skt{द्वा}\ A H}}\skt{ }\edtext{\skt{चान्त्यौ}}{\Afootnote{K H; \skt{वन्त्यौ}\ A}}\skt{ }\edtext{\skt{वर्ज\-ये\-त्क्ला\-म्य\-मानौ}}{\Afootnote{K H; \skt{वर्ज\-ये\-द्धि\-क्क\-मानौ ॥}\ A}}\skt{ ।}\pend

\pstart \ledsidenote{\passref{6.50.16}}\edtext{\skt{प्राणा\-या\-मो\-द्वे\-ज\-नो\-त्रा\-स\-नानि}}{\Afootnote{K H; \skt{प्राणा\-या\-मो\-द्वे\-ज\-न\-त्रा\-स\-नानि}\ A}}\skt{ }\edtext{\edtext{\skt{पिपी\-लि\-कै\-र्दा\-श\-न\-ञ्चात्र}}{\lemma{\skt{पिपीलिकैर्}}\Afootnote{K; \skt{सूचीतोदैः}\ A; \skt{पिपीलिकै}\ H}}}{\lemma{\skt{दाशनञ्}}\Afootnote{K; \skt{संभ्रमश्}\ A; \skt{र्द्दाशनं}\ H}}\skt{ }\edtext{\skt{शस्तं}}{\Afootnote{K H; \skt{शस्तः । यष्ट्याह्वं वा माक्षि\-के\-णा\-व\-पीडे पिप्पल्यो वा शर्क\-रा\-चू\-र्ण\-युक्ताः ॥}\ A}}\skt{ ।}\pend

\pstart \ledsidenote{\passref{6.50.16a}}\edtext{\skt{सश\-र्क\-र\-म्म\-धु\-क\-ञ्चा\-व\-पी\-डे\-त्क्षौ\-द्रा\-न्वि\-ता\-मा\-ग\-धि\-का\-स्त\-थैव}}{\lemma{\skt{मधु\-क\-ञ्चा\-व\-पी\-डे\-त्क्षौ\-द्रा\-न्वि\-ता\-मा\-ग\-धि\-कास्}}\Afootnote{K; \skt{मधुकं चावपीडे क्षौद्रायुता मागधिका}\ H}}\skt{ ।}\pend

\pstart \ledsidenote{\passref{6.50.16b}}\edtext{\skt{ईक्षोरसाः}}{\Afootnote{K; \skt{इक्षोरसः}\ H}}\skt{ क्षीरमुष्णं जलम्वा । बलान्विते च्छर्दनम्वा प्रशस्तं ।}\pend

\pstart \ledsidenote{\passref{6.50.16c}}\skt{धूमं }\edtext{\skt{पिबे\-त्स\-र्ज\-र\-स्यस्य}}{\Afootnote{K; \skt{पिबे\-त्स\-र्ज्ज\-र\-सस्य}\ H}}\skt{ चापि }\edtext{\skt{नेपाल्यम्वा}}{\Afootnote{K; \skt{नेपाल्याम्वा}\ H}}\skt{ }\edtext{\skt{गोवि\-षा\-णो\-द्भ\-वस्ता}}{\Afootnote{K; \skt{गोवि\-षा\-णो\-द्भ\-वम्वा ।}\ H}}\skt{ ।}\pend

\pstart\versepar \ledsidenote{\passref{6.50.19}}\skt{सर्पिः }\edtext{\skt{स्निग्धं}}{\Afootnote{K H; \skt{स्निग्धैश्}\ A}}\skt{ चर्मबालैः }\edtext{\skt{हृतम्वा}}{\Afootnote{K; \skt{कृतं वा}\ A; \skt{कृताम्वा}\ H}}\skt{ हिक्कास्थाने स्वेदनं }\edtext{\skt{वापि}}{\Afootnote{K; \skt{चापि}\ A H}}\skt{ कार्यं । }\pend
\pstart\versepar \skt{}\edtext{\skt{क्षौद्राप्लुतं}}{\Afootnote{K H; \skt{क्षौद्रोपेतं}\ A}}\skt{ गैरिकं }\edtext{\skt{काञ्चनाख्यं}}{\Afootnote{K H; \skt{काञ्चनाह्वं}\ A}}\skt{ }\edtext{\skt{लिह्या\-द्भ\-स्म\-ग्रा\-म्य\-स\-त्वा\-स्थि\-जम्वा}}{\lemma{\skt{लिह्या\-द्भ\-स्म\-ग्रा\-म्य\-स\-त्वा\-स्थि\-जम्}}\Afootnote{K; \skt{लिह्याद्भस्म ग्राम्यसत्त्वास्थिजं}\ A; \skt{नि… …स्य ग्राम्यसत्त्वास्थिजं}\ H}}\skt{ ।}\pend

\pstart \ledsidenote{\passref{6.50.20a}}\edtext{\skt{रोम्णां}}{\Afootnote{K; \skt{रोम्नां}\ H}}\skt{ }\edtext{\skt{भस्मश्वाविधां}}{\Afootnote{K; \skt{भस्मश्वा विधां}\ H}}\skt{ शल्यकानां }\edtext{\skt{संहृत्य}}{\Afootnote{K; \skt{संहृत्यम्}\ H}}\skt{ }\edtext{\skt{ववस्तगोभ्यां}}{\Afootnote{K; \skt{वा वस्तगोभ्यां}\ H}}\skt{ यथावत् ।}\pend

\pstart \ledsidenote{\passref{6.50.20b}}\skt{सर्पिर्मधुभ्यां }\edtext{\skt{शिग्विपत्रजम्वा}}{\Afootnote{K; \skt{गि…पत्रजम्वा}\ H}}\skt{ }\edtext{\skt{भस्मो\-ल्लि\-हे\-त्पि\-प्प\-लि\-चू\-र्ण्ण\-युक्तं}}{\Afootnote{K; \skt{भस्मो त्पिप्प\-लि\-चू\-र्ण्ण\-युक्तं ।}\ H\textsuperscript{ac}; \skt{भस्मो लिहे\-त्पि\-प्प\-लि\-चू\-र्ण्ण\-युक्तं ।}\ H\textsuperscript{pc}}}\skt{ ।}\pend

\pstart \ledsidenote{\passref{6.50.20c}}\skt{दग्ध्वा }\edtext{\edtext{\skt{फले\-ति\-न्दु\-को\-दु\-म्ब\-रा\-भ्या\-मेवं}}{\lemma{\skt{फले\-ति\-न्दु\-को\-दु\-म्ब\-रा\-भ्याम्}}\Afootnote{K; \skt{फले\-ति\-न्दु\-कौ\-दु\-म्ब\-राब्या}\ H}}}{\lemma{\skt{एवं}}\Afootnote{K; \skt{मेवं}\ H}}\skt{ }\edtext{\skt{लिहेद्वापि}}{\Afootnote{K; \skt{लिहेद्वापि}\ H}}\skt{ गदं जिघांसुः ।}\pend

\pstart \ledsidenote{\passref{6.50.21}}\edtext{\skt{सुवर्चिका}}{\Afootnote{K H; \skt{स्वर्जिक्षारं}\ A}}\skt{ }\edtext{\skt{बीज\-पू\-र्ण्णा\-द्र\-सेन}}{\Afootnote{K H; \skt{बीज\-पू\-रा\-द्र\-सेन}\ A}}\skt{ }\edtext{\skt{क्षौद्रोपेतां}}{\Afootnote{K H; \skt{क्षौद्रोपेतं}\ A}}\skt{ हन्ति लीढ्वाशु }\edtext{\skt{हिक्कां}}{\Afootnote{K H; \skt{हिक्काम् । सर्पिःस्निग्धा घ्नन्ति हिक्कां यवाग्वः कोष्णग्रासाः पायसो वा सुखोष्णः ॥}\ A}}\skt{ ।}\pend

\pstart \ledsidenote{\passref{6.50.21a}}\edtext{\skt{चतुर्गुणेष्वध्मु}}{\lemma{\skt{अध्मु}}\Afootnote{K; \skt{अद्मु}\ H}}\skt{ पिबेत्सुसिद्धं सनागरं }\edtext{\skt{गौडि\-क\-दु\-ग्ध\-माजं}}{\Afootnote{K; \skt{गौडि\-क\-दु\-ग्ध\-माजां ॥}\ H}}\skt{ ।}\pend

\pstart \ledsidenote{\passref{6.50.21cd}}\skt{सर्पिः स्निग्धा घ्नन्ति हिक्कां }\edtext{\skt{सुखोष्णा}}{\Afootnote{K; \skt{सुखोष्णां}\ H}}\skt{ ग्रासायवाग्वः }\edtext{\skt{पयसः}}{\Afootnote{K; \emph{om.}\ H}}\skt{ सुखोष्णाः ॥}\pend

\pstart \ledsidenote{\passref{6.50.22}}\edtext{\skt{याव\-त्तृ\-प्तिं\-चो\-भयं}}{\Afootnote{K H; \skt{शुण्ठीतोये साधितं क्षीरमाजं तद्वत्पीतं शर्करासंयुतं वा ।}\ A}}\skt{ }\edtext{\skt{सेव्यमानं}}{\Afootnote{K H; \skt{आतृप्तेर्वा सेव्यमानं}\ A}}\skt{ }\edtext{\skt{घोरं}}{\Afootnote{K; \skt{निहन्त्याद्घ्रातं}\ A; \skt{घोरां}\ H}}\skt{ }\edtext{\skt{हिक्कां}}{\Afootnote{K H; \skt{हिक्का}\ A}}\skt{ }\edtext{\skt{हन्ति}}{\Afootnote{K H; \skt{माशु}\ A}}\skt{ मूत्रं }\edtext{\skt{त्वजानां}}{\Afootnote{K H; \skt{त्वजाव्योः ॥}\ A}}\skt{ ।}\pend

\pstart\versepar \ledsidenote{\passref{6.50.23}}\edtext{\skt{पूतिः}}{\Afootnote{K; \skt{सपूति}\ A; \skt{पूति}\ H}}\skt{ }\edtext{\skt{कीटां}}{\Afootnote{K; \skt{कीटं}\ A H}}\skt{ }\edtext{\skt{लशुनं}}{\Afootnote{K H; \skt{लशुनो}\ A}}\skt{ }\edtext{\skt{हिंगुमत्स्यां}}{\Afootnote{K H; \skt{ग्रगन्धाहिङ्}\ A}}\skt{ }\edtext{\edtext{\skt{त्वचञ्च}}{\lemma{\skt{त्वचञ्}}\Afootnote{K H; \skt{ग्वब्ज}\ A}}}{\lemma{\skt{च}}\Afootnote{K H; \skt{मा}\ A}}\skt{ }\edtext{\skt{सम्चूर्ण्ण}}{\Afootnote{K H; \skt{चूर्ण्य}\ A}}\skt{ सुभावितन्तत् । }\pend
\pstart\versepar \skt{क्षौद्रं सितां }\edtext{\skt{नांगपुष्पञ्च}}{\Afootnote{K; \skt{वारणकेशरं च}\ A; \skt{नागपुषं च}\ H}}\skt{ }\edtext{\skt{तुल्यं}}{\Afootnote{K H; \emph{om.}\ A}}\skt{ }\edtext{\skt{पिबे\-द्र\-से\-ने\-क्षु\-म\-धू\-क\-जेन}}{\Afootnote{A K; \skt{पिबे\-द्र\-से\-ने\-क्षु\-म\-धू\-क\-जेन ॥}\ H}}\skt{ ।}\pend

\pstart \ledsidenote{\passref{6.50.24}}\skt{पिबेत्पलं वा }\edtext{\skt{लव\-ण\-ओ\-त्त\-मस्य}}{\Afootnote{K; \skt{लवणोत्तमस्य}\ A H}}\skt{ द्वाभ्यां पलाभ्यां हविषः }\edtext{\skt{समग्रं}}{\Afootnote{K H; \skt{समग्रम् । हरीतकीं कोष्ण\-ज\-ला\-नु\-पानां पिबेद्घृतं क्षार\-म\-धू\-प\-प\-न्नम् ॥}\ A}}\skt{ ।}\pend

\pstart\versepar \ledsidenote{\passref{6.50.29}}\edtext{\skt{कपो\-त\-पा\-रा\-व\-त\-श\-ल्य\-को\-द्भवां}}{\Afootnote{K H; \skt{कपो\-त\-पा\-रा\-व\-त\-ला\-व\-श\-क(ल्ल)}\ A}}\skt{ }\edtext{\skt{क्रव्या\-द\-गो\-धा\-वृ\-ष\-दं\-श\-जा\-म्रसां}}{\Afootnote{K; \skt{श्वदं\-ष्ट्र\-गो\-धा\-वृ\-ष\-दं\-श\-जा\-न्रसा}\ A; \skt{क्र्व्याद\-गो\-धा\-वृ\-ष\-दं\-शजां रसां ॥}\ H}}\skt{ ॥ }\pend
\pstart\versepar \skt{}\edtext{\skt{पिबेद्फलाम्व्ला न हि मान्ससैन्धवां}}{\Afootnote{K; \skt{न् । पिबे\-त्फ\-ला\-म्ला\-न\-हि\-मा\-न्स\-सै\-न्धवा}\ A; \skt{पिबे\-त्फ\-ला\-म्लां\-न\-हि\-मां\-स\-सै\-न्धवां}\ H\textsuperscript{ac}; \skt{पिबे\-त्फ\-ला\-म्ला\-न\-हि\-मां\-स\-सै\-न्धवां}\ H\textsuperscript{pc}}}\skt{ }\edtext{\skt{स्निग्धां\-स्था\-र\-ण्य\-मृ\-ग\-द्वि\-जो\-द्भवां}}{\Afootnote{K; \skt{न्स्निग्धांस्तथैव र्ष्यमृ\-ग\-द्वि\-जो\-द्भ\-वान् ॥}\ A; \skt{स्निग्धांस्तथा रण्य\-मृ\-ग\-द्वि\-जो\-द्भवां ।}\ H}}\skt{ ।}\pend

\pstart \ledsidenote{\passref{6.50.24cd}}\skt{हरी\-त\-की\-ञ्चो\-ष्ण\-ज\-ला\-नु\-पानं पिबेत्घृतं क्षार\-र\-जो\-व\-कीर्ण्णं ।}\pend

\pstart \ledsidenote{\passref{6.50.25}}\skt{रसं }\edtext{\skt{कपि\-त्था\-त्म\-धु\-पि\-प्प\-लीञ्च}}{\Afootnote{K; \skt{कपि\-त्था\-न्म\-धु\-पि\-प्प\-लीभ्यां}\ A; \skt{कपि\-त्था\-त्म\-धु\-पि\-प्प\-लीञ्च …}\ H}}\skt{ }\edtext{\skt{पाणिप्रमाणं}}{\Afootnote{K; \skt{शुक्तिप्रमाणं}\ A; \skt{पाणि}\ H}}\skt{ प्रपिबेत्सुखाय ।}\pend

\pstart \ledsidenote{\passref{6.50.26}}\skt{कृष्णां }\edtext{\skt{सिता\-मा\-म\-ल\-कञ्च}}{\lemma{\skt{आमलकञ्च}}\Afootnote{K H; \skt{चामलकं च}\ A}}\skt{ }\edtext{\skt{लिह्यात्स}}{\Afootnote{K H; \skt{लीढं स}\ A}}\skt{ शृङ्गवेरं }\edtext{\skt{मधुनाथवापि}}{\Afootnote{K H; \skt{मधुनाथवापि । कोला\-स्थि\-म\-ज्जा\-ञ्ज\-न\-ला\-ज\-चूर्णं हिक्का निह\-न्या\-न्म\-धु\-ना\-व\-ली\-ढम् ॥}\ A}}\skt{ ।}\pend

\pstart \ledsidenote{\passref{6.50.26a}}\edtext{\edtext{\edtext{\skt{लाजा\-ञ्ज\-ने\-क\-म\-ल\-म\-ध्यञ्च}}{\lemma{\skt{लाजाञ्जनेकम्}}\Afootnote{K; \skt{लाजाञ्जनेको}\ H}}}{\lemma{\skt{अलम्}}\Afootnote{K; \skt{ल}\ H}}}{\lemma{\skt{अध्यञ्च}}\Afootnote{K; \skt{मध्यञ्च}\ H}}\skt{ }\edtext{\skt{तुल्यं}}{\Afootnote{K; \skt{तुलयं}\ H}}\skt{ हिक्का हन्या\-त्पु\-ष्प\-र\-सेन लीढं ॥}\pend

\pstart\versepar \ledsidenote{\passref{6.50.27}}\edtext{\skt{फलं}}{\Afootnote{K H; \skt{पाटलायाः फलं}\ A}}\skt{ }\edtext{\skt{पुष्पञ्च}}{\Afootnote{K H; \skt{पुष्पं}\ A}}\skt{ }\edtext{\skt{पाटल्याः}}{\Afootnote{K; \emph{om.}\ A; \skt{पाटलायाः}\ H}}\skt{ गैरिकं }\edtext{\skt{कटुरोहिणीं}}{\Afootnote{K; \skt{कटुरोहिणी ।}\ A H}}\skt{ । }\pend
\pstart\versepar \skt{}\edtext{\skt{खर्जुरमध्यं}}{\Afootnote{K; \skt{खर्जूरमध्यं}\ A H}}\skt{ मागध्यः }\edtext{\skt{कासीसं}}{\Afootnote{K H; \skt{काशीशं}\ A}}\skt{ दधिनाम च ।}\pend

\pstart\versepar \ledsidenote{\passref{6.50.28}}\skt{चत्वार एते }\edtext{\skt{योगा स्युः}}{\Afootnote{K; \skt{योगाः स्युः}\ A H}}\skt{ }\edtext{\edtext{\skt{पादेष्वेवं}}{\lemma{\skt{पादेष्व्}}\Afootnote{K H; \skt{प्रतिपाद}\ A}}}{\lemma{\skt{एवं}}\Afootnote{K; \skt{प्रद}\ A; \skt{एव}\ H}}\skt{ }\edtext{\edtext{\skt{चतुर्ष्विह}}{\lemma{\skt{चतुर्ष्व्}}\Afootnote{K H; \skt{र्शिताः ।}\ A}}}{\lemma{\skt{इह}}\Afootnote{K H; \emph{om.}\ A}}\skt{ }\pend
\pstart\versepar \skt{}\edtext{\skt{मधुद्वितीयाः}}{\Afootnote{A K; \skt{मधुद्वितीया}\ H}}\skt{ }\edtext{\skt{कर्त्तव्यास्ते}}{\Afootnote{K; \skt{कर्तव्यास्ते}\ A; \skt{कर्त्तव्या सा}\ H}}\skt{ हिक्का सुविजानता ॥}\pend

\pstart\versepar \ledsidenote{\passref{6.50.30}}\skt{विरेचनं पथ्यतमं ससैन्धवं }\pend
\pstart\versepar \skt{}\edtext{\skt{हितं}}{\Afootnote{K H; \skt{घृतं}\ A}}\skt{ सुखोष्णं }\edtext{\skt{प्रवदन्ति}}{\Afootnote{K H; \skt{च सितोप}\ A}}\skt{ }\edtext{\skt{हिक्किने}}{\Afootnote{K H; \skt{लायुतम् ।}\ A}}\skt{ । }\pend
\pstart\versepar \skt{}\edtext{\skt{समी\-र\-णू\-र्ध्व\-गते}}{\Afootnote{K; \skt{सदा\-ग\-ता\-वू\-र्ध्व\-गते}\ A; \skt{समीरणे प्यूर्ध्वगते}\ H}}\skt{ }\edtext{\edtext{\skt{प्यथापरे}}{\lemma{\skt{प्य्}}\Afootnote{K H; \skt{ऽ}\ A}}}{\lemma{\skt{अथापरे}}\Afootnote{K H; \skt{नुवासनं}\ A}}\skt{ }\pend
\pstart\versepar \skt{वदन्ति }\edtext{\skt{नस्यं}}{\Afootnote{K H; \skt{केचि}\ A}}\skt{ }\edtext{\skt{सुसुखाय}}{\Afootnote{K H; \skt{च्च हिताय}\ A}}\skt{ }\edtext{\skt{हिक्किनामिति}}{\lemma{\skt{इति}}\Afootnote{K H; \emph{om.}\ A}}\skt{ ॥}\pend

\pstart \ledsidenote{\passref{6.50.trailer}}\edtext{\skt{काय}}{\Afootnote{K; \skt{इति सुश्रु\-त\-सं\-हि\-ता\-या\-मु\-त्त\-र\-त\-न्त्रा\-न्त\-र्गते काय}\ A; \skt{॥ उत्तर}\ H}}\skt{ }\edtext{\skt{चि}}{\Afootnote{A K; \emph{om.}\ H}}\skt{ ॥ }\edtext{\skt{लृ}}{\Afootnote{K; \skt{कि}\ A; \emph{om.}\ H}}\skt{ ॥ }\edtext{\skt{ओ}}{\Afootnote{K; \skt{त्सातन्त्रे हिक्काप्रतिषेधो नाम (द्वादशोऽध्यायः, आदितः) पञ्चा\-श\-त्त\-मोऽ\-ध्यायः ॥५० ॥}\ A; \skt{तन्त्रे हिक्काचिकित्सा सप्त\-च\-त्वा\-रिं\-श\-त्तमः कायचिकित्सायां द्वादशमो ध्यायः ॥}\ H}}\skt{ ॥}\pend

\endnumbering
\beginnumbering
\eledsection*{Adhyāya 51}
\pstart \ledsidenote{\passref{6.51.1}}\skt{अथातः श्वासप्रतिषेधं }\edtext{\skt{व्याख्यास्यामः}}{\Afootnote{A H; \emph{om.}\ K}}\skt{ ॥}\pend

\pstart\versepar \ledsidenote{\passref{6.51.3}}\edtext{\edtext{\skt{यैरेव}}{\lemma{\skt{यैर्}}\Afootnote{A K; \emph{om.}\ H}}}{\lemma{\skt{एव}}\Afootnote{A K; \emph{om.}\ H}}\skt{ }\edtext{\skt{कारणै}}{\Afootnote{A K; \skt{…}\ H}}\skt{ }\edtext{\skt{हिक्का}}{\Afootnote{K; \skt{र्हिक्का}\ A; \skt{…र्हिक्का}\ H}}\skt{ बहुभिः }\edtext{\skt{संप्रपद्यते}}{\Afootnote{K H; \skt{संप्रवर्तते ।}\ A}}\skt{ । }\pend
\pstart\versepar \skt{}\edtext{\skt{तैरेवकारणैः}}{\Afootnote{K; \skt{तैरेव कारणैः}\ A H}}\skt{ }\edtext{\skt{श्वासो}}{\Afootnote{A K; \skt{श्वासौ}\ H}}\skt{ }\edtext{\skt{घोरो}}{\Afootnote{A K; \skt{घोरौ}\ H}}\skt{ भवति देहिनां ।}\pend

\pstart\versepar \ledsidenote{\passref{6.51.4}}\skt{विहाय प्रकृतिं }\edtext{\skt{वायुरपानः}}{\Afootnote{K H; \skt{वायुः प्राणोऽथ}\ A}}\skt{ कफसंयुतः }\pend
\pstart\versepar \skt{}\edtext{\skt{श्वासयत्यूर्ध्वगो}}{\Afootnote{A K; \skt{श्वासयत्यूर्धगो}\ H}}\skt{ भूत्वा तं श्वासं परिचक्षते ।}\pend

\pstart\versepar \ledsidenote{\passref{6.51.5}}\edtext{\skt{क्षुद्र\-क\-स्त\-म\-क\-स्छिन्नो}}{\Afootnote{K; \skt{क्षुद्र\-क\-स्त\-म\-क\-श्छिन्नो}\ A; \skt{क्षुद्र\-क\-स्त\-म\-क\-च्छिन्नो}\ H}}\skt{ महानूर्ध्वश्च पञ्चधा । }\pend
\pstart\versepar \skt{भिद्यते स }\edtext{\skt{महाव्याधिः}}{\Afootnote{A K; \skt{महा…}\ H}}\skt{ }\edtext{\skt{श्वसः}}{\Afootnote{K; \skt{श्वास}\ A; \emph{om.}\ H}}\skt{ }\edtext{\skt{कोप}}{\Afootnote{K; \skt{एको}\ A; \emph{om.}\ H}}\skt{ }\edtext{\skt{विशेषतः}}{\Afootnote{A K; \emph{om.}\ H}}\skt{ ।}\pend

\pstart\versepar \ledsidenote{\passref{6.51.6}}\edtext{\skt{प्राग्रूपं}}{\Afootnote{A K; \emph{om.}\ H}}\skt{ }\edtext{\skt{तस्य}}{\Afootnote{A K; \skt{…}\ H}}\skt{ हृत्पीडा }\edtext{\edtext{\skt{शूल\-मा\-ध्मा\-न\-मेव}}{\lemma{\skt{शूलमाध्मानम्}}\Afootnote{K H; \skt{भक्तद्वेषो}\ A}}}{\lemma{\skt{एव}}\Afootnote{K H; \skt{ऽरतिः}\ A}}\skt{ }\edtext{\skt{च}}{\Afootnote{K H; \skt{परा ।}\ A}}\skt{ । }\pend
\pstart\versepar \skt{आनाहो }\edtext{\skt{वक्त्रवैरस्यं}}{\Afootnote{K H; \skt{पार्श्वयोः शूलं वैरस्यं}\ A}}\skt{ }\edtext{\edtext{\skt{शंख\-नि\-स्तो\-द\-मेव}}{\lemma{\skt{शंखनिस्तोदम्}}\Afootnote{K; \skt{वदनस्य च ॥}\ A; \skt{शंखनिस्तोद}\ H}}}{\lemma{\skt{एव}}\Afootnote{K H; \emph{om.}\ A}}\skt{ }\edtext{\skt{च}}{\Afootnote{K H; \emph{om.}\ A}}\skt{ ।}\pend

\pstart\versepar \ledsidenote{\passref{6.51.7}}\edtext{\skt{किञ्चिदारभतः}}{\Afootnote{K; \skt{किञ्चिदारभमा}\ A; \skt{किञ्चिदारभते}\ H}}\skt{ }\edtext{\skt{कर्म}}{\Afootnote{K H; \skt{णस्य}\ A}}\skt{ यस्य श्वासः प्रवर्त्तते । }\pend
\pstart\versepar \skt{निष\-ण्ण\-स्यै\-ति\-शा\-न्तिञ्च सक्षुद्र इति संज्ञितः ।}\pend

\pstart \ledsidenote{\passref{6.51.8}}\edtext{\skt{तृट्स्वे\-द\-व\-म\-थु\-प्रायः}}{\Afootnote{A K; \skt{तृट्भे\-द\-व\-म\-थु\-प्रायः}\ H}}\skt{ }\edtext{\skt{कण्ठे}}{\Afootnote{K; \skt{कण्ठ}\ A H}}\skt{ }\edtext{\skt{घुरघुरान्वितः}}{\Afootnote{K; \skt{घुर्घुरिकान्वितः । विशेषाद्दुर्दिने ताम्येच्छ्वासः स तमको मतः ॥}\ A; \skt{घुरि\-घु\-रि\-का\-न्वितः ।}\ H}}\skt{ ।}\pend

\pstart \ledsidenote{\passref{6.51.9}}\skt{घोषेण महता }\edtext{\skt{ताम्यं}}{\Afootnote{K H; \skt{विष्टः}\ A}}\skt{ }\edtext{\skt{सकासं}}{\Afootnote{K H; \skt{सकासः}\ A}}\skt{ }\edtext{\edtext{\skt{सकफन्नरः}}{\lemma{\skt{सकफन्}}\Afootnote{K H; \skt{सकफो}\ A}}}{\lemma{\skt{नरः}}\Afootnote{K H; \skt{नरः । यः श्वसि\-त्य\-ब\-लोऽ\-न्न\-द्वि\-ट्सु\-प्त\-स्त\-म\-क\-पी\-डितः ॥}\ A}}\skt{ ।}\pend

\pstart \ledsidenote{\passref{6.51.9a}}\skt{आनाहो वक्त्रवैरस्यं }\edtext{\skt{शंखनिस्तोद}}{\Afootnote{K; \skt{शङ्खनिस्तोद}\ H}}\skt{ एव च ।}\pend

\pstart \ledsidenote{\passref{6.51.9cd}}\edtext{\skt{यः}}{\Afootnote{K; \skt{यः स्वमे}\ H}}\skt{ }\edtext{\skt{श्वसे\-द्दु\-र्ब\-लो\-न्न\-द्विट्स}}{\Afootnote{K; \skt{सेद्दु\-र्ब्ब\-लो\-न्न\-द्विट्य}\ H\textsuperscript{ac}; \skt{श्वसे\-द्दु\-र्ब्ब\-लो\-न्न\-द्विट्य}\ H\textsuperscript{pc}}}\skt{ }\edtext{\skt{वै}}{\Afootnote{K; \skt{भवै}\ H}}\skt{ तमकसंज्ञितः ।}\pend

\pstart \ledsidenote{\passref{6.51.10}}\edtext{\skt{संशाम्यति}}{\Afootnote{K; \skt{स शाम्यति}\ A H}}\skt{ }\edtext{\skt{कफोच्छित्तौ}}{\Afootnote{K; \skt{कफे हीने}\ A; \skt{कटोच्छिन्त्रौ}\ H}}\skt{ स्वपतश्च }\edtext{\skt{विवर्धते}}{\Afootnote{K H; \skt{विधर्धते । मूर्च्छा\-ज्व\-रा\-भि\-भू\-तस्य ज्ञेयः प्रतमकस्तु सः ॥}\ A}}\skt{ ।}\pend

\pstart\versepar \ledsidenote{\passref{6.51.11}}\edtext{\skt{आध्माति}}{\Afootnote{K H; \skt{आध्मातो}\ A}}\skt{ दह्यमानेन बस्तिना सरुजन्नरः । }\pend
\pstart\versepar \skt{सर्वप्राणेन विच्छिन्नं }\edtext{\skt{श्वासंतं}}{\Afootnote{K; \skt{श्वस्याच्छिन्}\ A; \skt{श्वसन्तं}\ H}}\skt{ }\edtext{\skt{च्छिन्नमादिशेत्}}{\Afootnote{K; \skt{नं तमादिशेत् ॥}\ A; \skt{छिन्नमादिशेत् ।}\ H}}\skt{ ।}\pend

\pstart\versepar \ledsidenote{\passref{6.51.12}}\skt{निःसंज्ञः }\edtext{\skt{पार्श्वशूलार्त्तः}}{\Afootnote{A K; \skt{पार्श्वशूलार्त्तं}\ H}}\skt{ शुष्क\-क\-ण्ठो\-ति\-घो\-ष\-वान् । }\pend
\pstart\versepar \skt{संर\-ब्ध\-ने\-त्र\-स्त्वा\-यम्य यः }\edtext{\skt{श्वसेत्स}}{\Afootnote{K H; \skt{श्वस्यात्स}\ A}}\skt{ }\edtext{\skt{महां}}{\Afootnote{K; \skt{महा}\ A H}}\skt{ }\edtext{\skt{स्मृतः}}{\Afootnote{K; \skt{न्स्मृतः}\ A H}}\skt{ ।}\edtext{}{\lemma{\skt{शुष्क\-क\-ण्ठो\-ति\-घो\-ष\-वान् । संर\-ब्ध\-ने\-त्र\-स्त्वा\-यम्य यः}}\Afootnote{K H; \skt{शुष्क\-क\-ण्ठोऽ\-ति\-घो\-ष\-वान् । संर\-ब्ध\-ने\-त्र\-स्त्वा\-यम्य यः}\ A}}\pend

\pstart\versepar \ledsidenote{\passref{6.51.13}}\edtext{\skt{मर्म\-स्वा\-य\-म्य\-म\-नेषु}}{\Afootnote{K; \skt{मर्म\-स्वा\-य\-म्य\-मा\-नेषु}\ A; \skt{मर्म्म\-स्वा\-य\-म्य\-मा\-नेषु}\ H}}\skt{ }\edtext{\skt{मुहुर्मूढः}}{\Afootnote{K; \skt{श्वसन्मूढो मुहुश्}\ A; \skt{मुहुर्म्मूढः}\ H}}\skt{ }\edtext{\skt{श्वसेत्तु यः}}{\Afootnote{K H; \skt{च यः ।}\ A}}\skt{ }\pend
\pstart\versepar \skt{}\edtext{\skt{ऊर्ध्व\-प्रे\-क्षी\-ह\-त\-र\-ब\-स्त\-मू\-र्ध्व\-श्वा\-स\-मा\-दि\-शेत्}}{\Afootnote{K; \skt{ऊर्ध्वप्रेक्षी हत\-र\-व\-स्त\-मू\-र्ध्व\-श्वा\-स\-मा\-दि\-शेत् ॥}\ A; \skt{ऊर्द्ध्वप्रेक्षी हृत\-र\-व\-स्त\-मू\-र्द्ध्व\-श्वा\-स\-मा\-दि\-शेत् ॥}\ H}}\skt{ ।}\pend

\pstart\versepar \ledsidenote{\passref{6.51.14}}\skt{क्षुद्रः }\edtext{\skt{साध्यतमस्तेषां}}{\Afootnote{K; \skt{साध्यतमस्तेषां}\ A H}}\skt{ तमकः कृच्छ्र उच्यते । }\pend
\pstart\versepar \skt{}\edtext{\skt{त्रय}}{\Afootnote{K; \skt{त्रयः}\ A H}}\skt{ श्वासा न सिध्यन्ति तमको दुर्बलस्य च ।}\pend

\pstart\versepar \ledsidenote{\passref{6.51.15}}\edtext{\skt{स्नेहबस्तिक्रमं}}{\Afootnote{K H; \skt{स्नेहबस्तिं विना}\ A}}\skt{ }\edtext{\skt{केचि\-दू\-र्ध्व\-श्चा\-धश्च}}{\lemma{\skt{केचिदूर्ध्वश्}}\Afootnote{K; \skt{केचिदूर्ध्वं}\ A H}}\skt{ }\edtext{\skt{शोधणं}}{\Afootnote{K; \skt{शोधनम्}\ A H}}\skt{ । }\pend
\pstart\versepar \skt{}\edtext{\skt{मृदुः}}{\Afootnote{K; \skt{मृदु}\ A H}}\skt{ }\edtext{\skt{प्राणवतां}}{\Afootnote{K H; \skt{राणवतां}\ A}}\skt{ श्रेष्ठं }\edtext{\edtext{\skt{श्वासि\-ना\-मा\-दि\-शन्ति}}{\lemma{\skt{श्वासिनाम्}}\Afootnote{K; \skt{श्वासिना}\ A H}}}{\lemma{\skt{आदिशन्ति}}\Afootnote{K; \skt{मादिशन्ति}\ A; \skt{मादिशक्ति}\ H}}\skt{ }\edtext{\skt{ह}}{\Afootnote{K H; \skt{हि ॥}\ A}}\skt{ ।}\pend

\pstart\versepar \ledsidenote{\passref{6.51.16}}\edtext{\skt{कासे}}{\Afootnote{K H; \skt{श्वासे}\ A}}\skt{ }\edtext{\skt{श्वासे च}}{\Afootnote{K H; \skt{कासे च}\ A}}\skt{ हिक्कायां हृद्रोगे चापि पूजितं । }\pend
\pstart\versepar \skt{घृतं पुराणं }\edtext{\skt{संसिद्धं}}{\Afootnote{K H; \skt{संसिद्ध}\ A}}\skt{ }\edtext{\skt{विड\-सौ\-व\-र्च\-ला\-भयैः}}{\Afootnote{K H; \skt{मभ\-या\-वि\-ड\-रा\-मठैः ॥}\ A}}\skt{ ।}\pend

\pstart \ledsidenote{\passref{6.51.17}}\edtext{\skt{पिप्पल्यादि}}{\Afootnote{K; \skt{सौव\-र्च\-ला\-भ\-या\-बिल्वैः संस्कृतं वानवं घृतम् । पिप्पल्यादि}\ A; \skt{पिप्पल्यादिः}\ H}}\skt{ }\edtext{\skt{प्रतीवापं}}{\Afootnote{A K; \skt{प्रतीवापः}\ H}}\skt{ सिद्धम्वा प्रथमे गणे ।}\pend

\pstart\versepar \ledsidenote{\passref{6.51.18}}\skt{सपञ्चलवणं सर्पिः श्वासकासौ व्यपोहति ॥ }\pend
\pstart\versepar \skt{}\edtext{\skt{हिंस्रा\-विं\-डं\-ग\-पू\-ती\-क\-त्रि\-फ\-ला\-व्यो\-ष\-चि\-त्रकैः}}{\Afootnote{K; \skt{हिंस्रा\-वि\-ड\-ङ्ग\-पू\-ती\-क\-त्रि\-फ\-ला\-व्यो\-ष\-चि\-त्रकैः}\ A H}}\skt{ ।}\pend

\pstart\versepar \ledsidenote{\passref{6.51.19}}\edtext{\skt{द्विक्षीरं}}{\Afootnote{A K; \skt{द्विःक्षीरं}\ H}}\skt{ }\edtext{\skt{सर्पिषः}}{\Afootnote{K H; \skt{साधितं}\ A}}\skt{ }\edtext{\skt{प्रस्थ\-ञ्च\-तु\-र्गु\-ण\-ज\-ला\-न्वितं}}{\Afootnote{K H; \skt{सर्पि\-श्च\-तु\-र्गु\-ण\-ज\-ला\-प्लु\-तम् ।}\ A}}\skt{ । }\pend
\pstart\versepar \skt{}\edtext{\skt{कोलमात्रैः}}{\Afootnote{A K; \skt{कोलमात्रं}\ H}}\skt{ }\edtext{\skt{पचेदेभिः}}{\Afootnote{K; \skt{पिबेत्तद्धि}\ A; \skt{पचेदेभिः}\ H}}\skt{ श्वासकासौ व्यपोहति ।}\pend

\pstart\versepar \ledsidenote{\passref{6.51.20}}\edtext{\skt{अर्शांस्यरोचकं}}{\Afootnote{A K; \skt{अशा…रोचकं}\ H}}\skt{ गुल्मं }\edtext{\skt{शकृ\-द्भे\-द\-ञ्क्ष\-य\-न्तथा}}{\Afootnote{A K; \skt{शकृद्भेदं जयेत्तथा ॥}\ H}}\skt{ ॥ }\pend
\pstart\versepar \skt{कृत्स्ने }\edtext{\skt{वृषकषाये तु}}{\Afootnote{K H; \skt{वृषकषाये वा}\ A}}\skt{ }\edtext{\skt{पचे\-त्स\-र्पि\-श्च\-तु\-र्गुणं}}{\Afootnote{K; \skt{पचे\-त्स\-र्पि\-श्च\-तु\-र्गुणे ॥}\ A; \skt{पचे\-त्स\-र्प्पि\-श्च\-तु\-र्ग्गुणं ।}\ H}}\skt{ ।}\pend

\pstart \ledsidenote{\passref{6.51.21}}\edtext{\skt{तन्मू\-ल\-कु\-स\-मा\-वापं}}{\Afootnote{K; \skt{तन्मू\-ल\-कु\-सु\-मा\-वापं}\ A H}}\skt{ }\edtext{\skt{शीतं}}{\Afootnote{A K; \skt{शीत}\ H}}\skt{ क्षौद्रेण }\edtext{\skt{योजयेत्}}{\Afootnote{K H; \skt{योजयेत् । शृङ्गी\-म\-धू\-लि\-का\-भा\-र्गी\-शु\-ण्ठी\-ता\-र्क्ष्य\-सि\-ता\-म्बुदैः ॥}\ A}}\skt{ ॥}\pend

\pstart \ledsidenote{\passref{6.51.21a}}\edtext{\skt{कर्कटाह्वं}}{\Afootnote{K; \skt{कर्क्कटाख्यां}\ H}}\skt{ सितां सुस्तां }\edtext{\skt{भार्गीं}}{\Afootnote{K; \skt{भार्ग्गी}\ H}}\skt{ }\edtext{\skt{शूण्ठीम्मधूलिकां}}{\lemma{\skt{शूण्ठीम्}}\Afootnote{K; \skt{शुण्ठी}\ H}}\skt{ ॥}\pend

\pstart \ledsidenote{\passref{6.51.21b}}\skt{रसाञ्जनं समधुकं }\edtext{\skt{समान्या}}{\Afootnote{K; \skt{सामान्या}\ H}}\skt{ वाप्य योगतः ।}\pend

\pstart \ledsidenote{\passref{6.51.22}}\edtext{\skt{घृतप्रस्थं}}{\Afootnote{K H; \skt{सहरिद्रैः सयष्ट्याह्वैः समैरावाप्य योगतः । घृतप्रस्थं}\ A}}\skt{ }\edtext{\skt{पचेद्धीमां}}{\Afootnote{K; \skt{पचेद्धीमा}\ A; \skt{पचेद्धीमा}\ H}}\skt{ }\edtext{\skt{शीततोये}}{\Afootnote{K; \skt{न्शीततोये}\ A H}}\skt{ चतुर्गुणे ॥}\pend

\pstart \ledsidenote{\passref{6.51.23}}\edtext{\skt{सुवहां}}{\Afootnote{K H; \skt{श्वासं कासं तथा हिक्कां सर्पि\-रे\-त\-न्नि\-य\-च्छति । सुवहा}\ A}}\skt{ }\edtext{\skt{कालिकां}}{\Afootnote{K H; \skt{कालिका}\ A}}\skt{ }\edtext{\skt{भार्गीशुकाख्यां}}{\Afootnote{K; \skt{भर्गी शुकाख्या}\ A; \skt{भर्ग्गीं शुकाख्यां}\ H}}\skt{ नैचुलं फलं ।}\pend

\pstart\versepar \ledsidenote{\passref{6.51.24}}\edtext{\skt{काकादनीं}}{\Afootnote{K H; \skt{काकादनी}\ A}}\skt{ शृंगवेरं }\edtext{\skt{वर्षा\-भू\-बृ\-ह\-ती\-द्वयं}}{\Afootnote{K H; \skt{वर्ष\-भू\-र्बृ\-ह\-ती\-द्व\-यम् ।}\ A}}\skt{ । }\pend
\pstart\versepar \skt{}\edtext{\edtext{\skt{कोल\-मा\-त्रै\-र्घृ\-त\-प्रस्थं}}{\lemma{\skt{कोलमात्रैर्}}\Afootnote{K H; \skt{कोलमात्रै}\ A}}}{\lemma{\skt{घृतप्रस्थं}}\Afootnote{K H; \skt{र्घृतप्रस्थं}\ A}}\skt{ }\edtext{\edtext{\skt{पचे\-दे\-तै\-र्ज\-ल\-द्विकं}}{\lemma{\skt{पचेदेतैर्}}\Afootnote{K; \skt{पचेदेभिर्}\ A; \skt{पचेदेतैर्}\ H}}}{\lemma{\skt{जलद्विकं}}\Afootnote{A K; \skt{जलादिकं ॥}\ H\textsuperscript{ac}; \skt{जलद्विलादिकं ॥}\ H\textsuperscript{pc}}}\skt{ ॥}\pend

\pstart \ledsidenote{\passref{6.51.25}}\edtext{\skt{सौव\-र्च\-ल\-य\-व\-क्षा\-र\-क\-टु\-का\-व्यो\-ष\-सै\-न्धवैः}}{\Afootnote{K H; \skt{कटूष्णं पीतमेतद्धि श्वासा\-म\-य\-वि\-ना\-श\-नम् । सौव\-र्च\-ल\-य\-व\-क्षा\-र\-क\-टु\-का\-व्यो\-ष\-चि\-त्रकैः ॥}\ A}}\skt{ ।}\pend

\pstart\versepar \ledsidenote{\passref{6.51.26}}\skt{वचा\-भ\-या\-वि\-ड\-ङ्गैश्च }\edtext{\skt{पचेद्वा}}{\Afootnote{K H; \skt{साधितं}\ A}}\skt{ }\edtext{\skt{विधिवद्घृतं}}{\Afootnote{K H; \skt{श्वासशान्तये ।}\ A}}\skt{ ॥ }\pend
\pstart\versepar \skt{}\edtext{\skt{गोपवल्युदके}}{\Afootnote{K; \skt{गोपवल्ल्युदके}\ A; \skt{गोपवल्ल्युदके}\ H}}\skt{ सिद्धं }\edtext{\skt{स्यादन्य}}{\Afootnote{K; \skt{स्यादन्य}\ A H}}\skt{ }\edtext{\skt{द्विगुणे}}{\Afootnote{K; \skt{द्द्विगुणे}\ A H}}\skt{ घृतं}\pend

\pstart \ledsidenote{\passref{6.51.27}}\skt{पञ्चैतानि }\edtext{\edtext{\edtext{\skt{हवीं\-ष्या\-हु\-र्भि\-षजः}}{\lemma{\skt{हवींष्य्}}\Afootnote{K; \skt{हवीं}\ A H}}}{\lemma{\skt{आहुर्}}\Afootnote{K; \skt{ष्याहु}\ A; \skt{ष्याहुर्}\ H}}}{\lemma{\skt{भिषजः}}\Afootnote{K; \skt{र्भिषजः}\ A; \skt{भ्भिषजः}\ H}}\skt{ }\edtext{\skt{श्वासकासयोः}}{\Afootnote{K H; \skt{श्वासकासयोः । ताली\-श\-ता\-म\-ल\-क्यु\-ग्रा\-जी\-व\-न्ती\-कु\-ष्ठ\-सै\-न्धवैः ॥}\ A}}\skt{ ॥}\pend

\pstart\versepar \ledsidenote{\passref{6.51.30}}\skt{तैलं }\edtext{\skt{दशगुणे}}{\Afootnote{A K; \skt{दशगुणैः}\ H}}\skt{ सिद्धं }\edtext{\skt{केसराजरसे}}{\Afootnote{K; \skt{भृङ्गराजरसे}\ A H}}\skt{ शुभे । }\pend
\pstart\versepar \skt{}\edtext{\skt{पीयमानं}}{\Afootnote{K H; \skt{सेव्यमानं}\ A}}\skt{ यथान्यायं श्वासकासौ व्यपोहति ॥}\pend

\pstart\versepar \ledsidenote{\passref{6.51.31}}\edtext{\skt{फला\-म्व्ला\-वि\-ष्कि\-र\-रसा}}{\Afootnote{K; \skt{फलाम्ला विष्किररसाः}\ A H}}\skt{ }\edtext{\skt{स्निग्धाः}}{\Afootnote{A K; \skt{…}\ H}}\skt{ प्रव्यक्तसैन्धवाः । }\pend
\pstart\versepar \skt{एणादीनां शिरोभिर्वा }\edtext{\skt{सकुलत्थाः}}{\Afootnote{K H; \skt{कौलत्था}\ A}}\skt{ }\edtext{\skt{ससैन्धवाः}}{\Afootnote{K H; \skt{वा सुसंस्कृताः ॥}\ A}}\skt{ ।}\pend

\pstart\versepar \ledsidenote{\passref{6.51.32}}\skt{हन्युः श्वासञ्च }\edtext{\skt{कासश्च}}{\Afootnote{K; \skt{कासं च}\ A H}}\skt{ }\edtext{\skt{संस्कृतानि}}{\Afootnote{A K; \skt{सकृतानि}\ H}}\skt{ पयांसि च ॥ }\pend
\pstart\versepar \skt{तिमिरस्य च बीजानि कर्कटाख्या च }\edtext{\skt{चूर्ण्णिताः}}{\Afootnote{K; \skt{चूर्णिता}\ A H}}\skt{ ।}\pend

\pstart\versepar \ledsidenote{\passref{6.51.33}}\skt{दुरा\-ल\-भा\-थ\-पि\-प्पल्यः }\edtext{\skt{कटुकाख्य}}{\Afootnote{K; \skt{कटुकाख्या}\ A H}}\skt{ हरीतकी । }\pend
\pstart\versepar \skt{}\edtext{\skt{श्वावि\-ड्म\-यू\-र\-रो\-माणि}}{\Afootnote{K H; \skt{श्वावि\-न्म\-यू\-र\-रो\-माणि}\ A}}\skt{ }\edtext{\skt{कोल\-पि\-प्प\-लि\-त\-ण्डुलाः}}{\Afootnote{K H; \skt{कोला मागधिकाकणाः ॥}\ A}}\skt{ ।}\pend

\pstart\versepar \ledsidenote{\passref{6.51.34}}\skt{भार्गी }\edtext{\skt{त्वक्शृङ्गवेरञ्च}}{\Afootnote{A K; \skt{त्व …वेरण्ड}\ H}}\skt{ }\edtext{\skt{शर्कराशल्यकं}}{\Afootnote{K; \skt{शर्करा शल्लकाङ्}\ A; \skt{शर्क्करा शल्य्कं}\ H}}\skt{ }\edtext{\skt{गजं}}{\Afootnote{A K; \skt{गजा ।}\ H}}\skt{ । }\pend
\pstart\versepar \skt{}\edtext{\skt{वृत्त\-का\-ण्ट\-क\-बी\-जानि}}{\Afootnote{K; \skt{नृत्त\-कौ\-ण्ड\-क\-बी\-जानि}\ A; \skt{कटु\-त्रि\-क\-स्य\-बी\-जानि}\ H}}\skt{ चूर्ण्णितानि तु केवलं ।}\pend

\pstart\versepar \ledsidenote{\passref{6.51.35}}\edtext{\skt{पञ्चश्लोकाधिका}}{\Afootnote{K; \skt{पञ्च श्लोकार्धिका}\ A H}}\skt{ }\edtext{\skt{ह्येते}}{\Afootnote{K; \skt{स्त्वेते}\ A; \skt{ह्येते}\ H}}\skt{ लेहा ये सम्यगीरिताः । }\pend
\pstart\versepar \skt{}\edtext{\skt{सर्पिर्मधुभ्यां}}{\Afootnote{K; \skt{सर्पिर्मधुभ्यां ते}\ A H}}\skt{ }\edtext{\skt{संलेहा}}{\Afootnote{K; \skt{लेह्याः}\ A; \skt{संलेह्य}\ H}}\skt{ }\edtext{\skt{श्वास\-का\-सा\-र्दि\-तै\-र्निरैः}}{\Afootnote{K; \skt{कास\-श्वा\-सा\-र्दि\-तै\-र्नरैः ॥}\ A; \skt{श्वास\-का\-सा\-र्द्दि\-तै\-र्न्नरैः ॥}\ H}}\skt{ ।}\pend

\pstart\versepar \ledsidenote{\passref{6.51.36}}\skt{सप्तच्छदस्य पुष्पाणि }\edtext{\skt{पिप्पल्यश्चापि}}{\lemma{\skt{पिप्पल्यश्}}\Afootnote{K H; \skt{पिप्पलीश्}\ A}}\skt{ मस्तुना । }\pend
\pstart\versepar \skt{}\edtext{\skt{पिबेत्सम्पिष्य}}{\Afootnote{K H; \skt{पिबेत्संचूर्ण्य}\ A}}\skt{ मधुना }\edtext{\edtext{\edtext{\skt{धानाम्वाप्यथ}}{\lemma{\skt{धानाम्}}\Afootnote{K; \skt{धानाश्}\ A; \skt{… ॥}\ H}}}{\lemma{\skt{वाप्य्}}\Afootnote{K; \skt{चाप्य्}\ A; \emph{om.}\ H}}}{\lemma{\skt{अथ}}\Afootnote{A K; \emph{om.}\ H}}\skt{ }\edtext{\skt{भक्षयेत्}}{\Afootnote{A K; \emph{om.}\ H}}\skt{ ।}\pend

\pstart\versepar \ledsidenote{\passref{6.51.37}}\edtext{\edtext{\skt{अर्का\-ङ्कु\-रै\-र्भा\-वि\-तानां}}{\lemma{\skt{अर्काङ्कुरैर्}}\Afootnote{K; \skt{अर्काङ्कुरै}\ A; \emph{om.}\ H}}}{\lemma{\skt{भावितानां}}\Afootnote{K; \skt{र्भावितानां}\ A; \emph{om.}\ H}}\skt{ }\edtext{\skt{यवानां}}{\Afootnote{A K; \skt{…}\ H}}\skt{ }\edtext{\skt{साध्वनेकशः}}{\Afootnote{A K; \skt{…ध्वनेकशः ।}\ H}}\skt{ । }\pend
\pstart\versepar \skt{}\edtext{\skt{सर्पिणा}}{\Afootnote{K; \skt{तर्पणं}\ A; \skt{सर्प्पिषा}\ H}}\skt{ वापि }\edtext{\skt{वेदेषां}}{\Afootnote{K H; \skt{बेदेषां}\ A}}\skt{ }\edtext{\skt{सक्षौद्रां}}{\Afootnote{K H; \skt{सक्षौद्रं}\ A}}\skt{ श्वासपीडितः ।}\pend

\pstart\versepar \ledsidenote{\passref{6.51.38}}\edtext{\skt{शिरी\-ष\-कु\-न्द\-क\-द\-ली\-पुष्पं}}{\Afootnote{K H; \skt{शिरी\-ष\-क\-द\-ली\-कु\-न्द\-पुष्पं}\ A}}\skt{ }\edtext{\skt{पिप्पलिसंयुतं}}{\Afootnote{K H; \skt{माग\-धि\-का\-यु\-तम् ।}\ A}}\skt{ । }\pend
\pstart\versepar \skt{तण्डुलाम्बुयुतं पीत्वा जये\-च्छ्वा\-सा\-न\-शे\-षतः ।}\pend

\pstart \ledsidenote{\passref{6.51.40cd}}\skt{द्राक्षां }\edtext{\skt{हरीतकी}}{\Afootnote{K; \skt{हरीतकीङ्}\ H}}\skt{ }\edtext{\skt{कृष्णआं}}{\Afootnote{K; \skt{कृष्णां}\ H}}\skt{ कर्क\-टा\-ख्या\-न्दु\-रा\-लभां}\pend

\pstart \ledsidenote{\passref{6.51.41ab}}\skt{सर्पिर्मधुभ्यां विलिहं }\edtext{\skt{च्छ्वासां}}{\Afootnote{K; \skt{छ्वासा}\ H}}\skt{ }\edtext{\skt{हन्ति}}{\Afootnote{K; \skt{न्हन्ति}\ H}}\skt{ }\edtext{\skt{सुदुस्तरां}}{\Afootnote{K; \skt{सुदुस्तरान् ॥}\ H}}\skt{ ।}\pend

\pstart\versepar \ledsidenote{\passref{6.51.39}}\edtext{\skt{कोलिमज्जा}}{\Afootnote{K H; \skt{कोलमज्जां}\ A}}\skt{ }\edtext{\skt{तालुमूलं}}{\Afootnote{K; \skt{तालमूल}\ A; \skt{तालमूलं}\ H}}\skt{ }\edtext{\skt{पिष्ट्वा}}{\Afootnote{K H; \skt{मृष्य}\ A}}\skt{ }\edtext{\skt{चर्ममशिन्तथा}}{\Afootnote{K; \skt{चर्ममसीमपि ।}\ A; \skt{चर्म्ममसीं तथा ।}\ H}}\skt{ । }\pend
\pstart\versepar \skt{लिह्यात्क्षौद्रेण }\edtext{\skt{पामा\-र्ग\-स\-र्पि\-र्म\-धु\-स\-मा\-युतं}}{\Afootnote{K; \skt{भार्गीं वा सर्पि\-र्म\-धु\-स\-मा\-यु\-ताम् ॥}\ A; \skt{पामा… सर्प्पि\-र्म्म\-धु\-स\-मा\-युतं ॥}\ H}}\skt{ ।}\pend

\pstart \ledsidenote{\passref{6.51.40}}\edtext{\skt{नीपः}}{\Afootnote{K; \skt{नीचैः}\ A; \skt{नी}\ H}}\skt{ }\edtext{\skt{कदम्बो}}{\Afootnote{K H; \skt{कदम्ब}\ A}}\skt{ }\edtext{\skt{रक्तघ्नः तं}}{\Afootnote{K H; \skt{बीजं वा सक्षौद्रं}\ A}}\skt{ }\edtext{\skt{पीत्वा}}{\Afootnote{K H; \emph{om.}\ A}}\skt{ }\edtext{\skt{तण्डुलाम्बुना}}{\Afootnote{K H; \skt{तण्डुलाम्बुना । द्राक्षां हरीतकीं कृष्णां कर्कटाख्यां दुरालभाम् ॥}\ A}}\skt{ ।}\pend

\pstart \ledsidenote{\passref{6.51.41}}\edtext{\skt{हरिद्रां}}{\Afootnote{K H; \skt{सर्पिर्मधुभ्यां विलिहन्हन्ति श्वासा\-न्सु\-दा\-रु\-णान् । हरिद्रां}\ A}}\skt{ मरिचं }\edtext{\skt{रास्नां}}{\Afootnote{K H; \skt{द्राक्षां}\ A}}\skt{ गुडं }\edtext{\skt{द्राक्षां}}{\Afootnote{K H; \skt{रास्नां}\ A}}\skt{ }\edtext{\skt{सपिप्पलीं}}{\Afootnote{K H; \skt{कणां शटीम् ॥}\ A}}\skt{ ।}\pend

\pstart \ledsidenote{\passref{6.51.42}}\edtext{\skt{लिह्यात्क्षौद्रेण}}{\Afootnote{K H; \skt{लिह्यात्तैलेन}\ A}}\skt{ तुल्यानि श्वासार्त्तो }\edtext{\skt{हितभोजनः}}{\Afootnote{K H; \skt{हितभोजनः । गवां पुरीषस्वरसं मधु\-स\-र्पिः\-क\-णा\-यु\-तम् ॥}\ A}}\skt{ ॥}\pend

\pstart \ledsidenote{\passref{6.51.44}}\edtext{\edtext{\skt{भार्गीन्त्रिकण्टकं}}{\lemma{\skt{भार्गीन्}}\Afootnote{K H; \skt{श्वास\-का\-सा\-प\-हा\-स्तेऽपि कासघ्ना ये च कीर्तिताः । भार्गी}\ A}}}{\lemma{\skt{त्रिकण्टकं}}\Afootnote{K; \skt{त्वक्त्र्यूषणं}\ A; \skt{त्रिकटुकं}\ H}}\skt{ तैलं हरिद्रां कटुरोहिणीं ।}\pend

\pstart\versepar \ledsidenote{\passref{6.51.45}}\edtext{\skt{पिप्पल्यो}}{\Afootnote{K H; \skt{पिप्पलीं}\ A}}\skt{ मरिचं चण्डां }\edtext{\skt{गोशकृद्रस}}{\Afootnote{K H; \skt{गोशकृद्रसम्}\ A}}\skt{ एव च । }\pend
\pstart\versepar \skt{}\edtext{\skt{तटकोलस्य}}{\Afootnote{K; \skt{तलकोटस्य}\ A; \skt{तटकोटिस्य}\ H}}\skt{ }\edtext{\skt{बीजानि}}{\Afootnote{K H; \skt{बीजेषु}\ A}}\skt{ }\edtext{\skt{पचेदुत्कारिकां}}{\Afootnote{K H; \skt{पचेदुत्कारिकां}\ A}}\skt{ शुभां ।}\pend

\pstart \ledsidenote{\passref{6.51.46}}\skt{सेव्यमाना }\edtext{\edtext{\skt{निहन्त्येषा}}{\lemma{\skt{निहन्त्य्}}\Afootnote{K; \skt{निहन्}\ A; \skt{निहन्ते}\ H}}}{\lemma{\skt{एषा}}\Afootnote{K; \skt{त्येषा}\ A; \skt{षां}\ H}}\skt{ }\edtext{\skt{श्वासानाशु}}{\Afootnote{A K; \skt{श्वासान्हन्ति}\ H}}\skt{ }\edtext{\edtext{\skt{सुदुर्जयामिति}}{\lemma{\skt{सुदुर्जयाम्}}\Afootnote{K; \skt{सुदुस्तरा}\ A; \skt{सुदुर्ज्जया}\ H}}}{\lemma{\skt{इति}}\Afootnote{K; \skt{न् । पुराणसर्पिः पिप्पल्यः कौलत्था जाङ्गला रसाः ॥}\ A; \skt{निति ॥}\ H}}\skt{ ॥}\pend

\pstart \ledsidenote{\passref{6.51.trailer}}\edtext{\skt{कायचि}}{\Afootnote{K; \skt{इति सुश्रु\-त\-सं\-हि\-ता\-या\-मु\-त्त\-र\-त\-न्त्रा\-न्त\-र्गते कायचि}\ A; \emph{om.}\ H}}\skt{ ॥ }\edtext{\skt{लृ}}{\Afootnote{K; \skt{कि}\ A; \emph{om.}\ H}}\skt{ ॥ }\edtext{\skt{ओ}}{\Afootnote{K; \skt{त्सातन्त्रे हिक्का\-श्वा\-स\-प्र\-ति\-षेधो नाम (त्र्योदशोऽध्यायः, आदितः) एक\-प\-ञ्चा\-श\-त्त\-मोऽ\-ध्यायः ॥५१ ॥}\ A; \skt{उत्तरतन्त्रे श्वास\-चि\-कि\-त्सा\-ष्ट\-च\-त्वा\-रिं\-श\-त्तमः कायचिकित्सायां त्रयोदश ॥}\ H}}\skt{ ॥}\pend

\endnumbering
\beginnumbering
\eledsection*{Adhyāya 52}
\pstart \ledsidenote{\passref{6.52.1}}\skt{अथातः कासप्रतिषेधं }\edtext{\skt{व्याख्यास्यामः}}{\Afootnote{K H; \skt{अध्यायं व्याख्यास्यामः ॥}\ A}}\skt{ ॥}\pend

\pstart\versepar \ledsidenote{\passref{6.52.4A}}\edtext{\skt{यैरेव}}{\lemma{\skt{एव}}\Afootnote{K; \skt{एवं}\ H}}\skt{ कारणैर्हिक्का श्वास\-श्चै\-वो\-प\-जा\-यते । }\pend
\pstart\versepar \skt{}\edtext{\skt{तैरेव}}{\lemma{\skt{एव}}\Afootnote{K; \emph{om.}\ H}}\skt{ }\edtext{\edtext{\skt{कारणैर्नृणां}}{\lemma{\skt{कारणैर्}}\Afootnote{K; \emph{om.}\ H}}}{\lemma{\skt{नृणां}}\Afootnote{K; \emph{om.}\ H}}\skt{ }\edtext{\skt{कासः}}{\Afootnote{K; \emph{om.}\ H}}\skt{ }\edtext{\skt{समुपजायते}}{\Afootnote{K; \emph{om.}\ H}}\skt{ ।}\edtext{}{\lemma{\skt{चैवोपजायते । तैर्}}\Afootnote{K; \skt{चैवोपजायते ॥}\ H}}\pend

\pstart\versepar \ledsidenote{\passref{6.52.5}}\edtext{\skt{प्राणो}}{\Afootnote{A K; \skt{प्राणो ह्य्}\ H}}\skt{ }\edtext{\skt{ह्युदानानुगतः}}{\Afootnote{A K; \skt{उदानानुगतः}\ H}}\skt{ }\edtext{\skt{प्रदुष्टः}}{\Afootnote{A K; \skt{प्रदिष्टः}\ H}}\skt{ }\edtext{\skt{संभि\-न्न\-कां\-स\-स्व\-न\-तु\-ल्य\-घो\-न\-तु\-ल्य\-घोषः}}{\Afootnote{K H; \skt{संभि\-न्न\-कां\-स्य\-स्व\-न\-तु\-ल्य\-घोषः ।}\ A}}\skt{ । }\pend
\pstart\versepar \skt{निरेति वक्त्रात्सहसा }\edtext{\skt{सघोषः}}{\Afootnote{K; \skt{सदोषः}\ A; \skt{सघोषं}\ H}}\skt{ कासः स विद्व\-द्भि\-रु\-दा\-हृ\-तस्तु ।}\pend

\pstart\versepar \ledsidenote{\passref{6.52.6}}\skt{स वात\-पि\-त्त\-प्र\-भवः कफाच्च क्षतात्तथान्यः }\edtext{\skt{क्षयजोपरश्च}}{\Afootnote{K; \skt{क्षयजोऽपरश्च}\ A H}}\skt{ । }\pend
\pstart\versepar \skt{पञ्चप्रकारः }\edtext{\skt{पठितो}}{\Afootnote{K H; \skt{कथितो}\ A}}\skt{ }\edtext{\skt{भिष\-ग्भि\-र्वि\-शे\-षतो}}{\lemma{\skt{विशेषतो}}\Afootnote{K H; \skt{विवर्धितो}\ A}}\skt{ }\edtext{\skt{लक्षणतः}}{\Afootnote{K H; \skt{यक्ष्म}\ A}}\skt{ }\edtext{\skt{पुरा}}{\Afootnote{K; \skt{विकारकृत्स्यात् ॥}\ A; \skt{पुराणैः ॥}\ H}}\skt{ ॥}\pend

\pstart\versepar \ledsidenote{\passref{6.52.8}}\edtext{\skt{हृच्शं\-ख\-मू\-र्द्धो\-द\-र\-पा\-र्श्व\-शूली}}{\Afootnote{K H; \skt{हृच्छ\-ङ्ख\-मू\-र्धो\-द\-र\-पा\-र्श्व\-शूली}\ A}}\skt{ क्षामाननः }\edtext{\skt{क्षीणवलः}}{\Afootnote{K; \skt{क्षीणबल}\ A; \skt{क्षीणबलः}\ H}}\skt{ }\edtext{\skt{क्षतौजा}}{\Afootnote{K; \skt{स्वरौजाः ।}\ A; \skt{क्षतौजाः ।}\ H}}\skt{ । }\pend
\pstart\versepar \skt{}\edtext{\skt{प्रसक्तमन्तः}}{\Afootnote{A K; \skt{प्रसक्तमन्तः}\ H}}\skt{ }\edtext{\edtext{\skt{कफ\-अ\-मी\-र\-णेन}}{\lemma{\skt{कफअम्}}\Afootnote{K; \skt{कफ}\ A; \skt{सस}\ H}}}{\lemma{\skt{ईरणेन}}\Afootnote{K; \skt{मीरणे}\ A; \skt{मीरणेन}\ H}}\skt{ }\edtext{\skt{भिन्नस्वरः}}{\Afootnote{K H; \skt{न}\ A}}\skt{ }\edtext{\skt{कासति}}{\Afootnote{K H; \skt{कासेत्तु}\ A}}\skt{ }\edtext{\skt{शुष्कमेव}}{\lemma{\skt{एव}}\Afootnote{K H; \skt{स्वरभेदयुक्तः ॥}\ A}}\skt{ ।}\pend

\pstart\versepar \ledsidenote{\passref{6.52.9}}\edtext{\edtext{\skt{उरो\-वि\-दा\-ह\-ज्व\-र\-व\-क्त्र\-शो\-षै\-र\-भ्य\-र्दि\-त\-स्ति\-क्त\-मु\-ख\-स्तृ\-षार्तः}}{\lemma{\skt{उरो\-वि\-दा\-ह\-ज्व\-र\-व\-क्त्र\-शो\-षैर्}}\Afootnote{K H; \skt{उरो\-वि\-दा\-ह\-ज्व\-र\-व\-क्त्र\-शोषै}\ A}}}{\lemma{\skt{अभ्य\-र्दि\-त\-स्ति\-क्त\-मु\-ख\-स्तृ\-षार्तः}}\Afootnote{K; \skt{रभ्य\-र्दि\-त\-स्ति\-क्त\-मु\-ख\-स्तृ\-षार्तः ।}\ A; \skt{अभ्य\-र्दि\-त\-स्ति\-क्त\-मु\-ख\-स्तृ\-षार्त्तः ।}\ H}}\skt{ । }\pend
\pstart\versepar \skt{पित्तेन पीतानि वमेत्कटूनि }\edtext{\skt{कासेत}}{\Afootnote{K H; \skt{कासेत्स}\ A}}\skt{ पाण्डुः परिदह्य मानः ।}\pend

\pstart\versepar \ledsidenote{\passref{6.52.10}}\edtext{\skt{प्रलिप्यमानेन}}{\Afootnote{K H; \skt{प्रलि\-प्य\-मा\-ने\-न(\-वि\-लि\-प्य\-मा\-नेन)}\ A}}\skt{ मुखेन }\edtext{\skt{सीद\-न्शि\-रो\-रु\-जार्तः}}{\Afootnote{A K; \skt{सीद\-न्शि\-रो\-रु\-जार्त्त}\ H}}\skt{ कफपूर्ण्णदेहः । }\pend
\pstart\versepar \skt{}\edtext{\skt{अभ\-क्त\-रु\-ग्गौ\-र\-व\-पा\-ण्डु\-युक्तः}}{\Afootnote{K H; \skt{अभ\-क्त\-रु\-ग्गौ\-र\-व\-सा\-द\-युक्तः}\ A}}\skt{ }\edtext{\skt{कासेद्भृशं}}{\Afootnote{K; \skt{कासेत ना}\ A; \skt{कासेद्भृशं}\ H}}\skt{ }\edtext{\skt{सान्द्रकफः}}{\Afootnote{K H; \skt{सान्द्रकफं}\ A}}\skt{ कफेन ।}\pend

\pstart\versepar \ledsidenote{\passref{6.52.11}}\edtext{\skt{वक्षोतिमात्रं}}{\Afootnote{K; \skt{वक्षोऽतिमात्रं}\ A H}}\skt{ }\edtext{\skt{विहतं तु}}{\Afootnote{A K; \skt{विहितं}\ H\textsuperscript{ac}; \skt{विहितं तु}\ H\textsuperscript{pc}}}\skt{ यस्य व्याया\-म\-भा\-रा\-ध्य\-य\-ना\-भि\-घातैः }\pend
\pstart\versepar \skt{विश्लिष्टवक्षाः स नरः सरक्तं }\edtext{\skt{ष्ठीवत्यभीक्ष्णं}}{\Afootnote{A K; \skt{ष्ठीवत्यभीक्ष्णं}\ H}}\skt{ क्षतजं }\edtext{\edtext{\skt{तमाहुः}}{\lemma{\skt{तम्}}\Afootnote{K H; \skt{त}\ A}}}{\lemma{\skt{आहुः}}\Afootnote{K H; \skt{मूहुः ॥}\ A}}\skt{ ।}\pend

\pstart\versepar \ledsidenote{\passref{6.52.12}}\skt{स }\edtext{\skt{गात्र\-शू।\-लं\-ज्व\-र\-दा\-ह\-मोहं}}{\Afootnote{K; \skt{गात्र\-शू\-ल\-ज्व\-र\-दा\-ह\-मोहा}\ A; \skt{गात्र\-शू\-ल\-ज्व\-र\-दा\-ह\-मोहं}\ H}}\skt{ }\edtext{\skt{प्राणक्षयं}}{\Afootnote{K H; \skt{न्प्राणक्षयं}\ A}}\skt{ चोपलभेत कासात् । }\pend
\pstart\versepar \skt{शुष्य\-न्वि\-नि\-ष्ठी\-वति दुर्ब्बलस्तु प्रक्षीणमांसो रुधिरं सपूयं ॥}\edtext{}{\lemma{\skt{कासात् । शुष्य\-न्वि\-नि\-ष्ठी\-वति}}\Afootnote{K H; \skt{कासी । शुष्य\-न्वि\-नि\-ष्ठी\-वति}\ A}}\pend

\pstart \ledsidenote{\passref{6.52.13}}\edtext{\skt{तं}}{\Afootnote{K H; \skt{स}\ A}}\skt{ सर्वलिङ्गं }\edtext{\skt{भृशदुश्चिकित्स्यं}}{\Afootnote{K; \skt{भृशदुश्चिकित्स्यं}\ A; \skt{भृशं दुश्चिकित्स्यं}\ H}}\skt{ }\edtext{\skt{चिकित्सितज्ञे}}{\Afootnote{K; \skt{चिकित्सितज्ञाः}\ A H}}\skt{ }\edtext{\edtext{\skt{क्षयकासमाहुः}}{\lemma{\skt{क्षयकासम्}}\Afootnote{K H; \skt{क्षयजं वदन्ति । वृद्धत्वमासाद्य भवेत्तु यो वै याप्यं तमा\-हु\-र्भि\-ष\-जस्तु कासम् ॥}\ A}}}{\lemma{\skt{आहुः}}\Afootnote{K H; \emph{om.}\ A}}\skt{ ॥}\pend

\pstart\versepar \ledsidenote{\passref{6.52.16}}\edtext{\skt{पथ्यां}}{\Afootnote{A K; \skt{पथ्या}\ H}}\skt{ }\edtext{\edtext{\skt{सिता\-मा\-म\-ल\-कानि}}{\lemma{\skt{सिताम्}}\Afootnote{K; \skt{सिता}\ A H}}}{\lemma{\skt{आमलकानि}}\Afootnote{K; \skt{मामलकानि}\ A; \skt{माम…लकानि}\ H}}\skt{ }\edtext{\skt{लाजाः}}{\Afootnote{K H; \skt{लाजां}\ A}}\skt{ }\edtext{\skt{समागधञ्चापि}}{\lemma{\skt{समागधञ्}}\Afootnote{K H; \skt{समागधीं}\ A}}\skt{ विचूर्ण्ण्य }\edtext{\skt{शुण्ठी}}{\Afootnote{K H; \skt{शुण्ठीम् ।}\ A}}\skt{ । }\pend
\pstart\versepar \skt{}\edtext{\edtext{\skt{सर्पिर्मधुभ्यां}}{\lemma{\skt{सर्पिर्}}\Afootnote{K; \skt{सर्पि}\ A H}}}{\lemma{\skt{मधुभ्यां}}\Afootnote{K; \skt{र्मधुभ्यां}\ A H}}\skt{ }\edtext{\skt{विलिहेद्धि}}{\Afootnote{K; \skt{विलिहीत}\ A; \skt{विन्निहेद्धि}\ H}}\skt{ कासी }\edtext{\skt{ससैन्धवं}}{\Afootnote{K H; \skt{ससैन्धवां}\ A}}\skt{ }\edtext{\skt{चोष्णजलेन}}{\Afootnote{K H; \skt{वोष्णजलेन}\ A}}\skt{ कृष्णां ।}\pend

\pstart\versepar \ledsidenote{\passref{6.52.17}}\edtext{\skt{पिबेद्गुडं}}{\Afootnote{K H; \skt{खादेद्गुडं}\ A}}\skt{ }\edtext{\skt{पिप्प\-लि\-शृं\-ग\-वेरं}}{\Afootnote{K H; \skt{नाग\-र\-पि\-प्प\-लीभ्यां}\ A}}\skt{ द्राक्षाञ्च सर्प्पिर्मधुना }\edtext{\skt{लिहेच्च}}{\Afootnote{K H; \skt{वलिह्या}\ A}}\skt{ । }\pend
\pstart\versepar \skt{}\edtext{\skt{द्राक्षां}}{\Afootnote{K H; \skt{त् । द्राक्षां}\ A}}\skt{ }\edtext{\skt{सितां}}{\Afootnote{A K; \skt{सिता}\ H}}\skt{ मागधिकाञ्च तुल्यां सशृंगवेरं मधुकं तुगाञ्च ॥}\pend

\pstart\versepar \ledsidenote{\passref{6.52.18}}\edtext{\skt{सर्पिमधुभ्यां}}{\Afootnote{K; \emph{om.}\ A H\textsuperscript{ac}; \skt{सर्प्पिर्म्म्धुभ्यां}\ H\textsuperscript{pc}}}\skt{ }\edtext{\skt{विलिहेत्समांशां}}{\Afootnote{K H\textsuperscript{pc}; \skt{लिह्या\-द्घृ\-त\-क्षौ\-द्र\-युतां समांशां}\ A; \skt{सर्प्पिर्म्म्धुभ्यां विलिहेत्समांशां}\ H\textsuperscript{ac}}}\skt{ }\edtext{\skt{क्षौद्रेण तुल्यं}}{\Afootnote{K; \skt{सितोपलां वा}\ A; \emph{om.}\ H\textsuperscript{ac}; \skt{क्षौ}\ H\textsuperscript{pc}}}\skt{ }\edtext{\skt{मरिचं}}{\Afootnote{K; \skt{मरि\-चां\-श\-यु\-क्ताम् । धात्रीकणाविश्व}\ A; \skt{क्षौद्रे\-ण\-तु\-ल्य\-म्म्रिचं}\ H\textsuperscript{ac}; \skt{द्रेणतुल्यम्म्रिचं}\ H\textsuperscript{pc}}}\skt{ }\edtext{\skt{सिताञ्च}}{\Afootnote{K H; \skt{सितोपलाश्च}\ A}}\skt{ । }\pend
\pstart\versepar \skt{संचूर्ण्ण्य मण्डेन पिबेच्च }\edtext{\skt{द ध्नोः}}{\Afootnote{K H\textsuperscript{ac}; \skt{दध्नः}\ A H\textsuperscript{pc}}}\skt{ }\edtext{\skt{हरेण वः}}{\Afootnote{K H; \emph{om.}\ A}}\skt{ }\edtext{\skt{पिप्पलिकाञ्च}}{\Afootnote{K H; \emph{om.}\ A}}\skt{ }\edtext{\skt{तुल्यां}}{\Afootnote{K H; \emph{om.}\ A}}\skt{ ॥}\pend

\pstart\versepar \ledsidenote{\passref{6.52.19}}\edtext{\skt{दध्ना}}{\Afootnote{K; \skt{हरेणुकां मागधिकां च तुल्यां दध्ना}\ A; \skt{… दध्ना}\ H}}\skt{ }\edtext{\skt{लिहे\-त्का\-स\-क\-र\-ग्नि\-मग्रः}}{\Afootnote{K; \skt{पिबे\-त्का\-स\-ग\-दा\-भि\-भूतः ।}\ A; \skt{लिहे\-त्का\-स\-रु\-जा\-भि\-भूतः}\ H\textsuperscript{ac} H\textsuperscript{pc}}}\skt{ }\pend
\pstart\versepar \skt{उभे हरिद्रे }\edtext{\skt{सुरदारुशुण्ठी}}{\Afootnote{K; \skt{सुरदारुशुण्ठीं}\ A H}}\skt{ । गायत्रिसारञ्च पिबेत्समांशं}\pend

\pstart\versepar \ledsidenote{\passref{6.52.20}}\edtext{\skt{वस्तस्य}}{\Afootnote{K; \skt{बस्तस्य}\ A H}}\skt{ }\edtext{\skt{मूत्रेण}}{\Afootnote{A K; \skt{मूत्रे}\ H}}\skt{ }\edtext{\skt{सुखाम्वुना वा}}{\Afootnote{K; \skt{सुखाम्बुना वा}\ A H}}\skt{ }\edtext{\skt{दन्तीद्रवन्ती च}}{\Afootnote{K; \skt{दन्तीं द्रवन्तीं च}\ A; \skt{दन्ती द्रवन्तीञ्च}\ H}}\skt{ }\edtext{\skt{सतिल्वकानां}}{\Afootnote{K H; \skt{सतिल्वकाख्याम् ।}\ A}}\skt{ }\pend
\pstart\versepar \skt{भृष्टानि }\edtext{\skt{सर्पीः}}{\Afootnote{K H; \skt{सर्पीं}\ A}}\skt{ }\edtext{\skt{ष्वथ}}{\Afootnote{K H; \skt{ष्यथ}\ A}}\skt{ बादराणि ॥ }\edtext{\skt{खादेत्पलांशानि}}{\Afootnote{K H; \skt{खादेत्पलाशानि}\ A}}\skt{ ससैन्धवानि}\pend

\pstart\versepar \ledsidenote{\passref{6.52.21}}\edtext{\skt{हिङ्गोः}}{\Afootnote{K H; \skt{कोलप्रमाणं प्र}\ A}}\skt{ }\edtext{\skt{पिबेत्कोलसमं}}{\Afootnote{K H; \skt{पिबे}\ A}}\skt{ }\edtext{\skt{हिताशी ।सौ\-वी\-र\-के\-णा\-म्ल\-र\-सेन}}{\Afootnote{K H; \skt{द्धि हिङ्गु सौवी\-र\-के\-णा\-म्ल\-र\-सेन}\ A}}\skt{ वापि । }\pend
\pstart\versepar \skt{क्षौद्रेण }\edtext{\skt{लिह्यात्मरिचानि}}{\Afootnote{K H; \skt{लिह्यान्मरिचानि}\ A}}\skt{ }\edtext{\skt{चा । भार्गी\-व\-चा\-हिं\-गु\-कृता च}}{\Afootnote{K; \skt{वापि भार्गी\-व\-चा\-हि\-ङ्गु\-कृता च}\ A; \skt{चापि । भार्ग्गी\-व\-चा\-हि\-ङ्गु\-कृता तु}\ H}}\skt{ }\edtext{\skt{वर्तिः}}{\Afootnote{A K H\textsuperscript{ac}; \skt{वर्त्ति}\ H\textsuperscript{pc}}}\pend

\pstart \ledsidenote{\passref{6.52.22}}\edtext{\skt{धूमे}}{\Afootnote{A K; \skt{र्धूमे}\ H}}\skt{ प्रशस्ता }\edtext{\skt{घृतसंप्रयुक्ता}}{\Afootnote{K; \skt{घृतसंप्रयुक्ता वेणु\-त्व\-गे\-ला\-ल\-वणैः कृता वा । मुस्ते\-ङ्गु\-दी\-त्व\-ङ्म\-धु\-का\-ह्व\-मां\-सी\-म\-नः\-शि\-ला\-लै\-श्छ\-ग\-ला\-म्बु\-पिष्टैः ॥}\ A; \skt{घृतसंप्रयुक्ता ॥ …}\ H}}\skt{ ।}\pend

\pstart \ledsidenote{\passref{6.52.23}}\edtext{\skt{पिबेच्च}}{\Afootnote{K H; \skt{विधाय वर्तीश्च पयोऽनुपानं धूमं पिबे\-द्वा\-त\-ब\-ला\-स\-कासी । पिबेच्च}\ A}}\skt{ }\edtext{\skt{शीधुं}}{\Afootnote{K H; \skt{सीधुं}\ A}}\skt{ }\edtext{\skt{मरिचान्वितं वा}}{\Afootnote{K; \skt{मरिचान्वितं वा तेनाशु कासं जयति प्रसह्य ॥}\ A; \skt{मरिचान्वितम्वा … ।}\ H}}\pend

\pstart\versepar \ledsidenote{\passref{6.52.24}}\edtext{\skt{मञ्जि\-ष्ठ\-शु\-ण्ठी\-ज\-ल\-मृ\-ति\-काभिः}}{\Afootnote{K H; \skt{द्राक्षा\-म्बु\-म\-ञ्जि\-ष्ठ\-पु\-रा\-ह्व\-याभिः}\ A}}\skt{ । क्षीरं }\edtext{\edtext{\skt{श्रित\-म्मा\-क्षि\-क\-सं\-यु\-तम्वा}}{\lemma{\skt{श्रितम्}}\Afootnote{K; \skt{शृतं}\ A H}}}{\lemma{\skt{माक्षि\-क\-सं\-यु\-तम्वा}}\Afootnote{K H; \skt{माक्षि\-क\-सं\-प्र\-यु\-क्तम् ।}\ A}}\skt{ }\pend
\pstart\versepar \skt{}\edtext{\skt{निदि\-ग्धि\-का\-मू\-ल\-स\-मां\-श\-सिद्धां}}{\Afootnote{K H; \skt{निदि\-ग्धि\-का\-ना\-ग\-र\-पि\-प्प\-लीभिः}\ A}}\skt{ । खादेच्च }\edtext{\skt{मुद्गां}}{\Afootnote{A K; \skt{मुद्गा}\ H}}\skt{ }\edtext{\skt{मरि\-चो\-प\-दं\-शाम्}}{\Afootnote{K; \skt{मधुना सुसिद्धान् ॥}\ A; \skt{न्मरि\-चो\-प\-दं\-शाम्}\ H}}\pend

\pstart \ledsidenote{\passref{6.52.25}}\skt{उत्कारिकां सर्पिषि }\edtext{\skt{नागराढ्यां}}{\Afootnote{K; \skt{नागराढ्यां पक्त्वा समू\-लै\-स्त्रु\-टि\-को\-ल\-पत्रैः । एभिर्निषेवेत कृतां च पेयां तन्वीं सुशीतां मधुना विमिश्राम् ॥}\ A; \skt{नागराढ्यां । …}\ H}}\skt{ ।}\pend

\pstart\versepar \ledsidenote{\passref{6.52.26}}\skt{यत्प्लीह्नि सर्पिर्विहितं षडङ्गं }\edtext{\skt{तद्वातकासं}}{\Afootnote{H; \skt{तद्वातकासं}\ A K}}\skt{ जयति }\edtext{\skt{प्रसह्यं}}{\Afootnote{K H; \skt{प्रसह्य ।}\ A}}\skt{ । }\pend
\pstart\versepar \skt{}\edtext{\skt{विदा\-रि\-ग\-न्धा\-दि\-ग\-ण\-श्रितं वा}}{\Afootnote{K H; \skt{विदा\-रि\-ग\-न्धा\-दि\-कृतं घृतं वा}\ A}}\skt{ }\edtext{\skt{रसेन}}{\Afootnote{K H; \skt{रसेन वा}\ A}}\skt{ वा }\edtext{\skt{यत्खलु}}{\Afootnote{K H; \skt{सकजेन पक्वम् ॥}\ A}}\skt{ }\edtext{\skt{कण्ठकार्याः}}{\Afootnote{K; \emph{om.}\ A; \skt{कणृकार्याः ।}\ H}}\skt{ ।}\pend

\pstart\versepar \ledsidenote{\passref{6.52.27}}\skt{विरेचनं }\edtext{\edtext{\skt{स्नैहिकमत्र}}{\lemma{\skt{स्नैहिकम्}}\Afootnote{K H; \skt{स्नैहिक}\ A}}}{\lemma{\skt{अत्र}}\Afootnote{K; \skt{मत्र}\ A; \skt{अक्त्र}\ H}}\skt{ }\edtext{\edtext{\skt{चोक्तमाथापनं}}{\lemma{\skt{चोक्तम्}}\Afootnote{K; \skt{चोक्त}\ A H}}}{\lemma{\skt{आथापनं}}\Afootnote{K; \skt{मास्थापनं}\ A H}}\skt{ चापि }\edtext{\skt{वदन्ति}}{\Afootnote{K H; \skt{अनुवा}\ A}}\skt{ }\edtext{\skt{पथ्यं}}{\Afootnote{K H; \skt{सनं च ।}\ A}}\skt{ । }\pend
\pstart\versepar \skt{}\edtext{\edtext{\skt{धूम\-म्पि\-बे\-त्स्नै\-हि\-क\-म\-प्र\-मत्तः}}{\lemma{\skt{पिबेत्स्नैहिकम्}}\Afootnote{H; \skt{पिबेत्स्नैहिक}\ A K}}}{\lemma{\skt{अप्रमत्तः}}\Afootnote{H; \skt{मप्रमत्तः}\ A K}}\skt{ पिबेत्सुखोष्णं घृतमेव चात्र ॥}\pend

\pstart\versepar \ledsidenote{\passref{6.52.28}}\skt{हिता }\edtext{\skt{यवाग्वश्च}}{\Afootnote{A K; \skt{यवाग्वः स्व}\ H}}\skt{ रसेषु सिद्धाः पयांसि }\edtext{\skt{लेह्याः}}{\Afootnote{K H; \skt{लेहाः}\ A}}\skt{ सघृतास्तथैव ॥ }\pend
\pstart\versepar \skt{प्रछर्दनं काय\-शि\-रो\-वि\-रे\-का\-स्त\-थैव धूमाः }\edtext{\skt{कवडग्रहाश्च}}{\lemma{\skt{कवडग्रहाश्}}\Afootnote{K H; \skt{कवलग्रहाश्}\ A}}\skt{ ।}\pend

\pstart\versepar \ledsidenote{\passref{6.52.29}}\edtext{\skt{उष्णाम्ल}}{\Afootnote{K H; \skt{उष्णाश्च}\ A}}\skt{ }\edtext{\skt{लेह्याः}}{\Afootnote{K H; \skt{लेहाः}\ A}}\skt{ }\edtext{\edtext{\skt{कटुकाश्च}}{\lemma{\skt{कटुकाश्}}\Afootnote{K H; \skt{कटुका}\ A}}}{\lemma{\skt{च}}\Afootnote{K H; \skt{नि}\ A}}\skt{ हन्युः कफम्विशेषेण विशोषणञ्च । }\pend
\pstart\versepar \skt{कटुत्रिकञ्चापि वदन्ति }\edtext{\skt{मुख्यं}}{\Afootnote{K H; \skt{पथ्यं}\ A}}\skt{ घृतं }\edtext{\skt{क्रिघ्नस्वरसे}}{\Afootnote{K; \skt{कृमिघ्नस्वरसे}\ A; \skt{क्रिमिघ्नस्वरसे}\ H}}\skt{ विपक्वं ।}\pend

\pstart \ledsidenote{\passref{6.52.30}}\edtext{\skt{निर्गु\-ण्डि\-प\-त्र\-स्व\-र\-सेन}}{\Afootnote{K H; \skt{निर्गु\-ण्डि\-प\-त्र\-स्व\-रसे च}\ A}}\skt{ }\edtext{\skt{सिद्धं}}{\Afootnote{K H; \skt{पक्वं}\ A}}\skt{ सर्पिः कफो\-त्थ\-म्वि\-नि\-हन्ति }\edtext{\skt{कासं}}{\Afootnote{K H; \skt{कासम् । पाठा\-वि\-ड\-व्यो\-ष\-वि\-ड\-ङ्ग\-सि\-न्धु\-त्रि\-क\-ण्ट\-रा\-स्ना\-हु\-त\-भु\-ग्ब\-लाभिः ॥}\ A}}\skt{ ।}\pend

\pstart \ledsidenote{\passref{6.52.32}}\edtext{\skt{विदा\-रि\-ग\-न्धो\-त्प\-ल\-सा\-रि\-वादी}}{\Afootnote{K H; \skt{श्वासा\-ग्नि\-सा\-द\-स्व\-र\-भे\-द\-भि\-न्ना\-न्नि\-ह\-न्त्यु\-दी\-र्णा\-नपि पञ्च कासान् । विदा\-रि\-ग\-न्धो\-त्प\-ल\-सा\-रि\-वा\-दीन्}\ A}}\skt{ }\edtext{\skt{निःक्वाथ्य}}{\Afootnote{K H; \skt{निष्क्वाथ्य}\ A}}\skt{ वर्गम्मधुरञ्च कृत्स्नं ।}\pend

\pstart\versepar \ledsidenote{\passref{6.52.33}}\edtext{\skt{घृत\-म्पि\-बे\-दि\-क्षु\-र\-सा\-म्बु\-दुग्धैः}}{\lemma{\skt{पिबे\-दि\-क्षु\-र\-सा\-म्बु\-दुग्धैः}}\Afootnote{K H; \skt{पचे\-दि\-क्षु\-र\-सा\-म्बु\-दुग्धैः}\ A}}\skt{ }\edtext{\skt{काकोलिवर्गञ्च}}{\Afootnote{K H; \skt{काकोलिवर्गे च}\ A}}\skt{ सशर्कञ्तत् । }\pend
\pstart\versepar \skt{प्रातः }\edtext{\skt{पिबेत्पित्तकृते तु}}{\Afootnote{K H; \skt{पिबेत्पित्तकृते च}\ A}}\skt{ कासे }\edtext{\skt{क्षतो}}{\Afootnote{K H; \skt{रति}\ A}}\skt{ }\edtext{\skt{हितो यः}}{\Afootnote{K H; \skt{प्रसूते}\ A}}\skt{ }\edtext{\skt{क्षतजश्च}}{\Afootnote{K H; \skt{क्षतजे च}\ A}}\skt{ }\edtext{\skt{कासः}}{\Afootnote{K H; \skt{कासे ॥}\ A}}\skt{ ।}\edtext{}{\lemma{\skt{सशर्कञ्तत् । प्रातः}}\Afootnote{K; \skt{सशर्करं तत् । प्रातः}\ A; \skt{सशर्क्करे तत् ॥ प्रातः}\ H}}\pend

\pstart\versepar \ledsidenote{\passref{6.52.34}}\edtext{\skt{खर्जू\-र\-मु\-स्ता\-म\-धुकं}}{\Afootnote{K H; \skt{खर्जू\-र\-भा\-र्गी\-म\-गधा}\ A}}\skt{ }\edtext{\skt{पियालं}}{\Afootnote{K H; \skt{प्रियाल}\ A}}\skt{ }\edtext{\skt{मधूलिका}}{\Afootnote{K H; \skt{मधूलिकै}\ A}}\skt{ }\edtext{\skt{पिप्प\-लि\-भा\-र्गि\-चूर्णं}}{\Afootnote{K; \skt{लामलकैः समांशैः । चूर्णं}\ A; \skt{पिप्प\-लि\-भा\-र्गि\-चूर्ण्णं ॥ …}\ H}}\skt{ ॥ }\pend
\pstart\versepar \skt{}\edtext{\skt{सर्पिः}}{\Afootnote{K H; \emph{om.}\ A}}\skt{ सिता }\edtext{\skt{माक्षिसंप्रयुक्तं}}{\Afootnote{K; \skt{क्षौद्र\-घृ\-त\-प्रा\-गाढं}\ A; \skt{माक्षि\-क\-स\-म्प्र\-युक्तं}\ H}}\skt{ }\edtext{\skt{त्रीं}}{\Afootnote{K; \skt{त्री}\ A H}}\skt{ }\edtext{\skt{हन्ति}}{\Afootnote{K; \skt{न्हन्ति}\ A H}}\skt{ कासा\-नु\-प\-यु\-ज्य\-मानं ।}\pend

\pstart \ledsidenote{\passref{6.52.34a}}\skt{मञ्जि\-ष्ठ\-मू\-र्वा\-ज\-न\-व\-ह्नि\-पाथां कृष्णां हरिद्राञ्च तथा विचूर्ण्य ॥}\pend

\pstart \ledsidenote{\passref{6.52.35}}\edtext{\skt{क्षौद्रेण}}{\Afootnote{K H; \skt{रक्ता\-ह\-रि\-द्रा\-ञ्ज\-न\-व\-ह्नि\-पा\-ठा\-मू\-र्वो\-प\-कुल्या विलिहेत्समांशाः । क्षौद्रेण}\ A}}\skt{ }\edtext{\skt{कासे}}{\Afootnote{A K; \skt{कासो}\ H}}\skt{ }\edtext{\skt{क्षतजो}}{\Afootnote{K H; \skt{क्षतजे}\ A}}\skt{ क्षयोत्थे पिबेद्घृतं चेक्षुरसे विपक्वं ।}\pend

\pstart \ledsidenote{\passref{6.52.37cd}}\skt{गुडोदकं वा क्वथितं }\edtext{\skt{पिबेत}}{\Afootnote{\skt{पिबेद्धि}\ A}}\skt{ क्षौद्रेण शीतं मरिचोपदंशं ॥}\pend

\pstart\versepar \ledsidenote{\passref{6.52.36}}\skt{चूर्णं }\edtext{\skt{पिबे\-दा\-म\-ल\-कस्य}}{\Afootnote{A K; \skt{पिबे\-दा\-म\-ल\-कस्य}\ H}}\skt{ }\edtext{\skt{चापि}}{\Afootnote{K H; \skt{वापि}\ A}}\skt{ क्षीरे }\edtext{\skt{विपक्वं}}{\Afootnote{K H; \skt{ण पक्वं}\ A}}\skt{ सघृतं }\edtext{\skt{हिताय}}{\Afootnote{K H; \skt{हिताशी ।}\ A}}\skt{ ॥ }\pend
\pstart\versepar \skt{चूर्णानि गोधूम }\edtext{\skt{यथोद्भवानि}}{\Afootnote{K H; \skt{यवोद्भवानि}\ A}}\skt{ }\edtext{\skt{काकोलिवर्गञ्च}}{\Afootnote{K; \skt{काकोलिवर्गश्च}\ A H}}\skt{ }\edtext{\skt{कृतः}}{\Afootnote{K; \skt{कृतः सुसूक्ष्मः ॥}\ A H}}\skt{ ।}\pend

\pstart \ledsidenote{\passref{6.52.37ab}}\edtext{\skt{षु}}{\Afootnote{K; \skt{कासेषु}\ A H}}\skt{ }\edtext{\skt{पेयस्तृषु}}{\Afootnote{K; \skt{पेयस्त्रिषु}\ A H}}\skt{ कासवद्भिः }\edtext{\skt{क्षौद्रेण}}{\Afootnote{K H; \skt{क्षीरेण}\ A}}\skt{ }\edtext{\skt{दुग्धेन}}{\Afootnote{K H; \skt{सक्षौद्र}\ A}}\skt{ घृतेन वापि ।}\pend

\pstart\versepar \ledsidenote{\passref{6.52.47}}\edtext{\skt{कुली\-र\-शु\-क्ती\-श्च\-ट\-कैण}}{\lemma{\skt{कुलीरशुक्तीश्}}\Afootnote{K H; \skt{कुलीरशुक्ती}\ A}}\skt{ }\edtext{\skt{लावान्निः}}{\Afootnote{K H; \skt{लावान्नि}\ A}}\skt{ }\edtext{\skt{क्वाथ्य}}{\Afootnote{K H; \skt{ष्क्वाथ्य}\ A}}\skt{ }\edtext{\edtext{\skt{वर्ग\-म्म\-धु\-रै\-स्त\-थान्यैः}}{\lemma{\skt{मधुरैस्}}\Afootnote{K H; \skt{मधुरं}\ A}}}{\lemma{\skt{तथान्यैः}}\Afootnote{K H; \skt{च कृत्स्नम् ।}\ A}}\skt{ । }\pend
\pstart\versepar \skt{}\edtext{\skt{पिबेद्घृतं }\edtext{\skt{तन्तु}}{\lemma{\skt{तन्}}\Afootnote{K H; \skt{उ}\ A}}}{\lemma{\skt{पिबेद्घृतं तु}}\Afootnote{K H; \skt{पचेद्घृतं तत्त्}\ A}}\skt{ }\edtext{\skt{निषेव्यमानं}}{\Afootnote{K H; \skt{निषेव्यमाणं}\ A}}\skt{ }\edtext{\skt{हन्यात्क्षयोत्थं}}{\Afootnote{K H; \skt{हन्यात्क्षतोत्थं}\ A}}\skt{ }\edtext{\skt{क्षतजञ्च}}{\Afootnote{K H; \skt{क्षयजं च}\ A}}\skt{ कासं ॥ }\edtext{\skt{शता\-व\-री\-ना\-ग\-व\-ला\-वि\-पक्वं}}{\Afootnote{K; \skt{शता\-व\-री\-ना\-ग\-ब\-ला\-वि\-पक्वं}\ A; \skt{शता\-व\-री\-ना\-ग\-व\-री\-वि\-पक्वं}\ H}}\skt{ घृतं }\edtext{\skt{विदध्याच्च}}{\Afootnote{K H; \skt{विधेयं च}\ A}}\skt{ हिताय }\edtext{\skt{तस्य}}{\Afootnote{K H; \skt{कासिनाम् ।}\ A}}\skt{ ॥}\pend

\pstart \ledsidenote{\passref{6.52.37a}}\skt{जन्तु\-घ्न\-दा\-रु\-त्रि\-फ\-लो\-ष\-ण\-त्रयं सपद्मकं }\edtext{\edtext{\skt{संभृ\-त\-मे\-क\-त\-स्ततः}}{\lemma{\skt{संभृतम्}}\Afootnote{K; \skt{संभृत}\ H}}}{\lemma{\skt{एकतस्ततः}}\Afootnote{K; \skt{मेकतस्ततः ।}\ H}}\skt{ ।}\pend

\pstart \ledsidenote{\passref{6.52.37b}}\edtext{\skt{सिता\-स\-मां\-शै\-र\-व\-लि\-ह्य\-मानो ः}}{\lemma{\skt{अवलिह्यमानो ः}}\Afootnote{K; \skt{अवलिह्यमानो}\ H}}\skt{ }\edtext{\skt{सकासानुदति}}{\Afootnote{K; \skt{लहः सकासान्नुदति}\ H}}\skt{ प्रसह्यं ॥}\pend

\pstart \ledsidenote{\passref{6.52.37c}}\skt{शटी\-प\-ला\-ग्र\-न्थि\-क\-पु\-ष्क\-राग्नि }\edtext{\skt{चव्या\-त्म\-गु\-प्ता\-ख\-र\-शं\-ख\-पुष्पी}}{\Afootnote{K; \skt{चव्या\-त्म\-गु\-प्ता\-ख\-र\-श\-ङ्ख\-पुष्पी ।}\ H}}\skt{ ।}\pend

\pstart \ledsidenote{\passref{6.52.37d}}\skt{द्विपञ्चमूली च }\edtext{\skt{पलद्विभागे}}{\Afootnote{K; \skt{पल्द्विभागे}\ H}}\skt{ भागेन }\edtext{\skt{सम्भृत्य}}{\Afootnote{K; \skt{सन्भृत्य}\ H}}\skt{ विपाचयीत ।}\pend

\pstart \ledsidenote{\passref{6.52.37e}}\edtext{\edtext{\skt{यवा\-ढ\-क\-ञ्चा\-प्य\-भया}}{\lemma{\skt{चाप्य्}}\Afootnote{K; \skt{चा}\ H}}}{\lemma{\skt{अभया}}\Afootnote{K; \skt{प्यभया}\ H}}\skt{ शतञ्च वार्याढकैः }\edtext{\skt{पञ्चभिरं ह्रि}}{\Afootnote{K; \skt{पञ्चभिरं}\ H}}\skt{ शेषं ।}\pend

\pstart \ledsidenote{\passref{6.52.37f}}\skt{क्वाथ\-म्प\-चे\-च्छौ\-ण्ठि\-क\-तै\-ल\-भागे सर्पीं षिदत्वा कुडवोन्मितानि ।}\pend

\pstart \ledsidenote{\passref{6.52.37g}}\skt{गुडस्य भागार्द्धशतं विपक्वं }\edtext{\skt{शीतं}}{\Afootnote{K; \skt{यव\-स्थ\-ल्यां\-त\-न्नि\-धेयं सुशीतं ।}\ H}}\skt{ ।}\pend

\pstart \ledsidenote{\passref{6.52.37h}}\edtext{\skt{पलान्य}}{\Afootnote{K; \skt{पलान्यष्टौ}\ H}}\skt{ मधुनस्तत्र दत्वा त्वगेलशुण्ठी मरिचञ्च युञ्ज्यात् ।}\pend

\pstart \ledsidenote{\passref{6.52.37i}}\skt{द्वे द्वे चाद्या\-द\-व\-लि\-ह्याच्च }\edtext{\skt{लेहां}}{\Afootnote{K; \skt{लेहा}\ H}}\skt{ }\edtext{\skt{पथ्ये}}{\Afootnote{K; \skt{न्पथ्ये}\ H}}\skt{ पथ्याशी }\edtext{\skt{सर्वकासान्ज्ञ्घांसुः}}{\Afootnote{K; \skt{सर्व्व\-का\-सा\-न्जि\-घांसुः ।}\ H}}\skt{ ॥}\pend

\pstart \ledsidenote{\passref{6.52.37j}}\skt{श्वासं हिक्कां }\edtext{\skt{स्वरभेदं}}{\Afootnote{K; \skt{स्वरभेद}\ H}}\skt{ क्षयञ्च कार्ष्यं छर्दिं }\edtext{\skt{वातर्क्तं}}{\Afootnote{K; \skt{वातर्क्त}\ H}}\skt{ ज्वरञ्च ॥}\pend

\pstart \ledsidenote{\passref{6.52.37k}}\skt{हृत्पाण्डुरोगं श्वयथुंकामलां मूर्च्छां }\edtext{\skt{दोष्मूत्रकृच्छ्रार्त व}}{\Afootnote{K; \skt{दोषां मूत्र\-कृ\-च्छ्रा\-र्त्त\-वञ्}\ H}}\skt{ च ।}\pend

\pstart \ledsidenote{\passref{6.52.37l}}\skt{हन्यादगस्त्येन निषेवितो यं }\edtext{\skt{लेहारुजां}}{\Afootnote{K; \skt{लेहे रुजां}\ H}}\skt{ वै नुदति प्रसह्यमिति ॥}\pend

\pstart \ledsidenote{\passref{6.52.trailer}}\edtext{\skt{कासचि ॥ ॥ {\symfont ❈}}}{\Afootnote{\skt{इति सुश्रु\-त\-सं\-हि\-ता\-या\-मु\-त्त\-र\-त\-न्त्रा\-न्त\-र्गते काय\-चि\-कि\-त्सा\-तन्त्रे कासप्रतिषेधो नाम (च\-तु\-र्द\-शोऽ\-ध्यायः, आदितः) द्विप\-ञ्चा\-श\-त्त\-मोऽ\-ध्यायः ॥५२ ॥}\ A; \skt{कासचि ॥लोष्ठ ॥ {\symfont ❈} ॥}\ K; \skt{उत्तरतन्त्रे कासचिकित्सा ऊन\-प\-ञ्चा\-श\-त्तमः कायचिकित्सा चतुर्दश ॥}\ H}}\skt{ ॥}\pend

\endnumbering
\beginnumbering
\eledsection*{Adhyāya 53}
\pstart \ledsidenote{\passref{6.53.1}}\skt{अथातः }\edtext{\skt{स्वरो\-प\-घा\-त\-प्र\-ति\-षेधं}}{\Afootnote{K H; \skt{स्वर\-भे\-द\-प्र\-ति\-षे\-धम्}\ A}}\skt{ }\edtext{\skt{व्याख्यास्यामः}}{\Afootnote{H; \skt{अध्यायं व्याख्यास्यामः ॥}\ A; \emph{om.}\ K}}\skt{ ॥}\pend

\pstart\versepar \ledsidenote{\passref{6.53.3}}\edtext{\skt{अत्यु\-च्च\-श\-ब्द\-क\-र\-ण्य\-ध्य\-य\-ना\-भि\-घा\-त\-स\-न्दू\-षणैः}}{\Afootnote{\skt{अत्यु\-च्च\-भा\-ष\-ण\-वि\-षा\-ध्य\-य\-ना\-ति\-गी\-त\-शी\-ता\-दिभिः}\ A; \skt{अत्यु\-च्च\-श\-ब्द\-क\-र\-ण्य\-ध्य\-य\-ना\-भि\-घा\-त\-स\-न्दू\-षणैः}\ K; \skt{अत्यु\-च्च\-श\-ब्द\-क\-र\-णा\-ध्य\-य\-ना\-भि\-घा\-त\-स\-न्दू\-षणैः}\ H}}\skt{ प्रकुपिताः }\edtext{\skt{पठनादयस्तु}}{\Afootnote{K H; \skt{पवनादयस्तु ।}\ A}}\skt{ । }\pend
\pstart\versepar \skt{}\edtext{\skt{ते}}{\Afootnote{K H; \skt{स्रोतःसु ते}\ A}}\skt{ }\edtext{\skt{शब्दवाहीषु}}{\Afootnote{K H; \skt{स्वरवहेषु}\ A}}\skt{ }\edtext{\skt{सिरासुगताः}}{\Afootnote{K; \skt{गताः}\ A; \skt{सिसुगता}\ H}}\skt{ }\edtext{\skt{प्रतिष्ठं}}{\Afootnote{K; \skt{प्रतिष्ठां}\ A; \skt{प्रतिष्ठ}\ H}}\skt{ }\edtext{\skt{हन्युः}}{\Afootnote{A K; \skt{… …}\ H}}\skt{ स्वरं भवति चापि हि षड्विधः सः ।}\pend

\pstart\versepar \ledsidenote{\passref{6.53.4}}\skt{वातेन कृष्ण\-न\-य\-ना\-न\-न\-मू\-त्र\-वर्चा }\edtext{\skt{भिन्नस्वरं}}{\Afootnote{K H; \skt{भिन्नं शनैर्}\ A}}\skt{ वदति }\edtext{\skt{गर्द\-भ\-व\-त्ख\-रञ्च}}{\Afootnote{K H; \skt{गद्गदवत्स्वरं च ।}\ A}}\skt{ । }\pend
\pstart\versepar \skt{पित्तेन }\edtext{\skt{पीत\-न\-य\-ना\-न\-न\-मू\-त्र\-वर्चा}}{\Afootnote{K; \skt{पीत\-व\-द\-ना\-क्षि\-पु\-री\-ष\-मूत्रो}\ A; \skt{पीत\-न\-य\-ना\-न\-न\-मू\-त्र\-वर्च्छा}\ H}}\skt{ }\edtext{\skt{ब्रूयात्गलैन}}{\Afootnote{K; \skt{ब्रूयाद्गलेन}\ A; \skt{ब्रूयाद्भलेन}\ H}}\skt{ }\edtext{\skt{स}}{\Afootnote{K H; \skt{प}\ A}}\skt{ }\edtext{\skt{च}}{\Afootnote{K; \skt{रि}\ A; \skt{ह}\ H}}\skt{ दाह\-स\-म\-न्वि\-तेन ।}\pend

\pstart\versepar \ledsidenote{\passref{6.53.5}}\edtext{\skt{ब्रूयात्कफेन}}{\Afootnote{K H; \skt{कृच्छ्रात्कफेन}\ A}}\skt{ सततं }\edtext{\skt{कफरुद्धकण्ठौ}}{\Afootnote{K; \skt{कफरुद्धकण्ठो}\ A H}}\skt{ मन्दं शनैर्वदति चापि दिवा विशेषात् । }\pend
\pstart\versepar \skt{सर्वात्मके भवति }\edtext{\skt{सर्व\-वि\-का\-र\-सं\-पत् । }\edtext{\edtext{\edtext{\skt{तञ्चा\-प्य\-सा\-ध्य\-मृष्यः}}{\lemma{\skt{चाप्य्}}\Afootnote{K; \emph{om.}\ A; \skt{चा}\ H}}}{\lemma{\skt{असाध्यम्}}\Afootnote{K; \skt{दव्य}\ A; \skt{प्यसाध्य}\ H}}}{\lemma{\skt{ऋष्यः}}\Afootnote{K; \skt{क्तता}\ A; \skt{मृष्यः}\ H}}}{\lemma{\skt{सर्व\-वि\-का\-र\-सं\-पत् । तञ्}}\Afootnote{K; \skt{सर्व\-वि\-का\-र\-संप}\ A; \skt{सर्व्व\-वि\-का\-र\-सं\-पत्तं}\ H}}\skt{ }\edtext{\edtext{\skt{स्वरभेदमाहुः}}{\lemma{\skt{स्वरभेदम्}}\Afootnote{K; \skt{च वच\-स\-स्त\-म\-सा\-ध्यम्}\ A; \skt{स्वरभेद}\ H}}}{\lemma{\skt{आहुः}}\Afootnote{A K; \skt{माहुः ॥}\ H}}\skt{ ।}\edtext{}{\lemma{\skt{विशेषात् । सर्वात्मके}}\Afootnote{K H; \skt{विशेषः । सर्वात्मके}\ A}}\pend

\pstart\versepar \ledsidenote{\passref{6.53.6}}\edtext{\skt{धूमायति}}{\Afootnote{K H; \skt{धूप्येत वा}\ A}}\skt{ }\edtext{\skt{क्षतकृते}}{\Afootnote{K H; \skt{क्क्षयकृते}\ A}}\skt{ }\edtext{\skt{क्षयमाप्नुयाच्च}}{\Afootnote{\skt{क्षयमाप्नुयाच्च}\ A K; \skt{… क्षयाच्च}\ H}}\skt{ । }\edtext{\skt{वागेष}}{\Afootnote{A K; \skt{वाशेष}\ H}}\skt{ }\edtext{\skt{चापि}}{\Afootnote{K H; \skt{चापि(वापि)}\ A}}\skt{ }\edtext{\skt{हत\-वा\-क्प\-रि\-व\-र्ज\-नीयः}}{\Afootnote{A K; \skt{हत\-वा\-क्प\-रि\-व\-र्ज्ज\-यीत ।}\ H}}\skt{ । }\pend
\pstart\versepar \skt{}\edtext{\edtext{\edtext{\skt{अन्त\-र्ग\-त\-स्व\-र\-म\-ल\-क्ष्य\-म\-वा\-क्प्र\-वीणं}}{\lemma{\skt{अन्त\-र्ग\-त\-स्व\-रम्}}\Afootnote{K; \skt{अन्तर्गलं स्वर}\ A; \skt{अन्तर्ग्गतस्वर}\ H}}}{\lemma{\skt{अलक्ष्यम्}}\Afootnote{K; \skt{मलक्ष्य}\ A; \skt{मलज्य}\ H}}}{\lemma{\skt{अवाक्प्रवीणं}}\Afootnote{K; \skt{पदं चिरेण}\ A; \skt{मवाक्प्रवीणं}\ H}}\skt{ }\edtext{\skt{मेदो\-च्छ्र\-या\-द्व\-दति}}{\Afootnote{K; \skt{मेद\-श्च\-या\-द्व\-दति}\ A; \skt{मेदो\-च्छ्र\-या\-द्व\-दति}\ H}}\skt{ }\edtext{\skt{दिग्ध\-ग\-ल\-स्तृ\-षालुः}}{\Afootnote{K H; \skt{दिग्ध\-ग\-लौ\-ष्ठ\-तालुः ॥}\ A}}\skt{ ।}\pend

\pstart\versepar \ledsidenote{\passref{6.53.8}}\edtext{\skt{स्निग्धां}}{\Afootnote{K H; \skt{स्निग्धा}\ A}}\skt{ }\edtext{\skt{स्वरा\-तु\-र\-न\-रा\-न\-व\-ब\-द्ध\-दोषां}}{\Afootnote{K H; \skt{न्स्वरा\-तु\-र\-न\-रा\-न\-प\-कृ\-ष्ट\-दो\-षान्}\ A}}\skt{ न्यायेन }\edtext{\skt{च्छर्द\-न\-वि\-रे\-च\-न\-ब\-स्ति\-भिश्तु}}{\Afootnote{K; \skt{तान्व\-म\-न\-रे\-च\-न\-ब\-स्ति\-भिश्च ।}\ A; \skt{छर्द्द\-न\-वि\-रे\-च\-न\-ब\-स्ति\-भिश्च ।}\ H}}\skt{ । }\pend
\pstart\versepar \skt{}\edtext{\skt{नस्या\-व\-पी\-ड\-मु\-ख\-शो\-ध\-न\-धू\-म\-लेहैः}}{\Afootnote{K H; \skt{नस्या\-व\-पी\-ड\-मु\-ख\-धा\-व\-न\-धू\-म\-लेहैः}\ A}}\skt{ }\edtext{\skt{संपादयेत}}{\Afootnote{K H; \skt{संपादयेच्च}\ A}}\skt{ विविधैः }\edtext{\skt{कवḑअग्रहैश्च}}{\lemma{\skt{कवḑअग्रहैश्}}\Afootnote{K H; \skt{कवलग्रहैश्}\ A}}\skt{ ।}\pend

\pstart \ledsidenote{\passref{6.53.9ab}}\skt{यः }\edtext{\skt{श्वास\-का\-स\-वि\-धि\-रा\-दित}}{\Afootnote{K H; \skt{श्वास\-का\-स\-वि\-धि\-रा\-दित}\ A}}\skt{ एव }\edtext{\skt{चोक्तं }\edtext{\edtext{\edtext{\skt{श्तम्चा\-प्य\-शे\-ष\-म\-व\-ता\-र\-यितुं}}{\lemma{\skt{चाप्य्}}\Afootnote{A K; \skt{चा}\ H}}}{\lemma{\skt{अशेषम्}}\Afootnote{A K; \skt{प्यशेष}\ H}}}{\lemma{\skt{अवतारयितुं}}\Afootnote{A K; \skt{मवतारयितुं}\ H}}}{\lemma{\skt{चोक्तं श्तम्}}\Afootnote{K; \skt{चोक्तस्तं}\ A H}}\skt{ यतेत ॥}\pend

\pstart\versepar \ledsidenote{\passref{6.53.10}}\skt{स्वरोपघाते }\edtext{\skt{निलजे}}{\Afootnote{K; \skt{ऽनिलजे}\ A H}}\skt{ }\edtext{\skt{भक्तोपरि}}{\Afootnote{K H; \skt{भुक्तोपरि}\ A}}\skt{ घृतं पिबेत् । }\pend
\pstart\versepar \skt{कास\-म\-र्द\-क\-वा\-र्त्ता\-कु\-मा\-र्क\-व\-स्व\-रसे }\edtext{\skt{युतं}}{\Afootnote{K; \skt{शृतम् ॥}\ A; \skt{र्युतं ॥}\ H}}\skt{ ॥}\edtext{}{\lemma{\skt{पिबेत् । कास\-म\-र्द\-क\-वा\-र्त्ता\-कु\-मा\-र्क\-व\-स्व\-रसे}}\Afootnote{K; \skt{पिबेत् । कास\-म\-र्द\-क\-वा\-ता\-र्क\-मा\-र्क\-व\-स्व\-रसे}\ A; \skt{पिबेत् । कास\-म\-र्द्द\-क\-वा\-र्त्ता\-कु\-मा\-र्क्क\-व\-स्व\-रसै}\ H}}\pend

\pstart\versepar \ledsidenote{\passref{6.53.11}}\edtext{\skt{सिद्धं}}{\Afootnote{K H; \skt{पीतं}\ A}}\skt{ घृतं }\edtext{\edtext{\skt{हन्त्यनिलं}}{\lemma{\skt{हन्त्य्}}\Afootnote{A K; \skt{हन्}\ H}}}{\lemma{\skt{अनिलं}}\Afootnote{A K; \skt{त्यनिलं}\ H}}\skt{ सिद्धं }\edtext{\skt{वार्कगले}}{\Afootnote{K; \skt{आर्तगले}\ A; \skt{चार्त्तगले}\ H}}\skt{ रसे । }\pend
\pstart\versepar \skt{}\edtext{\skt{यव\-क्षा\-रा\-ज\-मो\-दाभ्यां}}{\Afootnote{A K\textsuperscript{pc} H; \skt{यव\-आ\-क्षा\-रा\-ज\-मो\-दाभ्यां}\ K\textsuperscript{ac}}}\skt{ चित्र\-का\-म\-ल\-केषु वा ।}\pend

\pstart\versepar \ledsidenote{\passref{6.53.12}}\edtext{\skt{देव\-दा\-र्व्य\-ग्नि\-मभ्यां}}{\Afootnote{K; \skt{देव\-दा\-र्व\-ग्नि\-काभ्यां वा}\ A; \skt{देव\-दा\-र्व्य\-ग्नि\-म\-न्थाभ्यां}\ H}}\skt{ }\edtext{\edtext{\skt{सिद्धमाजं}}{\lemma{\skt{सिद्धम्}}\Afootnote{K; \skt{सिद्धा}\ A H}}}{\lemma{\skt{आजं}}\Afootnote{K; \skt{माजं}\ A H}}\skt{ समाक्षिकं । }\pend
\pstart\versepar \skt{सुखो\-द\-का\-नु\-पानो वा }\edtext{\skt{ससर्पिष्को}}{\Afootnote{A K; \skt{ससर्प्पिष्टो}\ H}}\skt{ }\edtext{\skt{गुडोदनः}}{\Afootnote{K; \skt{गुडौदनः ॥}\ A; \skt{गुडोदकः ॥}\ H}}\skt{ ।}\pend

\pstart\versepar \ledsidenote{\passref{6.53.13}}\edtext{\skt{क्षीरानुपानं}}{\Afootnote{A K; \skt{क्षीरान्नपानं}\ H}}\skt{ }\edtext{\skt{पैत्ते तु}}{\Afootnote{K H; \skt{पित्ते तु}\ A}}\skt{ पिबे\-त्स\-र्पि\-र\-त\-न्द्रितः । }\pend
\pstart\versepar \skt{}\edtext{\skt{अश्नुयाच्च}}{\Afootnote{K H; \skt{अश्नीयाच्च}\ A}}\skt{ ससर्पिष्कं यष्टी\-म\-धु\-क\-पा\-यसं ।}\pend

\pstart\versepar \ledsidenote{\passref{6.53.14}}\edtext{\skt{लिहे\-त्म\-धु\-र\-काण्यं वा}}{\Afootnote{K; \skt{लिह्या\-न्म\-धु\-र\-काणां वा}\ A; \skt{लिहे\-त्म\-धु\-र\-का\-नाम्वा}\ H}}\skt{ चूर्ण्णं }\edtext{\skt{मधुसमायुतं}}{\Afootnote{K H; \skt{मधु\-घृ\-ता\-प्लु\-तम् ।}\ A}}\skt{ । }\pend
\pstart\versepar \skt{शता\-व\-री\-चू\-र्ण\-योगं बला\-चू\-र्ण्ण\-म\-थापि वा ।}\pend

\pstart\versepar \ledsidenote{\passref{6.53.15}}\skt{पिबेत्कटूनि मूत्रेण कफजे स्वरसं क्षये । }\pend
\pstart\versepar \skt{}\edtext{\skt{लिहेद्वा}}{\Afootnote{K H; \skt{लिह्याद्वा}\ A}}\skt{ मधुतैलाभ्यां भुक्त्वा खादेत्कटूनि वा ॥}\pend

\pstart\versepar \ledsidenote{\passref{6.53.16}}\edtext{\skt{स्वरोपघातै}}{\Afootnote{K; \skt{स्वरोपघाते}\ A H}}\skt{ मेदोजे }\edtext{\skt{कफ\-व\-द्वि\-धि\-रि\-ष्यते}}{\Afootnote{A K; \emph{om.}\ H}}\skt{ ॥ }\pend
\pstart\versepar \skt{}\edtext{\skt{रक्तजे}}{\Afootnote{K; \skt{सर्वजे}\ A; \skt{… ।}\ H}}\skt{ क्षयजे }\edtext{\skt{वापि}}{\Afootnote{K H; \skt{चापि}\ A}}\skt{ }\edtext{\skt{प्रत्या\-ख्या\-या\-च\-रे\-त्क्रियां}}{\Afootnote{A K; \skt{प्रत्या\-क्षा\-या\-र\-त्क्रियां ।}\ H}}\skt{ ।}\pend

\pstart \ledsidenote{\passref{6.53.16a}}\edtext{\skt{सर्वजे}}{\Afootnote{K; \emph{om.}\ A; \skt{सरव्वजे}\ H}}\skt{ }\edtext{\skt{क्षतजे}}{\Afootnote{K H; \emph{om.}\ A}}\skt{ }\edtext{\skt{तद्व\-त्भि\-ष\-ग्वि\-द्या\-द्वि\-च\-क्षणः}}{\Afootnote{K H; \emph{om.}\ A}}\skt{ ॥}\pend

\pstart\versepar \ledsidenote{\passref{6.53.17}}\skt{शर्क\-रा\-म\-धु\-मि\-श्राणि }\edtext{\skt{श्रितानि}}{\Afootnote{K H; \skt{शृतानि}\ A}}\skt{ मधुरैः सह । }\pend
\pstart\versepar \skt{पिबेत्पयांसि }\edtext{\skt{यस्योच्चैर्वदतो}}{\Afootnote{\skt{यस्योच्चैर्वदतो}\ A; \skt{यस्यो च्चैर्वदतो}\ K; \skt{यस्यो द्वे वदतो}\ H}}\skt{ }\edtext{\skt{भिहतस्वर}}{\Afootnote{K; \skt{ऽभिहतः स्वरः ॥}\ A; \skt{ऽभिहतस्वर}\ H}}\skt{ }\edtext{\skt{इति}}{\Afootnote{K H; \emph{om.}\ A}}\skt{ ॥}\pend

\pstart \ledsidenote{\passref{6.53.trailer}}\edtext{\skt{कायचिकित्सा}}{\Afootnote{\skt{इति सुश्रु\-त\-सं\-हि\-ता\-या\-मु\-त्त\-र\-त\-न्त्रा\-न्त\-र्गते काय\-चि\-कि\-त्सा\-तन्त्रे छर्दिप्रतिषेधो नाम (प\-ञ्च\-द\-शोऽ\-ध्यायः, आदितः) त्रिप\-ञ्चा\-श\-त्त\-मोऽ\-ध्यायः ॥५३ ॥}\ A; \skt{…कायचिकित्सा ॥ १५ ॥ ॥ {\symfont ❈} ॥}\ K; \skt{उत्तरतन्त्रे पञ्चाशत्तमः कायचिकित्सायां स्वरो\-प\-घा\-त\-चि\-कि\-त्सा\-प\-ञ्च\-द\-श\-मो\-ध्यायः ॥}\ H}}\skt{ ॥ ॥}\pend

\endnumbering
\beginnumbering
\eledsection*{Adhyāya 24}
\pstart \ledsidenote{\passref{6.24.1}}\skt{अथातः }\edtext{\skt{प्रति\-श्या\-य\-प्र\-ति\-षे\-ध\-म\-ध्यायं व्याख्यास्यामः}}{\lemma{\skt{अध्यायं व्याख्यास्यामः}}\Afootnote{\skt{व्याख्यास्यामः ॥}\ A H; \skt{अध्यायं व्याख्यास्यामः ॥}\ K}}\skt{ ॥}\pend

\pstart\versepar \ledsidenote{\passref{6.24.3}}\edtext{\skt{नारी\-प्र\-सं\-ग\-च्छि\-र\-सोऽ\-भि\-घा\-ताद्}}{\lemma{\skt{नारी\-प्र\-सं\-ग\-च्छि\-रसो}}\Afootnote{K H; \skt{नारीप्रसङ्गः शिरसो}\ A}}\skt{ }\pend
\pstart\versepar \skt{}\edtext{\skt{धूमा\-द्र\-जा\-च्छी\-त\-म\-व\-श्य\-याच्च}}{\lemma{\skt{अवश्ययाच्च}}\Afootnote{K; \skt{मतिप्रतापः ।}\ A; \skt{अवश्य यच्च ।}\ H}}\skt{ । }\pend
\pstart\versepar \skt{}\edtext{\skt{विनि\-ग्र\-हा\-न्मू\-त्र\-पु\-री\-ष\-योश्च}}{\Afootnote{K H; \skt{संधारणं मूत्र\-पु\-री\-ष\-योश्च}\ A}}\skt{ }\pend
\pstart\versepar \skt{सद्यः }\edtext{\edtext{\skt{प्रति\-श्या\-य\-नि\-मि\-त्त\-मुक्तं}}{\lemma{\skt{प्रति\-श्या\-य\-नि\-मि\-त्तम्}}\Afootnote{K H; \skt{प्रति\-श्या\-य\-नि\-दान}\ A}}}{\lemma{\skt{उक्तं}}\Afootnote{K H; \skt{मुक्तम् ॥}\ A}}\skt{ ।}\edtext{}{\lemma{\skt{ऽभि\-घा\-ता\-द्धू\-मा\-द्र\-जा\-च्छी\-तम्}}\Afootnote{\skt{ऽभितापो धूमो रजः शीत}\ A; \skt{भिघा\-ता\-द्धू\-मा\-द्र\-जा\-च्छी\-तम्}\ K H}}\pend

\pstart\versepar \ledsidenote{\passref{6.24.4}}\skt{चयं गता मूर्धनि मारुतादयः }\pend
\pstart\versepar \skt{पृथक्समस्ताश्च तथैव शोणितं । }\pend
\pstart\versepar \skt{}\edtext{\skt{प्रकोप्यमाना}}{\Afootnote{K H; \skt{प्रकोप्यमाणा}\ A}}\skt{ विविधैः }\edtext{\skt{प्रकोपनैर्}}{\Afootnote{K; \skt{प्रकोपणै}\ A; \skt{प्रकापनैर्}\ H\textsuperscript{ac}; \skt{प्रकारणैर्}\ H\textsuperscript{pc}}}\skt{ }\pend
\pstart\versepar \skt{}\edtext{\skt{नॄणां}}{\Afootnote{K; \skt{र्नृणां}\ A; \skt{न्नृणां}\ H}}\skt{ प्रतिश्यायकरा भवंति हि ।}\pend

\pstart\versepar \ledsidenote{\passref{6.24.5}}\edtext{\skt{शिरो}}{\Afootnote{K H; \skt{शिरो}\ A}}\skt{ गुरुत्वं }\edtext{\skt{क्षवथुः}}{\Afootnote{K; \skt{क्षवथोः}\ A; \skt{क्षवथु}\ H}}\skt{ प्रवर्त्तनं }\pend
\pstart\versepar \skt{तथांगमर्दः । परि\-हृ\-ष्ट\-रो\-मता }\pend
\pstart\versepar \skt{}\edtext{\skt{उप\-द्र\-वा\-श्चा\-प्य\-परे}}{\lemma{\skt{अपरे}}\Afootnote{A K; \skt{अपरो}\ H}}\skt{ }\edtext{\skt{पृथग्विधाः}}{\Afootnote{A K\textsuperscript{pc}; \skt{थग्विधाः ।}\ K\textsuperscript{ac}; \skt{पृथग्विधा न्}\ H}}\skt{ । }\pend
\pstart\versepar \skt{}\edtext{\skt{नॄणा\-म्प्र\-ति\-श्या\-य\-पु\-रः\-सराः}}{\lemma{\skt{नॄणाम्}}\Afootnote{K; \skt{नृणां}\ A H}}\skt{ }\edtext{\skt{स्मृताः}}{\Afootnote{A K; \skt{स्मृता}\ H}}\skt{ ।}\pend

\pstart\versepar \ledsidenote{\passref{6.24.6}}\edtext{\skt{आनद्ध}}{\Afootnote{K\textsuperscript{pc} H; \skt{आनद्धा}\ A; \skt{आद्ध}\ K\textsuperscript{ac}}}\skt{ }\edtext{\skt{पिथिता}}{\Afootnote{K; \skt{पिहिता}\ A H}}\skt{ नासा }\edtext{\skt{तनु\-स्रा\-व\-प्र\-से\-वनी}}{\Afootnote{K; \skt{तनु\-स्रा\-व\-प्र\-व\-र्तिनी ।}\ A; \skt{तनु\-स्रा\-व\-प्र\-से\-किनी ।}\ H}}\skt{ । }\pend
\pstart\versepar \skt{}\edtext{\skt{गल\-ता\-ल्वो\-ष्ठ\-घो\-षश्च}}{\Afootnote{K H; \skt{गल\-ता\-ल्वो\-ष्ठ\-शो\-षश्च}\ A}}\skt{ निस्तोदः }\edtext{\skt{शंखयोस्तथा}}{\lemma{\skt{शंखयोस्}}\Afootnote{K; \skt{शङ्खयोस्}\ A H}}\skt{ ।}\pend

\pstart\versepar \ledsidenote{\passref{6.24.7}}\skt{स्वरोपघातश्च भवेत्प्रतिश्याये }\edtext{\skt{निलात्मके}}{\Afootnote{K H; \skt{ऽनिलात्मके ।}\ A}}\skt{ । }\pend
\pstart\versepar \skt{उष्णः सपीतकः स्रावो घ्राणात्स्रवति पैत्तिके ।}\pend

\pstart\versepar \ledsidenote{\passref{6.24.8}}\edtext{\skt{कृषो}}{\Afootnote{K H; \skt{कृशो}\ A}}\skt{ }\edtext{\skt{भिपाण्डूः}}{\Afootnote{K; \skt{ऽतिपाण्डुः}\ A; \skt{ऽभिपाण्डूः}\ H}}\skt{ संतप्तो }\edtext{\skt{भवे\-त्तृ\-ष्णा\-नि\-पि\-डीतः}}{\Afootnote{K H; \skt{भवे\-त्तृ\-ष्णा\-नि\-पी\-डि\-तः(\-तृ\-ष्णा\-भि\-पी\-डितः) ।}\ A}}\skt{ । }\pend
\pstart\versepar \skt{}\edtext{\skt{सधूममग्निं}}{\lemma{\skt{अग्निं}}\Afootnote{K H; \emph{om.}\ A}}\skt{ सहसा }\edtext{\skt{वमतीव च}}{\Afootnote{K H; \skt{वह्निं वमतीव च}\ A}}\skt{ मानवः ।}\pend

\pstart\versepar \ledsidenote{\passref{6.24.9}}\edtext{\skt{घ्राणोत्कफः}}{\Afootnote{K H; \skt{कफः}\ A}}\skt{ कफकृते }\edtext{\skt{शुक्लः}}{\Afootnote{K H; \skt{घ्राणाच्छुक्लः}\ A}}\skt{ }\edtext{\edtext{\skt{पाण्डुर्व्रजेद्घनः}}{\lemma{\skt{पाण्डुर्}}\Afootnote{K; \skt{शीतः}\ A; \skt{पाण्डुं}\ H}}}{\lemma{\skt{व्रजेद्घनः}}\Afootnote{K\textsuperscript{pc}; \skt{स्रवेन्मुहुः ।}\ A; \skt{व्रजेद्घानः ।}\ K\textsuperscript{ac}; \skt{व्रजेद्घनः ।}\ H}}\skt{ । }\pend
\pstart\versepar \skt{शुक्लावभासः शूनाक्षो }\edtext{\skt{भवे\-द्गु\-रु\-शि\-रो\-मुखः}}{\Afootnote{K H; \skt{भवे\-द्गु\-रु\-शि\-रो\-मुखः ॥}\ A}}\skt{ ।}\pend

\pstart\versepar \ledsidenote{\passref{6.24.10}}\edtext{\skt{गल\-ता\-ल्वो\-ष्ठ\-शि\-रसः}}{\Afootnote{K H; \skt{शिरो\-ग\-लौ\-ष्ठ\-ता\-लूनां}\ A}}\skt{ }\edtext{\edtext{\skt{कण्डू\-य\-न\-म\-तीव}}{\lemma{\skt{कण्डूयनम्}}\Afootnote{K H; \skt{कण्डूयन}\ A}}\skt{ च}}{\lemma{\skt{अतीव च}}\Afootnote{K H; \skt{मतीव च ।}\ A}}\skt{ । }\pend
\pstart\versepar \skt{भूत्वा भूत्वा प्रतिश्यायो यो }\edtext{\skt{कस्मा\-द्वि\-नि\-व\-र्त्तते}}{\Afootnote{K H; \skt{ऽक\-स्मा\-द्वि\-नि\-व\-र्तते ॥}\ A}}\skt{ ।}\pend

\pstart \ledsidenote{\passref{6.24.11}}\skt{संपक्वो }\edtext{\edtext{\skt{वाप्यपक्वो}}{\lemma{\skt{वाप्य्}}\Afootnote{K; \skt{वा}\ A; \skt{वा वि}\ H}}\skt{ वा}}{\lemma{\skt{अपक्वो वा}}\Afootnote{K; \skt{प्यपक्वो वा}\ A; \skt{पक्वो वा}\ H}}\skt{ स }\edtext{\skt{सर्वप्रभवो}}{\Afootnote{A K; \skt{सत्वप्रभवो}\ H}}\skt{ }\edtext{\skt{मतः}}{\Afootnote{K H; \skt{स्मृतः । लिङ्गानि चैव सर्वेषां पीनसानां च सर्वजे ॥}\ A}}\skt{ ।}\pend

\pstart \ledsidenote{\passref{6.24.14}}\edtext{\skt{प्रक्लिद्यते}}{\Afootnote{K H; \skt{कृमि\-मू\-र्ध\-वि\-का\-रेण समानं चास्य लक्षणम् । प्रक्लिद्यति}\ A}}\skt{ }\edtext{\edtext{\skt{पुनर्न्नासा}}{\lemma{\skt{पुनर्}}\Afootnote{K H; \skt{पुन}\ A}}}{\lemma{\skt{न्नासा}}\Afootnote{K H; \skt{र्नासा}\ A}}\skt{ }\edtext{\skt{पुनश्च}}{\Afootnote{K H; \skt{पुअनश्च}\ A}}\skt{ परिशुष्यति ।}\pend

\pstart\versepar \ledsidenote{\passref{6.24.15}}\edtext{\edtext{\skt{मुहुरानह्यते}}{\lemma{\skt{मुहुर्}}\Afootnote{K H; \skt{मुहु}\ A}}}{\lemma{\skt{आनह्यते}}\Afootnote{K H; \skt{रानह्यते}\ A}}\skt{ चापि }\edtext{\edtext{\skt{मुहुर्विव्रियते}}{\lemma{\skt{मुहुर्}}\Afootnote{K H; \skt{मुहु}\ A}}\skt{ थ}}{\lemma{\skt{विव्रियते वा}}\Afootnote{K; \skt{र्विव्रियते तथा}\ A; \skt{व्विप्रियते थवा}\ H}}\skt{ }\edtext{\skt{वा ।}}{\Afootnote{K H; \skt{।}\ A}}\skt{ }\pend
\pstart\versepar \skt{}\edtext{\skt{निश्वासो}}{\Afootnote{K; \skt{निःश्वासो}\ A; \skt{निश्वासो}\ H}}\skt{ }\edtext{\skt{वति\-दु\-र्ग\-न्ध\-स्त\-था\-ग\-न्धान्न}}{\lemma{\skt{वतिदुर्गन्धस्}}\Afootnote{K; \skt{च्छ्वासदौर्गन्ध्यं}\ A; \skt{वतिनिर्ग्गन्धस्}\ H}}\skt{ वेत्ति च ।}\pend

\pstart \ledsidenote{\passref{6.24.16}}\skt{एवं दुष्टप्रतिश्यायं }\edtext{\skt{जानी\-या\-त्स\-र्व\-मेव}}{\lemma{\skt{जानीयात्सर्वम्}}\Afootnote{K H; \skt{जानी\-या\-त्कृ\-च्छ्र\-सा\-ध\-नम् । सर्व}\ A}}\skt{ }\edtext{\skt{तु}}{\Afootnote{K H; \skt{प्रतिश्याया नर\-स्या\-प्र\-ति\-का\-रिणः ॥}\ A}}\skt{ ।}\pend

\pstart\versepar \ledsidenote{\passref{6.24.16-1}}\edtext{\skt{रक्तजे तु}}{\Afootnote{K H; \emph{om.}\ A}}\skt{ }\edtext{\skt{प्रतिश्याये}}{\Afootnote{K H; \emph{om.}\ A}}\skt{ }\edtext{\skt{रक्तास्रावः}}{\Afootnote{K H; \emph{om.}\ A}}\skt{ }\edtext{\skt{प्रवर्त्तते}}{\Afootnote{K H; \emph{om.}\ A}}\skt{ । }\pend
\pstart\versepar \skt{}\edtext{\skt{ताम्राक्षश्च}}{\Afootnote{K; \emph{om.}\ A; \skt{ताम्राक्षस्य}\ H}}\skt{ }\edtext{\edtext{\skt{भवे\-ज्जं\-तु\-रु\-रो\-घा\-त\-प्र\-पी\-डितः}}{\lemma{\skt{भवेज्जंतुर्}}\Afootnote{K H; \emph{om.}\ A}}}{\lemma{\skt{उरो\-घा\-त\-प्र\-पी\-डितः}}\Afootnote{K H; \emph{om.}\ A}}\skt{ ।}\pend

\pstart\versepar \ledsidenote{\passref{6.24.16-2}}\edtext{\skt{दुर्ग्ग\-न्धो\-च्छ्वा\-स\-व\-दनो}}{\Afootnote{K H; \emph{om.}\ A}}\skt{ }\edtext{\skt{गन्धानपि न}}{\Afootnote{K\textsuperscript{pc} H; \emph{om.}\ A; \skt{गन्धानपि}\ K\textsuperscript{ac}}}\skt{ }\edtext{\skt{वेत्ति च}}{\Afootnote{K H; \emph{om.}\ A}}\skt{ ॥ }\pend
\pstart\versepar \skt{}\edtext{\skt{मूर्च्छन्ति}}{\Afootnote{K H; \emph{om.}\ A}}\skt{ }\edtext{\edtext{\skt{क्रिमयश्चात्र}}{\lemma{\skt{क्रिमयश्}}\Afootnote{K H; \emph{om.}\ A}}}{\lemma{\skt{चात्र}}\Afootnote{K H; \emph{om.}\ A}}\skt{ }\edtext{\skt{स्वेता}}{\Afootnote{K H; \emph{om.}\ A}}\skt{ }\edtext{\edtext{\skt{स्निग्धाणवस्तथा}}{\lemma{\skt{स्निग्धाणवस्}}\Afootnote{K; \emph{om.}\ A; \skt{स्निग्धानवस्}\ H}}}{\lemma{\skt{तथा}}\Afootnote{K H; \emph{om.}\ A}}\skt{ ।}\pend

\pstart \ledsidenote{\passref{6.24.16-3}}\edtext{\skt{क्रिमि\-शी\-र्ष\-वि\-का\-रेण}}{\Afootnote{K H; \emph{om.}\ A}}\skt{ }\edtext{\edtext{\skt{समानञ्चास्य}}{\lemma{\skt{समानञ्}}\Afootnote{K H; \emph{om.}\ A}}}{\lemma{\skt{चास्य}}\Afootnote{K H; \emph{om.}\ A}}\skt{ }\edtext{\skt{लक्षणं}}{\Afootnote{K H; \emph{om.}\ A}}\skt{ ।}\pend

\pstart\versepar \ledsidenote{\passref{6.24.17}}\edtext{\edtext{\edtext{\skt{बाधि\-र्य\-मा\-न्ध्य\-म\-घ्रत्वं}}{\lemma{\skt{बाधिर्यम्}}\Afootnote{K H; \skt{कालेन रोगजनना जायन्ते दुष्टपीनसाः । बाधिर्य}\ A}}}{\lemma{\skt{आन्ध्यम्}}\Afootnote{K H; \skt{मान्ध्य}\ A}}}{\lemma{\skt{अघ्रत्वं}}\Afootnote{K H; \skt{मघ्राणं}\ A}}\skt{ }\edtext{\skt{घोरान्वा}}{\Afootnote{K H; \skt{घोरांश्च}\ A}}\skt{ }\edtext{\skt{नयनामयां}}{\Afootnote{K; \skt{नयनामया}\ A H}}\skt{ । }\pend
\pstart\versepar \skt{}\edtext{\skt{कार्या\-ग्नि\-सा\-द\-शो\-षान्वा}}{\Afootnote{K; \skt{न् । कासा\-ग्नि\-सा\-द\-शो\-फांश्च}\ A; \skt{न् । कार्या\-ग्नि\-सा\-द\-शो\-षान्वा}\ H}}\skt{ वृद्धाः कुर्वन्ति पीनसाः ॥}\pend

\pstart\versepar \ledsidenote{\passref{6.24.18}}\skt{नवं }\edtext{\edtext{\skt{प्रति\-श्या\-य\-म\-पास्य}}{\lemma{\skt{प्रतिश्यायम्}}\Afootnote{K H; \skt{प्रतिश्याय}\ A}}}{\lemma{\skt{अपास्य}}\Afootnote{K H; \skt{मपास्य}\ A}}\skt{ सर्वान् }\pend
\pstart\versepar \skt{उपा\-च\-रे\-त्स\-र्पिष एव पानैः । }\pend
\pstart\versepar \skt{}\edtext{\edtext{\edtext{\skt{स्वेदै\-र्वि\-चि\-त्रै\-र्व\-म\-नैश्च}}{\lemma{\skt{स्वेदैर्}}\Afootnote{K H; \skt{स्वेदै}\ A}}}{\lemma{\skt{विचित्रैर्}}\Afootnote{K H; \skt{र्विचित्रै}\ A}}}{\lemma{\skt{वमनैश्च}}\Afootnote{K H; \skt{र्वमनैश्च}\ A}}\skt{ युक्तैः }\pend
\pstart\versepar \skt{}\edtext{\skt{कालो\-प\-प\-न्नै\-र\-व\-पी\-ड\-नैश्च}}{\lemma{\skt{अवपीडनैश्च}}\Afootnote{A K; \skt{अवपीडितैश्च ॥}\ H}}\skt{ ॥ }\pend
\pstart\versepar \skt{}\edtext{\skt{शिरो}}{\Afootnote{K H\textsuperscript{pc}; \emph{om.}\ A H\textsuperscript{ac}}}\skt{ }\edtext{\skt{विरेकं}}{\Afootnote{K H\textsuperscript{pc}; \emph{om.}\ A H\textsuperscript{ac}}}\skt{ }\edtext{\skt{त्वपरे तु}}{\Afootnote{K H\textsuperscript{pc}; \emph{om.}\ A H\textsuperscript{ac}}}\skt{ केचिद् }\pend
\pstart\versepar \skt{आमेन }\edtext{\edtext{\skt{दातव्यमिति}}{\lemma{\skt{दातव्यम्}}\Afootnote{K H\textsuperscript{pc}; \emph{om.}\ A H\textsuperscript{ac}}}}{\lemma{\skt{इति}}\Afootnote{K H\textsuperscript{pc}; \emph{om.}\ A H\textsuperscript{ac}}}\skt{ }\edtext{\skt{ब्रुवन्ति}}{\Afootnote{K H\textsuperscript{pc}; \emph{om.}\ A H\textsuperscript{ac}}}\skt{ ।}\edtext{}{\lemma{\skt{सर्वा\-नु\-पा\-च\-रे\-त्स\-र्पिष}}\Afootnote{K H; \skt{सर्व\-मु\-पा\-च\-रे\-त्स\-र्पिष}\ A}}\edtext{}{\lemma{\skt{केचिदामेन}}\Afootnote{K H\textsuperscript{pc}; \emph{om.}\ A H\textsuperscript{ac}}}\pend

\pstart\versepar \ledsidenote{\passref{6.24.19}}\edtext{\skt{अपच्यमानस्य च}}{\Afootnote{K H; \skt{अपच्यमानस्य हि}\ A}}\skt{ पाचनार्थं }\pend
\pstart\versepar \skt{स्वेदो हितो }\edtext{\skt{म्ब्लं}}{\Afootnote{K; \skt{ऽम्लै}\ A; \skt{म्लं}\ H}}\skt{ }\edtext{\skt{विहिमञ्च}}{\Afootnote{K H; \skt{रहिमं च}\ A}}\skt{ भोज्यं । }\pend
\pstart\versepar \skt{निषे\-व्य\-मा\-ण\-म्प\-य\-सा\-र्द्र\-कम्वा }\pend
\pstart\versepar \skt{}\edtext{\skt{सम्पा\-च\-ये\-च्च\-क्षु\-वि\-का\-र\-योगैः}}{\Afootnote{K H; \skt{संपा\-च\-ये\-दि\-क्षु\-वि\-का\-र\-योगैः ॥}\ A}}\skt{ ।}\pend

\pstart\versepar \ledsidenote{\passref{6.24.20}}\skt{पक्वं }\edtext{\skt{कफ\-ञ्चा\-प्य\-व\-ल\-म्ब\-मानं}}{\lemma{\skt{कफञ्}}\Afootnote{K; \skt{घनं}\ A; \skt{फलञ्}\ H}}\skt{ }\edtext{\skt{शिरो}}{\Afootnote{K H; \skt{शिरो}\ A}}\skt{ }\edtext{\edtext{\skt{विरे\-कै\-र\-प\-क\-र्ष\-येच्च}}{\lemma{\skt{विरेकैर्}}\Afootnote{K H; \skt{विरेकै}\ A}}}{\lemma{\skt{अपकर्षयेच्च}}\Afootnote{K H; \skt{रप\-क\-र्ष\-ये\-त्तम् ।}\ A}}\skt{ । }\pend
\pstart\versepar \skt{}\edtext{\skt{विरे\-च\-ना\-स्था\-प\-न\-धू\-म\-पा\-नैर्}}{\Afootnote{K\textsuperscript{pc} H; \skt{विरे\-च\-ना\-स्था\-प\-न\-धू\-म\-पानै}\ A; \skt{विरे\-च\-ना\-स्था\-प\-न\-धू\-मा\-पा\-नैर्}\ K\textsuperscript{ac}}}\skt{ }\pend
\pstart\versepar \skt{}\edtext{\skt{अवेक्ष्य}}{\Afootnote{K H; \skt{रवेक्ष्य}\ A}}\skt{ }\edtext{\skt{दोषां}}{\Afootnote{K; \skt{दोषा}\ A H}}\skt{ }\edtext{\skt{कवलग्रहैश्च}}{\lemma{\skt{कवलग्रहैश्}}\Afootnote{K; \skt{न्कवलग्रहैश्}\ A; \skt{न्कवडग्रहैश्}\ H}}\skt{ ।}\pend

\pstart\versepar \ledsidenote{\passref{6.24.24}}\edtext{\skt{जयेद्विदित्वा}}{\Afootnote{K; \skt{बहु\-द्र\-वै\-र्वा\-त\-क\-फो\-प\-सृष्टं प्रच्छ\-र्द\-ये\-त्पी\-न\-सिनं वयःस्थम् । उपद्रवांश्चापि यथोपदेशं स्वैर्भे\-ष\-जै\-र्भो\-ज\-न\-सं\-वि\-धानैः । जयेद्विदित्वा}\ A; \skt{जयेद्विरि…}\ H}}\skt{ }\edtext{\skt{मृदुताञ्च}}{\lemma{\skt{च}}\Afootnote{K H; \skt{ग}\ A}}\skt{ }\edtext{\skt{तेषां}}{\Afootnote{K H; \skt{तेषु}\ A}}\skt{ }\pend
\pstart\versepar \skt{}\edtext{\skt{प्राग्लक्षणोक्तं}}{\Afootnote{K H; \skt{प्राग्लक्षणेषू}\ A}}\skt{ }\edtext{\edtext{\skt{विधिमादिशेच्च}}{\lemma{\skt{विधिम्}}\Afootnote{K H; \skt{क्तम}\ A}}}{\lemma{\skt{आदिशेच्च}}\Afootnote{K H; \skt{थादिशेच्च ॥}\ A}}\skt{ ॥}\pend

\pstart\versepar \ledsidenote{\passref{6.24.25}}\skt{वातिके तु प्रतिश्याये }\edtext{\edtext{\skt{पिबे\-त्स\-र्पि\-र्य\-था\-क्रमं}}{\lemma{\skt{पिबेत्सर्पिर्}}\Afootnote{K H; \skt{पिबेत्सर्पि}\ A}}}{\lemma{\skt{यथाक्रमं}}\Afootnote{K H; \skt{र्यथाक्रमम् ।}\ A}}\skt{ । }\pend
\pstart\versepar \skt{पञ्चभिर्ल्लवणैः सिद्धम्प्रथमेन गणेन च ।}\pend

\pstart \ledsidenote{\passref{6.24.26}}\edtext{\skt{पित्तरक्तोत्थयोः}}{\Afootnote{K H; \skt{नस्यादिषु विधिं कृत्स्न\-म\-वे\-क्षे\-ता\-र्दि\-ते\-रि\-तम् । पित्तरक्तोत्थयोः}\ A}}\skt{ पेयं }\edtext{\edtext{\skt{सर्पिर्मधुरकैः}}{\lemma{\skt{सर्पिर्}}\Afootnote{K H; \skt{सर्पि}\ A}}}{\lemma{\skt{मधुरकैः}}\Afootnote{K H; \skt{र्मधुरकैः}\ A}}\skt{ }\edtext{\skt{श्रितं}}{\Afootnote{K H; \skt{शृतम् ॥}\ A}}\skt{ ।}\pend

\pstart \ledsidenote{\passref{6.24.27ab}}\edtext{\skt{परि\-शे\-का\-न्प्र\-दे\-हांश्च}}{\Afootnote{H; \skt{परि\-षे\-का\-न्प्र\-दे\-हांश्च}\ A; \skt{परिशेकां प्रदेहांश्च}\ K}}\skt{ }\edtext{\skt{कुर्यादपि च}}{\Afootnote{A H; \skt{कुर्यादपि}\ K}}\skt{ }\edtext{\skt{शीतलान्}}{\Afootnote{A H; \emph{om.}\ K}}\pend

\pstart \ledsidenote{\passref{6.24.30cd}}\skt{कफजे सर्प्पिषा }\edtext{\skt{स्निग्धं}}{\Afootnote{A K; \skt{स्निग्ध}\ H}}\skt{ तिल\-मा\-ष\-वि\-प\-क्वया ।}\pend

\pstart \ledsidenote{\passref{6.24.31ab}}\skt{यवाग्वा वामयेद्वान्तः कफघ्नं }\edtext{\edtext{\skt{क्रममाचरेत्}}{\lemma{\skt{क्रमम्}}\Afootnote{K H; \skt{क्रम}\ A}}}{\lemma{\skt{आचरेत्}}\Afootnote{K H; \skt{माचरेत् ।}\ A}}\skt{ ।}\pend

\pstart\versepar \ledsidenote{\passref{6.24.34}}\skt{सर्पींषि }\edtext{\skt{कटुसिद्धानि}}{\Afootnote{K H; \skt{कटुतिक्तानि}\ A}}\skt{ }\edtext{\skt{तीक्ष्णाधूमाः}}{\Afootnote{K H; \skt{तीक्ष्णधूमाः}\ A}}\skt{ कटूनि च । }\pend
\pstart\versepar \skt{}\edtext{\edtext{\skt{भेष\-जा\-न्यु\-प\-यु\-क्तानि}}{\lemma{\skt{भेषजान्य्}}\Afootnote{K H; \skt{भेषजा}\ A}}}{\lemma{\skt{उपयुक्तानि}}\Afootnote{K; \skt{न्युपयुक्तानि}\ A; \skt{उपयुज्यानि}\ H}}\skt{ हन्युः सर्वप्रकोपजं ।}\pend

\pstart\versepar \ledsidenote{\passref{6.24.41}}\edtext{\skt{यथास्वं}}{\Afootnote{K H; \emph{om.}\ A}}\skt{ }\edtext{\edtext{\edtext{\skt{दोष\-श\-म\-नै\-स्तै\-ल\-प\-क्व\-म्प्र\-मा\-णतः}}{\lemma{\skt{दोषशमनैस्}}\Afootnote{K H; \emph{om.}\ A}}}{\lemma{\skt{तैलपक्वम्}}\Afootnote{K H; \emph{om.}\ A}}}{\lemma{\skt{प्रमाणतः}}\Afootnote{K H; \emph{om.}\ A}}\skt{ ॥ }\pend
\pstart\versepar \skt{}\edtext{\edtext{\skt{नस्त\-प्र\-यु\-क्त\-मु\-द्वृत्तां}}{\lemma{\skt{नस्तप्रयुक्तम्}}\Afootnote{K H; \skt{नस्ये प्रयुक्त}\ A}}}{\lemma{\skt{उद्वृत्तां}}\Afootnote{K H; \skt{मुद्रिक्ता}\ A}}\skt{ }\edtext{\skt{प्रति\-श्या\-या\-न\-पो\-हति}}{\Afootnote{K H; \skt{न्प्रति\-श्या\-या\-न्व्य\-पो\-हति । यथास्वं दोष\-श\-म\-नै\-स्तैलं कुर्याच्च यत्नतः ॥}\ A}}\skt{ ।}\pend

\pstart\versepar \ledsidenote{\passref{6.24.42}}\edtext{\skt{मूत्रयुक्ताः}}{\Afootnote{K H; \skt{समूत्रपित्ताश्}\ A}}\skt{ }\edtext{\skt{समुद्दिष्टाः}}{\Afootnote{K H; \skt{चोद्दिष्टाः}\ A}}\skt{ क्रियाः }\edtext{\skt{क्रिमिषु}}{\Afootnote{K H; \skt{कृमिषु}\ A}}\skt{ योजयेत् । }\pend
\pstart\versepar \skt{यापनार्थं }\edtext{\skt{क्रिमिघ्नानि}}{\Afootnote{K H; \skt{कृमिघ्नानि}\ A}}\skt{ }\edtext{\skt{भोजनानि च}}{\Afootnote{K H; \skt{भेषजानि च}\ A}}\skt{ }\edtext{\skt{बुद्धिमानिति}}{\Afootnote{K; \skt{बुद्धिमान् ॥}\ A; \skt{बुद्धिमानिति ॥ {\symfont ❈} ॥}\ H}}\skt{ ॥ ॥}\pend

\pstart \ledsidenote{\passref{6.24.trailer}}\edtext{\skt{कायचिकित्सायां}}{\Afootnote{\skt{इति सुश्रु\-त\-सं\-हि\-ता\-या\-मु\-त्त\-र\-त\-न्त्रा\-न्त\-र्गते शालाक्यतन्त्रे प्रति\-श्या\-य\-प्र\-ति\-षेधो नाम चतु\-र्विं\-श\-ति\-त\-मोऽ\-ध्यायः ॥२४ ॥}\ A; \skt{कायचिकित्सायां ॥ १६ ॥ ० ॥}\ K; \skt{उत्तरतन्त्रे एक\-प\-ञ्चा\-श\-त्तमः कायचिकित्सायां प्रति\-श्या\-य\-चि\-कित्सा षोडशमो ध्यायः ॥}\ H}}\skt{ ॥ १६ ॥}\pend

\endnumbering
\beginnumbering
\eledsection*{Adhyāya 54}
\pstart \ledsidenote{\passref{6.54.1}}\skt{अथातः }\edtext{\skt{कृइ\-मि\-प्र\-ति\-षेधं}}{\Afootnote{K H; \skt{कृमि\-रो\-ग\-प्र\-ति\-षे\-धम्}\ A}}\skt{ }\edtext{\skt{व्याख्यास्यामः}}{\Afootnote{H; \skt{अध्यायं व्याख्यास्यामः ॥}\ A; \emph{om.}\ K}}\skt{ ॥}\pend

\pstart \ledsidenote{\passref{6.54.3}}\edtext{\skt{अजीर्ण्णा}}{\Afootnote{H; \skt{अजी\-र्णा\-ध्य\-श\-ना\-सा\-त्म्य\-वि\-रु\-द्ध\-म\-लि\-ना\-शनैः । अव्या\-या\-म\-दि\-वा\-स्व\-प्न\-गु\-र्व\-ति\-स्नि\-ग्ध\-शी\-तलैः ॥}\ A; \skt{अजीर्ण्णा …}\ K}}\pend

\pstart \ledsidenote{\passref{6.54.3a}}\edtext{\skt{भोजी\-म\-धू\-रा\-म्व्ल\-नि\-त्यो\-द्र\-व\-प्रियः}}{\Afootnote{K; \emph{om.}\ A}}\skt{ }\edtext{\skt{पिष्त\-गु\-डो\-प\-योक्ता।}}{\lemma{\skt{पिष्त\-गु\-डो\-प\-योक्ता}}\Afootnote{K; \emph{om.}\ A}}\pend

\pstart \ledsidenote{\passref{6.54.3b}}\edtext{\skt{व्यायामजातानि}}{\Afootnote{K; \emph{om.}\ A}}\pend

\pstart\versepar \ledsidenote{\passref{6.54.7}}\edtext{\skt{विंशते}}{\Afootnote{K; \skt{विंशतेः}\ A H}}\skt{ }\edtext{\skt{कृइमिजातीनां}}{\Afootnote{K H; \skt{कृमिजातीनां}\ A}}\skt{ त्रिविधः संभवः स्मृतः । }\pend
\pstart\versepar \skt{}\edtext{\skt{पुरी\-श\-क\-फ\-र\-क्तानि}}{\Afootnote{K; \skt{पुरी\-ष\-क\-फ\-र\-क्तानि}\ A H}}\skt{ }\edtext{\skt{तेसां}}{\Afootnote{K H; \skt{तासां}\ A}}\skt{ वक्ष्यामि }\edtext{\skt{लक्षणं}}{\Afootnote{K H; \skt{विस्तरम् ॥}\ A}}\skt{ ।}\pend

\pstart\versepar \ledsidenote{\passref{6.54.8}}\edtext{\skt{आजौपाः}}{\Afootnote{K; \skt{अजवा}\ A; \skt{अजौपाः}\ H}}\skt{ }\edtext{\skt{पीबनाः}}{\Afootnote{K H; \skt{विजवाः}\ A}}\skt{ }\edtext{\skt{कम्पाः}}{\Afootnote{K H; \skt{किप्याश्}\ A}}\skt{ }\edtext{\skt{श्वेताः}}{\Afootnote{K; \skt{चिप्या}\ A; \skt{स्वेता}\ H}}\skt{ गण्डूपदास्तथा । }\pend
\pstart\versepar \skt{}\edtext{\skt{चुरवो}}{\Afootnote{K; \skt{चूरवो}\ A H}}\skt{ द्विमुखाश्चैव }\edtext{\edtext{\skt{विज्ञेयास्तु}}{\lemma{\skt{विज्ञेयास्}}\Afootnote{K H; \skt{ज्ञेयाः}\ A}}}{\lemma{\skt{तु}}\Afootnote{K H; \skt{सप्त}\ A}}\skt{ }\edtext{\skt{पुरीशजाः}}{\Afootnote{K; \skt{पुरीषजाः ॥}\ A H}}\skt{ ॥}\pend

\pstart\versepar \ledsidenote{\passref{6.54.9}}\edtext{\skt{श्वेताः}}{\Afootnote{A K; \skt{स्वेताः}\ H}}\skt{ सूक्ष्मा स्तुदन्त्येते गुदं प्रतिसरन्ति च । }\pend
\pstart\versepar \skt{}\edtext{\edtext{\skt{तेषामेवा}}{\lemma{\skt{तेषाम्}}\Afootnote{K; \skt{तेषा}\ A H}}}{\lemma{\skt{एवा}}\Afootnote{K; \skt{मेवा}\ A H}}\skt{ }\edtext{\skt{परे}}{\Afootnote{A K; \skt{परेः}\ H}}\skt{ पुच्छैः पृथवश्च भवंति हि ।}\pend

\pstart \ledsidenote{\passref{6.54.11cd}}\skt{शूला\-टो\-प\-श\-कृ\-द्भे\-द\-प\-क्तिना शकराश्च ते ।}\pend

\pstart\versepar \ledsidenote{\passref{6.54.12}}\edtext{\skt{द}}{\Afootnote{K; \skt{दर्भपुष्पा महापुष्पाः}\ A; \skt{दर्भपुष्यो महापुष्पा}\ H}}\skt{ }\edtext{\skt{प्रलू\-ता\-श्चि\-पि\-टा\-स्तथा}}{\lemma{\skt{प्रलू\-ता\-श्चि\-पि\-टास्}}\Afootnote{K; \skt{प्रलू\-ना\-श्चि\-पि\-टास्}\ A; \skt{प्रलू\-ता\-न्य\-श्चि\-पि\-टास्}\ H}}\skt{ । }\pend
\pstart\versepar \skt{पिपीलिका }\edtext{\skt{काररुहा}}{\Afootnote{K H\textsuperscript{pc}; \skt{दारु}\ A; \skt{कारुहा}\ H\textsuperscript{ac}}}\skt{ }\edtext{\skt{विज्ञेयाः}}{\Afootnote{K H; \skt{णाश्}\ A}}\skt{ }\edtext{\skt{कफसंभवाः}}{\Afootnote{K; \skt{च कफ\-को\-प\-स\-मु\-द्भवाः ॥}\ A; \skt{कफ सम्भवाः ॥}\ H}}\skt{ ।}\pend

\pstart\versepar \ledsidenote{\passref{6.54.14}}\edtext{\skt{मज्जा\-दो\-ने\-त्र\-ले\-ढा\-र\-स्ता\-लु\-श्रो\-त्र\-भु\-ज\-स्तथा}}{\lemma{\skt{मज्जा\-दो\-ने\-त्र\-ले\-ढा\-र\-स्ता\-लु\-श्रो\-त्र\-भु\-जस्}}\Afootnote{K; \skt{मज्जादा नेत्र\-ले\-ढा\-र\-स्ता\-लु\-श्रो\-त्र\-भु\-जस्}\ A; \skt{मज्जा\-दो\-ने\-त्र\-ले\-ढा\-र\-स्ता\-लु\-श्तो\-त्र\-भु\-जस्}\ H}}\skt{ । }\pend
\pstart\versepar \skt{शिरो\-हृ\-द्रो\-ग\-व\-म\-थु\-प्र\-ति\-श्या\-य\-क\-राश्च ते ।}\pend

\pstart\versepar \ledsidenote{\passref{6.54.15}}\edtext{\skt{कर्ण्ण\-रो\-म\-न\-खा\-दाश्च}}{\lemma{\skt{कर्ण्ण\-रो\-म\-न\-खा\-दाश्}}\Afootnote{K H; \skt{केश\-रो\-म\-न\-खा\-दाश्}\ A}}\skt{ दन्तादाः }\edtext{\skt{किक्किसास्तथा}}{\lemma{\skt{किक्किसास्}}\Afootnote{K H; \skt{किक्किशास्}\ A}}\skt{ । }\pend
\pstart\versepar \skt{कुष्ठजाः सपरीसर्पाः ज्ञेयाः }\edtext{\skt{शोणितसम्भवाः}}{\Afootnote{K; \skt{शोणितसंभवाः}\ A H}}\skt{ ।}\pend

\pstart\versepar \ledsidenote{\passref{6.54.16}}\skt{ते सरक्ताश्च कृष्णाश्च }\edtext{\edtext{\skt{तेशामन्ते}}{\lemma{\skt{तेशाम्}}\Afootnote{K H; \skt{स्निग्धाश्}\ A}}}{\lemma{\skt{अन्ते}}\Afootnote{K H; \skt{च पृ}\ A}}\skt{ }\edtext{\skt{च}}{\Afootnote{K H; \skt{थ}\ A}}\skt{ }\edtext{\skt{दुर्दृशः}}{\Afootnote{K; \skt{वस्तथा ।}\ A; \skt{दु ट्टषाः ।}\ H}}\skt{ ॥ }\pend
\pstart\versepar \skt{}\edtext{\skt{रक्ताधिष्ठानजां}}{\Afootnote{K; \skt{रक्ताधिष्ठानजा}\ A H}}\skt{ }\edtext{\skt{प्रायो}}{\Afootnote{K; \skt{न्प्रायो}\ A H}}\skt{ }\edtext{\skt{विका\-रा\-ञ्ज\-न\-यन्ति}}{\Afootnote{K; \skt{विका\-रा\-न्ज\-न\-यन्ति}\ A; \skt{विका\-रा\-ज\-न\-यन्तिः ॥}\ H}}\skt{ }\edtext{\skt{ह}}{\Afootnote{K; \skt{ते ॥}\ A; \emph{om.}\ H}}\skt{ ।}\pend

\pstart\versepar \ledsidenote{\passref{6.54.17}}\skt{माष\-पि\-ष्टा\-न्न\-वि\-द\-ल\-प\-र्ण्ण\-शाकैः पुरीषजाः । }\pend
\pstart\versepar \skt{}\edtext{\skt{मांस\-मा\-ष\-गु\-ड\-क्षी\-र\-द\-धि\-शुक्तैः}}{\Afootnote{K H; \skt{मांस\-मा\-ष\-गु\-ड\-क्षी\-र\-द\-धि\-तैलैः}\ A}}\skt{ कफोद्भवाः ।}\pend

\pstart \ledsidenote{\passref{6.54.18ab}}\skt{विरु\-द्धा\-जी\-र्ण्ण\-शा\-काद्यैः }\edtext{\skt{शोणिताभा}}{\Afootnote{K; \skt{शोणितोत्था}\ A; \skt{शोणितभा}\ H}}\skt{ भवन्ति }\edtext{\skt{च}}{\Afootnote{K H; \skt{हि ।}\ A}}\skt{ ।}\pend

\pstart \ledsidenote{\passref{6.54.20cd}}\edtext{\skt{एषामन्यतमां}}{\Afootnote{K; \skt{एषामन्यतमं}\ A; \skt{एषामन्यतम}\ H}}\skt{ }\edtext{\skt{वैद्यो}}{\Afootnote{K H; \skt{ज्ञात्वा}\ A}}\skt{ जिघांसुः }\edtext{\skt{स्निग्धमातुरां}}{\Afootnote{K; \skt{स्निग्धमातुरम् ॥}\ A; \skt{स्निग्धमात्तरम्॥}\ H}}\skt{ ।}\pend

\pstart\versepar \ledsidenote{\passref{6.54.21}}\edtext{\skt{सुर\-सा\-दि\-वि\-प\-वेन}}{\Afootnote{K; \skt{सुर\-सा\-दि\-वि\-प\-क्वेन}\ A H}}\skt{ }\edtext{\skt{सर्पि\-षा\-वा\-न्त\-मा\-दितः}}{\Afootnote{A K; \skt{सर्प्पि\-षा\-वा\-न्त\-मा\-दिशः ।}\ H}}\skt{ । }\pend
\pstart\versepar \skt{}\edtext{\skt{विरे\-च\-यै\-त्ती\-क्ष्ण\-त\-रै\-र्यो\-गै\-रा\-स्था\-प\-यीत}}{\Afootnote{K; \skt{विरे\-च\-ये\-त्ती\-क्ष्ण\-त\-रै\-र्यो\-गै\-रा\-स्था\-प\-येच्च}\ A; \skt{विरे\-च\-ये\-त्ती\-क्ष्ण\-त\-रै\-र्यो\-गै\-रा\-स्था\-प\-यीत}\ H}}\skt{ }\edtext{\skt{च}}{\Afootnote{K H; \skt{तम् ॥}\ A}}\skt{ ।}\pend

\pstart\versepar \ledsidenote{\passref{6.54.22}}\edtext{\skt{यव\-को\-ल\-कु\-ल\-त्थानां}}{\Afootnote{A K; \skt{यव\-को\-ल\-क्श\-ल\-त्थानां}\ H}}\skt{ सुरसादेर्ग }\edtext{\skt{च}}{\Afootnote{K; \skt{णस्य च ।}\ A H}}\skt{ । }\pend
\pstart\versepar \skt{}\edtext{\skt{विड\-ङ्ग\-स्ने\-ह\-युक्तौ न}}{\Afootnote{K; \skt{विड\-ङ्ग\-स्ने\-ह\-यु\-क्तेन}\ A H}}\skt{ क्वाथेन लवणेन च ।}\pend

\pstart \ledsidenote{\passref{6.54.23}}\edtext{\edtext{\edtext{\skt{निरू\-ढ\-म\-थ\-तै\-ले\-न्त\-त्सि\-द्धे\-ना\-नुवा}}{\lemma{\skt{निरूढम्}}\Afootnote{K; \skt{प्रत्यागते निरूहे तु नरं स्नातं सुखाम्बुना । युञ्ज्या\-त्कृ\-मि\-घ्नै\-र\-श\-नै\-स्ततः शीघ्रं भिषग्वरः ॥}\ A; \skt{निरु}\ H}}}{\lemma{\skt{अथतैलेन्}}\Afootnote{K; \emph{om.}\ A; \skt{ट्मथतैलेन्}\ H}}}{\lemma{\skt{तत्सिद्धेनानुवा}}\Afootnote{K H; \emph{om.}\ A}}\skt{ }\edtext{\skt{सयेत्}}{\Afootnote{K H; \emph{om.}\ A}}\pend

\pstart \ledsidenote{\passref{6.54.24cd}}\skt{ततः शिरी\-ष\-कि\-णि\-ही\-रसं }\edtext{\skt{क्षौद्रयुतं}}{\Afootnote{A K; \skt{क्षैद्रयुतं}\ H}}\skt{ पिबेत् ॥}\pend

\pstart \ledsidenote{\passref{6.54.25ab}}\edtext{\skt{क्रमूकः}}{\Afootnote{K H; \skt{केवूक}\ A}}\skt{ }\edtext{\skt{स्वरसम्वापि}}{\lemma{\skt{स्वरसम्}}\Afootnote{A K; \skt{सुरसं}\ H}}\skt{ }\edtext{\skt{पूर्व\-व\-त्ती\-क्ष\-ण\-भो\-जनः}}{\Afootnote{K; \skt{पूर्व\-व\-त्ती\-क्ष्ण\-भो\-जनः ।}\ A; \skt{पूर्व\-व\-त्ती\-क्ष\-त\-भो\-जनः ।}\ H}}\skt{ ।}\pend

\pstart \ledsidenote{\passref{6.54.26ab}}\skt{पारि\-भ\-द्र\-क\-प\-त्राणां }\edtext{\skt{पुर्वत्सुरसं}}{\Afootnote{K H; \skt{क्षौद्रेण स्वरसं}\ A}}\skt{ पिबेत् ।}\pend

\pstart \ledsidenote{\passref{6.54.25cd}}\edtext{\skt{पलाशबीजस्य}}{\Afootnote{K; \skt{पलाशबीजस्व}\ A; \skt{पलाशबीज}\ H}}\skt{ }\edtext{\skt{रसं}}{\Afootnote{A K; \skt{…}\ H}}\skt{ }\edtext{\skt{पिबेद्वा}}{\Afootnote{K H; \emph{om.}\ A}}\skt{ }\edtext{\skt{क्शौद्र}}{\Afootnote{K; \skt{क}\ A; \skt{क्शैद्र}\ H}}\skt{ }\edtext{\skt{संयुतं।}}{\lemma{\skt{संयुतं}}\Afootnote{K H; \skt{ल्कं वा तण्डुलाम्बुना ॥}\ A}}\pend

\pstart \ledsidenote{\passref{6.54.26a}}\skt{पृथग्वा }\edtext{\skt{सुरसादीनां}}{\Afootnote{K; \skt{सुरसादीनं}\ H}}\skt{ }\edtext{\skt{पत्तु\-र\-स्या\-थ\-वा\-रसं}}{\Afootnote{K; \skt{पतु\-र\-स्था\-वा\-र\-सम्}\ H}}\skt{ ।}\pend

\pstart\versepar \ledsidenote{\passref{6.54.27}}\edtext{\edtext{\skt{लिहेदश्च}}{\lemma{\skt{लिहेदश्}}\Afootnote{K H; \skt{लिह्याद}\ A}}}{\lemma{\skt{च}}\Afootnote{K H; \skt{श्व}\ A}}\skt{ }\edtext{\skt{शकृच्चूर्ण्णं}}{\Afootnote{A K; \skt{सकृच्चूर्णं}\ H}}\skt{ वैडंगं वा }\edtext{\skt{समिक्षिकम्}}{\Afootnote{K H; \skt{समाक्षिकम् ।}\ A}}\skt{ । }\pend
\pstart\versepar \skt{}\edtext{\edtext{\skt{पत्रै\-र्मू\-षि\-क\-पर्ण्या}}{\lemma{\skt{पत्रैर्}}\Afootnote{K H; \skt{पत्रै}\ A}}\skt{ वा}}{\lemma{\skt{मूषिकपर्ण्या वा}}\Afootnote{K H; \skt{र्मूषिकपर्ण्या वा}\ A}}\skt{ }\edtext{\skt{संपिष्टैः}}{\Afootnote{K; \skt{सुपिष्टैः}\ A H}}\skt{ पिष्टमिश्रितैः ।}\pend

\pstart\versepar \ledsidenote{\passref{6.54.28}}\edtext{\skt{खादेत्पूपलिकां}}{\Afootnote{K H; \skt{खादेत्पूपलिकाः}\ A}}\skt{ पक्वा }\edtext{\skt{धान्याम्व्लञ्च}}{\Afootnote{K; \skt{धान्याम्लं च}\ A; \skt{धान्याम्लाञ्च}\ H}}\skt{ }\edtext{\skt{पिबेदनु}}{\Afootnote{K H; \skt{पिबेदनु ।}\ A}}\skt{ । }\pend
\pstart\versepar \skt{}\edtext{\edtext{\skt{सुरसाद्यैर्गणैः}}{\lemma{\skt{सुरसाद्यैर्}}\Afootnote{K; \skt{सुरसादि}\ A; \skt{सुरसाधैर्}\ H}}}{\lemma{\skt{गणैः}}\Afootnote{K; \skt{गणे}\ A; \skt{गनणैः}\ H}}\skt{ पक्वं }\edtext{\edtext{\skt{तैल\-न्त\-त्पा\-न\-मि\-ष्यते}}{\lemma{\skt{तत्पानम्}}\Afootnote{K H; \skt{वा पान}\ A}}}{\lemma{\skt{इष्यते}}\Afootnote{K H; \skt{मिष्यते ॥}\ A}}\skt{ ।}\pend

\pstart\versepar \ledsidenote{\passref{6.54.29}}\edtext{\skt{विडं\-ग\-चू\-र्ण्ण\-मि\-स्रैर्वा}}{\Afootnote{K; \skt{विड\-ङ्ग\-चू\-र्ण\-यु\-क्तैर्वा}\ A; \skt{विड\-ङ्ग\-चू\-र्ण\-मि\-शैर्व्वा}\ H}}\skt{ }\edtext{\skt{पिष्टैर्भक्षां}}{\Afootnote{K; \skt{पिष्टैर्भक्ष्यांस्}\ A H}}\skt{ प्रयोजयेत् । }\pend
\pstart\versepar \skt{तत्क\-षा\-य\-प्र\-पी\-तानां तिलानां }\edtext{\skt{तैल्पाचितां}}{\Afootnote{K H; \skt{स्नेहमेव वा ॥}\ A}}\skt{ ।}\edtext{}{\lemma{\skt{प्रयोजयेत् । तत्क\-षा\-य\-प्र\-पी\-तानां}}\Afootnote{K H; \skt{तु कारयेत् । तत्क\-षा\-य\-प्र\-पी\-तानां}\ A}}\pend

\pstart\versepar \ledsidenote{\passref{6.54.30}}\skt{श्वाविधः }\edtext{\skt{सकृतश्चूर्ण्णं}}{\lemma{\skt{सकृतश्}}\Afootnote{K; \skt{शकृतश्}\ A H}}\skt{ सप्तकृत्वः }\edtext{\skt{सुभावितः}}{\Afootnote{K; \skt{सुभावितम् ।}\ A; \skt{सुभाविमः।}\ H}}\skt{ । }\pend
\pstart\versepar \skt{विडंगानां कषायेण त्रैफलेन }\edtext{\skt{रसेन}}{\Afootnote{K H; \skt{तथैव}\ A}}\skt{ च ।}\pend

\pstart \ledsidenote{\passref{6.54.31ab}}\edtext{\skt{क्षौद्रेण}}{\Afootnote{A K; \skt{क्षैद्रेण}\ H}}\skt{ }\edtext{\edtext{\skt{लीढ्वा\-नु\-पि\-बे\-द्र\-स\-मा\-म\-लकं}}{\lemma{\skt{लीढ्वा\-नु\-पि\-बे\-द्र\-स\-माम्}}\Afootnote{K; \skt{लीढ्वा\-नु\-पि\-बे\-द्र\-समा}\ A; \skt{लीढ्वा\-नु\-पि\-बे\-द्र\-समा}\ H}}}{\lemma{\skt{अलकं}}\Afootnote{K; \skt{मलको}\ A; \skt{म्लकम्}\ H}}\skt{ }\edtext{\skt{शुभं}}{\Afootnote{K H; \skt{द्भवम् ।}\ A}}\skt{ ।}\pend

\pstart \ledsidenote{\passref{6.54.32cd}}\edtext{\skt{पिबेद्वा}}{\Afootnote{K H; \skt{पिबेद्वा}\ A}}\skt{ }\edtext{\edtext{\skt{पिप्प\-ली\-चू\-र्ण्ण\-म\-ज\-मू\-त्रेण}}{\lemma{\skt{पिप्पलीचूर्ण्णम्}}\Afootnote{K; \skt{पिप्पलीमूल}\ A; \skt{पिप्पलीचूर्न}\ H}}}{\lemma{\skt{अजमूत्रेण}}\Afootnote{K; \skt{मजामूत्रेण}\ A; \skt{म जमूत्रेण}\ H}}\skt{ संयुतं।}\pend

\pstart \ledsidenote{\passref{6.54.31cd}}\edtext{\skt{अक्षाभयारसे}}{\Afootnote{K; \skt{अक्षाभयारसं}\ A; \skt{अक्षाभयोरसे}\ H}}\skt{ वापि }\edtext{\edtext{\skt{विधिरेषो}}{\lemma{\skt{विधिर्}}\Afootnote{K H; \skt{विधि}\ A}}}{\lemma{\skt{एषो}}\Afootnote{K H; \skt{रेषो}\ A}}\skt{ }\edtext{\skt{यसादपि}}{\Afootnote{K H; \skt{ऽयसामपि ॥}\ A}}\skt{ ।}\pend

\pstart \ledsidenote{\passref{6.54.32ab}}\edtext{\skt{पूतीकस्य}}{\Afootnote{K H; \skt{पूतीकस्व}\ A}}\skt{ रसम्वापि }\edtext{\skt{सुरस्यां वा}}{\Afootnote{K; \emph{om.}\ A; \skt{सुरस्या…}\ H}}\skt{ }\edtext{\skt{पिबेत्नरः}}{\Afootnote{K; \skt{पिबेद्वा मधुना सह ।}\ A; \emph{om.}\ H}}\skt{ ।}\pend

\pstart\versepar \ledsidenote{\passref{6.54.33}}\skt{सप्तरात्रं पिबेद्घृष्टं }\edtext{\skt{त्रपु वा}}{\Afootnote{A K; \skt{त्रपु वो}\ H}}\skt{ दधिमस्तुना । }\pend
\pstart\versepar \skt{}\edtext{\skt{पुरीशजां}}{\Afootnote{K; \skt{पुरीषजा}\ A H}}\skt{ }\edtext{\skt{श्लेष्मजांश्च}}{\Afootnote{K; \skt{न्कफोत्थांश्च}\ A; \skt{न्श्शेष्मजा श्च}\ H}}\skt{ }\edtext{\skt{हन्यादेवं}}{\Afootnote{K H; \skt{हन्यादेवं}\ A}}\skt{ }\edtext{\skt{क्रिमीं}}{\Afootnote{K; \skt{कृमी}\ A H}}\skt{ }\edtext{\skt{भिषक्}}{\Afootnote{K; \skt{न्भिषक् ॥}\ A; \skt{न्भिषक ॥}\ H}}\skt{ ।}\pend

\pstart\versepar \ledsidenote{\passref{6.54.34}}\edtext{\skt{शिरो\-हृ\-त्क\-र्ण्ण\-ना\-सा\-क्षि\-सं\-श्रि\-तांश्च}}{\Afootnote{K; \skt{शिरो\-हृ\-द्घ्रा\-ण\-क\-र्णा\-क्षि\-सं\-श्रि\-तांश्च}\ A; \skt{शिरो\-हृ\-त्क\-र्ण\-ना\-सा\-ष्कि\-स\-म्श्रि\-ताश्च}\ H}}\skt{ पृथक्पृथक् । }\pend
\pstart\versepar \skt{विशे\-ष\-तो\-ञ्ज\-नै\-व\-र्ज\-मा\-नादि वास\-या\-ना\-ल\-डते क्रिमीर्नस्यैर व }\edtext{\skt{पीडैश्च}}{\Afootnote{A K; \skt{पीन्डैश्च}\ H}}\skt{ साधयेत् ।}\edtext{}{\lemma{\skt{पृथक्पृथक् । विशे\-ष\-तो\-ञ्ज\-नै\-व\-र्ज\-मा\-नादि वास\-या\-ना\-ल\-डते क्रिमीर्नस्यैर}}\Afootnote{K\textsuperscript{pc}; \skt{पृथग्विधान् । विशे\-षे\-णा\-ञ्ज\-नै\-र्न\-स्यैर}\ A; \skt{पृथक्पृथक् । विशे\-ष\-तो\-ञ्ज\-नै\-र्न\-स्यैर}\ K\textsuperscript{ac}; \skt{पृथक्पृथक् । विषे\-श\-तो\-ञ्ज\-नै\-र्न्न\-स्यैर}\ H}}\pend

\pstart\versepar \ledsidenote{\passref{6.54.35}}\edtext{\skt{सकृद्रसं}}{\Afootnote{K; \skt{शकृद्रसं}\ A H}}\skt{ }\edtext{\skt{शारङ्गस्य}}{\Afootnote{K; \skt{तुरङ्गस्य}\ A; \skt{सरङ्गस्य}\ H}}\skt{ }\edtext{\skt{शोष्यन्तुमौ}}{\Afootnote{K; \skt{सुशु}\ A; \skt{शोष्यतुमौ}\ H}}\skt{ विभावयेत् । }\pend
\pstart\versepar \skt{निःक्वाथेन विडंगानां चूर्ण्णं प्रधमनन्तु तत् ।}\edtext{}{\lemma{\skt{विभावयेत् । निःक्वाथेन}}\Afootnote{K H; \skt{ष्कं भावयेदति । निष्क्वाथेन}\ A}}\pend

\pstart\versepar \ledsidenote{\passref{6.54.36}}\edtext{\skt{अय\-श्चू\-र्ण्णा\-न्य\-ने\-रैव}}{\lemma{\skt{चूर्ण्णान्यनेरैव}}\Afootnote{K; \skt{चूर्णान्यनेनैव}\ A; \skt{चूर्णास्यत्रै लेवा}\ H}}\skt{ }\edtext{\skt{विधि\-ना\-यो\-ज\-यीत}}{\Afootnote{K; \skt{विधिना योजयेद्भि}\ A; \skt{वि धिनायोजयीत्च ।}\ H}}\skt{ }\edtext{\skt{च}}{\Afootnote{K; \skt{ष}\ A; \emph{om.}\ H}}\skt{ । }\pend
\pstart\versepar \skt{}\edtext{\skt{सकांसनीलं}}{\Afootnote{K H; \skt{क् । सकांस्यनीलं}\ A}}\skt{ तैलं च नस्यं }\edtext{\skt{स्यात्सुरसादिके}}{\Afootnote{A K; \skt{स्यात्सुरनादिके ॥}\ H}}\skt{ ।}\pend

\pstart\versepar \ledsidenote{\passref{6.54.37}}\edtext{\skt{इन्द्र\-लु\-प्ते\-वि\-धि\-श्चापि}}{\lemma{\skt{इन्द्र\-लु\-प्ते\-वि\-धिश्}}\Afootnote{K; \skt{इन्द्र\-लु\-प्त\-वि\-धिश्}\ A H}}\skt{ विधेयो }\edtext{\skt{रोमराजिषु}}{\Afootnote{K; \skt{रोमभोजिषु ।}\ A; \skt{रोमवाजिषु ।}\ H}}\skt{ । }\pend
\pstart\versepar \skt{}\edtext{\skt{दन्ताजानां}}{\Afootnote{K H; \skt{दन्तादानां}\ A}}\skt{ समुद्दिष्टं विधानं }\edtext{\skt{मुखरौगिके}}{\Afootnote{K; \skt{मुखरोगिकम् ॥}\ A; \skt{मुखरोगिके ॥}\ H}}\skt{ ।}\pend

\pstart\versepar \ledsidenote{\passref{6.54.38}}\skt{रक्तजानां प्रतीकारं कुर्या\-त्कु\-ष्ठ\-चि\-कि\-त्सि\-तात् । }\pend
\pstart\versepar \skt{सुरसादिं च }\edtext{\skt{सर्वेषु}}{\Afootnote{A K; \skt{सर्व्वषु}\ H}}\skt{ }\edtext{\skt{सर्व\-थै\-वो\-प\-यो\-ज\-येत्}}{\Afootnote{K H; \skt{सर्व\-थौ\-वो\-प\-यो\-ज\-येत् ॥}\ A}}\skt{ ।}\edtext{}{\lemma{\skt{कुर्या\-त्कु\-ष्ठ\-चि\-कि\-त्सि\-तात् । सुरसादिं च}}\Afootnote{K; \skt{कुर्या\-त्कु\-ष्ठ\-चि\-कि\-त्सिते । सुरसादिं तु}\ A; \skt{कुर्या\-त्कु\-ष्ठ\-चि\-कि\-त्सि\-तत् । सुरसादिश्च}\ H}}\pend

\pstart\versepar \ledsidenote{\passref{6.54.39}}\edtext{\skt{प्रव्य\-क्त\-ति\-क्त\-क\-टुकं}}{\Afootnote{K; \skt{प्रव्य\-क्त\-ति\-क\-त\-क\-टुकं}\ A H}}\skt{ भोजनञ्च हितं }\edtext{\skt{सदा}}{\Afootnote{K; \skt{भवे}\ A; \skt{सद ।}\ H}}\skt{ । }\pend
\pstart\versepar \skt{}\edtext{\skt{कुल\-थ\-र\-स\-सं\-युक्तं}}{\Afootnote{K; \skt{त् । कुल\-त्थ\-क्षा\-र\-सं\-सृष्टं}\ A; \skt{कु लथ\-र\-स\-सं\-युक्तं…}\ H}}\skt{ क्षारपानञ्च पूजितं ।}\pend

\pstart\versepar \ledsidenote{\passref{6.54.40}}\edtext{\skt{क्षीराणि}}{\Afootnote{A K; \skt{क्षीराषि}\ H}}\skt{ मांसानि घृतानि चैव दधीनि शाकानि च पर्णवन्ति । }\pend
\pstart\versepar \skt{}\edtext{\skt{माषान}}{\Afootnote{K H; \skt{समास}\ A}}\skt{ }\edtext{\skt{तोम्व्लां}}{\Afootnote{K; \skt{तोऽम्लान्}\ A; \skt{तौऽम्लान्}\ H}}\skt{ }\edtext{\skt{मधूरां}}{\Afootnote{K; \skt{मधुरा}\ A; \skt{मधूरा}\ H}}\skt{ }\edtext{\skt{रसांश्च}}{\Afootnote{K; \skt{न्हिमांश्च}\ A; \skt{एअसा च}\ H}}\skt{ }\edtext{\skt{क्र्मीं}}{\Afootnote{K; \skt{कृमी}\ A H}}\skt{ }\edtext{\skt{जिघांसुः}}{\Afootnote{K; \skt{न्जिघांसुः}\ A H}}\skt{ }\edtext{\skt{परिवर्जयीत}}{\Afootnote{K; \skt{परिवर्जयेत्तु ॥}\ A; \skt{परिवर्जयीत् ॥}\ H}}\skt{ ।}\pend

\pstart \ledsidenote{\passref{6.54.trailer}}\edtext{\skt{कायचिकित्सा}}{\Afootnote{\skt{इति सुश्रु\-त\-सं\-हि\-ता\-या\-मु\-त्त\-र\-त\-न्त्रा\-न्त\-र्गते काय\-चि\-कि\-त्सा\-तन्त्रे कृमिप्रतिषेधो नाम (षोडशोऽध्यायः, आदितः) चतुः\-प\-ञ्चा\-श\-त्त\-मोऽ\-ध्यायः ॥५४ ॥}\ A; \skt{कायचि कित्सा ॥ ॥}\ K; \skt{उत्तरतन्त्रे द्वपन्चशत्तमः कायचिकित्सायं कृमिचिकित्स नाम सप्त\-द\-ष\-मो\-ध्यायः ॥५४ ॥}\ H}}\skt{ ॥ ॥}\pend

\endnumbering
\beginnumbering
\eledsection*{Adhyāya 55}
\pstart \ledsidenote{\passref{6.55.1}}\skt{अथात उदा\-व\-र्त्त\-प्र\-ति\-षेधं }\edtext{\skt{व्याव्याख्यास्यामः}}{\lemma{\skt{व्या}}\Afootnote{K; \skt{अध्यायं व्याख्यास्यामः ॥}\ A; \skt{व्याख्यास्यामः ॥}\ H}}\skt{ ॥}\pend

\pstart\versepar \ledsidenote{\passref{6.55.3}}\skt{अधश्चोर्ध्वञ्च भावानां प्रवृत्तानां स्वभावतः । }\pend
\pstart\versepar \skt{न वेगा\-न्धा\-र\-ये\-त्प्राज्ञो वाता\-दी\-ना\-ञ्जि\-जी\-विषुः ।}\pend

\pstart\versepar \ledsidenote{\passref{6.55.4}}\edtext{\edtext{\skt{वात\-वि\-ड्मू\-त्र\-जृ\-म्भा\-णा\-म\-स्रुणः}}{\lemma{\skt{वात\-वि\-ड्मू\-त्र\-जृ\-म्भा\-णाम्}}\Afootnote{K\textsuperscript{pc}; \skt{वात\-वि\-ण्मू\-त्र\-जृ\-म्भाश्रु}\ A; \skt{वात\-आ\-वि\-ड्मू\-त्र\-जृ\-म्भा\-णाम्}\ K\textsuperscript{ac}; \skt{वात\-वि\-ण्मू\-त्र\-जृ\-म्भा\-णाम्}\ H}}}{\lemma{\skt{अस्रुणः}}\Afootnote{K H; \skt{क्षवो\-द्गा\-र\-व\-मी\-न्द्रियैः ।}\ A}}\skt{ }\edtext{\edtext{\skt{क्षवथोरपि}}{\lemma{\skt{क्षवथोर्}}\Afootnote{K H; \skt{व्याहन्यमानैर्}\ A}}}{\lemma{\skt{अपि}}\Afootnote{K H; \skt{उदितैर्}\ A}}\skt{ । }\pend
\pstart\versepar \skt{}\edtext{\skt{उद्गारस्य}}{\Afootnote{K H; \skt{उदावर्तो निरुच्य}\ A}}\skt{ }\edtext{\skt{तथा}}{\Afootnote{K H; \skt{ते ॥}\ A}}\skt{ }\edtext{\skt{र्च्छद्या}}{\Afootnote{K; \emph{om.}\ A; \skt{च्छर्द्द्या}\ H}}\skt{ }\edtext{\edtext{\skt{रेतसश्चापि}}{\lemma{\skt{रेतसश्}}\Afootnote{K H; \emph{om.}\ A}}}{\lemma{\skt{चापि}}\Afootnote{K H; \emph{om.}\ A}}\skt{ }\edtext{\skt{धारणात्}}{\Afootnote{K H; \emph{om.}\ A}}\skt{ ॥}\pend

\pstart\versepar \ledsidenote{\passref{6.55.5}}\edtext{\edtext{\skt{क्षुत्तृ\-ष्णो\-च्छ्वा\-स\-नि\-द्रा\-णा\-मु\-दा\-वर्त्तो}}{\lemma{\skt{क्षुत्तृ\-ष्णो\-च्छ्वा\-स\-नि\-द्रा\-णाम्}}\Afootnote{K H; \skt{क्षुत्तृ\-ष्णा\-श्वा\-स\-नि\-द्रा\-णाम्}\ A}}}{\lemma{\skt{उदावर्त्तो}}\Afootnote{K H; \emph{om.}\ A}}\skt{ }\edtext{\skt{भवेदिह}}{\Afootnote{K H; \emph{om.}\ A}}\skt{ । }\pend
\pstart\versepar \skt{}\edtext{\skt{व्याह\-न्य\-मा\-न\-वे\-गार्त्त}}{\Afootnote{K H; \emph{om.}\ A}}\skt{ उदावर्त्तो }\edtext{\skt{निरुच्यते}}{\Afootnote{K H; \skt{विधारणा}\ A}}\skt{ । }\pend
\pstart\versepar \skt{}\edtext{\skt{तस्याभिधास्ये}}{\Afootnote{K H; \skt{त् । तस्याभिधास्ये}\ A}}\skt{ व्यासेन लक्षणं }\edtext{\skt{सचिकित्सितं}}{\Afootnote{K H; \skt{च चिकित्सितम् ॥}\ A}}\skt{ ।}\pend

\pstart \ledsidenote{\passref{6.55.6}}\edtext{\skt{त्रयो\-द\-श\-वि\-ध\-स्यास्य}}{\Afootnote{K H; \skt{त्रयो\-द\-श\-वि\-ध\-श्चासौ}\ A}}\skt{ }\edtext{\skt{भिन्नस्येतैस्तु}}{\Afootnote{K; \skt{भिन्न एतैस्तु}\ A; \skt{भितस्यैतैस्तु}\ H}}\skt{ }\edtext{\skt{कारणैः}}{\Afootnote{K H; \skt{कारणैः । अप\-थ्य\-भो\-ज\-ना\-च्चापि वक्ष्यते च तथापरः ॥}\ A}}\skt{ ।}\pend

\pstart\versepar \ledsidenote{\passref{6.55.7}}\edtext{\skt{आध्मानशूलो}}{\Afootnote{K H; \skt{आध्मानशूलौ}\ A}}\skt{ }\edtext{\skt{हृदयो}}{\Afootnote{K; \skt{हृदयो}\ A H}}\skt{ परोधं शिरोरुजं }\edtext{\edtext{\skt{श्वासमतीव}}{\lemma{\skt{श्वासम्}}\Afootnote{K; \skt{श्वास}\ A; \skt{वास}\ H}}}{\lemma{\skt{अतीव}}\Afootnote{K; \skt{मतीव}\ A H}}\skt{ हिक्कां । }\pend
\pstart\versepar \skt{}\edtext{\skt{कास\-प्र\-ति\-श्या\-य\-ग\-ल\-ग्र\-ह\-त्वा\-द्व\-ला\-स\-पि\-त्त\-प्र\-स\-रञ्च}}{\Afootnote{K; \skt{कास\-प्र\-ति\-श्या\-य\-ग\-ल\-ग्र\-हांश्च बला\-स\-पि\-त्त\-प्र\-सरं च}\ A; \skt{कास\-प्र\-ति\-श्या\-य\-ग\-ल\-ग्र\-ह\-त्वा\-व\-ला\-स\-पि\-त्त\-प्र\-स\-रञ्च}\ H}}\skt{ घोरं ।}\pend

\pstart\versepar \ledsidenote{\passref{6.55.8}}\skt{कुर्यादपानो }\edtext{\skt{भिहतः}}{\Afootnote{K; \skt{ऽभिहतः}\ A H}}\skt{ स्वमार्गे }\edtext{\skt{हन्यात्पुरीशं}}{\Afootnote{K; \skt{हन्यात्पुरीषं}\ A H}}\skt{ }\edtext{\skt{मुखतः}}{\Afootnote{A K H\textsuperscript{pc}; \skt{मुखत}\ H\textsuperscript{ac}}}\skt{ क्षिपेद्वा । }\pend
\pstart\versepar \skt{आटोपशूलौ }\edtext{\skt{परिवर्त्तनञ्च}}{\Afootnote{K H; \skt{परिकर्तनं च}\ A}}\skt{ संगः }\edtext{\skt{पुरीशस्य}}{\Afootnote{K; \skt{पुरीषस्य}\ A H}}\skt{ तथोर्ध्ववातः ।}\pend

\pstart\versepar \ledsidenote{\passref{6.55.9}}\edtext{\skt{पुरी\-श\-मा\-स्या\-दपि वा}}{\Afootnote{K; \skt{पुरी\-ष\-मा\-स्या\-दपि वा}\ A H}}\skt{ निरेति }\edtext{\skt{पुरीशवेगे}}{\Afootnote{K; \skt{पुरीषवेगे}\ A H}}\skt{ }\edtext{\skt{भिहिते}}{\Afootnote{K; \skt{ऽभिहते}\ A; \skt{भिहते}\ H}}\skt{ नरस्य । }\pend
\pstart\versepar \skt{मूत्रस्य वेगे }\edtext{\skt{भिहते}}{\Afootnote{K\textsuperscript{pc} H; \skt{ऽभिहते}\ A; \skt{भिहि ते}\ K\textsuperscript{ac}}}\skt{ नरस्तु कृच्छ्रेण मूत्रं }\edtext{\skt{स्रवते}}{\Afootnote{K H; \skt{कुरुते}\ A}}\skt{ }\edtext{\skt{ल्पमल्पं}}{\lemma{\skt{ल्पम्}}\Afootnote{K H; \skt{ऽल्पम्}\ A}}\skt{ ।}\pend

\pstart\versepar \ledsidenote{\passref{6.55.10}}\skt{मेढ्रे गुदे }\edtext{\skt{वंक्ष\-ण\-ब\-स्ति\-मु\-ष्क\-ना\-भि\-प्र\-दे\-शे\-ष्वथ}}{\lemma{\skt{वंक्ष\-ण\-ब\-स्ति\-मु\-ष्क\-ना\-भि\-प्र\-दे\-शेष्व्}}\Afootnote{A K; \skt{वंक्ष\-ण\-मु\-ष्क\-ना\-भि\-प्र\-दे\-शेष्व्}\ H}}\skt{ }\edtext{\skt{मूर्ध्नि}}{\Afootnote{K H; \skt{वा+अपि मूर्ध्नि ।}\ A}}\skt{ }\edtext{\skt{चापि}}{\Afootnote{K H; \emph{om.}\ A}}\skt{ । }\pend
\pstart\versepar \skt{}\edtext{\skt{आनद्धबस्तेश्च}}{\Afootnote{K; \skt{आनद्धबस्तिश्च}\ A H}}\skt{ भवन्ति }\edtext{\skt{तीव्रा}}{\Afootnote{K H; \skt{तीव्राः}\ A}}\skt{ }\edtext{\skt{रुजश्च}}{\Afootnote{K H; \skt{शूलाश्च}\ A}}\skt{ }\edtext{\skt{शूलैरिव}}{\lemma{\skt{शूलैर्}}\Afootnote{A K\textsuperscript{pc} H; \skt{शूलौर्}\ K\textsuperscript{ac}}}\skt{ }\edtext{\skt{तुद्यते च}}{\Afootnote{K H; \skt{भिन्नमूर्तेः ॥}\ A}}\skt{ ।}\pend

\pstart\versepar \ledsidenote{\passref{6.55.11}}\skt{मन्या\-ग\-ल\-स्त\-म्भ\-शि\-रो\-वि\-काराः । }\edtext{\skt{जृम्भो\-प\-घा\-ता\-त्प\-व\-ना\-त्मिका स्युः}}{\Afootnote{K; \skt{जृम्भो\-प\-घा\-ता\-त्प\-व\-ना\-त्मकाः स्युः ।}\ A; \skt{जृम्भो\-प\-घा\-ता\-त्प\-व\-ना\-त्मिकाः स्युः ।}\ H}}\skt{ । }\pend
\pstart\versepar \skt{}\edtext{\skt{तथा\-क्षि\-ना\-सा\-व\-द\-ना\-म\-याश्च}}{\Afootnote{K H; \skt{श्रोत्रा\-न\-न\-घ्रा\-ण\-वि\-लो\-च\-नोत्था}\ A}}\skt{ भवन्ति तीव्राः }\edtext{\skt{सह}}{\Afootnote{K H; \skt{च तथा}\ A}}\skt{ }\edtext{\skt{कर्ण्णरोगैः}}{\Afootnote{K H; \skt{विकाराः ॥}\ A}}\skt{ ।}\pend

\pstart\versepar \ledsidenote{\passref{6.55.12}}\edtext{\edtext{\skt{आन\-न्द\-ज\-म्वा\-प्यथ}}{\lemma{\skt{वाप्य्}}\Afootnote{K H; \emph{om.}\ A}}}{\lemma{\skt{अथ}}\Afootnote{K H; \emph{om.}\ A}}\skt{ }\edtext{\skt{शोकजम्वा}}{\Afootnote{K H; \skt{शोकसमुद्भवं वा}\ A}}\skt{ नेत्रोदकं प्राप्तममुञ्चतो हि । }\pend
\pstart\versepar \skt{शिरोगुरुत्वं नयनामयाश्च भवन्ति तीव्राः सह }\edtext{\skt{पीनसेन}}{\Afootnote{A K\textsuperscript{pc} H; \skt{पीनसेनः}\ K\textsuperscript{ac}}}\pend

\pstart\versepar \ledsidenote{\passref{6.55.13}}\skt{भवन्ति गाढं क्षव\-थो\-र्वि\-घा\-ता\-च्छि\-रो\-क्षि\-ना\-सा\-श्र\-व\-णेषु रोगाः । }\pend
\pstart\versepar \skt{}\edtext{\skt{कण्ठस्य}}{\Afootnote{K H; \skt{कण्ठास्य}\ A}}\skt{ पूर्ण्णत्वमतीव तोदः }\edtext{\skt{कुक्षौ तु}}{\Afootnote{K H; \skt{कूजश्च}\ A}}\skt{ }\edtext{\skt{वायोरुभयोः}}{\lemma{\skt{उभयोः}}\Afootnote{K H; \skt{उत वा}\ A}}\skt{ प्रवृत्तिः ।}\pend

\pstart\versepar \ledsidenote{\passref{6.55.14}}\skt{उद्गारवेगे }\edtext{\skt{भिहते}}{\Afootnote{K H; \skt{ऽभिहते}\ A}}\skt{ भवन्ति }\edtext{\skt{जन्तोर्विकाराः}}{\lemma{\skt{जन्तोर्}}\Afootnote{K H; \skt{घोरा}\ A}}\skt{ पवनप्रसूताः । }\pend
\pstart\versepar \skt{}\edtext{\skt{छर्द्याभिघातेन}}{\Afootnote{K H; \skt{छर्देर्विघातेन}\ A}}\skt{ भवेच्च कुष्ठं येनैव दोषेण विदग्धमन्नं ।}\pend

\pstart\versepar \ledsidenote{\passref{6.55.15}}\edtext{\skt{मूत्रायणे}}{\Afootnote{K; \skt{मूत्राशये}\ A; \skt{मूत्रायने}\ H}}\skt{ पायुनि मुष्कयोश्च }\edtext{\skt{शोथो}}{\Afootnote{K H; \skt{शोफो}\ A}}\skt{ }\edtext{\skt{रुजा}}{\Afootnote{K H; \skt{रुजो}\ A}}\skt{ मूत्र\-वि\-नि\-ग्र\-हश्च । }\pend
\pstart\versepar \skt{शुक्राश्मरी तत्स्रवणं भवेद्वा ते ते }\edtext{\skt{विकाराभिहते तु}}{\Afootnote{K H; \skt{विकारा विहते तु}\ A}}\skt{ शुक्रे ।}\pend

\pstart\versepar \ledsidenote{\passref{6.55.16}}\edtext{\skt{तन्द्रा\-ङ्ग\-म\-र्दा\-व\-रुचि}}{\Afootnote{K; \skt{तन्द्रा\-ङ्ग\-म\-र्दा\-रु\-चिवि}\ A; \skt{तन्द्रा\-ङ्ग\-म\-र्द्दा\-व\-रुचिः}\ H}}\skt{ }\edtext{\skt{श्रमश्च}}{\Afootnote{K H; \skt{भ्रमाः स्युः}\ A}}\skt{ }\edtext{\edtext{\skt{क्षुधोर्विकारा}}{\lemma{\skt{क्षुधोर्}}\Afootnote{K; \skt{क्षुधो}\ A H}}}{\lemma{\skt{विकारा}}\Afootnote{K; \skt{ऽभिघाता}\ A; \skt{विकाराः}\ H}}\skt{ }\edtext{\skt{कृशता च}}{\Afootnote{K H; \skt{त्कृशता च}\ A}}\skt{ }\edtext{\skt{दृष्टा}}{\Afootnote{K H; \skt{दृष्टेः ।}\ A}}\skt{ । }\pend
\pstart\versepar \skt{कण्ठास्यशोषः }\edtext{\edtext{\skt{श्रव\-णा\-व\-रो\-ध\-स्तृ\-ष्णा\-भि\-घा\-ता\-द्धृ\-द\-य\-व्यथा}}{\lemma{\skt{श्रव\-णा\-व\-रो\-धस्}}\Afootnote{K; \skt{श्रवणावरोध}\ A; \skt{स्रव\-णा\-व\-रो\-धस्}\ H}}}{\lemma{\skt{तृष्णा\-भि\-घा\-ता\-द्धृ\-द\-य\-व्यथा}}\Afootnote{K H; \skt{स्तृष्णा\-भि\-घा\-ता\-द्धृ\-दये व्यथा}\ A}}\skt{ च ।}\pend

\pstart\versepar \ledsidenote{\passref{6.55.17}}\skt{श्रान्तस्य }\edtext{\skt{निश्वा\-स\-वि\-नि\-ग्र\-हेण}}{\Afootnote{K H; \skt{निःश्वा\-स\-वि\-नि\-ग्र\-हेण}\ A}}\skt{ हृद्रो\-ग\-मो\-हा\-वथ वापि गुल्मः । }\pend
\pstart\versepar \skt{}\edtext{\edtext{\skt{जृम्भोऽङ्गमर्दो}}{\lemma{\skt{जृम्भो}}\Afootnote{K H; \skt{जृम्भाङ्}\ A}}}{\lemma{\skt{ऽङ्गमर्दो}}\Afootnote{K H; \skt{गमर्दो}\ A}}\skt{ }\edtext{\skt{ङ्गशिरोथ}}{\Afootnote{K H; \skt{ऽङ्गशिरोक्षि}\ A}}\skt{ जाड्यं निद्रा\-भि\-घा\-ता\-द\-थ\-वापि }\edtext{\skt{तन्द्री}}{\Afootnote{K H; \skt{तन्द्रा ॥}\ A}}\skt{ ॥}\pend

\pstart\versepar \ledsidenote{\passref{6.55.18}}\skt{तृष्णार्द्दितं परिक्लिष्टं क्षीणं शूलैरभिद्रुतं । }\pend
\pstart\versepar \skt{शकृद्वमन्तं मति\-मा\-नु\-दा\-व\-र्त्ति\-न\-मु\-त्सृ\-जेत् ।}\pend

\pstart\versepar \ledsidenote{\passref{6.55.19}}\edtext{\skt{सर्स्वेष्वेतेषु}}{\lemma{\skt{सर्स्वेष्व्}}\Afootnote{K; \skt{सर्वेष्व्}\ A H}}\skt{ }\edtext{\skt{विधि\-व\-दु\-दा\-व\-र्त्तेषु}}{\Afootnote{A K H\textsuperscript{pc}; \skt{विधि\-व\-दु\-दा\-वर्त्ते}\ H\textsuperscript{ac}}}\skt{ कृत्स्नशः । }\pend
\pstart\versepar \skt{वायोः क्रिया }\edtext{\skt{विधातव्या}}{\Afootnote{K H; \skt{विधातव्याः}\ A}}\skt{ स्वमा\-र्ग\-प्र\-ति\-प\-त्तये ।}\pend

\pstart\versepar \ledsidenote{\passref{6.55.20}}\skt{सामान्यतः पृथक्त्वेन क्रियां भूयो निबोध मे । }\pend
\pstart\versepar \skt{आस्थापनं मारुतजे }\edtext{\edtext{\skt{स्निग्ध\-स्वि\-न्न\-म्वि\-शे\-षतः}}{\lemma{\skt{स्निग्धस्विन्नम्}}\Afootnote{K H; \skt{स्निग्धस्विन्ने}\ A}}}{\lemma{\skt{विशेषतः}}\Afootnote{K H; \skt{विशिष्यते ॥}\ A}}\pend

\pstart\versepar \ledsidenote{\passref{6.55.21}}\edtext{\skt{पुरीशजे तु}}{\Afootnote{K H; \skt{पुरीषजे तु}\ A}}\skt{ कर्त्तव्यो }\edtext{\skt{विधि\-रा\-ना\-हि\-कस्तु}}{\lemma{\skt{तु}}\Afootnote{K H; \skt{भ}\ A}}\skt{ }\edtext{\skt{यः}}{\Afootnote{K H; \skt{वे}\ A}}\skt{ । }\pend
\pstart\versepar \skt{}\edtext{\skt{सौवर्चलाढ्यां}}{\Afootnote{K H; \skt{त् । सौवर्चलाढ्यां}\ A}}\skt{ मदिरां मूत्रे त्वभिहते पिबेत् ।}\pend

\pstart\versepar \ledsidenote{\passref{6.55.22}}\skt{एलाम्वाप्यथ मद्येन }\edtext{\skt{क्षीरवारिं}}{\Afootnote{K; \skt{क्षीरं वापि}\ A; \skt{क्षीरवारि}\ H}}\skt{ }\edtext{\skt{पिबेत वा}}{\Afootnote{K H; \skt{पिबेन्नरः ।}\ A}}\skt{ । }\pend
\pstart\versepar \skt{}\edtext{\skt{धात्रीफलानाञ्च}}{\lemma{\skt{च}}\Afootnote{K H; \skt{स्व}\ A}}\skt{ रसं }\edtext{\skt{सतैलम्वा}}{\Afootnote{K H; \skt{सजलं वा}\ A}}\skt{ पिबेत्त्र्यहं ।}\pend

\pstart\versepar \ledsidenote{\passref{6.55.23}}\edtext{\edtext{\skt{रस\-म\-श्व\-पु\-री\-शस्य}}{\lemma{\skt{रसम्}}\Afootnote{K H; \skt{रस}\ A}}}{\lemma{\skt{अश्वपुरीशस्य}}\Afootnote{K; \skt{मश्वपुरीषस्य}\ A; \skt{अश्वपुरीषस्य}\ H}}\skt{ }\edtext{\skt{गर्दभस्य}}{\Afootnote{K H; \skt{गर्दभस्या}\ A}}\skt{ }\edtext{\skt{पिबेत वा}}{\Afootnote{K; \skt{थवा पिबे}\ A; \skt{पिबेत्तु वा ।}\ H}}\skt{ । }\pend
\pstart\versepar \skt{}\edtext{\skt{मांसोपदन्शं}}{\Afootnote{K; \skt{त् । मांसोपदंशं}\ A; \skt{मासोपदंशम्}\ H}}\skt{ }\edtext{\skt{मद्यन्तु}}{\Afootnote{K H; \skt{मधु वा पिबेद्वा सीधु गौडिकम् ॥}\ A}}\skt{ }\edtext{\skt{मुख्यं}}{\Afootnote{K H; \emph{om.}\ A}}\skt{ }\edtext{\skt{सेवेत}}{\Afootnote{K H; \emph{om.}\ A}}\skt{ }\edtext{\skt{योगवित्}}{\Afootnote{K H; \emph{om.}\ A}}\skt{ ।}\pend

\pstart\versepar \ledsidenote{\passref{6.55.24}}\skt{भद्रदारु }\edtext{\skt{घणं}}{\Afootnote{K; \skt{घनं}\ A H}}\skt{ }\edtext{\skt{मूर्वां}}{\Afootnote{K; \skt{मूर्वा}\ A H}}\skt{ }\edtext{\skt{हरिद्रां}}{\Afootnote{K H; \skt{हरिद्रा}\ A}}\skt{ मधुकन्तथा । }\pend
\pstart\versepar \skt{कोलप्रमाणानि }\edtext{\skt{पिबे\-न्मा\-ष\-क्वा\-थेन}}{\Afootnote{K; \skt{पिबे}\ A; \skt{पिबे\-त्मा\-ष\-क्वा\-थेन}\ H}}\skt{ }\edtext{\skt{योगवित्}}{\Afootnote{K H; \skt{दान्तरिक्षेण वारिणा ॥}\ A}}\skt{ ।}\pend

\pstart\versepar \ledsidenote{\passref{6.55.25}}\edtext{\skt{दुस्पशायाः}}{\Afootnote{K; \skt{दुःस्पर्शा}\ A; \skt{दुर्स्पर्शायाः}\ H\textsuperscript{ac}; \skt{दुस्पर्शायाः}\ H\textsuperscript{pc}}}\skt{ }\edtext{\skt{स्वरसम्वा}}{\lemma{\skt{स्वरसम्}}\Afootnote{K H; \skt{स्वरसं वाप्}\ A}}\skt{ कषायं }\edtext{\skt{कुंकुमस्य वा}}{\Afootnote{K; \skt{कुङ्कुमस्य च ।}\ A; \skt{कुंक्षमस्य वा ।}\ H}}\skt{ ॥ }\pend
\pstart\versepar \skt{एर्वारुबीजं तोयेन पिबेद्वा लवणीकृतं ।}\pend

\pstart\versepar \ledsidenote{\passref{6.55.26}}\edtext{\skt{पञ्चमूलश्रितं}}{\Afootnote{K; \skt{पञ्चमूलीशृतं}\ A; \skt{कणामूलश्रितं}\ H}}\skt{ क्षीरं द्राक्षा\-र\-स\-म\-थापि वा । }\pend
\pstart\versepar \skt{योगांश्च वितरेदत्र }\edtext{\skt{पूर्वो\-क्ता\-न\-श्म\-री\-भिदः}}{\Afootnote{A K; \skt{पूर्वो\-क्ता\-न\-श्म\-री\-र्भिदः ॥}\ H}}\skt{ ।}\pend

\pstart\versepar \ledsidenote{\passref{6.55.27}}\skt{मूत्र\-कृ\-च्छ्र\-क्र\-म\-ञ्चापि कुर्या\-न्नि\-र\-व\-शे\-षतः ॥ }\pend
\pstart\versepar \skt{भूयो वक्ष्यामि }\edtext{\edtext{\skt{योगांश्च}}{\lemma{\skt{योगांश्}}\Afootnote{K H; \skt{योगा}\ A}}}{\lemma{\skt{च}}\Afootnote{K H; \skt{न्या}\ A}}\skt{ }\edtext{\skt{मूत्रा\-घा\-तो\-प\-शा\-न्तये}}{\Afootnote{K H; \skt{न्मूत्रा\-घा\-तो\-प\-शा\-न्तये ॥}\ A}}\skt{ ।}\pend

\pstart\versepar \ledsidenote{\passref{6.55.28}}\edtext{\edtext{\skt{स्नेह\-स्वे\-दै\-रु\-दा\-वर्त्तं}}{\lemma{\skt{स्नेहस्वेदैर्}}\Afootnote{K H; \skt{स्नेहैः स्वेदैर्}\ A}}}{\lemma{\skt{उदावर्त्तं}}\Afootnote{A K; \skt{उदावर्त्त}\ H}}\skt{ }\edtext{\skt{जृम्भजं}}{\Afootnote{K H; \skt{जृम्भाजं}\ A}}\skt{ समुपाचरेत् । }\pend
\pstart\versepar \skt{अश्रुमोक्षास्रुजे कार्यः स्निग्धस्विन्नस्य देहिनः ।}\edtext{}{\lemma{\skt{समुपाचरेत् । अश्रुमोक्षास्रुजे}}\Afootnote{K H; \skt{समुपाचरेत् । अश्रुमोक्षोऽश्रुजे}\ A}}\pend

\pstart\versepar \ledsidenote{\passref{6.55.29}}\edtext{\skt{तीक्ष्णां\-ज\-ना\-व\-पी\-डाभ्यां}}{\Afootnote{A K; \skt{तीक्ष्णं\-ज\-ना\-व\-पी\-डाभ्यां}\ H}}\skt{ }\edtext{\skt{तीक्ष्ण\-ग\-न्धो\-प\-जि\-घ्रणैः}}{\Afootnote{K H; \skt{तीक्ष्ण\-ग\-न्धो\-प\-शि\-ङ्घनैः ।}\ A}}\skt{ । }\pend
\pstart\versepar \skt{वर्त्ति\-प्र\-यो\-गै\-रथ वा }\edtext{\skt{क्षवसक्तिं}}{\Afootnote{A K; \skt{क्षवशक्तिम्}\ H}}\skt{ प्रवर्त्तयेत् ।}\pend

\pstart \ledsidenote{\passref{6.55.30}}\edtext{\skt{उद्गारजे}}{\Afootnote{K H; \skt{(ती\-क्ष्णै\-ष\-ध\-प्र\-ध\-म\-नै\-र\-थ\-वा\-दि\-त्य\-र\-श्मिभिः । ) उद्गारजे}\ A}}\skt{ }\edtext{\skt{क्रमोपेतं}}{\Afootnote{A K; \skt{क्रमोपेत}\ H}}\skt{ स्नैहिकं धूममाचरेत् ।}\pend

\pstart \ledsidenote{\passref{6.55.31}}\edtext{\skt{च्छर्द्याघातं}}{\Afootnote{K H; \skt{सुरां सौवर्चलवतीं बीज\-पू\-र\-र\-सा\-न्वि\-ताम् । छर्द्याघातं}\ A}}\skt{ यथादोषं सम्य\-क्स्ने\-हा\-दि\-भि\-र्ज\-येत् ।}\pend

\pstart \ledsidenote{\passref{6.55.32}}\edtext{\skt{बस्ति\-शु\-द्धि\-क\-रा\-वापं}}{\Afootnote{K H; \skt{सक्षा\-र\-ल\-व\-णो\-पे\-त\-म\-भ्यङ्गं चात्र दापयेत् । बस्ति\-शु\-द्धि\-क\-रा\-वापं}\ A}}\skt{ चतुर्गुणजलं पयः ।}\pend

\pstart\versepar \ledsidenote{\passref{6.55.33}}\edtext{\skt{आवा\-रि\-ना\-शा\-क्व\-थितं}}{\Afootnote{K; \skt{आवा\-रि\-ना\-शा\-त्क्व\-थितं}\ A; \skt{आवा\-रि\-ना\-स\-क्व\-थितं}\ H}}\skt{ पीतवन्तं प्रकामतः । }\pend
\pstart\versepar \skt{}\edtext{\skt{कामयेरन्प्रियां}}{\Afootnote{K H; \skt{रमयेयुः प्रिया}\ A}}\skt{ }\edtext{\skt{नार्यां}}{\Afootnote{K H; \skt{नार्यः}\ A}}\skt{ शुक्रो\-दा\-व\-र्त्ति\-न\-न्नरं ।}\pend

\pstart\versepar \ledsidenote{\passref{6.55.34}}\skt{क्षुद्विघाते हितं }\edtext{\skt{स्निग्ध\-मु\-ष्ण\-म\-ल्पञ्च}}{\lemma{\skt{अल्पञ्च}}\Afootnote{A K; \skt{अल्पन च}\ H}}\skt{ भोजनं । }\pend
\pstart\versepar \skt{तृष्णाघाते }\edtext{\skt{पिबेन्मन्थं}}{\Afootnote{A K; \skt{पिबेत्मन्थं}\ H}}\skt{ यवागूम्वा }\edtext{\skt{सुशीतलां}}{\Afootnote{K H; \skt{पि शीतलाम् ॥}\ A}}\skt{ ।}\pend

\pstart\versepar \ledsidenote{\passref{6.55.35}}\edtext{\skt{भोज्येद्रसेन}}{\Afootnote{K H; \skt{भोज्यो रसेन}\ A}}\skt{ विश्रान्तः श्रमश्वासातुरो नरः । }\pend
\pstart\versepar \skt{निद्राघाते पिबेत्क्षीरं }\edtext{\skt{स्वपे\-च्चे\-ष्ट\-क\-था\-रतः}}{\Afootnote{K; \skt{स्वप्याच्चेष्टकथा नरः ॥}\ A; \skt{स्वपे चेष्टकथारतः ।}\ H}}\skt{ ।}\pend

\pstart\versepar \ledsidenote{\passref{6.55.36}}\skt{आध्मानाद्येषु रोगेषु यथास्वं प्रयतेत हि । }\pend
\pstart\versepar \skt{यच्च यत्र भवेत्प्राप्तं तच्च }\edtext{\skt{तस्मिं}}{\Afootnote{K; \skt{तस्मि}\ A H}}\skt{ }\edtext{\skt{प्रयोजयेत्}}{\Afootnote{K; \skt{न्प्रयोजयेत्}\ A H}}\skt{ ।}\pend

\pstart\versepar \ledsidenote{\passref{6.55.37}}\skt{वायुः कोष्ठानुगो रूक्षैः कषा\-य\-क\-टु\-ति\-क्तकैः । }\pend
\pstart\versepar \skt{भोजनैः कुपितः सद्य उदावर्त्तं करोति }\edtext{\skt{ह}}{\Afootnote{K H; \skt{हि ॥}\ A}}\skt{ ।}\pend

\pstart\versepar \ledsidenote{\passref{6.55.38}}\edtext{\skt{वात\-मू\-त्र\-पु\-री\-शा\-सृ\-क्क\-फ\-मे\-दो\-व\-हानि वै}}{\Afootnote{K; \skt{वात\-मू\-त्र\-पु\-री\-षा\-सृ\-क्क\-फ\-मे\-दो\-व\-हानि वै ।}\ A; \skt{वात\-मू\-त्र\-पु\-री\-षा\-सृ\-क्क\-फ\-मो\-दो\-व\-हानि वै ॥}\ H}}\skt{ । }\pend
\pstart\versepar \skt{स्रोतां\-स्यु\-दा\-व\-र्त्त\-यति }\edtext{\edtext{\skt{पुरीशम्वापि}}{\lemma{\skt{पुरीशम्}}\Afootnote{K; \skt{पुरीषं}\ A; \skt{पुरीशश्}\ H}}}{\lemma{\skt{वापि}}\Afootnote{K; \skt{चाति}\ A; \skt{चापि}\ H}}\skt{ वर्त्तयेत् ।}\pend

\pstart\versepar \ledsidenote{\passref{6.55.39}}\skt{ततो हृद्बस्तिशूलार्त्तो }\edtext{\skt{हृल्ला\-सा\-रु\-चि\-पी\-डितः}}{\Afootnote{K H; \skt{गौर\-वा\-रु\-चि\-पी\-डितः ।}\ A}}\skt{ । }\pend
\pstart\versepar \skt{}\edtext{\skt{वात\-मू\-त्र\-पु\-री\-शाणि}}{\Afootnote{K; \skt{वात\-मू\-त्र\-पु\-री\-षाणि}\ A H}}\skt{ कृच्छ्रेण }\edtext{\skt{लभते न्नरः}}{\Afootnote{K; \skt{कुरुते नरः ॥}\ A; \skt{लभते नरः ।}\ H}}\skt{ ।}\pend

\pstart\versepar \ledsidenote{\passref{6.55.40}}\edtext{\skt{श्वास\-का\-स\-प्र\-ति\-श्या\-य\-दा\-ह\-मो\-ह\-तृ\-षा\-ज्वरं}}{\Afootnote{K; \skt{श्वास\-का\-स\-प्र\-ति\-श्या\-य\-दा\-ह\-मो\-ह\-व\-मि\-ज्व\-रान् ।}\ A; \skt{श्वास\-का\-स\-प्र\-ति\-श्या\-य\-दा\-ह\-सा\-द\-तृ\-षा\-ज्वरं ॥}\ H}}\skt{ । }\pend
\pstart\versepar \skt{}\edtext{\skt{वमीं}}{\Afootnote{K H; \skt{तृष्णा}\ A}}\skt{ }\edtext{\skt{हिक्कां}}{\Afootnote{K H; \skt{हिक्का}\ A}}\skt{ }\edtext{\skt{शिरोरोगं}}{\Afootnote{K H; \skt{शिरोरोग}\ A}}\skt{ }\edtext{\skt{मनः\-श्र\-व\-ण\-वि\-भ्रमां}}{\Afootnote{K; \skt{मनः\-श्र\-व\-ण\-वि\-भ्र\-मान् ॥}\ A; \skt{मन\-स्स्र\-व\-ण\-वि\-भ्रमां ।}\ H}}\skt{ ।}\pend

\pstart\versepar \ledsidenote{\passref{6.55.41}}\edtext{\skt{बहूनन्यांश्च}}{\Afootnote{K H; \skt{लभते च बहूनन्या}\ A}}\skt{ }\edtext{\skt{लभते}}{\Afootnote{K H; \emph{om.}\ A}}\skt{ }\edtext{\edtext{\skt{विका\-रा\-म्वा\-त\-को\-पजां}}{\lemma{\skt{विकाराम्}}\Afootnote{K; \skt{न्विकारा}\ A; \skt{विकारा}\ H}}}{\lemma{\skt{वातकोपजां}}\Afootnote{K; \skt{न्वातकोपजान्}\ A H}}\skt{ । }\pend
\pstart\versepar \skt{}\edtext{\skt{तत्तै\-ल\-ल\-व\-णा\-भ्यक्तं}}{\Afootnote{K H; \skt{तं तैल\-ल\-व\-णा\-भ्यक्तं}\ A}}\skt{ स्निग्धं स्विन्नं निरूहयेत् ।}\pend

\pstart\versepar \ledsidenote{\passref{6.55.42}}\skt{दोषतो }\edtext{\skt{भिन्नवर्चांसि}}{\Afootnote{K H; \skt{भिन्नवर्चस्कं}\ A}}\skt{ भुक्तं चाप्य\-नु\-वा\-स\-येत् । }\pend
\pstart\versepar \skt{}\edtext{\skt{भवे\-च्छा\-न्ति\-म्व्र\-ज\-त्ये\-व\-मु\-दा\-व\-र्त्तस्तु}}{\lemma{\skt{तु}}\Afootnote{K H; \skt{सु}\ A}}\skt{ दारुणः ।}\edtext{}{\lemma{\skt{अनुवासयेत् । भवेच्छान्तिम्}}\Afootnote{K H; \skt{अनुवासयेत् । न चेच्छान्तिं}\ A}}\pend

\pstart\versepar \ledsidenote{\passref{6.55.43}}\skt{अथैनं बहुशः स्विन्नं युञ्ज्या\-त्स्ने\-ह\-वि\-रे\-चनैः । }\pend
\pstart\versepar \skt{}\edtext{\skt{पाययेद्वा}}{\Afootnote{K H; \skt{पाययेत}\ A}}\skt{ }\edtext{\edtext{\skt{तृवृ\-त्पी\-लु\-य\-वा\-ग्वा\-म\-म्ब्ल\-पा\-नकैः}}{\lemma{\skt{तृवृ\-त्पी\-लु\-य\-वा\-ग्वाम्}}\Afootnote{K H; \skt{त्रिवृ\-त्पी\-लु\-य\-वानी}\ A}}}{\lemma{\skt{अम्ब्लपानकैः}}\Afootnote{K; \skt{रम्लपाचनैः ॥}\ A; \skt{अम्लपानकैः ।}\ H}}\skt{ ॥}\pend

\pstart\versepar \ledsidenote{\passref{6.55.44}}\edtext{\skt{हिंगु\-कु\-ष्ठ\-व\-चा\-स्व\-र्जि\-वि\-ड\-ञ्चापि}}{\lemma{\skt{चापि}}\Afootnote{K H; \skt{गं वा}\ A}}\skt{ द्विरुत्तरं । }\pend
\pstart\versepar \skt{}\edtext{\edtext{\skt{योगा\-वे\-ता\-वु\-दा\-वर्त्तं}}{\lemma{\skt{एताव्}}\Afootnote{A K; \skt{एता}\ H}}}{\lemma{\skt{उदावर्त्तं}}\Afootnote{A K; \skt{नुदावर्त्तं}\ H}}\skt{ शूलं }\edtext{\skt{चातिबलं}}{\Afootnote{K H; \skt{चानिलजं}\ A}}\skt{ }\edtext{\skt{जेयेत्}}{\Afootnote{K H; \skt{हतः ॥}\ A}}\skt{ ।}\pend

\pstart\versepar \ledsidenote{\passref{6.55.45}}\edtext{\skt{देवदार्व्यग्निकौ}}{\Afootnote{K H; \skt{देवदार्वग्निकौ}\ A}}\skt{ कुष्ठं शुण्ठीं पथ्यां }\edtext{\skt{पलंकशां}}{\Afootnote{K; \skt{पलङ्कषाम्}\ A H}}\skt{ । }\pend
\pstart\versepar \skt{}\edtext{\skt{पौष्करस्य च}}{\Afootnote{K H; \skt{पौष्कराणि च}\ A}}\skt{ मूलानि }\edtext{\skt{तोयस्यार्धाढके}}{\Afootnote{K H; \skt{तोयस्यर्धाढके}\ A}}\skt{ पचेत् ।}\pend

\pstart\versepar \ledsidenote{\passref{6.55.46}}\edtext{\skt{पादावशिष्टं}}{\Afootnote{A K; \skt{पादावसिष्टन्}\ H}}\skt{ तत्पी\-त\-मु\-दा\-वर्त्तं }\edtext{\skt{व्यपोहति}}{\Afootnote{K H; \skt{अपोहति ।}\ A}}\skt{ । }\pend
\pstart\versepar \skt{मूलकं शुष्कमार्द्रञ्च }\edtext{\skt{वर्षा\-भू\-मू\-ल\-प\-ञ्चकं}}{\Afootnote{K; \skt{वर्षाभूः पञ्चमूलकम् ॥}\ A; \skt{वर्षा\-भू\-ल\-प\-ञ्चकं ।}\ H}}\skt{ ।}\pend

\pstart\versepar \ledsidenote{\passref{6.55.47}}\skt{आरे\-व\-त\-फ\-ल\-ञ्चाप्सु }\edtext{\skt{पक्त्वा}}{\Afootnote{A K; \skt{यत्क्ष}\ H}}\skt{ तेन घृतम्पचेत् । }\pend
\pstart\versepar \skt{तत्पीयमानं }\edtext{\edtext{\skt{हंत्यु\-ग्र\-मु\-दा\-व\-र्त्त\-म\-शे\-षतः}}{\lemma{\skt{हंत्य्}}\Afootnote{K; \skt{शास्त्य्}\ A; \skt{कं}\ H}}}{\lemma{\skt{उग्रम्}}\Afootnote{A K; \skt{त्युग्रंम्}\ H\textsuperscript{ac} H\textsuperscript{pc}}}\skt{ ।}\edtext{}{\lemma{\skt{पचेत् । तत्पीयमानं}}\Afootnote{A K; \skt{पचेत् ॥ तत्पीयमान}\ H}}\pend

\pstart\versepar \ledsidenote{\passref{6.55.48}}\edtext{\edtext{\skt{वचा\-म\-ति\-वि\-ष\-म्पाठां}}{\lemma{\skt{अतिविषम्}}\Afootnote{K; \skt{अतिविषां}\ A H}}}{\lemma{\skt{पाठां}}\Afootnote{K H; \skt{कुष्ठं}\ A}}\skt{ यवक्षारं हरीतकीं । }\pend
\pstart\versepar \skt{}\edtext{\skt{कृष्णा\-न्नि\-र्द\-ह\-नी\-ञ्चैव}}{\lemma{\skt{चैव}}\Afootnote{K H; \skt{चापि}\ A}}\skt{ पिबेदुष्णेन वारिणा ।}\pend

\pstart\versepar \ledsidenote{\passref{6.55.49}}\skt{इक्ष्वाकुमूलं मदनं विशल्यातिविषे वचां । }\pend
\pstart\versepar \skt{}\edtext{\skt{कुष्ठं}}{\Afootnote{A K; \skt{कुष्ठ}\ H}}\skt{ किण्वाग्निकौ }\edtext{\skt{चापि}}{\Afootnote{K H; \skt{चैव}\ A}}\skt{ पिबेत्तुल्यानि पूर्ववत् ।}\pend

\pstart\versepar \ledsidenote{\passref{6.55.50}}\skt{मूत्रेण }\edtext{\skt{देव\-दा\-र्व\-ग्नि\-त्रि\-फ\-ला\-बृ\-हतीं}}{\Afootnote{K H; \skt{देव\-दा\-र्व\-ग्नि\-त्रि\-फ\-ला\-बृ\-हतीः}\ A}}\skt{ पिबेत् । }\pend
\pstart\versepar \skt{यच्च प्रस्थं }\edtext{\skt{पले द्वे}}{\lemma{\skt{पले च}}\Afootnote{K; \skt{फलैः सार्ध}\ A; \skt{बले द्वे}\ H}}\skt{ च कण्टकार्या जलाढके ।}\edtext{}{\lemma{\skt{पिबेत् । यच्च}}\Afootnote{K H; \skt{पिबेत् । यव}\ A}}\edtext{}{\lemma{\skt{द्वे}}\Afootnote{K H; \skt{ं}\ A}}\pend

\pstart\versepar \ledsidenote{\passref{6.55.51}}\edtext{\skt{पक्त्वा\-र्ध\-प्र\-स्थ\-शेषं तु}}{\Afootnote{A K; \skt{पक्त्वा\-र्द्ध\-म्प्र\-स्थ\-शे\-षन्तु}\ H}}\skt{ }\edtext{\skt{पिबे\-द्धिं\-गु\-स\-मा\-यतं}}{\Afootnote{K H; \skt{पिबे\-द्धि\-ङ्गु\-स\-म\-न्वि\-तम् ।}\ A}}\skt{ । }\pend
\pstart\versepar \skt{मद\-ना\-ला\-बु\-बी\-जानि }\edtext{\skt{पिप्पलीं}}{\Afootnote{A K; \skt{पिप्पली}\ H}}\skt{ सनिदिग्धिकां ।}\pend

\pstart\versepar \ledsidenote{\passref{6.55.52}}\skt{सञ्चूर्ण्ण्य }\edtext{\skt{नाड्यप्रधमे}}{\Afootnote{K H; \skt{प्रधमेन्}\ A}}\skt{ }\edtext{\skt{द्विषेत्येतद्यथा}}{\lemma{\skt{द्विषेत्य्}}\Afootnote{K H; \skt{नाड्या विशत्य्}\ A}}\skt{ गुदं । }\pend
\pstart\versepar \skt{}\edtext{\skt{चूर्ण्णं न्निकु\-म्भ\-क\-म्पि\-ल्य\-श्या\-मे\-क्त्वा\-क्व\-ग्नि\-को\-द्भवं}}{\Afootnote{K; \skt{चूर्णं निकु\-म्भ\-क\-म्पि\-ल्ल\-श्या\-मे\-क्ष्वा\-क्व\-ग्नि\-को\-द्भ\-वम् ॥}\ A; \skt{चूर्ण्णं निकु\-म्भ\-क\-म्पि\-ल्य\-श्या\-मे\-क्वा\-न्था\-ग्नि\-को\-द्भवं ।}\ H}}\skt{ ।}\pend

\pstart\versepar \ledsidenote{\passref{6.55.53}}\edtext{\skt{कृतवेधनस्य}}{\Afootnote{K H; \skt{कृतवेधन}\ A}}\skt{ }\edtext{\skt{कृष्णाया}}{\Afootnote{K H; \skt{मागध्योर्}\ A}}\skt{ लवणानाञ्च साधयेत् । }\pend
\pstart\versepar \skt{गवां मूत्रेण ता }\edtext{\skt{वर्त्यः}}{\Afootnote{K H; \skt{वर्तीः}\ A}}\skt{ }\edtext{\skt{कारयेत}}{\Afootnote{K H; \skt{कारयेत्तु}\ A}}\skt{ }\edtext{\skt{गुदंगमाः}}{\Afootnote{K H; \skt{गुदानुगाः ।}\ A}}\skt{ । }\pend
\pstart\versepar \skt{सद्यः शर्मकरावेतौ }\edtext{\edtext{\skt{योगा\-व\-मृ\-त\-स\-म्भ\-वा\-विति}}{\lemma{\skt{अमृ\-त\-स\-म्भ\-वाव्}}\Afootnote{K H; \skt{अमृतसंमतौ ॥}\ A}}\skt{ ॥ {\symfont ❈}}}{\lemma{\skt{इति ॥ {\symfont ❈}}}\Afootnote{K H; \emph{om.}\ A}}\skt{ ॥}\pend

\pstart \ledsidenote{\passref{6.55.trailer}}\edtext{\skt{कायचि ॥ १८ ॥ {\symfont ❈}}}{\Afootnote{K; \skt{इति सुश्रु\-त\-सं\-हि\-ता\-या\-मु\-त्त\-र\-त\-न्त्रा\-न्त\-र्गते काय\-चि\-कि\-त्सा\-तन्त्रे उदा\-व\-र्त\-प्र\-ति\-षेधो नाम (स\-प्त\-द\-शो+\-आ\-ध्यायः, आदितः) पञ्च\-प\-ञ्चा\-श\-त्त\-मोऽ\-ध्यायः ॥५५ ॥}\ A; \skt{उत्तरतन्त्रे त्रयः\-प\-ञ्चा\-श\-त्तमः काय\-चि\-कि\-त्सा\-या\-मु\-दा\-व\-र्त्त\-चि\-कि\-त्सा\-ष्टा\-द\-श\-मोऽ\-ध्यायः ॥ {\symfont ❈} ॥}\ H}}\skt{ ॥}\pend

\endnumbering
\beginnumbering
\eledsection*{Adhyāya 56}
\pstart \ledsidenote{\passref{6.56.1}}\skt{अथातो }\edtext{\skt{विषू\-चि\-का\-प्र\-ति\-षेधं}}{\Afootnote{K; \skt{विसू\-चि\-का\-प्र\-ति\-षे\-धम्}\ A; \skt{विशू\-चि\-का\-प्र\-ति\-षे\-धम्}\ H}}\skt{ }\edtext{\skt{व्याव्याख्यास्यामः}}{\lemma{\skt{व्या}}\Afootnote{K; \skt{अध्यायं व्याख्यास्यामः ॥}\ A; \skt{व्याख्यास्यामः ॥}\ H}}\skt{ ॥}\pend

\pstart\versepar \ledsidenote{\passref{6.56.3}}\skt{अजीर्ण्णमामं विष्टब्धं विदग्धञ्च यदीरितं । }\pend
\pstart\versepar \skt{}\edtext{\skt{विषूच्यलसकौ}}{\Afootnote{K; \skt{विसूच्यलसकौ}\ A; \skt{विशूच्यलसकौ}\ H}}\skt{ तस्माद्भवेच्चापि }\edtext{\skt{विडंबिका}}{\Afootnote{K H; \skt{विलम्बिका ॥}\ A}}\skt{ ।}\pend

\pstart\versepar \ledsidenote{\passref{6.56.4}}\edtext{\skt{सूचीभिरिव}}{\lemma{\skt{इव}}\Afootnote{A K\textsuperscript{pc} H; \skt{इरव}\ K\textsuperscript{ac}}}\skt{ गात्राणि }\edtext{\skt{तुदं}}{\Afootnote{K; \skt{तुद}\ A H}}\skt{ }\edtext{\skt{सन्तिष्ठते}}{\Afootnote{K; \skt{न्संतिष्ठते}\ A H}}\skt{ }\edtext{\skt{निलः}}{\Afootnote{K; \skt{ऽनिलः ।}\ A H}}\skt{ । }\pend
\pstart\versepar \skt{}\edtext{\edtext{\skt{यस्यामजीर्ण्णात्सा}}{\lemma{\skt{यस्याम्}}\Afootnote{K H; \skt{यस्या}\ A}}\skt{ }\edtext{\skt{वैद्यैर्विषूचीति}}{\lemma{\skt{विषूचीति}}\Afootnote{K; \skt{विसूचिका ॥}\ A; \skt{व्विसूचीति}\ H}}}{\lemma{\skt{अजीर्ण्णात्सा वैद्यैर्}}\Afootnote{K H; \skt{जीर्णेन सा वैद्यैरुच्यते ति}\ A}}\skt{ }\edtext{\skt{निरुच्यते}}{\Afootnote{K H; \emph{om.}\ A}}\skt{ ।}\pend

\pstart\versepar \ledsidenote{\passref{6.56.5}}\skt{न ताम्प\-रि\-मि\-ता\-हारा लभन्ते विदितागमाः । }\pend
\pstart\versepar \skt{मूढा\-स्ता\-म\-जि\-ता\-त्मानो लभन्ते }\edtext{\skt{शनलोलुपाः}}{\Afootnote{K H\textsuperscript{pc}; \skt{कलुषाशयाः ॥}\ A; \skt{लोलुपाः ॥}\ H\textsuperscript{ac}}}\skt{ ॥}\pend

\pstart\versepar \ledsidenote{\passref{6.56.6}}\skt{मूर्च्छातिसारौ वमथुः पिपासा }\edtext{\skt{शूल\-भ्र\-मो\-द्वे\-ष्ट\-न\-जृ\-म्भ\-दाहाः}}{\Afootnote{K H; \skt{शूलं भ्रमो\-द्वे\-ष्ट\-न\-जृ\-म्भ\-दाहाः ।}\ A}}\skt{ । }\pend
\pstart\versepar \skt{वैवर्ण्ण्यकम्पौ हृदये रुजश्च भवन्ति तस्यां शिरसश्च भेदः ।}\pend

\pstart\versepar \ledsidenote{\passref{6.56.7}}\skt{कुक्षिरानह्यते }\edtext{\skt{त्यर्थं ताम्यते}}{\Afootnote{K; \skt{ऽत्यर्थं प्रताम्यति}\ A; \skt{ऽत्यर्थं ताम्यते}\ H}}\skt{ }\edtext{\skt{परिकूजति}}{\Afootnote{K\textsuperscript{pc} H; \skt{विकूजति ।}\ A; \skt{परिकूजिति ।}\ K\textsuperscript{ac}}}\skt{ । }\pend
\pstart\versepar \skt{निरुद्धो }\edtext{\skt{मारुतश्चैव}}{\lemma{\skt{चैव}}\Afootnote{K H; \skt{चापि}\ A}}\skt{ }\edtext{\skt{कुक्षावुपरि}}{\lemma{\skt{उपरि}}\Afootnote{K H; \skt{विपरि}\ A}}\skt{ धावति ।}\pend

\pstart\versepar \ledsidenote{\passref{6.56.8}}\edtext{\skt{वातवर्चो}}{\Afootnote{K H; \skt{वातवर्चो}\ A}}\skt{ निरोधश्च कुक्षौ यस्य भृशम्भवेत् । }\pend
\pstart\versepar \skt{तस्या\-ल\-स\-क\-मा\-चष्टे तृष्णोद्गारौ }\edtext{\skt{च}}{\Afootnote{K H; \skt{अ}\ A}}\skt{ }\edtext{\skt{यस्य}}{\Afootnote{K H; \skt{रोध}\ A}}\skt{ }\edtext{\skt{तु}}{\Afootnote{K H; \skt{कौ ॥}\ A}}\skt{ ।}\pend

\pstart\versepar \ledsidenote{\passref{6.56.9}}\edtext{\skt{तुष्टन्तु}}{\Afootnote{K; \skt{दुष्टं तु}\ A; \skt{दुष्टञ्च}\ H}}\skt{ भुक्त.\-म्क\-फ\-मा\-रु\-ताभ्यां प्रवर्त्तते नोर्ध्वमधश्च यस्य । }\pend
\pstart\versepar \skt{}\edtext{\skt{विडम्बिकां}}{\Afootnote{K H; \skt{विलम्बिकां}\ A}}\skt{ }\edtext{\skt{तस्य}}{\Afootnote{K H; \skt{तां भृ}\ A}}\skt{ }\edtext{\skt{सुदु\-श्चि\-त्सा\-मा\-च\-क्षते}}{\lemma{\skt{सुदुश्चित्साम्}}\Afootnote{K; \skt{शदुश्चिकित्स्याम्}\ A; \skt{सुदुश्चिकित्स्याम्}\ H}}\skt{ शास्त्रविदः पुराणाः ।}\pend

\pstart\versepar \ledsidenote{\passref{6.56.10}}\edtext{\edtext{\skt{यत्रस्थमामं}}{\lemma{\skt{यत्रस्थम्}}\Afootnote{K H; \skt{यत्रस्थ}\ A}}}{\lemma{\skt{आमं}}\Afootnote{K H; \skt{मामं}\ A}}\skt{ }\edtext{\skt{हि}}{\Afootnote{K H; \skt{वि}\ A}}\skt{ }\edtext{\skt{रुजार्त्तमेव}}{\lemma{\skt{रुजार्त्तम्}}\Afootnote{K H; \skt{रुजेत्तम्}\ A}}\skt{ देशं विशेषेण विकारजातैः । }\pend
\pstart\versepar \skt{दोषेण }\edtext{\skt{येनानुगतं}}{\Afootnote{K H; \skt{येनावततं}\ A}}\skt{ }\edtext{\skt{च}}{\Afootnote{K H; \skt{स्व}\ A}}\skt{ }\edtext{\skt{गाढं तं}}{\Afootnote{K H; \skt{लिङ्गैस्तं}\ A}}\skt{ }\edtext{\edtext{\skt{लक्ष\-णै\-रा\-म\-स\-मु\-द्भ\-वैश्च}}{\lemma{\skt{लक्षणैर्}}\Afootnote{K H; \skt{लक्षये}\ A}}}{\lemma{\skt{आम\-स\-मु\-द्भ\-वैश्}}\Afootnote{K H; \skt{दाम\-स\-मु\-द्भ\-वैश्}\ A}}\skt{ ।}\pend

\pstart\versepar \ledsidenote{\passref{6.56.11}}\skt{यः }\edtext{\skt{स्याव\-द\-न्तौ\-ष्ठ\-नखो}}{\Afootnote{K; \skt{श्याव\-द\-न्तौ\-ष्ठ\-नखो}\ A H}}\skt{ }\edtext{\skt{ल्पसंज्ञः}}{\Afootnote{K H; \skt{ऽल्पसंज्ञश्}\ A}}\skt{ च्छर्द्यर्दितो }\edtext{\skt{भ्यन्त\-र\-या\-त\-नेत्रः}}{\Afootnote{K; \skt{ऽभ्य\-न्त\-र\-या\-त\-नेत्रः}\ A H}}\skt{ }\pend
\pstart\versepar \skt{क्षामस्वरः । }\edtext{\skt{सर्व\-वि\-मु\-क्त\-स\-न्धि\-र्या\-या\-न्नरो सोऽपु\-न\-रा\-ग\-माय}}{\lemma{\skt{नरो सो}}\Afootnote{K; \skt{नरः सो}\ A H}}\skt{ ।}\pend

\pstart\versepar \ledsidenote{\passref{6.56.12}}\edtext{\skt{साध्यस्य}}{\Afootnote{K H; \skt{साध्यासु}\ A}}\skt{ }\edtext{\skt{पार्ष्ण्योर्दहनं}}{\lemma{\skt{पार्ष्ण्योर्}}\Afootnote{A K; \skt{पार्ष्ण्यो}\ H}}\skt{ प्रश\-स्त\-म\-ग्नि\-प्र\-तापो वमनञ्च तीक्ष्णं । }\pend
\pstart\versepar \skt{पक्वे ततो }\edtext{\skt{न्ने तु}}{\Afootnote{A K; \skt{ऽन्नं तु}\ H}}\skt{ }\edtext{\skt{विलंघनं}}{\Afootnote{K; \skt{विलङ्घनं}\ A; \skt{विलंघनः}\ H}}\skt{ }\edtext{\skt{स्यात्स\-म्पा\-च\-न\-म्वापि}}{\lemma{\skt{वापि}}\Afootnote{K H; \skt{चापि}\ A}}\skt{ }\edtext{\skt{विरेचनम्वा}}{\lemma{\skt{वा}}\Afootnote{K H; \skt{च ॥}\ A}}\skt{ ।}\pend

\pstart\versepar \ledsidenote{\passref{6.56.13}}\skt{विशुद्धदेहस्य हि सद्य एव मूर्च्छा\-ति\-सा\-रा\-दि\-रु\-पैति शान्तिं । }\pend
\pstart\versepar \skt{आस्थापनञ्चापि }\edtext{\skt{हितम्वदन्ति}}{\lemma{\skt{हितम्}}\Afootnote{K H; \emph{om.}\ A}}\skt{ }\edtext{\skt{सर्वासु}}{\Afootnote{K H; \skt{पथ्यं सर्वासु}\ A}}\skt{ }\edtext{\edtext{\skt{योगा\-न\-प\-रा\-न्नि\-बोध}}{\lemma{\skt{योगानपरान्}}\Afootnote{A K; \skt{योगान्न परान्}\ H}}}{\lemma{\skt{निबोध}}\Afootnote{A K; \skt{निबोधः ॥}\ H}}\skt{ ॥}\pend

\pstart \ledsidenote{\passref{6.56.14}}\skt{पथ्या वचा हिंगु }\edtext{\skt{विडं}}{\Afootnote{K H; \skt{कलि\-ङ्ग\-गृ\-ञ्ज\-सौ\-व\-र्चलैः सातिविषैश्}\ A}}\skt{ }\edtext{\skt{विचू\-र्ण्ण्य\-सु\-खा\-म्बु\-ना\-सा\-ति\-विषं}}{\Afootnote{K; \skt{च चूर्णम् । सुखाम्बुपीतं विनिहन्}\ A; \skt{विचू\-र्ण्ण\-सु\-खा\-म्बु\-ना\-सा\-ति\-वि\-षन्}\ H}}\skt{ }\edtext{\skt{निषेवेत्}}{\Afootnote{K H; \skt{त्यजीर्णं शूलं विसूचीमरुचिं च सद्यः ॥}\ A}}\skt{ ।}\pend

\pstart\versepar \ledsidenote{\passref{6.56.15}}\skt{क्षारागदम्वा लवणं }\edtext{\skt{विडम्वा}}{\Afootnote{K H; \skt{विषं वा}\ A}}\skt{ गुड\-प्र\-गा\-ढा\-नथ }\edtext{\skt{सर्षपम्वा}}{\Afootnote{K H; \skt{सर्षपान्वा ।}\ A}}\skt{ । }\pend
\pstart\versepar \skt{अम्लेन वा }\edtext{\skt{सैन्ध\-व\-हिं\-गु\-युक्तौ}}{\Afootnote{A K; \skt{सैन्ध\-व\-हिं\-गु\-युक्तैः}\ H}}\skt{ सबीजपूर्ण्णौ }\edtext{\skt{सुकृतौ}}{\Afootnote{K H; \skt{सघृतौ}\ A}}\skt{ त्रिवर्गौ ।}\pend

\pstart \ledsidenote{\passref{6.56.16}}\skt{कटुत्रिकम्वा लवणैरुपेतं }\edtext{\skt{पिबे\-त्स्नु\-ही\-क्षी\-र\-वि\-मि\-श्रि\-तम्वा}}{\lemma{\skt{वा}}\Afootnote{K H; \skt{तु । कल्याणकं वा लवणं पिबेत्तु यदु\-क्त\-मा\-दा\-व\-नि\-ला\-म\-येषु ॥}\ A}}\skt{ ।}\pend

\pstart \ledsidenote{\passref{6.56.17}}\edtext{\skt{कृष्णा\-ज\-मो\-दा\-क्ष\-स\-मा\-ग्नि\-कानि}}{\Afootnote{K H; \skt{कृष्णा\-ज\-मो\-द\-क्ष\-व\-काणि वापि}\ A}}\skt{ }\edtext{\skt{तुल्यं}}{\Afootnote{K H; \skt{तुल्यौ}\ A}}\skt{ पिबेद्वा }\edtext{\skt{मगधानिकुम्भौ}}{\Afootnote{K H; \skt{मगधानिकुम्भौ । दन्तीयुत वा मगधोद्भवानां कल्कं पिबे\-त्को\-ष\-व\-ती\-र\-सेन ॥}\ A}}\skt{ ।}\pend

\pstart\versepar \ledsidenote{\passref{6.56.18}}\skt{उष्णा\-भि\-र\-द्भि\-र्म\-ग\-धो\-द्भ\-वानां }\edtext{\edtext{\skt{कल्क\-म्पि\-बे\-द्वा\-प्यथ}}{\lemma{\skt{पिबेद्वाप्य्}}\Afootnote{K H; \skt{पिबेन्}\ A}}}{\lemma{\skt{अथ}}\Afootnote{K H; \emph{om.}\ A}}\skt{ }\edtext{\skt{नागदन्त्या}}{\Afootnote{K H; \skt{नाग\-र\-क\-ल्क\-यु\-क्तम् ।}\ A}}\skt{ । }\pend
\pstart\versepar \skt{व्योषं }\edtext{\skt{करञ्जस्य}}{\Afootnote{A K; \skt{करञ्जास्य}\ H}}\skt{ फलं }\edtext{\skt{हरिद्रां}}{\Afootnote{K H; \skt{हरिद्रे}\ A}}\skt{ मूलं समं }\edtext{\skt{वाप्यथ}}{\lemma{\skt{वाप्य्}}\Afootnote{K H; \skt{चाप्य्}\ A}}\skt{ मातुलुङ्ग्याः ।}\pend

\pstart\versepar \ledsidenote{\passref{6.56.19}}\skt{च्छायाविशुष्का }\edtext{\skt{गुडिका}}{\Afootnote{K H; \skt{गुटिकाः}\ A}}\skt{ कृतास्ताः }\edtext{\skt{हन्युर्विसूचिं}}{\lemma{\skt{विसूचिं}}\Afootnote{K; \skt{विसूचीं}\ A H}}\skt{ }\edtext{\skt{नयनाञ्जनेन}}{\Afootnote{A K\textsuperscript{pc} H; \skt{नयनाञ्जनेन ः}\ K\textsuperscript{ac}}}\skt{ । }\pend
\pstart\versepar \skt{सुवामितं साधु विरेचितं वा }\edtext{\skt{सुलङ्घितम्वा}}{\Afootnote{K; \skt{सुलङ्घतं वा}\ A; \skt{सुलंघितम्वा}\ H}}\skt{ मनुजं विदित्वा ।}\pend

\pstart\versepar \ledsidenote{\passref{6.56.20}}\skt{पेया\-दि\-भि\-र्दी\-प\-न\-पा\-च\-नीयैः सम्यक्क्षुधार्त्तं समुपक्रमेत ॥ }\pend
\pstart\versepar \skt{आमं शकृद्वा }\edtext{\skt{निचि भूयो}}{\Afootnote{\skt{निचितं क्रमेण भूयो}\ A H; \skt{निचि… भूयो}\ K}}\skt{ विबद्धं विगुणानिलेन ।}\pend

\pstart\versepar \ledsidenote{\passref{6.56.21}}\skt{प्रवर्त्तमानं न }\edtext{\edtext{\skt{यथा\-स्त\-मे\-त\-द्वि\-का\-र\-मा\-ना\-ह\-मु\-दा\-ह\-रन्ति}}{\lemma{\skt{यथास्तम्}}\Afootnote{K H; \skt{यथास्व}\ A}}}{\lemma{\skt{एतद्विकारम्}}\Afootnote{K H; \skt{मेनं विकारम्}\ A}}\skt{ । }\pend
\pstart\versepar \skt{}\edtext{\skt{तस्मिं}}{\Afootnote{K; \skt{तस्मि}\ A H}}\skt{ }\edtext{\edtext{\skt{भव\-त्या\-म\-स\-मु\-द्भवे}}{\lemma{\skt{भवत्य्}}\Afootnote{K; \skt{न्भवन्}\ A; \skt{न्भवन्त्य्}\ H}}\skt{ तु}}{\lemma{\skt{आमसमुद्भवे तु}}\Afootnote{K H; \skt{त्यामसमुद्भवे तु}\ A}}\skt{ }\edtext{\skt{भ्रम\-प्र\-ति\-श्या\-य\-शि\-रो\-वि\-दाहाः}}{\Afootnote{K H; \skt{तृष्णा\-प्र\-ति\-श्या\-य\-शि\-रो\-वि\-दाहाः ॥}\ A}}\skt{ ।}\pend

\pstart\versepar \ledsidenote{\passref{6.56.22}}\skt{आमाशये शूलमथो गुरुत्वं }\edtext{\skt{हृत्स्त\-म्भ\-मु\-द्गा\-र\-वि\-घा\-त\-नञ्च}}{\lemma{\skt{हृत्स्तम्भम्}}\Afootnote{K H; \skt{हृल्लास}\ A}}\skt{ । }\pend
\pstart\versepar \skt{स्तम्भः }\edtext{\skt{कटी\-पृ\-ष्ठ\-पु\-री\-श\-मूत्रे}}{\Afootnote{K; \skt{कटी\-पृ\-ष्ठ\-पु\-री\-ष\-मूत्रे}\ A H}}\skt{ }\edtext{\skt{शूलो}}{\Afootnote{K H; \skt{शूलो}\ A}}\skt{ }\edtext{\skt{थ}}{\Afootnote{K H; \skt{ऽथ}\ A}}\skt{ मूर्च्छा स शकृद्वमेच्च ।}\pend

\pstart\versepar \ledsidenote{\passref{6.56.23}}\skt{श्वासश्च पक्वाशयजे भवन्ति लिङ्गानि चात्रा\-ल\-स\-को\-द्भ\-वानि । }\pend
\pstart\versepar \skt{आमोद्भवे }\edtext{\skt{वान्त\-मु\-प\-क्र\-मेत}}{\lemma{\skt{वान्तम्}}\Afootnote{A K; \skt{द्वान्तम्}\ H}}\skt{ }\edtext{\skt{संसक्तभक्तं}}{\Afootnote{K H; \skt{संसर्गभक्त}\ A}}\skt{ क्रमदीपनीयैः ।}\pend

\pstart\versepar \ledsidenote{\passref{6.56.24}}\skt{अथेतरं यो न }\edtext{\skt{शकृद्वमेत}}{\Afootnote{K H; \skt{शकृद्वमेत्तम्}\ A}}\skt{ }\edtext{\skt{तं}}{\Afootnote{K H; \skt{आ}\ A}}\skt{ }\edtext{\edtext{\skt{साध\-ये\-त्स्वि\-न्न\-मिमैः}}{\lemma{\skt{साधयेत्स्विन्नम्}}\Afootnote{K H; \skt{मं जयेत्स्वेद}\ A}}}{\lemma{\skt{इमैः}}\Afootnote{K H; \skt{नपा}\ A}}\skt{ }\edtext{\skt{प्रयोगैः}}{\Afootnote{K H; \skt{चनैश्च ।}\ A}}\skt{ । }\pend
\pstart\versepar \skt{}\edtext{\skt{विषूचिकायां}}{\Afootnote{K; \skt{विसूचिकायां}\ A H}}\skt{ परिकीर्त्तितानि द्रव्याणि वैरेचनिकानि यानि ।}\pend

\pstart\versepar \ledsidenote{\passref{6.56.25}}\edtext{\skt{संचूर्ण्ण्य}}{\Afootnote{K H; \skt{तान्येव}\ A}}\skt{ }\edtext{\edtext{\skt{वर्त्ति\-म्वि\-त\-रे\-द्वि\-धिज्ञो}}{\lemma{\skt{वर्त्तिम्}}\Afootnote{K H; \skt{वर्तीर्}\ A}}}{\lemma{\skt{वितरेद्विधिज्ञो}}\Afootnote{K H; \skt{विरचेद्विचूर्ण्य}\ A}}\skt{ }\edtext{\skt{महि\-ष्य\-जा\-व्यु\-ष्ट्र\-ग\-वांश्च}}{\Afootnote{K; \skt{महि\-ष्य\-जा\-वी\-भ\-गवां तु}\ A; \skt{महि\-ष्य\-जा\-व्यु\-ष्ट्र\-ग\-वाञ्च}\ H}}\skt{ मूत्रैः । }\pend
\pstart\versepar \skt{}\edtext{\skt{स्विन्नस्य}}{\Afootnote{A K; \skt{स्विन्न}\ H}}\skt{ पायौ विनिवेश्य ताश्च चूर्ण्णानि चैषां प्रधमेत्तु नाड्या ।}\pend

\pstart\versepar \ledsidenote{\passref{6.56.26}}\edtext{\skt{मूत्रेण}}{\Afootnote{K H; \skt{मूत्रेषु}\ A}}\skt{ }\edtext{\skt{वा}}{\Afootnote{K H; \skt{सं}\ A}}\skt{ }\edtext{\skt{साध्य}}{\Afootnote{A K; \skt{साध्या}\ H}}\skt{ }\edtext{\skt{यथाविधाणं}}{\Afootnote{K; \skt{यथाविधानं}\ A H}}\skt{ }\edtext{\skt{द्रव्यानि}}{\Afootnote{K H; \skt{द्रव्याणि}\ A}}\skt{ यान्यूर्ध्वमधश्च यान्ति । }\pend
\pstart\versepar \skt{क्वाथेन }\edtext{\skt{तेनाशु}}{\Afootnote{A K H\textsuperscript{pc}; \skt{तेनाशु ते}\ H\textsuperscript{ac}}}\skt{ }\edtext{\skt{निरूहयीत}}{\Afootnote{K H; \skt{निरूहयेच्च}\ A}}\skt{ }\edtext{\skt{मात्रार्धयुक्तेन}}{\Afootnote{K H; \skt{मूत्रार्धयुक्तेन}\ A}}\skt{ समाक्षिकेण ।}\pend

\pstart\versepar \ledsidenote{\passref{6.56.27}}\edtext{\skt{त्रिभण्डिपिष्टं}}{\Afootnote{K H; \skt{त्रिभण्डियुक्तं}\ A}}\skt{ लवणप्रकुञ्चं दत्वा }\edtext{\skt{विरिक्तं}}{\Afootnote{K H; \skt{विरिक्त}\ A}}\skt{ }\edtext{\skt{क्रममादिशेच्च}}{\lemma{\skt{आदिशेच्च}}\Afootnote{K H; \skt{आचरेच्च ।}\ A}}\skt{ । }\pend
\pstart\versepar \skt{}\edtext{\edtext{\skt{एतेष्वेव}}{\lemma{\skt{एतेष्व्}}\Afootnote{K H; \skt{एष्व्}\ A}}\skt{ च}}{\lemma{\skt{एव च}}\Afootnote{K H; \skt{एव}\ A}}\skt{ }\edtext{\skt{तैलेन}}{\Afootnote{K H; \skt{तैलेन च}\ A}}\skt{ साधितेन }\edtext{\skt{सुखाम्बुना}}{\Afootnote{K H; \emph{om.}\ A}}\skt{ । प्राप्तं }\edtext{\skt{यदि न च}}{\Afootnote{K H; \skt{यदि}\ A}}\skt{ }\edtext{\skt{स्याद्वा}}{\Afootnote{K H; \emph{om.}\ A}}\skt{ }\edtext{\edtext{\edtext{\edtext{\skt{भागे\-ष्वे\-ष्व\-नु\-वा\-स\-न\-मिति}}{\lemma{\skt{भागेष्व्}}\Afootnote{K H; \emph{om.}\ A}}}{\lemma{\skt{एष्व्}}\Afootnote{K; \skt{स्या}\ A; \skt{ए}\ H}}}{\lemma{\skt{अनुवासनम्}}\Afootnote{K; \skt{दनुवासये}\ A; \skt{व तु वासनम्}\ H}}}{\lemma{\skt{इति}}\Afootnote{K; \skt{च्च ॥}\ A; \skt{इति ॥ {\symfont ❈} ॥}\ H}}\skt{ ॥}\pend

\pstart \ledsidenote{\passref{6.56.trailer}}\edtext{\skt{कायचि}}{\Afootnote{K; \skt{इति सुश्रु\-त\-सं\-हि\-ता\-या\-मु\-त्त\-र\-तन्त्रे काय\-चि\-कि\-त्सा\-तन्त्रे विसू\-चि\-का\-प्र\-ति\-षेधो नाम (अ\-ष्टा\-द\-शोऽ\-ध्यायः, आदितः) षट्प\-ञ्चा\-श\-त्त\-मोऽ\-ध्यायः ॥५६ ॥}\ A; \skt{उत्तरतन्त्रे चतुः\-प\-ञ्चा\-श\-त्तमः कायचिकित्सायां विसूचिचिकित्सा ऊन\-विं\-श\-ति\-त\-मोऽ\-ध्यायः ॥ {\symfont ❈} ॥}\ H}}\skt{ ॥ १९ ॥ ० ॥}\pend

\endnumbering
\beginnumbering
\eledsection*{Adhyāya 57}
\pstart \ledsidenote{\passref{6.57.1}}\skt{अथातो }\edtext{\skt{रोच\-क\-प्र\-ति\-षेधं}}{\Afootnote{K; \skt{ऽरो\-च\-क\-प्र\-ति\-षे\-धम्}\ A H}}\skt{ }\edtext{\skt{व्याख्यास्यामः}}{\Afootnote{K H; \skt{अध्यायं व्याख्यास्यामः ॥}\ A}}\skt{ ॥}\pend

\pstart\versepar \ledsidenote{\passref{6.57.3}}\skt{दोषैः }\edtext{\skt{पृथ\-क्त्रि\-वि\-ध\-सो\-क\-स\-मु\-च्छ्र\-याच्च}}{\Afootnote{K; \skt{पृथक्सह च चित्त\-वि\-प\-र्य\-याच्च}\ A; \skt{पृथ\-क्त्रि\-वि\-ध\-शो\-क\-स\-मु\-च्छ्र\-याच्च}\ H}}\skt{ भक्तायनेषु हृदि चावतते प्रगाढं । }\pend
\pstart\versepar \skt{नान्ने रुचिर्भवति तं }\edtext{\skt{भिषजो}}{\Afootnote{A K\textsuperscript{pc} H; \skt{भिजो}\ K\textsuperscript{ac}}}\skt{ विकारं भक्तो\-प\-घा\-त\-मिह पञ्चविधं वदन्ति । ।}\pend

\pstart\versepar \ledsidenote{\passref{6.57.4}}\edtext{\skt{हृच्छू\-ल\-शो\-ष\-व\-मथुः}}{\Afootnote{K H; \skt{हृच्छू\-ल\-पी\-ड\-न\-युतं}\ A}}\skt{ }\edtext{\skt{स्वरसंक्षयाच्च}}{\Afootnote{K H; \skt{विरसाननत्वं}\ A}}\skt{ वाता }\edtext{\skt{सदा\-ह\-म\-द\-तृ\-ड्व\-मने}}{\Afootnote{K; \skt{त्मके भवति लिङ्गमरोचके तु । हृद्दा\-ह\-चो\-ष\-ब\-हुता मुखतिक्तता च मूर्च्छा सतृड्भ}\ A; \skt{त्सदा\-ह\-म\-द\-तृ\-ड्व\-मने}\ H}}\skt{ }\edtext{\skt{च}}{\Afootnote{K H; \skt{व}\ A}}\skt{ }\edtext{\skt{पित्तं}}{\Afootnote{K H; \skt{ति}\ A}}\skt{ । }\pend
\pstart\versepar \skt{}\edtext{\skt{पित्तात्कफेन}}{\Afootnote{K H; \skt{पित्तकृते त}\ A}}\skt{ }\edtext{\skt{हृदये}}{\Afootnote{K H; \skt{थैव ॥}\ A}}\skt{ }\edtext{\skt{गुरुताविलेपः}}{\Afootnote{K H; \emph{om.}\ A}}\skt{ }\edtext{\skt{श्लेष्म\-प्र\-से\-क\-व\-मथुः}}{\Afootnote{K H; \emph{om.}\ A}}\skt{ }\edtext{\skt{शिरसो}}{\Afootnote{K; \emph{om.}\ A; \skt{सशिरो}\ H}}\skt{ }\edtext{\edtext{\skt{रुजञ्च}}{\lemma{\skt{रुजञ्}}\Afootnote{K H; \emph{om.}\ A}}}{\lemma{\skt{च}}\Afootnote{K H; \emph{om.}\ A}}\skt{ ।}\pend

\pstart\versepar \ledsidenote{\passref{6.57.5}}\edtext{\skt{सर्वात्मकं}}{\Afootnote{K H; \emph{om.}\ A}}\skt{ }\edtext{\edtext{\skt{त्रिवि\-ध\-लिं\-ग\-मथो}}{\lemma{\skt{त्रिविधलिंगम्}}\Afootnote{K H; \emph{om.}\ A}}}{\lemma{\skt{अथो}}\Afootnote{K H; \emph{om.}\ A}}\skt{ }\edtext{\skt{वदन्ति}}{\Afootnote{K H; \emph{om.}\ A}}\skt{ ॥ }\pend
\pstart\versepar \skt{}\edtext{\edtext{\edtext{\skt{शोको\-प\-घा\-त\-ज\-म\-था\-नि\-ल\-ल\-क्ष\-णैस्तु}}{\lemma{\skt{शोको\-प\-घा\-त\-जम्}}\Afootnote{K H; \skt{कण्}\ A}}}{\lemma{\skt{अथा\-नि\-ल\-ल\-क्ष\-णैस्}}\Afootnote{K; \skt{डूगु\-त्व\-क\-फ\-संस्र}\ A; \skt{अथानिलल}\ H}}}{\lemma{\skt{तु}}\Afootnote{K; \skt{व}\ A; \emph{om.}\ H}}\skt{ । }\pend
\pstart\versepar \skt{}\edtext{\skt{वातात्मके}}{\Afootnote{K; \skt{सादतन्द्राः श्लेष्मात्मके}\ A; \skt{ङैश्च । वातात्मकैर्}\ H}}\skt{ }\edtext{\edtext{\edtext{\skt{विर\-स\-मा\-स्य\-म\-रो\-चके}}{\lemma{\skt{विरसम्}}\Afootnote{K H; \skt{मधुर}\ A}}}{\lemma{\skt{आस्यम्}}\Afootnote{K H; \skt{मास्य}\ A}}\skt{ तु}}{\lemma{\skt{अरोचके तु}}\Afootnote{K H; \skt{मरोचके तु ।}\ A}}\skt{ }\edtext{\skt{पित्तेन}}{\Afootnote{K H; \skt{सर्वात्मके पवन}\ A}}\skt{ }\edtext{\skt{तिक्तकटुकं}}{\Afootnote{K H; \skt{पि}\ A}}\skt{ }\edtext{\skt{मधुरं}}{\Afootnote{K H; \skt{त्त}\ A}}\skt{ }\edtext{\skt{कफेन}}{\Afootnote{K H; \skt{कफा बहूनि रूपाण्यथास्य हृदये समुदीरयन्ति ॥}\ A}}\skt{ ।}\pend

\pstart\versepar \ledsidenote{\passref{6.57.6}}\edtext{\edtext{\edtext{\skt{सर्वै\-रु\-पे\-त\-मु\-प\-धा\-र\-य\-स\-न्नि\-पाते}}{\lemma{\skt{सर्वैर्}}\Afootnote{K H; \emph{om.}\ A}}}{\lemma{\skt{उपेतम्}}\Afootnote{K H; \skt{सं}\ A}}}{\lemma{\skt{उप\-धा\-र\-य\-स\-न्नि\-पाते}}\Afootnote{K H; \skt{राग\-शो\-क\-भ\-य\-वि\-प्लु\-तचे}\ A}}\skt{ । }\edtext{\skt{दैन्यं}}{\Afootnote{K; \skt{तसस्तु चिन्}\ A; \skt{दै}\ H}}\skt{ }\edtext{\skt{भृशम्भवति}}{\lemma{\skt{भृशम्}}\Afootnote{K; \skt{ताकृतो}\ A; \skt{न्यं भृशम्}\ H}}\skt{ }\edtext{\skt{शोकसमुद्भवे तु}}{\Afootnote{K H; \emph{om.}\ A}}\skt{ । }\pend
\pstart\versepar \skt{}\edtext{\edtext{\skt{वांतो\-व\-चा\-दि\-र\-निले}}{\lemma{\skt{वांतोवचादिर्}}\Afootnote{K; \skt{सोऽशुचिद}\ A; \skt{वां}\ H}}}{\lemma{\skt{अनिले}}\Afootnote{K; \skt{र्शनाच्च । वाते वचाम्बुव}\ A; \skt{तोव\-चा\-दि\-र\-निले}\ H}}\skt{ }\edtext{\skt{विधिवत्पिबेच्च}}{\Afootnote{K; \skt{मनं कृतवान्पिबेच्च}\ A; \skt{विधिवत्पित्तेन}\ H}}\skt{ }\edtext{\skt{स्नेहो\-ष्ण\-तो\-य\-म\-दि\-रा\-न्य\-त\-मेन}}{\Afootnote{K; \skt{स्नेहैः सुरा\-भि\-र\-थ\-वो\-ष्ण\-ज\-लेन}\ A; \skt{कोष्ण\-तो\-य\-म\-दि\-रा\-न्य\-त\-मेन}\ H}}\skt{ चूर्ण्णं ।}\pend

\pstart\versepar \ledsidenote{\passref{6.57.7}}\skt{कृष्णा\-वि\-डं\-ग\-य\-व\-भ\-स्म\-ह\-रे\-णु\-भा\-र्गी\-रा\-स्नै\-ल\-हिं\-गु\-ल\-व\-णो\-त्त\-म\-ना\-ग\-राणां । }\pend
\pstart\versepar \skt{पित्ते गुडा\-म्बु\-म\-धु\-रै\-र्व\-मनं प्रशस्तं }\edtext{\edtext{\skt{लेहश्च}}{\lemma{\skt{लेहश्}}\Afootnote{K H; \skt{स्नेहः}\ A}}}{\lemma{\skt{च}}\Afootnote{K; \skt{स}\ A; \skt{तु}\ H}}\skt{ सैन्ध\-व\-सि\-ता\-म\-धु\-स\-र्पि\-रिष्टः ।}\pend

\pstart\versepar \ledsidenote{\passref{6.57.8}}\edtext{\skt{निम्बा\-म्बु\-च्छ\-र्दि\-त\-वतः}}{\Afootnote{K H; \skt{निम्बा\-म्बु\-वा\-मि\-त\-वतः}\ A}}\skt{ कफजे }\edtext{\skt{नुपानं}}{\Afootnote{K H; \skt{ऽनुपानं}\ A}}\skt{ राजद्रुमाम्बु मधुना तु }\edtext{\skt{सदीप्यकाढ्यं}}{\Afootnote{K H; \skt{सदीप्यकं स्या}\ A}}\skt{ । }\pend
\pstart\versepar \skt{}\edtext{\skt{चूर्ण्णं}}{\Afootnote{K H; \skt{त् । चूर्णं}\ A}}\skt{ यदुक्तमथ वा }\edtext{\skt{तिलजे}}{\Afootnote{K H; \skt{निलजे}\ A}}\skt{ }\edtext{\skt{तदेव}}{\Afootnote{A K; \skt{तरे च}\ H}}\skt{ सर्वैश्च सर्व\-कृ\-त\-मे\-व\-मु\-प\-क्र\-मेत ।}\pend

\pstart\versepar \ledsidenote{\passref{6.57.9}}\edtext{\skt{द्राक्षा\-प\-टो\-ल\-वि\-ड\-वे\-त्र\-क\-री\-र\-नि\-म्ब\-मू\-र्वा\-भ\-या\-क्ष\-ब\-द\-रा\-ग\-ल\-के\-न्द्र\-वृक्षाः}}{\Afootnote{K; \skt{द्राक्षा\-प\-टो\-ल\-वि\-ड\-वे\-त्र\-क\-री\-र\-नि\-म्ब\-मू\-र्वा\-भ\-या\-क्ष\-ब\-द\-रा\-म\-ल\-के\-न्द्र\-वृक्षैः ।}\ A; \skt{द्राक्षा\-प\-टो\-ल\-वि\-ड\-वे\-त्र\-क\-री\-र\-नि\-म्ब\-मू\-र्व्वा\-भ\-या\-क्ष\-ब\-द\-रा\-म\-ल\-के\-न्द्र\-वृक्षाः ।}\ H}}\skt{ । }\pend
\pstart\versepar \skt{बीजैः कर\-ञ्ज\-नृ\-प\-वृ\-क्ष\-भ\-वैश्च }\edtext{\skt{पिष्टै}}{\Afootnote{K; \skt{पिष्टैर्}\ A H}}\skt{ लेहं पचे\-त्सु\-र\-भि\-मू\-त्र\-युतं यथावत् ।}\pend

\pstart\versepar \ledsidenote{\passref{6.57.10}}\edtext{\edtext{\skt{मुस्ता\-व\-चा\-क\-टु\-क\-रो\-हि\-णि\-नि\-र्द\-ह\-न्य\-स्तु\-ल्याथ}}{\lemma{\skt{मुस्ता\-व\-चा\-क\-टु\-क\-रो\-हि\-णि\-नि\-र्द\-ह\-न्यस्}}\Afootnote{K H; \skt{मुस्तां वचां त्रिकटुकं}\ A}}}{\lemma{\skt{तुल्याथ}}\Afootnote{K H; \emph{om.}\ A}}\skt{ }\edtext{\skt{वापि}}{\Afootnote{K H; \emph{om.}\ A}}\skt{ }\edtext{\skt{रज\-ना\-भ\-य\-सं\-प्र\-युक्ता}}{\Afootnote{K H; \skt{रजनीद्वयं च भार्गीं च कुष्ठमथ निर्दहनीं च पिष्ट्वा ।}\ A}}\skt{ । }\pend
\pstart\versepar \skt{मूत्रे }\edtext{\skt{विके}}{\Afootnote{K H; \skt{ऽविजे}\ A}}\skt{ }\edtext{\skt{द्विर\-द\-मू\-त्र\-युताः}}{\Afootnote{K H; \skt{द्विर\-द\-मू\-त्र\-युते}\ A}}\skt{ पचेद्वा }\edtext{\edtext{\skt{पाठा\-स\-मा\-म\-ति\-विषां}}{\lemma{\skt{पाठासमाम्}}\Afootnote{K H; \skt{पाठां तुगा}\ A}}}{\lemma{\skt{अतिविषां}}\Afootnote{K H; \skt{मतिविषां}\ A}}\skt{ }\edtext{\skt{रजनाञ्च}}{\Afootnote{K H; \skt{रजनीं च}\ A}}\skt{ मुख्यां ।}\pend

\pstart\versepar \ledsidenote{\passref{6.57.11}}\edtext{\edtext{\edtext{\skt{मण्डू\-कि\-म\-र्क\-म\-मृतां}}{\lemma{\skt{मण्डूकिम्}}\Afootnote{K H; \skt{मण्डूकि}\ A}}}{\lemma{\skt{अर्कम्}}\Afootnote{K H; \skt{मर्क}\ A}}}{\lemma{\skt{अमृतां}}\Afootnote{K H; \skt{ममृतां}\ A}}\skt{ }\edtext{\skt{रशलांगलाख्यां}}{\Afootnote{K H; \skt{च सलाङ्गलाख्यां}\ A}}\skt{ मूत्रे }\edtext{\skt{पचेत}}{\Afootnote{K H; \skt{पचेत्तु}\ A}}\skt{ महिषस्य विधाणविद्वान् । }\pend
\pstart\versepar \skt{एतान्न सन्ति चतुरो }\edtext{\skt{भ्यसतस्तु}}{\Afootnote{K H; \skt{लिहतस्तु}\ A}}\skt{ }\edtext{\edtext{\skt{लेहा\-न्गु\-ल्मा\-रु\-चिश्च}}{\lemma{\skt{लेहा\-न्गु\-ल्मा\-रु\-चिश्}}\Afootnote{K H; \skt{लेहान्गुल्मारुचि}\ A}}}{\lemma{\skt{च}}\Afootnote{K H; \skt{श्व}\ A}}\skt{ }\edtext{\skt{शिखि\-सा\-द\-हृ\-दा\-म\-याश्च}}{\lemma{\skt{शिखि\-सा\-द\-हृ\-दा\-म\-याश्}}\Afootnote{K; \skt{सन\-क\-ण्ठ\-हृ\-दा\-म\-याश्}\ A; \skt{शिखि\-सा\-द\-हृ\-दा\-म\-याश्}\ H}}\skt{ ।}\edtext{}{\lemma{\skt{विधाणविद्वान् । एतान्न}}\Afootnote{K; \skt{विधानविद्वा । एतान्न}\ A; \skt{विधानविद्वान् । एतान्न}\ H}}\pend

\pstart\versepar \ledsidenote{\passref{6.57.12}}\edtext{\skt{सात्म्या\-स्व\-दे\-श\-च\-रितां विविधांश्च}}{\Afootnote{K; \skt{सात्म्या\-न्स्व\-दे\-श\-र\-चि\-ता\-न्वि\-वि\-धांश्च}\ A; \skt{सात्म्या\-स्व\-दे\-श\-च\-रि\-ता\-न्वि\-वि\-धाश्च}\ H}}\skt{ }\edtext{\skt{भक्ष्यां पानानि}}{\Afootnote{K; \skt{भक्ष्यान्पानानि}\ A H}}\skt{ }\edtext{\skt{मूल\-फ\-ल\-षा\-ड\-व\-रो\-ग\-योगां}}{\Afootnote{K; \skt{मूल\-फ\-ल\-षा\-ड\-व\-रा\-ग\-योगा}\ A; \skt{मूल\-फ\-ल\-पा\-च\-रा\-योगा ।}\ H}}\skt{ । }\pend
\pstart\versepar \skt{}\edtext{\skt{सेवेद्रसांश्च}}{\Afootnote{K H; \skt{न् । अद्याद्रसांश्च}\ A}}\skt{ }\edtext{\skt{विविधां विविधैः}}{\Afootnote{K; \skt{विविधान्विविधैः}\ A H}}\skt{ प्रकारैर्भुञ्जीत }\edtext{\skt{वापि}}{\Afootnote{K H; \skt{चापि}\ A}}\skt{ लघु\-रू\-क्ष\-म\-नः\-सु\-खानि ।}\pend

\pstart\versepar \ledsidenote{\passref{6.57.13}}\skt{आस्थापनं }\edtext{\edtext{\skt{विविधमत्र}}{\lemma{\skt{विविधम्}}\Afootnote{K H; \skt{विधिव}\ A}}}{\lemma{\skt{अत्र}}\Afootnote{K H; \skt{दत्र}\ A}}\skt{ विरेचनञ्च }\edtext{\skt{कुर्यान्मृदूनि}}{\Afootnote{A K; \skt{कुर्यात्मृदूनि}\ H}}\skt{ शिरसश्च विरेचनानि । }\pend
\pstart\versepar \skt{त्रीण्यूषणानि }\edtext{\skt{त्रिफलां}}{\Afootnote{K H; \emph{om.}\ A}}\skt{ }\edtext{\skt{रजनीद्वयञ्च}}{\Afootnote{K H; \skt{रज\-नी\-त्रि\-फ\-ला\-यु\-तानि}\ A}}\skt{ चूर्ण्णीकृतानि यव\-शू\-क\-वि\-मि\-श्रि\-तानि ।}\pend

\pstart\versepar \ledsidenote{\passref{6.57.14}}\skt{क्षौद्रायुतानि }\edtext{\skt{वित\-रे\-न्मु\-ख\-धा\-व\-ना\-र्थं\-म\-न्यानि}}{\lemma{\skt{वित\-रे\-न्मु\-ख\-धा\-व\-ना\-र्थंम्}}\Afootnote{K; \skt{वित\-रे\-न्मु\-ख\-बो\-ध\-ना\-र्थम्}\ A; \skt{वित\-रे\-त्मु\-ख\-धा\-व\-ना\-र्थम्}\ H}}\skt{ तिक्तकटुकानि च भेषजानि । }\pend
\pstart\versepar \skt{}\edtext{\skt{मुस्ता\-दि\-रा\-ज\-फ\-ल\-व\-र्ग\-द\-शां\-ग\-सिद्धः}}{\Afootnote{K; \skt{मुस्ता\-दि\-रा\-ज\-त\-रु\-व\-र्ग\-द\-शा\-ङ्ग\-सिद्धैः}\ A; \skt{मुस्तानि कफ\-व\-र्ग्ग\-द\-शां\-श\-सिद्धैः}\ H}}\skt{ }\edtext{\edtext{\skt{क्वाथै\-र्ज\-ये\-न्म\-धु\-यु\-तै\-र्वि\-वि\-धैश्च}}{\lemma{\skt{क्वाथैर्}}\Afootnote{A K; \skt{क्वाथेर्}\ H}}}{\lemma{\skt{जये\-न्म\-धु\-यु\-तैर्}}\Afootnote{A K; \skt{जये\-त्म\-धु\-यु\-तैर्}\ H}}\skt{ लेहैः ॥}\pend

\pstart\versepar \ledsidenote{\passref{6.57.15}}\edtext{\edtext{\skt{मूत्रा\-स\-वै\-र्गु\-ड\-कृ\-तैश्च}}{\lemma{\skt{मूत्रासवैर्}}\Afootnote{K H; \skt{मूत्रासवै}\ A}}}{\lemma{\skt{गुडकृतैश्च}}\Afootnote{K H; \skt{र्गुडकृतैश्च}\ A}}\skt{ }\edtext{\edtext{\skt{तथाप्यरिष्टैः}}{\lemma{\skt{तथाप्य्}}\Afootnote{K H; \skt{तथा}\ A}}}{\lemma{\skt{अरिष्टैः}}\Afootnote{K H; \skt{त्वरिष्टैः}\ A}}\skt{ क्षारासवैश्च }\edtext{\skt{मधु\-मा\-ध\-व\-तु\-ल्य\-गन्धैः}}{\Afootnote{A K; \skt{मधु\-मा\-ध\-व\-तु\-गंधैः ।}\ H}}\skt{ । }\pend
\pstart\versepar \skt{}\edtext{\skt{स्यादेष}}{\Afootnote{K; \skt{स्यादेष}\ A; \skt{स्याद्देश}\ H}}\skt{ एव }\edtext{\skt{कफवातहरो}}{\Afootnote{K; \skt{कफवातहते}\ A; \skt{कफवातहरे}\ H}}\skt{ }\edtext{\skt{विधिश्च}}{\Afootnote{A K; \skt{द्विधिश्च}\ H}}\skt{ । }\edtext{\skt{शातिं}}{\Afootnote{K; \skt{शान्तिं}\ A H}}\skt{ गते हुतभुजि प्रशमाय तस्य ।}\pend

\pstart\versepar \ledsidenote{\passref{6.57.16}}\edtext{\skt{इच्छा\-वि\-ना\-श\-भ\-य\-जेषु}}{\Afootnote{K H; \skt{इच्छा\-भि\-घा\-त\-भ\-य\-शो\-क\-हत्}\ A}}\skt{ }\edtext{\skt{च}}{\Afootnote{K H; \skt{ए}\ A}}\skt{ }\edtext{\skt{बाधकेषु}}{\Afootnote{K H; \skt{ऽन्तरग्नौ}\ A}}\skt{ । }\edtext{\skt{भावाद्भवाय}}{\Afootnote{K H; \skt{भावान्भवाय}\ A}}\skt{ वितरेत्खलु शल्यतापान् । }\pend
\pstart\versepar \skt{अर्थेषु }\edtext{\skt{चातिपतितेषु}}{\Afootnote{K H; \skt{चाप्यपचितेषु}\ A}}\skt{ पुनर्भवाय }\edtext{\edtext{\skt{पौरा\-णि\-क\-श्रु\-ति\-मु\-खै\-र\-व\-मा\-न\-येत}}{\lemma{\skt{पौरा\-णि\-क\-श्रु\-ति\-मु\-खैर्}}\Afootnote{K; \skt{पौराणिकैः श्रुतिशतै}\ A; \skt{पौरा\-णि\-क\-स्रु\-ति\-मु\-खैर्}\ H}}}{\lemma{\skt{अवमानयेत}}\Afootnote{K H; \skt{रनु\-मा\-न\-ये\-त्तम् ॥}\ A}}\skt{ }\edtext{\skt{इति}}{\Afootnote{K H; \emph{om.}\ A}}\skt{ ॥ }\edtext{\skt{कायचि}}{\Afootnote{K; \emph{om.}\ A H}}\skt{ ॥ २० ॥ ० ॥}\edtext{}{\lemma{\skt{शल्यतापान् । अर्थेषु}}\Afootnote{K; \skt{शक्यरूपान् । अर्थेषु}\ A; \skt{शल्यतायान् । अर्थेषु}\ H}}\pend

\pstart\versepar \ledsidenote{\passref{6.57.17}}\edtext{\skt{च्छर्दि}}{\Afootnote{K; \emph{om.}\ A; \skt{छर्द्दिर्}\ H}}\skt{ }\edtext{\skt{हिक्का}}{\Afootnote{K H; \emph{om.}\ A}}\skt{ }\edtext{\skt{श्वासकासः}}{\Afootnote{K; \emph{om.}\ A; \skt{श्वासकासौ}\ H}}\skt{ }\edtext{\edtext{\skt{स्वर\-घा\-त\-म\-था\-परं}}{\lemma{\skt{स्वरघातम्}}\Afootnote{K H; \emph{om.}\ A}}}{\lemma{\skt{अथापरं}}\Afootnote{K H; \emph{om.}\ A}}\skt{ । }\pend
\pstart\versepar \skt{}\edtext{\skt{प्रतिश्यायः}}{\Afootnote{K H; \emph{om.}\ A}}\skt{ }\edtext{\edtext{\skt{क्रिमिश्चैव}}{\lemma{\skt{क्रिमिश्}}\Afootnote{K H; \emph{om.}\ A}}}{\lemma{\skt{चैव}}\Afootnote{K H; \skt{दैन्यं गते मनसि बोधनमत्र शस्तं यद्यत्प्रियं तदुप}\ A}}\skt{ }\edtext{\skt{सोदावर्त्ता}}{\Afootnote{K H; \skt{सेव्यम}\ A}}\skt{ }\edtext{\skt{विसूचिका}}{\Afootnote{K; \emph{om.}\ A; \skt{विशूचिका ॥}\ H}}\skt{ । }\pend
\pstart\versepar \skt{}\edtext{\skt{अरोचकेन च}}{\Afootnote{K H; \skt{रोचके तु ॥}\ A}}\skt{ }\edtext{\skt{तथा}}{\Afootnote{K H; \emph{om.}\ A}}\skt{ }\edtext{\skt{पूर्यते}}{\Afootnote{K H; \emph{om.}\ A}}\skt{ }\edtext{\skt{द्वितीयो}}{\Afootnote{K\textsuperscript{pc} H; \emph{om.}\ A; \skt{द्वितीः ॥ {\symfont ❈} ॥}\ K\textsuperscript{ac}}}\skt{ }\edtext{\skt{दशः ॥ {\symfont ❈}}}{\Afootnote{K\textsuperscript{pc}; \emph{om.}\ A K\textsuperscript{ac}; \skt{दश ॥ {\symfont ❈} ॥}\ H}}\skt{ ॥}\pend

\endnumbering
